\documentclass[11pt,a4paper]{article}
\usepackage[T1]{fontenc}
\usepackage{lmodern}
\usepackage[margin=24mm,headheight=15pt]{geometry}
\usepackage{amsmath,amssymb,amsthm,mathtools,mathrsfs}
\usepackage{booktabs,array,longtable,tabularx,enumitem,graphicx}
\usepackage{microtype,xcolor,listings,upquote,fancyhdr,needspace}
\usepackage[colorlinks=true,linkcolor=blue!40!black,citecolor=blue!40!black,urlcolor=blue!40!black]{hyperref}
\hypersetup{pdftitle={Excellent Frobenius Lifts and Supercongruences: Complete Edition},pdfauthor={Research manuscript prepared for Keio},pdfsubject={All-dimensional Beukers--Vlasenko congruences, modular extensions, alternative arguments, sharpness, and complete verification source}}
\newtheorem{theorem}{Theorem}[section]
\newtheorem{lemma}[theorem]{Lemma}
\newtheorem{proposition}[theorem]{Proposition}
\newtheorem{corollary}[theorem]{Corollary}
\newtheorem{conjecture}[theorem]{Conjecture}
\newtheorem*{bvconj}{Conjecture BV (Beukers--Vlasenko {\cite[Conj.~7.5]{BV3}})}
\theoremstyle{definition}\newtheorem{definition}[theorem]{Definition}
\theoremstyle{remark}\newtheorem{remark}[theorem]{Remark}
\newcommand{\Z}{\mathbb Z}
\newcommand{\Q}{\mathbb Q}
\newcommand{\C}{\mathbb C}
\newcommand{\Zp}{\mathbb Z_p}
\newcommand{\Fp}{\mathbb F_p}
\newcommand{\Hb}{\mathbb H}
\newcommand{\Hh}{\mathbb H}
\DeclareMathOperator{\CT}{CT}
\DeclareMathOperator{\Tr}{Tr}
\DeclareMathOperator{\ord}{ord}
\DeclareMathOperator{\SL}{SL}
\DeclareMathOperator{\GL}{GL}
\DeclareMathOperator{\adj}{adj}
\newcommand{\vp}{v_p}
\newcommand{\bmat}[4]{\left(\begin{smallmatrix}#1&#2\\#3&#4\end{smallmatrix}\right)}
\numberwithin{equation}{section}
\allowdisplaybreaks[2]
\setlength{\emergencystretch}{3em}
\setlength{\parskip}{2pt}
\setcounter{tocdepth}{1}
\setlist{topsep=4pt,itemsep=2pt}
\lstset{basicstyle=\ttfamily\scriptsize,breaklines=true,columns=fullflexible,
 frame=single,rulecolor=\color{black!20},keepspaces=true,showstringspaces=false,
 upquote=true,aboveskip=8pt,belowskip=8pt,breakatwhitespace=false,
 literate={²}{{\textsuperscript{2}}}1}
\pagestyle{fancy}\fancyhf{}
\fancyhead[L]{\small Excellent Frobenius Lifts and Supercongruences}
\fancyhead[R]{\small Complete edition}
\fancyfoot[C]{\thepage}
\renewcommand{\headrulewidth}{0.3pt}
\newcommand{\partstart}[3]{\clearpage\part{#1}\label{#2}\noindent\textit{#3}\par\medskip}
\newcommand{\ZZ}{\mathbb Z}\newcommand{\QQ}{\mathbb Q}\newcommand{\tht}{\theta}
\hypersetup{bookmarksopen=true,bookmarksopenlevel=0}
\begin{document}
\hypersetup{pageanchor=false}
\begin{titlepage}
\centering
\vspace*{12mm}
{\large NUMBER THEORY\par}
\vspace{11mm}
{\LARGE\bfseries Excellent Frobenius Lifts\par}
\vspace{2mm}
{\LARGE\bfseries and Supercongruences\par}
\vspace{7mm}
{\large The three Beukers--Vlasenko families\par}
{\large in all dimensions at every odd prime\par}
\vspace{9mm}
{\large Complete collected edition\par}
\vspace{2mm}
{Prepared for Keio\par}
{9 November 2026\par}
\vspace{8mm}
\begin{minipage}{0.89\textwidth}\small
\textbf{Abstract.}
We assemble the complete arguments for excellent-Frobenius-lift
truncation supercongruences for the hypercubic, simplicial, and
hyperoctahedral families in every dimension and at every odd prime.
For every positive cutoff $K$ divisible by $p$, the normalized
truncation quotient descends modulo $p^{2v_p(K)}$, coefficientwise
in the full formal power-series ring. The prescribed multipliers of
Beukers--Vlasenko, Conjecture~7.5, are included, together with every
positive normalized multiplier.

The proof combines a directly constructed compatible orbit module,
positive-vertex interpolation, balanced factorial blocks, and
support-weighted symmetric Bessel reduction. At the remaining prime
$3$, an integral two-variable cutoff kernel gives a quadratic bound
in the ordinary dimension-polynomial ring. A universal weak descent
loses one factor of $3$, and scalar residue dimensions recover it
at every integer dimension. A separate saturation argument transfers
the prescribed congruence to the standard analytic coefficient ring
at the primes covered by the published excellent-lift theorem.

The volume retains the earlier independent proofs, modular trace
extensions, the elliptic-zero families $27t$, $256t$, and $108t$,
the all-level Zagier D arguments, negative-vertex and differential
methods, and sharpness results. Exact finite certificates through
level five and their complete verification sources accompany the
algebraic proofs. The all-level results do not rely on extrapolation
from finite calculations. The scope of numerical specialization and
the remaining extension problems is stated explicitly.
\end{minipage}
\vfill
{\small Research manuscript; internal checking completed, external review pending.\par}
{\small MSC 2020: 11F33, 11B65, 14F30, 33C20.\par}
\end{titlepage}
\hypersetup{pageanchor=true}
\pagenumbering{roman}
\tableofcontents
\clearpage
\pagenumbering{arabic}
\section{Global statement, proof dependencies, and exact scope}
\label{sec:unified-guide}\label{G:sec:guide}

This complete edition combines the original four-part manuscript with
three subsequent supplements on small primes and hyperoctahedral descent.
The resulting formal-series theorem covers the three standard
Beukers--Vlasenko families in every dimension at every odd prime.
Earlier proofs with more restrictive prime ranges are retained as
independent arguments; their local hypotheses are not restrictions on
the global theorem below. The modular extensions and their separate
prime conditions are retained as well.

\subsection{Families, coordinates, and the global theorem}
For an integer $n\ge1$, let $g$ be one of the Laurent polynomials
\[
\begin{array}{c|c|c}
\text{family}&g(\mathbf x)&h\\\hline
\text{hypercubic}&\displaystyle\prod_{i=1}^n(x_i+x_i^{-1})&2\\[3pt]
\text{simplicial}&\displaystyle x_1+\cdots+x_n+(x_1\cdots x_n)^{-1}&n+1\\[3pt]
\text{hyperoctahedral}&\displaystyle\sum_{i=1}^n(x_i+x_i^{-1})&2.
\end{array}
\]
Write
\[
 \mathcal F(z)=\CT\frac1{1-zg}
       =\sum_{j\ge0}\CT(g^j)z^j=F(u),\qquad u=z^h,
 \qquad F(u)=\sum_{j\ge0}a_ju^j,
\]
and define the ordinary truncations
\[
 \mathcal F_L(z)=\sum_{0\le j<L}\CT(g^j)z^j,
 \qquad F_K(u)=\sum_{0\le j<K}a_ju^j.
\]
The formal excellent lift is
\[
 \tau(z)=q_0^{-1}\bigl(q_0(z)^p\bigr),\qquad
 q_0(z)=z\exp\bigl(G_0(z)/\mathcal F(z)\bigr),
\]
where $G_0$ is the vertex logarithmic companion of
\cite[Proposition~9.1]{BV3}. The orbit-module construction in
Proposition~\ref{SP:prop:orbit} verifies its integrality and the excellent
Cartier relation for every odd prime. Its normalized form is
$\sigma(u)=\tau(z)^h$. This is a well-defined member of
$u^p(1+p\mathbb Z_p[[u]])$, including the simplicial case $p\mid h$.
Every period or truncation series used as a formal denominator has
constant term one.

\begin{theorem}[Global formal truncation congruence]
\label{G:thm:global}
For each of the three families, every dimension $n\ge1$, every odd
prime $p$, and every positive integer $K$ divisible by $p$,
\begin{equation}\label{G:eq:strong}
 \frac{F_K(u)}{F(u)}
 \equiv
 \frac{F_{K/p}(\sigma(u))}{F(\sigma(u))}
 \pmod{p^{2v_p(K)}\mathbb Z_p[[u]]}.
\end{equation}
Consequently, for every pair of positive integers $s,m$,
\begin{align}
 E_{p,s,m}(u)
 &:={F(u)F_{mp^{s-1}}(\sigma(u))
      -F(\sigma(u))F_{mp^s}(u)}\notag\\
 &\in p^{2s+2v_p(m)}\mathbb Z_p[[u]],
 \label{G:eq:defect}
\end{align}
and in particular
\begin{equation}\label{G:eq:ratio}
 \frac{F(u)}{F(\sigma(u))}
 \equiv\frac{F_{mp^s}(u)}{F_{mp^{s-1}}(\sigma(u))}
       \pmod{p^{2s}\mathbb Z_p[[u]]}.
\end{equation}
All coefficients of these infinite formal series are included.
\end{theorem}

\begin{proof}
The hypercubic and simplicial cases, including $p\mid n+1$ in the
latter, are Theorem~\ref{SP:thm:scope}(1). For the hyperoctahedral
family in dimension $n\ge3$, Theorem~\ref{OP:thm:all-level} gives
\eqref{G:eq:strong} at every prime $p\ge5$, and
Theorem~\ref{AP:thm:main} gives it at $p=3$. The dimensions one and
two reduce to the corresponding scalar hypercubic positive-vertex
matrices, as recorded in those parts. Thus these cases exhaust all
dimensions and all odd primes. Substitute $K=mp^s$ into
\eqref{G:eq:strong} and multiply by the formal units $F(u)$ and
$F(\sigma(u))$ to obtain \eqref{G:eq:defect}. Division by the
appropriate formal units gives \eqref{G:eq:ratio}.
\end{proof}

The stronger modulus in \eqref{G:eq:defect} comes from the single
cutoff $K=mp^s$; no assumption $p\nmid m$ is made.
The relation
\begin{equation}\label{G:eq:normalization}
 \mathcal F_{hK}(z)=F_K(z^h)
\end{equation}
shows exactly how this statement meets the original conjecture.

\begin{corollary}[The prescribed Beukers--Vlasenko truncations]
\label{G:cor:BV}
For every dimension and every odd prime, the formal congruence
\[
 \frac{\mathcal F(z)}{\mathcal F(\tau(z))}
 \equiv
 \frac{\mathcal F_{hp^s}(z)}
      {\mathcal F_{hp^{s-1}}(\tau(z))}
 \pmod{p^{2s}\mathbb Z_p[[z]]}
 \qquad(s\ge1)
\]
holds with $h=2$ for the hypercubic and hyperoctahedral families and
$h=n+1$ for the simplicial family. These are the multipliers specified
in \cite[Conjecture~7.5]{BV3}.
\end{corollary}

\begin{proof}
Use $m=1$ in Theorem~\ref{G:thm:global}, substitute $u=z^h$, and
apply \eqref{G:eq:normalization} together with $\sigma(u)=\tau(z)^h$.
\end{proof}

In particular, the formerly missing hyperoctahedral sector
$p=3$, $n\ge4$, $3\nmid n$, $s\ge5$ is covered. The theorem is
stronger than its $m=1$ specialization, but the expression ``all
multipliers'' throughout the standard-family arguments means
\emph{all normalized multipliers}: the raw lengths are $hmp^s$.
No assertion for every arbitrary raw multiplier is inferred.

\subsection{Proof architecture and retained extensions}
\begin{center}\small
\renewcommand{\arraystretch}{1.18}
\begin{tabularx}{\textwidth}{@{}>{\raggedright\arraybackslash}p{.25\textwidth}>{\raggedright\arraybackslash}p{.19\textwidth}>{\raggedright\arraybackslash}X@{}}
\toprule
Family or sector&Prime range&Conclusion and proof\\\midrule
Hypercubic, all $n\ge1$&Every odd $p$&All $s,m$ and all coefficients;
Theorem~\ref{SP:thm:scope}(1).\\
Simplicial, all $n\ge1$&Every odd $p$, including $p\mid n+1$&All $s,m$ and all coefficients;
Theorem~\ref{SP:thm:scope}(1).\\
Hyperoctahedral, all $n\ge1$&Every $p\ge5$&All $s,m$ and all coefficients;
Theorem~\ref{OP:thm:all-level} and the scalar reductions.\\
Hyperoctahedral, all $n\ge1$&$p=3$&All $s,m$ and all coefficients;
Theorem~\ref{AP:thm:main} and the scalar reductions.\\
Thirteen cusp-type modular families&$p\nmid N$, stated Hecke condition&Exponent $(w+1)s$, $w\in\{1,2\}$;
Theorem~\ref{B:thm:main}.\\
Elliptic-zero $27t,256t,108t$&$p\ge5$&Exponents $2s,3s,3s$, respectively;
Theorems~\ref{B:thm:ell} and~\ref{E:thm:ell-all-s}.\\
Zagier D&Every $p\ne5$&Exponent $2s$ for all $s,m$, including $p=2,3$;
Theorems~\ref{B:thm:D} and~\ref{T:thm:zagier-nonsplit-complete}.\\
\bottomrule
\end{tabularx}
\end{center}

The original positive-vertex part gives the large-prime proof using
the published excellent Cartier theorem. The small-prime supplement
constructs a compatible orbit module directly, so no division by a
symmetry-group order is used. The next hyperoctahedral supplement
controls factorial denominators by the number of differentiated
variables and proves descent for every $p\ge5$. The final supplement
establishes an integral two-variable Bessel kernel, a universal
coefficientwise descent with one factor of $3$ lost, and a dimension
residue argument that recovers this factor at every integer dimension.
These are algebraic arguments for all indices. The finite certificates
at specified levels are independent checks and alternative proofs of
those individual congruences, not premises of an extrapolation to
arbitrary levels.

The original modular, elliptic, two-cusp, negative-vertex, differential,
and Domb arguments remain included. A universal replacement of $2s$
by $3s$ is false: the exact sharpness statement for three-dimensional
hyperoctahedra at $p=5$ is Theorem~\ref{E:thm:sh-octa3}. The stronger
$3s$ conclusions in the table have their own hypotheses. The broader
elliptic-zero Conjecture~\ref{B:conj:ell}, beyond the proved list,
remains an extension problem.

\subsection{Transfer to the original analytic coefficient ring}
\label{G:sec:analytic-transfer}
The all-odd-prime theorem is formulated in $\mathbb Z_p[[z]]$.
For primes satisfying the hypotheses of \cite[Theorem~7.3]{BV3},
one can also state the prescribed-truncation congruence in the
standard analytic coefficient ring. The following elementary
argument makes this transfer explicit; it does not assert that the
period $\mathcal F$ itself belongs to that ring.

\begin{lemma}[Saturation of the standard coefficient ring]
\label{G:lem:saturation}
Let $P\in\mathbb Z_p[z]$ with $P(0)\in\mathbb Z_p^\times$, and let
$R_P$ be the $p$-adic completion of $\mathbb Z_p[z,1/P(z)]$.
Expansion at zero embeds $R_P$ into $\mathbb Z_p[[z]]$, and, for
every integer $k\ge1$,
\[
 R_P\cap p^k\mathbb Z_p[[z]]=p^kR_P.
\]
\end{lemma}

\begin{proof}
For every $k$, the polynomial $P$ is a unit in
$(\mathbb Z/p^k\mathbb Z)[[z]]$. The expansion map
\[
 (\mathbb Z/p^k\mathbb Z)[z,1/P]
       \longrightarrow(\mathbb Z/p^k\mathbb Z)[[z]]
\]
is injective: if a fraction has zero expansion, multiplication by
its denominator shows that its polynomial numerator is zero modulo
$p^k$. The quotient of $R_P$ by $p^kR_P$ is this localized
polynomial ring, since $\mathbb Z_p[z,1/P]$ is Noetherian and its
$p$-adic completion has the same quotients modulo $p^k$.
The compatible expansion maps therefore embed the inverse limit
$R_P$ into $\mathbb Z_p[[z]]$. Their injectivity modulo $p^k$
gives the stated intersection equality.
\end{proof}

\begin{lemma}[The prescribed truncations are analytic units]
\label{G:lem:truncation-units}
For each of the three standard families, put
$c_j=\CT(g^j)$ and $A(z)=\mathcal F_p(z)$.
For every prime $p$ and every $r\ge0$,
\begin{equation}\label{G:eq:Lucas-truncation}
 \mathcal F_{hp^r}(z)
 \equiv\prod_{i=0}^{r-1} A(z^{p^i})\pmod p,
\end{equation}
where an empty product is one. Consequently these truncations and
their compositions with any Frobenius lift are units in a
$p$-adically complete coefficient ring in which $A$ is a unit.
\end{lemma}

\begin{proof}
The Newton polytope $\Delta$ is reflexive, and zero is its only
interior lattice point. If $0\le a<p$ and $b\ge0$, the congruence
$g^{a+pb}\equiv g^a g(\mathbf x^p)^b\pmod p$ gives
\[
 c_{a+pb}\equiv c_a c_b\pmod p.
\]
Indeed an exponent of $g^a$ that can pair with an exponent of
$g(\mathbf x^p)^b$ is divisible by $p$. Dividing it by $p$ places
it in $(a/p)\Delta\subset\Delta^\circ$, so it must be zero.
Iterating this digit relation and summing the indices less than
$hp^r$ yields
\[
 \mathcal F_{hp^r}(z)
 \equiv \mathcal F_h(z^{p^r})
          \prod_{i=0}^{r-1} A(z^{p^i})\pmod p.
\]
In these families a positive constant-term degree is divisible by
$h$, hence $\mathcal F_h=1$. This proves
\eqref{G:eq:Lucas-truncation}. Modulo $p$ each factor is
$A(z)^{p^i}$; it is therefore a unit whenever $A$ is a unit.
Units modulo $p$ lift to units in a $p$-adically complete ring.
A ring endomorphism sends a unit to a unit, which gives the
composition assertion.
\end{proof}

\begin{corollary}[Congruence in the BV3 analytic ring]
\label{G:cor:analytic}
Assume that $p$ and the symmetry data satisfy the hypotheses of
\cite[Theorem~7.3]{BV3}. Compute the Hasse--Witt polynomials with
the initial lift $z\mapsto z^p$, and set
\[
 R_0=\widehat{\mathbb Z_p
       [z,1/hw^{(1)}(z),1/hw^{(2)}(z)]},\qquad
 \lambda(z)=\frac{\mathcal F(z)}{\mathcal F(\tau(z))}.
\]
For each $s\ge1$, the prescribed-truncation congruence holds in
$R_0$:
\[
 \lambda(z)
 -\frac{\mathcal F_{hp^s}(z)}
        {\mathcal F_{hp^{s-1}}(\tau(z))}
       \in p^{2s}R_0.
\]
\end{corollary}

\begin{proof}
The cited theorem and \cite[Proposition~7.7 and Theorem~7.15]{BV3}
give $\tau,\lambda\in R_0$ and a Frobenius endomorphism of $R_0$.
The polynomials $hw^{(1)},hw^{(2)}$ have unit constant terms.
Also $A=\mathcal F_p\equiv hw^{(1)}\pmod p$, so $A$ is a unit
of $R_0$. Lemma~\ref{G:lem:truncation-units} shows that
$\mathcal F_{hp^{s-1}}(\tau(z))$ is a unit of $R_0$.
Corollary~\ref{G:cor:BV}, after multiplication by this denominator,
says that the element
\[
 \lambda(z)\mathcal F_{hp^{s-1}}(\tau(z))
            -\mathcal F_{hp^s}(z)\in R_0
\]
belongs to $p^{2s}\mathbb Z_p[[z]]$. Apply
Lemma~\ref{G:lem:saturation} with
$P=hw^{(1)}hw^{(2)}$, and then divide by the displayed unit.
\end{proof}

Thus the original analytic-ring interpretation in the published
admissible-prime setting follows from the formal theorem. This
does not by itself establish an analogous analytic realization at
every additional small prime admitted by the formal orbit-module
construction. A numerical specialization also requires an evaluation
map on the relevant ring. If such a continuous map is given, the
ring congruence specializes; identifying a fixed point of the lift
and identifying the resulting scalar with a geometric unit root
are separate assertions.

\subsection{Provenance, verification, and limits}
This edition retains the mathematical arguments of the four
original manuscripts and incorporates the three subsequent
supplements. The global theorem is their explicit combination.
The analytic-ring transfer above is recorded here as a short
additional argument to clarify the relationship with the original
BV3 coefficient ring. Published Cartier, connection, and coordinate
results are cited at their points of use; their complete published
proofs are not represented as new results of this manuscript.

Three kinds of evidence are distinguished: algebraic proofs for
all indices; bounded exact-arithmetic checks, with any historical
floating-point checks expressly identified; and finite polynomial
certificates whose proved degree and denominator properties imply
an entire formal-series congruence. Finite samples alone are never
used to infer the global theorem. The manuscript and accompanying
calculations have been checked internally; external peer review
and proof-assistant formalization have not been completed.

The scope of the global theorem excludes the prime $2$, arbitrary
raw truncation multipliers, arbitrary Laurent-polynomial families,
and unrestricted numerical specialization. These exclusions do not
remove any of the prescribed formal truncations of the three-family
odd-prime statement. Separately, the modular extension at Zagier D
does include $p=2$, under its own theorem. No claim is made that the
open elliptic-zero generalization or every related special-value
problem is resolved.

The single \texttt{.tex} source contains the mathematical text,
references, and embedded program listings of this edition. It
requires no companion chapter, bibliography, image, or data file
to compile. Standard installed LaTeX packages are required.
Computational certificate data used for independent checks are
not hypotheses of the algebraic all-level proofs.

\subsection{Reading the complete edition}
Parts~\ref{part:positive}--\ref{part:octa-three} contain the full
three-family argument, including the original restricted-prime proof
and the subsequent extensions. Parts~\ref{part:modular}--\ref{part:twocusp}
retain the complementary modular and geometric arguments. Repeated
formal-input proofs are kept with their local notation so that no
proof is replaced by an editorial synopsis. The source-to-section
map in Appendix~\ref{app:coverage} identifies every retained source
section. Earlier code comments and recorded outputs are historical
records of those particular computations; current theorem scope is
always the scope stated in this introduction.

The algebraic global proof is independent of the optional certificate
data. All certificate algorithms needed for reproduction are printed
in the appendices, including the new fifth-level generator and an
independent evaluator. Large coefficient arrays are reproducible
outputs of those programs. They are not external inputs needed to
read, compile, or prove the global theorem.

Compile this one source with \texttt{latexmk -pdf}, or run
\texttt{pdflatex} until references stabilize. No separate bibliography,
chapter, figure, or data file is needed for compilation.

\partstart{Positive vertex interpolation: the uniform three-family theorem}{part:positive}{All dimensions, all levels, all multipliers, and all coefficients in the standard full-symmetry prime ranges.}
% BEGIN RETAINED SOURCE P
\section{The target and the resulting scope}\label{P:source-main-1}
\label{P:sec:main-scope}

This part proves the uniform ordinary-truncation theorem for the three standard families, including the three-dimensional hyperoctahedral family and every higher dimension. The modular and Zagier D arguments in the other parts are independent extensions. The common input
is the excellent Cartier theorem of Beukers and Vlasenko
\cite[Theorem~7.3]{BV3}, together with its coefficient interpretation
immediately after \cite[Corollary~3.4]{BV3}.

Let $g$ be one of the following Laurent polynomials in $n$ variables,
and put $f=1-zg$:
\[
\begin{array}{c|c|c}
\text{family}&g({\bf x})&h\\ \hline
\text{hypercubic}&\displaystyle\prod_{i=1}^n(x_i+x_i^{-1})&2\\[3pt]
\text{simplicial}&\displaystyle x_1+\cdots+x_n+(x_1\cdots x_n)^{-1}&n+1\\[3pt]
\text{hyperoctahedral}&\displaystyle\sum_{i=1}^n(x_i+x_i^{-1})&2
\end{array}
\]
Normalize $u=z^h$ and write
\[
 \mathcal F(z)=\CT\frac1{1-zg}=F(u),\qquad
 F(u)=\sum_{j\ge0}a_ju^j,\qquad
 F_K(u)=\sum_{0\le j<K}a_ju^j.
\]
Let $z\mapsto\tau(z)$ be the actual excellent lift of
\cite[Theorem~7.3]{BV3}. Its normalized form is
$\sigma(u)=\tau(z)^h$. The canonical coordinate formulas in
\cite[Definition~7.14 and Proposition~9.1]{BV3} show that this
is a series in $u$, with
$\sigma/u^p\in1+p\mathbb Z_p[[u]]$.
All period and truncation series appearing in denominators have
constant term one.

Define the cleared defect
\begin{equation}\label{P:eq:main-defect}
 E_{p,s,m}(u)=F(u)F_{mp^{s-1}}(\sigma(u))
                   -F(\sigma(u))F_{mp^s}(u).
\end{equation}
The assertion $E_{p,s,m}\in p^{2s}\mathbb Z_p[[u]]$ means that
\emph{every} coefficient is divisible by $p^{2s}$, with no
bound on its index.

\begin{theorem}[The three standard families]\label{P:thm:main-all}
Let $n\ge1$, and use the standard full symmetry groups of
\cite[Section~7]{BV3}. Assume that $p$ is an odd prime and
\[
 p>n\quad\text{for the hypercubic and hyperoctahedral families},
 \qquad
 p>n+1\quad\text{for the simplicial family}.
\]
Then, for every positive integer $K$ divisible by $p$,
\begin{equation}\label{P:eq:main-strong}
 \frac{F_K(u)}{F(u)}
 \equiv\frac{F_{K/p}(\sigma(u))}{F(\sigma(u))}
           \pmod{p^{2v_p(K)}\mathbb Z_p[[u]]}.
\end{equation}
In particular, for all $m,s\ge1$,
\[
 E_{p,s,m}\in p^{2s}\mathbb Z_p[[u]],\qquad
 \frac{F(u)}{F(\sigma(u))}
 \equiv\frac{F_{mp^s}(u)}{F_{mp^{s-1}}(\sigma(u))}
                       \pmod{p^{2s}}.
\]
In the original variable, the result covers all truncation lengths
$hmp^s$. The prescribed lengths $2p^s$ for cubes and octahedra
and $(n+1)p^s$ for simplices in Conjecture~7.5 are the case $m=1$.
\end{theorem}

The cubic and simplicial statements are proved in
Theorem~\ref{P:thm:positive-cube-simplex}; the octahedral statement
in dimension $n\ge3$ is Theorem~\ref{P:thm:octa-all-levels}.
The low-dimensional octahedral reduction is recorded in that section.
The prime conditions of this original proof are retained.
Theorems~\ref{SP:thm:scope}, \ref{OP:thm:all-level}, and
\ref{AP:thm:main} extend the conclusion to every odd prime by
additional arguments, combined in Theorem~\ref{G:thm:global}.

\begin{center}\small
\renewcommand{\arraystretch}{1.25}
\begin{tabular}{@{}p{.37\textwidth}p{.18\textwidth}p{.35\textwidth}@{}}
\toprule
Family and dimension&Prime range&Conclusion\\ \midrule
Hypercubic, $n\ge4$&$p>n$&All $s,m$, every coefficient\\
Simplicial, $n\ge4$&$p>n+1$&All $s,m$, every coefficient\\
Hyperoctahedral, $n=3$&$p\ge5$&All $s,m$, every coefficient\\
Hyperoctahedral, $n\ge4$&$p>n$&All $s,m$, every coefficient\\
\bottomrule
\end{tabular}
\end{center}

\subsection{How the excellent Cartier relation reaches ordinary truncations}

For each family, a positive vertex coefficient of $1/f$ has the
form
\[
 \Omega_N(u)=\sum_{j\ge N}B(j,N)u^j,
 \quad B(j,N)\in\mathbb Z,\quad B(j,0)=a_j,\quad B(j,j)=1.
\]
Set $C_N=\Omega_N/F$. The excellent Cartier relation gives
\begin{equation}\label{P:eq:main-Cartier-roadmap}
 C_{pN}-C_N\circ\sigma\in
                p^{2+2v_p(N)}\mathbb Z_p[[u]].
\end{equation}
The unit triangular matrix $B$ defines integral coefficients
$c_{K,N}$ by
\[
 a_j\mathbf1_{j\ge K}=\sum_{N=K}^j c_{K,N}B(j,N),
 \qquad
 1-\frac{F_K}{F}=\sum_{N\ge K}c_{K,N}C_N.
\]
The essential new estimates are
\[
 v_p(c_{K,N})\ge\max\{0,2v_p(K)-2v_p(N)\},
\]
and, when $p\mid K$,
\[
 c_{K,N}\equiv\mathbf1_{p\mid N}\,c_{K/p,N/p}
                              \pmod{p^{2v_p(K)}}.
\]
These two estimates let the ordinary tail descend term by term
through \eqref{P:eq:main-Cartier-roadmap} at the required precision.
All sums converge $u$-adically, so the argument covers the entire
series. The next two sections prove the estimates without assuming
the desired truncation congruence.

\section{Positive vertex interpolation: complete hypercubic and simplicial congruences}\label{P:source-main-2}
\label{P:sec:positive-vertex-interpolation}

All congruences of formal series in this section are coefficientwise.
Let \(n\geq1\), and consider one of the following two families:
\[
\begin{array}{c|c|c|c}
 &g({\bf x})&h&a_j\\ \hline
\text{hypercubic}
 &\prod_{i=1}^n(x_i+x_i^{-1})&2&\binom{2j}{j}^{n}\\[2pt]
\text{simplicial}
 &x_1+\cdots+x_n+(x_1\cdots x_n)^{-1}
 &n+1&((n+1)j)!/(j!)^{n+1}
\end{array}
\]
Put \(f=1-zg\), \(u=z^h\),
\[
 \mathcal F(z)=\operatorname{CT}(1/f)=F(u),\qquad
 F(u)=\sum_{j\geq0}a_ju^j,\qquad
 F_K(u)=\sum_{0\leq j<K}a_ju^j.
\]
Throughout, \(p\) is an odd prime satisfying the symmetry and
support-lattice hypotheses of \cite[Theorem 7.3]{BV3}.
For the full symmetry groups specified in that paper, these conditions
are \(p>n\) in the cubic case and \(p>n+1\) in the simplicial case,
with the additional odd-prime requirement in dimension one.
In particular, \(p\nmid h\).

Let \(z\mapsto\tau\) be the actual excellent Frobenius lift of that
theorem.  The normalized lift \(u\mapsto\sigma(u)=\tau^h\) is the
one characterized by \(q(\sigma(u))=q(u)^p\), where
\(q(z^h)=q_0(z)^h\) and \(q_0\) is the canonical coordinate
of \cite[Definition 7.14 and Proposition 9.1]{BV3}.
No coordinate change different from this lift is introduced.

\begin{theorem}[All levels for hypercubes and simplices]
\label{P:thm:positive-cube-simplex}
For either of the above families, and every positive integer \(K\)
divisible by \(p\),
\begin{equation}\label{P:eq:positive-main}
 \frac{F_K(u)}{F(u)}
 \equiv
 \frac{F_{K/p}(\sigma(u))}{F(\sigma(u))}
 \pmod{p^{2v_p(K)}\mathbb Z_p[[u]]}.
\end{equation}
Hence, for all integers \(m,s\geq1\),
\[
 F(u)F_{mp^{s-1}}(\sigma(u))
 -F(\sigma(u))F_{mp^s}(u)
 \in p^{2s}\mathbb Z_p[[u]].
\]
Equivalently,
\[
 \frac{F(u)}{F(\sigma(u))}
 \equiv
 \frac{F_{mp^s}(u)}{F_{mp^{s-1}}(\sigma(u))}
 \pmod{p^{2s}}.
\]
In the original variable, this proves Conjecture 7.5 for every
hypercubic and every simplicial dimension, at all primes in its
stated symmetry range.  In fact all original truncation lengths
\(hmp^s\), \(m,s\geq1\), are covered, not only \(hp^s\).
\end{theorem}

The proof uses positive vertex coefficient series whose leading
terms form an integral triangular basis.  It does not use the
negative-power canonical basis or any duality between opposite
expansions.

\subsection{A binomial and multinomial block congruence}

Set \(r=n+1\) in the simplicial case and define, for
\(0\leq N\leq j\),
\[
 B(j,N)=
 \begin{cases}
 \displaystyle\binom{2j}{j-N}^{n},&\text{hypercubic},\\[7pt]
 \displaystyle\frac{(rj)!}
 {(j+(r-1)N)!(j-N)!^{r-1}},&\text{simplicial}.
 \end{cases}
\]
Put \(B(j,N)=0\) for \(N>j\).
In both cases these are nonnegative integers with
\(B(j,0)=a_j\) and \(B(j,j)=1\).

\begin{lemma}[Balanced factorial block congruence]
\label{P:lem:positive-block}
Let \(p\) be odd and \(l=\lfloor j/p\rfloor\).
Then \(d_j=a_j/a_l\in\mathbb Z_p\), and, for every
\(1\leq N\leq l\),
\begin{equation}\label{P:eq:positive-block}
 B(j,pN)-d_j B(l,N)
 \in p^{2+2v_p(N)}\mathbb Z_p.
\end{equation}
This lemma itself holds for every positive dimension and every
odd prime; the larger prime exclusions in the theorem are used
only for the excellent Cartier relation.
\end{lemma}

\begin{proof}
For \(a\geq0\), write
\[
 U(a)=\prod_{\substack{1\leq i\leq a\\p\nmid i}}i.
\]
We shall use the exact identity
\[
 a!=p^{\lfloor a/p\rfloor}\lfloor a/p\rfloor!U(a).
\]
Write \(j=pl+b\), \(0\leq b<p\), and put \(H=pN\) and
\(v=v_p(H)=1+v_p(N)\).
For the unit index set
\[
 I=\{i:j-H<i\leq j,\ p\nmid i\},
\]
one has
\begin{equation}\label{P:eq:positive-reciprocal-block}
 \sum_{i\in I}\frac1i\in p^v\mathbb Z_p.
\end{equation}
Indeed, the interval has length \(H\), a multiple of \(p^v\).
Its residues modulo \(p^v\) consist of \(H/p^v\) complete
copies of the residue classes.  Among the unit classes, inversion
permutes the residues, and opposite residues cancel because
\(p\) is odd.

For the cubic case, first omit the power \(n\) and write
\(b(j,N)=\binom{2j}{j-N}\).
With \(\epsilon=\lfloor2b/p\rfloor\in\{0,1\}\), the
factorial identity gives
\[
 d_j^{(1)}=\frac{b(j,0)}{b(l,0)}
 =p^\epsilon(2l+1)^\epsilon\frac{U(2j)}{U(j)^2}
 \in\mathbb Z_p,
\]
and
\[
 b(j,H)=d_j^{(1)}b(l,N)S,\qquad
 S=\frac{U(j)^2}{U(j-H)U(j+H)}
   =\prod_{i\in I}\left(1+\frac H i\right)^{-1}.
\]
Since each \(H/i\) is in \(p^v\mathbb Z_p\),
\[
 S\equiv1-H\sum_{i\in I}\frac1i\equiv1\pmod{p^{2v}}.
\]
Taking the \(n\)-th power proves \eqref{P:eq:positive-block}.

For the simplicial case set \(\epsilon=\lfloor rb/p\rfloor\).
Then
\[
 d_j=\frac{a_j}{a_l}
 =p^\epsilon\frac{(rl+\epsilon)!}{(rl)!}
             \frac{U(rj)}{U(j)^r}\in\mathbb Z_p.
\]
The balanced shifts in the factorial denominator give the exact
identity
\[
 B(j,H)=d_jB(l,N)S,\qquad
 S=\frac{U(j)^r}
 {U(j+(r-1)H)U(j-H)^{r-1}}.
\]
Decomposing the long upper interval into \(r-1\) blocks of
length \(H\) yields
\[
 S=\prod_{i\in I}\prod_{k=1}^{r-1}
          \left(1+\frac{kH}{i}\right)^{-1}.
\]
Therefore
\[
 S\equiv
 1-H\frac{r(r-1)}2\sum_{i\in I}\frac1i
 \equiv1\pmod{p^{2v}},
\]
again by \eqref{P:eq:positive-reciprocal-block}.
All prefactors multiplying \(S-1\) are \(p\)-integral,
which proves the lemma.
\end{proof}

\subsection{An integral triangular descent}

The following argument is stated in the notation of the two
families, but uses only integrality of \(B\), its diagonal entries
one, and Lemma~\ref{P:lem:positive-block}.
For each positive integer \(K\), define integers \(c_{K,N}\) by
\begin{equation}\label{P:eq:positive-triangular}
 a_j\mathbf1_{j\geq K}
 =\sum_{N=K}^{j}c_{K,N}B(j,N),\qquad j\geq0,
\end{equation}
and put \(c_{K,N}=0\) for \(N<K\).
Existence, uniqueness and integrality follow by successive solution
in \(j\), since \(B(j,j)=1\).

\begin{lemma}[Weighted interpolation and descent]
\label{P:lem:positive-descent}
For every \(K,N\geq1\),
\begin{equation}\label{P:eq:positive-c-weight}
 v_p(c_{K,N})\geq
 \max\{0,\,2v_p(K)-2v_p(N)\}.
\end{equation}
If \(p\mid K\), then
\begin{equation}\label{P:eq:positive-c-descent}
 c_{K,N}\equiv
 \begin{cases}
 c_{K/p,N/p},&p\mid N,\\
 0,&p\nmid N
 \end{cases}
 \pmod{p^{2v_p(K)}}.
\end{equation}
\end{lemma}

\begin{proof}
Induct on \(v=v_p(K)\).  If \(v=0\), the weighted bound
is integrality and there is no descent assertion.
Suppose \(v\geq1\), set \(L=K/p\), and assume the weighted
bound for \(L\).  Let \(l=\lfloor j/p\rfloor\).
Multiplying the defining identity for \(L\) by \(d_j=a_j/a_l\)
gives
\[
 a_j\mathbf1_{j\geq K}
 =d_ja_l\mathbf1_{l\geq L}
 =\sum_{N=L}^{l}c_{L,N}d_jB(l,N).
\]
By Lemma~\ref{P:lem:positive-block}, replacing each
\(d_jB(l,N)\) by \(B(j,pN)\) makes an error of valuation at
least
\[
 2+2v_p(N)+\max\{0,\,2(v-1)-2v_p(N)\}\geq2v.
\]
Thus the candidate coefficient vector, equal to \(c_{L,N}\)
at index \(pN\) and zero at indices not divisible by \(p\),
satisfies \eqref{P:eq:positive-triangular} modulo \(p^{2v}\).
Successively solving this unit triangular system proves
\eqref{P:eq:positive-c-descent}.
For \(p\nmid N\), it also gives the weighted bound directly.
For \(p\mid N\), the induction hypothesis applied to
\(c_{L,N/p}\) gives the bound
\(\max\{0,2(v-v_p(N))\}\), and the error \(p^{2v}\)
does not weaken it.
\end{proof}

\subsection{Positive vertex coefficients and the excellent lift}

For the cubic family choose the vertex \({\bf v}=(1,\ldots,1)\).
For the simplex choose \({\bf v}=(-1,\ldots,-1)\).
The coefficient at \({\bf x}^{hN{\bf v}}\) in the
\(z\)-adic expansion of \(1/f\) is
\begin{equation}\label{P:eq:positive-omega}
 \Omega_N(z)=\sum_{j\geq N}B(j,N)z^{hj}.
\end{equation}
For the cube this follows by choosing \(j-N\) negative steps in
each of the \(n\) independent factors of \(g^{2j}\).
For the simplex, the required term of \(g^{rj}\) uses the
chosen vertex \(j+(r-1)N\) times and each of the other
\(r-1\) vertices \(j-N\) times, giving exactly the displayed
multinomial coefficient.

Set
\[
 C_N(u)=\frac{\Omega_N(z)}{F(u)}\in u^N\mathbb Z_p[[u]].
\]
The coefficientwise interpretation of the module of second formal
derivatives, including in the positive \(z\)-adic expansion, is
\cite[Section 3, the paragraph following Corollary 3.4]{BV3}.
Read \cite[Theorem 7.3]{BV3} at the exponent vector
\(hN{\bf v}\).  The left coefficient is \(\Omega_{pN}(z)\),
and the right coefficient is
\(\mathcal F(z)\Omega_N(\tau)/\mathcal F(\tau)\).
The error is divisible by \(p^2(hN)^2\) in \(\mathbb Z_p[[z]]\).
Since \(p\nmid h\) and \(F(0)=1\), division by \(F(u)\)
gives
\begin{equation}\label{P:eq:positive-vertex-cartier}
 C_{pN}(u)-C_N(\sigma(u))
 \in p^{2+2v_p(N)}\mathbb Z_p[[u]].
\end{equation}

\begin{proof}[Proof of Theorem~\ref{P:thm:positive-cube-simplex}]
Summing \eqref{P:eq:positive-triangular} in \(j\), using
\eqref{P:eq:positive-omega}, and dividing by \(F(u)\), gives
the exact identity
\[
 1-\frac{F_K(u)}{F(u)}
 =\sum_{N\geq K}c_{K,N}C_N(u).
\]
It converges \(u\)-adically because \(C_N\) has order at
least \(N\).
Put \(v=v_p(K)\geq1\) and \(L=K/p\).
Lemma~\ref{P:lem:positive-descent} first gives
\[
 1-\frac{F_K(u)}{F(u)}
 \equiv\sum_{N\geq L}c_{L,N}C_{pN}(u)
 \pmod{p^{2v}}.
\]
By \eqref{P:eq:positive-vertex-cartier} and
\eqref{P:eq:positive-c-weight}, each summand may be replaced by
\(c_{L,N}C_N(\sigma(u))\) with error divisible by \(p^{2v}\).
Indeed, the error has valuation at least
\[
 2+2v_p(N)+\max\{0,2(v-1)-2v_p(N)\}\geq2v.
\]
Both sums still converge: \(C_{pN}(u)\) and
\(C_N(\sigma(u))\) have order at least \(pN\).
Consequently
\[
 1-\frac{F_K(u)}{F(u)}
 \equiv\sum_{N\geq L}c_{L,N}C_N(\sigma(u))
 =1-\frac{F_L(\sigma(u))}{F(\sigma(u))}
 \pmod{p^{2v}},
\]
which is \eqref{P:eq:positive-main}.
Multiplication by the period units gives the cleared-defect and
quotient formulations.  Finally,
\(\mathcal F_{hM}(z)=F_M(z^h)\) proves the assertion in the
original normalization.
\end{proof}

\paragraph{Scope.}
The proof proves an all-level statement for the ordinary positive
truncations, without a restriction on the series coefficient index.
It is independent of the earlier first-level and negative-basis
arguments.  It does not by itself prove the hyperoctahedral case:
that positive vertex matrix contains sums of multinomial
coefficients, and Lemma~\ref{P:lem:positive-block} in its present
scalar form does not apply to those sums.


\section{A matrix block argument for the hyperoctahedral family}\label{P:source-main-3}
\label{P:sec:octa-matrix-interpolation}

Let \(n\geq3\) and let \(p>n\) be prime, so in particular
\(p\geq5\).  Put
\[
 g({\bf x})=\sum_{i=1}^n(x_i+x_i^{-1}),\quad
 f=1-zg,\quad u=z^2,
\]
and write
\[
 F(u)=\operatorname{CT}(1/f)=\sum_{j\geq0}a_ju^j,
 \qquad F_K(u)=\sum_{j<K}a_ju^j.
\]
The normalized excellent lift is denoted by \(\sigma\), exactly
as in the preceding section.  Define the positive vertex matrix
\begin{equation}\label{P:eq:octa-B}
 B(j,N)=(2j)!
 \sum_{\substack{k_1+\cdots+k_n=j\\k_1\geq N}}
 \frac1{(k_1-N)!(k_1+N)!\prod_{i=2}^n(k_i!)^2},
 \qquad 0\leq N\leq j.
\end{equation}
Set \(B(j,N)=0\) for \(N>j\).
These are nonnegative integers, \(B(j,0)=a_j\), and
\(B(j,j)=1\).  The coefficient at \(x_1^{2N}\), with all
other exponents zero, in \(1/f\) is
\[
 \Omega_N(z)=\sum_{j\geq N}B(j,N)z^{2j}.
\]
Consequently the positive vertex functions
\[
 C_N(u)=\Omega_N(z)/F(u)\in u^N\mathbb Z_p[[u]]
\]
satisfy, by the excellent Cartier relation and the positive
coefficient interpretation of second formal derivatives,
\begin{equation}\label{P:eq:octa-positive-cartier}
 C_{pN}(u)-C_N(\sigma(u))
 \in p^{2+2v_p(N)}\mathbb Z_p[[u]].
\end{equation}
Here one reads the coefficient at \(x_1^{2N}\) in
\cite[Theorem 7.3]{BV3}; see also the paragraph following
\cite[Corollary 3.4]{BV3}.  Thus the remaining problem is
integral interpolation of the ordinary period tail in the
\(C_N\).

\begin{theorem}[All-level hyperoctahedral truncations]
\label{P:thm:octa-all-levels}
For every \(n\geq3\), every prime \(p>n\), and every positive
integer \(K\) divisible by \(p\),
\[
 \frac{F_K(u)}{F(u)}
 \equiv\frac{F_{K/p}(\sigma(u))}{F(\sigma(u))}
 \pmod{p^{2v_p(K)}\mathbb Z_p[[u]]}.
\]
In particular the cleared excellent-lift defect is in
\(p^{2s}\mathbb Z_p[[u]]\) for every normalized \(m,s\geq1\).
In the original variable, this includes all truncation lengths
\(2mp^s\), and hence the prescribed lengths in Conjecture 7.5.
\end{theorem}

\paragraph{Dimensions one and two.}
For \(n=1\), the octahedral and cubic Laurent polynomials are
identical.  For \(n=2\), the substitution
\(x_1=y_1y_2\), \(x_2=y_1/y_2\) gives
\[
 x_1+x_1^{-1}+x_2+x_2^{-1}
 =(y_1+y_1^{-1})(y_2+y_2^{-1}).
\]
Reading the coefficient at \(x_1^{2N}x_2^0\) gives
\(B(j,N)=\binom{2j}{j-N}^2\).  The positive vertex matrix
is therefore exactly the cubic matrix in dimension two.
Applying the scalar block proof with the octahedral excellent
Cartier relation proves the same conclusion at every odd prime,
including \(p=3\).  Thus the low-dimensional cases introduce
no additional exception to the original symmetry range.

The scalar ratio \(a_j/a_{\lfloor j/p\rfloor}\) need not give
the block congruence used for cubes and simplices.  We construct
an integral lower triangular row operator in its place.

\subsection{An integral analytic unit-factorial expansion}

Write \(\mathbb Z_p\langle X_1,\ldots,X_d\rangle\) for the
restricted power series ring: its coefficients tend to zero
\(p\)-adically as their total degree tends to infinity.
For a nonnegative integer \(a\), put
\[
 U(a)=\prod_{\substack{1\leq i\leq a\\p\nmid i}}i,
 \qquad c=(p-1)!.
\]

\begin{lemma}[Unit-factorial interpolation]
\label{P:lem:octa-unit-factorial}
For every \(0\leq r<p\), there is
\(T_r(X)\in\mathbb Z_p\langle X\rangle\) such that, for all
integers \(\ell\geq0\),
\begin{equation}\label{P:eq:octa-unit-expansion}
 U(p\ell+r)^{-2}
 =c^{-2\ell}(r!)^{-2}
   \bigl(1-2pH_r\ell+p^2T_r(\ell)\bigr),
 \qquad H_r=\sum_{a=1}^r\frac1a.
\end{equation}
\end{lemma}

\begin{proof}
The unit factorial over one block has the exact paired form
\[
 Q(k):=\frac{U(p(k+1))}{cU(pk)}
 =\prod_{a=1}^{(p-1)/2}
       \left(1+\frac{p^2k(k+1)}{a(p-a)}\right).
\]
Thus \(\log Q(X)\) is a sum, indexed by \(m\geq1\), of
polynomials of degree at most \(2m\) whose coefficients lie in
\(p^{2m-v_p(m)}\mathbb Z_p\).
The polynomial \(S_d(X)\), equal to
\(\sum_{0\leq k<\ell}k^d\) at nonnegative integers
\(X=\ell\), has coefficient denominators dividing
\(\operatorname{lcm}(1,\ldots,d+1)\).  Indeed, expansion in
falling factorials gives
\[
 k^d=\sum_{a=0}^d
 \left\{\begin{matrix}d\\a\end{matrix}\right\}
 k(k-1)\cdots(k-a+1),
\]
and the sum of the falling factorial of degree \(a\) is the
falling factorial of degree \(a+1\), divided by \(a+1\).
Taking this discrete antiderivative term by term in
\(\log Q\), the \(m\)-th contribution therefore has
coefficient valuation at least
\[
 2m-v_p(m)-\lfloor\log_p(2m+1)\rfloor\geq2.
\]
The bound tends to infinity with \(m\), since \(p\geq5\).
We obtain a restricted series
\(A(X)\in p^2\mathbb Z_p\langle X\rangle\) such that
\[
 A(\ell)=\sum_{k=0}^{\ell-1}\log Q(k)
 =\log\bigl(U(p\ell)/c^\ell\bigr).
\]
Exponentiation converges in the restricted series ring and gives
\(\exp(-2A)\in1+p^2\mathbb Z_p\langle X\rangle\).
Finally,
\[
 U(p\ell+r)=U(p\ell)\,r!
                \prod_{a=1}^r(1+p\ell/a),
\]
and
\[
 \prod_{a=1}^r(1+pX/a)^{-2}
 \in1-2pH_rX+p^2\mathbb Z_p\langle X\rangle.
\]
Combining these identities proves the lemma.
\end{proof}

\subsection{Symmetric Bessel operators}

For an integer \(N\geq0\), put
\[
 I_N(x)=\sum_{k\geq N}\frac{x^k}{(k-N)!(k+N)!},\qquad
 \Phi_N(x)=I_N(x)I_0(x)^{n-1}.
\]
For \(\theta=x\,d/dx\), one has
\[
 (\theta^2-N^2)I_N=xI_N.
\]
Use the factorial coefficient lattice
\[
 \mathcal H=\left\{\sum_{j\geq0}h_jx^j:
                          (2j)!h_j\in\mathbb Z_p\right\}.
\]
It is a complete normed \(\mathbb Z_p\)-algebra under the coefficient
norm \(\sup_j|(2j)!h_j|_p\); multiplication, \(\theta\),
and multiplication by \(x\) preserve its unit ball.
In particular, \(I_N\), \(\Phi_N\), and all their
\(\theta\)-derivatives belong to \(\mathcal H\).

\begin{lemma}[Integral symmetric operator reduction]
\label{P:lem:octa-symmetric-reduction}
Let
\(P\in\mathbb Z_p\langle X_1,\ldots,X_n\rangle\)
be symmetric in its \(n\) variables, and let
\(\theta_i=x_i\,\partial/\partial x_i\).
There are restricted series
\(A_0(x),\ldots,A_n(x)\in\mathbb Z_p\langle x\rangle\),
independent of \(N\), such that
\begin{align}\label{P:eq:octa-reduction}
 &\left.
 P(\theta_1,\ldots,\theta_n)
 I_N(x_1)\prod_{i=2}^n I_0(x_i)
 \right|_{x_1=\cdots=x_n=x}
 \nonumber\\
 &\hspace{25mm}
 -\sum_{k=0}^n A_k(x)\theta^k\Phi_N(x)
 \in N^2\mathcal H.
\end{align}
The construction is \(\mathbb Z_p\)-linear and bounded for
the coefficient norms.  If
\[
 P=\alpha+\beta(X_1+\cdots+X_n)+p^2R,
 \qquad R\in\mathbb Z_p\langle X_1,\ldots,X_n\rangle,
\]
the error in \eqref{P:eq:octa-reduction} belongs to
\(p^2N^2\mathcal H\), and every positive-degree coefficient
of every \(A_k\) is divisible by \(p^2\).
At \(N=0\), \eqref{P:eq:octa-reduction} is an exact identity.
\end{lemma}

\begin{proof}
First take a symmetric polynomial \(P\).
Using the equations \(\theta_i^2I_0=x_iI_0\), and using
\(\theta_1^2I_N=x_1I_N+N^2I_N\), reduce each differentiated
factor to an integral polynomial combination of the factor and
its first derivative, modulo \(N^2\).
All commutations \(\theta_i x_i=x_i(\theta_i+1)\) have
integer coefficients.  After setting all \(x_i=x\), symmetry
of the reduction with the \(N^2\) terms suppressed gives a
combination, with coefficients in \(\mathbb Z_p[x]\), of
\[
 E_k=\sum_{\substack{S\subseteq\{1,\ldots,n\}\\|S|=k}}
       \prod_{i\in S}\theta_i f_i(x)
       \prod_{i\notin S}f_i(x),
 \quad f_1=I_N,\quad f_i=I_0\ (i>1).
\]
Here \(E_0=\Phi_N\), \(E_1=\theta\Phi_N\), and
\[
 \theta E_k=(k+1)E_{k+1}+(n-k+1)xE_{k-1}
                   \pmod{N^2\mathcal H}.
\]
The correction comes only from differentiating a first derivative
of \(I_N\) and is \(N^2\) times an integral product of
the indicated functions.  Solving successively for
\(E_2,\ldots,E_n\) divides only by \(2,\ldots,n\),
which are units because \(p>n\).
It expresses the result in the form asserted in
\eqref{P:eq:octa-reduction}, with \(A_k\in\mathbb Z_p[x]\)
independent of \(N\).

Every error lies in \(N^2\mathcal H\): a coefficient of
degree \(j\), multiplied by \((2j)!\), is an integral
multinomial coefficient times integral powers of summation
indices and coefficients in \(\mathbb Z_p\).
Multiplication by \(x\) and differentiation by \(\theta\)
preserve this property.  The reduction never decreases a
coefficient's \(p\)-adic valuation, since its only denominators
are products of units from \(1,\ldots,n\).
It is therefore bounded and extends continuously from symmetric
polynomials to symmetric restricted series.  The output
coefficients also tend to zero with degree, giving the asserted
restricted \(A_k\).
For the affine symmetric polynomial
\(\alpha+\beta\sum X_i\), take
\(A_0=\alpha\), \(A_1=\beta\), and all other
\(A_k=0\); this is exact for every \(N\).
Apply the bounded linear construction to \(p^2R\) for the
last assertion.  The exactness at \(N=0\) is immediate from
the original differential equations.
\end{proof}

\clearpage
\subsection{The lower triangular row operator}

\begin{lemma}[Matrix block congruence with a boundary factor]
\label{P:lem:octa-matrix-block}
For each integer \(j\geq0\), put \(L=\lfloor j/p\rfloor\).
There exist \(D_{j,i}\in\mathbb Z_p\), \(0\leq i\leq L\),
such that
\begin{align}
 a_j&=\sum_{i=0}^{L}D_{j,i}a_i,
                    \label{P:eq:octa-central-exact}\\
 B(j,pN)-\sum_{i=0}^{L}D_{j,i}B(i,N)
    &\in p^{2+2v_p(N)}\mathbb Z_p\quad(N\geq1),
                    \label{P:eq:octa-row-block}\\
 D_{j,i}a_i&\in
       p^2\left(\frac{L!}{i!}\right)^2\mathbb Z_p
              \quad(0\leq i<L).
                    \label{P:eq:octa-row-boundary}
\end{align}
\end{lemma}

\begin{proof}
Write \(j=pL+b\), \(0\leq b<p\), and
\(\epsilon=\lfloor2b/p\rfloor\).
In \eqref{P:eq:octa-B}, decompose
\(k_i=p\ell_i+r_i\), \(0\leq r_i<p\).
There is a carry \(d\), \(0\leq d\leq n-1\), such that
\[
 \sum_i r_i=b+pd,\qquad \sum_i\ell_i=L-d.
\]
Terms with \(d>L\) are absent.  Let \(H=pN\).
For the first factorial pair, one has
\[
 \frac{U(k_1)^2}{U(k_1-H)U(k_1+H)}
 \equiv1\pmod{p^{2+2v_p(N)}}.
\]
This is the symmetric unit-factorial block calculation of
Lemma~\ref{P:lem:positive-block}, applied to the interval of
length \(H\) ending at \(k_1\).
The corresponding factorial summand is a \(p\)-integer and
the displayed quotient is a unit, so replacing that quotient
by one preserves the claimed modulus term by term.

For each \(d\), define the symmetric restricted series
\[
 P_{b,d}(X_1,\ldots,X_n)
 =\sum_{\substack{0\leq r_i<p\\\sum r_i=b+pd}}
       \prod_{i=1}^n
       (r_i!)^{-2}
       \bigl(1-2pH_{r_i}X_i+p^2T_{r_i}(X_i)\bigr).
\]
Symmetry of the residue-index set implies
\[
 P_{b,d}=\alpha_{b,d}
       +\beta_{b,d}(X_1+\cdots+X_n)+p^2R_{b,d},
 \qquad R_{b,d}\in\mathbb Z_p\langle X_1,\ldots,X_n\rangle.
\]
In particular, the coefficients of the linear terms are all equal;
no division by the period \(a_L\) is used.
The exact factorial identity and
Lemma~\ref{P:lem:octa-unit-factorial} now give
\begin{align}\label{P:eq:octa-block-Bessel}
 B(j,pN)\equiv{}
 &p^\epsilon U(2j)(2L+\epsilon)!
 \sum_{0\leq d\leq\min(n-1,L)}
 p^{2d}c^{-2(L-d)}\,[x^{L-d}]
 \nonumber\\[-2pt]
 &\left.
 P_{b,d}(\theta_1,\ldots,\theta_n)
 I_N(x_1)\prod_{i=2}^n I_0(x_i)
 \right|_{x_1=\cdots=x_n=x}
 \pmod{p^{2+2v_p(N)}}.
\end{align}
For \(N=0\), this is an exact equality.
Indeed, after taking the \(p\)-parts of the factorials, the
power of \(p\) is \(2d+\epsilon\), their numerator is
\((2L+\epsilon)!\), and their denominator is
\((\ell_1-N)!(\ell_1+N)!\prod_{i>1}(\ell_i!)^2\),
which verifies all factors in \eqref{P:eq:octa-block-Bessel}.

Apply Lemma~\ref{P:lem:octa-symmetric-reduction} to each
\(P_{b,d}\), writing its output as
\[
 A_{d,k}(x)=\sum_{a\geq0}\alpha_{d,a,k}x^a,
 \qquad 0\leq k\leq n.
\]
For \(a>0\), all \(\alpha_{d,a,k}\) are in
\(p^2\mathbb Z_p\).  The reduction error is in
\(p^2N^2\mathcal H\).  Multiplying its coefficient of
degree \(L-d\) by \((2L+\epsilon)!\) preserves this
divisibility, since
\((2L+\epsilon)!/(2(L-d))!\) is an integer.
As \([x^i]\Phi_N=B(i,N)/(2i)!\), define
\begin{align}\label{P:eq:octa-D-definition}
 D_{j,i}={}&
 p^\epsilon U(2j)\frac{(2L+\epsilon)!}{(2i)!}
 \sum_{d=0}^{\min(n-1,L-i)}
 p^{2d}c^{-2(L-d)}
 \sum_{k=0}^{n}\alpha_{d,L-d-i,k}\,i^k.
\end{align}
The convention \(i^0=1\) applies also at \(i=0\).
All terms are \(p\)-integral.  Equations
\eqref{P:eq:octa-central-exact} and \eqref{P:eq:octa-row-block}
follow from the exactness at \(N=0\) and the just-proved
error bound for \(N\geq1\).

If \(i<L\), each summand in the inner expression of
\eqref{P:eq:octa-D-definition} contains \(p^2\): either
\(d\geq1\), or \(d=0\) and \(L-d-i>0\).
Therefore
\[
 D_{j,i}\in p^2\frac{(2L+\epsilon)!}{(2i)!}\mathbb Z_p.
\]
Finally,
\[
 a_i=\frac{(2i)!}{(i!)^2}
       \sum_{k_1+\cdots+k_n=i}
       \binom{i}{k_1,\ldots,k_n}^{\!2}.
\]
The sum is an integer, and
\[
 \frac{(2L+\epsilon)!}{(i!)^2}
 =\frac{(2L+\epsilon)!}{(L!)^2}
       \left(\frac{L!}{i!}\right)^2
\]
is an integer multiple of \((L!/i!)^2\).
This proves \eqref{P:eq:octa-row-boundary}.
\end{proof}

\subsection{Descent across the truncation boundary}

Define integers \(c_{K,N}\), zero for \(N<K\), by
\[
 a_j\mathbf1_{j\geq K}
    =\sum_{N=K}^j c_{K,N}B(j,N).
\]
The diagonal entries of \(B\) are one, so these integers
exist uniquely.

\begin{lemma}[Octahedral positive interpolation]
\label{P:lem:octa-positive-descent}
For all positive \(K,N\),
\[
 v_p(c_{K,N})\geq\max\{0,2v_p(K)-2v_p(N)\}.
\]
If \(p\mid K\), then
\[
 c_{K,N}\equiv
 \begin{cases}
 c_{K/p,N/p},&p\mid N,\\
 0,&p\nmid N
 \end{cases}
 \pmod{p^{2v_p(K)}}.
\]
\end{lemma}

\begin{proof}
Induct on \(v=v_p(K)\), with integrality as the base case.
Let \(v\geq1\), \(M=K/p\), and assume the weighted bound
for \(c_{M,N}\).  Consider
\(\sum_{N\geq M}c_{M,N}B(j,pN)\).
For \(j<K\), it is zero, as required.
For \(j\geq K\), put \(L=\lfloor j/p\rfloor\geq M\).
Lemma~\ref{P:lem:octa-matrix-block}, together with the induction
hypothesis, gives
\begin{align*}
 \sum_{N=M}^{L}c_{M,N}B(j,pN)
 &\equiv\sum_{i=0}^{L}D_{j,i}
                \sum_{N=M}^{i}c_{M,N}B(i,N)\\
 &=\sum_{i=M}^{L}D_{j,i}a_i
 =a_j-\sum_{i<M}D_{j,i}a_i
 \pmod{p^{2v}}.
\end{align*}
The errors in the first congruence have valuation at least
\[
 2+2v_p(N)+\max\{0,2(v-1)-2v_p(N)\}\geq2v.
\]
For \(i<M\leq L\), the integer \(L!/i!\) is divisible by
\(M\).  The boundary estimate
\eqref{P:eq:octa-row-boundary} thus puts each omitted term in
\(p^2M^2\mathbb Z_p\subseteq p^{2v}\mathbb Z_p\).
The candidate coefficient vector, supported on multiples of
\(p\), therefore solves the defining triangular system for
\(c_{K,N}\) modulo \(p^{2v}\).
Unit triangular inversion proves the descent congruence.
The weighted bound follows from it exactly as before: for
\(p\nmid N\) use the zero candidate, and for \(p\mid N\)
apply the induction hypothesis to \(c_{M,N/p}\).
\end{proof}

\begin{proof}[Proof of Theorem~\ref{P:thm:octa-all-levels}]
The defining triangular identity gives the exact expansion
\[
 1-\frac{F_K(u)}{F(u)}
    =\sum_{N\geq K}c_{K,N}C_N(u).
\]
Let \(v=v_p(K)\geq1\) and \(M=K/p\).
Lemma~\ref{P:lem:octa-positive-descent} gives
\[
 1-\frac{F_K(u)}{F(u)}
 \equiv\sum_{N\geq M}c_{M,N}C_{pN}(u)
 \pmod{p^{2v}}.
\]
Use \eqref{P:eq:octa-positive-cartier} and the weighted bound
on \(c_{M,N}\) to replace \(C_{pN}\) by
\(C_N\circ\sigma\) at the same precision.  The result is
\[
 1-\frac{F_K(u)}{F(u)}
 \equiv\sum_{N\geq M}c_{M,N}C_N(\sigma(u))
 =1-\frac{F_M(\sigma(u))}{F(\sigma(u))}
 \pmod{p^{2v}}.
\]
All sums are legitimate coefficientwise, since their summands
have order at least \(N\), respectively \(pN\).
Multiplication by the period units and the identity
\(\mathcal F_{2M}(z)=F_M(z^2)\) finish the proof.
\end{proof}



\section{Verification, sharpness, and scope}\label{P:source-main-4}
\label{P:sec:verification}

\subsection{What has been closed}

Theorem~\ref{P:thm:main-all} supplies the full lower bound that was
missing for each of the targets in the table of
Section~\ref{P:sec:main-scope}. It proves the ordinary positive
truncation statement for the actual excellent lift, with no
restriction on the coefficient index and no upper bound on $s$.
In particular the proof for the three-dimensional octahedral
family is uniform in $p\ge5$ and does not stop at a fixed level.

The theorem establishes the exponent $2s$. A uniform exponent
$3s$ already fails for the three-dimensional octahedral family:
exact integer evaluation gives
\[
 [u^{125}]E_{5,3,1}=a_{125}-a_{25}
       \equiv2\cdot5^6\pmod{5^7}.
\]
Thus this coefficient has valuation $6=2s$, which is less than $3s=9$.
Any previously established matching upper bounds for a coefficient
can now be combined with Theorem~\ref{P:thm:main-all} to give
optimality in their precise parameter ranges.

This original argument has the prime range displayed in its theorem;
the complete edition has the larger range of Theorem~\ref{G:thm:global}.
The elliptic-zero and Zagier D arguments are retained in Parts~\ref{part:modular}--\ref{part:twocusp}. They are separate results and are not dependencies of this theorem. Within the three standard families and the displayed prime conditions, the theorem covers every dimension, level, multiplier, and coefficient.

\subsection{Why a scalar octahedral block is insufficient}

Here is an exact counterexample to using the cubic scalar block
argument unchanged for octahedra. In dimension three, take
$p=5$, $j=10$, $L=2$, and $N=1$. Then
\[
\begin{aligned}
 a_{10}&=25\,989\,269\,017\,140,& a_2&=90,\\
 B(10,5)&=20\,650\,661\,328,& B(2,1)&=28.
\end{aligned}
\]
The ratio $a_{10}/a_2=288\,769\,655\,746$ is integral, but
\[
 B(10,5)-(a_{10}/a_2)B(2,1)
 =-8\,064\,899\,699\,560\equiv15\pmod{25}.
\]
Its $5$-adic valuation is one. The matrix block lemma avoids
this failed scalar assertion and controls the necessary sum
over lower rows, including the boundary terms.

\subsection{Exact arithmetic accompanying the proofs}

Appendix~\ref{app:all-programs} contains the standalone verification programs in full, using integer, rational, and modular polynomial arithmetic.
The general proofs above do not depend on the number of
coefficients tested. The computations check signs, indices,
normalizations, and the actual constructions used in the proof.

\begin{itemize}
\item \texttt{verify\_positive\_interpolation.py} checks the
balanced factorial block congruence, the weighted triangular
descent, and independently constructed excellent-lift defects
for cubes and simplices. Its report records eight block cases
(12,949 matrix entries), eleven interpolation cases, and seven
direct-defect cases.
The pure factorial tests intentionally include some primes
outside the geometric theorem's range; these are not reported
as extensions of the theorem.
\item \texttt{verify\_octahedral\_interpolation.py} checks the restricted
unit-factorial and symmetric-operator construction, the resulting
row matrix, and the interpolation descent. At $p=5$ and precision
$p^6$, it constructs every row through $j=145$ in dimensions
three and four, checking 4,118 shifted columns and 4,118 boundary
terms in total. It also checks five interpolation cases up to
column 300 and three independently constructed actual-BV defects,
including $s=3$ in dimensions three and four. Exact parameters and precisions are specified in the embedded program, which regenerates its JSON report.
\item The separate program
\texttt{certificate\_octa4\_p5\_s2.py} constructs the actual
lift modulo $625$ from negative vertex polynomials and checks
a finite polynomial identity of degree at most $2819$.
Appendix~\ref{P:sec:octa4-certificate} proves why that identity
certifies every coefficient at $n=4,p=5,s=2$. This provides a
different exact route to one case of the general theorem.
\end{itemize}

All output reports distinguish bounded validation from a theorem
about infinitely many coefficients. The independently certified
whole-series case is labelled as such because its proof includes
a finite polynomial identity with an explicit denominator that
is a unit.


% END RETAINED SOURCE P

\partstart{All odd primes for cubes and simplices; the common formal input}{part:small-primes}{The orbit module, ramified normalization, scalar descent, independent first levels, and the period-functional obstruction.}
% BEGIN RETAINED SOURCE SP
\section{Scope and results}\label{SP:source-main-1}
For the cube, simplex, and octahedron put respectively
\[
g=\prod_{i=1}^n(x_i+x_i^{-1}),\qquad
g=x_1+\cdots+x_n+(x_1\cdots x_n)^{-1},\qquad
g=\sum_{i=1}^n(x_i+x_i^{-1}).
\]
Let $f=1-zg$, $h=2,n+1,2$, $u=z^h$, and
\[
F_0(z)=\CT(1/f)=F(u),\quad F(u)=\sum_{j\ge0}a_ju^j,
\quad F_K(u)=\sum_{j<K}a_ju^j.
\]
We also write
\[
\mathcal F(z)=F_0(z),\qquad
\mathcal F_L(z)=\sum_{0\le j<L}\CT(g^j)z^j.
\]
The canonical excellent lift is denoted by $\tau(z)$ and its normalized
form by $\sigma(u)=\tau(z)^h$. Define
\[
E_{p,s,m}=F(u)F_{mp^{s-1}}(\sigma(u))
             -F(\sigma(u))F_{mp^s}(u).
\]
Every assertion below is coefficientwise in the indicated formal-series
ring. No arbitrary parameter specialization is asserted.

\begin{theorem}[Odd-prime scope of this addendum]\label{SP:thm:scope}
For every dimension $n\ge1$ and every odd prime $p$:
\begin{enumerate}
\item In the cube and simplex families, for every $p\mid K$,
\[
\frac{F_K}{F}\equiv\frac{F_{K/p}\circ\sigma}{F\circ\sigma}
\pmod{p^{2v_p(K)}\Zp[[u]]}.
\]
Thus $E_{p,s,m}\in p^{2s}\Zp[[u]]$ for all $m,s\ge1$.
The simplex statement includes $p\mid n+1$.
\item In the octahedral family, $E_{p,1,1}\in p^2\Zp[[u]]$.
\item The following additional octahedral statements hold:
\[
\begin{array}{ll}
n=4:&E_{3,2,1}\in81\Z_3[[u]],\quad
       E_{3,3,1}\in729\Z_3[[u]],\\
n=5:&E_{5,2,1}\in625\Z_5[[u]].
\end{array}
\]
The last three statements are proved by exact finite polynomial
certificates whose complete program is included in Appendix~\ref{SP:app:code}.
\end{enumerate}
\end{theorem}

The octahedral conclusions of this part are the independent first-level
argument and three exact certificates. The all-level strengthening at
every odd prime is supplied by Parts~\ref{part:octa-small} and~\ref{part:octa-three}.
The finite certificates here are not used to infer untested levels or
general multipliers.
The orbit-module argument below is a new verification using the general
results of \cite[Sections 2--5]{BV3}; it is not an application of
\cite[Theorem 7.3]{BV3} after dropping its stated hypotheses.
The argument is restricted to $R=\Zp[[z]]$. Its analytic continuation to
other base rings is not needed or claimed here.

\section{The excellent Cartier input without group averaging}\label{SP:source-main-2}
\begin{proposition}[Formal excellent input from an orbit module]\label{SP:prop:orbit}
Let $g$ be the hypercubic, simplicial, or hyperoctahedral Laurent
polynomial, with vertex coefficients one, and let $p$ be an odd prime.
Put $R=\mathbb Z_p[[z]]$ and $f=1-zg$.
Let $F_0(z)$ and $G_0(z)$ be the period and logarithmic companion
given by the vertex-coefficient formulas of BV3, Proposition~9.1.
Then
\[
 q_0(z)=z\exp(G_0(z)/F_0(z))\in z+z^2\mathbb Z_p[[z]].
\]
The lift $\tau=q_0^{-1}(q_0(z)^p)$ is integral and satisfies
\[
 \mathscr C_p(1/f)-\frac{F_0(z)}{F_0(\tau)}\frac1{f^\tau}
 \in p^2\mathscr F_2^\tau.
\]
In particular, no exclusion $p\mid n$ for the octahedron, or
$p\mid n+1$ for the simplex, is needed for this formal input.
\end{proposition}

\begin{proof}
Write $\Gamma$ for the support lattice and $H$ for the full monomial
symmetry group. Its order will never be inverted. The support lattice
is $\mathbb Z^n$ in the simplicial and octahedral cases; for a cube its
index is $2^{n-1}$. Thus its index is prime to $p$ in all three cases.
The symmetry group preserves the reflexive polytope degree
$d(u)=\deg_\Delta(u)$ and the support lattice. For an $H$-orbit
$O\subset\Gamma$, put
\[
 S_O=\sum_{u\in O}x^u,\qquad d(O)=d(u)\ (u\in O),\qquad
 L=\bigoplus_O R z^{d(O)}S_O.
\]
Here a Laurent polynomial uses only finitely many orbits.
This is a direct construction of a module of polynomials. We do not
obtain it by applying an averaging projector to another crystal.

Fix temporarily any lift $\sigma$ with
$z^\sigma=z^pv$, $v\in1+pR$. We check BV3, Definition~2.1.
The orbit sums have disjoint supports and coefficient one, so $L$ is
$p$-saturated. The same is true of every $L^{\sigma^i}$: its orbit
coefficient generator is $(z^{\sigma^i})^{d(O)}$, a power of $z$
times a unit, which is a non-zero-divisor modulo $p$.
Multiplication by $f$ preserves invariance and admissibility, hence
$fL\subset L$.

The subset $p\Gamma$ is $H$-stable. An orbit is either disjoint from
$p\Gamma$ or has the form $pO'$ for one orbit $O'$; the equality uses
the $p$-saturation of the support lattice. Since $d(pu)=pd(u)$,
\[
 \mathscr C_p(z^{d(pO')}S_{pO'})
   =z^{pd(O')}S_{O'}
   =\left(\frac{z^p}{z^\sigma}\right)^{d(O')}
        (z^\sigma)^{d(O')}S_{O'}\in L^\sigma.
\]
The image of any other orbit is zero. This proves compatibility
without dividing by either $|H|$ or $|O|$.

Form the Dwork crystal
$M=\widehat\Omega_{L,f}(\Delta^\circ)$ directly, as in BV3,
Definitions~2.6 and~2.8. We do not identify it with the invariant
part of an ambient completed module; in particular no exactness of
an invariant functor is being assumed. In all three families,
$\Delta\cap\Gamma$ consists of the origin and the vertices.
Reflexivity then implies that the first two levels have bases
\[
 M(1)=R/f,\qquad M(2)=R/f^2\oplus Rzg/f^2.
\]
Indeed $2\Delta^\circ\cap\Gamma=\Delta\cap\Gamma$ and the
vertices form one orbit. The pair $1/f,zg/f^2$ is an extended basis,
since $1/f=(1-zg)/f^2$ and the change of basis has determinant one.
These levels are $p$-saturated in the ambient integral formal
expansion: multiplying an expansion by $f$ or $f^2$ reduces the
claim to divisibility of the numerator coefficients. This also
shows that the two basis vectors remain independent modulo $p$.

The Hasse--Witt computations of BV3, Lemma~7.2, now apply to this
explicit compatible module. Their coefficient extraction uses a
single vertex coefficient of a polynomial known to be constant on
the vertex orbit; it does not average over that orbit. With the
initial lift $z^\sigma=z^p$, the first determinant has constant
term one, and
\[
 \left.\frac{\det HW^{(2)}(M)}p\right|_{z=0}
 =-\frac1p\binom{2p-2}{p-1}\in\mathbb Z_p^\times.
\]
The latter computation reads the coefficient of $z^px^{pv}$ at one
vertex $v$ in $(zg)F^{(2)}$; the extreme vertex forces the unique
contribution $(-z)^p\binom{2p-2}{p-1}$. The normalized determinant
modulo $p$ is independent of the lift. Consequently both
Hasse--Witt conditions hold over $R$ for every lift under consideration.
The general BV3, Corollary~5.9 and Corollary~4.3, therefore give
\[
 M=M(2)\oplus\mathscr F_2M,
 \qquad
 \mathscr C_p(M)\subset M^\sigma(2)+p^2\mathscr F_2M^\sigma.
\]
Only $p>2$ is required by these general results.

For completeness, the remaining canonical-coordinate argument
uses these statements rather than a group-order hypothesis.
The derivation $\theta=z\,d/dz$ preserves $M$. Thus the rank-two
quotient has the connection and Cartier matrices of BV3,
Propositions~7.6 and~7.7. The positive vertex coefficient calculation
of Proposition~9.1 gives exactly $F_0$ and
$F_0\log z+G_0$ as its two solutions. The same coefficient
calculation at exponents $p^sv$ gives $A(0)=B(0)=0$, the
constant Cartier matrix $\operatorname{diag}(1,p)$ for the
initial lift, and the upper-right coefficient
\[
 \lambda_1=
 \frac{pF_0(z)F_0(z^\sigma)}{W(z^\sigma)}
 \left(-\frac1p\log\frac{z^\sigma}{z^p}
 +\frac{G_0(z)}{F_0(z)}
 -\frac1p\frac{G_0(z^\sigma)}{F_0(z^\sigma)}\right),
\]
where $W=F_0^2+F_0\theta G_0-(\theta F_0)G_0$.
The rank-one reduction modulo $p$ implies $\lambda_1\in pR$.
The determinant identity
$\det\Lambda=pW/W^\sigma$ and $W(0)=1$ give $W\in1+zR$
by the same iteration as BV3, Proposition~7.9.
For $z^\sigma=z^p$, the displayed formula and the
Dieudonne--Dwork lemma yield $q_0\in z+z^2R$.
Its formal inverse is integral, so
$\tau=q_0^{-1}(q_0^p)$ defines a lift with
$\tau/z^p\in1+pR$. Reapply the compatible-module construction
with this lift. The displayed expression has $\lambda_1=0$,
and the remaining coefficient is $F_0/F_0^\tau$.
This proves the required excellent Cartier identity in
$p^2\mathscr F_2M^\tau$.

Finally, membership in $\mathscr F_2M$ means that every $s$-fold
Cartier iterate lies in $p^{2s}M^{\sigma^s}$. Mapping the directly
constructed module to the ambient Dwork module therefore sends
$\mathscr F_2M$ into the ambient $\mathscr F_2$.
Only this inclusion is needed; no converse assertion about
invariants or completed modules is used.
\end{proof}

\begin{remark}[Normalization even when $p\mid h$]
Let $h=2$ for cubes and octahedra and $h=n+1$ for simplices.
The explicit vertex formulas give $F_0(z)=F(z^h)$ and
$G_0(z)=G(z^h)$, so $q_0(z)=zQ(z^h)$. For every $h$th root of
unity $\zeta$ over an algebraic extension of $\mathbb Q_p$,
uniqueness of the formal inverse gives
$\tau(\zeta z)=\zeta^p\tau(z)$. It follows that
$\tau(z)=z^pV(z^h)$ with $V\in1+p\mathbb Z_p[[u]]$.
Consequently $\sigma(u)=\tau(z)^h=u^pV(u)^h$ is integral even
when $p\mid h$; no $h$th root is extracted. The positive-vertex
Cartier error is divisible by $p^2(hN)^2$, which is at least
as strong as $p^{2+2v_p(N)}$. The scalar block descent in the next section therefore applies to
cubes and simplices at every odd prime.
\end{remark}

\section{All odd primes for cubes and simplices}\label{SP:source-main-3}
This section reproduces the scalar interpolation argument, with its
external input replaced by Proposition~\ref{SP:prop:orbit}. This makes the
new prime range explicit and records all algebraic steps needed for it.
The coefficients are $a_j=\binom{2j}{j}^n$ for cubes and
$a_j=((n+1)j)!/(j!)^{n+1}$ for simplices. Throughout $p$ is any odd prime.
\begin{theorem}\label{SP:S:thm:positive-cube-simplex}
For these two families and every $K$ divisible by $p$,
\begin{equation}\label{SP:S:eq:positive-main}
\frac{F_K(u)}{F(u)}\equiv
\frac{F_{K/p}(\sigma(u))}{F(\sigma(u))}
\pmod{p^{2v_p(K)}\Zp[[u]]}.
\end{equation}
\end{theorem}
\subsection{A binomial and multinomial block congruence}

Set \(r=n+1\) in the simplicial case and define, for
\(0\leq N\leq j\),
\[
 B(j,N)=
 \begin{cases}
 \displaystyle\binom{2j}{j-N}^{n},&\text{hypercubic},\\[7pt]
 \displaystyle\frac{(rj)!}
 {(j+(r-1)N)!(j-N)!^{r-1}},&\text{simplicial}.
 \end{cases}
\]
Put \(B(j,N)=0\) for \(N>j\).
In both cases these are nonnegative integers with
\(B(j,0)=a_j\) and \(B(j,j)=1\).

\begin{lemma}[Balanced factorial block congruence]
\label{SP:S:lem:positive-block}
Let \(p\) be odd and \(l=\lfloor j/p\rfloor\).
Then \(d_j=a_j/a_l\in\mathbb Z_p\), and, for every
\(1\leq N\leq l\),
\begin{equation}\label{SP:S:eq:positive-block}
 B(j,pN)-d_j B(l,N)
 \in p^{2+2v_p(N)}\mathbb Z_p.
\end{equation}
This lemma itself holds for every positive dimension and every
odd prime; the excellent Cartier input is now supplied by Proposition~\ref{SP:prop:orbit}.
\end{lemma}

\begin{proof}
For \(a\geq0\), write
\[
 U(a)=\prod_{\substack{1\leq i\leq a\\p\nmid i}}i.
\]
We shall use the exact identity
\[
 a!=p^{\lfloor a/p\rfloor}\lfloor a/p\rfloor!U(a).
\]
Write \(j=pl+b\), \(0\leq b<p\), and put \(H=pN\) and
\(v=v_p(H)=1+v_p(N)\).
For the unit index set
\[
 I=\{i:j-H<i\leq j,\ p\nmid i\},
\]
one has
\begin{equation}\label{SP:S:eq:positive-reciprocal-block}
 \sum_{i\in I}\frac1i\in p^v\mathbb Z_p.
\end{equation}
Indeed, the interval has length \(H\), a multiple of \(p^v\).
Its residues modulo \(p^v\) consist of \(H/p^v\) complete
copies of the residue classes.  Among the unit classes, inversion
permutes the residues, and opposite residues cancel because
\(p\) is odd.

For the cubic case, first omit the power \(n\) and write
\(b(j,N)=\binom{2j}{j-N}\).
With \(\epsilon=\lfloor2b/p\rfloor\in\{0,1\}\), the
factorial identity gives
\[
 d_j^{(1)}=\frac{b(j,0)}{b(l,0)}
 =p^\epsilon(2l+1)^\epsilon\frac{U(2j)}{U(j)^2}
 \in\mathbb Z_p,
\]
and
\[
 b(j,H)=d_j^{(1)}b(l,N)S,\qquad
 S=\frac{U(j)^2}{U(j-H)U(j+H)}
   =\prod_{i\in I}\left(1+\frac H i\right)^{-1}.
\]
Since each \(H/i\) is in \(p^v\mathbb Z_p\),
\[
 S\equiv1-H\sum_{i\in I}\frac1i\equiv1\pmod{p^{2v}}.
\]
Taking the \(n\)-th power proves \eqref{SP:S:eq:positive-block}.

For the simplicial case set \(\epsilon=\lfloor rb/p\rfloor\).
Then
\[
 d_j=\frac{a_j}{a_l}
 =p^\epsilon\frac{(rl+\epsilon)!}{(rl)!}
             \frac{U(rj)}{U(j)^r}\in\mathbb Z_p.
\]
The balanced shifts in the factorial denominator give the exact
identity
\[
 B(j,H)=d_jB(l,N)S,\qquad
 S=\frac{U(j)^r}
 {U(j+(r-1)H)U(j-H)^{r-1}}.
\]
Decomposing the long upper interval into \(r-1\) blocks of
length \(H\) yields
\[
 S=\prod_{i\in I}\prod_{k=1}^{r-1}
          \left(1+\frac{kH}{i}\right)^{-1}.
\]
Therefore
\[
 S\equiv
 1-H\frac{r(r-1)}2\sum_{i\in I}\frac1i
 \equiv1\pmod{p^{2v}},
\]
again by \eqref{SP:S:eq:positive-reciprocal-block}.
All prefactors multiplying \(S-1\) are \(p\)-integral,
which proves the lemma.
\end{proof}

\subsection{An integral triangular descent}

The following argument is stated in the notation of the two
families, but uses only integrality of \(B\), its diagonal entries
one, and Lemma~\ref{SP:S:lem:positive-block}.
For each positive integer \(K\), define integers \(c_{K,N}\) by
\begin{equation}\label{SP:S:eq:positive-triangular}
 a_j\mathbf1_{j\geq K}
 =\sum_{N=K}^{j}c_{K,N}B(j,N),\qquad j\geq0,
\end{equation}
and put \(c_{K,N}=0\) for \(N<K\).
Existence, uniqueness and integrality follow by successive solution
in \(j\), since \(B(j,j)=1\).

\begin{lemma}[Weighted interpolation and descent]
\label{SP:S:lem:positive-descent}
For every \(K,N\geq1\),
\begin{equation}\label{SP:S:eq:positive-c-weight}
 v_p(c_{K,N})\geq
 \max\{0,\,2v_p(K)-2v_p(N)\}.
\end{equation}
If \(p\mid K\), then
\begin{equation}\label{SP:S:eq:positive-c-descent}
 c_{K,N}\equiv
 \begin{cases}
 c_{K/p,N/p},&p\mid N,\\
 0,&p\nmid N
 \end{cases}
 \pmod{p^{2v_p(K)}}.
\end{equation}
\end{lemma}

\begin{proof}
Induct on \(v=v_p(K)\).  If \(v=0\), the weighted bound
is integrality and there is no descent assertion.
Suppose \(v\geq1\), set \(L=K/p\), and assume the weighted
bound for \(L\).  Let \(l=\lfloor j/p\rfloor\).
Multiplying the defining identity for \(L\) by \(d_j=a_j/a_l\)
gives
\[
 a_j\mathbf1_{j\geq K}
 =d_ja_l\mathbf1_{l\geq L}
 =\sum_{N=L}^{l}c_{L,N}d_jB(l,N).
\]
By Lemma~\ref{SP:S:lem:positive-block}, replacing each
\(d_jB(l,N)\) by \(B(j,pN)\) makes an error of valuation at
least
\[
 2+2v_p(N)+\max\{0,\,2(v-1)-2v_p(N)\}\geq2v.
\]
Thus the candidate coefficient vector, equal to \(c_{L,N}\)
at index \(pN\) and zero at indices not divisible by \(p\),
satisfies \eqref{SP:S:eq:positive-triangular} modulo \(p^{2v}\).
Successively solving this unit triangular system proves
\eqref{SP:S:eq:positive-c-descent}.
For \(p\nmid N\), it also gives the weighted bound directly.
For \(p\mid N\), the induction hypothesis applied to
\(c_{L,N/p}\) gives the bound
\(\max\{0,2(v-v_p(N))\}\), and the error \(p^{2v}\)
does not weaken it.
\end{proof}

\subsection{Positive vertex coefficients and the excellent lift}

For the cubic family choose the vertex \({\bf v}=(1,\ldots,1)\).
For the simplex choose \({\bf v}=(-1,\ldots,-1)\).
The coefficient at \({\bf x}^{hN{\bf v}}\) in the
\(z\)-adic expansion of \(1/f\) is
\begin{equation}\label{SP:S:eq:positive-omega}
 \Omega_N(z)=\sum_{j\geq N}B(j,N)z^{hj}.
\end{equation}
For the cube this follows by choosing \(j-N\) negative steps in
each of the \(n\) independent factors of \(g^{2j}\).
For the simplex, the required term of \(g^{rj}\) uses the
chosen vertex \(j+(r-1)N\) times and each of the other
\(r-1\) vertices \(j-N\) times, giving exactly the displayed
multinomial coefficient.

Set
\[
 C_N(u)=\frac{\Omega_N(z)}{F(u)}\in u^N\mathbb Z_p[[u]].
\]
The coefficientwise interpretation of the module of second formal
derivatives, including in the positive \(z\)-adic expansion, is
\cite[Section 3, the paragraph following Corollary 3.4]{BV3}.
Read Proposition~\ref{SP:prop:orbit} at the exponent vector
\(hN{\bf v}\).  The left coefficient is \(\Omega_{pN}(z)\),
and the right coefficient is
\(\mathcal F(z)\Omega_N(\tau)/\mathcal F(\tau)\).
The error is divisible by \(p^2(hN)^2\) in \(\mathbb Z_p[[z]]\).
Since \(h\in\mathbb Z\) and \(F(0)=1\), division by \(F(u)\)
gives
\begin{equation}\label{SP:S:eq:positive-vertex-cartier}
 C_{pN}(u)-C_N(\sigma(u))
 \in p^{2+2v_p(N)}\mathbb Z_p[[u]].
\end{equation}

\begin{proof}[Proof of Theorem~\ref{SP:S:thm:positive-cube-simplex}]
Summing \eqref{SP:S:eq:positive-triangular} in \(j\), using
\eqref{SP:S:eq:positive-omega}, and dividing by \(F(u)\), gives
the exact identity
\[
 1-\frac{F_K(u)}{F(u)}
 =\sum_{N\geq K}c_{K,N}C_N(u).
\]
It converges \(u\)-adically because \(C_N\) has order at
least \(N\).
Put \(v=v_p(K)\geq1\) and \(L=K/p\).
Lemma~\ref{SP:S:lem:positive-descent} first gives
\[
 1-\frac{F_K(u)}{F(u)}
 \equiv\sum_{N\geq L}c_{L,N}C_{pN}(u)
 \pmod{p^{2v}}.
\]
By \eqref{SP:S:eq:positive-vertex-cartier} and
\eqref{SP:S:eq:positive-c-weight}, each summand may be replaced by
\(c_{L,N}C_N(\sigma(u))\) with error divisible by \(p^{2v}\).
Indeed, the error has valuation at least
\[
 2+2v_p(N)+\max\{0,2(v-1)-2v_p(N)\}\geq2v.
\]
Both sums still converge: \(C_{pN}(u)\) and
\(C_N(\sigma(u))\) have order at least \(pN\).
Consequently
\[
 1-\frac{F_K(u)}{F(u)}
 \equiv\sum_{N\geq L}c_{L,N}C_N(\sigma(u))
 =1-\frac{F_L(\sigma(u))}{F(\sigma(u))}
 \pmod{p^{2v}},
\]
which is \eqref{SP:S:eq:positive-main}.
Multiplication by the period units gives the cleared-defect and
quotient formulations.  Finally,
\(\mathcal F_{hM}(z)=F_M(z^h)\) proves the assertion in the
original normalization.
\end{proof}


\section{The octahedral first level at every odd prime}\label{SP:source-main-4}
We retain the geometric first-level proof with the new input. It also
applies to the other two families, but only the octahedral conclusion is
new beyond the preceding section. Write
\[
c_j=\CT(g^j),\quad \mathcal F(z)=\sum_{j\ge0}c_jz^j,
\quad H_N(z)=\CT(f^{N-1}),\quad
\lambda=\mathcal F(z)/\mathcal F(\tau(z)).
\]
In this section $\sigma_p(z)=\tau(z)$ and $\sigma_p(u)=\sigma(u)$
denote the two normalizations. Proposition~\ref{SP:prop:orbit} supplies
the excellent Cartier identity at every odd prime.
\begin{lemma}[Second-order geometric approximation]
\label{SP:L:lem:second-geometric}
For the above excellent lift,
\[
 \lambda\equiv2H_p-H_{2p}\pmod{p^2\mathbb Z_p[[z]]}.
\]
More precisely, write
\[
 \mathscr C_p(f^{2p-1})=H_{2p}(z)+B_p(z)g(\mathbf{x}).
\]
Then
\[
 \tau\equiv-\frac{B_p(z)}{H_{2p}(z)}\pmod{p^2\mathbb Z_p[[z]]}.
\]
\end{lemma}

\begin{proof}
All the calculations below may be made in the ring of Laurent series
whose coefficients are \(z\)-adically convergent integral power series.
Put \(h(\mathbf{x})=f^\sigma(\mathbf{x}^p)\).  Since \(h-f^p\) is
divisible by \(p\), the exact identity
\[
 \frac1f-\frac{2hf^{p-1}-f^{2p-1}}{h^2}
 =\frac{(h-f^p)^2}{fh^2}
\]
has an integral remainder divisible by \(p^2\).  Applying the Cartier
operator and using
\(\mathscr C_p(A(\mathbf{x})/h(\mathbf{x})^2)
 =\mathscr C_p(A) / f^\sigma(\mathbf{x})^2\), we obtain
\[
 \mathscr C_p(1/f)
 \equiv
 \frac{2f^\sigma\mathscr C_p(f^{p-1})
              -\mathscr C_p(f^{2p-1})}{(f^\sigma)^2}
 \pmod{p^2}.
\]

We justify the two numerator supports explicitly.  Let \(\Gamma\)
be the lattice spanned by the support of \(g\).  Since its index is
prime to \(p\), \(p\mathbf{u}\in\Gamma\) implies
\(\mathbf{u}\in\Gamma\).  Reflexivity gives integral facet functions
\(L_j\) with \(\Delta=\{L_j\leq1\}\).  A lattice point in
\((1-1/p)\Delta\) has \(L_j\leq0\) for every \(j\), and hence is
zero.  A lattice point in \((2-1/p)\Delta\) has \(L_j\leq1\), and
hence belongs to \(\Delta\).  The only points of
\(\Delta\cap\Gamma\) are zero and the vertices.  Complete symmetry
therefore yields the exact identities
\[
 \mathscr C_p(f^{p-1})=H_p,
 \qquad
 \mathscr C_p(f^{2p-1})=H_{2p}+B_p g.
\]
For the hypercube, the restriction to \(\Gamma\) is essential; its
points have coordinates all even or all odd.

The excellent Cartier relation
\(\mathscr C_p(1/f)\equiv\lambda/f^\sigma
 \pmod{p^2\mathscr F_2^\sigma}\)
is in particular a coefficientwise congruence modulo \(p^2\).
Multiplying by \((f^\sigma)^2\) gives
\[
 \lambda(1-\tau g)
 \equiv 2H_p(1-\tau g)-H_{2p}-B_p g\pmod{p^2}.
\]
Since \(\operatorname{CT}(g)=0\), comparison of the coefficients of
\(1\) and \(g\) first gives
\(\lambda\equiv2H_p-H_{2p}\), and then
\(\tau H_{2p}\equiv-B_p\).
The latter can be divided by the unit \(H_{2p}\), whose constant
term is one.
\end{proof}

\begin{theorem}[The first level at every odd prime]
\label{SP:L:thm:first-all-dim}
Set \(h=2\) for the hypercubic and hyperoctahedral families and
\(h=n+1\) for the simplicial family.  At every odd prime,
\[
 \frac{\mathcal F(z)}{\mathcal F(\sigma_p(z))}
 \equiv\mathcal F_{hp}(z)\pmod{p^2\mathbb Z_p[[z]]}.
\]
Equivalently, in the normalized variable \(u=z^h\),
\[
 \frac{F(u)}{F(\sigma_p(u))}
 \equiv F_p(u)\pmod{p^2\mathbb Z_p[[u]]}.
\]
These are the \(s=1\) cases of the specified original truncation
lengths; \(\mathcal F_h=1\).
\end{theorem}

\begin{proof}
The coefficient of \(z^j\) in \(2H_p-H_{2p}\) is
\[
 c_jw_j,\qquad
 w_j=(-1)^j\left(2\binom{p-1}{j}-\binom{2p-1}{j}\right),
\]
with binomial coefficients zero past the top index.  If \(0\leq j<p\),
the product expansion
\[
 (-1)^j\binom{ap-1}{j}
 =\prod_{i=1}^j\left(1-\frac{ap}{i}\right)
 \equiv1-apH_j\pmod{p^2}
\]
shows that \(w_j\equiv1\pmod{p^2}\).
For \(p\leq j<2p\), Lucas's congruence gives
\(w_j\equiv1\pmod p\).

The remaining required divisibilities follow directly from the
factorial formulas.  In the cubic case,
\(c_{2k}=\binom{2k}{k}^n\), and for \(k<p\) this is divisible by
\(p\) whenever \(2k\geq p\).
In the octahedral case,
\[
 c_{2k}=\sum_{k_1+\cdots+k_n=k}
              \frac{(2k)!}{(k_1!\cdots k_n!)^2},
\]
and every denominator is prime to \(p\) when \(k<p\); hence the
same divisibility holds term by term.
In the simplicial case, putting \(r=n+1\),
\[
 c_{rk}=\frac{(rk)!}{(k!)^r},\qquad 0\leq k<p,
 \qquad v_p(c_{rk})\geq\left\lfloor\frac{rk}{p}\right\rfloor.
\]
Thus in all three cases \(c_j\) is divisible by \(p\) for
\(p\leq j<\min(2p,hp)\), and in the simplicial case it is divisible
by \(p^2\) for \(2p\leq j<hp\).  The other two families have
\(h=2\), so there is no latter range.
It follows coefficientwise that
\[
 2H_p-H_{2p}\equiv\mathcal F_{hp}\pmod{p^2}.
\]
Apply Lemma~\ref{SP:L:lem:second-geometric} and then the normalized
coordinate relation \(q_0(z)^h=q(z^h)\).
\end{proof}


\section{Three exact whole-series certificates}\label{SP:source-main-5}
Here $g=\sum_{i=1}^n(x_i+x_i^{-1})$, $u=z^2$ and
\[
a_j=\binom{2j}{j}
\sum_{k_1+\cdots+k_n=j}\binom{j}{k_1,\ldots,k_n}^{\!2}.
\]
We use $(n,p,s)=(4,3,2),(4,3,3),(5,5,2)$ and normalized multiplier
$m=1$. In the last case Proposition~\ref{SP:prop:orbit} supplies the
Cartier input and its rank-two differential quotient, including the
integral Wronskian, despite $p\mid n$. For the first two cases the
smaller group $C_8$ also suffices for the published BV input.
The following argument is common to the three certificates.
\subsection{Vertex polynomials and their derivative estimate}
Define
\[
 b_j=\sum_{k_2+\cdots+k_n=j}\binom{j}{k_2,\ldots,k_n}^{\!2},\qquad
 A_K(u)=\sum_{j=0}^{K-1}u^j\sum_{\ell=0}^j(-1)^\ell
 \frac{(2K-\ell-1)!b_{j-\ell}}
 {(2K-2j-1)!\ell!(j-\ell)!^2}.
\]
The factorial ratio is an integer multinomial coefficient and $A_K(0)=1$.
Expanding $1/f$ in the chamber with dominant vertex $x_1$, the negative
of the coefficient of $x_1^{-2K}$, all other exponents being zero, is
\[
 P_{2K}(u)=u^{-K}A_K(u).
\]
Let $\lambda=F(u)/F(\sigma(u))$. Reading this vertex coefficient in
the excellent Cartier congruence gives
\begin{equation}\label{SP:eq:vertex}
 P_{2pK}(u)-\lambda P_{2K}(\sigma(u))
       \in p^{2+2\vp(K)}\Z_p((u)).
\end{equation}
The displayed elements have bounded negative support. No claim that every
element of the completed coefficient ring is a Laurent series with bounded
negative support is needed.

\begin{lemma}\label{SP:lem:derivative}
For $C_K=A_{2K}/A_K$ and $\theta=u\,d/du$,
\[
 \theta C_K\in p^{\vp(K)}\Z_p[[u]].
\]
\end{lemma}
\begin{proof}
Set $R_K=P_{2K}/F$ and use the normalized canonical coordinate $q$, with
$\partial=q\,d/dq$. The second-order quotient relation
\cite[Proposition 7.6]{BV3}, read at the same vertex, gives after division
by $F$ and passage to $q$
\[
 \partial^2R_K\in p^{2\vp(K)}\Z_p((q)).
\]
More explicitly, in normalized coordinates let
$\mathcal L=\theta^2-B_u\theta-A_u$ and $H=\theta q/q=W/F^2$.
Then $\theta H/H=B_u-2\theta F/F$ and the exact identity
\[
\mathcal L(FR)=FH^2\partial^2R
\]
shows that this passage preserves the required precision.
The gauge and coordinate factors are integral units by
\cite[Proposition 7.9 and Corollary 7.11]{BV3}; the normalization factor
$2$ is a $p$-adic unit. For $r=\vp(K)$, integrality of $R_K$ together
with $j^2[q^j]R_K\in p^{2r}\Z_p$ implies
$j[q^j]R_K\in p^r\Z_p$ for every nonzero integer $j$.
Thus $\partial R_K$, and then $\theta R_K$, belong to
$p^r\Z_p((q))$ and $p^r\Z_p((u))$, respectively.
Differentiate $C_K=u^K R_{2K}/R_K$; the factors after removing their
displayed powers of $u$ are integral units. Every term has the factor
$p^r$, giving the assertion.
\end{proof}

\subsection{The scaled Hensel equation}
Fix one of the three triples $(n,p,s)$ in Theorem~\ref{SP:thm:scope}, and set
\[
 M=p^{2s},\quad M_y=p^{2s-1},\quad N=p^{2s-2},\quad X=u^p,
 \qquad \sigma(u)=X(1+py(u)),\quad y\in\Z_p[[u]].
\]
Eliminating $\lambda$ between \eqref{SP:eq:vertex} for $K=N,2N$ gives
\[
 (1+py)^N C_{pN}(u)-C_N(X(1+py))
       \in p^{4s-2}\Z_p[[u]].
\]
Divide this expression by $pN$. Its coefficient of $y^r$, for $r\ge1$,
has valuation at least
\begin{equation}\label{SP:eq:coefficient-valuation}
 r-1-\vp(r).
\end{equation}
Indeed Lemma~\ref{SP:lem:derivative} implies that the $r$th Taylor
coefficient of $C_N(X(1+h))$ lies in $(N/r)\Z_p[[u]]$;
also $\binom Nr\in(N/r)\Z_p$. Multiplication by $p^r/(pN)$ gives
\eqref{SP:eq:coefficient-valuation}. The constant coefficient is integral
as well, since the actual $y$ solves the preceding congruence and all
positive-degree coefficients are integral.

Take $d=3$ for the two level-two cases, and $d=6$ for $(n,p,s)=(4,3,3)$.
The bound \eqref{SP:eq:coefficient-valuation} is at least $2s-1$ for every
$r>d$. Thus the equation for $y$ modulo $M_y$ is a degree-$d$ polynomial.
Its linear coefficient has constant term $1$, and all nonlinear
coefficients are divisible by $p$.

Here are completely explicit common-denominator numerators for that
polynomial. Compute the four vertex polynomials at indices
$N,2N,pN,2pN$ modulo $p^{4s-2}$ and set
\[
 a_r=[h^r]A_N(X(1+h)),\quad b_r=[h^r]A_{2N}(X(1+h)),
 \quad a_0=A_N(X).
\]
Define
\[
 T_0=b_0,\qquad
 T_r=b_ra_0^r-\sum_{i=1}^r a_iT_{r-i}a_0^{i-1}.
\]
Then $T_r/a_0^{r+1}=[h^r]C_N(X(1+h))$. For $0\le r\le d$ put
\begin{equation}\label{SP:eq:Qs}
 Q_r=\frac{p^r}{pN}
 \left(\binom Nr A_{2pN}a_0^{d+1}
             -A_{pN}T_ra_0^{d-r}\right)\pmod{M_y}.
\end{equation}
The quotients are coefficientwise integral by the preceding argument;
the certificate independently checks this exact divisibility.
Their common denominator $A_{pN}a_0^{d+1}$ is a unit. Thus
\begin{equation}\label{SP:eq:Hensel}
 \sum_{r=0}^d Q_r y^r=0\pmod{M_y},\qquad
 Q_0(0)=0,\quad Q_1(0)=1,\quad Q_r\in p\Z_p[u]\ (r\ge2).
\end{equation}
Two roots of this equation have difference multiplied by a unit,
because after subtracting their equations the factor is congruent to
$Q_1$ modulo $p$. Therefore the root modulo $M_y$ is unique.

The program constructs a rational root $Y/Z$ from these $Q_r$ using
the finite multivariate Lagrange expansion, with $Z(0)=1$, and then
checks the entire polynomial identity
\begin{equation}\label{SP:eq:root-check}
 \sum_{r=0}^d Q_rY^r Z^{d-r}=0\pmod{M_y}.
\end{equation}
This identity and uniqueness, rather than an assumption about the
target truncations, prove $y=Y/Z\pmod{M_y}$ and hence
$\sigma=X(1+pY/Z)\pmod M$.

For clarity, the construction is explicit. Write $A=Q_0$, $B=Q_1$.
For a profile $(m_2,\ldots,m_d)$ of nonnegative integers, put
$m=\sum m_r\ge1$, $R=\sum r m_r$ and $\ell=R-m+1$. Its contribution to
$Y/Z$ is
\[
 (-1)^{R+1}\frac{R!}{\ell!\prod m_r!}
 \frac{A^\ell\prod_{r=2}^d Q_r^{m_r}}{B^{R+1}}.
\]
Include also the term $-A/B$, and retain only profiles of total
coefficient valuation less than $2s-1$.
The integer factorial quotient is the usual rooted-plane-tree
coefficient in Lagrange inversion. The code checks its integrality.
The verified common denominators are $B^7$, $B^{13}$, and $B^5$ for $(n,p,s)=(4,3,2),(4,3,3),(5,5,2)$ respectively.
The direct check \eqref{SP:eq:root-check} makes the root certification
independent of any truncation of the Lagrange expansion.

\subsection{The target polynomial identities}
The vertex relation at $K=N$ yields
\[
 \lambda=\frac{(1+py)^N A_{pN}(u)}{A_N(X(1+py))}\pmod M.
\]
Since $\vp\binom Nr\ge\vp(N)-\vp(r)$, for every $r\ge2$ the term
$\binom Nr(py)^r$ is divisible by $M$. Hence
$(1+h)^N=1+Nh\pmod M$ on $h\in p\Z_p[[u]]$.

Set $f=F_{p^s}$, $f_0=F_{p^{s-1}}$ and, for $0\le r<2s$, define
\[
 L_r=A_{pN}(u)\bigl([h^r]f_0(X(1+h))
               +N[h^{r-1}]f_0(X(1+h))\bigr)
          -f(u)[h^r]A_N(X(1+h)),
\]
where a coefficient with negative index is zero.
The target congruence is equivalent to
$\sum_{r<2s}L_rh^r=0\pmod M$.
Substitute the already certified $h=pY/Z$ and clear the unit denominator:
\begin{equation}\label{SP:eq:target-check}
 \boxed{\ \sum_{r=0}^{2s-1}L_r(pY)^r Z^{2s-1-r}=0\pmod{p^{2s}}.\ }
\end{equation}
Only at this last step do the ordinary period truncations enter the
calculation. The program verifies every coefficient of
\eqref{SP:eq:root-check} and \eqref{SP:eq:target-check} exactly.

\begin{center}\small
\begin{tabular}{rrrrrrr}
\toprule
$n$&$p$&$s$&$N$&working modulus&$(\deg Y,\deg Z)$&target bound\\
\midrule
4&3&2&9&729&(521,520)&1562\\
4&3&3&81&59049&(14297,14296)&70519\\
5&5&2&25&15625&(301,300)&506\\
\bottomrule
\end{tabular}
\end{center}
Both polynomial identities are zero. Multiplying the verified
$\lambda F_{p^{s-1}}(\sigma)-F_{p^s}$ congruence by the unit
$F(\sigma)$ proves the certificate assertions of Theorem~\ref{SP:thm:scope}.

\section{A sharp obstruction for the uncorrected period functional}\label{SP:source-main-6}
\label{SP:sec:period-obstruction-small-primes}

This section concerns the transfer functional used in \cite[Section~4]{BV2}, rather than a counterexample to an
excellent-lift truncation congruence.  All constant terms below are taken
in the positive $t$-adic expansion.  Let $p$ be odd, let $s\geq1$, and put
$q=p^s$ and $M=2q$.  For the $n$-dimensional octahedral polynomial set
\[
 g_n({\bf x})=\sum_{i=1}^n(x_i+x_i^{-1}),\qquad
 f=1-tg_n,\qquad
 B_n(j,1)=[x_1^2]g_n({\bf x})^{2j},
\]
where the coefficient at every other variable is zero.

\begin{lemma}[A coefficient of exact first-order precision]
\label{SP:lem:octa-near-central-sharp}
For every $n\geq1$, every odd prime $p$, and every $s\geq1$,
\[
 B_n(p^s,1)\equiv 2(2-n)p^s\pmod{p^{s+1}}.
\]
In particular, if $p\nmid n-2$, then $v_p(B_n(p^s,1))=s$.
\end{lemma}

\begin{proof}
For a composition $k_1+\cdots+k_n=q$ with $k_i\geq0$ write
$V({\bf k})=q!/(k_1!\cdots k_n!)$.  Pairing the positive and negative
steps in each coordinate gives
\[
 B_n(q,1)=\binom{2q}{q}
 \sum_{\substack{k_1+\cdots+k_n=q\\ k_1\geq1}}
 V({\bf k})^2\frac{k_1}{k_1+1}.
\]
The identity
\[
 k_1 V({\bf k})
 =q\frac{(q-1)!}{(k_1-1)!k_2!\cdots k_n!}
\]
implies $v_p(V({\bf k}))\geq s-v_p(k_1)$.  Thus each summand has
valuation at least
\[
 2s-v_p(k_1)-v_p(k_1+1).
\]
If $1\leq k_1\leq q-2$, at most one of the consecutive integers $k_1$,
$k_1+1$ is divisible by $p$, and its valuation is at most $s-1$.
Consequently all such summands vanish modulo $p^{s+1}$.
The summand with $k_1=q$ equals $q/(q+1)$.
For $k_1=q-1$ there are $n-1$ compositions, each contributing
$q(q-1)$.  Hence
\[
 \sum V({\bf k})^2\frac{k_1}{k_1+1}
 \equiv\frac q{q+1}+(n-1)q(q-1)
 \equiv(2-n)q\pmod{p^{s+1}}.
\]
Finally, Lucas's congruence gives $\binom{2q}{q}\equiv2\pmod p$,
which proves the claim.
\end{proof}

Define the period functional and its Cartier compatibility defect by
\[
 \rho_{M,g_n}(\omega)=\operatorname{CT}(g_n^M\omega),\qquad
 \mathcal D_M(\omega)=\rho_{M,g_n}(\omega)
       -\rho_{M/p,g_n}(\mathscr C_p\omega).
\]
No coefficient Frobenius changes $g_n$, since its coefficients are integers.
The Cartier operator on positive $t$-adic expansions keeps the coefficients
of ${\bf x}^{p\mathbf u}$ and divides their exponent vectors by $p$.

\begin{proposition}[Sharpness of the general transfer precision]
\label{SP:prop:period-functional-obstruction}
There is a fixed second formal derivative $\eta\in\mathscr F_2$,
independent of $s$, for which
\[
 [t^3]\rho_{M,g_n}(\eta)
 =[t^3]\mathcal D_M(\eta)=4B_n(p^s,1).
\]
Thus, if $p\nmid n-2$, both displayed coefficients have valuation exactly
$s$.  Moreover there is a fixed signed-permutation-invariant
$\eta_{\rm sym}\in\mathscr F_2$ for which
\[
 [t^3]\rho_{M,g_n}(\eta_{\rm sym})
 =[t^3]\mathcal D_M(\eta_{\rm sym})=8nB_n(p^s,1).
\]
If $p\nmid n-2$, the latter valuation is exactly $s+v_p(n)$.
\end{proposition}

\begin{proof}
Put $\theta_i=x_i\partial/\partial x_i$ and choose
\[
 \eta=\theta_1^2\left(\frac{t^3x_1^{-2}}{f^3}\right),\qquad
 \eta_{\rm sym}=\sum_{i=1}^n\theta_i^2
       \left(\frac{t^3(x_i^2+x_i^{-2})}{f^3}\right).
\]
These elements lie in the integral Dwork module at odd $p$: the
numerator exponents $\pm2e_i$ lie in the interior of $3\Delta$, their
$t$-adic orders satisfy admissibility, and the pole-order-three
normalization factor $2!=2$ is a $p$-adic unit.  Applying Euler
derivations preserves this module.  Being second derivatives, the
displayed elements belong to $\mathscr F_2$; equivalently their
coefficients have the required squared exponent divisibility.
The symmetric element is fixed by every coordinate permutation and
coordinate inversion.  Its construction uses a sum, with no division
by the group order.

Their initial positive expansions are
\[
 \eta=4t^3x_1^{-2}+O(t^4),\qquad
 \eta_{\rm sym}=4t^3\sum_i(x_i^2+x_i^{-2})+O(t^4).
\]
Since $p$ is odd, none of the exponent vectors $\pm2e_i$ is divisible
by $p$.  Therefore both Cartier images are $O(t^4)$.  Taking constant
terms after multiplying by $g_n^M$, and using signed-permutation
symmetry of $g_n^M$, proves the coefficient identities.  The valuation
claims now follow from Lemma~\ref{SP:lem:octa-near-central-sharp}.
\end{proof}

The same obstruction applies to the truncation kernel $h=tg_n$.
For a lift $\sigma(t)=t^pv(t)$ with $v(0)$ a unit, replace the defect by
\[
 \rho_{M,h}(\omega)-\rho_{M/p,h^\sigma}(\mathscr C_p\omega).
\]
The first term has initial coefficient $4B_n(p^s,1)$ (respectively
$8nB_n(p^s,1)$) at $t^{M+3}$, whereas the second has order at least
$M+4$.  Thus the same exact valuation holds.

For $(n,p)=(4,3)$, even the symmetric example has valuation exactly
$s$, so the raw period-functional Cartier compatibility cannot be
uniformly upgraded to $p^{s+1}$ on all second formal derivatives.
For $(n,p)=(5,5)$ the symmetric example has valuation $s+1$, which is
less than $2s$ for $s\geq2$.  If a factor $p^2$ is prefixed to the
test element, the corresponding valuation is $s+v_p(n)+2$;
this still rules out a uniform $p^{2s}$ conclusion for all sufficiently
large $s$ from membership in $p^2\mathscr F_2$ alone.

These are operator-level obstructions, not failures of the original
supercongruence.  The actual excellent Cartier error is a particular
element of $p^2\mathscr F_2$ and could possess additional cancellation.
The calculation shows precisely why such additional information, or a
different transfer functional, is needed.  It does not prove that no
corrected period method can work.

\paragraph{Source of the transfer statement.}
Beukers--Vlasenko, \emph{Dwork crystals II}, arXiv:1907.10390v3,
Section~4, Proposition~4.1, states the period-functional compatibility
modulo $p^{v_p(M)}$.  Corollary~4.4 applies it to truncation matrices.
The propositions above exhibit an explicit sharpness mechanism for
this general functional estimate in the octahedral setting.

\paragraph{What the excellent error alone contributes.}
Suppose $\mathscr C_p(\omega)=\lambda\omega^\sigma+p^2\epsilon$
with $\epsilon\in\mathscr F_2$.  Put $h=tg_n$ and $m=M/p$.
Then the exact transfer decomposition is
\[
 \rho_{M,h}(\omega)-\lambda\rho_{m,h^\sigma}(\omega^\sigma)
 =\bigl(\rho_{M,h}(\omega)-\rho_{m,h^\sigma}(\mathscr C_p\omega)\bigr)
   +p^2\rho_{m,h^\sigma}(\epsilon).
\]
The standard period-modulo-$m$ estimate gives the second summand
divisibility by $p^{s+1}$, since $\mathscr F_2\subseteq\mathscr F_1$
and $v_p(m)=s-1$.  The first summand is only guaranteed to be divisible
by $p^s$ by BV2, Proposition~4.1.  Thus merely inserting the excellent
Cartier identity into that general proof does not prove a gain beyond
$p^s$.  The exact coefficient examples above show why a uniform upgrade
of the first functional estimate on $\mathscr F_2$ fails.  They do not
exclude cancellations for the distinguished element $\omega=1/f$.

\section{The interpolation obstruction and its subsequent resolution}\label{SP:source-main-7}
The original matrix proof in Part~\ref{part:positive} uses $p>n$
in its symmetric Bessel reduction. The support-weighted argument in
Part~\ref{part:octa-small} removes that restriction at $p\ge5$.
At $p=3$ the normalized leading correction is
$1-3(X^3-X)\pmod9$, which is pointwise $1\pmod9$ on integers
but not coefficientwise $1\pmod9$. The dimension-polynomial
argument in Part~\ref{part:octa-three} resolves this remaining
obstruction at every level and every normalized multiplier.

The three certificates in this part bypass these steps at the stated
parameter choices. Merely having an algorithm that constructs a
finite target polynomial at each level would not prove its vanishing
at unexecuted levels. Their role as individual whole-series proofs
is unchanged by the later algebraic theorem.

One elementary source of potential compensation at primes dividing the
dimension is the exact identity, for $j\ge1$,
\[
a_j=2n\,\CT(x_1g^{2j-1})\in2n\mathbb Z.
\]
Indeed the $2n$ monomials of $g$ contribute equal constant terms by
signed-permutation symmetry. This does not by itself control the zero
boundary index or the shifted $N^2$ errors in the row construction.

\section{Verification and relation to the earlier manuscript}\label{SP:source-main-8}
This part incorporates the first small-prime supplement into the
complete edition. It supplies the all-odd-prime cube and simplex
results, the common formal orbit-module input, and independent
octahedral first-level and certificate arguments. Subsequent parts
supply the full octahedral descent.

The new orbit-module and scalar results are symbolic arguments.
The three specified octahedral levels are computer-assisted proofs:
their reduction to finite polynomial identities is proved above, and
the complete exact verifier is printed below. Finite power-series
checks used during exploration are not premises of these theorems.
The results remain subject to independent mathematical review; no
historical priority or formal proof-assistant certification is claimed.

For reproduction, use Python 3 and execute
\begin{verbatim}
python3 verify_small_prime_certificates.py --dimension 4 --p 3 --s 2
python3 verify_small_prime_certificates.py --dimension 4 --p 3 --s 3
python3 verify_small_prime_certificates.py --dimension 5 --p 5 --s 2
\end{verbatim}
The program writes each coefficient payload and a report. An assertion
failure invalidates that run; a successful run verifies both the lift
root identity and the cleared target identity. The embedded source regenerates the three full payloads and their
reports; the independently archived payloads are not needed to
compile this single-source edition. No third-party Python packages are needed.



% END RETAINED SOURCE SP
\partstart{Hyperoctahedral descent in arbitrary dimension}{part:octa-small}{All primes at least five, the mixed-precision proof at three, and independent certificates through level four.}
% BEGIN RETAINED SOURCE OP
\section{Arbitrary-dimensional descent and independent finite-level results}\label{OP:source-main-1}\label{OP:sec:scope}
Put
\[
 g_n(\mathbf x)=\sum_{i=1}^n(x_i+x_i^{-1}),\quad f=1-zg_n,
 \quad u=z^2,
\]
\[
 F_n(u)=\sum_{j\ge0}a_{n,j}u^j,\qquad
 a_{n,j}=\frac{(2j)!}{(j!)^2}
       \sum_{k_1+\cdots+k_n=j}\binom{j}{k_1,\ldots,k_n}^{\!2},
 \qquad F_{n,K}(u)=\sum_{j<K}a_{n,j}u^j.
\]
Let $\sigma_n$ be the normalized excellent lift of the canonical
coordinate specified in Section~\ref{OP:sec:formal-input}. Define
\[
 E_{n,p,s,m}(u)=F_n(u)F_{n,mp^{s-1}}(\sigma_n(u))
             -F_n(\sigma_n(u))F_{n,mp^s}(u).
\]
All congruences in this paper are coefficientwise. The normalized
length $mp^s$ corresponds to the original length $2mp^s$ in the
variable $z$; normalized $m=1$ therefore corresponds to the multiplier
$2$ specified for octahedra in BV3, Conjecture~7.5.

\begin{theorem}[All-level small-prime descent]\label{OP:thm:all-level}
For every dimension $n\ge3$ and prime $p$ satisfying either
\[
 p\ge5\qquad\hbox{or}\qquad(p=3\ \hbox{and}\ 3\mid n),
\]
and every positive integer $K$ divisible by $p$, one has
\[
 \frac{F_{n,K}(u)}{F_n(u)}
 \equiv\frac{F_{n,K/p}(\sigma_n(u))}{F_n(\sigma_n(u))}
       \pmod{p^{2v_p(K)}\mathbb Z_p[[u]]}.
\]
In particular $E_{n,p,s,m}\in p^{2s}\mathbb Z_p[[u]]$ for every
$s,m\ge1$, with no restriction on the coefficient index.
\end{theorem}

\begin{theorem}[Dimension-uniform remaining levels]\label{OP:thm:finite-scope}
For every $n\ge3$, at $p=3$ and normalized multiplier $m=1$,
\[
 E_{n,3,s,1}\in3^{2s}\mathbb Z_3[[u]]\qquad(s=1,2,3,4).
\]
The first level is proved algebraically in Section~\ref{OP:sec:first-level}.
The three higher levels are proved by finite exact polynomial identities
and the dimension argument of Sections~\ref{OP:sec:dimension-newton} and~\ref{OP:sec:dimension-newton-level4}.
They do not constitute an induction over all $s$.
\end{theorem}

The following table records the scope of the arguments in this part;
the global scope is Theorem~\ref{G:thm:global}.
\begin{center}\small
\begin{tabular}{c|p{57mm}|p{48mm}}
\toprule
Dimension & Prime $3$ & Every prime $p\ge5$\\\midrule
$3,6,9,12,\ldots$ & Every $s,m$; every coefficient & Every $s,m$; every coefficient\\
$4,5,7,8,10,11,\ldots$ & $m=1$, levels $s=1,2,3,4$; every coefficient & Every $s,m$; every coefficient\\
\bottomrule
\end{tabular}
\end{center}
The two all-level methods of this part leave the sector
$p=3$, $n\ge4$, $3\nmid n$ untreated. The finite-level
theorem supplies normalized $m=1$ through $s=4$. The complete
all-level, all-multiplier statement in that sector is proved in
Part~\ref{part:octa-three}; it is not inferred from these tests.
Dimensions one and two already have all-odd-prime scalar
proofs, since they agree with the corresponding cubic positive vertex
matrices. This paper supplies new descent arguments beyond the prime
range in the published BV3 theorem; its formal input is checked using
the general machinery of that paper, not by silently removing a
hypothesis from its Theorem~7.3.
\section{The formal excellent Cartier input}\label{OP:source-main-2}\label{OP:sec:formal-input}
\begin{proposition}[Formal excellent input from an orbit module]\label{OP:prop:orbit}
Let $g=g_n$ be the hyperoctahedral Laurent polynomial and let $p$ be an odd prime.
Put $R=\mathbb Z_p[[z]]$ and $f=1-zg$.
Let $F_0(z)$ and $G_0(z)$ be the period and logarithmic companion
given by the vertex-coefficient formulas of BV3, Proposition~9.1.
Then
\[
 q_0(z)=z\exp(G_0(z)/F_0(z))\in z+z^2\mathbb Z_p[[z]].
\]
The lift $\tau=q_0^{-1}(q_0(z)^p)$ is integral and satisfies
\[
 \mathscr C_p(1/f)-\frac{F_0(z)}{F_0(\tau)}\frac1{f^\tau}
 \in p^2\mathscr F_2^\tau.
\]
In particular, no exclusion $p\mid n$ is needed for this formal input.
\end{proposition}

\begin{proof}
Write $\Gamma$ for the support lattice and $H$ for the full monomial
symmetry group. Its order will never be inverted. The support lattice is $\mathbb Z^n$, so it is $p$-saturated.
The symmetry group preserves the reflexive polytope degree
$d(u)=\deg_\Delta(u)$ and the support lattice. For an $H$-orbit
$O\subset\Gamma$, put
\[
 S_O=\sum_{u\in O}x^u,\qquad d(O)=d(u)\ (u\in O),\qquad
 L=\bigoplus_O R z^{d(O)}S_O.
\]
Here a Laurent polynomial uses only finitely many orbits.
This is a direct construction of a module of polynomials. We do not
obtain it by applying an averaging projector to another crystal.

Fix temporarily any lift $\sigma$ with
$z^\sigma=z^pv$, $v\in1+pR$. We check BV3, Definition~2.1.
The orbit sums have disjoint supports and coefficient one, so $L$ is
$p$-saturated. The same is true of every $L^{\sigma^i}$: its orbit
coefficient generator is $(z^{\sigma^i})^{d(O)}$, a power of $z$
times a unit, which is a non-zero-divisor modulo $p$.
Multiplication by $f$ preserves invariance and admissibility, hence
$fL\subset L$.

The subset $p\Gamma$ is $H$-stable. An orbit is either disjoint from
$p\Gamma$ or has the form $pO'$ for one orbit $O'$; the equality uses
the $p$-saturation of the support lattice. Since $d(pu)=pd(u)$,
\[
 \mathscr C_p(z^{d(pO')}S_{pO'})
   =z^{pd(O')}S_{O'}
   =\left(\frac{z^p}{z^\sigma}\right)^{d(O')}
        (z^\sigma)^{d(O')}S_{O'}\in L^\sigma.
\]
The image of any other orbit is zero. This proves compatibility
without dividing by either $|H|$ or $|O|$.

Form the Dwork crystal
$M=\widehat\Omega_{L,f}(\Delta^\circ)$ directly, as in BV3,
Definitions~2.6 and~2.8. We do not identify it with the invariant
part of an ambient completed module; in particular no exactness of
an invariant functor is being assumed. In this family,
$\Delta\cap\Gamma$ consists of the origin and the vertices.
Reflexivity then implies that the first two levels have bases
\[
 M(1)=R/f,\qquad M(2)=R/f^2\oplus Rzg/f^2.
\]
Indeed $2\Delta^\circ\cap\Gamma=\Delta\cap\Gamma$ and the
vertices form one orbit. The pair $1/f,zg/f^2$ is an extended basis,
since $1/f=(1-zg)/f^2$ and the change of basis has determinant one.
These levels are $p$-saturated in the ambient integral formal
expansion: multiplying an expansion by $f$ or $f^2$ reduces the
claim to divisibility of the numerator coefficients. This also
shows that the two basis vectors remain independent modulo $p$.

The Hasse--Witt computations of BV3, Lemma~7.2, now apply to this
explicit compatible module. Their coefficient extraction uses a
single vertex coefficient of a polynomial known to be constant on
the vertex orbit; it does not average over that orbit. With the
initial lift $z^\sigma=z^p$, the first determinant has constant
term one, and
\[
 \left.\frac{\det HW^{(2)}(M)}p\right|_{z=0}
 =-\frac1p\binom{2p-2}{p-1}\in\mathbb Z_p^\times.
\]
The latter computation reads the coefficient of $z^px^{pv}$ at one
vertex $v$ in $(zg)F^{(2)}$; the extreme vertex forces the unique
contribution $(-z)^p\binom{2p-2}{p-1}$. The normalized determinant
modulo $p$ is independent of the lift. Consequently both
Hasse--Witt conditions hold over $R$ for every lift under consideration.
The general BV3, Corollary~5.9 and Corollary~4.3, therefore give
\[
 M=M(2)\oplus\mathscr F_2M,
 \qquad
 \mathscr C_p(M)\subset M^\sigma(2)+p^2\mathscr F_2M^\sigma.
\]
Only $p>2$ is required by these general results.

For completeness, the remaining canonical-coordinate argument
uses these statements rather than a group-order hypothesis.
The derivation $\theta=z\,d/dz$ preserves $M$. Thus the rank-two
quotient has the connection and Cartier matrices of BV3,
Propositions~7.6 and~7.7. The positive vertex coefficient calculation
of Proposition~9.1 gives exactly $F_0$ and
$F_0\log z+G_0$ as its two solutions. The same coefficient
calculation at exponents $p^sv$ gives $A(0)=B(0)=0$, the
constant Cartier matrix $\operatorname{diag}(1,p)$ for the
initial lift, and the upper-right coefficient
\[
 \lambda_1=
 \frac{pF_0(z)F_0(z^\sigma)}{W(z^\sigma)}
 \left(-\frac1p\log\frac{z^\sigma}{z^p}
 +\frac{G_0(z)}{F_0(z)}
 -\frac1p\frac{G_0(z^\sigma)}{F_0(z^\sigma)}\right),
\]
where $W=F_0^2+F_0\theta G_0-(\theta F_0)G_0$.
The rank-one reduction modulo $p$ implies $\lambda_1\in pR$.
The determinant identity
$\det\Lambda=pW/W^\sigma$ and $W(0)=1$ give $W\in1+zR$
by the same iteration as BV3, Proposition~7.9.
For $z^\sigma=z^p$, the displayed formula and the
Dieudonne--Dwork lemma yield $q_0\in z+z^2R$.
Its formal inverse is integral, so
$\tau=q_0^{-1}(q_0^p)$ defines a lift with
$\tau/z^p\in1+pR$. Reapply the compatible-module construction
with this lift. The displayed expression has $\lambda_1=0$,
and the remaining coefficient is $F_0/F_0^\tau$.
This proves the required excellent Cartier identity in
$p^2\mathscr F_2M^\tau$.

Finally, membership in $\mathscr F_2M$ means that every $s$-fold
Cartier iterate lies in $p^{2s}M^{\sigma^s}$. Mapping the directly
constructed module to the ambient Dwork module therefore sends
$\mathscr F_2M$ into the ambient $\mathscr F_2$.
Only this inclusion is needed; no converse assertion about
invariants or completed modules is used.
\end{proof}


The sign symmetry gives $q_0(z)=zQ(z^2)$ and
$\tau(-z)=-\tau(z)$. Thus $\tau=z^pV(u)$ with
$V\in1+p\mathbb Z_p[[u]]$, and
\[
 \sigma_n(u)=\tau(z)^2=u^pV(u)^2.
\]
No division by the dimension occurs in this normalization. The
corresponding canonical coordinate in $u$ is $q(u)=q_0(z)^2$.

\section{Positive vertices and the exact factorial block}\label{OP:source-main-3}
Suppress the subscript $n$ when it is fixed. For $0\le N\le j$ put
\[
 B(j,N)=(2j)!\sum_{\substack{k_1+\cdots+k_n=j\\k_1\ge N}}
 \frac{1}{(k_1-N)!(k_1+N)!\prod_{i>1}(k_i!)^2},
\]
and set $B(j,N)=0$ if $N>j$. These are integers, $B(j,j)=1$, and
$B(j,0)=a_j$. Let
\[
 \Omega_N(u)=\sum_{j\ge N}B(j,N)u^j,\qquad
 C_N=\Omega_N/F\in u^N\mathbb Z_p[[u]].
\]
The series $\Omega_N$ is the coefficient at $x_1^{2N}$, with all
other exponents zero, in the positive expansion of $1/f$. Reading
Proposition~\ref{OP:prop:orbit} at that exponent and using the
coefficient characterization of second formal derivatives gives
\begin{equation}\label{OP:eq:positive-cartier}
 C_{pN}-C_N\circ\sigma_n
       \in p^{2+2v_p(N)}\mathbb Z_p[[u]].
\end{equation}
Here the exponent factor is $(2N)^2$ and $2$ is a $p$-adic unit.

The unit triangular system uniquely defines integers $c_{K,N}$ by
\[
 a_j\mathbf1_{j\ge K}=\sum_{N=K}^jc_{K,N}B(j,N),
 \qquad c_{K,N}=0\quad(N<K).
\]
Consequently
\begin{equation}\label{OP:eq:exact-tail}
 1-F_K/F=\sum_{N\ge K}c_{K,N}C_N.
\end{equation}
This identity converges $u$-adically. It is the connection between
Cartier's coefficient congruence and ordinary truncations.

Use the Bessel series and the factorial coefficient lattice
\[
 I_N(x)=\sum_{k\ge N}\frac{x^k}{(k-N)!(k+N)!},\quad
 \Phi_N=I_NI_0^{n-1},\quad \theta=x\frac d{dx},\quad
 \mathcal H=\left\{\sum h_jx^j:(2j)!h_j\in\mathbb Z_p\right\}.
\]
Then $(\theta^2-N^2)I_N=xI_N$ and
$(2j)![x^j]\Phi_N=B(j,N)$. The lattice is complete for the factorial
coefficient norm; multiplication, $x$, and $\theta$ preserve it.
This follows from the integrality of the relevant multinomial factors.
For $f_1=I_N$, $f_i=I_0$ for $i>1$, put
\[
 E_l=\sum_{|S|=l}\prod_{i\in S}\theta f_i\prod_{i\notin S}f_i.
\]
Then $E_0=\Phi_N$, $E_1=\theta\Phi_N$, and
\[
 \theta E_l=(l+1)E_{l+1}+(n-l+1)xE_{l-1}
                    \pmod{N^2\mathcal H}.
\]
The omitted term comes from differentiating the first factor and is
$N^2$ times an integral product of the displayed factors.

\subsection{A balanced unit-factorial estimate at every odd prime}
Write $U(a)=\prod_{1\le i\le a,\,p\nmid i}i$. For $H=pN$ and $k\ge H$,
\[
 \frac{U(k)^2}{U(k-H)U(k+H)}
 =\prod_{\substack{k-H<i\le k\\p\nmid i}}(1+H/i)^{-1}
 \equiv1\pmod{p^{2+2v_p(N)}}.
\]
Indeed, for $v=v_p(H)$ the unit residues in this interval are complete
copies modulo $p^v$. Their reciprocals sum to zero modulo $p^v$ by
pairing opposite residues. Expanding the product to first order proves
the estimate.
\subsection{An integral analytic unit-factorial expansion}

Write \(\mathbb Z_p\langle X_1,\ldots,X_d\rangle\) for the
restricted power series ring: its coefficients tend to zero
\(p\)-adically as their total degree tends to infinity.
For a nonnegative integer \(a\), put
\[
 U(a)=\prod_{\substack{1\leq i\leq a\\p\nmid i}}i,
 \qquad c=(p-1)!.
\]

\begin{lemma}[Unit-factorial interpolation]
\label{OP:lem:unit-factorial}
For $p\ge5$ and every \(0\leq r<p\), there is
\(T_r(X)\in\mathbb Z_p\langle X\rangle\) such that, for all
integers \(\ell\geq0\),
\begin{equation}\label{OP:eq:unit-expansion}
 U(p\ell+r)^{-2}
 =c^{-2\ell}(r!)^{-2}
   \bigl(1-2pH_r\ell+p^2T_r(\ell)\bigr),
 \qquad H_r=\sum_{a=1}^r\frac1a.
\end{equation}
\end{lemma}

\begin{proof}
The unit factorial over one block has the exact paired form
\[
 Q(k):=\frac{U(p(k+1))}{cU(pk)}
 =\prod_{a=1}^{(p-1)/2}
       \left(1+\frac{p^2k(k+1)}{a(p-a)}\right).
\]
Thus \(\log Q(X)\) is a sum, indexed by \(m\geq1\), of
polynomials of degree at most \(2m\) whose coefficients lie in
\(p^{2m-v_p(m)}\mathbb Z_p\).
The polynomial \(S_d(X)\), equal to
\(\sum_{0\leq k<\ell}k^d\) at nonnegative integers
\(X=\ell\), has coefficient denominators dividing
\(\operatorname{lcm}(1,\ldots,d+1)\).  Indeed, expansion in
falling factorials gives
\[
 k^d=\sum_{a=0}^d
 \left\{\begin{matrix}d\\a\end{matrix}\right\}
 k(k-1)\cdots(k-a+1),
\]
and the sum of the falling factorial of degree \(a\) is the
falling factorial of degree \(a+1\), divided by \(a+1\).
Taking this discrete antiderivative term by term in
\(\log Q\), the \(m\)-th contribution therefore has
coefficient valuation at least
\[
 2m-v_p(m)-\lfloor\log_p(2m+1)\rfloor\geq2.
\]
The bound tends to infinity with \(m\), since \(p\geq5\).
We obtain a restricted series
\(A(X)\in p^2\mathbb Z_p\langle X\rangle\) such that
\[
 A(\ell)=\sum_{k=0}^{\ell-1}\log Q(k)
 =\log\bigl(U(p\ell)/c^\ell\bigr).
\]
Exponentiation converges in the restricted series ring and gives
\(\exp(-2A)\in1+p^2\mathbb Z_p\langle X\rangle\).
Finally,
\[
 U(p\ell+r)=U(p\ell)\,r!
                \prod_{a=1}^r(1+p\ell/a),
\]
and
\[
 \prod_{a=1}^r(1+pX/a)^{-2}
 \in1-2pH_rX+p^2\mathbb Z_p\langle X\rangle.
\]
Combining these identities proves the lemma.
\end{proof}


\subsection{The block formula used in both arguments}\label{OP:sec:block}
For the prime in use, put $c=(p-1)!$ and choose restricted functions
$W_r$ interpolating
\[
 U(p\ell+r)^{-2}=c^{-2\ell}(r!)^{-2}W_r(\ell).
\]
The preceding lemma gives them for $p\ge5$; Section~\ref{OP:sec:mixed}
gives the correct version for $p=3$. Write $j=pL+b$, $0\le b<p$,
$\epsilon=\lfloor2b/p\rfloor$. Splitting each $k_i=p\ell_i+r_i$
in the formula for $B(j,pN)$ gives a carry $d$ with
\[
 \sum r_i=b+pd,\quad\sum\ell_i=L-d,\quad0\le d\le n-1.
\]
Define
\[
 P_{b,d}(\mathbf X)=\sum_{\substack{0\le r_i<p\\\sum r_i=b+pd}}
       \prod_i(r_i!)^{-2}W_{r_i}(X_i).
\]
The factorial identity $k!=p^{\lfloor k/p\rfloor}\lfloor k/p\rfloor!U(k)$
and the balanced estimate give
\begin{align}\label{OP:eq:block-bessel}
 B(j,pN)\equiv{}&p^\epsilon U(2j)(2L+\epsilon)!
 \sum_{d=0}^{\min(n-1,L)}p^{2d}c^{-2(L-d)}[x^{L-d}]\\[-4pt]
 &\left.P_{b,d}(\theta_1,\ldots,\theta_n)
 I_N(x_1)\prod_{i>1}I_0(x_i)\right|_{x_i=x}
 \pmod{p^{2+2v_p(N)}}.\nonumber
\end{align}
At $N=0$ the equality is exact. The power of $p$ extracted from the
factorials is $2d+\epsilon$, and the remaining numerator is
$(2L+\epsilon)!$. Each original summand and the modified summand
is $p$-integral, so the balanced replacement is valid term by term.
The reduction of this symmetric operator is the only part that
previously required $p>n$.
% Complete extension of the octahedral matrix proof to every prime p>=5.
\section{Support-weighted symmetric reduction in arbitrary dimension}\label{OP:source-main-4}
\label{OP:sec:support-weighted-all-dimensions}

Fix any dimension $n\ge3$ and any prime $p\ge5$. No inequality between
$p$ and $n$ will be imposed. We use the unit-factorial interpolation and
the positive vertex matrix from the preceding section.

\subsection{The precise denominator of scalar reduction}
For a monomial in $X_1,\ldots,X_n$, its support is the set of variables
with positive exponent. Consider a symmetric restricted series $P_k$
all of whose nonzero monomials have support of size exactly $k$.
Applying its differential operator to
$I_N(x_1)\prod_{i>1}I_0(x_i)$, and using
\[
 \theta_i^2 I_0=x_iI_0,\qquad
 \theta_1^2 I_N=(x_1+N^2)I_N,
\]
reduces it, modulo $N^2\mathcal H$, to a combination of $E_l$ for
$0\le l\le k$. This support bound holds because reducing one factor
cannot differentiate a formerly undifferentiated factor. All reduction
coefficients and commutator coefficients are integral.

Set $D_l=l!E_l$. The elementary-derivative recurrence gives
\[
 D_0=\Phi_N,\quad D_1=\theta\Phi_N,\qquad
 D_{l+1}=\theta D_l-l(n-l+1)xD_{l-1}
                         \pmod{N^2\mathcal H}.
\]
More precisely, the omitted term is $N^2$ times an integral combination
of products of the indicated Bessel factors and their first
derivatives. Induction therefore expresses $D_l$ by an integral scalar
differential polynomial in $x,\theta$, with an error in
$N^2\mathcal H$. Hence expressing $E_l$ loses at most $v_p(l!)$,
both in its scalar coefficients and in the error. It does not lose
an additional factorial at each recurrence step.

It follows that a symmetric support-$k$ input in
$p^k\mathbb Z_p\langle\mathbf X\rangle$, for $k\ge2$, has a scalar
reduction with all coefficients and its $N^2$ error divisible by $p^2$:
indeed
\[
 k-v_p(l!)\ge k-v_p(k!)\ge2\qquad(0\le l\le k, p\ge3, k\ge2).
\]
A support-one input in $p^2\mathbb Z_p\langle\mathbf X\rangle$
has the same property, since only $E_0,E_1$ occur and no denominator is
introduced. Constants and a multiple of $\sum_iX_i$ reduce exactly to
the corresponding constant and multiple of $\theta$, with no positive
power of $x$.

These statements extend to restricted series: for fixed $n$, the only
denominators are among $1!,\ldots,n!$, and all other operations are
integral and bounded. Coefficients tending to zero before reduction
still tend to zero afterwards. The scalar coefficient series in $x$
are therefore restricted as well.

\subsection{Application to the actual unit-factorial operator}
For $0\le r<p$, unit-factorial interpolation gives
\[
 W_r(X)=1-2pH_rX+p^2T_r(X),\qquad
 T_r\in\mathbb Z_p\langle X\rangle.
\]
Evaluation at $\ell=0$ gives $T_r(0)=0$, hence $W_r(0)=1$.
For $b<p$ and a carry $d$, the operator used in the block formula is
\[
 P_{b,d}(\mathbf X)=
 \sum_{\substack{0\le r_i<p\\\sum r_i=b+pd}}
 \prod_i(r_i!)^{-2}W_{r_i}(X_i).
\]
After splitting each $W_r$ into its constant term and its positive-degree
part, a term with exactly $k$ positive-degree factors has support $k$
and a factor $p^k$. The symmetry of the residue-index set makes each
support component symmetric. For support one, its linear part is an
affine symmetric operator and every term of degree at least two has
a factor $p^2$. These are exactly the cases just covered.

Consequently there are restricted scalar coefficient series
$A_{d,0}(x),\ldots,A_{d,n}(x)\in\mathbb Z_p\langle x\rangle$
such that
\[
 \left.P_{b,d}(\theta_1,\ldots,\theta_n)
 I_N(x_1)\prod_{i>1}I_0(x_i)\right|_{x_i=x}
 -\sum_{k=0}^nA_{d,k}(x)\theta^k\Phi_N
 \in p^2N^2\mathcal H.
\]
All positive-degree coefficients of the $A_{d,k}$ lie in
$p^2\mathbb Z_p$. At $N=0$ the identity is exact.
This conclusion is asserted only for the displayed factorial operators;
it need not hold for an arbitrary symmetric input of the form
$\alpha+\beta\sum X_i+p^2R$ when $p\le n$.

\subsection{The row operator and the complete descent}
Writing $j=pL+b$, $\epsilon=\lfloor2b/p\rfloor$, and
$A_{d,k}(x)=\sum_a\alpha_{d,a,k}x^a$, define
\[
 D_{j,i}=p^\epsilon U(2j)\frac{(2L+\epsilon)!}{(2i)!}
 \sum_{d=0}^{\min(n-1,L-i)}p^{2d}c^{-2(L-d)}
 \sum_{k=0}^n\alpha_{d,L-d-i,k}i^k.
\]
The exact factorial-block decomposition and the preceding reduction
give, without any further use of $p>n$,
\begin{align*}
 a_j&=\sum_{i\le L}D_{j,i}a_i,\\
 B_n(j,pN)-\sum_{i\le L}D_{j,i}B_n(i,N)
   &\in p^{2+2v_p(N)}\mathbb Z_p.
\end{align*}
The coefficients $D_{j,i}$ are integral. If $i<L$, each summand either
has a carry $d\ge1$ and thus a factor $p^2$, or has positive $x$-degree
and thus a factor $p^2$ from the scalar reduction. Therefore
\[
 D_{j,i}\in p^2\frac{(2L+\epsilon)!}{(2i)!}\mathbb Z_p,
 \qquad
 D_{j,i}a_i\in p^2(L!/i!)^2\mathbb Z_p.
\]

Define the integral coefficients $c_{K,N}$ by
\[
 a_j\mathbf1_{j\ge K}=\sum_{N=K}^j c_{K,N}B_n(j,N).
\]
For completeness, the descent now proceeds exactly as follows. Induct
on $v=v_p(K)$. At $v=0$ integrality suffices. If $p\mid K$, put
$M=K/p$. By the induction hypothesis,
\[
 v_p(c_{M,N})\ge\max(0,2(v-1)-2v_p(N)).
\]
Hence every row error in replacing $B_n(j,pN)$ by its $D$-combination,
after multiplication by $c_{M,N}$, has valuation at least $2v$.
For $j\ge K$, the omitted boundary terms satisfy
\[
 \sum_{i<M}D_{j,i}a_i\in p^{2v}\mathbb Z_p,
\]
because $M\mid L!/i!$. For $j<K$ both triangular sums vanish. Thus
the vector supported on multiples of $p$, with entries $c_{M,N/p}$,
solves the defining triangular system for $c_{K,N}$ modulo $p^{2v}$.
Since $B_n$ is unit triangular,
\[
 c_{K,N}\equiv\mathbf1_{p\mid N}c_{M,N/p}\pmod{p^{2v}}.
\]
This also gives
$v_p(c_{K,N})\ge\max(0,2v_p(K)-2v_p(N))$, completing the induction.

The exact identity
\[
 1-F_K/F=\sum_{N\ge K}c_{K,N}C_N
\]
and excellent Cartier's coefficient congruence
\[
 C_{pN}-C_N\circ\sigma
 \in p^{2+2v_p(N)}\mathbb Z_p[[u]]
\]
now imply
\[
 \frac{F_K(u)}{F(u)}
 \equiv\frac{F_{K/p}(\sigma(u))}{F(\sigma(u))}
 \pmod{p^{2v_p(K)}}.
\]
Together with Proposition~\ref{OP:prop:orbit}, this proves that every octahedral dimension
$n\ge3$ satisfies the all-level, all-multiplier, every-coefficient
truncation congruence at every prime $p\ge5$.
% Research proof fragment: p=3, 3 divides the octahedral dimension.
% Uses the formal excellent Cartier input from Small_Prime_Advances.
\section{Mixed-precision descent at the prime three}\label{OP:source-main-5}\label{OP:sec:mixed}

Throughout this section $3\mid n$, and $B_n(j,N)$, $a_j=B_n(j,0)$,
$I_N$, $\Phi_N=I_NI_0^{n-1}$ and $c_{K,N}$ have their usual meanings.
We write $v=v_3$.

\subsection{The support-weighted reduction}
Let $\mathcal H=\{\sum h_jx^j:(2j)!h_j\in\mathbb Z_3\}$.
For a symmetric operator whose monomials involve exactly $k$ distinct
variables, reduction by $\theta_i^2I_N=(x_i+\delta_{i1}N^2)I_N$
uses elementary derivative products $E_l$ only for $0\le l\le k$.
Put $R_l=l!E_l$. Suppressing the $N^2$ terms, their recurrence is
\[
 R_0=\Phi_N,\quad R_1=\theta\Phi_N,\qquad
 R_{l+1}=\theta R_l-l(n-l+1)xR_{l-1}.
\]
Consequently conversion to scalar derivatives loses at most $v(l!)$,
including in the error term. If the input of support $k$ is multiplied
by $3^k$, then for $k\ge2$ its output and error retain a factor $9$,
because $k-v(l!)\ge k-v(k!)\ge2$. This reduction is bounded for fixed
$n$ and extends from polynomials to restricted series.

An additional observation will control the coefficient of $\theta^0$.
For $l\ge1$, the scalar differential polynomial expressing $R_l$
has all its positive-$x$-degree $\theta^0$ coefficients divisible by $n$.
This follows directly from the recurrence above, evaluated on $1$.
For $l=0$, a symmetric support-$k$ input contributes a multiple of
$\binom nk$. Since
\[
 \binom nk=\frac nk\binom{n-1}{k-1},
 \qquad k-v(k)\ge2\quad(k\ge2),
\]
the positive-$x$-degree scalar constant coefficient from a $3^k$
support-$k$ input is divisible by $9n$ in $\mathbb Z_3$.
For a support-one input with coefficients in $9\mathbb Z_3$, the same
conclusion is immediate. Thus all these good terms have positive-$x$
coefficients divisible by $9$, and their $\theta^0$ coefficients are
divisible by $9n$.

\subsection{Isolating the cubic term}
Write $U(a)=\prod_{1\le i\le a,\,3\nmid i}i$ and $c=2$.
For $0\le r<3$ there are restricted integral series $T_r$ such that
\[
 U(3\ell+r)^{-2}
 =c^{-2\ell}(r!)^{-2}W_r(\ell),\qquad
 W_r(X)=1-6H_rX-3(X^3-X)+9T_r(X).
\]
Indeed, the first term in the discrete antiderivative of the paired
unit-factorial logarithm is $\frac32(X^3-X)$; all subsequent terms
have coefficient valuation at least $3$. In exponentiating, every
power at least two has valuation at least $2$ after its factorial
denominator. Multiplication by
$\prod_{a\le r}(1+3X/a)^{-2}$ gives the formula.

For $b\in\{0,1,2\}$ and a carry $d$, put
\[
 P_{b,d}(\mathbf X)=
 \sum_{\substack{0\le r_i<3\\\sum r_i=b+3d}}
 \prod_i(r_i!)^{-2}W_{r_i}(X_i),
 \qquad
 \alpha_{b,d}=[y^{b+3d}](1+y+y^2/4)^n.
\]
Apart from an affine symmetric operator and the single term
\[
 -3\alpha_{b,d}\sum_i(X_i^3-X_i),
\]
all terms are good in the support-weighted sense just proved:
support one has a factor $9$, and support $k\ge2$ has a factor $3^k$.
The exceptional term has the exact reduction
\[
 \sum_i(\theta_i^3-\theta_i)\prod_i f_i
 = nx\Phi_N+(x-1)\theta\Phi_N
    +N^2(\theta I_N)I_0^{n-1},
\]
where $f_1=I_N$ and $f_i=I_0$ for $i>1$.
It therefore loses only one factor of $3$, explicitly and controllably.

\subsection{A row operator with two precisions}
There are integral coefficients $D_{j,i}$ for $0\le i\le L=\lfloor j/3\rfloor$
with the following properties:
\begin{align*}
 a_j&=\sum_{i\le L}D_{j,i}a_i,\\
 B_n(j,3N)-\sum_{i\le L}D_{j,i}B_n(i,N)
   &\in3^{1+2v(N)+\mathbf1_{3\nmid j}}\mathbb Z_3,\\
 D_{j,i}a_i&\in
 3^{2+\mathbf1_{3\nmid j}}(L!/i!)^2\mathbb Z_3
 \qquad(i<L).
\end{align*}

To prove this, use the residue decomposition $j=3L+b$ and the exact
factorial-block expression \eqref{OP:eq:block-bessel}. Its carry-$d$
term has the prefactor
$3^{\epsilon+2d}U(2j)(2L+\epsilon)!c^{-2(L-d)}$,
where $\epsilon=\lfloor2b/3\rfloor$.
Apply the preceding support-weighted reduction and the displayed exact
cubic reduction, and define $D_{j,i}$ by collecting the coefficients
of $\theta^k\Phi_N$ as in the ordinary row construction.
The balanced unit-factorial replacement has precision
$3^{2+2v(N)}$ and hence is sufficient throughout.

The exceptional error only matters for $d=0$. For $b=1$ one has
$\alpha_{1,0}=n\in3\mathbb Z_3$, and for $b=2$ the outer factor
$3^\epsilon$ supplies the missing $3$. This proves the stated error.
At $N=0$ all reductions are exact.

For the boundary bound, first take $i>0$. Write
\[
 a_i=\frac{(2i)!}{(i!)^2}b_n(i),\qquad
 b_n(i)=\sum_{k_1+\cdots+k_n=i}\binom{i}{k_1,\ldots,k_n}^2.
\]
The integer $b_n(i)$ is divisible by $n$: it counts pairs of words
with the same content, and cyclically relabelling all $n$ letters acts
freely on such pairs when $i>0$.
Good positive-$x$ terms provide $9$, and the exceptional term provides
$3\alpha_{b,0}$; the factor $n\mid b_n(i)$ supplies the extra $3$.
For $b=1$ its $\alpha_{1,0}=n$, and for $b=2$ the outer $3$ supplies
the additional precision. Terms with $d\ge1$ already have $3^{2d}$.
After multiplication by $a_i$, the factorial factor is an integer
multiple of $(L!/i!)^2$, exactly as in the original boundary argument.

For $i=0$ and positive $x$-degree, use the stronger divisibility of the
scalar constant coefficient: good terms have $9n$, and the exceptional
one is $-3\alpha_{b,d}nx$.
The only remaining terms have $d=L\ge1$ and $x$-degree zero. If $L\ge2$,
$3^{2L}$ already supplies more than the required extra power. If
$L=d=1$, the required $9$ is present; for $b=2$ there is the outer $3$;
for $b=1$, the coefficient
$\alpha_{1,1}=[y^4](1+y+y^2/4)^n$ is divisible by $3$, since $3\mid n$
and this polynomial has, modulo $3$, only degrees divisible by $3$.
This completes the boundary estimate.

\subsection{The preserved triangular lattice}
For $3\mid n$, Lucas' theorem applied to the factorial summands gives
\[
 B_n(j,3N)\equiv a_b B_n(\lfloor j/3\rfloor,N)\pmod3,
 \qquad b=j\bmod3.
\]
Here $a_1=2n\equiv0$ and $a_2=6n(2n-1)\equiv0$ modulo $3$.
In particular $B_n(j,N)\in3\mathbb Z_3$ when $3\nmid j$ and $3\mid N$.
Therefore the unit-triangular matrix $B_n$ and its inverse preserve
the lattice
\[
 \{(e_j):e_j\in3^{q+\mathbf1_{3\nmid j}}\mathbb Z_3\}.
\]

We prove simultaneously, by induction on $v(K)$, that
\[
 v(c_{K,N})\ge1+\max\{0,2v(K)-2v(N)\},
\]
and, if $3\mid K$,
\[
 c_{K,N}-\mathbf1_{3\mid N}c_{K/3,N/3}
 \in3^{2v(K)+\mathbf1_{3\nmid N}}\mathbb Z_3.
\]
The base case uses $3\mid a_j$ for $j>0$ and unit-triangular inversion.
Set $M=K/3$. The row operator and the induction hypothesis imply that
\[
 a_j\mathbf1_{j\ge K}-\sum_{N\ge M}c_{M,N}B_n(j,3N)
 \in3^{2v(K)+\mathbf1_{3\nmid j}}\mathbb Z_3.
\]
The boundary terms have this precision because $M\mid L!/i!$ for
$i<M\le L$; the row errors have this precision because
\[
 1+2v(N)+1+\max(0,2v(M)-2v(N))\ge2v(K).
\]
The additional power for $3\nmid j$ is supplied by the row lemma.
The preserved lattice gives the asserted refined descent.
For $3\nmid N$ it gives the enhanced weighted bound directly. For
$3\mid N$, the induction bound on $c_{M,N/3}$ supplies the desired
bound, and the error of valuation at least $2v(K)$ is stronger than
$1+\max(0,2v(K)-2v(N))$.

Finally combine ordinary descent modulo $3^{2v(K)}$, the exact tail
expansion, and excellent Cartier's congruence
$C_{3N}-C_N\circ\sigma\in3^{2+2v(N)}\mathbb Z_3[[u]]$.
This proves the complete all-level, all-multiplier, all-coefficient
octahedral truncation congruence at $p=3$ whenever $3\mid n$.
\section{The first level}\label{OP:source-main-6}\label{OP:sec:first-level}
We record the geometric first-level argument for completeness. Write
\[
c_j=\CT(g^j),\quad \mathcal F(z)=\sum_{j\ge0}c_jz^j,
\quad H_N(z)=\CT(f^{N-1}),\quad
\lambda=\mathcal F(z)/\mathcal F(\tau(z)).
\]
In this section $\sigma_p(z)=\tau(z)$ and $\sigma_p(u)=\sigma(u)$
denote the two normalizations. Proposition~\ref{OP:prop:orbit} supplies
the excellent Cartier identity at every odd prime.
\begin{lemma}[Second-order geometric approximation]
\label{OP:L:lem:second-geometric}
For the above excellent lift,
\[
 \lambda\equiv2H_p-H_{2p}\pmod{p^2\mathbb Z_p[[z]]}.
\]
More precisely, write
\[
 \mathscr C_p(f^{2p-1})=H_{2p}(z)+B_p(z)g(\mathbf{x}).
\]
Then
\[
 \tau\equiv-\frac{B_p(z)}{H_{2p}(z)}\pmod{p^2\mathbb Z_p[[z]]}.
\]
\end{lemma}

\begin{proof}
All the calculations below may be made in the ring of Laurent series
whose coefficients are \(z\)-adically convergent integral power series.
Put \(h(\mathbf{x})=f^\sigma(\mathbf{x}^p)\).  Since \(h-f^p\) is
divisible by \(p\), the exact identity
\[
 \frac1f-\frac{2hf^{p-1}-f^{2p-1}}{h^2}
 =\frac{(h-f^p)^2}{fh^2}
\]
has an integral remainder divisible by \(p^2\).  Applying the Cartier
operator and using
\(\mathscr C_p(A(\mathbf{x})/h(\mathbf{x})^2)
 =\mathscr C_p(A) / f^\sigma(\mathbf{x})^2\), we obtain
\[
 \mathscr C_p(1/f)
 \equiv
 \frac{2f^\sigma\mathscr C_p(f^{p-1})
              -\mathscr C_p(f^{2p-1})}{(f^\sigma)^2}
 \pmod{p^2}.
\]

We justify the two numerator supports explicitly.  Let \(\Gamma\)
be the lattice spanned by the support of \(g\).  Since its index is
prime to \(p\), \(p\mathbf{u}\in\Gamma\) implies
\(\mathbf{u}\in\Gamma\).  Reflexivity gives integral facet functions
\(L_j\) with \(\Delta=\{L_j\leq1\}\).  A lattice point in
\((1-1/p)\Delta\) has \(L_j\leq0\) for every \(j\), and hence is
zero.  A lattice point in \((2-1/p)\Delta\) has \(L_j\leq1\), and
hence belongs to \(\Delta\).  The only points of
\(\Delta\cap\Gamma\) are zero and the vertices.  Complete symmetry
therefore yields the exact identities
\[
 \mathscr C_p(f^{p-1})=H_p,
 \qquad
 \mathscr C_p(f^{2p-1})=H_{2p}+B_p g.
\]


The excellent Cartier relation
\(\mathscr C_p(1/f)\equiv\lambda/f^\sigma
 \pmod{p^2\mathscr F_2^\sigma}\)
is in particular a coefficientwise congruence modulo \(p^2\).
Multiplying by \((f^\sigma)^2\) gives
\[
 \lambda(1-\tau g)
 \equiv 2H_p(1-\tau g)-H_{2p}-B_p g\pmod{p^2}.
\]
Since \(\operatorname{CT}(g)=0\), comparison of the coefficients of
\(1\) and \(g\) first gives
\(\lambda\equiv2H_p-H_{2p}\), and then
\(\tau H_{2p}\equiv-B_p\).
The latter can be divided by the unit \(H_{2p}\), whose constant
term is one.
\end{proof}

\begin{theorem}[The first level at every odd prime]
\label{OP:L:thm:first-all-dim}
Set $h=2$ for the hyperoctahedral family.  At every odd prime,
\[
 \frac{\mathcal F(z)}{\mathcal F(\sigma_p(z))}
 \equiv\mathcal F_{hp}(z)\pmod{p^2\mathbb Z_p[[z]]}.
\]
Equivalently, in the normalized variable \(u=z^h\),
\[
 \frac{F(u)}{F(\sigma_p(u))}
 \equiv F_p(u)\pmod{p^2\mathbb Z_p[[u]]}.
\]
These are the \(s=1\) cases of the specified original truncation
lengths; \(\mathcal F_h=1\).
\end{theorem}

\begin{proof}
The coefficient of \(z^j\) in \(2H_p-H_{2p}\) is
\[
 c_jw_j,\qquad
 w_j=(-1)^j\left(2\binom{p-1}{j}-\binom{2p-1}{j}\right),
\]
with binomial coefficients zero past the top index.  If \(0\leq j<p\),
the product expansion
\[
 (-1)^j\binom{ap-1}{j}
 =\prod_{i=1}^j\left(1-\frac{ap}{i}\right)
 \equiv1-apH_j\pmod{p^2}
\]
shows that \(w_j\equiv1\pmod{p^2}\).
For \(p\leq j<2p\), Lucas's congruence gives
\(w_j\equiv1\pmod p\).

The remaining required divisibilities follow directly from the
factorial formulas.  In the octahedral case,
\[
 c_{2k}=\sum_{k_1+\cdots+k_n=k}
              \frac{(2k)!}{(k_1!\cdots k_n!)^2},
\]
and every denominator is prime to \(p\) when \(k<p\); hence the
numerator is divisible by $p$ whenever $2k\ge p$.
Thus $p\mid c_j$ for $p\le j<2p$: odd $j$ give zero,
and for even $j=2k$ in this range the numerator $(2k)!$ is divisible
by $p$ while every denominator is a unit.
It follows coefficientwise that
\[
 2H_p-H_{2p}\equiv\mathcal F_{hp}\pmod{p^2}.
\]
Apply Lemma~\ref{OP:L:lem:second-geometric} and then the normalized
coordinate relation \(q_0(z)^h=q(z^h)\).
\end{proof}


\section{Finite polynomial certificates for entire series}\label{OP:source-main-7}\label{OP:sec:cert-framework}
\subsection{Vertex polynomials and their derivative estimate}
Define
\[
 b_j=\sum_{k_2+\cdots+k_n=j}\binom{j}{k_2,\ldots,k_n}^{\!2},\qquad
 A_K(u)=\sum_{j=0}^{K-1}u^j\sum_{\ell=0}^j(-1)^\ell
 \frac{(2K-\ell-1)!b_{j-\ell}}
 {(2K-2j-1)!\ell!(j-\ell)!^2}.
\]
The factorial ratio is an integer multinomial coefficient and $A_K(0)=1$.
Expanding $1/f$ in the chamber with dominant vertex $x_1$, the negative
of the coefficient of $x_1^{-2K}$, all other exponents being zero, is
\[
 P_{2K}(u)=u^{-K}A_K(u).
\]
Let $\lambda=F(u)/F(\sigma(u))$. Reading this vertex coefficient in
the excellent Cartier congruence gives
\begin{equation}\label{OP:eq:vertex}
 P_{2pK}(u)-\lambda P_{2K}(\sigma(u))
       \in p^{2+2\vp(K)}\Z_p((u)).
\end{equation}
The displayed elements have bounded negative support. No claim that every
element of the completed coefficient ring is a Laurent series with bounded
negative support is needed.

\begin{lemma}\label{OP:lem:derivative}
For $C_K=A_{2K}/A_K$ and $\theta=u\,d/du$,
\[
 \theta C_K\in p^{\vp(K)}\Z_p[[u]].
\]
\end{lemma}
\begin{proof}
Set $R_K=P_{2K}/F$ and use the normalized canonical coordinate $q$, with
$\partial=q\,d/dq$. The second-order quotient relation
\cite[Proposition 7.6]{BV3}, read at the same vertex, gives after division
by $F$ and passage to $q$
\[
 \partial^2R_K\in p^{2\vp(K)}\Z_p((q)).
\]
More explicitly, in normalized coordinates let
$\mathcal L=\theta^2-B_u\theta-A_u$ and $H=\theta q/q=W/F^2$.
Then $\theta H/H=B_u-2\theta F/F$ and the exact identity
\[
\mathcal L(FR)=FH^2\partial^2R
\]
shows that this passage preserves the required precision.
The gauge and coordinate factors are integral units by
\cite[Proposition 7.9 and Corollary 7.11]{BV3}; the normalization factor
$2$ is a $p$-adic unit. For $r=\vp(K)$, integrality of $R_K$ together
with $j^2[q^j]R_K\in p^{2r}\Z_p$ implies
$j[q^j]R_K\in p^r\Z_p$ for every nonzero integer $j$.
Thus $\partial R_K$, and then $\theta R_K$, belong to
$p^r\Z_p((q))$ and $p^r\Z_p((u))$, respectively.
Differentiate $C_K=u^K R_{2K}/R_K$; the factors after removing their
displayed powers of $u$ are integral units. Every term has the factor
$p^r$, giving the assertion.
\end{proof}

\subsection{The scaled Hensel equation}
Fix an odd prime $p$, dimension $n$, and level $s\ge2$, and set
\[
 M=p^{2s},\quad M_y=p^{2s-1},\quad N=p^{2s-2},\quad X=u^p,
 \qquad \sigma(u)=X(1+py(u)),\quad y\in\Z_p[[u]].
\]
Eliminating $\lambda$ between \eqref{OP:eq:vertex} for $K=N,2N$ gives
\[
 (1+py)^N C_{pN}(u)-C_N(X(1+py))
       \in p^{4s-2}\Z_p[[u]].
\]
Divide this expression by $pN$. Its coefficient of $y^r$, for $r\ge1$,
has valuation at least
\begin{equation}\label{OP:eq:coefficient-valuation}
 r-1-\vp(r).
\end{equation}
Indeed Lemma~\ref{OP:lem:derivative} implies that the $r$th Taylor
coefficient of $C_N(X(1+h))$ lies in $(N/r)\Z_p[[u]]$;
also $\binom Nr\in(N/r)\Z_p$. Multiplication by $p^r/(pN)$ gives
\eqref{OP:eq:coefficient-valuation}. The constant coefficient is integral
as well, since the actual $y$ solves the preceding congruence and all
positive-degree coefficients are integral.

For $p=3$, take $d=3$ at $s=2$ and $d=6$ at $s=3$.
The bound \eqref{OP:eq:coefficient-valuation} is at least $2s-1$ for every
$r>d$. Thus the equation for $y$ modulo $M_y$ is a degree-$d$ polynomial.
Its linear coefficient has constant term $1$, and all nonlinear
coefficients are divisible by $p$.

Here are completely explicit common-denominator numerators for that
polynomial. Compute the four vertex polynomials at indices
$N,2N,pN,2pN$ modulo $p^{4s-2}$ and set
\[
 a_r=[h^r]A_N(X(1+h)),\quad b_r=[h^r]A_{2N}(X(1+h)),
 \quad a_0=A_N(X).
\]
Define
\[
 T_0=b_0,\qquad
 T_r=b_ra_0^r-\sum_{i=1}^r a_iT_{r-i}a_0^{i-1}.
\]
Then $T_r/a_0^{r+1}=[h^r]C_N(X(1+h))$. For $0\le r\le d$ put
\begin{equation}\label{OP:eq:Qs}
 Q_r=\frac{p^r}{pN}
 \left(\binom Nr A_{2pN}a_0^{d+1}
             -A_{pN}T_ra_0^{d-r}\right)\pmod{M_y}.
\end{equation}
The quotients are coefficientwise integral by the preceding argument;
the certificate independently checks this exact divisibility.
Their common denominator $A_{pN}a_0^{d+1}$ is a unit. Thus
\begin{equation}\label{OP:eq:Hensel}
 \sum_{r=0}^d Q_r y^r=0\pmod{M_y},\qquad
 Q_0(0)=0,\quad Q_1(0)=1,\quad Q_r\in p\Z_p[u]\ (r\ge2).
\end{equation}
Two roots of this equation have difference multiplied by a unit,
because after subtracting their equations the factor is congruent to
$Q_1$ modulo $p$. Therefore the root modulo $M_y$ is unique.

The program constructs a rational root $Y/Z$ from these $Q_r$ using
the finite multivariate Lagrange expansion, with $Z(0)=1$, and then
checks the entire polynomial identity
\begin{equation}\label{OP:eq:root-check}
 \sum_{r=0}^d Q_rY^r Z^{d-r}=0\pmod{M_y}.
\end{equation}
This identity and uniqueness, rather than an assumption about the
target truncations, prove $y=Y/Z\pmod{M_y}$ and hence
$\sigma=X(1+pY/Z)\pmod M$.

For clarity, the construction is explicit. Write $A=Q_0$, $B=Q_1$.
For a profile $(m_2,\ldots,m_d)$ of nonnegative integers, put
$m=\sum m_r\ge1$, $R=\sum r m_r$ and $\ell=R-m+1$. Its contribution to
$Y/Z$ is
\[
 (-1)^{R+1}\frac{R!}{\ell!\prod m_r!}
 \frac{A^\ell\prod_{r=2}^d Q_r^{m_r}}{B^{R+1}}.
\]
Include also the term $-A/B$, and retain only profiles of total
coefficient valuation less than $2s-1$.
The integer factorial quotient is the usual rooted-plane-tree
coefficient in Lagrange inversion. The code checks its integrality.
The fixed common denominators used below are $B^7$ and $B^{13}$ at $p=3$, levels $2$ and $3$, respectively.
The direct check \eqref{OP:eq:root-check} makes the root certification
independent of any truncation of the Lagrange expansion.

\subsection{The target polynomial identities}
The vertex relation at $K=N$ yields
\[
 \lambda=\frac{(1+py)^N A_{pN}(u)}{A_N(X(1+py))}\pmod M.
\]
Since $\vp\binom Nr\ge\vp(N)-\vp(r)$, for every $r\ge2$ the term
$\binom Nr(py)^r$ is divisible by $M$. Hence
$(1+h)^N=1+Nh\pmod M$ on $h\in p\Z_p[[u]]$.

Set $f=F_{p^s}$, $f_0=F_{p^{s-1}}$ and, for $0\le r<2s$, define
\[
 L_r=A_{pN}(u)\bigl([h^r]f_0(X(1+h))
               +N[h^{r-1}]f_0(X(1+h))\bigr)
          -f(u)[h^r]A_N(X(1+h)),
\]
where a coefficient with negative index is zero.
The target congruence is equivalent to
$\sum_{r<2s}L_rh^r=0\pmod M$.
Substitute the already certified $h=pY/Z$ and clear the unit denominator:
\begin{equation}\label{OP:eq:target-check}
 \boxed{\ \sum_{r=0}^{2s-1}L_r(pY)^r Z^{2s-1-r}=0\pmod{p^{2s}}.\ }
\end{equation}
Only at this last step do the ordinary period truncations enter the
calculation. The program verifies every coefficient of
\eqref{OP:eq:root-check} and \eqref{OP:eq:target-check} exactly.


If the two identities hold, multiplication by the unit $F(\sigma)$
gives $E_{n,p,s,1}\in p^{2s}\mathbb Z_p[[u]]$. Thus a finite exact
polynomial certificate proves a statement about the entire series.
\par
\section{Newton certificates uniform in the octahedral dimension}\label{OP:source-main-8}
\label{OP:sec:dimension-newton}

The dimension parameter can be handled by integer polynomials rather
than by a finite search with an unproved upper bound.  This gives the
following additional result at the wild prime $3$.

\begin{theorem}\label{OP:thm:uniform-dimension-levels23}
For every octahedral dimension $n\geq2$, let $\sigma_n$ be its normalized
excellent lift.  Then, for $s=2,3$,
\[
 F_n(u)F_{n,3^{s-1}}(\sigma_n(u))-
 F_n(\sigma_n(u))F_{n,3^s}(u)
                 \in3^{2s}\mathbb Z_3[[u]].
\]
This is an assertion about all dimensions and every formal-series
coefficient at each of the two specified levels.  The multiplier is
$m=1$; no induction over all levels is asserted here.
\end{theorem}

We use the negative-vertex and scaled-Hensel certificate construction.
Its formal Cartier and derivative inputs are those already established
by the integral orbit-module argument.  The new ingredient is a finite
difference criterion that makes the dependence on $n$ rigorous.

\subsection{Integer polynomials in the dimension}
Put
\[
 b_n(j)=\sum_{k_1+\cdots+k_n=j}
       \binom{j}{k_1,\ldots,k_n}^{\!2},\qquad
 a_{n,j}=\binom{2j}{j}b_n(j).
\]
For each $j$, $b_n(j)$ is the value at $n$ of a polynomial in
$\mathbb Z[T]$.  To see this, count ordered pairs of length-$j$ words
with the same letter multiplicities.  If precisely $r$ letters occur,
let $d_{j,r}$ be the number of pairs on a fixed $r$-letter alphabet
using every letter.  Simultaneous renaming of letters is a free action
of $S_r$, so $r!\mid d_{j,r}$, and
\[
 b_n(j)=\sum_{r=0}^j\binom nr d_{j,r}
       =\sum_{r=0}^j\frac{d_{j,r}}{r!}
             n(n-1)\cdots(n-r+1).
\]
Thus the assertion concerns ordinary integer coefficients, not merely
integer-valued polynomials.

In particular the coefficients of every period truncation and of
\[
 A_{n,K}(u)=\sum_{j<K}u^j\sum_{\ell=0}^j(-1)^\ell
 \frac{(2K-\ell-1)!\,b_{n-1}(j-\ell)}
 {(2K-2j-1)!\ell!(j-\ell)!^2}
\]
are integer polynomials in $n$.  The factorial ratio is an integer
multinomial coefficient.

For $e\geq1$ set
\[
 \nu_3(e)=\min\{r\geq0:v_3(r!)\geq e\},\qquad
 \Delta P(T)=P(T+1)-P(T).
\]

\begin{lemma}[Finite difference criterion]\label{OP:lem:newton-test}
Let $c$ be an integer.
\begin{enumerate}
\item If $P(T)\in3^{-a}\mathbb Z_3[T]$, then
\[
 \Delta^rP(c)\in3^{-a}r!\mathbb Z_3
 \qquad(r\geq0).
\]
\item If $P\in\mathbb Z_3[T]$ and
$P(c),\ldots,P(c+\nu_3(e)-1)$ are all divisible by $3^e$, then
$P(n)\in3^e\mathbb Z_3$ for every integer $n$.
\end{enumerate}
The same conclusions hold coefficientwise with an additional
polynomial variable $u$.
\end{lemma}
\begin{proof}
Finite differences of monomials are multiples of $r!$, by the usual
Stirling-number formula, and translation by $c$ preserves ordinary
integer coefficients.  This proves the first assertion.  For the
second, use the finite Newton expansion
\[
 P(c+X)=\sum_{r\geq0}\Delta^rP(c)\binom Xr.
\]
For $r<\nu_3(e)$, the difference is an integer combination of the
assumed values.  For $r\geq\nu_3(e)$ the first assertion makes it
divisible by $3^e$.  Finally $\binom Xr$ is an integer at every integer
$X$, including negative integers.
\end{proof}

\subsection{Ordinary-polynomial models for the Hensel coefficients}
Fix $s=2$ or $3$, and put
\[
 k=2s-1,\quad N=3^{2s-2},\quad M=3^{2s},\quad
 M_y=3^k,\quad M_{\rm work}=3^{4s-2},\quad X=u^3.
\]
Let $d=3$ when $s=2$, and $d=6$ when $s=3$.
The four vertex polynomials are evaluated at $N,2N,3N,6N$.
Writing
\[
 a_r=[h^r]A_{n,N}(X(1+h)),\quad
 b_r=[h^r]A_{n,2N}(X(1+h)),
\]
define
\[
 T_0=b_0,\qquad
 T_r=b_ra_0^r-\sum_{i=1}^r a_iT_{r-i}a_0^{i-1}
\]
and
\[
 Q_r(n,u)=\frac{3^r}{3N}
 \left(\binom Nr A_{n,6N}a_0^{d+1}
       -A_{n,3N}T_ra_0^{d-r}\right),\qquad0\leq r\leq d.
\]
Before reduction, these are polynomials satisfying
\[
 Q_r(T,u)\in3^{-k}\mathbb Z_3[T,u].
\]
For every integer dimension $n\geq2$, their coefficient values are
integral.  The scaled-Hensel equation is
\[
 \sum_{r=0}^dQ_r(n,u)y_n(u)^r=0\pmod{3^k},\qquad
 \sigma_n(u)=u^3(1+3y_n(u)).
\]
The root is unique because $Q_1(n,0)=1$ and the nonlinear coefficients
are divisible by $3$.  The omitted Taylor coefficients have valuation
at least $r-1-v_3(r)\geq k$.

We now construct polynomials
\[
 \widetilde Q_r(T,u)\in\mathbb Z_3[T,u]
\]
whose values agree with $Q_r(n,u)$ modulo $3^k$ for every $n\geq2$.
Set
\[
 D=\nu_3(2k),\qquad d_0=\nu_3(k),\qquad c=2.
\]
The program computes the $D$ input values at $n=2,\ldots,D+1$ modulo
$3^k$, and all their differences of orders $0\leq j<D$.  It verifies,
coefficientwise in $u$, that
\begin{equation}\label{OP:eq:newton-integrality-check}
 \Delta^j Q_r(2,u)\in
 3^{\min(k,v_3(j!))}\mathbb Z_3[u]
 \qquad(0\leq j<D).
\end{equation}
For $j\geq D$, Lemma~\ref{OP:lem:newton-test} gives
$\Delta^jQ_r(2,u)\in3^k\mathbb Z_3[u]$ automatically, since the only
possible common denominator is $3^k$.  The finite calculation therefore
controls the entire Newton expansion in the dimension variable.

For $0\leq j<d_0$, choose the coefficientwise integer representative
$\delta_{r,j}(u)$ of $\Delta^jQ_r(2,u)$ modulo $3^k$ with coefficients
between $0$ and $3^k-1$.  The verified divisibility implies that
$\delta_{r,j}/j!\in\mathbb Z_3[u]$.  Define
\[
 \widetilde Q_r(T,u)=
 \sum_{j<d_0}\frac{\delta_{r,j}(u)}{j!}
                (T-2)(T-3)\cdots(T-j-1).
\]
For $j=0$ the product is $1$.  Every discarded Newton coefficient is
zero modulo $3^k$, either by the calculation or by the factorial bound.
For the retained coefficients, replacing the exact difference by
$\delta_{r,j}$ changes its value by a multiple of $3^k\binom{n-2}{j}$.
Hence
\[
 \widetilde Q_r(n,u)\equiv Q_r(n,u)\pmod{3^k}
 \qquad(n\geq2).
\]
This argument proves the model for all dimensions; it is not an
extrapolation from the computed dimensions.

\subsection{Two polynomial identities imply all dimensions}
Use the same fixed finite Lagrange expression in the
$\widetilde Q_r$ at every dimension to construct
\[
 Y(T,u),Z(T,u)\in\mathbb Z_3[T,u].
\]
Concretely, let $A=\widetilde Q_0$, $B=\widetilde Q_1$ and use the
universal nonlinear valuation weights
$w_r=r-1-v_3(r)$, $2\leq r\leq d$.
For a profile $(m_2,\ldots,m_d)$ with
$\sum w_rm_r<k$, put $R=\sum rm_r$, $m=\sum m_r$, and
$\ell=R-m+1$.  Besides $-A/B$, its rational-root term is
\[
 (-1)^{R+1}\frac{R!}{\ell!\prod m_r!}
       \frac{A^\ell\prod_{r=2}^d\widetilde Q_r^{m_r}}{B^{R+1}}.
\]
The factorial quotient is integral.  Clear all these terms with the
fixed denominator $Z=B^b$, where $b=7$ for $s=2$ and $b=13$ for $s=3$.
The resulting $Y$ and $Z$ have ordinary $3$-integral coefficients in
the dimension variable.  The program directly verifies the root
identity below, so its conclusion does not rely on an unverified
truncation of the Lagrange expression.

The two finite polynomials tested are
\[
 \mathcal R(T,u)=\sum_{r=0}^d
       \widetilde Q_r(T,u)Y(T,u)^rZ(T,u)^{d-r}
\]
and
\[
 \mathcal T(T,u)=\sum_{r=0}^{2s-1}
       L_r(T,u)(3Y(T,u))^rZ(T,u)^{2s-1-r},
\]
where $L_r$ is the target numerator
\[
 L_r=A_{n,3N}\bigl([h^r]F_{n,3^{s-1}}(X(1+h))
       +N[h^{r-1}]F_{n,3^{s-1}}(X(1+h))\bigr)
       -F_{n,3^s}[h^r]A_{n,N}(X(1+h)),
\]
viewed as an ordinary integer polynomial in $T=n$.
Every coefficient of both $\mathcal R$ and $\mathcal T$ belongs to
$\mathbb Z_3[T]$.

For each of the $\nu_3(2s)$ consecutive dimensions
\[
 n=2,\ldots,\nu_3(2s)+1,
\]
the program checks the entire identities
\[
 \mathcal R(n,u)=0\pmod{3^{2s-1}},\qquad
 \mathcal T(n,u)=0\pmod{3^{2s}}.
\]
Lemma~\ref{OP:lem:newton-test}, applied to each coefficient of $u$,
proves both identities in every dimension.  Also $Z(n,0)$ is a unit
for all dimensions, because the modeled linear coefficient agrees
modulo $3$ with the actual coefficient whose constant term is $1$.
The root identity and uniqueness give $y_n=Y(n,u)/Z(n,u)\pmod{3^k}$.
The target identity consequently gives the asserted excellent-lift
truncation congruence.  This proves
Theorem~\ref{OP:thm:uniform-dimension-levels23}.

\subsection{Reproducible exact verification}
The complete verifier is
\texttt{certify\_all\_dimensions\_newton.py}; it imports the exact
polynomial arithmetic from \texttt{verify\_dimension\_families.py}.
Run it with \texttt{--s 2} or \texttt{--s 3}.
The input and output counts are as follows.
\begin{center}
\setlength{\tabcolsep}{4pt}\small
\begin{tabular}{c|r|r|r|r|r}
$s$&working modulus&input dimensions&model degree&target dimensions&max. target degree\\
\hline
2&$729$&15&8&9&1562\\
3&$59049$&24&11&15&70519
\end{tabular}
\end{center}
All difference-integrality checks and all coefficients of both
polynomial identities pass.  The target modulus is $81$ in the first
row and $729$ in the second row.  No formal-series coefficient cutoff
is used.

The deterministic SHA-256 certificate identifiers are
\begin{center}\ttfamily\small
\begin{tabular}{c|l}
$s=2$&ca7ba583179e93370023d03b848ddf0f52b63fa39dcc57b3dd9d4c078e79938f\\
$s=3$&cbe30563e24bebbd768fce3ce9eda9513d18ec87e8fa9aa07814c062e1a83712
\end{tabular}
\end{center}
The archived certificate includes the input Hensel coefficients, all
Newton differences, the ordinary-model coefficients, the fixed
Lagrange profiles, and the target verification reports.  The digests
identify those data; the proof is the exact algebra above together with
the reproducible finite polynomial calculations.
\section{The fourth level by compact rational Hensel certificates}\label{OP:source-main-9}
\label{OP:sec:dimension-newton-level4}

The Newton method extends one level further.  A compact rational
Hensel calculation avoids explicitly expanding the large universal
numerator and denominator in the dimension parameter.

\begin{theorem}\label{OP:thm:uniform-dimension-level4}
For every octahedral dimension $n\geq2$,
\[
 F_n(u)F_{n,27}(\sigma_n(u))-
 F_n(\sigma_n(u))F_{n,81}(u)
                         \in6561\mathbb Z_3[[u]].
\]
Together with the earlier first-level argument and
Theorem~\ref{OP:thm:uniform-dimension-levels23}, this proves the normalized
multiplier-$1$ congruence at levels $s=1,2,3,4$ in every dimension.
It does not give an induction to all levels.
\end{theorem}

\subsection{The dimension model at precision $3^7$}
Here the parameters in the preceding construction are
\[
 s=4,\quad k=7,\quad N=729,\quad M=6561,\quad M_y=2187,
 \quad M_{\rm work}=4782969=3^{14},\quad d=9.
\]
The Taylor valuation bound $r-1-v_3(r)$ is at least $7$ for every
$r\geq10$, so a degree-$9$ scaled Hensel equation suffices.  Its
coefficients, before reduction, satisfy
\[
 Q_r(T,u)\in3^{-7}\mathbb Z_3[T,u],\qquad0\leq r\leq9.
\]
The four negative-vertex polynomials used are
$A_{n,729}$, $A_{n,1458}$, $A_{n,2187}$ and $A_{n,4374}$, all computed
modulo $3^{14}$.  The division by $3N=3^7$ therefore determines the
Hensel coefficients modulo $3^7$ without loss of required precision.

We have
\[
 \nu_3(14)=30,\qquad \nu_3(7)=18,\qquad \nu_3(8)=18.
\]
The exact calculation at the $30$ input dimensions $n=2,\ldots,31$
verifies all the finite-difference conditions
\[
 \Delta^j Q_r(2,u)\in
 3^{\min(7,v_3(j!))}\mathbb Z_3[u]
       \qquad(0\leq j<30,\ 0\leq r\leq9).
\]
For $j\geq30$, the common-denominator bound and
Lemma~\ref{OP:lem:newton-test} give divisibility by $3^7$ automatically.
The preceding Newton construction thus supplies fixed polynomials
\[
 \widetilde Q_r(T,u)\in\mathbb Z_3[T,u],\qquad
 \deg_T\widetilde Q_r\leq17,
\]
such that
\[
 \widetilde Q_r(n,u)=Q_r(n,u)\pmod{2187}
                     \qquad\hbox{for every }n\geq2.
\]
This statement has already been proved before the target truncations
are tested.

\subsection{A universal polynomial numerator exists}
The following argument makes the change of rational denominators in
individual compact certificates harmless.  Put
$B=\widetilde Q_1$ and define rational functions in $T,u$ by
\[
 y_0=0,\qquad
 y_{j+1}=-\frac{\widetilde Q_0}{B}
       -\sum_{r=2}^9\frac{\widetilde Q_r}{B}y_j^r
                        \qquad(0\leq j<7).
\]
There are fixed polynomials $Y_*(T,u),Z_*(T,u)\in\mathbb Z_3[T,u]$
with $y_7=Y_*/Z_*$.  Indeed one may take $Z_*=B^e$, where the
successive denominator exponents satisfy
\[
 e_0=0,\qquad e_{j+1}=9e_j+1;
 \qquad e_7=597871.
\]
Only the existence of these finite polynomials is needed; they are not
expanded by the verifier.

For each integer dimension $n\geq2$, $B(n,u)$ is a formal-series unit,
and $\widetilde Q_r(n,u)\in3\mathbb Z_3[[u]]$ for $r\geq2$.
These divisibilities are pointwise in the dimension; no divisibility
of every ordinary polynomial coefficient in $T$ is asserted.  Thus the
map on the right of the recurrence contracts differences by a factor
of $3$ on $\mathbb Z_3[[u]]$.  Starting from $0$, seven iterations
agree modulo $3^7$ with the unique integral Hensel root.  Consequently
\begin{equation}\label{OP:eq:universal-rational-root4}
 \frac{Y_*(n,u)}{Z_*(n,u)}=y_n(u)\pmod{2187}
                          \qquad(n\geq2),
\end{equation}
and $Z_*(n,0)$ is a unit.

Let $L_0(T,u),\ldots,L_7(T,u)\in\mathbb Z_3[T,u]$ be the same target
polynomials as before, now with $N=729$, $F_{n,81}$ and $F_{n,27}$.
The single fixed polynomial
\[
 \mathcal T_*(T,u)=
 \sum_{r=0}^7L_r(T,u)(3Y_*(T,u))^rZ_*(T,u)^{7-r}
                                  \in\mathbb Z_3[T,u]
\]
will be used for the dimension argument.  It is important that this
polynomial is fixed independently of $n$, even though the compact
rational functions computed below use different denominators at
different dimensions.

\subsection{Compact roots at the eighteen testing dimensions}
For each $n=2,\ldots,19$, the program constructs a small rational
root $\widehat Y_n/\widehat Z_n$ of the degree-$9$ Hensel equation
modulo $2187$.  Its denominator has unit constant term.
At the initial digit, reduce $-Q_0/Q_1$ in $\mathbb F_3(u)$ by a
polynomial gcd.  Suppose $A/D$ is already a root modulo $3^j$, and put
\[
 S=\sum_{r=0}^9Q_rA^rD^{9-r}\pmod{3^{j+1}}.
\]
Every coefficient of $S$ is divisible by $3^j$.  Over $\mathbb F_3(u)$
the correction is
\[
 \delta=-\frac{S/3^j}{\overline D^{9}\overline Q_1}.
\]
Reduce this fraction by a polynomial gcd and replace $A/D$ by
$A/D+3^j\delta$.  Repeating for $j=1,\ldots,6$ gives the desired root.
At every step all divisions by powers of $3$ are checked exactly.
The final calculation independently verifies the entire identity
\[
 \sum_{r=0}^9Q_r\widehat Y_n^{r}
                     \widehat Z_n^{9-r}=0\pmod{2187}.
\]
Thus correctness of the final certificate does not depend on merely
trusting the gcd algorithm or the digit construction.

The program then verifies every coefficient of
\begin{equation}\label{OP:eq:compact-target4}
 \sum_{r=0}^7L_r(n,u)(3\widehat Y_n)^r
                        \widehat Z_n^{7-r}=0\pmod{6561}
                   \qquad(n=2,\ldots,19).
\end{equation}
No formal-series cutoff is used.  The largest of these compact target
polynomials has degree $10471$.

The actual Hensel polynomial and its ordinary dimension model agree
modulo $2187$.  Uniqueness therefore identifies each compact root
with the universal root \eqref{OP:eq:universal-rational-root4} modulo
$2187$.  Multiplication by $3$ improves the precision to $6561$ for
the variable $h=3y$ in the target polynomial.  Because both
denominators are units, \eqref{OP:eq:compact-target4} implies
\[
 \mathcal T_*(n,u)=0\pmod{6561}
                       \qquad(n=2,\ldots,19).
\]
Apply Lemma~\ref{OP:lem:newton-test} coefficientwise to the fixed
ordinary polynomial $\mathcal T_*(T,u)$.  The $18=\nu_3(8)$ consecutive
values imply the same divisibility for every integer dimension
$n\geq2$.  By \eqref{OP:eq:universal-rational-root4}, clearing and then
restoring the unit denominators gives
$\lambda_nF_{n,27}(\sigma_n)-F_{n,81}\in6561\mathbb Z_3[[u]]$.
Multiplication by the unit $F_n(\sigma_n)$ proves
Theorem~\ref{OP:thm:uniform-dimension-level4}.

\subsection{Exact implementation and reproducibility}
The model calculation uses factorial unit parts and valuations.  For
an integral factorial ratio, subtract the factorial valuations and
multiply the unit parts and their inverses modulo the working prime
power; then restore the nonnegative power of $3$.  This computes the
same integer ratio modulo that prime power.  The implementation
checks nonnegativity of every resulting valuation.
NumPy integer matrix operations are exact here: the explicit bound
\[
 (3^{14})^2\cdot4374<2^{63}
\]
guarantees that no signed $64$-bit dot product can overflow.  All
polynomial arithmetic in the root and target identities uses exact
integer arithmetic.

The model stage has a \texttt{--model-only} option, so no attempt to
expand a large universal numerator is necessary.  A complete
reproduction consists of
\begin{verbatim}
python3 certify_all_dimensions_newton_fast.py --s 4 --model-only
python3 compact_hensel.py --s 4 --n \
  2 3 4 5 6 7 8 9 10 11 12 13 14 15 16 17 18 19
python3 summarize_level4_certificate.py
\end{verbatim}
The first script saves the $30$ Hensel inputs, their Newton
differences, and the ordinary models.  The second saves the $18$
compact roots and their complete target numerator polynomials.  The
last script records deterministic identifiers for these already
verified data.

The aggregate SHA-256 identifier is
\begin{center}\ttfamily\small
795a410dff5b06382a482373dc2e2451689b3c643634ac13b66e384ab0ca313c.
\end{center}
It hashes the model checkpoint identifier together with all eighteen
individual compact-certificate identifiers.  As before, the proof is
the exact algebra and the reproducible polynomial identities; the
hash only identifies the archived data.
\section{Limits of the auxiliary rows and their resolution}\label{OP:source-main-10}
The two arguments in this part do not cover $p=3$, $n\ge4$,
$3\nmid n$ at arbitrary levels. Part~\ref{part:octa-three}
supplies that case, including arbitrary normalized multipliers.
The following failures of stronger auxiliary estimates explain
why a different argument is needed; they remain valid after the
target congruence has been proved.

The obstruction at $3$ is concrete. The coefficient
$-3(X^3-X)$ in unit-factorial interpolation cannot be placed in
$9\mathbb Z_3\langle X\rangle$. Its values on integers do have an
additional factor of $3$, but ordinary restricted-series coefficient
divisibility is stronger than integer-valued divisibility. The
mixed-precision proof compensates for this using $3\mid n$; it does
not justify applying the same induction when $3\nmid n$.

A seemingly stronger alternative row construction is also insufficient.
Define the unique exact integral row by
\[
 B(j,3N)=\sum_{i\le L}D^{\rm ex}_{j,i}B(i,N),\qquad L=\lfloor j/3\rfloor,
\]
including $N=0$. It need not have the individual boundary divisibility
used above. For example, in dimension $5$, at $j=12,L=4,i=0$,
\[
 D^{\rm ex}_{12,0}\equiv27\pmod{81},
\]
whereas $9(4!)^2$ has $3$-adic valuation $4$. These are failures of
auxiliary proposed estimates, not counterexamples to the target
excellent-lift congruence.

The numerical checks of the mixed induction in dimensions
$3,6,9,12,15,18$, for $K\le40$ and $N\le120$, are consistency checks
only. The all-level theorems rest on the written induction. In contrast,
the three finite-level dimension conclusions explicitly include exact finite
computation as part of their proofs; their finite test bounds are
proved by the Newton lemma. No step treats a power-series cutoff or a
finite list of dimensions as an unproved substitute for a universal
statement.




% END RETAINED SOURCE OP
\partstart{All-level three-adic descent}{part:octa-three}{The integral cutoff kernel, universal weak descent, dimension residues, and an independent fifth-level certificate.}
% BEGIN RETAINED SOURCE AP
\section{Statement and proof architecture}\label{AP:source-main-1}\label{AP:sec:scope}
Let
\[
 g_n(\mathbf x)=\sum_{i=1}^n(x_i+x_i^{-1}),\qquad
 f=1-zg_n,\qquad u=z^2,
\]
and put
\[
 b_n(j)=\sum_{k_1+\cdots+k_n=j}
             \binom{j}{k_1,\ldots,k_n}^{\!2},\qquad
 a_{n,j}=\binom{2j}{j}b_n(j),
\]
\[
 F_n(u)=\sum_{j\ge0}a_{n,j}u^j,\qquad
 F_{n,K}(u)=\sum_{0\le j<K}a_{n,j}u^j.
\]
The excellent lift $\sigma_n(u)$ is defined in
Section~\ref{AP:sec:formal-input}.  All divisibilities below are
coefficientwise in the displayed formal-series or ordinary polynomial
ring.

\begin{theorem}[All levels at the remaining prime]\label{AP:thm:main}
For every integer $n\ge3$ and every positive integer $K$ divisible by $3$,
\[
 \frac{F_{n,K}(u)}{F_n(u)}
 \equiv
 \frac{F_{n,K/3}(\sigma_n(u))}{F_n(\sigma_n(u))}
       \pmod{3^{2v_3(K)}\mathbb Z_3[[u]]}.
\]
In particular, for every $s,m\ge1$,
\begin{align*}
 E_{n,3,s,m}(u)
 &:=F_n(u)F_{n,m3^{s-1}}(\sigma_n(u))
       -F_n(\sigma_n(u))F_{n,m3^s}(u)\\
 &\in3^{2s}\mathbb Z_3[[u]].
\end{align*}
When $3\mid m$, the proof gives the stronger precision
$3^{2s+2v_3(m)}$.
\end{theorem}

The same arithmetic proof works in dimensions one and two; they are
also used as residue representatives in the dimension argument.
Normalized length $m3^s$ in $u$ is length $2m3^s$ in $z$.
Thus $m=1$ is the hyperoctahedral multiplier $2$ in
BV3, Conjecture~7.5~\cite{BV3}.

The proof has four independent mathematical components:
\begin{enumerate}
\item An orbit module supplies the formal excellent Cartier congruence
at every odd prime, including $3\mid n$.
\item An integral Bessel kernel proves the uniform quadratic cutoff
bound in the ordinary coefficient ring $\mathbb Z[T]$.
\item A universal row over $\mathbb Z_3\langle T\rangle$ gives
descent modulo $3^{2v_3(K)-1}$ coefficientwise in $T$.
\item The scalar dimensions $0,1,2$ recover the missing factor at
every integer dimension.  The exact tail expansion transfers the
result to ordinary truncations.
\end{enumerate}
The level-five certificates are separate checks of entire series.
They are not premises of the all-level argument.  This part
establishes the remaining formal $3$-adic hyperoctahedral case. The
analytic transfer at the original BV3-admissible primes is stated
separately in Corollary~\ref{G:cor:analytic}.

\section{The formal excellent Cartier input}\label{AP:source-main-2}\label{AP:sec:formal-input}
\begin{proposition}[Formal excellent input from an orbit module]\label{AP:prop:orbit}
Let $g=g_n$ be the hyperoctahedral Laurent polynomial and let $p$ be an odd prime.
Put $R=\mathbb Z_p[[z]]$ and $f=1-zg$.
Let $F_0(z)$ and $G_0(z)$ be the period and logarithmic companion
given by the vertex-coefficient formulas of BV3, Proposition~9.1~\cite{BV3}.
Explicitly, $F_0(z)=F_n(z^2)$ and, with $H_r=\sum_{a=1}^r1/a$,
\[
 G_0(z)=\sum_{j\ge1}z^{2j}
 \sum_{k_1+\cdots+k_n=j}
 (H_{2j}-H_{k_1})\frac{(2j)!}{\prod_i(k_i!)^2}.
\]
Then
\[
 q_0(z)=z\exp(G_0(z)/F_0(z))\in z+z^2\mathbb Z_p[[z]].
\]
The lift $\tau=q_0^{-1}(q_0(z)^p)$ is integral and satisfies
\[
 \mathscr C_p(1/f)-\frac{F_0(z)}{F_0(\tau)}\frac1{f^\tau}
 \in p^2\mathscr F_2^\tau.
\]
In particular, no exclusion $p\mid n$ is needed for this formal input.
\end{proposition}

\begin{proof}
Write $\Gamma$ for the support lattice and $H$ for the full monomial
symmetry group. Its order will never be inverted. The support lattice is $\mathbb Z^n$, so it is $p$-saturated.
The symmetry group preserves the reflexive polytope degree
$d(u)=\deg_\Delta(u)$ and the support lattice. For an $H$-orbit
$O\subset\Gamma$, put
\[
 S_O=\sum_{u\in O}x^u,\qquad d(O)=d(u)\ (u\in O),\qquad
 L=\bigoplus_O R z^{d(O)}S_O.
\]
Here a Laurent polynomial uses only finitely many orbits.
This is a direct construction of a module of polynomials. We do not
obtain it by applying an averaging projector to another crystal.

Fix temporarily any lift $\sigma$ with
$z^\sigma=z^pv$, $v\in1+pR$. We check BV3, Definition~2.1.
The orbit sums have disjoint supports and coefficient one, so $L$ is
$p$-saturated. The same is true of every $L^{\sigma^i}$: its orbit
coefficient generator is $(z^{\sigma^i})^{d(O)}$, a power of $z$
times a unit, which is a non-zero-divisor modulo $p$.
Multiplication by $f$ preserves invariance and admissibility, hence
$fL\subset L$.

The subset $p\Gamma$ is $H$-stable. An orbit is either disjoint from
$p\Gamma$ or has the form $pO'$ for one orbit $O'$; the equality uses
the $p$-saturation of the support lattice. Since $d(pu)=pd(u)$,
\[
 \mathscr C_p(z^{d(pO')}S_{pO'})
   =z^{pd(O')}S_{O'}
   =\left(\frac{z^p}{z^\sigma}\right)^{d(O')}
        (z^\sigma)^{d(O')}S_{O'}\in L^\sigma.
\]
The image of any other orbit is zero. This proves compatibility
without dividing by either $|H|$ or $|O|$.

Form the Dwork crystal
$M=\widehat\Omega_{L,f}(\Delta^\circ)$ directly, as in BV3,
Definitions~2.6 and~2.8. We do not identify it with the invariant
part of an ambient completed module; in particular no exactness of
an invariant functor is being assumed. In this family,
$\Delta\cap\Gamma$ consists of the origin and the vertices.
Reflexivity then implies that the first two levels have bases
\[
 M(1)=R/f,\qquad M(2)=R/f^2\oplus Rzg/f^2.
\]
Indeed $2\Delta^\circ\cap\Gamma=\Delta\cap\Gamma$ and the
vertices form one orbit. The pair $1/f,zg/f^2$ is an extended basis,
since $1/f=(1-zg)/f^2$ and the change of basis has determinant one.
These levels are $p$-saturated in the ambient integral formal
expansion: multiplying an expansion by $f$ or $f^2$ reduces the
claim to divisibility of the numerator coefficients. This also
shows that the two basis vectors remain independent modulo $p$.

The Hasse--Witt computations of BV3, Lemma~7.2, now apply to this
explicit compatible module. Their coefficient extraction uses a
single vertex coefficient of a polynomial known to be constant on
the vertex orbit; it does not average over that orbit. With the
initial lift $z^\sigma=z^p$, the first determinant has constant
term one, and
\[
 \left.\frac{\det HW^{(2)}(M)}p\right|_{z=0}
 =-\frac1p\binom{2p-2}{p-1}\in\mathbb Z_p^\times.
\]
The latter computation reads the coefficient of $z^px^{pv}$ at one
vertex $v$ in $(zg)F^{(2)}$; the extreme vertex forces the unique
contribution $(-z)^p\binom{2p-2}{p-1}$. The normalized determinant
modulo $p$ is independent of the lift. Consequently both
Hasse--Witt conditions hold over $R$ for every lift under consideration.
The general BV3, Corollary~5.9 and Corollary~4.3, therefore give
\[
 M=M(2)\oplus\mathscr F_2M,
 \qquad
 \mathscr C_p(M)\subset M^\sigma(2)+p^2\mathscr F_2M^\sigma.
\]
Only $p>2$ is required by these general results.

For completeness, the remaining canonical-coordinate argument
uses these statements rather than a group-order hypothesis.
The derivation $\theta=z\,d/dz$ preserves $M$: for an admissible
numerator $A$,
\[
 \theta\left((r-1)!\frac{A}{f^r}\right)
 =(r-1)!\frac{\theta A}{f^r}
       +r!\frac{Azg}{f^{r+1}}.
\]
Thus the rank-two
quotient has the connection and Cartier matrices of BV3,
Propositions~7.6 and~7.7. The positive vertex coefficient calculation
of Proposition~9.1 gives exactly $F_0$ and
$F_0\log z+G_0$ as its two solutions. The same coefficient
calculation at exponents $p^sv$ gives $A(0)=B(0)=0$, the
constant Cartier matrix $\operatorname{diag}(1,p)$ for the
initial lift, and the upper-right coefficient
\[
 \lambda_1=
 \frac{pF_0(z)F_0(z^\sigma)}{W(z^\sigma)}
 \left(-\frac1p\log\frac{z^\sigma}{z^p}
 +\frac{G_0(z)}{F_0(z)}
 -\frac1p\frac{G_0(z^\sigma)}{F_0(z^\sigma)}\right),
\]
where $W=F_0^2+F_0\theta G_0-(\theta F_0)G_0$.
The rank-one reduction modulo $p$ implies $\lambda_1\in pR$.
In the Cartier matrix $\Lambda=(\lambda_i;\mu_i)$, horizontality
gives
\[
 \mu_1=\theta\lambda_1+
       \frac{\theta(z^\sigma)}{z^\sigma}
              (\lambda_0+B^\sigma\lambda_1)\in pR,
\]
because $\theta(z^\sigma)/z^\sigma\in pR$ and the connection
coefficient $B$ is integral.  Hence $\det\Lambda/p\in R$.
The determinant identity
$\det\Lambda=pW/W^\sigma$ and $W(0)=1$ give $W\in1+zR$
by the same iteration as BV3, Proposition~7.9.
For $z^\sigma=z^p$, the displayed formula and the
Dieudonne--Dwork lemma yield $q_0\in z+z^2R$.
Its formal inverse is integral, so
$\tau=q_0^{-1}(q_0^p)$ defines a lift with
$\tau/z^p\in1+pR$. Reapply the compatible-module construction
with this lift. The displayed expression has $\lambda_1=0$,
and the remaining coefficient is $F_0/F_0^\tau$.
This proves the required excellent Cartier identity in
$p^2\mathscr F_2M^\tau$.

Finally, membership in $\mathscr F_2M$ means that every $s$-fold
Cartier iterate lies in $p^{2s}M^{\sigma^s}$. Mapping the directly
constructed module to the ambient Dwork module therefore sends
$\mathscr F_2M$ into the ambient $\mathscr F_2$.
Only this inclusion is needed; no converse assertion about
invariants or completed modules is used.
\end{proof}


The sign symmetry gives $q_0(z)=zQ(z^2)$ and
$\tau(-z)=-\tau(z)$. Thus $\tau=z^pV(u)$ with
$V\in1+p\mathbb Z_p[[u]]$, and
\[
 \sigma_n(u)=\tau(z)^2=u^pV(u)^2.
\]
No division by the dimension occurs in this normalization. The
corresponding canonical coordinate in $u$ is $q(u)=q_0(z)^2$.

\section{Positive vertices and exact tails}\label{AP:source-main-3}\label{AP:sec:positive}
Suppress the subscript $n$ when it is fixed. For $0\le N\le j$ put
\[
 B(j,N)=(2j)!\sum_{\substack{k_1+\cdots+k_n=j\\k_1\ge N}}
 \frac{1}{(k_1-N)!(k_1+N)!\prod_{i>1}(k_i!)^2},
\]
and set $B(j,N)=0$ if $N>j$. These are integers, $B(j,j)=1$, and
$B(j,0)=a_j$. Let
\[
 \Omega_N(u)=\sum_{j\ge N}B(j,N)u^j,\qquad
 C_N=\Omega_N/F\in u^N\mathbb Z_p[[u]].
\]
The series $\Omega_N$ is the coefficient at $x_1^{2N}$, with all
other exponents zero, in the positive expansion of $1/f$. Reading
Proposition~\ref{AP:prop:orbit} at that exponent and using the
coefficient characterization of second formal derivatives gives
\begin{equation}\label{AP:eq:positive-cartier}
 C_{pN}-C_N\circ\sigma_n
       \in p^{2+2v_p(N)}\mathbb Z_p[[u]].
\end{equation}
Here the exponent factor is $(2N)^2$ and $2$ is a $p$-adic unit.

The unit triangular system uniquely defines integers $c_{K,N}$ by
\[
 a_j\mathbf1_{j\ge K}=\sum_{N=K}^jc_{K,N}B(j,N),
 \qquad c_{K,N}=0\quad(N<K).
\]
Consequently
\begin{equation}\label{AP:eq:exact-tail}
 1-F_K/F=\sum_{N\ge K}c_{K,N}C_N.
\end{equation}
This identity converges $u$-adically. It is the connection between
Cartier's coefficient congruence and ordinary truncations.

\subsection{A balanced unit-factorial estimate at every odd prime}
Write $U(a)=\prod_{1\le i\le a,\,p\nmid i}i$. For $H=pN$ and $k\ge H$,
\[
 \frac{U(k)^2}{U(k-H)U(k+H)}
 =\prod_{\substack{k-H<i\le k\\p\nmid i}}(1+H/i)^{-1}
 \equiv1\pmod{p^{2+2v_p(N)}}.
\]
Indeed, for $v=v_p(H)$ the unit residues in this interval are complete
copies modulo $p^v$. Their reciprocals sum to zero modulo $p^v$ by
pairing opposite residues. Expanding the product to first order proves
the estimate.
\section{A two-variable kernel and quadratic cutoff divisibility}\label{AP:source-main-4}\label{AP:sec:kernel}

Put
\[
 J(x)=\sum_{j\geq0}\frac{x^j}{(j!)^2},\qquad
 I_N(x)=\sum_{j\geq N}\frac{x^j}{(j-N)!(j+N)!},\qquad
 \tht=x\frac{d}{dx}.
\]
Then $\tht^2J=xJ$ and $(\tht^2-x)I_N=N^2I_N$.
In this section, $n$ is an indeterminate.  Define
\[
 J(x)^n=\sum_{j\geq0}b_n(j)\frac{x^j}{(j!)^2},\qquad
 a_{n,j}=\binom{2j}{j}b_n(j),
\]
and
\[
 B_n(j,N)=(2j)![x^j]\bigl(I_N(x)J(x)^{n-1}\bigr).
\]
The matrix $B_n$ is lower triangular with diagonal entries $1$.  For $K\geq1$
define $c_{K,N}(n)$, $N\geq K$, by the formal identity
\begin{equation}\label{AP:eq:tail-definition}
 (J(x)^n)_{\geq K}
   =\sum_{N\geq K}c_{K,N}(n)I_N(x)J(x)^{n-1}.
\end{equation}
Equivalently,
\[
 a_{n,j}\,\boldsymbol 1_{j\geq K}
   =\sum_{N=K}^{j}c_{K,N}(n)B_n(j,N).
\]

\begin{theorem}[Quadratic cutoff divisibility]\label{AP:thm:weighted}
For all integers $1\leq K\leq N$,
\begin{equation}\label{AP:eq:strong-weighted}
 \boxed{\displaystyle
  \frac{N^2}{2nK^2}\,c_{K,N}(n)\in\ZZ[n].}
\end{equation}
The quotient in \eqref{AP:eq:strong-weighted} is a polynomial: in particular,
the assertion includes divisibility of $c_{K,N}$ by $n$.
\end{theorem}

We prove a stronger statement at the level of the Bessel differential
operator.  The proof uses the following two elementary integrality facts.

\begin{lemma}[Integral dimension polynomials]\label{AP:lem:dimension}
For every $j\geq0$, $b_n(j)\in\ZZ[n]$.  Consequently $J(x)^n$ and
$J(x)^{1-n}$ belong to the factorial-square series ring
\[
 \mathcal A_{\ZZ[n]}=
 \left\{\sum_{j\geq0}q_j\frac{x^j}{(j!)^2}:q_j\in\ZZ[n]\right\}.
\]
The analogous two-variable ring, with denominators $(a!b!)^2$, is also closed
under multiplication.
\end{lemma}
\begin{proof}
The binomial expansion of $J^n$ gives
\[
 b_n(j)=\sum_{r=0}^{j}(n)_r T_{j,r},\qquad
 T_{j,r}=\frac1{r!}
 \sum_{\substack{k_1+\cdots+k_r=j\\k_i\geq1}}
       \binom{j}{k_1,\ldots,k_r}^{\!2},
\]
where $(n)_r=n(n-1)\cdots(n-r+1)$.  The integer in the sum before division by
$r!$ counts pairs of ordered partitions of two labelled $j$-element sets,
with corresponding blocks nonempty and of equal size.  Simultaneously
permuting the $r$ block labels is a free action of the symmetric group.
Thus $T_{j,r}\in\ZZ$.  The case $j=r=0$ has its usual value $1$.
Closure of the series rings follows because their multiplication constants
are squares of binomial coefficients.  Substitution of $1-n$ for $n$ gives
the second assertion.
\end{proof}

\begin{lemma}[An integral ballot quotient]\label{AP:lem:ballot}
For nonnegative integers $a,b$,
\begin{equation}\label{AP:eq:ballot}
 \frac{a-b}{a+b+1}\binom{2a}{a}\binom{2b}{b}\in\ZZ.
\end{equation}
\end{lemma}
\begin{proof}
The case $a=b$ is immediate.  Fix a prime $p$ and put
$h=v_p(a+b+1)$.  For $1\leq r\leq h$, the residues $A,B$ of $a,b$ modulo
$q=p^r$ satisfy $A+B=q-1$.  The $r$th contribution to the valuation of the
product of central binomial coefficients is
\[
 \left\lfloor\frac{2A}{q}\right\rfloor+
 \left\lfloor\frac{2B}{q}\right\rfloor.
\]
It is $1$ except when $q$ is odd and $A=B=(q-1)/2$, when it is $0$.
For odd $p$, this exception is equivalent to $p^r\mid 2a+1$, and, since
$p^r\mid a+b+1$, to $p^r\mid a-b$.  For $p=2$ every contribution is $1$;
when $h>0$, $a-b$ is then odd.  Summing the first $h$ contributions therefore
gives at least $\max(0,h-v_p(a-b))$.  All remaining contributions are
nonnegative.  This proves that the numerator in \eqref{AP:eq:ballot} contains
every prime power in its denominator.
\end{proof}

Put $r(x)=\tht J(x)/J(x)$.  This is a well-defined series with zero constant
term.  Define the divided-difference kernel
\begin{equation}\label{AP:eq:g-def}
 g(X,Y)=\frac{(r(X)-r(Y))^2}{X-Y}.
\end{equation}
The apparent division is harmless: the numerator is divisible by
$(X-Y)^2$ in $\QQ[[X,Y]]$.

\begin{lemma}[Integral Bessel kernel]\label{AP:lem:kernel}
The kernel $g$ belongs to the two-variable factorial-square series ring
$\mathcal A^{(2)}_{\ZZ}$; explicitly,
\[
 (a!b!)^2[X^aY^b]g(X,Y)\in\ZZ\qquad(a,b\geq0).
\]
\end{lemma}
\begin{proof}
Write
\[
 A=J(X)J(Y),\qquad
 W=(\tht_XJ(X))J(Y)-J(X)(\tht_YJ(Y)),\qquad
 D=\frac{W^2}{X-Y}.
\]
Both $W/(X-Y)$ and $D$ are formal power series.  For
$T=\tht_X+\tht_Y$, the Bessel equation gives
\[
 TW=(X-Y)A.
\]
Since $T(X-Y)=X-Y$, it follows that
\[
 (T+1)D=2WA=(\tht_X-\tht_Y)A^2.
\]
Vandermonde's identity gives
\[
 J(X)^2=\sum_{a\geq0}\binom{2a}{a}\frac{X^a}{(a!)^2}.
\]
Comparing coefficients therefore yields the exact formula
\begin{equation}\label{AP:eq:D-coeff}
 (a!b!)^2[X^aY^b]D
  =\frac{a-b}{a+b+1}\binom{2a}{a}\binom{2b}{b}.
\end{equation}
Lemma~\ref{AP:lem:ballot} proves $D\in\mathcal A^{(2)}_{\ZZ}$.
Finally,
\[
 g=J(X)^{-2}J(Y)^{-2}D,
\]
and the inverse powers lie in the same integral ring by
Lemma~\ref{AP:lem:dimension} (or by geometric inversion in that ring).
\end{proof}

\begin{theorem}[Integral cutoff kernel]\label{AP:thm:cutoff}
Let
\[
 h_K(x)=J(x)^{1-n}(J(x)^n)_{\geq K},\qquad
 R_K(x)=(\tht^2-x)h_K(x).
\]
Then
\begin{equation}\label{AP:eq:R-integrality}
 \frac{((K-1)!L!)^2}{n}[x^{K+L}]R_K(x)\in\ZZ[n]
 \qquad(K\geq1,\ L\geq0).
\end{equation}
More precisely, the following exact two-variable generating identity holds:
\begin{equation}\label{AP:eq:bivariate-identity}
 \boxed{\displaystyle
 \sum_{K\geq1,\,L\geq0}
   \frac{[x^{K+L}]R_K(x)}{n}X^LY^{K-1}
  =J(X)^{1-n}J(Y)^n\bigl(1-(n-1)g(X,Y)\bigr).}
\end{equation}
Here division by $n$ denotes an identity in $\QQ[n][[X,Y]]$; the proof
itself establishes that every coefficient of $R_K$ is divisible by $n$.
\end{theorem}
\begin{proof}
Summing the cutoffs first gives
\[
 \sum_{K\geq1}h_K(x)z^{K-1}
  =\frac{J(x)-J(x)^{1-n}J(zx)^n}{1-z}.
\]
Apply $\tht_x^2-x$ with $z$ fixed.  Set $X=x$, $Y=zx$ and
$P=J(X)^{1-n}J(Y)^n$.  The operator $\tht_x$ becomes
$T=\tht_X+\tht_Y$.  Using
$\tht r=x-r^2$, direct logarithmic differentiation gives
\[
 (T^2-X)P
  =n(Y-X)P+n(n-1)(r(X)-r(Y))^2P.
\]
The term $J(x)$ is annihilated.  Thus
\[
 \sum_{K\geq1}R_K(X)(Y/X)^{K-1}
  =nX P\left(1-(n-1)
             \frac{(r(X)-r(Y))^2}{X-Y}\right).
\]
Since $R_K$ starts in degree at least $K$, dividing by $nX$ and comparing
nonnegative powers of $X,Y$ gives \eqref{AP:eq:bivariate-identity}.
Its right side is in $\mathcal A^{(2)}_{\ZZ[n]}$ by
Lemmas~\ref{AP:lem:dimension} and~\ref{AP:lem:kernel}.  This proves
\eqref{AP:eq:R-integrality}.
\end{proof}

\begin{proof}[Proof of Theorem~\ref{AP:thm:weighted}]
Multiply \eqref{AP:eq:tail-definition} by $J^{1-n}$ and apply
$\tht^2-x$.  This gives the exact locally finite identity
\[
 R_K(x)=\sum_{N\geq K}N^2c_{K,N}(n)I_N(x).
\]
For $j=K+L$, Theorem~\ref{AP:thm:cutoff} implies
\[
 \frac{(2j)!}{2nK^2}[x^j]R_K(x)
 =\frac12\binom{2j}{j}\binom{j}{K}^{\!2}
   \left(\frac{((K-1)!L!)^2}{n}[x^j]R_K(x)\right)
 \in\ZZ[n].
\]
Here $j\geq K\geq1$, so the central binomial coefficient is even.
On the other hand,
\[
 \frac{(2j)!}{2nK^2}[x^j]R_K(x)
  =\sum_{N=K}^{j}
      \frac{N^2c_{K,N}(n)}{2nK^2}\binom{2j}{j-N}.
\]
The matrix $\bigl(\binom{2j}{j-N}\bigr)_{j,N\geq K}$ is lower triangular
over $\ZZ$, with diagonal $1$.  Induction on $j$ therefore proves that
every quotient $N^2c_{K,N}(n)/(2nK^2)$ lies in $\ZZ[n]$.
\end{proof}

\begin{corollary}[A stronger truncated-convolution divisibility]
Define
\[
 d_j(K)=\sum_{k=K}^{j}\binom{j}{k}^{\!2}
                   b_n(k)b_{1-n}(j-k),\qquad d_{K-1}(K)=0.
\]
Then, for every $j\geq K\geq1$,
\[
 \frac{d_j(K)-d_{j-1}(K)}{n\binom{j-1}{K-1}^{\!2}}\in\ZZ[n].
\]
\end{corollary}
\begin{proof}
The definition gives $h_K(x)=\sum_{j\geq K}d_j(K)x^j/(j!)^2$.
Consequently
\[
 [x^j]R_K(x)=\frac{d_j(K)-d_{j-1}(K)}{((j-1)!)^2}.
\]
Apply \eqref{AP:eq:R-integrality} with $L=j-K$.
\end{proof}

\paragraph{$3$-adic consequence.}
For every $K,N\geq1$ with $N\geq K$, the theorem supplies the
coefficientwise bound
\[
 c_{K,N}(n)\in
  3^{\max(0,\,2v_3(K)-2v_3(N))}\ZZ_3[n].
\]
It also supplies an additional factor $n$.  No evaluation at an integer
dimension is needed to obtain these statements.  This coefficientwise
form is the feature required by a residue-class argument in the dimension.

\section{Universal weak descent in the dimension}\label{AP:source-main-5}\label{AP:sec:universal}
Put
\[
 J(x)=\sum_{j\ge0}\frac{x^j}{(j!)^2},\qquad
 I_N(x)=\sum_{j\ge N}\frac{x^j}{(j-N)!(j+N)!},\qquad
 \theta=x\frac d{dx}.
\]
For an indeterminate $T$ define
\[
 B_T(j,N)=(2j)![x^j]I_NJ^{T-1},\qquad
 a_{T,j}=B_T(j,0).
\]
These belong to $\mathbb Z[T]$, and $B_T(j,j)=1$.
Indeed, $(j!)^2[x^j]J^T$ is the integer polynomial obtained by counting
pairs of words with equal letter multiplicities, and multiplication by
$I_N$ gives integer multinomial factors.  The inverse triangular system
\[
 a_{T,j}\mathbf1_{j\ge K}
   =\sum_{N=K}^j c_{K,N}(T)B_T(j,N)
\]
therefore defines $c_{K,N}(T)\in\mathbb Z[T]$.

Theorem~\ref{AP:thm:weighted}, with $T$ in place of $n$, proves
the following ordinary-polynomial bound for every $1\le K\le N$:
\[
                    N^2c_{K,N}(T)\in K^2\mathbb Z_3[T].
\tag{UW}
\]
We use this proved bound throughout the rest of the argument.

\subsection{A universal factorial coefficient lattice}
Let $A=\mathbb Z_3\langle T\rangle$ be the ring of restricted power
series with integral coefficients, with its coefficient norm, and put
\[
 \mathcal H_A=\left\{\sum_{j\ge0}h_jx^j:(2j)!h_j\in A\right\}.
\]
Products, $x$, and $\theta$ preserve $\mathcal H_A$ and have norm at
most one.  For every integer $r$, $J^{T-r}\in\mathcal H_A$.  All
statements in this subsection are coefficient statements in $T$;
divisibility of values at integer dimensions alone would not suffice.

\begin{lemma}[A single-slot unit factorial interpolation]
There are $V_r(X)\in1+3\mathbb Z_3\langle X\rangle$, for
$r=0,1,2$, such that $V_r(0)=1$ and
\[
 U(3\ell+r)^{-1}=2^{-\ell}(r!)^{-1}V_r(\ell),\qquad
 U(a)=\prod_{\substack{1\le i\le a\\3\nmid i}}i.
\]
Writing $W_r=V_r^2$, there are restricted integral series $Z_r(X,Y)$
with
\[
 V_r(X-Y)V_r(X+Y)-W_r(X)=3Y^2Z_r(X,Y).
\tag{1}
\]
\end{lemma}
\begin{proof}
The paired block ratio is
\[
 \frac{U(3(\ell+1))}{2U(3\ell)}
       =1+\frac92\ell(\ell+1).
\]
The logarithm and its discrete antiderivative converge in the
restricted series ring.  The $m$-th logarithmic term, after discrete
antidifferentiation, has coefficient valuation at least
\[
 2m-v_3(m)-\lfloor\log_3(2m+1)\rfloor\ge1,
\]
and the bound tends to infinity.  Here is the denominator estimate used
in this step.  If an integral polynomial $P$ has degree at most $d$,
write $P(X)=\sum_{a=0}^d e_a(X)_a$ with integral $e_a$.
Its discrete antiderivative vanishing at zero is
$\sum_{a=0}^d e_a(X)_{a+1}/(a+1)$, since
$\Delta(X)_{a+1}=(a+1)(X)_a$.  Its coefficient denominator therefore
loses at most $\lfloor\log_3(d+1)\rfloor$ powers of $3$.
Apply this with $d=2m$ to the $m$th logarithmic term.
Exponentiation gives the interpolation
for $r=0$; multiply by $\prod_{a=1}^r(1+3X/a)^{-1}$ for the other
two residues.  The logarithmic interpolant vanishes at $X=0$.
For (1), write $V_r=1+3Z$.  The difference on its left has coefficients
divisible by $3$, is even in $Y$, and vanishes at $Y=0$.  Division by
$3Y^2$ thus remains in the integral restricted series ring.
\end{proof}

\subsection{Keeping dimension denominators in one factorial}
For a fixed positive integer dimension $n$, let $f_1=I_N$ and
$f_i=J$ for $i>1$.  For $\ell\ge0$ set
\[
 E_\ell=\sum_{|S|=\ell}\prod_{i\in S}\theta f_i
                              \prod_{i\notin S}f_i,
 \qquad D_\ell=\ell!E_\ell.
\]
Their polynomial-dimension continuation is explicit:
\[
 D_\ell=(T-1)_\ell I_NJ^{T-1-\ell}(\theta J)^\ell
 +\ell(T-1)_{\ell-1}(\theta I_N)J^{T-\ell}
                                  (\theta J)^{\ell-1},
\]
where the second term is omitted when $\ell=0$ and $(T)_r$ denotes
a falling factorial.  In particular $D_\ell\in\mathcal H_A$.
For $\ell\ge1$ the exact differential recurrence is
\[
 \theta D_\ell=D_{\ell+1}
  +\ell(T-\ell+1)xD_{\ell-1}
  +N^2\ell(T-1)_{\ell-1}I_NJ^{T-\ell}(\theta J)^{\ell-1}.
\tag{2}
\]
Together with $\theta D_0=D_1$, this shows that $D_\ell$ equals an integral scalar differential operator
in $x,\theta,T$ applied to $I_NJ^{T-1}$, modulo
$N^2\mathcal H_A$.  No factorial division occurs in this conversion.

An active-slot pattern has $R$ labelled occupied positions.  At each
position prescribe a residue $r\in\{0,1,2\}$ and either a constant
factor or a positive-degree factor of $W_r$.  The pair consisting of
residue zero and a constant factor is inactive and is omitted.
If precisely $k$ positions have positive-degree factors and the residue
sum is $b+3d$, then
\[
                        R\le k+b+3d.
\tag{3}
\]
Expand each $W_r$ as $1+3T_r$, where $T_r(0)=0$ and
$T_r\in\mathbb Z_3\langle X\rangle$.
The corresponding symmetric operator is the sum over ordered injections
of the $R$ positions into the $n$ variables, divided by $R!$, with a
factor $3^k$ and the units $(r!)^{-2}$.

For a monomial in a pattern, reduce each one-variable derivative using
$\theta^2J=xJ$ and $\theta^2I_N=(x+N^2)I_N$.  Modulo
$N^2\mathcal H_A$, all single-variable reduction coefficients are
integral and independent of the chosen slot.  If exactly $\ell$ of
the $R$ positions retain first derivatives, the injection sum is
\[
              (T-\ell)_{R-\ell}\ell!E_\ell
                 =(T-\ell)_{R-\ell}D_\ell.
\tag{4}
\]
Thus (2) gives scalar reduction with the \emph{single} denominator
$R!$.  The same denominator suffices for the error.  In fact, a
single-variable error can only arise when one of the labelled positions
is the distinguished first slot; the other positions contribute
$(T-1)_{R-1}$ times an integral product of Bessel factors and their
derivatives.  The error from (2) is explicitly integral in the same
lattice.  All these assertions extend continuously to the restricted
series in a fixed pattern.

Here are explicit error and convergence details.  There are polynomials
$A_q,B_q\in\mathbb Z[x]$ and $H_{q,N}\in\mathcal H_{\mathbb Z_3}$,
all of norm at most one, such that
\[
 \theta^qJ=A_qJ+B_q\theta J,\qquad
 \theta^qI_N=A_qI_N+B_q\theta I_N+N^2H_{q,N}.
\]
Start with $A_0=1$, $B_0=H_{0,N}=0$, and use
\[
 A_{q+1}=\theta A_q+xB_q,\quad
 B_{q+1}=A_q+\theta B_q,\quad
 H_{q+1,N}=\theta H_{q,N}+B_qI_N.
\]
If the derivative orders of the $R$ labelled positions are
$e_1,\ldots,e_R$, the error before division by $R!$ is exactly
\[
 N^2(T-1)_{R-1}J^{T-R}
       \sum_{h=1}^R H_{e_h,N}
                         \prod_{q\ne h}\theta^{e_q}J.
\tag{4a}
\]
There is only one distinguished Bessel factor, so errors from two
different slots never multiply.  Formula (4a) proves directly that
no dimension denominator is hidden in the error.
Similarly, iterating (2) gives
$D_\ell=P_\ell(x,\theta,T)\Phi_N+N^2Q_{\ell,N}$ with
$P_\ell\in\mathbb Z[x,\theta,T]$ in normal order and
$Q_{\ell,N}\in\mathcal H_A$, both of norm at most one.
Consequently the entire ordered-injection reduction is bounded by one
before its single division by $R!$.

For fixed $R$, coefficients of the restricted input tend to zero with
the total derivative order; the preceding bounded maps therefore
converge to scalar coefficient series in $A\langle x\rangle$ and
errors in $\mathcal H_A$.  The scalar coefficients have degree at
most $R$ in $T$, before this restricted limit.  For fixed residue
sum $B=b+3d$, one also has $k\ge R-B$.  Thus the norm of the
contribution of all patterns of size $R$, after their prefactor, is
bounded by
\[
 3^{-\epsilon-2d-k+v_3(R!)}
       \le 3^{-\epsilon-2d-R+B+v_3(R!)},
\]
which tends to zero as $R\longrightarrow\infty$.
There are only finitely many residue/constant patterns at each $R$,
and their number causes no loss in a nonarchimedean norm.  These bounds
justify summing over all $R$, all derivative orders, and the bounded
coefficient extractions used below.  In particular the resulting
dimension dependence lies in the ordinary-coefficient Tate algebra
$A$, not merely in the larger ring of integer-valued functions.

Put $\epsilon=\lfloor2b/3\rfloor$.  For $k+d\ge1$,
\[
 \epsilon+2d+k-v_3((k+b+3d)!)\ge1.
\tag{5}
\]
Indeed, for $r\ge1$, $v_3(r!)\le(r-1)/2$; apply this to the three
values of $b$.  The resulting real lower bounds are respectively
$(k+d+1)/2$, $(k+d)/2$, and $(k+d+1)/2$, and the left side is an
integer.  For $k=d=0$, the denominator divides $b!$ with $b<3$.
It follows that the prefactor $3^{\epsilon+2d+k}$ absorbs the entire
dimension denominator $R!$, and retains one factor of $3$ whenever
$k+d\ge1$.  For fixed residue sum, the lower bound tends to infinity
as $R$ tends to infinity.  Hence the sum over all patterns converges
in the coefficient norm of $A$, including its scalar coefficients
and its $N^2$ error.

\begin{lemma}[Universal weak row]
\label{AP:lem:universal-weak-row}
Write $j=3L+b$, with $0\le b<3$.  There are
$D_{j,i}(T)\in A$, $0\le i\le L$, such that
\begin{align*}
 a_{T,j}&=\sum_{i\le L}D_{j,i}(T)a_{T,i},\\
 B_T(j,3N)-\sum_{i\le L}D_{j,i}(T)B_T(i,N)
                         &\in3N^2A,\\
 D_{j,i}(T)a_{T,i}&\in3(L!/i!)^2A\qquad(i<L).
\end{align*}
\end{lemma}
\begin{proof}
Split every index as $k_i=3\ell_i+r_i$.  If the residue carry is $d$,
the outside factorial factor is
\[
 3^{\epsilon+2d}U(2j)(2L+\epsilon)!2^{-2(L-d)}.
\]
After the balanced unit replacement, the operator is
\[
 P_{b,d}=\sum_{\sum r_i=b+3d}\prod_i(r_i!)^{-2}W_{r_i}(\theta_i)
\]
acting on $I_N(x_1)\prod_{i>1}J(x_i)$, followed by diagonal
specialization and extraction of $[x^{L-d}]$.
Use the active-slot reduction just proved, including the prefactor
$3^{\epsilon+2d}$.  Its scalar part is integral in $A$ and its
error is in $3N^2\mathcal H_A$: the case $k=d=0$ is a constant
operator and has no reduction error, and every other case has the
factor of $3$ in (5).

The balanced unit replacement also has an error in $3N^2A$
coefficientwise in $T$.  This point cannot be inferred merely by
interpolating its values at integer dimensions.  Instead use (1) at
the distinguished first slot.  It supplies $3N^2$ times an integral
restricted derivative operator.  With its residue $r_1$ fixed, the
remaining slots have residue sum $b+3d-r_1$.  Their active patterns
have $R\le k+b+3d-r_1\le k+b+3d$ and contribute at most the
denominator $R!$, without requiring scalar reduction.  The factor
$3^{1+\epsilon+2d+k}$ leaves at least one $3$ after this denominator,
by (5) and its $k=d=0$ case.  The universal factorial lattice makes
the remaining coefficient extraction integral.
More explicitly, an ordered pattern on these $R$ ordinary slots,
with derivative orders $e_q$, has error numerator
\[
 (T-1)_R\bigl(Z_{r_1}(\theta,N)I_N\bigr)
                     J^{T-1-R}\prod_{q=1}^R\theta^{e_q}J,
\]
apart from its displayed factors $3N^2$, $3^k$, and residue units.
This expression is in $\mathcal H_A$ with norm at most one before
division by $R!$, and the same convergence bounds apply.

Collecting the scalar operators defines the row $D_{j,i}(T)$.
At $N=0$ the balanced identity and all Bessel reductions are exact,
so the first identity follows in $A$.  Every row summand has the
factor $(2L+\epsilon)!/(2i)!$.  If $i<L$, it must either have a
positive carry $d$ or positive scalar $x$-degree.  The case $k=d=0$
has neither; thus (5) supplies one factor of $3$.  Therefore
\[
 D_{j,i}(T)\in3\frac{(2L+\epsilon)!}{(2i)!}A\quad(i<L).
\]
Multiplication by
$a_{T,i}=(2i)!b_T(i)/(i!)^2$ proves the boundary assertion, since
$(2L+\epsilon)!/(L!)^2$ is an integer.
\end{proof}

\paragraph{A finite choice of active patterns for each row.}
The preceding convergence argument can be shortened for the actual
row construction, and the rows can in fact be chosen in
$\mathbb Z_3[T]$.  Make the following truncation \emph{before}
scalar reduction.  At carry $d$ the coefficient being extracted is
$[x^{L-d}]$.  Every positive-degree factor of $W_r(\theta)-1$
contributes at least one power of $x$: for a $J$ factor this follows
from $\theta^a(1)=0$ for $a>0$, and for the distinguished $I_N$
factor its initial order is $N\ge1$ when $N>0$, whereas at $N=0$
it is again $J$.  Thus a pattern with $k>L-d$ differentiated slots
has zero extracted coefficient, simultaneously for every $N\ge0$
and every dimension.  The same conclusion holds directly in the
universal dimension formulation: the remaining power of $J$ has
constant term one and cannot lower these orders.

Retain only $k\le L-d$.  Then
\[
 R\le k+b+3d\le L+b+2d\le3L+b=j.
\]
There are therefore only finitely many active-slot patterns for a
fixed row.  The derivative orders within a pattern remain infinite,
but their restricted convergence is bounded as above, and their
scalar coefficients have degree at most $R\le j$ in $T$.
Their limit is consequently an ordinary polynomial over
$\mathbb Z_3$, after the already proved denominator compensation.
This gives $D_{j,i}(T)\in\mathbb Z_3[T]$ with degree at most $j$.
For the balanced error the fixed distinguished operator need not
have positive derivative order; applying the same cutoff to the
$k$ differentiated \emph{remaining} slots still gives a zero
coefficient when $k>L-d$, and bounds their active union by $R\le j$.
Thus even the infinite-$R$ limit is dispensable for each individual
row.  The weaker $A$-valued formulation of the lemma was retained
to avoid relying on this optional refinement.

\subsection{The coefficientwise one-digit loss}
\begin{proposition}[Coefficientwise weak descent]
If $3\mid K$, $N\ge K$, and $v=v_3(K)$, then
\[
 c_{K,N}(T)-\mathbf1_{3\mid N}c_{K/3,N/3}(T)
                       \in3^{2v-1}\mathbb Z_3[T].
\tag{6}
\]
\end{proposition}
\begin{proof}
Put $M=K/3$.  By (UW), multiplying the universal row error $3N^2A$
by $c_{M,N}(T)$ gives an error in $3M^2A$.  For $j\ge K$ the
boundary terms with $i<M\le L$ are also in $3M^2A$, because
$M\mid L!/i!$.  Thus
\[
 a_{T,j}\mathbf1_{j\ge K}
 -\sum_{N\ge M}c_{M,N}(T)B_T(j,3N)\in3M^2A.
\]
For $j<K$ both triangular expressions vanish.  Unit triangular
inversion over $\mathbb Z[T]$ now gives (6) in $A$.  Its left side
is an ordinary integer polynomial; divisibility in the coefficient
norm of $A$ is exactly divisibility of its ordinary coefficients,
so the assertion holds in $\mathbb Z_3[T]$.
\end{proof}

\begin{lemma}[The three residue representatives]
For dimensions $n=0,1,2$, every $3\mid K$ and $N\ge K$ satisfy
\[
 c_{K,N}(n)-\mathbf1_{3\mid N}c_{K/3,N/3}(n)
                         \in3^{2v_3(K)}\mathbb Z_3.
\tag{7}
\]
\end{lemma}
\begin{proof}
At $n=0$, $a_{0,j}=0$ for $j>0$, and hence $c_{K,N}(0)=0$ for
$K\ge1$ by triangular inversion.  At $n=1,2$, respectively,
\[
 B_n(j,N)=\binom{2j}{j-N}^{\rho},\qquad \rho=1,2.
\]
For completeness, these two matrices have a one-term row with
precision $9N^2$.  If $j=3L+b$ and
$\epsilon=\lfloor2b/3\rfloor$, the balanced unit-factorial
congruence gives
\[
 B_n(j,3N)\equiv d_j B_n(L,N)\pmod{9N^2\mathbb Z_3},
\]
where
\[
 d_j=\left(3^\epsilon(2L+1)^\epsilon
                          \frac{U(2j)}{U(j)^2}\right)^{\rho}.
\]
At $N=0$ this is the exact identity $a_{n,j}=d_j a_{n,L}$.
The needed balanced estimate is elementary and independent of the
weaker universal interpolation above.  If $H=3N$ and $j\ge H$, then
\[
 \frac{U(j)^2}{U(j-H)U(j+H)}
  =\prod_{\substack{j-H<i\le j\\3\nmid i}}(1+H/i)^{-1}
  \equiv1\pmod{H^2\mathbb Z_3}.
\]
Indeed the reciprocal sum over the interval is zero modulo
$3^{v_3(H)}$, because its units comprise complete copies of the
unit residue classes modulo $3^{v_3(H)}$, paired with their negatives;
all terms of order at least two already have $H^2$.
Now put $M=K/3$.  By (UW), multiplying the one-term row error by
$c_{M,N}(n)$ gives an error in $9M^2\mathbb Z_3=K^2\mathbb Z_3$.
There are no off-diagonal boundary terms.  The same unit-triangular
inversion as above proves (7) directly, without induction on the level.
\end{proof}

\begin{theorem}[Recovery of all levels at $3$]
For every integer dimension $n\ge1$, every
$3\mid K$, and every $N\ge K$,
\[
 c_{K,N}(n)-\mathbf1_{3\mid N}c_{K/3,N/3}(n)
                         \in3^{2v_3(K)}\mathbb Z_3.
\tag{8}
\]
Consequently, using the excellent Cartier coefficient congruence
\eqref{AP:eq:positive-cartier}, the truncation supercongruence
holds at $p=3$ for every dimension, every level, and every multiplier.
\end{theorem}
\begin{proof}
By (6), the polynomial
\[
 Q(T)=3^{1-2v_3(K)}
       \bigl(c_{K,N}(T)-\mathbf1_{3\mid N}c_{K/3,N/3}(T)\bigr)
\]
belongs to $\mathbb Z_3[T]$.  By (7), $Q(0),Q(1),Q(2)$ are all
zero modulo $3$.  An ordinary integral polynomial has the same
value modulo $3$ at integers in the same residue class; thus
$Q(n)\equiv0\pmod3$ for every integer $n$.  This proves (8).
Combine (8), the weighted value bound supplied by (UW), the exact
tail identity, and
\[
 C_{3N}-C_N\circ\sigma_n\in3^{2+2v_3(N)}\mathbb Z_3[[u]].
\]
Taking $K=m3^s$ gives the asserted truncation precision (indeed the
stronger $3^{2v_3(K)}$ precision when $3\mid m$).
\end{proof}

\begin{corollary}[The dimension Frobenius ideal]
Let $3\mid K$, let $N\ge K$, and put $v=v_3(K)$.  As an ordinary
polynomial in the dimension variable, the descent error satisfies
\[
 \boxed{\;
 c_{K,N}(T)-\mathbf1_{3\mid N}c_{K/3,N/3}(T)
 \in3^{2v-1}(T^3-T)\mathbb Z_3[T]
                     +3^{2v}\mathbb Z_3[T].\;}
\tag{9}
\]
\end{corollary}
\begin{proof}
The reduction modulo $3$ of the integral quotient $Q(T)$ in the
preceding proof vanishes at $0,1,2$.  Over $\mathbb F_3$, its three
distinct linear factors therefore divide $Q$:
\[
 T(T-1)(T-2)=T^3-T.
\]
Lift the polynomial quotient to $\mathbb Z_3[T]$ and multiply by
$3^{2v-1}$ to obtain (9).  Since $n^3-n$ is divisible by $3$ at
every integer dimension, the ideal on the right evaluates into
$3^{2v}\mathbb Z_3$.  This records exactly where the recovered
factor of $3$ resides.
\end{proof}

\subsection{Why the formulation is sharp}
The full coefficientwise analogue of (8) is false already at
$K=N=3$.  Direct triangular inversion gives
\[
 c_{1,1}(T)=2T,\qquad c_{3,3}(T)=80T-180T^2+120T^3,
\]
and hence
\[
 \frac{c_{3,3}(T)-c_{1,1}(T)}3\equiv T^3-T\pmod3.
\]
The quotient is a nonzero polynomial over $\mathbb F_3$ even though
it vanishes at all integer dimensions modulo $3$.  The recovery step
therefore proves full descent \emph{at all dimension values}; it does
not strengthen (6) to coefficientwise full descent.  In particular,
it cannot be recycled to prove (UW) without an additional argument.

\section{Completion of the formal-series congruence}\label{AP:source-main-6}\label{AP:sec:completion}
We spell out the last transfer to ensure that no coefficient cutoff or
level cutoff is implicit.  Fix a positive integer dimension $n$ and a
positive $K$ divisible by $3$, and write $M=K/3$ and $v=v_3(K)$.
The proved arithmetic descent gives
\[
 c_{K,N}(n)\equiv
    \mathbf1_{3\mid N}c_{M,N/3}(n)\pmod{3^{2v}}.
\]
Therefore the exact tail identity of Section~\ref{AP:sec:positive} yields
\[
 1-\frac{F_{n,K}}{F_n}
 \equiv\sum_{R\ge M}c_{M,R}(n)C_{3R}\pmod{3^{2v}}.
\]
By the quadratic cutoff theorem,
\[
 v_3(c_{M,R}(n))\ge
          \max\{0,2v_3(M)-2v_3(R)\}.
\]
The excellent Cartier error in replacing $C_{3R}$ by
$C_R\circ\sigma_n$ thus has valuation at least
\[
 2+2v_3(R)+\max\{0,2v_3(M)-2v_3(R)\}\ge2v.
\]
Every displayed sum is $u$-adically locally finite, since $C_R$ starts
in degree $R$.  It follows that
\[
 1-\frac{F_{n,K}}{F_n}
 \equiv\sum_{R\ge M}c_{M,R}(n)(C_R\circ\sigma_n)
 =1-\frac{F_{n,M}\circ\sigma_n}{F_n\circ\sigma_n}
                      \pmod{3^{2v}\mathbb Z_3[[u]]}.
\]
This proves Theorem~\ref{AP:thm:main}.  Both $F_n$ and
$F_n\circ\sigma_n$ have constant term $1$, so multiplying by them
preserves this coefficientwise ideal and gives the cross-multiplied
statement for $E_{n,3,s,m}$.

\begin{center}\small\renewcommand{\arraystretch}{1.2}
\begin{tabular}{p{53mm}p{39mm}p{51mm}}
\toprule
Previously remaining sector & Result here & Scope\\\midrule
$p=3$, $n\ge4$, $3\nmid n$, $m=1$ & All $s\ge5$ & Every coefficient; no dimension bound\\
$p=3$, $n\ge4$, $3\nmid n$, normalized $m$ & All $s,m\ge1$ & Precision $3^{2s+2v_3(m)}$\\
$p=3$, $3\mid n$ & Included in the same theorem & No extra argument or exclusion\\
\bottomrule
\end{tabular}
\end{center}

\section{An independent fifth-level certificate in every dimension}\label{AP:source-main-7}
\label{AP:sec:dimension-newton-level5}

The following finite calculation proves an entire formal-series
congruence at level five.  It is independent of any induction from
level five to higher levels.  Its theoretical inputs are the formal
excellent Cartier identity of Proposition~\ref{AP:prop:orbit}, the
negative-vertex formula, the derivative bound for $A_{2K}/A_K$, and
ordinary-polynomial finite differences in the dimension parameter.
No target truncation congruence is used to construct the lift.

\begin{theorem}\label{AP:thm:uniform-dimension-level5}
For every hyperoctahedral dimension $n\geq2$, with the normalized
excellent lift $\sigma_n$, one has
\[
 F_n(u)F_{n,81}(\sigma_n(u))-
 F_n(\sigma_n(u))F_{n,243}(u)
                  \in59049\mathbb Z_3[[u]].
\]
Thus every coefficient of the normalized multiplier-$1$ error at
$p=3$, $s=5$ is divisible by $3^{10}$.  The dimension is unrestricted.
\end{theorem}

\subsection{The two formal vertex inputs}
Here is a justification of the inputs used by the finite verifier.
Expand $1/(1-zg_n)$ in the chamber with dominant vertex $x_1$.
For $u=z^2$, the negative of its coefficient of $x_1^{-2K}$, with
all other exponents zero, is
\[
 P_{2K}(u)=u^{-K}A_{n,K}(u),\qquad A_{n,K}(0)=1.
\]
Indeed, write $S=\sum_{i=2}^n(x_i+x_i^{-1})$.  The identity
\[
 \frac1{1-zg_n}=-\frac1{zx_1}
 \frac1{1-(z^{-1}-S)x_1^{-1}+x_1^{-2}}
\]
and the recurrence $U_{m+1}(v/2)=vU_m(v/2)-U_{m-1}(v/2)$ give
\[
 P_{2K}=z^{-1}\operatorname{CT}_{x_2,\ldots,x_n}
 U_{2K-1}((z^{-1}-S)/2).
\]
Expanding
$U_m(v/2)=\sum_{\ell\geq0}(-1)^\ell\binom{m-\ell}{\ell}v^{m-2\ell}$
and using $\operatorname{CT}S^{2r}=\binom{2r}r b_{n-1}(r)$ yields
exactly the finite multinomial formula for $A_{n,K}$ below.
The excellent Cartier identity of
Proposition~\ref{AP:prop:orbit}, and the coefficient characterization of
second formal derivatives, give
\begin{equation}\label{AP:eq:vertex5}
 P_{2pK}(u)-\lambda_n(u)P_{2K}(\sigma_n(u))
        \in p^{2+2v_p(K)}\mathbb Z_p((u)),\qquad
 \lambda_n=F_n/F_n^{\sigma_n}.
\end{equation}
This reads one vertex coefficient after inverting $z$; each of the
displayed coefficients has bounded negative support.  No assertion
that all coefficients of a completed ambient ring have bounded
negative support is required.

We also need $\theta(A_{2K}/A_K)\in p^{v_p(K)}\mathbb Z_p[[u]]$,
where $\theta=u\,d/du$.  To see this, the rank-two quotient of the
same orbit module gives an operator
$\mathcal L=\theta^2-B_u\theta-A_u$ with integral coefficients.
Reading its relation at the vertex gives
$\mathcal L(P_{2K})\in p^{2v_p(K)}\mathbb Z_p((u))$.
Let $R_K=P_{2K}/F_n$, use the integral canonical coordinate $q$, and
put $\partial=q\,d/dq$, $H=\theta q/q$.  The period and logarithmic
solution of that quotient imply
\[
 H=W/F_n^2\in\mathbb Z_p[[u]]^\times,\qquad
 \theta H/H=B_u-2\theta F_n/F_n,\qquad
 \mathcal L(F_nR)=F_nH^2\partial^2R.
\]
The canonical-coordinate and Wronskian integrality were established
in the proof of Proposition~\ref{AP:prop:orbit}; the normalization from
$z$ to $u$ divides only by the unit $2$.  Hence
$\partial^2R_K\in p^{2v_p(K)}\mathbb Z_p((q))$.
If $r=v_p(K)$, the two facts
$[q^j]R_K\in\mathbb Z_p$ and $j^2[q^j]R_K\in p^{2r}\mathbb Z_p$
imply $j[q^j]R_K\in p^r\mathbb Z_p$ for every integer $j$.
Thus $\theta R_K\in p^r\mathbb Z_p((u))$.
Finally differentiate
$A_{2K}/A_K=u^KR_{2K}/R_K$: the term from $u^K$ contains $K$,
and the other two terms contain the just-proved factor $p^r$.
After their displayed powers of $u$ are removed, all remaining
factors are integral units.  This proves the derivative estimate.

\subsection{A directly cleared lift equation}
Put $\mathcal C_K=A_{2K}/A_K$.  For a general level $s\geq2$, set
\[
 k=2s-1,\qquad N=3^{2s-2},\qquad X=u^3,\qquad
 \sigma=X(1+3y).
\]
The negative-vertex Cartier identities, at indices $N$ and $2N$,
give the scaled equation
\[
 \Phi(y)=\frac{(1+3y)^N \mathcal C_{3N}(u)-\mathcal C_N(X(1+3y))}{3N}
                       =0\pmod{3^k}.
\]
For its positive-degree coefficients, the derivative estimate gives
$[h^r]C_N(X(1+h))\in(N/r)\mathbb Z_3[[u]]$; also
$\binom Nr\in(N/r)\mathbb Z_3$.  Division by $3N$ therefore makes
every positive-degree coefficient of $\Phi$ integral.  The constant
coefficient is integral because the actual integral $y$ satisfies
the unscaled equation modulo $3^{4s-2}$.  Its coefficient
of $y$ is a formal-series unit with constant term one; its coefficients
of $y^r$, $r\geq2$, are divisible by $3$.  More precisely, for $r\geq1$
the coefficient has valuation at least $r-1-v_3(r)$.

Multiplication by the integral unit
$A_{3N}(u)A_N(X(1+3y))$ gives a substantially smaller polynomial
calculation than repeatedly differentiating a quotient.  Define
\begin{equation}\label{AP:eq:direct-hensel5}
 H(y)=\frac{(1+3y)^N A_{6N}(u)A_N(X(1+3y))-
                  A_{3N}(u)A_{2N}(X(1+3y))}{3N}
       =\sum_{r\geq0}Q_r(u)y^r.
\end{equation}
If
\[
 a_i=[h^i]A_N(X(1+h)),\qquad
 b_i=[h^i]A_{2N}(X(1+h)),
\]
then these coefficients are calculated by the exact formula
\begin{equation}\label{AP:eq:direct-coeff5}
 Q_r=\frac{3^r}{3N}
 \left(A_{6N}\sum_{i=0}^r\binom N{r-i}a_i-A_{3N}b_r\right).
\end{equation}
In particular, $Q_0(0)=0$, $Q_1(0)=1$, and
$Q_r\in3\mathbb Z_3[[u]]$ for $r\geq2$.

For completeness, multiplication by the clearing unit changes the
Taylor tail bound.  Its coefficient of $y^i$ is divisible by $3^i$.
For $r\geq1$, a term involving the $j$th coefficient of $\Phi$,
$1\leq j\leq r$, therefore has valuation at least
$r-1-v_3(j)$.  The term from $j=0$ has valuation at least $r$.
Consequently
\begin{equation}\label{AP:eq:direct-tail5}
 v_3(Q_r)\geq r-1-\lfloor\log_3 r\rfloor\qquad(r\geq1).
\end{equation}
At $s=5$ the parameters are
\[
 k=9,\quad N=6561,\quad M_y=3^9=19683,\quad
 M=3^{10}=59049,\quad M_{\rm work}=3^{18}=387420489.
\]
The bound in \eqref{AP:eq:direct-tail5} is at least nine for every
$r\geq12$.  Hence the degree-$11$ equation
$\sum_{r=0}^{11}Q_r y^r=0\pmod{19683}$ is sufficient.  Formula
\eqref{AP:eq:direct-coeff5} divides by $3N=3^9$, so computing its
numerator modulo $3^{18}$ retains all nine required digits.

\subsection{Exact scaled Bessel convolutions}
The large vertex indices are
\[
 N=6561,\quad2N=13122,\quad3N=19683,\quad6N=39366.
\]
A dense matrix of multinomial coefficients is unnecessary.  Write
\[
 b_d(j)=\sum_{k_1+\cdots+k_d=j}\binom{j}{k_1,\ldots,k_d}^{\!2},
 \qquad
 J_3(x)=\sum_{j\geq0}\frac{3^j}{(j!)^2}x^j.
\]
Legendre's formula gives
\[
 v_3\!\left(\frac{3^j}{(j!)^2}\right)=s_3(j)\geq0,
\]
where $s_3(j)$ is the ternary digit sum.  Thus $J_3\in\mathbb Z_3[[x]]$
and
\[
 [x^j]J_3(x)^d=\frac{3^j b_d(j)}{(j!)^2}.
\]
The negative-vertex polynomial has the finite expression
\[
 A_{n,K}(u)=\sum_{j=0}^{K-1}\frac{u^j}{(2K-2j-1)!}
 \sum_{\ell=0}^j
 \frac{(-1)^\ell(2K-\ell-1)!}{\ell!(j-\ell)!^2}
                   b_{n-1}(j-\ell).
\]
Define
\[
 h_K=\max_{0\leq\ell<K}
 \{K-\ell-v_3((2K-\ell-1)!)+v_3(\ell!)\},
 \qquad
 \alpha_{K,\ell}=(-1)^\ell3^{\ell-K+h_K}
                     \frac{(2K-\ell-1)!}{\ell!}.
\]
Every $\alpha_{K,\ell}$ is $3$-integral by the definition of $h_K$.
The coefficient formula becomes one ordinary convolution:
\begin{equation}\label{AP:eq:scaled-vertex5}
 [u^j]A_{n,K}=
 \frac{3^{K-j-h_K}}{(2K-2j-1)!}
 [x^j]\left(\sum_{\ell=0}^{K-1}\alpha_{K,\ell}x^\ell\right)
                  J_3(x)^{n-1},\qquad0\leq j<K.
\end{equation}
If the factorial in front is separated into its unit part and its
power of $3$, the valuation of the prefactor is
\[
 K-j-h_K-v_3((2K-2j-1)!)
       =\frac{1+s_3(2K-2j-1)}2-h_K\geq1-h_K.
\]
Therefore convolution modulo $3^{E+H}$, with
$H=\max_Kh_K$, determines every required vertex coefficient modulo
$3^E$.  At level five $E=18$ and $H=10$, so the implementation uses
$3^{28}$.  All divisions by powers of $3$ are exact and are checked.
Only coefficients below the relevant finite vertex degree are
required in \eqref{AP:eq:scaled-vertex5}; this finite polynomial
calculation imposes no cutoff on the final formal-series assertion.

All convolutions use exact integer multiplication.  Coefficients
are packed in a base $B=2^{8w}$ satisfying
\[
 B>\min(a,b)(m-1)^2,
\]
where $a,b$ are the input lengths and $m$ is the modulus.  Thus no
convolution coefficient can carry into the next packed block.
GMP multiplies the resulting nonnegative integers exactly, after
which the blocks are reduced modulo $m$.  No floating-point
arithmetic or unchecked fixed-width multiplication is used.
The scaled formula was also checked against every coefficient of
the original factorial formula at several small vertex indices and
against the archived complete level-two and level-three vertex
certificates.

\subsection{The ordinary dimension model}
The number $b_n(j)$ counts ordered pairs of length-$j$ words in an
$n$-letter alphabet with the same letter multiplicities.  For words
using exactly $r$ prescribed letters, simultaneous relabelling by
$S_r$ acts freely.  Their number is therefore $r!c_{j,r}$ with
$c_{j,r}\in\mathbb Z$.  It follows that
$b_n(j)=\sum_r c_{j,r}(n)_r\in\mathbb Z[n]$.
Every factorial ratio in the vertex formula is an integer
multinomial coefficient, so $A_{n,K}\in\mathbb Z[n,u]$ as well.  Before reduction modulo $3^9$, formula
\eqref{AP:eq:direct-coeff5} has the uniform denominator bound
\[
 Q_r(T,u)\in3^{-9}\mathbb Z_3[T,u].
\]
The finite-difference fact needed here is elementary.  If
$P(T)\in3^{-a}\mathbb Z_3[T]$, then
\[
 \Delta^rP(c)\in3^{-a}r!\mathbb Z_3,
\]
since finite differences of monomials are multiples of $r!$.
Consequently, if $P\in\mathbb Z_3[T]$ takes values divisible by $3^e$
at $\nu_3(e)$ consecutive integers, it does so at every integer:
in the Newton expansion
$P(c+X)=\sum_r\Delta^rP(c)\binom Xr$, the initial differences are
integer combinations of those values and the later differences
are divisible by $3^e$ because of $r!$.  This holds coefficientwise
with an extra polynomial variable $u$.

Recall $\nu_3(e)=\min\{r:v_3(r!)\geq e\}$.  Here
\[
 \nu_3(18)=39,\qquad\nu_3(9)=21,\qquad\nu_3(10)=24.
\]
The calculation uses $39$ consecutive input dimensions
$n=2,\ldots,40$.  For every $0\leq r\leq11$ and $0\leq j<39$ it
checks every coefficient of
\begin{equation}\label{AP:eq:newton-check5}
 \Delta^jQ_r(2,u)\in
 3^{\min(9,v_3(j!))}\mathbb Z_3[u].
\end{equation}
For $j\geq39$, the ordinary-polynomial denominator bound makes the
Newton difference divisible by $3^9$ automatically.  The retained
differences of orders $0,\ldots,20$, divided by $j!$ and expressed in
the falling-factorial basis, give fixed polynomials
\[
 \widetilde Q_r(T,u)\in\mathbb Z_3[T,u],\qquad
 \deg_T\widetilde Q_r\leq20,
\]
with
\[
 \widetilde Q_r(n,u)\equiv Q_r(n,u)\pmod{3^9}
                \qquad\text{for every integer }n\geq2.
\]
All $39$ input evaluations of this model are additionally checked
against their original coefficients.

Put $B=\widetilde Q_1$ and define
\[
 y_0=0,\qquad
 y_{j+1}=-\frac{\widetilde Q_0+
               \sum_{r=2}^{11}\widetilde Q_r y_j^r}{B}
                              \qquad(0\leq j<9).
\]
This gives fixed finite polynomials $Y_*,Z_*\in\mathbb Z_3[T,u]$
with $y_9=Y_*/Z_*$.  For example one can take
$Z_*=B^{235794769}$, using $e_0=0$ and $e_{j+1}=11e_j+1$.
It is unnecessary to expand these enormous universal polynomials.
For each integer $n$, $B(n,u)$ is a unit and every specialized
nonlinear coefficient is divisible by $3$.  The iteration therefore
contracts by $3$ after specialization and gives
\[
 Y_*(n,u)/Z_*(n,u)\equiv y_n(u)\pmod{3^9}.
\]
This argument uses divisibility at integer dimensions; it does not
assume that every ordinary coefficient in the indeterminate $T$ of
$\widetilde Q_r$ is divisible by $3$.

\subsection{Twenty-four complete target identities}
For each $n=2,\ldots,25$, the program constructs a compact rational
root $Y_n/Z_n$ with $Z_n(0)$ a unit, and independently verifies
\[
 \sum_{r=0}^{11}Q_r(n,u)Y_n^rZ_n^{11-r}=0\pmod{3^9}.
\]
The digit construction is rational Hensel lifting in $\mathbb F_3(u)$.
It starts with $-Q_0/Q_1$ modulo $3$ and corrects one digit at a time.
Each residual division by the current power of $3$ is checked
exactly.  Rational cancellation is accelerated, when applicable,
by factoring the denominator into powers of $1+u$ and $1-u$ using
ternary-digit Taylor translation.  The final root identity makes
correctness independent of trusting that acceleration.

Let $F=F_{n,243}$, $F^-=F_{n,81}$ and $X=u^3$.  Define integral
polynomials, ordinary in $n$, by
\[
 \begin{aligned}
 L_r(n,u)={}&A_{n,3N}(u)
 \bigl([h^r]F^-(X(1+h))+N[h^{r-1}]F^-(X(1+h))\bigr)\\
 &-F(u)[h^r]A_{n,N}(X(1+h)),\qquad0\leq r\leq9.
 \end{aligned}
\]
where the coefficient with index $-1$ is zero.  This uses
$(1+3y)^N\equiv1+3Ny\pmod{3^{10}}$.
At every one of the $24$ dimensions, every coefficient of the
following finite polynomial is checked:
\begin{equation}\label{AP:eq:compact-target5}
 \sum_{r=0}^9L_r(n,u)(3Y_n)^rZ_n^{9-r}=0\pmod{59049}.
\end{equation}
The largest degree bound among these exact target calculations is
$141406$.  For example, dimension four has rational degrees
$(\deg Y_4,\deg Z_4)=(9865,9856)$ and target degree bound $108664$;
dimension five has degrees $(13548,10263)$ and target degree bound
$141397$.

Now use the single universal ordinary polynomial
\[
 \mathcal T_*(T,u)=\sum_{r=0}^9
 L_r(T,u)(3Y_*(T,u))^rZ_*(T,u)^{9-r}
                         \in\mathbb Z_3[T,u].
\]
Hensel uniqueness identifies the compact root and the universal
root modulo $3^9$ at each testing dimension.  Multiplication by $3$
improves that precision to $3^{10}$ in the target variable.  Since
all denominators have unit constant term,
\eqref{AP:eq:compact-target5} gives
$\mathcal T_*(n,u)\equiv0\pmod{3^{10}}$ for $n=2,\ldots,25$.
The finite-difference criterion just proved, applied coefficientwise
to this fixed ordinary polynomial, propagates the congruence from
$24=\nu_3(10)$ consecutive dimensions to every $n\geq2$.
Restoring the unit denominators and the negative-vertex expression
for the Cartier multiplier proves
$\lambda_nF_{n,81}(\sigma_n)-F_{n,243}\in3^{10}\mathbb Z_3[[u]]$.
Multiplication by the unit $F_n(\sigma_n)$ proves
Theorem~\ref{AP:thm:uniform-dimension-level5}.

\subsection{Reproduction and archived data}
The complete source programs are printed in Appendix~\ref{app:level5-programs}.
They can be saved in \texttt{research/octa\_level5}.  A reproduction
which recalculates even previously cached data is
\begin{verbatim}
python3 scaled_convolution.py
python3 linear_reduce.py
python3 direct_certificate.py --test
python3 uniform_level.py --s 5 --stage all --recompute
\end{verbatim}
The scripts import the elementary polynomial and earlier certificate
modules in the sibling \texttt{octa\_new\_certificates} directory.
The archive contains $39$ exact input files, the Newton differences
and ordinary model, $24$ full compact certificates, and a summary
with \texttt{power\_series\_cutoff: null}.  The aggregate identifier
of the model and the $24$ target certificates is
\begin{center}\ttfamily\small
013f86b4dc5cb189b9be96b71c0eac752cffb93dbf0428d0f04b3f50674b0529.
\end{center}
The identifier records the archived objects; the proof consists of
the exact algebraic reductions and the complete finite polynomial
identities, not the identifier itself.

\clearpage
\section{Verification and exact scope}\label{AP:source-main-8}\label{AP:sec:scope-end}
The main theorem is an algebraic all-level proof.  It depends on the
formal excellent Cartier construction in Section~\ref{AP:sec:formal-input},
the exact cutoff kernel, the universal coefficientwise weak row, and
the scalar residue representatives.  It does not assume that a finite
set of tested levels controls subsequent levels.  In particular, the
transition from $s=5$ to arbitrary $s$ uses no extrapolation.

The ordinary-polynomial qualification is essential.  Integer-valued
polynomials alone would not justify reduction of dimensions modulo $3$.
Conversely, the final descent is asserted at every integer dimension:
its full-precision coefficientwise version in the indeterminate $T$
is false, as the explicit $T^3-T$ example demonstrates.

The formal base ring of this part is $\mathbb Z_3[[u]]$.
Corollary~\ref{G:cor:analytic} supplies the transfer to the standard
analytic ring at BV3-admissible primes. An analytic realization at
every additional small prime, evaluation at a specified fixed point,
and identification with a geometric unit root require their own
analytic and geometric verification.  The prime $2$ is outside the
Cartier input used here.  This manuscript and its exact certificates
have been checked internally; they have not undergone external peer
review or proof-assistant formalization.



% END RETAINED SOURCE AP

\partstart{Modular Hecke traces and stronger congruences}{part:modular}{The complete modular framework, thirteen cusp-type families, elliptic zeros, and a primal/dual Laurent proof for Zagier D.}
% BEGIN RETAINED SOURCE B
\section{Modular trace methods and their conclusions}\label{B:source-main-1}

\subsection{Dwork congruences and the excellent lift}
For a power series $F(t)=\sum_{n\ge0}a_nt^n\in1+t\Z[[t]]$ and an integer $M\ge1$ we write
$F_M(t)=\sum_{n<M}a_nt^n$ for its truncation. Let $p$ be a prime. A \emph{Frobenius lift} is a
power series $t^\sigma\in t\,\Z_p[[t]]$ with $t^\sigma\equiv t^p\pmod p$. When $a_n$ is the
constant term of the $n$-th power of a Laurent polynomial, Mellit--Vlasenko \cite{MV16} and
Beukers--Vlasenko \cite{BV1,BV3} proved that for every Frobenius lift
\begin{equation}\label{B:eq:dwork}
\frac{F(t)}{F(t^\sigma)}\equiv\frac{F_{mp^s}(t)}{F_{mp^{s-1}}(t^\sigma)}\pmod{p^s}
\qquad(m,s\ge1)
\end{equation}
as power series in $t$. For the lift $t\mapsto t^p$ and the Legendre family this goes back to
Dwork~\cite{Dw69}. In the situations of interest $F$ satisfies a differential equation with maximal
unipotent monodromy at $t=0$. Let $F(t)\log t+G(t)$ be its second Frobenius solution and let
$q(t)=t\exp(G(t)/F(t))$ be the canonical coordinate; its inverse $t(q)$ is the mirror map. Dwork
singled out the \emph{excellent} Frobenius lift, which in these terms is
\[
t^\sigma=t\bigl(\gamma^{p-1}q(t)^p\bigr)\qquad\text{(\cite[Def.~7.14]{BV3})},
\]
where $\gamma$ is the vertex coefficient of the Laurent polynomial; in all cases below
$\gamma=1$, so $t^\sigma(t(q))=t(q^p)$. Beukers and Vlasenko \cite[Thm.~7.3]{BV3} showed that for
completely symmetric Calabi--Yau families $f=1-tg(x)$ the excellent lift is the unique Frobenius
lift for which the Cartier operator acts on $1/f$ by the scalar $F(t)/F(t^\sigma)$ modulo
$p^2F_2$ rather than only modulo $pF_1$. They then wrote: \emph{``Our second main objective was to
show that (19) holds modulo $p^{2s}$ when $t\mapsto t^\sigma$ is an excellent Frobenius lift.
Unfortunately we were unable to prove this for any example of $g$.''} Their (19) is our
\eqref{B:eq:dwork}, and they formulated the following conjecture (see also the survey \cite{Vl24}).

\begin{bvconj}
Congruence \eqref{B:eq:dwork} with the excellent lift holds modulo $p^{2s}$ for $m=2$ in the
hypercubic and hyperoctahedral families and for $m=n+1$ in the simplicial families.
\end{bvconj}

Here $n$ is the number of variables, $F(t)=\sum_k\CT(g^k)t^k$, and
\[
g_{\rm hypercubic}=\prod_{i=1}^n\Bigl(x_i+\frac1{x_i}\Bigr),\quad
g_{\rm hyperoctahedral}=\sum_{i=1}^n\Bigl(x_i+\frac1{x_i}\Bigr),\quad
g_{\rm simplicial}=x_1+\dots+x_n+\frac1{x_1\cdots x_n}.
\]
The primes are those of \cite[Thm.~7.3]{BV3}: $p\nmid \#G\cdot[\Z^n:\Gamma_g]$, with $G$ the
symmetry group of $g$ and $\Gamma_g$ the lattice spanned by its support. Beukers and Vlasenko also
observe \cite[\S8]{BV3} that for $n=2$ the hypercubic family is modular: its mirror map satisfies
$t(q)^2=\lambda(q^2)/16$.

Beukers \cite{Beu26} proves first-level congruences modulo $p^2$ or $p^3$ for many modular families and asks about their extension to higher levels. The trace inductions in this part supply the higher-level statements under the hypotheses stated below. Fu \cite{Fu26} studies a cohomological criterion for supercongruences associated with rank-two factors and excellent-lift fixed points; that criterion is not an input to any proof in this volume.

\subsection{Results}
The precise hypotheses are collected in Definition~\ref{B:def:cusp}. Informally, a modular family of
\emph{cusp type} consists of a Hauptmodul $t\in q+q^2\Z[[q]]$ for a genus-zero group
$\Gamma\supseteq\Gamma_0(N)$ and a holomorphic modular form $\Theta\in1+q\Z[[q]]$ of weight
$w\in\{1,2\}$ which has no zeros in the upper half plane and vanishes only at the cusps where $t$ has its
pole. Let $F\in1+t\Z[[t]]$ be defined by $F(t(q))=\Theta(q)$ and let
$t^\sigma\in t^p\Z[[t]]$ be the excellent lift, $t^\sigma(t(q))=t(q^p)$.

\begin{theorem}[Main Theorem]\label{B:thm:main}
Let $(N,\Gamma,w,\chi,\Theta,t)$ be a modular family of cusp type and let $p\nmid N$ be a prime for
which the Hecke condition $(\mathrm C_p)$ of Definition~\ref{B:def:cusp} holds (by
Lemma~\ref{B:lem:Cp} this is automatic in all examples of Table~\ref{B:tab:fam}). Then for all
integers $m,s\ge1$
\[
\frac{F(t)}{F(t^\sigma)}\equiv\frac{F_{mp^s}(t)}{F_{mp^{s-1}}(t^\sigma)}\pmod{p^{(w+1)s}\Z_p[[t]]}.
\]
Equivalently,
$\Theta(q)\,F_{mp^{s-1}}\bigl(t(q^p)\bigr)\equiv\Theta(q^p)\,F_{mp^s}\bigl(t(q)\bigr)
\pmod{p^{(w+1)s}\Z[[q]]}$.
\end{theorem}

\begin{table}[ht]
\centering\small
\renewcommand{\arraystretch}{1.25}
\begin{tabular}{@{}llllllll@{}}
\toprule
family & $w$ & $\Gamma$ & $\chi$ & $\Theta$ & $t$ & cusps & $r_\Theta$\\\midrule
Legendre, $\binom{2n}{n}^2$ & 1 & $\Gamma_0(4)$ & $\chi_{-4}$ & $\eta_2^{10}/\eta_1^4\eta_4^4$ &
  $\eta_1^8\eta_4^{16}/\eta_2^{24}$ & $\frac12$ & $\frac12$\\
Zagier A (Franel) & 1 & $\Gamma_0(6)$ & $\chi_{-3}$ & $\eta_2\eta_3^6/\eta_1^2\eta_6^3$ &
  $\eta_1^3\eta_6^9/\eta_2^3\eta_3^9$ & $\frac13$ & 1\\
Zagier B & 1 & $\Gamma_0(9)$ & $\chi_{-3}$ & $\eta_1^3/\eta_3$ & $\eta_9^3/\eta_1^3$ & $0$ & 1\\
Zagier C & 1 & $\Gamma_0(6)$ & $\chi_{-3}$ & $\eta_2^6\eta_3/\eta_1^3\eta_6^2$ &
  $\eta_1^4\eta_6^8/\eta_2^8\eta_3^4$ & $\frac12$ & 1\\
Zagier E & 1 & $\Gamma_0(8)$ & $\chi_{-4}$ & $\eta_2^{10}/\eta_1^4\eta_4^4$ &
  $\eta_1^4\eta_4^2\eta_8^4/\eta_2^{10}$ & $\frac12$ & 1\\
Zagier F & 1 & $\Gamma_0(6)$ & $\chi_{-3}$ & $\eta_1^6\eta_6/\eta_2^3\eta_3^2$ &
  $\eta_2\eta_6^5/\eta_1^5\eta_3$ & $0$ & 1\\\midrule
$\binom{2n}n^3$ & 2 & $\Gamma_0(4)+4$ & 1 & $\eta_2^{20}/\eta_1^8\eta_4^8$ &
  $\eta_1^{24}\eta_4^{24}/\eta_2^{48}$ & $\frac12$ & $\frac12$\\
Ap\'ery $(\gamma)$ & 2 & $\Gamma_0(6)+6$ & 1 & $\eta_2^7\eta_3^7/\eta_1^5\eta_6^5$ &
  $\eta_1^{12}\eta_6^{12}/\eta_2^{12}\eta_3^{12}$ & $\frac12,\frac13$ & 1\\
Almkvist $(\eta)$ & 2 & $\Gamma_0(5)$ & $(\frac{\cdot}{5})$ & $\eta_1^5/\eta_5$ &
  $\eta_5^6/\eta_1^6$ & $0$ & 1\\
Domb $(\alpha)$ & 2 & $\Gamma_0(6)+3$ & 1 & $\eta_1^4\eta_3^4/\eta_2^2\eta_6^2$ &
  $\eta_2^6\eta_6^6/\eta_1^6\eta_3^6$ & $0,\frac13$ & 1\\
Almkvist--Zudilin $(\delta)$ & 2 & $\Gamma_0(6)+2$ & 1 & $\eta_1^3\eta_2^3/\eta_3\eta_6$ &
  $\eta_3^4\eta_6^4/\eta_1^4\eta_2^4$ & $0,\frac12$ & 1\\
Almkvist $(\varepsilon)$ & 2 & $\Gamma_0(8)+8$ & 1 & $\eta_2^6\eta_4^6/\eta_1^4\eta_8^4$ &
  $\eta_1^8\eta_8^8/\eta_2^8\eta_4^8$ & $\frac12,\frac14$ & 1\\
Almkvist $(\zeta)$ & 2 & $\Gamma_0(9)+9$ & 1 & $\eta_3^{10}/\eta_1^3\eta_9^3$ &
  $\eta_1^6\eta_9^6/\eta_3^{12}$ & $\frac13,\frac23$ & 1\\\bottomrule
\end{tabular}
\caption{Modular families of cusp type, with $\eta_\delta=\eta(\delta\tau)$. The data are those of
\cite[App.~B, C]{Beu26}; ``cusps'' lists the $\Gamma_0(N)$-cusps in $\Gamma c$, where $\Theta$
vanishes and $t$ has its pole. The coefficients $a_n$ are listed in \S\ref{B:ss:examples}.}
\label{B:tab:fam}
\end{table}

The hypothesis that $\Theta$ vanishes only at cusps excludes several natural families. Our second and third results treat
the most prominent ones.

\begin{theorem}\label{B:thm:ell}
Let $F$ be one of
\[
F_{27}(t)=\sum_{n\ge0}\frac{(3n)!}{n!^3}t^n,\qquad F_{256}(t)=\sum_{n\ge0}\frac{(4n)!}{n!^4}t^n,\qquad
F_{108}(t)=\sum_{n\ge0}\frac{(2n)!\,(3n)!}{n!^5}t^n,
\]
that is ${}_2F_1(\frac13,\frac23;1;27t)$, ${}_3F_2(\frac14,\frac12,\frac34;1,1;256t)$ and
${}_3F_2(\frac13,\frac12,\frac23;1,1;108t)$, with $w=1,2,2$ respectively, and let $t^\sigma$ be the excellent lift for the
modular parametrisation of Table~\ref{B:tab:ell}. Then for every prime $p\ge5$ and all $m,s\ge1$
\[
\frac{F(t)}{F(t^\sigma)}\equiv\frac{F_{mp^s}(t)}{F_{mp^{s-1}}(t^\sigma)}\pmod{p^{(w+1)s}\Z_p[[t]]}.
\]
\end{theorem}

Here $\Theta$ vanishes at an elliptic point of $\Gamma$, and the trace of $Q_P$ is a rational function of $t$ with poles at CM
points of conductor $p$. For $s=m=1$, Theorem~\ref{B:thm:ell} is due to Beukers \cite[Thms.~1.4, 1.13 and Cor.~6.2]{Beu26}.

\begin{theorem}\label{B:thm:D}
Let $a_n=\sum_k\binom nk^2\binom{n+k}k$ be the Ap\'ery numbers for $\zeta(2)$ (Zagier's case D), $F(t)=\sum a_nt^n$, and
let $t^\sigma$ be the excellent lift, $t^\sigma(t(q))=t(q^p)$ with $t(q)=q\prod_{n\ge1}(1-q^n)^{5(\frac n5)}$. Then for
every prime $p\ne5$ and all $m,s\ge1$
\[
\frac{F(t)}{F(t^\sigma)}\equiv\frac{F_{mp^s}(t)}{F_{mp^{s-1}}(t^\sigma)}\pmod{p^{2s}\Z_p[[t]]}.
\]
\end{theorem}

This family lives on $\Gamma_1(5)$, and for $p\equiv\pm2\pmod5$ the Hecke operator involves the diamond operator
$\langle2\rangle$, which does not fix $t$. Even the case $s=m=1$ is not covered by \cite[Thm.~1.4]{Beu26}; Beukers
remarks that computer experiments support it \cite[Remark~1.3]{Beu26}. In Theorems~\ref{B:thm:ell} and \ref{B:thm:D} the
modular $t(q)$ is the normalized mirror map of the differential equation of $F$, as in the parametrizations of \cite[Appendices B--C]{Beu26}; hence $t^\sigma$ is the excellent lift. The exact coordinate normalization is treated in Lemma~\ref{B:lem:subst}; checks~[8g] and [9a] are supplementary finite checks.

\begin{corollary}[Conjecture BV in dimensions two and three]\label{B:cor:BV}
Conjecture BV holds for the hypercubic and the hyperoctahedral families with $n=2$ (for every odd prime $p$), for the
hypercubic family with $n=3$ (for every prime $p\ge5$), and for the simplicial families with $n=2$ and $n=3$ (for every
prime $p\ge5$). For $n=3$ the hypercubic and simplicial congruences even hold modulo $p^{3s}$. For every family and dimension explicitly listed in this corollary, the same assertion holds for every original truncation multiplier divisible by $2$ in the hypercubic or hyperoctahedral case, or by $n+1$ in the simplicial case. The exponent is $2s$ in dimension two and $3s$ for the hypercubic and simplicial families in dimension three.
\end{corollary}

The three-dimensional hyperoctahedral family has the $2s$ bound of Theorem~\ref{P:thm:main-all}; no $3s$ assertion for that family is made here. The same theorem supplies the uniform result in every dimension. Section~\ref{B:ss:hyp3} explains why the termwise modular estimate alone does not supply it.

\begin{corollary}[Beukers's question]\label{B:cor:Beu}
For every family of Table~\ref{B:tab:fam}, every prime $p\nmid N$ and all $m,s\ge1$, the excellent
lift congruence holds modulo $p^{2s}$ in weight one and modulo $p^{3s}$ in weight two. This
covers every family of \cite[App.~B, C]{Beu26} that satisfies Beukers's Assumptions (I) or (II)
with $\Theta$ an $\eta$-product.
\end{corollary}

\begin{corollary}\label{B:cor:coef}
For every family of Table~\ref{B:tab:fam}, $p\nmid N$ and $m,s\ge1$,
$a_{mp^s}\equiv a_{mp^{s-1}}\pmod{p^{(w+1)s}}$.
\end{corollary}

Corollary~\ref{B:cor:coef} is not new. For the Ap\'ery numbers it is Coster's theorem
\cite{Co88} (conjectured in \cite{Beu85}). For the sporadic sequences these two-term
supercongruences were conjectured by Osburn, Sahu and Straub \cite{OSS16} and proved through the
work of several authors, the last cases in \cite{St23,AO23}. Here it comes out of the same
argument as the power series congruences.

\begin{remark}
For $s=m=1$, Theorem~\ref{B:thm:main} is Beukers's Theorems~1.4 and 1.13 \cite{Beu26} in the
$\eta$-product cases. Our proof of the case $s=1$ is a trace version of his argument. It
needs only that the zeros of $\Theta$ are cusps, not the $\eta$-product identity
$\Theta(p\tau)\prod_k\Theta(\frac{\tau+k}p)=c\,\Theta(\tau)^{p+1}$ of \cite[Lemma~3.2]{Beu26}.
\end{remark}

\subsection{Method}
Let $P\in\Z[T]$ be a polynomial and put
\[
Q_P(\tau)=\frac{\Theta(\tau)\,P(t(p\tau))}{\Theta(p\tau)},\qquad
G_P(z)=\frac{\Theta(z^p)\,P(t(z))}{\Theta(z)}\in\Z[[z]].
\]
The function $Q_P$ is modular for $\Gamma_0(Np)$, and its $p+1$ conjugates under $\Gamma$ are
indexed by the Hecke cosets of level $p$. The conjugate for $\tau\mapsto p\tau$ is $Q_P$ itself.
The other $p$ conjugates, for $\tau\mapsto(\tau+k)/p$, add up to $p^{w+1}\chi(p)\,U_pG_P$, where
$U_p$ extracts the coefficients of $q^{pn}$. Hence the trace is
\begin{equation}\label{B:eq:key}
\Tr(Q_P)=Q_P+p^{w+1}\chi(p)\,U_pG_P .
\end{equation}
For a family of cusp type the trace has no poles in $\Hb$, so it is a polynomial
$\Pi_P(t)\in\Z[t]$. A local computation at the cusp bounds its degree by $p\deg P+(p-1)r_\Theta$
(Proposition~\ref{B:prop:poly}). Now take $P=F_{mp^{s-1}}$. The induction hypothesis says
$G_P\equiv F_{mp^{s-2}}(t(z^p))$ modulo $p^{(w+1)(s-1)}$, so
$p^{w+1}U_pG_P\equiv p^{w+1}F_{mp^{s-2}}(t)$ modulo $p^{(w+1)s}$. Thus $Q_P$ is congruent modulo
$p^{(w+1)s}$ to a polynomial in $t$ of degree $<mp^s$. Since $Q_P\equiv F(t)$ modulo $t^{mp^s}$,
this polynomial is $F_{mp^s}(t)$, and that is the theorem. In the language of Hecke operators on forms of
weight $-w$ the induction is the three-term recursion
\[
h_s\equiv p^{w+1}\chi(p)\bigl(T_ph_{s-1}-h_{s-2}\bigr)\pmod{p^{(w+1)s}},\qquad
h_s:=F_{mp^s}(t)/\Theta
\]
(Remark~\ref{B:rem:hecke}), in the spirit of Atkin--Swinnerton-Dyer congruences.

In the language of Dwork crystals, the excellent lift is the branch $\tau\mapsto p\tau$ of the
Hecke correspondence; Beukers--Vlasenko characterise it as the lift for which the off-diagonal
Cartier coefficient $\lambda_1$ vanishes. The remaining branches carry an explicit factor $p^w$,
and this factor doubles or triples the exponent of \eqref{B:eq:dwork}.

\emph{Elliptic zeros.} For the families of Theorem~\ref{B:thm:ell}, $\Theta$ vanishes at an elliptic point $e$ of order
$k_0$, like $t^{-r}$, and the trace has poles at the Hecke translates of $e$. We write the trace of $Q_P$ for
$P=F_{mp^{s-1}}=\sum_ia_it^i$ as $\sum_ia_iS_i$, where $S_i$ is the trace of $\Theta(\tau)t(p\tau)^i/\Theta(p\tau)$. The
functions $S_i$ satisfy a linear recursion in $i$ whose coefficients are the symmetric functions of the $p+1$ Hecke
translates $t\circ\beta$. These are rational functions in $t$ whose denominators are powers of a polynomial $D\equiv\pm1\pmod p$,
essentially Beukers's polynomial $T_{p,p+1}$ \cite{Beu26}. A Gauss norm estimate on the disc $|t|_p<p^{k_0/(p+1)}$ gives
\[
v_p\bigl([t^n]S_i\bigr)\ge\frac{k_0}{p+1}\bigl(n-pi-\lfloor pr\rfloor\bigr)
\]
(Proposition~\ref{B:prop:decay}). So the coefficients of $t^n$, $n\ge mp^s$, of the trace come essentially from monomials $t^i$
near the top of the truncation. For those, the base-$p$ digits of $i$ are $p-1$, and by Legendre's formula each such digit
contributes $w+1$ factors $p$ to $a_i$. The two effects balance exactly, and the induction goes through.

\emph{Zagier's case D.} Here $\Gamma=\Gamma_1(5)$; $\Theta$ vanishes at the cusp $\frac25$, where $t$ has its pole. For
$p\equiv\pm1\pmod5$ the argument for cusp type applies. For $p\equiv\pm2$ the Hecke correspondence on $X_1(5)$ contains the
diamond operator $\langle2\rangle$, which maps $t$ to $-1/t$ and $\Theta$ to $\pm t\Theta$. The trace of $Q_P$ is then a
Laurent polynomial in $t$, but $Q_P$ is a branch of the trace for the dual family $(t\Theta,-1/t)$, and conversely. We
prove the congruence and its dual together, by a joint induction (\S\ref{B:sec:D}).

\subsection{Limits of the termwise modular method}
The hyperoctahedral family with $n=3$ has a modular parametrisation with an elliptic zero on $\Gamma_0(6)+$, with $1/r=2<w+1=3$. The termwise decay estimate below is insufficient to recover the required cancellation (Section~\ref{B:ss:hyp3}). The full $2s$ bound follows from Theorem~\ref{P:thm:main-all}, and its sharpness is proved in Theorem~\ref{E:thm:sh-octa3}. The uniform higher-dimensional argument also comes from Part~\ref{part:positive}; it does not require modular parametrisations of the higher-dimensional periods. The broader elliptic-zero assertion in Conjecture~\ref{B:conj:ell} remains a separate extension beyond the listed families.

\subsection*{Organisation}
\S\ref{B:sec:setup} sets up notation and hypotheses; \S\ref{B:sec:trace} proves the trace identity and the
polynomiality; \S\ref{B:sec:proof} proves Theorem~\ref{B:thm:main}; \S\ref{B:sec:cor} proves the
corollaries; \S\ref{B:sec:ell} proves Theorem~\ref{B:thm:ell}; \S\ref{B:sec:D} proves Theorem~\ref{B:thm:D};
\S\ref{B:sec:open} discusses the limits of the modular method and a broader extension; \S\ref{B:sec:verif} describes the verification script, which is
reproduced in the appendix.

\section{Modular families of cusp type}\label{B:source-main-2}\label{B:sec:setup}

\subsection{Notation}
Let $\Hb$ be the upper half plane and $q=e^{2\pi i\tau}$. For $\gamma=\bmat abcd\in\GL_2^+(\Q)$ and a
function $f$ on $\Hb$ put $(f|_w\gamma)(\tau)=\det(\gamma)^{w/2}(c\tau+d)^{-w}f(\gamma\tau)$, so that
$(f|_w\gamma_1)|_w\gamma_2=f|_w(\gamma_1\gamma_2)$ and $f|_w(\lambda I)=\operatorname{sgn}(\lambda)^wf$
for $\lambda\in\Q^\times$; we write $j(\gamma,\tau)=c\tau+d$ and $T=\bmat1101$. For a power series
$h=\sum h_nq^n$ let $U_ph=\sum h_{pn}q^n$. If $h(z)\in\C[[z]]$ and $q^{1/p}:=e^{2\pi i\tau/p}$,
$\zeta_p=e^{2\pi i/p}$, then
\begin{equation}\label{B:eq:Up}
\sum_{k=0}^{p-1}h\bigl(\zeta_p^kq^{1/p}\bigr)=p\,(U_ph)(q),\qquad U_p\bigl(h_1(z^p)h_2(z)\bigr)=h_1(q)\,U_ph_2(q).
\end{equation}
For a cusp $c$ and $A\in\SL_2(\Z)$ with $A\infty=c$, and a function $f$ which is meromorphic at
$c$ (of weight $0$, or of weight $w$), the function $f\circ A$, respectively $f|_wA$, has an
expansion $\sum_{\nu}b_\nu e^{2\pi i\nu\tau}$ with $\nu$ in a lattice $\frac1H\Z$ bounded from below.
We write $\nu_A(f)=\min\{\nu:b_\nu\ne0\}$. Clearly $\nu_A(f_1f_2)=\nu_A(f_1)+\nu_A(f_2)$ and
$\nu_A(f_1+f_2)\ge\min(\nu_A(f_1),\nu_A(f_2))$, and $\nu_A$ does not change when $A$ is replaced by
$\pm AT^x$ or, for $f$ modular for $\Gamma_0(N)$, by $\gamma A$ with $\gamma\in\Gamma_0(N)$.

\subsection{Hecke cosets}
Let $p\nmid N$ and
$M(N,p)=\bigl\{\bmat ab{Nc}d\in M_2(\Z):ad-Nbc=p\bigr\}$; let $M_p$ be the set of all integral
matrices of determinant $p$.

\begin{lemma}\label{B:lem:hecke}
\begin{enumerate}[label=\textup{(H\arabic*)}]
\item $M(N,p)=\Gamma_0(N)\bmat100p\Gamma_0(N)$, and $\Gamma_0(N)\backslash M(N,p)$ has the $p+1$
representatives $\beta_\infty=\bmat p001$ and $\beta_k=\bmat1k0p$, $0\le k\le p-1$.
\item The map $\Gamma_0(N)\beta\mapsto\SL_2(\Z)\beta$ is a bijection
$\Gamma_0(N)\backslash M(N,p)\to\SL_2(\Z)\backslash M_p$.
\item $M(N,p)\gamma=M(N,p)$ for $\gamma\in\Gamma_0(N)$; right multiplication permutes the cosets.
\item If $g$ is an Atkin--Lehner matrix for $\Gamma_0(N)$, then $g^{-1}M(N,p)g=M(N,p)$, and
$\beta\mapsto g^{-1}\beta g$ permutes $\Gamma_0(N)\backslash M(N,p)$.
\item $\adj(\beta)=p\beta^{-1}\in M(N,p)$ for $\beta\in M(N,p)$.
\end{enumerate}
\end{lemma}
\begin{proof}
(H1) is standard, see \cite[\S5.2]{DS05} and \cite[\S3]{Beu26}. Note that
$\gcd(a,N)=\gcd(d,N)=1$ holds automatically for $\bmat ab{Nc}d\in M(N,p)$. For (H2): by Hermite normal form every element of $M_p$ is
$\SL_2(\Z)$-equivalent to exactly one $\bmat{a'}{b'}0{d'}$ with $a'd'=p$, $0\le b'<d'$. These
are the $p+1$ matrices of (H1), so the map is surjective between sets of the same cardinality.
(H3) and (H5) are clear. (H4) goes back to Atkin--Lehner \cite{AL70}; for completeness let $g=\bmat{ea}b{Nc}{ed}$ with $e\,\|\,N$ and $\det g=e$. A
direct computation gives
$g^{-1}\bmat100pg=\bmat{ead-p(N/e)bc}{(1-p)bd}{(p-1)Nac}{pead-(N/e)bc}\in M(N,p)$. Since $g$
normalises $\Gamma_0(N)$, (H1) gives $g^{-1}M(N,p)g\subseteq M(N,p)$; the same holds for $g^{-1}$,
which is a scalar multiple of an Atkin--Lehner matrix, so equality holds, and
$g^{-1}\Gamma_0(N)\beta g=\Gamma_0(N)g^{-1}\beta g$.
\end{proof}

For a Dirichlet character $\chi$ modulo $N$ with values in $\{\pm1\}$ and
$\beta=\bmat ab{Nc}d\in M(N,p)\cup\Gamma_0(N)$ we put $\chi(\beta)=\chi(d)$. Then
$\chi(\gamma\beta)=\chi(\gamma)\chi(\beta)=\chi(\beta\gamma)$ for $\gamma\in\Gamma_0(N)$, because the
lower right entries multiply modulo $N$.

\subsection{The hypotheses}

\begin{definition}\label{B:def:cusp}
A \emph{modular family of cusp type} is a tuple $(N,\Gamma,w,\chi,\Theta,t)$ with the following
properties.
\begin{enumerate}[label=\textup{(F\arabic*)}]
\item $w\in\{1,2\}$. Either $\Gamma=\Gamma_0(N)$ and $\chi$ is a character modulo $N$ with values
in $\{\pm1\}$ and $\chi(-1)=(-1)^w$; or $\Gamma=\Gamma_0(N)+S$ is an Atkin--Lehner extension and
$\chi$ is trivial.
\item $\Theta\in1+q\Z[[q]]$ is a holomorphic modular form of weight $w$ for $\Gamma_0(N)$ with
character $\chi$, i.e.\ $\Theta|_w\gamma=\chi(\gamma)\Theta$ for $\gamma\in\Gamma_0(N)$, and
$\Theta|_wg=\varepsilon_g\Theta$ with $\varepsilon_g\in\{\pm1\}$ for every Atkin--Lehner matrix
$g\in\Gamma$.
\item $t\in q+q^2\Z[[q]]$ is a Hauptmodul for $\Gamma$: it is $\Gamma$-invariant, holomorphic on
$\Hb$, meromorphic at the cusps, and induces an isomorphism $X(\Gamma)\to\mathbb P^1(\C)$.
\item ({\it cusp type}) $\Theta$ has no zeros in $\Hb$, and there is a cusp $c$ such that, for every
cusp $c'$, $\Theta$ vanishes at $c'$ if and only if $c'\in\Gamma c$, if and only if $t$ has a pole
at $c'$.
\item $r_\Theta:=\nu_A(\Theta)/(-\nu_A(t))\le1$ for $A\in\SL_2(\Z)$ with $A\infty=c$. (This is the
ratio of the orders of $\Theta$ and $1/t$ at $c$ in any local parameter; it does not depend on $A$.)
\end{enumerate}
For a prime $p\nmid N$ we consider the \emph{Hecke condition}
\begin{enumerate}[label=\textup{(C$_p$)}]
\item for every $\beta\in M(N,p)$ the cusp $\beta c$ lies in $\Gamma c$ and has the same width in
$\Gamma_0(N)$ as $c$.
\end{enumerate}
\end{definition}

\begin{lemma}\label{B:lem:Cp}
Let $c=a/s$ with $\gcd(a,s)=1$ and $e=\gcd(s,N)$. If $\gcd(e,N/e)\in\{1,2\}$, then $(\mathrm C_p)$
holds for every prime $p\nmid N$; in fact $\beta c\in\Gamma_0(N)c$. The same conclusion
$(\mathrm C_p)$ holds for the family $(\zeta)$ of Table~\ref{B:tab:fam} ($N=9$, $c=\frac13$).
\end{lemma}
\begin{proof}
Let $\beta=\bmat\alpha b{N\gamma}\delta\in M(N,p)$ and $(a_1,s_1)=(\alpha a+bs,N\gamma a+\delta s)$.
Since $\adj(\beta)(a_1,s_1)^{\mathsf T}=p(a,s)^{\mathsf T}$, the gcd $g$ of $a_1,s_1$ divides $p$,
and $\beta c=(a_1/g)/(s_1/g)$. As $p\nmid N$ and $\gcd(\delta,N)=1$ we get
$\gcd(s_1/g,N)=\gcd(\delta s,N)=e$. If $\gcd(e,N/e)\in\{1,2\}$, all cusps whose reduced denominator
has gcd $e$ with $N$ are $\Gamma_0(N)$-equivalent \cite[Prop.~3.8.3]{DS05},
\cite[Lemma~3.5]{Beu26}; equivalent cusps have the same width. For $N=9$ and $e=3$ there are two
such classes, represented by $\frac13$ and $\frac23$. They are interchanged by the Fricke matrix
$w_9=\bmat0{-1}90$ (as $w_9(\frac13)=-\frac13$), which lies in $\Gamma=\Gamma_0(9)+9$, and both have
width $1$.
\end{proof}

\subsection{\texorpdfstring{The power series $F$ and the excellent lift}{The power series F and the excellent lift}}
\begin{lemma}\label{B:lem:Ft}
Let $(N,\Gamma,w,\chi,\Theta,t)$ satisfy \textup{(F2)}, \textup{(F3)}. Let $q(t)\in t+t^2\Z[[t]]$ be
the inverse of $t(q)$, $F(t):=\Theta(q(t))\in1+t\Z[[t]]$, and $t^\sigma(t):=t(q(t)^p)$. Then
\begin{enumerate}[label=\textup{(\alph*)}]
\item $F(t(q))=\Theta(q)$ and $F(t^\sigma(t(q)))=\Theta(q^p)$;
\item $t^\sigma\in t^p+t^{p+1}\Z[[t]]$ and $t^\sigma\equiv t^p\pmod p$;
\item for $L\ge1$: $F_L(t^\sigma)/F(t^\sigma)\in1-a_Lt^{pL}+t^{pL+1}\Z[[t]]$.
\end{enumerate}
\end{lemma}
\begin{proof}
(a) is the definition. (b) $q(t)^p\in t^p+t^{p+1}\Z[[t]]$ and $t(q)\in q+q^2\Z[[q]]$; moreover
$t(q^p)\equiv t(q)^p\pmod p$ in $\Z[[q]]$, and substituting $q=q(t)$ gives $t^\sigma\equiv t^p$.
(c) $F_L(u)/F(u)=1-(F(u)-F_L(u))/F(u)\in1-a_Lu^L+u^{L+1}\Z[[u]]$; substitute $u=t^\sigma$.
\end{proof}

For all families of Table~\ref{B:tab:fam}, $t(q)$ is the mirror map of the Picard--Fuchs equation of
$F$ \cite{Za09,AZ06,Beu26}, so $t^\sigma$ is the excellent lift in the sense of \cite{BV3}. This
identification is used only in Corollary~\ref{B:cor:BV}; Theorem~\ref{B:thm:main} is a statement about
the modular $t(q)$.

\subsection{The examples}\label{B:ss:examples}
The coefficients $a_n$ of the families in Table~\ref{B:tab:fam} are
\begin{align*}
&\textstyle\binom{2n}n^2;\ \sum_k\binom nk^3;\ \sum_k(-3)^{n-3k}\frac{n!}{(n-3k)!k!^3};\
\sum_k\binom nk^2\binom{2k}k;\ \sum_k4^{n-2k}\binom n{2k}\binom{2k}k^2;\
\sum_k(-8)^{n-k}\binom nk\sum_j\binom kj^3;\\
&\textstyle\binom{2n}n^3;\ \sum_k\binom nk^2\binom{n+k}k^2;\
(-1)^n\sum_k(-1)^k\binom nk^3\bigl[\binom{4n-5k-1}{3n}+\binom{4n-5k}{3n}\bigr];\
(-1)^n\sum_k\binom nk^2\binom{2k}k\binom{2n-2k}{n-k};\\
&\textstyle\sum_k(-3)^{n-3k}\frac{(n+k)!}{k!^4(n-3k)!};\ \sum_k\binom nk^2\binom{2k}n^2;\
\sum_{k,l}\binom nk^2\binom nl\binom kl\binom{k+l}n,
\end{align*}
in the order of the table (binomial coefficients $\binom xy$ with $x<y$ are $0$); see
\cite[App.~B, C]{Beu26}, \cite{Za09,AZ06}. The hypotheses of Definition~\ref{B:def:cusp} are checked as follows.
Each $\Theta$ and $t$ is an $\eta$-quotient. By the theorem of Gordon, Hughes and Newman
\cite[Thm.~1.64]{Ono04} it is modular for $\Gamma_0(N)$ with character
$d\mapsto\bigl(\frac{(-1)^k\prod\delta^{r_\delta}}d\bigr)$, where $k$ is the weight. An
$\eta$-quotient has no zeros or poles in $\Hb$, and Ligozat's formula \cite[Thm.~1.65]{Ono04} gives
its order at every cusp. The Atkin--Lehner matrix $W_e$ maps $\eta_\delta$ to a constant times
$\eta_{e\delta/\gcd(e,\delta)^2}$, so invariance under $W_e$ can be read off the exponents. In every case
$t$ has a simple pole on $X(\Gamma)$ and no other poles, so it is a Hauptmodul of the genus-zero
group $\Gamma$. On $X_0(4)$ the family $\binom{2n}n^3$ has a double pole at $\frac12$, where
$X_0(4)\to X(\Gamma_0(4)+4)$ ramifies. All these data are verified by
checks~[1], [2] of the script in the appendix: the identity $F(t(q))=\Theta(q)$, the orders at all
cusps, $r_\Theta$, the characters and the Atkin--Lehner symmetries. Condition $(\mathrm C_p)$
holds for all $p\nmid N$ by Lemma~\ref{B:lem:Cp}, since $\gcd(e,N/e)\le2$ for each cusp in the
table except $\frac13$ in level $9$.

\section{The Hecke trace}\label{B:source-main-3}\label{B:sec:trace}
Throughout this section $(N,\Gamma,w,\chi,\Theta,t)$ is a modular family of cusp type, $p\nmid N$
satisfies $(\mathrm C_p)$, and $P\in\Z[T]$ is a nonzero polynomial. For $\beta\in M(N,p)$ put
\[
R_\beta(\tau):=p^{w/2}\chi(\beta)\,\frac{\Theta(\tau)}{(\Theta|_w\beta)(\tau)}\,P\bigl(t(\beta\tau)\bigr).
\]
Since $(\Theta|_w\beta)(\tau)=p^{w/2}j(\beta,\tau)^{-w}\Theta(\beta\tau)$ and $\Theta$ has no zeros in
$\Hb$, each $R_\beta$ is holomorphic on $\Hb$.

\begin{lemma}\label{B:lem:perm}
\begin{enumerate}[label=\textup{(\alph*)}]
\item $R_{\gamma\beta}=R_\beta$ for $\gamma\in\Gamma_0(N)$.
\item $R_\beta(\gamma\tau)=R_{\beta\gamma}(\tau)$ for $\gamma\in\Gamma_0(N)$.
\item $R_\beta(g\tau)=R_{g^{-1}\beta g}(\tau)$ for every Atkin--Lehner matrix $g\in\Gamma$.
\end{enumerate}
Consequently $\Tr_p(P):=\sum_{[\beta]\in\Gamma_0(N)\backslash M(N,p)}R_\beta$ is well defined and
$\Gamma$-invariant.
\end{lemma}
\begin{proof}
(a) $\Theta|_w\gamma\beta=\chi(\gamma)\Theta|_w\beta$, $\chi(\gamma\beta)=\chi(\gamma)\chi(\beta)$,
$\chi(\gamma)^2=1$ and $t\circ\gamma=t$. (b) We have
$\Theta(\gamma\tau)=\chi(\gamma)j(\gamma,\tau)^w\Theta(\tau)$ and
$(\Theta|_w\beta)(\gamma\tau)=j(\gamma,\tau)^w(\Theta|_w\beta\gamma)(\tau)$, so
$\chi(\beta)\Theta(\gamma\tau)/(\Theta|_w\beta)(\gamma\tau)=\chi(\beta\gamma)\Theta(\tau)/(\Theta|_w\beta\gamma)(\tau)$,
and $t(\beta\gamma\tau)$ is the argument of $P$ in $R_{\beta\gamma}$. (c) Here $\chi=1$. Put
$\beta'=g^{-1}\beta g\in M(N,p)$ (Lemma~\ref{B:lem:hecke}(H4)). Then
$\Theta|_w\beta g=\Theta|_wg\beta'=\varepsilon_g\Theta|_w\beta'$, and as in (b)
$\Theta(g\tau)/(\Theta|_w\beta)(g\tau)=\varepsilon_g\Theta(\tau)/(\Theta|_w\beta g)(\tau)
=\Theta(\tau)/(\Theta|_w\beta')(\tau)$, while $t(\beta g\tau)=t(g\beta'\tau)=t(\beta'\tau)$. The last
assertion follows from (H3) and (H4).
\end{proof}

\begin{lemma}\label{B:lem:qexp}
$R_{\beta_\infty}=Q_P$ with $Q_P(\tau)=\Theta(\tau)P(t(p\tau))/\Theta(p\tau)$, and
\[
\sum_{k=0}^{p-1}R_{\beta_k}=p^{w+1}\chi(p)\,(U_pG_P)(q),\qquad
G_P(z)=\frac{\Theta(z^p)\,P(t(z))}{\Theta(z)}\in\Z[[z]].
\]
In particular the $q$-expansion of $\Tr_p(P)$ is $Q_P+p^{w+1}\chi(p)U_pG_P\in\Z[[q]]$.
\end{lemma}
\begin{proof}
$\Theta|_w\beta_\infty=p^{w/2}\Theta(p\tau)$ and $\chi(\beta_\infty)=1$. Next,
$(\Theta|_w\beta_k)(\tau)=p^{-w/2}\Theta(\frac{\tau+k}p)$ and $\chi(\beta_k)=\chi(p)$, so
$R_{\beta_k}=p^w\chi(p)\Theta(\tau)P(t(\frac{\tau+k}p))/\Theta(\frac{\tau+k}p)=p^w\chi(p)G_P(\zeta_p^kq^{1/p})$,
because $\Theta(\tau)=\Theta\bigl((\zeta_p^kq^{1/p})^p\bigr)$. Now apply \eqref{B:eq:Up}. Integrality
holds because $\Theta\in1+q\Z[[q]]$ is a unit in $\Z[[q]]$.
\end{proof}

\begin{lemma}[Branches at the cusp]\label{B:lem:branch}
Let $A\in\SL_2(\Z)$ with $A\infty=c$. For $\beta\in M(N,p)$ write $\beta A=A_\beta\beta'$ with
$A_\beta\in\SL_2(\Z)$ and $\beta'=\bmat{a'}{b'}0{d'}$, $a'd'=p$ (Hermite normal form).
\begin{enumerate}[label=\textup{(\roman*)}]
\item Exactly one class $[\beta]$ has $(a',d')=(p,1)$; the other $p$ classes have $(a',d')=(1,p)$.
\item For every $\Gamma$-invariant $\varphi$ meromorphic at $c$:
$\nu_A(\varphi\circ\beta)=\frac{a'}{d'}\,\nu_A(\varphi)$. Moreover
$\nu_A(\Theta|_w\beta)=\frac{a'}{d'}\,\nu_A(\Theta)$.
\end{enumerate}
\end{lemma}
\begin{proof}
(i) By (H2) and since right multiplication by $A$ permutes $\SL_2(\Z)\backslash M_p$, the map
$[\beta]\mapsto\SL_2(\Z)\beta A$ is a bijection onto $\SL_2(\Z)\backslash M_p$. The diagonal
$(a',d')$ of the Hermite normal form is an invariant of the class, and $M_p$ has one class with
$(p,1)$ and $p$ classes with $(1,p)$.

(ii) $A_\beta\infty=\beta A\infty=\beta c$. By $(\mathrm C_p)$ there is $g\in\Gamma$ with $gc=\beta c$,
and the widths of $c$ and $\beta c$ in $\Gamma_0(N)$ are equal, say to $h$. The matrix
$\delta:=A_\beta^{-1}gA\in\GL_2^+(\Q)$ fixes $\infty$, so $\delta=\bmat ux0v$ with $uv>0$. Since $g$
normalises $\Gamma_0(N)$ and $\operatorname{Stab}_{\Gamma_0(N)}(c)=\pm A\langle T^h\rangle A^{-1}$,
$\operatorname{Stab}_{\Gamma_0(N)}(\beta c)=\pm A_\beta\langle T^h\rangle A_\beta^{-1}$, conjugation
by $\delta$ maps $\pm\langle T^h\rangle$ onto itself. As $\delta T^h\delta^{-1}=T^{hu/v}$, this
forces $u/v=\pm1$, hence $u=v$. Thus $A_\beta=u^{-1}gAT^{-x/u}$. For $\Gamma$-invariant $\varphi$ we get
$\varphi(\beta A\tau)=\varphi(A_\beta\beta'\tau)=\varphi\bigl(A(\beta'\tau-x/u)\bigr)$ with
$\beta'\tau-x/u=\frac{a'}{d'}\tau+\mathrm{const}$, which proves the first claim. Similarly
$(\Theta|_w\beta)|_wA=\Theta|_wA_\beta\beta'=\pm\varepsilon_g\,(\Theta|_wA)|_w(T^{-x/u}\beta')$, and
$(f|_w T^{y}\beta')(\tau)=p^{w/2}d'^{-w}f\bigl(\frac{a'\tau+b'}{d'}+y\bigr)$; here
$\varepsilon_g:=\chi(g)$ if $g\in\Gamma_0(N)$.
\end{proof}

\begin{proposition}\label{B:prop:poly}
There is a polynomial $\Pi_P\in\Z[T]$ with $\Tr_p(P)=\Pi_P(t)$ and
\[
\deg\Pi_P\le p\deg P+\lfloor r_\Theta(p-1)\rfloor .
\]
\end{proposition}
\begin{proof}
\emph{Cusps outside $\Gamma c$.} Let $c'\notin\Gamma c$ and $A'\in\SL_2(\Z)$ with $A'\infty=c'$. For
$\beta\in M(N,p)$ we have $\beta c'\notin\Gamma c$. Indeed, if $\beta c'=gc$ with $g\in\Gamma$, then
$c'=\adj(\beta)gc=g\,(g^{-1}\adj(\beta)g)\,c\in gM(N,p)c\subseteq\Gamma c$ by (H4), (H5) and
$(\mathrm C_p)$, a contradiction. Write $\beta A'=A'_\beta\beta'$ as in Lemma~\ref{B:lem:branch}. Then
$A'_\beta\infty=\beta c'$. Since $\Theta$ does not vanish at $c'$ and $\beta c'$, and $t$ has no pole there,
\[
R_\beta(A'\tau)=p^{w/2}\chi(\beta)\frac{(\Theta|_wA')(\tau)}{\bigl((\Theta|_wA'_\beta)|_w\beta'\bigr)(\tau)}
\,P\bigl(t(A'_\beta\beta'\tau)\bigr)
\]
is holomorphic at $\tau=i\infty$; the factors $j(A',\tau)^w$ cancel. So each $R_\beta$, and hence $\Tr_p(P)$, is holomorphic at $c'$.

\emph{The cusp $c$.} Put $\rho=\nu_A(\Theta)>0$ and $\kappa=-\nu_A(t)>0$, so $r_\Theta=\rho/\kappa$.
As above, $R_\beta(A\tau)=p^{w/2}\chi(\beta)\,(\Theta|_wA)(\tau)\,
\bigl((\Theta|_w\beta)|_wA\bigr)(\tau)^{-1}P(t(\beta A\tau))$. Since $\nu_A(t)<0$ we have
$\nu_A(P\circ t)=-\kappa\deg P$. Lemma~\ref{B:lem:branch} then gives
\[
\nu_A(R_\beta)=\rho-\tfrac{a'}{d'}\rho-\tfrac{a'}{d'}\kappa\deg P=
\begin{cases}
-(p-1)\rho-p\kappa\deg P, & (a',d')=(p,1),\\
(1-\frac1p)\rho-\frac1p\kappa\deg P, & (a',d')=(1,p).
\end{cases}
\]
Hence $\nu_A(\Tr_p(P))\ge-(p-1)\rho-p\kappa\deg P$.

\emph{Polynomiality.} By Lemma~\ref{B:lem:perm} and the two previous steps, $\Tr_p(P)$ is a
$\Gamma$-invariant function, holomorphic on $\Hb$ and at all cusps outside $\Gamma c$, and
meromorphic at $c$. It therefore defines a regular function on $X(\Gamma)\smallsetminus\{c\}$.
Since $t$ maps $X(\Gamma)$ isomorphically onto $\mathbb P^1$ with $t(c)=\infty$, we get
$\Tr_p(P)=\Pi_P(t)$ for some $\Pi_P\in\C[T]$. Comparing $\nu_A$ gives
$\kappa\deg\Pi_P=-\nu_A(\Tr_p(P))\le(p-1)\rho+p\kappa\deg P$, which is the degree bound. Finally, the
$q$-expansion of $\Pi_P(t)$ lies in $\Z[[q]]$ by Lemma~\ref{B:lem:qexp}, and $t\in q+q^2\Z[[q]]$.
Comparing coefficients of $q^0,q^1,\dots$ inductively shows $\Pi_P\in\Z[T]$.
\end{proof}

\begin{remark}\label{B:rem:hecke}
Let $T_p=U_p+\chi(p)p^{-w-1}V_p$ be the Hecke operator on forms of weight $-w$ and character
$\chi$, where $V_ph(q)=h(q^p)$. By \eqref{B:eq:Up}, $U_pG_P=\Theta\cdot U_p(P(t)/\Theta)$, so
Lemma~\ref{B:lem:qexp} says $\Tr_p(P)/\Theta=p^{w+1}\chi(p)\,T_p\bigl(P(t)/\Theta\bigr)$.
Proposition~\ref{B:prop:poly} is therefore the statement that $T_p$ maps
$\{P(t)/\Theta:\deg P\le d\}$ into $p^{-w-1}\{P(t)/\Theta:\deg P\le pd+\lfloor r_\Theta(p-1)\rfloor\}$.
For a family of cusp type these spaces of weakly holomorphic forms of weight $-w$ consist of
forms with poles only at $\Gamma c$.
\end{remark}

\section{Proof of the Main Theorem}\label{B:source-main-4}\label{B:sec:proof}
Fix $m\ge1$. Put $K_j=mp^j$ for $j\ge0$ and $F_{K_{-1}}:=0$. For $s\ge1$ consider
\[
(\mathrm H_s)\qquad \Theta(q)\,F_{K_{s-1}}\bigl(t(q^p)\bigr)\equiv\Theta(q^p)\,F_{K_s}\bigl(t(q)\bigr)\pmod{p^{(w+1)s}\Z[[q]]}.
\]
Dividing by the unit $\Theta(q^p)F_{K_{s-1}}(t(q^p))$ of $\Z[[q]]$ and substituting $q=q(t)$ shows that
$(\mathrm H_s)$ is equivalent to the congruence of Theorem~\ref{B:thm:main} (Lemma~\ref{B:lem:Ft}(a)).
We prove $(\mathrm H_s)$ by induction on $s$.

Let $s\ge1$, and assume $(\mathrm H_{s-1})$ if $s\ge2$. Apply Lemma~\ref{B:lem:qexp} and
Proposition~\ref{B:prop:poly} to $P=F_{K_{s-1}}$, a polynomial of degree $<K_{s-1}$:
\begin{equation}\label{B:eq:main}
Q:=\frac{\Theta(q)\,F_{K_{s-1}}(t(q^p))}{\Theta(q^p)}=\Pi(t)-p^{w+1}\chi(p)\,U_pG,\qquad
G(z)=\frac{\Theta(z^p)F_{K_{s-1}}(t(z))}{\Theta(z)},
\end{equation}
with $\Pi\in\Z[T]$ and $\deg\Pi\le p(K_{s-1}-1)+\lfloor r_\Theta(p-1)\rfloor\le K_s-p+(p-1)<K_s$, by (F5).

\emph{Claim:} $p^{w+1}U_pG\equiv p^{w+1}F_{K_{s-2}}(t(q))\pmod{p^{(w+1)s}}$. For $s=1$ both sides
are $\equiv0\pmod{p^{w+1}}$. For $s\ge2$, $(\mathrm H_{s-1})$ in the variable $z$ reads
$\Theta(z)F_{K_{s-2}}(t(z^p))\equiv\Theta(z^p)F_{K_{s-1}}(t(z))\pmod{p^{(w+1)(s-1)}}$. Dividing by the unit
$\Theta(z)$ gives $G(z)\equiv F_{K_{s-2}}(t(z^p))\pmod{p^{(w+1)(s-1)}}$. Applying $U_p$, which
preserves congruences, and \eqref{B:eq:Up} gives $U_pG\equiv F_{K_{s-2}}(t(q))$ modulo $p^{(w+1)(s-1)}$, and the claim follows.

By \eqref{B:eq:main} and the claim,
\[
Q\equiv\widetilde\Pi(t):=\Pi(t)-p^{w+1}\chi(p)F_{K_{s-2}}(t)\pmod{p^{(w+1)s}},\qquad
\widetilde\Pi\in\Z[T],\ \deg\widetilde\Pi<K_s .
\]
On the other hand, Lemma~\ref{B:lem:Ft}(c) gives
$Q=F(t)\cdot F_{K_{s-1}}(t^\sigma)/F(t^\sigma)\in F(t)+t^{K_s}\Z[[t]]$. Compare coefficients of
$t^n$ in $Q\equiv\widetilde\Pi(t)$. For $n<K_s$ we get $[t^n]\widetilde\Pi\equiv a_n$. For
$n\ge K_s$ we get $[t^n]Q\equiv0$. Hence $Q\equiv F_{K_s}(t)\pmod{p^{(w+1)s}}$, which is
$(\mathrm H_s)$. This completes the induction and the proof of Theorem~\ref{B:thm:main}.\qed

\begin{remark}
The proof gives the explicit polynomial congruence
\[
F_{mp^s}(t)\equiv\Pi_{F_{mp^{s-1}}}(t)-p^{w+1}\chi(p)F_{mp^{s-2}}(t)\pmod{p^{(w+1)s}},
\]
with $\Pi_{F_{mp^{s-1}}}$ the trace polynomial of Proposition~\ref{B:prop:poly}. In the examples the
exponent $(w+1)s$ is sharp: for most $(p,s,m)$ the difference in $(\mathrm H_s)$ has valuation exactly
$(w+1)s$ (check~[4]).
\end{remark}

\section{Proofs of the corollaries}\label{B:source-main-5}\label{B:sec:cor}

\subsection{Coefficients}
\begin{proof}[Proof of Corollary~\ref{B:cor:coef}]
By Lemma~\ref{B:lem:Ft}(b),(c) and $a_0=1$, the coefficient of $t^{K_s}$ in
\[
Q=F(t)\,\frac{F_{K_{s-1}}(t^\sigma)}{F(t^\sigma)}
\]
is $a_{K_s}-a_{K_{s-1}}$. Since
$Q\equiv F_{K_s}(t)\pmod{p^{(w+1)s}}$ and $\deg F_{K_s}<K_s$, it is divisible by $p^{(w+1)s}$.
\end{proof}

\subsection{Conjecture BV}
\begin{lemma}\label{B:lem:subst}
Let $F(u)\in1+u\Z[[u]]$ be annihilated by a differential operator with maximal unipotent
monodromy at $0$, with canonical coordinate $q_F(u)\in u+u^2\Z_p[[u]]$ and mirror map $t_F$, and
let $u^{\sigma}=t_F(q_F(u)^p)$. Let $k\ge1$ with $p\nmid k$, and $F_0(t):=F(t^k)$. Then:
\begin{enumerate}[label=\textup{(\alph*)}]
\item the canonical coordinate of $F_0$ is $q_0(t)=t\bigl(q_F(t^k)/t^k\bigr)^{1/k}\in t+t^2\Z_p[[t]]$;
\item the excellent lift $t^{\sigma_0}=t_0(q_0(t)^p)$ of $F_0$ satisfies $(t^{\sigma_0})^k=(t^k)^\sigma$;
\item $F_{0,kM}(t)=F_M(t^k)$ for all $M\ge1$.
\end{enumerate}
Consequently, for every $m',s$ and $e$: \eqref{B:eq:dwork} for $F_0$ with $m=km'$ holds modulo
$p^e$ if and only if \eqref{B:eq:dwork} for $F$ with $m'$ holds modulo $p^e$.
\end{lemma}
\begin{proof}
Let $L(u,\theta_u)$ be an operator annihilating $F$. Then $L(t^k,\theta_t/k)$ annihilates $F_0$.
Under $u=t^k$ the Frobenius solution $\log(u)F(u)+G(u)$ becomes
\[
\log(t^k)F(t^k)+G(t^k)=k\Bigl(\log(t)F_0(t)+\frac{G(t^k)}k\Bigr).
\]
So the normalised logarithmic solution of $F_0$ is $\log(t)F_0+G(t^k)/k$, and
\[
q_0(t)=t\exp\Bigl(\frac{G(t^k)}{kF(t^k)}\Bigr)=t\,\Bigl(\frac{q_F(u)}{u}\Bigr)^{1/k}\Big|_{u=t^k}.
\]
This proves (a); integrality holds because $p\nmid k$. Thus $q_0(t)^k=q_F(t^k)$, and inverting gives $t_0(Q)^k=t_F(Q^k)$. Hence
$(t^{\sigma_0})^k=t_F(q_0(t)^{pk})=t_F(q_F(t^k)^p)=(t^k)^\sigma$, which is (b). (c) holds because $F_0$
only contains powers $t^{kj}$. For the consequence, substitute $u=t^k$:
\[
\frac{F_0(t)}{F_0(t^{\sigma_0})}=\frac{F(u)}{F(u^\sigma)},\qquad F_{0,km'p^s}(t)=F_{m'p^s}(u),\qquad
F_{0,km'p^{s-1}}(t^{\sigma_0})=F_{m'p^{s-1}}(u^\sigma).
\] A power series in $u$ is divisible by
$p^e$ if and only if it is so after the substitution $u=t^k$.
\end{proof}

\begin{proof}[Proof of Corollary~\ref{B:cor:BV}]
For $n=2$ the hypercubic and hyperoctahedral Laurent polynomials have the same constant terms,
$\CT\bigl((x_1+x_1^{-1})^{j}(x_2+x_2^{-1})^{j}\bigr)=\CT\bigl((x_1+x_1^{-1}+x_2+x_2^{-1})^{j}\bigr)$,
equal to $\binom{2k}k^2$ for $j=2k$ and $0$ for $j$ odd. So
$F_0(t)=\sum_k\binom{2k}k^2t^{2k}=F(t^2)$ with $F$ the Legendre family of Table~\ref{B:tab:fam}. For
$n=3$, hypercubic, $F_0(t)=\sum_k\binom{2k}k^3t^{2k}=F(t^2)$ with $F$ the family $\binom{2n}n^3$. In all
three cases the vertex coefficient is $\gamma=1$, so the excellent lift of \cite[Def.~7.14]{BV3} is
$q_0\mapsto q_0^p$ in the canonical coordinate of $F_0$. Since the modular $t(q)$ of
Table~\ref{B:tab:fam} is the mirror map of $F$ (for Legendre see \cite[\S1]{Beu26} and the
discussion in \cite[\S8]{BV3}, where $t_0(q)^2=\lambda(q^2)/16$; for $\binom{2n}n^3$ see
\cite[App.~C, 9.1]{Beu26}), Lemma~\ref{B:lem:subst} with $k=2$ reduces \eqref{B:eq:dwork} for $F_0$ and
even $m=2m'$ to Theorem~\ref{B:thm:main} for $F$ and $m'$, with $w=1$, respectively $w=2$. The primes
allowed in \cite[Thm.~7.3]{BV3} are odd for $n=2$ and $\ge5$ for $n=3$ ($\#G=48$), and they do not
divide $N=4$. Check~[6] of the script recomputes the canonical coordinate of $F_0$ directly from
its Picard--Fuchs equation and verifies \eqref{B:eq:dwork} in the normalisation of \cite{BV3}.

For the simplicial family with $n$ variables, $\CT(g^k)=0$ unless $(n+1)\mid k$, and
$\CT(g^{(n+1)j})=((n+1)j)!/j!^{n+1}$. So $F_0(t)=F(t^{n+1})$ with $F=F_{27}$ for $n=2$ and $F=F_{256}$ for $n=3$. The vertex
coefficients of $g$ are $1$, so the excellent lift is again $q_0\mapsto q_0^p$. The modular $t(q)$ of Table~\ref{B:tab:ell} is the
mirror map of $F$; equivalently, by Lemma~\ref{B:lem:subst}(a), the mirror map $t_0$ of $F_0$ satisfies
$t_0(Q)^{n+1}=t(Q^{n+1})$, which check~[8g] verifies with $t_0$ computed from the Picard--Fuchs equation of $F_0$. Now
Lemma~\ref{B:lem:subst} with $k=n+1$ (note $p\nmid n+1$ for $p\ge5$) reduces \eqref{B:eq:dwork} for $F_0$ and $m=(n+1)m'$ to
Theorem~\ref{B:thm:ell} for $F$ and $m'$. The symmetry group $G$ of $g$ contains the $(n+1)!$ permutations of its monomials,
so the primes allowed in \cite[Thm.~7.3]{BV3} are $\ge5$. Check~[8g] also verifies \eqref{B:eq:dwork} for $F_0$ directly in
the normalisation of \cite{BV3}.
\end{proof}

\begin{proof}[Proof of Corollary~\ref{B:cor:Beu}]
Immediate from Theorem~\ref{B:thm:main}, Lemma~\ref{B:lem:Cp} and \S\ref{B:ss:examples}. Among the families
of \cite[App.~B, C]{Beu26} satisfying Assumption~(I) or (II), those with $\Theta$ an $\eta$-product
are exactly the families of Table~\ref{B:tab:fam}. For an $\eta$-product, (F4) holds automatically
because an $\eta$-quotient has no zeros in $\Hb$.
\end{proof}

\section{Families with an elliptic zero}\label{B:source-main-6}\label{B:sec:ell}

In this section we prove Theorem~\ref{B:thm:ell}. The trace of $Q_P$ is no longer a polynomial, so
the argument of \S\ref{B:sec:proof} has to be replaced by a $p$-adic estimate.

\subsection{The three families}
Let $a(q)=\sum_{x,y\in\Z}q^{x^2+xy+y^2}$ be the theta series of the hexagonal lattice, $E_2=1-24\sum_{n\ge1}\sigma(n)q^n$,
$h_2=\eta_2^{24}/\eta_1^{24}$ and $h_3=\eta_3^{12}/\eta_1^{12}$. The data of the three families are in
Table~\ref{B:tab:ell}; $W_N=\bmat0{-1}N0$ denotes the Fricke matrix.

\begin{table}[ht]
\centering\small
\renewcommand{\arraystretch}{1.3}
\begin{tabular}{@{}lllllllll@{}}
\toprule
$F$ & $w$ & $\Gamma$ & $\chi$ & $\Theta$ & $t$ & $e$ & $k_0$ & $o$\\\midrule
$F_{27}=\sum\frac{(3n)!}{n!^3}t^n$ & 1 & $\Gamma_0(3)$ & $\chi_{-3}$ & $a$ & $\eta_3^9/(\eta_1^3a^3)=\frac{h_3}{1+27h_3}$ &
$\frac{3+\sqrt{-3}}6$ & 3 & 1\\
$F_{256}=\sum\frac{(4n)!}{n!^4}t^n$ & 2 & $\Gamma_0(2)+2$ & 1 & $2E_2(2\tau)-E_2(\tau)$ & $h_2/(1+64h_2)^2$ & $\frac{1+i}2$ &
4 & 1\\
$F_{108}=\sum\frac{(2n)!(3n)!}{n!^5}t^n$ & 2 & $\Gamma_0(3)+3$ & 1 & $a^2$ & $h_3/(1+27h_3)^2$ & $\frac{3+\sqrt{-3}}6$ & 6
& 2\\\bottomrule
\end{tabular}
\caption{Families with an elliptic zero. $e$ is the zero of $\Theta$, $k_0$ the order of the elliptic point $e$ of $\Gamma$,
$o$ the order of vanishing of $\Theta$ at $e$ in the variable $\tau-e$, and $r=o/k_0\in\{\frac13,\frac14,\frac13\}$.}
\label{B:tab:ell}
\end{table}

We also use a root $\vartheta$ of $\Theta$: $\vartheta=\Theta$, $\chi_\vartheta=\chi$ in the first two cases and
$\vartheta=a$, $\chi_\vartheta=\chi_{-3}$ in the third case, so that $\Theta=\vartheta^{o}$ in all cases.

\begin{lemma}\label{B:lem:elldata}
For each family of Table~\ref{B:tab:ell}:
\begin{enumerate}[label=\textup{(\alph*)}]
\item $F(t(q))=\Theta(q)$; $\Theta$ is a modular form of weight $w$ for $\Gamma_0(N)$ with character $\chi$, and
$\Theta|_wW_N=-\Theta$ in the second and third case; $t\in q+q^2\Z[[q]]$ is a Hauptmodul for $\Gamma$. The form
$\vartheta$ is modular for $\Gamma_0(N)$ with character $\chi_\vartheta$, and in the third case $a|_1W_3=-ia$.
\item $e$ is an elliptic point of $\Gamma$ whose stabiliser in $\Gamma$, modulo scalar matrices, has order $k_0$.
\item The zeros of $\Theta$ in $\Hb\cup\mathbb P^1(\Q)$ form the orbit $\Gamma e$; $\Theta$ vanishes at $e$ to order $o$
and $\vartheta$ to order $1$, in the variable $\tau-e$. The function $t$ is holomorphic on $\Hb\smallsetminus\Gamma e$
and at all cusps, and has a pole of order $k_0$ at $e$ in the variable $\tau-e$.
\end{enumerate}
\end{lemma}

\begin{proof}
(a) The identities $F(t(q))=\Theta(q)$ are classical (Ramanujan's theories of elliptic functions to alternative bases,
see \cite{BBG94,Co17} and \cite[App.~B, C]{Beu26}); check~[8a] verifies them to $O(q^{60})$. The series $a(q)$ is a modular form of
weight one for $\Gamma_0(3)$ with character $\chi_{-3}$, and Poisson summation gives $a(-1/(3\tau))=-i\sqrt3\,\tau a(\tau)$,
that is $a|_1W_3=-ia$; hence $a^2|_2W_3=-a^2$. Further $2E_2(2\tau)-E_2(\tau)\in M_2(\Gamma_0(2))$ is anti-invariant under
$W_2$, and $(3E_2(3\tau)-E_2(\tau))/2=a^2$ (check~[8a]). The function $h_N$ is a Hauptmodul for $\Gamma_0(N)$ with
$h_N\circ W_N=N^{-12/(N-1)}/h_N$, so $h_N/(1+N^{6/(N-1)}h_N)^2$ (with $N^{6/(N-1)}=64,27$) is $W_N$-invariant of degree two in $h_N$, i.e.\ a
Hauptmodul for $\Gamma_0(N)+N$. In the first case $\eta_3^9/(\eta_1^3a^3)=h_3/(1+27h_3)$ is the identity
$a^3=\eta_1^9/\eta_3^3+27\eta_3^9/\eta_1^3$ of \cite{BBG94} (check~[8a]).

(b) Put $\gamma_e=\bmat1{-1}3{-2}$ in the first and third case and $\gamma_e=\bmat1{-1}2{-1}$ in the second. Then
$\gamma_e\in\Gamma_0(N)$ fixes $e$, and $\gamma_e$ has order $3$, respectively $2$, in $\SL_2(\Z)/\{\pm1\}$.
Moreover $W_3e=e-1$ and $W_2e=e-1$, so $T W_N$ also fixes $e$. Modulo scalars, the stabiliser of $e$ in $\Gamma$ is
therefore generated by $\gamma_e$ (first case), respectively by $\gamma_e$ and $TW_N$, and has order $3$, $4$, $6$.

(c) We have $j(\gamma_e,e)=3e-2=e^{2\pi i/3}$, respectively $2e-1=i$, and $\chi(\gamma_e)=\chi_{-3}(-2)=1$. Evaluating
$\Theta(\gamma_e\tau)=\chi(\gamma_e)j(\gamma_e,\tau)^w\Theta(\tau)$ at $\tau=e$ gives $\Theta(e)=e^{2\pi i/3}\Theta(e)$,
respectively $\Theta(e)=i^2\Theta(e)$, so $\Theta(e)=0$; likewise $a(e)=0$. Check~[8a] verifies the identities
\[
a^3t=\frac{\eta_3^9}{\eta_1^3},\qquad \Theta^4t=\eta_1^8\eta_2^8,\qquad \Theta^3t=\eta_1^6\eta_3^6,
\]
that is $\Theta^{1/r}t=$ an $\eta$-product, which has neither zeros nor poles in $\Hb$. As $\Theta(e)=0$, $t$ has a pole
at $e$. Since $t$ is a Hauptmodul, this is its only pole on $X(\Gamma)$, and it is simple on $X(\Gamma)$; as the local
parameter of $X(\Gamma)$ at $e$ is $((\tau-e)/(\tau-\bar e))^{k_0}$, the pole has order $k_0$ in $\tau-e$. The
$\eta$-identity now shows that $\Theta^{1/r}$ vanishes to order exactly $k_0$ at $e$, i.e.\ $\Theta$ to order $o=rk_0$,
and that $\Theta$ has no zeros in $\Hb\smallsetminus\Gamma e$. Finally, by the valence formula the zeros of a modular form
of weight $w$ on $X_0(N)$, each weighted by the inverse of the order of its stabiliser in
$\Gamma_0(N)/\{\pm1\}$, have total mass $w[\SL_2(\Z):\Gamma_0(N)]/12$, which is $\frac13,\frac12,\frac23$. In
$\Gamma_0(N)$ the point $e$ has order $3,2,3$, and $o/3,o/2,o/3$ are the same numbers. So $\Theta$ has no zeros at the
cusps. The statements for $\vartheta=a$ follow from $\Theta=a^2$.
\end{proof}

\paragraph{Symbolic verification of the modular identities and coordinate.}
The identity
\[
 a^2=\frac{3E_2(3\tau)-E_2(\tau)}2
\]
also follows directly from the valence formula: both sides are
holomorphic modular forms of weight two on $\Gamma_0(3)$ with constant
term $1$, whereas a nonzero difference cannot have order at least one
at infinity because the total valence is $2\cdot4/12=2/3$.
For the level-two form $\Theta=2E_2(2\tau)-E_2(\tau)$, the identity
\[
 \Theta^2=\frac{\eta_1^{16}}{\eta_2^8}
            +64\frac{\eta_2^{16}}{\eta_1^8}
\]
is proved in the same way. Both sides are holomorphic forms of weight
four on $\Gamma_0(2)$, and their constant and $q$ coefficients are
$1$ and $48$. A nonzero difference of order at least two would contradict
the total valence $4\cdot3/12=1$. For the two eta quotients, the exponent pairs at levels $1,2$ are
$(16,-8)$ and $(-8,16)$. Their characters are trivial, and their orders
at $(\infty,0)$ on $X_0(2)$ are respectively $(0,1)$ and $(1,0)$,
as follows from the eta-quotient criterion and cusp-order formula.
Thus both are holomorphic, as required for the valence argument. Substitution of the formulas for $t$ now gives
\[
 a^3t=\frac{\eta_3^9}{\eta_1^3},\qquad
 \Theta^4t=\eta_1^8\eta_2^8\quad(F_{256}),\qquad
 \Theta^3t=\eta_1^6\eta_3^6\quad(F_{108}).
\]
The first identity is the definition of $t$ in the $27t$ case; the last
follows algebraically from
$a^3=(\eta_1^9/\eta_3^3)(1+27h_3)$ and $\Theta=a^2$.

\smallskip
Write $D=q\,d/dq$. Logarithmic differentiation of the eta products gives
$Dh_3/h_3=a^2$ and $Dh_2/h_2=2E_2(2\tau)-E_2(\tau)$. Hence the three
modular parametrisations satisfy, respectively,
\[
 \frac{Dt}{t}=(1-27t)F_{27}(t)^2,\qquad
 \frac{Dt}{t}=\sqrt{1-256t}\,F_{256}(t),\qquad
 \frac{Dt}{t}=\sqrt{1-108t}\,F_{108}(t),
\]
where each square root has constant term $1$. For the second and third
identities use
\[
 \sqrt{1-256t}=\frac{1-64h_2}{1+64h_2},\qquad
 \sqrt{1-108t}=\frac{1-27h_3}{1+27h_3}.
\]
The Wronskian equation for the Gauss hypergeometric equation gives
the first differential identity for its normalized logarithmic
coordinate. Clausen's identity expresses the two ${}_3F_2$ periods as
the squares of
${}_2F_1(1/8,3/8;1;256t)$ and
${}_2F_1(1/6,1/3;1;108t)$, respectively, and the same Wronskian
calculation gives the other two identities. The normalized condition
$q=t+O(t^2)$ determines the formal coordinate uniquely. Thus these
identities identify the modular coordinate with the Frobenius
logarithmic coordinate coefficientwise; finite check [8g] is
supplementary.

\subsection{Hecke translates of the elliptic point}
Let $p\ge5$ be a prime. For $\beta\in M(N,p)$ put
\[
W_\beta=t\circ\beta,\qquad x_\beta=p^{w/2}\chi(\beta)\frac{\Theta}{\Theta|_w\beta},
\]
so that $R_\beta=x_\beta\,P(W_\beta)$ in the notation of \S\ref{B:sec:trace}. The proof of Lemma~\ref{B:lem:perm} applies
verbatim, because $\Theta|_wW_N=\pm\Theta$ by Lemma~\ref{B:lem:elldata}(a). So for $\gamma\in\Gamma$ the families
$(x_\beta\circ\gamma,W_\beta\circ\gamma)$ and $(x_\beta,W_\beta)$, indexed by $\Gamma_0(N)\backslash M(N,p)$, agree up
to a permutation of the index set. Hence every polynomial expression in the $p+1$ pairs $(x_\beta,W_\beta)$ that is symmetric
under permutations of the pairs is a $\Gamma$-invariant function. For $\tau_0\in\Hb$ let
\[
m(\tau_0)=\#\{[\beta]\in\Gamma_0(N)\backslash M(N,p):\ \beta\tau_0\in\Gamma e\}.
\]

\begin{lemma}\label{B:lem:loc}
Let $\tau_0\in\Hb$ and $\beta\in M(N,p)$.
\begin{enumerate}[label=\textup{(\alph*)}]
\item If $\beta\tau_0\in\Gamma e$, then in the variable $\tau-\tau_0$ the function $W_\beta$ has a pole of order $k_0$, and
$\Theta|_w\beta$, $\vartheta|\beta$ vanish to order $o$ and $1$. Consequently $x_\beta$ has a pole of order $o$ at $\tau_0$
if $\tau_0\notin\Gamma e$, and $x_\beta$ is holomorphic and nonzero at $\tau_0$ if $\tau_0\in\Gamma e$.
\item If $\beta\tau_0\notin\Gamma e$, then $W_\beta$ and $x_\beta$ are holomorphic at $\tau_0$, and if moreover
$\tau_0\in\Gamma e$, then $x_\beta$ vanishes to order $o$ at $\tau_0$.
\item $W_\beta$ and $x_\beta$ are holomorphic at every cusp.
\end{enumerate}
\end{lemma}

\begin{proof}
$\beta$ is a M\"obius transformation which is biholomorphic near $\tau_0$, and $(\Theta|_w\beta)(\tau)$ is
$\Theta(\beta\tau)$ times a nowhere vanishing factor. By modularity the orders of $t$, $\Theta$ and $\vartheta$ at a
point of $\Gamma e$ are those at $e$. Now (a) and (b) follow from Lemma~\ref{B:lem:elldata}(c). For (c), $\beta$ maps cusps
to cusps, where $t$ is holomorphic and $\Theta$ does not vanish.
\end{proof}

\subsection{Beukers's polynomial}
Let $N_p(\vartheta)(q)=\prod_{k=0}^{p-1}\vartheta(\zeta_p^kq^{1/p})$, a series in $q$, and
\[
D=\chi_\vartheta(p)\,\frac{\vartheta(q^p)\,N_p(\vartheta)(q)}{\vartheta(q)^{p+1}} .
\]
For $\vartheta=\Theta$ this is Beukers's polynomial $T_{p,p+1}$ \cite[\S\S4,5]{Beu26}.

\begin{lemma}\label{B:lem:D}
$D$ is a polynomial in $t$ with coefficients in $\Z$, $D\equiv\chi_\vartheta(p)\pmod p$, and $\deg D\le(p+1)/k_0$.
For $\tau_0\in\Hb\smallsetminus\Gamma e$, $D$ vanishes at $\tau_0$ to order $m(\tau_0)$ in the variable $\tau-\tau_0$.
\end{lemma}

\begin{proof}
As in Lemma~\ref{B:lem:qexp}, $D=p^{w_\vartheta(p-1)/2}\prod_{[\beta]}\chi_\vartheta(\beta)(\vartheta|\beta)/\vartheta$,
where $w_\vartheta$ is the weight of $\vartheta$. The argument of Lemma~\ref{B:lem:perm}(a),(b) shows that $D$ is invariant
under $\Gamma_0(N)$. In the second case $D$ is $W_2$-invariant by Lemma~\ref{B:lem:perm}(c). In the third case let
$g=W_3$ and $\beta=\bmat\alpha b{3c}\delta$. Then $\beta'=g^{-1}\beta g=\bmat\delta{-c}{-3b}\alpha\in M(3,p)$ and
$\chi_{-3}(\alpha)=\chi_{-3}(p)\chi_{-3}(\delta)$. From $a|_1g=-ia$ we get
$((a|\beta)/a)(g\tau)=((a|\beta')/a)(\tau)$, so the factor of $D$ indexed by $\beta$ at $g\tau$ equals
$\chi_{-3}(p)$ times the factor indexed by $\beta'$ at $\tau$. Since $\beta\mapsto\beta'$ permutes the classes
(Lemma~\ref{B:lem:hecke}(H4)) and $\chi_{-3}(p)^{p+1}=1$, $D$ is $W_3$-invariant. Thus $D$ is $\Gamma$-invariant. By
Lemma~\ref{B:lem:loc}, $D$ is holomorphic on $\Hb\smallsetminus\Gamma e$ and at the cusps and vanishes at
$\tau_0\notin\Gamma e$ to order $m(\tau_0)$; at $e$ the denominator $\vartheta^{p+1}$ vanishes to order $p+1$ and the
numerator to order $m(e)$, so $D$ has a pole of order $p+1-m(e)$ in $\tau-e$, i.e.\ of order $(p+1-m(e))/k_0$ in $t$.
Hence $D$ is a polynomial in $t$ of degree at most $(p+1)/k_0$. Its $q$-expansion lies in $\Z[[q]]$, so $D\in\Z[t]$.
Finally, $\prod_k\vartheta(\zeta_p^kz)\equiv\vartheta(z)^p\equiv\vartheta(z^p)$ modulo the ideal generated by
$1-\zeta_p$ in $\Z[\zeta_p][[z]]$. Both sides lie in $\Z[[z^p]]$, and $(1-\zeta_p)\Z[\zeta_p]\cap\Z=p\Z$, so
$N_p(\vartheta)(q)\equiv\vartheta(q)\pmod p$. Together with $\vartheta(q^p)\equiv\vartheta(q)^p$ this gives
$D\equiv\chi_\vartheta(p)\pmod p$.
\end{proof}

\subsection{The decay estimate}
Consider the $\Gamma$-invariant functions
\begin{gather*}
e_k=e_k\bigl(W_\beta\bigr)\ (0\le k\le p+1),\qquad
n_l=\sum_{[\beta]}x_\beta\,e_l\bigl(W_\gamma:[\gamma]\ne[\beta]\bigr)\ (0\le l\le p),\\
S_i=\sum_{[\beta]}x_\beta W_\beta^{\,i}\ (i\ge0),
\end{gather*}
where $e_k$ denotes the $k$-th elementary symmetric polynomial. Thus $S_i=\Tr_p(t^i)$, and by Lemma~\ref{B:lem:qexp}
\begin{equation}\label{B:eq:Si}
S_i=Q_{t^i}+p^{w+1}\chi(p)\,U_pG_{t^i}.
\end{equation}

\begin{lemma}\label{B:lem:sym}
\begin{enumerate}[label=\textup{(\alph*)}]
\item The $q$-expansions of $e_k$, $n_l$, $S_i$ lie in $\Z[[q]]$, and $e_1\equiv t^p\pmod p$.
\item $(-1)^ln_l=\sum_{k=0}^{l}(-1)^ke_kS_{l-k}$ for $0\le l\le p$, and $\sum_{k=0}^{p+1}(-1)^ke_kS_{l-k}=0$ for $l\ge p+1$.
\end{enumerate}
\end{lemma}

\begin{proof}
(a) The pair for $\beta_\infty$ is $(\Theta(q)/\Theta(q^p),t(q^p))$, and the pairs for $\beta_k$ are
$(\phi(\zeta_p^kq^{1/p}),t(\zeta_p^kq^{1/p}))$ with $\phi(z)=p^w\chi(p)\Theta(z^p)/\Theta(z)\in\Z[[z]]$
(proof of Lemma~\ref{B:lem:qexp}). A symmetric polynomial expression in the latter $p$ pairs is a power series in
$q^{1/p}$ with coefficients in $\Z[\zeta_p]$, invariant under $q^{1/p}\mapsto\zeta_pq^{1/p}$ and under
$\operatorname{Gal}(\Q(\zeta_p)/\Q)$, hence lies in $\Z[[q]]$. Moreover $e_1=t(q^p)+pU_pt\equiv t^p\pmod p$.
(b) Compare coefficients in $\bigl(\sum_k(-1)^ke_kY^{p+1-k}\bigr)\sum_{i\ge0}S_iY^{-i-1}=\prod_\beta(Y-W_\beta)\sum_\beta
\frac{x_\beta}{Y-W_\beta}=\sum_\beta x_\beta\prod_{\gamma\ne\beta}(Y-W_\gamma)$.
\end{proof}

\begin{proposition}[Decay of the Hecke trace]\label{B:prop:decay}
For all $i,n\ge0$,
\[
v_p\bigl([t^n]S_i\bigr)\ \ge\ \frac{k_0}{p+1}\,\bigl(n-pi-\lfloor pr\rfloor\bigr).
\]
\end{proposition}

\begin{proof}
\emph{Step 1: three families of polynomials.} Put $E_k=e_kD^{k_0}$, $M_l=n_lD^{k_0}$ and $A_0=S_0D^{o}$. We claim that
$E_k,M_l,A_0\in\Z[t]$ with
\[
\deg E_k\le p+1,\quad [t^{p+1}]E_1\equiv0\pmod p,\quad \deg M_l\le p,\quad \deg A_0\le\lfloor pr\rfloor .
\]
These functions are $\Gamma$-invariant and holomorphic at the cusps (Lemma~\ref{B:lem:loc}(c)). Let
$\tau_0\in\Hb\smallsetminus\Gamma e$ and $m=m(\tau_0)$. By Lemma~\ref{B:lem:loc}, exactly $m$ of the $W_\beta$ have a pole
at $\tau_0$, of order $k_0$, and the corresponding $x_\beta$ have poles of order $o\le k_0$; all other $W_\beta,x_\beta$ are
holomorphic. Hence $e_k$ and $n_l$ have poles of order at most $k_0m$ and $S_0$ of order at most $o$, while $D^{k_0}$
and $D^{o}$ vanish to orders $k_0m$ and $om\ge o$ (Lemma~\ref{B:lem:D}). So $E_k,M_l,A_0$ are holomorphic on
$\Hb\smallsetminus\Gamma e$, i.e.\ polynomials in $t$; their $q$-expansions are integral, so they lie in $\Z[t]$.
At $e$ let $m_e=m(e)$. By Lemma~\ref{B:lem:loc}, $m_e$ of the $W_\beta$ have a pole of order $k_0$ in $\tau-e$, i.e.\ of
order $1$ in $t$, and the corresponding $x_\beta$ are bounded; the other $W_\beta$ are bounded and the corresponding $x_\beta$
vanish to order $o$ in $\tau-e$, i.e.\ to order $r$ in $1/t$. The pole order of $D^{k_0}$ at $e$ is $p+1-m_e$ in $t$, and that
of $D^o$ is $r(p+1-m_e)$. Therefore, measuring pole orders in $t$:
$e_k$ has a pole of order $\le\min(k,m_e)$, so $\deg E_k\le\min(k,m_e)+p+1-m_e\le p+1$. Each term of $n_l$ has a pole of
order $\le\min(l,m_e-1)$ (if $\beta$ is one of the $m_e$ branches) or $\le\min(l,m_e)-r$ (otherwise), so
$\deg M_l\le p+1-m_e+\max(\min(l,m_e-1),\min(l,m_e)-r)\le p+1-r$, hence $\le p$. Finally $S_0$ is bounded at $e$, and vanishes
to order $r$ if $m_e=0$, so $\deg A_0\le r(p+1-m_e)-r[m_e=0]\le pr$. Lastly, $E_1\equiv\chi_\vartheta(p)^{k_0}e_1\equiv\pm t^p
\pmod p$ by Lemmas~\ref{B:lem:D} and \ref{B:lem:sym}(a), so $[t^{p+1}]E_1\equiv0\pmod p$.

\emph{Step 2: Gauss norms.} For $\rho\ge1$ and $f=\sum c_nt^n\in\Q_p[[t]]$ put $|f|_\rho=\sup_n|c_n|_p\rho^n$. This is
multiplicative and satisfies the ultrametric inequality. Let $R=p^{k_0/(p+1)}$ and $1\le\rho<R$. By
Lemma~\ref{B:lem:D}, $D=\pm(1+pC)$ with $C\in\Z[t]$ of degree $\le(p+1)/k_0$, so
$|pC|_\rho\le p^{-1}\rho^{(p+1)/k_0}<1$; hence $D^{-1}=\pm\sum_j(-pC)^j$ and $|D|_\rho=|D^{-1}|_\rho=1$. Consequently
\[
|e_1|_\rho=|E_1|_\rho\le\rho^p,\qquad |e_k|_\rho\le\rho^{p+1},\qquad|n_l|_\rho\le\rho^p,\qquad |S_0|_\rho\le\rho^{c_0},
\qquad c_0:=\lfloor pr\rfloor;
\]
for the first bound note $|[t^{p+1}]E_1|_p\rho^{p+1}\le p^{-1}\rho^{p+1}\le\rho^p$ because $\rho<R\le p$.

\emph{Step 3: induction.} We show $|S_l|_\rho\le\rho^{pl+c_0}$ for all $l\ge0$. For $l=0$ this is Step~2. For $l\ge1$,
Lemma~\ref{B:lem:sym}(b) gives $S_l=(-1)^ln_l-\sum_{k=1}^{\min(l,p+1)}(-1)^ke_kS_{l-k}$ with $n_l:=0$ for $l>p$. Here
$|n_l|_\rho\le\rho^p\le\rho^{pl+c_0}$, $|e_1S_{l-1}|_\rho\le\rho^{p+p(l-1)+c_0}$ and, for $k\ge2$,
$|e_kS_{l-k}|_\rho\le\rho^{p+1+p(l-k)+c_0}\le\rho^{pl+c_0}$.

\emph{Step 4.} Hence $|[t^n]S_i|_p\le\rho^{pi+c_0-n}$ for every $\rho\in[1,R)$. If $n\ge pi+c_0$ we let $\rho\to R$ and get
$|[t^n]S_i|_p\le p^{-k_0(n-pi-c_0)/(p+1)}$; if $n<pi+c_0$ the assertion is trivial since $S_i\in\Z[[t]]$.
\end{proof}

\begin{remark}
The bound is sharp in the sense that check~[8d] finds equality for many $(i,n)$. The rate $k_0/(p+1)$ is the
reciprocal of the number of Hecke translates of $e$ on $X(\Gamma)$ for $p$ inert in the CM field of $e$. These translates
are CM points of conductor $p$; for $p$ inert they lie on the circle $|t|_p=R$ (when $C$ has a unit leading coefficient),
which is where the trace has its poles. The proof does not need this description.
\end{remark}

\subsection{The truncation coefficients}
\begin{lemma}\label{B:lem:carry}
Let $p\ge5$ and $c=w+1$ ($c=2,3,3$ for $F_{27},F_{256},F_{108}$). If $n$ has the digit $p-1$ in exactly $d$ positions of
its base-$p$ expansion, then $v_p(a_n)\ge cd$.
\end{lemma}

\begin{proof}
By Legendre's formula $v_p\bigl((kn)!/n!^k\bigr)=\sum_{j\ge1}\lfloor k\{n/p^j\}\rfloor$, where $\{x\}$ is the fractional
part; all terms are $\ge0$. If the digit of $n$ at position $j-1$ is $p-1$, then $\{n/p^j\}\ge(p-1)/p$ and the $j$-th term is
$\ge\lfloor k(p-1)/p\rfloor=k-1$ for $p\ge k$. This gives $2d$ for $(3n)!/n!^3$ and $3d$ for $(4n)!/n!^4$. For
$(2n)!(3n)!/n!^5$ the $j$-th term is $\lfloor2\{n/p^j\}\rfloor+\lfloor3\{n/p^j\}\rfloor\ge1+2$.
\end{proof}

\begin{lemma}\label{B:lem:arith}
Let $p\ge5$, $(k_0,o,c)\in\{(3,1,2),(4,1,3),(6,2,3)\}$ and $r=o/k_0$. For $j\ge0$ let $\kappa=0$ if $j=0$ and
$\kappa=1+\lfloor\log_pj\rfloor$ (the number of base-$p$ digits of $j$) if $j\ge1$. Then
\[
\Bigl\lceil\frac{k_0}{p+1}\bigl(p(j+1)-\lfloor pr\rfloor\bigr)\Bigr\rceil\ \ge\ c(\kappa+1).
\]
\end{lemma}

\begin{proof}
The left side is at least $\lceil x\rceil$ with $x=\frac{p}{p+1}\bigl(k_0(j+1)-o\bigr)\ge\frac56\bigl(k_0(j+1)-o\bigr)$, and
it suffices that $x>c(\kappa+1)-1$. For $j=0$: $\frac56(k_0-o)=\frac53,\frac52,\frac{10}3$, which exceed $c-1=1,2,2$. For
$j\ge1$ we have $j+1\ge p^{\kappa-1}+1\ge5^{\kappa-1}+1$, and $\frac56\bigl(k_0(5^{\kappa-1}+1)-o\bigr)>c(\kappa+1)-1$ holds for
$\kappa=1$ ($\frac{25}6>3$, $\frac{35}6>5$, $\frac{50}6>5$). When $\kappa$ increases by one, the left side increases by
$\frac56k_0(5^{\kappa}-5^{\kappa-1})\ge\frac{10}3k_0\ge10$ and the right side by $c\le3$, so the inequality holds for all
$\kappa\ge1$.
\end{proof}

\subsection{Proof of Theorem~\ref{B:thm:ell}}
Fix $p\ge5$ and $m\ge1$, put $K_j=mp^j$, $F_{K_{-1}}=0$ and $c=w+1$. We prove
\[
(\mathrm H_s)\qquad \Theta(q)\,F_{K_{s-1}}\bigl(t(q^p)\bigr)\equiv\Theta(q^p)\,F_{K_s}\bigl(t(q)\bigr)\pmod{p^{cs}\Z[[q]]}
\]
by induction on $s\ge1$; as in \S\ref{B:sec:proof} it is equivalent to the theorem. Let $P=F_{K_{s-1}}=\sum_{i<K_{s-1}}a_it^i$
and $Q=Q_P$. By linearity and Lemma~\ref{B:lem:qexp},
\[
Q=\sum_{i<K_{s-1}}a_iS_i-p^{w+1}\chi(p)\,U_pG_P .
\]
Let $n\ge K_s$. As in the proof of the claim in \S\ref{B:sec:proof}, $(\mathrm H_{s-1})$ gives
$p^{w+1}U_pG_P\equiv p^{w+1}F_{K_{s-2}}(t)\pmod{p^{w+1+c(s-1)}}$, so $[t^n]\,p^{w+1}U_pG_P\equiv0\pmod{p^{cs}}$ (for $s=1$
trivially). Now let $i<K_{s-1}$ and $j=K_{s-1}-1-i$, with $\kappa$ as in Lemma~\ref{B:lem:arith}. Since
$n-pi\ge K_s-pi=p(j+1)$, Proposition~\ref{B:prop:decay} and Lemma~\ref{B:lem:arith} give
$v_p([t^n]S_i)\ge c(\kappa+1)$. If $\kappa\le s-1$, then $j<p^\kappa$ and
\[
i=(m-1)p^{s-1}+\bigl(p^{s-1}-p^{\kappa}\bigr)+\bigl(p^\kappa-1-j\bigr),\qquad 0\le p^\kappa-1-j<p^\kappa,
\]
so the digits of $i$ in the positions $\kappa,\dots,s-2$ are $p-1$, and $v_p(a_i)\ge c(s-1-\kappa)$ by
Lemma~\ref{B:lem:carry}. In all cases $v_p(a_i)+v_p([t^n]S_i)\ge cs$. Therefore $[t^n]Q\equiv0\pmod{p^{cs}}$ for all
$n\ge K_s$. On the other hand $[t^n]Q=a_n$ for $n<K_s$ by Lemma~\ref{B:lem:Ft}(c). Hence $Q\equiv F_{K_s}(t)\pmod{p^{cs}}$,
which is $(\mathrm H_s)$.\qed

\begin{remark}
The proof shows where the exponent comes from. The top monomials $t^i$, $i$ close to $mp^{s-1}$, contribute through
$[t^n]S_i$ with $n-pi\approx p$, which is only divisible by $p^{\lceil(k_0-o)p/(p+1)\rceil}$. This loss is exactly
compensated by the $c$ factors of $p$ per digit $p-1$ of $i$. For the families of Theorem~\ref{B:thm:main} the same
argument works trivially, because $S_i$ is then a polynomial of degree $<pi+p$.
\end{remark}

\section{Zagier's case D}\label{B:source-main-7}\label{B:sec:D}

In this section we prove Theorem~\ref{B:thm:D}. Let $a_n=\sum_k\binom nk^2\binom{n+k}k$ and $F=\sum a_nt^n$, and put
\[
t(\tau)=q\prod_{n\ge1}(1-q^n)^{5(\frac n5)},\qquad
\Theta(\tau)=\prod_{n\equiv\pm1\,(5)}(1-q^n)^{-3}\prod_{n\not\equiv\pm1\,(5)}(1-q^n)^{2},
\]
with $(\frac n5)$ the Legendre symbol. For $d\in(\Z/5\Z)^\times$ let $\sigma_d\in\Gamma_0(5)$ be a matrix with lower right
entry $\equiv d\pmod5$; the action of $\sigma_d$ on functions for $\Gamma_1(5)$ (the diamond operator) depends only on $d$.

\subsection{The modular data}
\begin{lemma}\label{B:lem:Ddata}
\begin{enumerate}[label=\textup{(\alph*)}]
\item $F(t(q))=\Theta(q)$, \ $\Theta^2t=\eta(5\tau)^5/\eta(\tau)$ \ and \ $1/t-11-t=\eta(\tau)^6/\eta(5\tau)^6$.
\item $t$ is a Hauptmodul for $\Gamma_1(5)$. It is holomorphic and nonzero on $\Hb$, has a simple zero at $\infty$, a simple
pole at the cusp $\frac25$, and the values $\frac{-11\pm5\sqrt5}2$ at the cusps $0,\frac12$. Moreover
$t\circ\sigma_{\pm2}=-1/t$ and $t\circ\sigma_{-1}=t$.
\item $\Theta$ is a modular form of weight one for $\Gamma_1(5)$, without zeros in $\Hb$, with a simple zero at $\frac25$
and no other zero at the cusps. There is $c\in\{\pm1\}$ with
\[
\Theta|_1\sigma_2=c\,t\Theta,\qquad \Theta|_1\sigma_3=-c\,t\Theta,\qquad \Theta|_1\sigma_{-1}=-\Theta .
\]
\end{enumerate}
\end{lemma}

\begin{proof}
(a) The first identity is Beukers's modular parametrisation of the Ap\'ery numbers for $\zeta(2)$ \cite{Beu87}, see also
\cite[case D]{Za09}; the third is Ramanujan's identity for the Rogers--Ramanujan continued fraction $r(\tau)$, since
$t=r^5$ \cite{Du05,Co17}. All three are checked to $O(q^{79})$ by check~[9a].

(b) That $t=r^5$ is a Hauptmodul for $\Gamma_1(5)$ is classical \cite{Du05}. It is holomorphic and nonzero on $\Hb$ by
its product expansion, and $t=q+O(q^2)$. The right side of the third identity in (a) is a Hauptmodul for $\Gamma_0(5)$,
with its only pole at the cusp $\infty$ of $X_0(5)$ and its only zero at the cusp $0$. Since
$X_1(5)\to X_0(5)$ has degree two and $t$ has degree one, $t$ is not $\Gamma_0(5)$-invariant; as $1/t-t$ is, we get
$t\circ\sigma_{\pm2}=-1/t$ and $t\circ\sigma_{\pm1}=t$. The fibre of $X_1(5)\to X_0(5)$ over $\infty$ consists of the
cusps $\infty$ and $\sigma_2\infty\sim\frac25$, both of width one, and over $0$ of $0$ and $\sigma_20=\frac12$, both of
width five. Hence $t$ has its pole at $\frac25$, and at $0$, $\frac12$ it takes the two roots of $t^2+11t-1=0$.

(c) $\eta(5\tau)^5/\eta(\tau)$ is a modular form of weight two for $\Gamma_0(5)$ with character $(\frac{\cdot}5)$
\cite[Thm.~1.64]{Ono04}, without zeros in $\Hb$; by Ligozat's formula \cite[Thm.~1.65]{Ono04} it vanishes to order one
(in $e^{2\pi i\tau}$) at the cusps $\infty,\frac25$ and does not vanish at $0,\frac12$. By (a) and (b),
$\Theta^2=\eta(5\tau)^5/(\eta(\tau)t)$ is therefore a meromorphic modular form of weight two for $\Gamma_1(5)$, holomorphic
and nonzero on $\Hb$, of order $0,2,0,0$ at $\infty,\frac25,0,\frac12$. The series $\Theta\in1+q\Z[[q]]$ is a holomorphic
square root of it on $\Hb$, so $\Theta|_1\gamma=\epsilon(\gamma)\Theta$ for $\gamma\in\Gamma_1(5)$, with a character
$\epsilon:\Gamma_1(5)\to\{\pm1\}$. The group $\Gamma_1(5)$ is torsion free, does not contain $-1$, and $X_1(5)$ has genus
zero with four cusps, all regular; so $\Gamma_1(5)$ is the free group generated by parabolic elements $A T^{h}A^{-1}$
fixing three of the cusps. Since the orders of $\Theta^2$ at the cusps are even, $\Theta|_1A$ is a power series in
$e^{2\pi i\tau/h}$ at each cusp $A\infty$ of width $h$, and $\epsilon$ is trivial on these generators. Thus $\epsilon=1$,
and $\Theta$ has the stated zeros.

For $d=\pm2$ we get from (a), (b) and $(\frac{\pm2}5)=-1$ that
$(\Theta|_1\sigma_d)^2=\Theta^2|_2\sigma_d=-\eta(5\tau)^5/(\eta(\tau)\cdot(-1/t))=t^2\Theta^2$, so
$\Theta|_1\sigma_d=\pm t\Theta$ with a constant sign. If $\Theta|_1\sigma_2=ct\Theta$, then, as $\sigma_3$ acts as
$\sigma_2^{-1}$, $\Theta=(ct\Theta)|_1\sigma_3=c\,(-1/t)\,\Theta|_1\sigma_3$, i.e.\ $\Theta|_1\sigma_3=-ct\Theta$. Finally
$\sigma_{-1}\in-\Gamma_1(5)$ and $\Theta|_1(-1)=-\Theta$. (Numerically $c=1$, check~[9a], but we do not need this.)
\end{proof}

\paragraph{Symbolic identification of the Zagier D coordinate.}
For the Zagier D parametrisation one similarly has the exact identity
\[
 \frac{Dt}{t}=(1-11t-t^2)\Theta^2.
\]
Indeed, the diamond relation $t\circ\sigma_2=-1/t$ shows that $Dt/t$
is a holomorphic weight-two form on $\Gamma_0(5)$ with quadratic
character $(\frac{\cdot}{5})$. The expression on the right has the
same transformation law: under $t\mapsto-1/t$ and
$\Theta^2\mapsto t^2\Theta^2$ it changes sign. The cusp orders in
Lemma~\ref{B:lem:Ddata} show that it is holomorphic at every cusp.
Both sides have expansion $1-5q+O(q^2)$. The total valence is
$2\cdot6/12=1$, so their difference, of order at least two at infinity,
vanishes identically. The Wronskian equation of
\[
 \bigl(\theta^2-t(11\theta^2+11\theta+3)
                    -t^2(\theta+1)^2\bigr)F=0
\]
gives $d\log q_F/dt=1/[t(1-11t-t^2)F(t)^2]$ for the normalized
Frobenius coordinate. Since $F(t(q))=\Theta(q)$ and $t=q+O(q^2)$,
the preceding modular identity proves $q_F(t(q))=q$. Thus the
mirror-map assertion does not rely on finite check [9a].

\subsection{Hecke operators on \texorpdfstring{$\Gamma_1(5)$}{Gamma1(5)}}
Let $p\ne5$ be a prime and $\Delta_p=\{\beta\in M_2(\Z):\det\beta=p,\ \beta\equiv\bmat1*0p\pmod5\}$. Then
$\Delta_p=\Gamma_1(5)\bmat100p\Gamma_1(5)$, and $\Gamma_1(5)\backslash\Delta_p$ has the $p+1$ representatives
$\beta_k=\bmat1k0p$, $0\le k<p$, and $\beta_\infty=\sigma_p\bmat p001$ \cite[Prop.~5.2.1]{DS05}. Clearly
$\gamma\Delta_p=\Delta_p\gamma=\Delta_p$ for $\gamma\in\Gamma_1(5)$, and $\sigma\Delta_p\sigma^{-1}=\Delta_p$ for
$\sigma\in\Gamma_0(5)$, because the diagonal entries modulo $5$ are unchanged by conjugation with a matrix that is upper
triangular modulo $5$.

Put $\check\Theta=t\Theta$ and $\check t=-1/t$. For $P\in\Z[T]$ of degree $D$ and $\beta\in\Delta_p$ let
\[
R_\beta=p^{1/2}\,\frac{\Theta\cdot P(t\circ\beta)}{\Theta|_1\beta},\qquad
\check R_\beta=p^{1/2}\,\frac{\check\Theta\cdot P(\check t\circ\beta)}{\check\Theta|_1\beta},\qquad
\Lambda_P=\sum_{[\beta]}R_\beta,\qquad \check\Lambda_P=\sum_{[\beta]}\check R_\beta,
\]
the sums over $\Gamma_1(5)\backslash\Delta_p$. Besides $Q_P$ and $G_P$ (with this $\Theta$ and $t$) we use
\[
\check Q_P(q)=\frac{\Theta(q)\,P\bigl(-1/t(q^p)\bigr)}{t(q^p)\,\Theta(q^p)},\qquad
\check G_P(z)=\frac{\Theta(z^p)\,P\bigl(-1/t(z)\bigr)}{t(z)\,\Theta(z)},
\]
which are Laurent series with integer coefficients.

\begin{lemma}\label{B:lem:Dtrace}
\begin{enumerate}[label=\textup{(\alph*)}]
\item $\Lambda_P$ and $\check\Lambda_P$ are well defined and $\Gamma_1(5)$-invariant, and $\check\Lambda_P=\Lambda_P\circ\sigma_2$.
\item If $p\equiv\pm1\pmod5$, then $\Lambda_P=\eta'_pQ_P+p^2U_pG_P$ with $\eta'_p=\pm1$ for $p\equiv\pm1$, and
$\Lambda_P$ is a polynomial in $t$ with integer coefficients of degree at most $p(D+1)-1$.
\item If $p\equiv\pm2\pmod5$, let $\eta_p=c$ for $p\equiv2$ and $\eta_p=-c$ for $p\equiv3\pmod5$. Then
\[
\Lambda_P=\eta_p\check Q_P+p^2U_pG_P,\qquad \check\Lambda_P=-\eta_p\,t\,Q_P+p^2\,t\,U_p\check G_P,
\]
$\Lambda_P$ and $\check\Lambda_P$ are Laurent polynomials in $t$ with integer coefficients, $\check\Lambda_P$ has degree at
most $p(D+1)$ in $t$, and $\check\Lambda_P(t)=\Lambda_P(-1/t)$.
\end{enumerate}
\end{lemma}

\begin{proof}
(a) Since $\Theta|_1\gamma=\Theta$ and $t\circ\gamma=t$ for $\gamma\in\Gamma_1(5)$, the proof of
Lemma~\ref{B:lem:perm}(a),(b) (with $\chi=1$) shows $R_{\gamma\beta}=R_\beta$ and $R_\beta\circ\gamma=R_{\beta\gamma}$, and
likewise for $\check R_\beta$. For the last claim let $\beta\in\Delta_p$ and $\beta'=\sigma_2^{-1}\beta\sigma_2\in\Delta_p$.
Then $t(\beta\sigma_2\tau)=t(\sigma_2\beta'\tau)=\check t(\beta'\tau)$, and with $\Theta|_1\sigma_2=c\check\Theta$,
\[
\frac{\Theta(\sigma_2\tau)}{(\Theta|_1\beta)(\sigma_2\tau)}=\frac{(\Theta|_1\sigma_2)(\tau)}{(\Theta|_1\beta\sigma_2)(\tau)}
=\frac{c\check\Theta(\tau)}{(\Theta|_1\sigma_2\beta')(\tau)}=\frac{\check\Theta(\tau)}{(\check\Theta|_1\beta')(\tau)} .
\]
Hence $R_\beta\circ\sigma_2=\check R_{\beta'}$, and $\beta\mapsto\beta'$ permutes $\Gamma_1(5)\backslash\Delta_p$.

(b), (c) As in Lemma~\ref{B:lem:qexp}, the classes $\beta_k$ contribute $p^2U_pG_P$ to $\Lambda_P$, and to
$\check\Lambda_P$ they contribute
\[
p^2U_p\Bigl(\frac{\check\Theta(z^p)\,P(\check t(z))}{\check\Theta(z)}\Bigr)=p^2\,t\,U_p\check G_P .
\]
For
$\beta_\infty$ we have $(\Theta|_1\beta_\infty)(\tau)=p^{1/2}(\Theta|_1\sigma_p)(p\tau)$ and
$t\circ\beta_\infty=(t\circ\sigma_p)(p\tau)$.
If $p\equiv\pm1$, Lemma~\ref{B:lem:Ddata} gives $R_{\beta_\infty}=\eta'_pQ_P$. If $p\equiv\pm2$, it gives
$R_{\beta_\infty}=\Theta(\tau)P(-1/t(p\tau))/(\eta_pt(p\tau)\Theta(p\tau))=\eta_p\check Q_P$, and, as
$\check\Theta|_1\sigma_p=(t\circ\sigma_p)\,\Theta|_1\sigma_p=-\eta_p\Theta$ and $\check t\circ\sigma_p=t$,
$\check R_{\beta_\infty}=t\Theta P(t(p\tau))/(-\eta_p\Theta(p\tau))=-\eta_ptQ_P$.

It remains to locate the poles. All $R_\beta$, $\check R_\beta$ are holomorphic on $\Hb$, because $\Theta$ has no zeros and
$t$ no poles or zeros there. Let $c'=A\infty$ be a cusp, $A\in\SL_2(\Z)$, and for $\beta\in\Delta_p$ write
$\beta A=A_\beta\bmat{a'}{b'}0{d'}$ with $A_\beta\in\SL_2(\Z)$, $a'd'=p$, as in Lemma~\ref{B:lem:branch}; then
$\beta c'=A_\beta\infty$, and exactly one class has $(a',d')=(p,1)$. For a $\Gamma_1(5)$-invariant $\varphi$ we have
$\nu_A(\varphi\circ\beta)=\frac{a'}{d'}\nu_{A_\beta}(\varphi)$, and $\nu_A(f|_1\beta)=\frac{a'}{d'}\nu_{A_\beta}(f)$ for a
form $f$ of weight one; here $\nu$ is the order in $e^{2\pi i\tau}$, not normalised by the width (proof of
Lemma~\ref{B:lem:branch}). If $c'\in\{0,\frac12\}$, the cusp $\beta c'$ again has denominator prime to $5$, so lies in the
class of $0$ or $\frac12$, where $t,\check t,\Theta,\check\Theta$ are finite and nonzero; all branches are holomorphic at $c'$.
Let $c'=x/(5y)$ with $\gcd(x,5y)=1$ be in the class of $\infty$ ($x\equiv\pm1$) or of $\frac25$ ($x\equiv\pm2$). Writing
$\beta=\bmat\alpha b{5\gamma}\delta$ with $\alpha\equiv1\pmod5$, we have $\beta(x,5y)^{\mathsf T}=a'(x',5y')^{\mathsf T}$
with $\gcd(x',5y')=1$ and $x'\equiv x/a'\pmod5$. So the branches with $a'=1$ preserve the classes, and the branch with
$a'=p$ multiplies the class by $p^{-1}$: it preserves the classes if $p\equiv\pm1$ and exchanges the classes of $\infty$
and $\frac25$ if $p\equiv\pm2$.

At $\frac25$ we have $\nu(t)=-1$, $\nu(\Theta)=1$, $\nu(\check t)=1$, $\nu(\check\Theta)=0$, and at $\infty$ the opposite:
$\nu(t)=1$, $\nu(\Theta)=0$, $\nu(\check t)=-1$, $\nu(\check\Theta)=1$. If $p\equiv\pm1$, all branches are holomorphic at
$\infty$, and at $\frac25$ the branch with $a'=p$ has $\nu(R_\beta)=1-p-pD$, the others $\nu(R_\beta)=1-\frac1p-\frac Dp$.
So $\Lambda_P$ is a polynomial in $t$ of degree at most $p(D+1)-1$. If $p\equiv\pm2$, then at $\frac25$ the branch with $a'=p$
has $\nu(\check R_\beta)=0-p\cdot1+pD\cdot(-1)=-p(D+1)$ and the others $\nu(\check R_\beta)\ge0$, while at
$\infty$ poles may occur. So $\check\Lambda_P$ is a Laurent polynomial in $t$ of degree at most $p(D+1)$. The
$\Gamma_1(5)$-invariant function $\Lambda_P$ is a rational function of $t$, and by (a) and $t\circ\sigma_2=-1/t$ we have
$\check\Lambda_P(t)=\Lambda_P(-1/t)$; hence $\Lambda_P(t)=\check\Lambda_P(-1/t)$ is a Laurent polynomial as well.
Integrality of the coefficients follows from the $q$-expansions, as $t\in q+q^2\Z[[q]]$.
\end{proof}

\subsection{Proof of Theorem~\ref{B:thm:D}}
Fix $p\ne5$ and $m\ge1$, and put $K_j=mp^j$, $F_{K_{-1}}=0$. Note that $\Theta\in1+q\Z[[q]]$ and $t\in q+q^2\Z[[q]]$, so
Lemma~\ref{B:lem:Ft} applies.

\emph{Case $p\equiv\pm1\pmod5$.} By Lemma~\ref{B:lem:Dtrace}(b), $Q_P=\eta'_p\Lambda_P-\eta'_pp^2U_pG_P$ with $\Lambda_P$ a
polynomial in $t$ of degree $\le p\deg P+p-1$. This is \eqref{B:eq:main} with $w=1$ and $\chi(p)$ replaced by $\eta'_p$, and the
proof of \S\ref{B:sec:proof} applies word for word.

\emph{Case $p\equiv\pm2\pmod5$.} We prove by simultaneous induction on $s\ge1$ the congruences
\begin{align*}
(\mathrm H_s)&\qquad \Theta(q)\,F_{K_{s-1}}\bigl(t(q^p)\bigr)\equiv\Theta(q^p)\,F_{K_s}\bigl(t(q)\bigr),\\
(\check{\mathrm H}_s)&\qquad t(q)\Theta(q)\,F_{K_{s-1}}\bigl(-1/t(q^p)\bigr)\equiv t(q^p)\Theta(q^p)\,F_{K_s}\bigl(-1/t(q)\bigr),
\end{align*}
modulo $p^{2s}$, as Laurent series in $q$. With $P=F_{K_{s-1}}$ they read $Q_P\equiv F_{K_s}(t)$ and
$t\check Q_P\equiv F_{K_s}(-1/t)$. Let $s\ge1$, assume $(\mathrm H_{s-1})$ and $(\check{\mathrm H}_{s-1})$ if $s\ge2$, and
let $P=F_{K_{s-1}}$, $D=K_{s-1}-1$.

\emph{Claim:} $p^2U_pG_P\equiv p^2F_{K_{s-2}}(t)$ and $p^2U_p\check G_P\equiv p^2F_{K_{s-2}}(-1/t)/t$ modulo $p^{2s}$. For
$s=1$ all four terms are $\equiv0\pmod{p^2}$. For $s\ge2$, the first congruence follows from $(\mathrm H_{s-1})$ as in
\S\ref{B:sec:proof}. For the second, $(\check{\mathrm H}_{s-1})$ in the variable $z$ says
$F_{K_{s-1}}(-1/t(z))\equiv t(z)\Theta(z)F_{K_{s-2}}(-1/t(z^p))/(t(z^p)\Theta(z^p))$ modulo $p^{2s-2}$, hence
$\check G_P(z)\equiv F_{K_{s-2}}(-1/t(z^p))/t(z^p)$, and we apply $U_p$.

\emph{Proof of $(\mathrm H_s)$.} By Lemma~\ref{B:lem:Dtrace}(c) and the claim,
\[
Q_P=-\eta_p\frac{\check\Lambda_P(t)}t+\eta_pp^2U_p\check G_P\equiv L(t):=-\eta_p\frac{\check\Lambda_P(t)}t+\eta_pp^2\,
\frac{F_{K_{s-2}}(-1/t)}t\pmod{p^{2s}},
\]
where $L$ is a Laurent polynomial whose terms of degree $\ge0$ have degree $\le p(D+1)-1=K_s-1$, and whose second summand
contains only negative powers of $t$. Since $Q_P\in F(t)+t^{K_s}\Z[[t]]$ (Lemma~\ref{B:lem:Ft}(c)), comparing coefficients
gives $[t^n]Q_P\equiv0$ for $n\ge K_s$, which is $(\mathrm H_s)$, and moreover the identity of Laurent polynomials
\begin{equation}\label{B:eq:Lcheck}
\check\Lambda_P(t)\equiv-\eta_p\,t\,F_{K_s}(t)+p^2F_{K_{s-2}}(-1/t)\pmod{p^{2s}}.
\end{equation}

\emph{Proof of $(\check{\mathrm H}_s)$.} By Lemma~\ref{B:lem:Dtrace}(c), \eqref{B:eq:Lcheck} with $t$ replaced by $-1/t$, and the
claim,
\begin{align*}
\eta_p\check Q_P&=\Lambda_P-p^2U_pG_P=\check\Lambda_P(-1/t)-p^2U_pG_P\\
&\equiv\frac{\eta_p}tF_{K_s}(-1/t)+p^2F_{K_{s-2}}(t)-p^2F_{K_{s-2}}(t)=\frac{\eta_p}tF_{K_s}(-1/t)
\end{align*}
modulo $p^{2s}$. Multiplying by $\eta_pt$ gives $(\check{\mathrm H}_s)$. This completes the induction.\qed

\begin{remark}
The two traces are the Hecke operators $T_p$ applied to the weight $-1$ forms $P(t)/\Theta$ and $P(-1/t)/(t\Theta)$, which
are exchanged by $\langle2\rangle$; for $p\equiv\pm2$ the operator $T_p$ involves $\langle p\rangle$ and therefore mixes the
family with its dual. The terms $p^2F_{K_{s-2}}$ cancel in the last display; this is why no condition on the signs is needed.
\end{remark}

\section{Limits of the modular method and an elliptic extension}\label{B:source-main-8}\label{B:sec:open}

\subsection{The hyperoctahedral family in dimension three}\label{B:ss:hyp3}
Let $c_k=\binom{2k}k\sum_j\binom kj^2\binom{2j}j$, the number of closed walks of length $2k$ on $\Z^3$, and
$F(u)=\sum c_ku^k$. The hyperoctahedral family with $n=3$ is $F(t^2)$, and by Lemma~\ref{B:lem:subst} Conjecture~BV for it
is \eqref{B:eq:dwork} for $F$ with $m=1$, modulo $p^{2s}$. The family is modular (check~[7b]): with the Hauptmodul
$v=(\eta_1\eta_6/\eta_2\eta_3)^{12}$ of $\Gamma_0(6)+6$ of the Ap\'ery family of Table~\ref{B:tab:fam},
\[
u=\frac v{(1+v)^2},\qquad F(u(q))=(1+v)\,\frac{\eta_2^7\eta_3^7}{\eta_1^5\eta_6^5},
\]
and $u(q)$ is the mirror map of $F$. The function $u$ is a Hauptmodul for the group $\Gamma_0(6)+$ generated by
$\Gamma_0(6)$ and all Atkin--Lehner involutions. Its pole is the fixed point $v=-1$ of the involution $v\mapsto1/v$, an
elliptic point of order $k_0=2$ of $\Gamma_0(6)+$, and $\Theta=(1+v)\eta_2^7\eta_3^7/\eta_1^5\eta_6^5$ has a simple zero there.
So this is a family of elliptic type with $w=2$, $o=1$, $r=\frac12$. (Numerically, the degree of the polynomial $D$ of
Lemma~\ref{B:lem:D} is $(p\mp1)/2$ according as $p\equiv\pm1\pmod3$ (check~[7d]), which indicates that the elliptic point has CM by
$\Q(\sqrt{-3})$.)

The proof of Proposition~\ref{B:prop:decay} applies and gives
$v_p([u^n]S_i)\ge\frac2{p+1}(n-pi-\lfloor p/2\rfloor)$ (check~[7e] confirms it for $p=5$). At the top monomials of the
truncation this yields only one factor $p$, and the coefficients $c_k$ gain only one factor $p$ per digit $p-1$, from
$\binom{2k}k$; the sum $\sum_j\binom kj^2\binom{2j}j$ is a $p$-adic unit for $k=p-1$ and $k=p^2-1$ (check~[7c],
$p\le13$). Hence, for the monomials at the top of the truncation, the lower bound for $v_p(c_i)+v_p([u^n]S_i)$ that the
argument of \S\ref{B:sec:ell} provides is only about $s$, not $2s$. In fact, for $p=5$ and $s=3$ the termwise bound
$\min_i\bigl(v_p(c_i)+v_p([u^n]S_i)\bigr)$ over $125\le n<150$ equals $5<6$, while the tail of the trace has valuation
$\ge6$ (check~[7e]). So the congruence modulo $p^{2s}$, which holds in all our tests (check~[7a]), requires cancellation
between the contributions of different monomials; this termwise estimate does not prove the congruence. The full bound is proved by Theorem~\ref{P:thm:main-all}. It is consistent with
\cite[Remark~1.14]{Beu26} that only the exponent $2s$, not $3s$, occurs: here $1/r=2<w+1$.

\subsection{Dimension four and higher}
For dimensions $n\ge4$, the modular arguments in this part do not provide a uniform proof for the three standard families. The positive-vertex and matrix arguments of Part~\ref{part:positive} prove the required coefficientwise congruence for every such dimension in the stated full-symmetry prime ranges, without assuming a modular parametrisation. The mirror-quintic simplicial period is included in that theorem.

\subsection{A conjecture}
Theorems~\ref{B:thm:main} and \ref{B:thm:ell} and the data for the hyperoctahedral family suggest the following.

\begin{conjecture}\label{B:conj:ell}
Let $(N,\Gamma,w,\chi,\Theta,t)$ satisfy \textup{(F1)--(F3)} and the conditions of Lemma~\ref{B:lem:elldata}\textup{(c)} for an
elliptic point $e$ of order $k_0$, with $r=o/k_0$. Then for every prime $p\nmid N$ with $p\ge5$ and all $m,s\ge1$,
congruence \eqref{B:eq:dwork} with the excellent lift holds modulo $p^{\min(w+1,\lfloor1/r\rfloor)s}$.
\end{conjecture}

For the families of Table~\ref{B:tab:ell}, this is Theorem~\ref{B:thm:ell}. Its three-dimensional hyperoctahedral specialization follows from Theorem~\ref{P:thm:main-all}. The assertion for arbitrary other elliptic-zero families is not claimed in this volume.

\section{Computational verification}\label{B:source-main-9}\label{B:sec:verif}
The script \texttt{verify\_excellent\_lift.py} (Appendix~\ref{B:app:script}; standard library only, about
two and a half minutes) works with exact integer and rational arithmetic; the three floating point checks are marked as
such. It performs the following checks.
\begin{enumerate}[label={[\arabic*]}]
\item $F(t(q))=\Theta(q)$ to $O(q^{60})$ for the 13 families of Table~\ref{B:tab:fam}, with $F$ given by
the closed formulas of \S\ref{B:ss:examples}.
\item Ligozat orders of $\Theta$ and $t$ at all cusps: zeros of $\Theta$ $=$ poles of $t$ $=$ the
listed cusps; the value of $r_\Theta$; condition $(\mathrm C_p)$ via Lemma~\ref{B:lem:Cp}; the character of $\Theta$; the
Atkin--Lehner symmetry of $\Theta$ and $t$.
\item Proposition~\ref{B:prop:poly}: for $P=F_{mp^{s-1}}$ and for random integer polynomials $P$, the
series $Q_P+p^{w+1}\chi(p)U_pG_P$, expanded in $t$, has all coefficients of degree
$>p\deg P+\lfloor r_\Theta(p-1)\rfloor$ equal to $0$, for $p\in\{2,3,5,7,11\}$, $p\nmid N$.
\item Theorem~\ref{B:thm:main}: the valuation of the difference in $(\mathrm H_s)$, for all primes
$p\le13$ with $p\nmid N$ and all pairs $(s,m)$ in $\{(1,1),(1,2),(2,1),(2,2),(3,1)\}$ with
$mp^s\le130$.
\item Corollary~\ref{B:cor:coef} for $mp^s\le250$.
\item Conjecture BV in its own normalisation for the two- and three-dimensional hypercubic
families. The canonical coordinate is computed from the Picard--Fuchs recurrence by the Frobenius
method, the excellent lift is $q\mapsto q^p$, and \eqref{B:eq:dwork} is checked with $m=2$.
\item The hyperoctahedral family with $n=3$ (\S\ref{B:ss:hyp3}): [7a] the congruence modulo $p^{2s}$ for small $p,s$, with
the canonical coordinate computed from the Picard--Fuchs recurrence; [7b] the modular parametrisation and the mirror map;
[7c] $v_p(c_{p^k-1})=k$; [7d] the degree of $D$; [7e] Proposition~\ref{B:prop:decay} and the failure of the termwise
bound for $p=5$, $s=3$.
\item Theorem~\ref{B:thm:ell}: [8a] the identities of Lemma~\ref{B:lem:elldata} and, in floating point, $\Theta(e)=0$;
[8b] Lemma~\ref{B:lem:D} for $p\le13$ (and $B=D^2$ for $F_{108}$); [8c] the polynomials $E_k,M_l,A_0$ of
Proposition~\ref{B:prop:decay} and their degrees for $p\in\{5,7\}$; [8d] Proposition~\ref{B:prop:decay} for $p\le13$; [8e]
the congruence $(\mathrm H_s)$ modulo $p^{(w+1)s}$ for $p\le13$ and $mp^s\le130$; [8f] Lemma~\ref{B:lem:arith} for $p<200$,
$\kappa\le12$, and Lemma~\ref{B:lem:carry} for sample $n$; [8g] Conjecture~BV for the simplicial families with $n=2,3$ in
the normalisation of \cite{BV3}, and the mirror map identity; [8h] the trace of $Q_1$ for $F_{27}$ and $p=5$ is not a
polynomial in $t$.
\item Theorem~\ref{B:thm:D}: [9a] Lemma~\ref{B:lem:Ddata}(a), the orders of $\eta(5\tau)^5/\eta(\tau)$ at the cusps, the
mirror map, and in floating point the action of the diamond operators; [9b] Lemma~\ref{B:lem:Dtrace} for
$p\in\{2,3,7,11,13,17,19\}$ and several $P$; [9c] $(\mathrm H_s)$ and $(\check{\mathrm H}_s)$ modulo $p^{2s}$ for
$p\in\{2,3,7,11,13,17\}$ and $mp^s\le125$.
\end{enumerate}
The supplied execution record reports 114 passing checks and is reproduced in Appendix~\ref{B:app:log}. These bounded checks, including the explicitly identified floating-point evaluations, supplement the symbolic proofs; they do not establish infinite identities by extrapolation.




% END RETAINED SOURCE B

\partstart{Alternative elliptic and cusp proofs; exact sharpness}{part:exceptional}{Norm-polynomial trace bounds, explicit Zagier D cusp analysis, and all retained coefficient obstructions.}
% BEGIN RETAINED SOURCE E
\section{Statements, normalization, and scope}\label{E:source-main-1}

Let
\[
 F(t)=\sum_{k\geq0}a_kt^k\in1+t\Z[[t]],\qquad
 F_M(t)=\sum_{0\leq k<M}a_kt^k.
\]
In a modular parametrization $F(t(q))=\Theta(q)$, with
$t(q)\in q+q^2\Z[[q]]$, define
\[
 \sigma_p(t(q))=t(q^p).
\]
The inverse series $q(t)$ is integral, so
$\sigma_p(t)\in t^p+t^{p+1}\Z[[t]]$ and
$\sigma_p(t)\equiv t^p\pmod p$.  Every congruence in this note is an
identity of complete formal power series, coefficient by coefficient.
In particular, it is stronger than checking one coefficient or evaluating
at one fixed parameter.

For later use define the cleared defect
\begin{equation}\label{E:eq:defect}
 E_{p,s,m}(t)=F(t)F_{mp^{s-1}}(\sigma_p(t))
             -F(\sigma_p(t))F_{mp^s}(t).
\end{equation}
Since all denominator series have constant term one, the two statements
\begin{equation}\label{E:eq:quotient}
 \frac{F(t)}{F(\sigma_p(t))}
 \equiv\frac{F_{mp^s}(t)}{F_{mp^{s-1}}(\sigma_p(t))}
 \pmod{p^{\rho s}\Zp[[t]]},\qquad
 E_{p,s,m}\in p^{\rho s}\Zp[[t]]
\end{equation}
are equivalent.

\subsection{The cases proved in this note}

\begin{theorem}[Main conclusions]\label{E:thm:summary}
For every $p\geq5$ and all $m,s\geq1$, \eqref{E:eq:quotient} holds for the
following three series, with the indicated exponent:
\[
\begin{array}{c|c|c}
F(t)&a_k&\rho\\ \hline
{}_2F_1(1/3,2/3;1;27t)&(3k)!/(k!)^3&2\\
{}_3F_2(1/4,1/2,3/4;1,1;256t)&(4k)!/(k!)^4&3\\
{}_3F_2(1/3,1/2,2/3;1,1;108t)&(3k)!(2k)!/(k!)^5&3
\end{array}
\]
For Zagier D,
\[
 a_k=\sum_{j=0}^k\binom kj^2\binom{k+j}{j},
\]
the same assertion holds with $\rho=2$ for every prime
$p\equiv\pm1\pmod5$ and all $m,s\geq1$.
For every prime $p\ne5$, its first level $s=1$ holds for all $m\geq1$.
\end{theorem}

The elliptic cases are proved in Section~\ref{E:sec:elliptic}; the proof for
Zagier D is in Section~\ref{E:sec:zagier-d}.  These assertions are supported by
the proofs, not by extrapolation from the accompanying calculations.
Part~\ref{part:modular} gives the modular framework and an alternative elliptic trace argument. The norm-polynomial estimate below is retained as a separate proof.

\subsection{Relation to the original simplicial conjecture}

For the simplicial Laurent polynomial in $d$ variables,
\[
 g(x)=x_1+\cdots+x_d+(x_1\cdots x_d)^{-1},\qquad r=d+1,
\]
the original period is
\[
 \mathcal F(z)=\sum_{k\geq0}\frac{(rk)!}{(k!)^r}z^{rk}=F(z^r).
\]
Beukers--Vlasenko's canonical coordinate is
$q_0(z)=z\exp(G_0(z)/\mathcal F(z))$, where
\[
 G_0(z)=\sum_{k\geq1}\frac{(rk)!}{(k!)^r}
                    (H_{rk}-H_k)z^{rk}.
\]
The normalized coordinate therefore satisfies
\[
 q(z^r)=q_0(z)^r.
\]
Consequently the lifts are compatible:
$(z^{\sigma})^r=\sigma_p(z^r)$, and
$\mathcal F_{rM}(z)=F_M(z^r)$.
Thus the normalized case $m=1$ gives exactly the original case $m=d+1$,
not a differently truncated assertion.  For $d=2,3$ the symmetry and
support-lattice hypotheses of \cite[Theorem 7.3]{BV3} exclude precisely
the primes $2,3$; the conclusions above cover every remaining prime.

\subsection{Relation to the unified results}

\begin{center}\small
\begin{tabular}{@{}p{0.29\textwidth}p{0.61\textwidth}@{}}
\toprule
Target&Role of this part and full conclusion\\ \midrule
Simplicial $d=2,3$&Full conjectured congruence for all admissible primes and all levels; the normalized $d=3$ result has exponent $3s$.\\[3pt]
The $108t$ family&Full formal-series congruence modulo $p^{3s}$ for $p\geq5$ and all $m,s$.\\[3pt]
Zagier D&Full $p^{2s}$ congruence for $p\equiv\pm1\pmod5$; first level for every $p\ne5$.  All classes and levels are treated in Theorem~\ref{T:thm:zagier-nonsplit-complete}.\\[3pt]
Hyperoctahedral $d=3$&An exact coefficient has valuation $2s$ at $p=5$ for every $s$.  The matching full lower bound is Theorem~\ref{P:thm:main-all}.\\[3pt]
Dimensions $d\geq4$&All-level sharpness obstructions for hypercubes, eventual sharpness for simplices, and a dimension-four octahedral obstruction.  The full $p^{2s}$ statement is Theorem~\ref{P:thm:main-all}.\\
\bottomrule
\end{tabular}
\end{center}

The last two rows give upper bounds on a possible uniform exponent.
The matching lower bound is supplied by Theorem~\ref{P:thm:main-all}; together they give optimality in the explicitly stated ranges.
No assertion about priority in the literature is needed for the
mathematical arguments below.

\section{Elliptic zeros: a rational trace estimate}\label{E:source-main-2}\label{E:sec:elliptic}

We prove the three elliptic-zero cases of Theorem~\ref{E:thm:summary}.
All congruences of series are coefficientwise.
Put $F_M(t)=\sum_{0\leq k<M}a_kt^k$ when $F(t)=\sum_{k\geq0}a_kt^k$.

\subsection{Modular input and notation}

Let $p\geq5$ be prime.  In the three cases below the prime is automatically
prime to the modular level.  We use the following modular data:
\[
\begin{array}{c|c|c|c|c}
 F(t)&w&e&\Gamma&a_k\\ \hline
 {}_2F_1(1/3,2/3;1;27t)&1&3&\Gamma_0(3)&(3k)!/(k!)^3\\
 {}_3F_2(1/4,1/2,3/4;1,1;256t)&2&4&\Gamma_0(2)+2&(4k)!/(k!)^4\\
 {}_3F_2(1/3,1/2,2/3;1,1;108t)&2&3&\Gamma_0(3)+3&(3k)!(2k)!/(k!)^5
\end{array}
\]
In each case $t(q)\in q+q^2\mathbb Z[[q]]$ is a Hauptmodul,
$\Theta(q)=F(t(q))\in1+q\mathbb Z[[q]]$ is a modular form of weight $w$,
and $\Theta$ has just one zero on the compactified modular curve, at the
unique pole of $t$.  The ratio of these orders is $1/e$.  Thus
$C=t\Theta^e$ is a holomorphic modular form.  More explicitly, writing
$h_N=(\eta_N/\eta_1)^{24/(N-1)}$ for $N=2,3$, the three pairs are
\[
\begin{array}{c|c|c|c}
 &t&\Theta&C=t\Theta^e\\ \hline
27t&h_3/(1+27h_3)&\theta_H&\eta_3^9/\eta_1^3\\
256t&h_2/(1+64h_2)^2&2E_2(2\tau)-E_2(\tau)&(\eta_1\eta_2)^8\\
108t&h_3/(1+27h_3)^2&(3E_2(3\tau)-E_2(\tau))/2&(\eta_1\eta_3)^6.
\end{array}
\]
These are the modular parametrizations in
\cite[Appendices B and C]{Beu26}.
The form $\theta_H$ has quadratic character $\chi_{-3}$;
in the two weight-two cases the character is trivial on $\Gamma_0(N)$
and may have a sign on the Atkin--Lehner extension.
Set $\epsilon=\chi_{-3}(p)$ in weight one and $\epsilon=1$ in weight two.

Use the $p+1$ Hecke branches
\[
 t_\infty=t(p\tau),\quad t_b=t((\tau+b)/p),\qquad
 \Theta_\infty=\Theta(p\tau),\quad
 \Theta_b=\Theta((\tau+b)/p)\quad(0\leq b<p),
\]
and similarly $C_i=C(\beta_i\tau)$, so $t_i=C_i/\Theta_i^e$.
The associated modular-form ratios are
\[
 x_\infty=\frac{\Theta}{\Theta_\infty},\qquad
 x_b=\epsilon p^w\frac{\Theta}{\Theta_b}.
\]
The Hecke action permutes the pairs $(x_i,t_i)$, including under the
Atkin--Lehner extensions.  This follows directly from the weight-$w$
slash operator, as in \cite[Lemmas 3.6 and 3.7]{Beu26}.

Define
\[
 D(t)=\frac{\Theta_\infty\prod_{b=0}^{p-1}\Theta_b}{\Theta^{p+1}},
 \qquad a=\left\lfloor\frac pe\right\rfloor,
 \qquad \delta=\left\lfloor\frac{p+1}e\right\rfloor.
\]
Here $D(0)=1$.

\begin{lemma}[Norm-polynomial facts]\label{E:lem:ell-norm}
The following assertions hold.
\begin{enumerate}
\item $D\in\mathbb Z[t]$, $\deg D\leq\delta$, and $D\equiv1\pmod p$.
\item The polynomial
\[
 T(X)=\left(1-\frac{\Theta_\infty}{\Theta}X\right)
       \prod_{b=0}^{p-1}\left(p^w-\epsilon\frac{\Theta_b}{\Theta}X\right)
      =\epsilon D\prod_i(X-x_i)
\]
has coefficients $T_j(t)\in\mathbb Z[t]$ satisfying
$\deg T_j\leq\lfloor j/e\rfloor$.
\item Writing
\[
 D^e\prod_i(X-t_i)=D^eX^{p+1}+A_1(t)X^p+\cdots+A_{p+1}(t),
\]
we have $A_j\in\mathbb Z[t]$, $\deg A_j\leq p+1$, and
$A_1(t)\equiv-t^p\pmod p$.
\item For every $k\geq0$ the rational trace
\[
 R_k(t)=\sum_i x_it_i^k
\]
belongs to $D^{-ek-1}\mathbb Z[t]$.  Its numerator after multiplication
by $D^{ek+1}$ has degree at most $(p+1)k+a$.
\end{enumerate}
\end{lemma}

\begin{proof}
Since $\prod_i x_i=\epsilon p^{wp}/D$, the Hecke permutation of the
$x_i$ first implies the invariance of $D$.  It then makes the
coefficients of $T=\epsilon D\prod_i(X-x_i)$ invariant under
$\Gamma$.  Each coefficient of $X^j$ is holomorphic away from the zero
of $\Theta$, and its pole order relative to $t$ is at most $j/e$.
There are no poles at the cusps, since $\Theta$ has no cusp zeros.
Consequently it is a polynomial in the Hauptmodul $t$ of the stated
degree.  Its $q$-expansion is integral: products of the branches with
arguments $\zeta_p^bq^{1/p}$ are invariant under the cyclic rotation and
under the Galois group of $\mathbb Q(\zeta_p)$, so their coefficients
are rational algebraic integers and their exponents are integers.
Division by $\Theta\in1+q\mathbb Z[[q]]$ preserves integrality.
Since the inverse series $q(t)$ has integral coefficients, a polynomial
in $t$ with integral $q$-expansion has integral coefficients.
The leading coefficient of $T$ is $\epsilon D$ because $p$ is odd.
This proves the polynomiality, degree and integrality assertions for
$D$ and $T$.

The elementary norm identity
\[
 \prod_{b=0}^{p-1}\Theta(\zeta_p^bq^{1/p})\equiv\Theta(q)\pmod p
\]
holds for every integral series $\Theta$: reduce first modulo
$1-\zeta_p$, use the freshman's dream, and then use
$(1-\zeta_p)\mathbb Z[\zeta_p]\cap\mathbb Z=p\mathbb Z$.
Together with $\Theta(q^p)\equiv\Theta(q)^p\pmod p$ this gives
$D\equiv1\pmod p$.

For the third assertion use the exact identity
\[
 D^e\prod_i(X-t_i)
   =\frac{\prod_i(X\Theta_i^e-C_i)}{\Theta^{e(p+1)}}.
\]
Its coefficients are invariant because $D$ and the elementary
symmetric functions in the $t_i$ are invariant.  Every numerator is
holomorphic (all branches of $\Theta$ and $C$ are holomorphic modular
forms for a common finite-index subgroup).  At each cusp this means
holomorphy after the appropriate weight slash transformation; the
weights of the numerator and denominator agree, so their automorphy
factors cancel.  The denominator has
its only zero at the pole of $t$.  Its pole order there, measured on
$X(\Gamma)$, is at most $e(p+1)/e=p+1$.  Thus the coefficients are
polynomials of degree at most $p+1$.  They have integral $q$-expansions
because the elementary symmetric functions in the $t_i$ do, and $D$
is integral; hence they lie in $\mathbb Z[t]$.  Finally,
\[
 \sum_i t_i=t(q^p)+p\,U_pt(q)\equiv t(q)^p\pmod p,
\]
so $A_1=-D^e\sum_i t_i\equiv-t^p\pmod p$.

For the last assertion multiply each summand of $R_k$ by
$D^{ek+1}$.  Up to the constant factor $1$ or $\epsilon p^w$,
the resulting summand is
\[
 \frac{C_i^k\prod_{j\ne i}\Theta_j^{ek+1}}
 {\Theta^{(p+1)(ek+1)-1}}.
\]
The same holomorphy and invariance argument bounds the degree of the
sum by
\[
 \left\lfloor\frac{(p+1)(ek+1)-1}{e}\right\rfloor
   =(p+1)k+a.
\]
Its coefficients are integral, since $R_k$ has the integral
$q$-expansion displayed in \eqref{E:eq:ell-trace} below and $D$ is integral.
\end{proof}

\begin{lemma}[Elliptic trace filtration]\label{E:lem:ell-filtration}
If $R_k(t)=\sum_{n\geq0}r_{k,n}t^n$, then
\[
 v_p(r_{k,n})\geq
 \left\lceil\frac{n-pk-a}{\delta}\right\rceil.
\]
In particular, modulo $p^u$ the trace $R_k$ is a polynomial of degree
at most $pk+a+\delta(u-1)$ for every integer $u\geq1$.
\end{lemma}

\begin{proof}
Equip $\mathbb Q_p(t)$ with a Gauss norm for which $\lVert t\rVert=R$,
where $1\leq R<p^{1/\delta}$.  It suffices to take $\log_pR$ rational.
Complete this field in its Gauss norm and extend the norm to a common
algebraic splitting field containing all paired branches $(x_i,t_i)$.
Because $D=1+pB$, $\deg B\leq\delta$, and $D(0)=1$, one has
$\lVert D\rVert=\lVert D^{-1}\rVert=1$.  In the monic polynomial
$T/(\epsilon D)$ for the $x_i$, the coefficient of $X^{p+1-j}$ has
norm at most $R^a\leq R^{aj}$, since its numerator has degree
$\leq\lfloor(p+1-j)/e\rfloor\leq a$.  The nonarchimedean root bound
therefore gives $|x_i|\leq R^a$ for every $i$ in any valued splitting
field.

The monic polynomial for the $t_i$ has coefficients $A_j/D^e$.
For $j\geq2$, their norms are at most $R^{p+1}\leq R^{pj}$.
For $j=1$ the possible term of degree $p+1$ in $A_1$ is divisible
by $p$, while all other terms have degree at most $p$.
Since $R<p^{1/\delta}\leq p$, this gives
$\lVert A_1/D^e\rVert\leq R^p$ as well.  Thus the root bound gives
$|t_i|\leq R^p$.  Consequently
\[
 \lVert R_k\rVert=\left|\sum_i x_it_i^k\right|
     \leq R^{pk+a}.
\]
By Lemma~\ref{E:lem:ell-norm}, $R_k$ has only a power of $D$ as
denominator, so its Taylor series at zero converges in this Gauss
norm.  Hence $|r_{k,n}|_pR^n\leq R^{pk+a}$.
Letting $R\uparrow p^{1/\delta}$ and using the integrality of
$v_p(r_{k,n})$ proves the bound.
\end{proof}

\subsection{Carries in the factorial coefficients}

Let $K=mp^{s-1}$, $0\leq k<K$, $j=K-k\geq1$.
For $e=3$ or $4$ put $h=\lfloor\log_p(ej-1)\rfloor$.
Legendre's formula gives
\begin{equation}\label{E:eq:ell-carry}
 v_p\left(\frac{(ek)!}{(k!)^e}\right)
 \geq(e-1)\max(0,s-1-h).
\end{equation}
Indeed, for $h<i\leq s-1$ we have $p^i\mid K$ and $ej\leq p^i$;
the residue of $k$ modulo $p^i$ is $p^i-j$, and the corresponding
summand in Legendre's formula is exactly $e-1$.  All other summands
are nonnegative.

For $e=3,4$ let $r=e-1$.  The elementary inequalities
\[
 h\leq j-1,\qquad a\leq\delta,\qquad r\delta\leq p-1
\]
hold for every prime $p\geq5$.  The first follows from
$p^j\geq pj>ej-1$; the last follows immediately by considering
$p$ modulo $e$.  Using \eqref{E:eq:ell-carry}, the degree bound from
Lemma~\ref{E:lem:ell-filtration} for $a_kR_k$ modulo $p^{rs}$ is at most
\[
 pk+a+\delta\{r(h+1)-1\}.
\]
When $a_k$ is already zero modulo $p^{rs}$ the summand is discarded.
The excess of this bound over $pK$ is strictly negative, because
\[
 a-pj+\delta\{r(h+1)-1\}
 =r\delta(h+1)+(a-\delta)-pj
 \leq(p-1)j-pj<0.
\]
It follows that
\begin{equation}\label{E:eq:ell-poly-27-256}
 \sum_{k<K}a_kR_k(t)\pmod{p^{rs}}
 \quad\hbox{is a polynomial of degree strictly less than }pK
\end{equation}
for the $27t$ case with $r=2$ and the $256t$ case with $r=3$.

For the $108t$ case put $h=\lfloor\log_p(3j-1)\rfloor$.
Each index $h<i\leq s-1$ contributes $2+1=3$ to the valuation of
$a_k=(3k)!(2k)!/(k!)^5$.  Thus
\[
 v_p(a_k)\geq3\max(0,s-1-h).
\]
The excess degree modulo $p^{3s}$ is consequently bounded by
\begin{equation}\label{E:eq:ell-excess108}
 E=a-pj+\delta(3h+2).
\end{equation}
If $p\equiv1\pmod3$, then $a=\delta=(p-1)/3$ and
$E=(p-1)(h+1)-pj<0$.
If $p\equiv2\pmod3$, write $p=3a+2$, so $\delta=a+1$ and
\[
 E=(p+1)h-p(j-1).
\]
For $h=0$ this is nonpositive, with equality only when $j=1$.
For $h=1$, the inequality $j\geq\lceil(p+1)/3\rceil$ makes it
strictly negative for $p\geq11$.  When $p=5$, it is $11-5j$;
the only exception to strict negativity is $j=2$.
For $h\geq2$, the inequality $j\geq(p^h+1)/3$ and
\[
 p(p^h-2)>3(p+1)h\qquad(p\geq5,\ h\geq2)
\]
give strict negativity.  The displayed inequality holds at $p=5,h=2$
and is preserved as either $p$ or $h$ increases.

It remains to repair $p=5,j=2$.  If $s=1$, use $v_p(a_k)\geq0$
directly: the excess degree is at most $1-10+2(3-1)=-5$.
If $s\geq2$, then $k=K-2\equiv3\pmod5$; the $i=1$ summand in
Legendre's formula is $\lfloor9/5\rfloor+\lfloor6/5\rfloor=2$,
and all indices $2\leq i\leq s-1$ contribute $3$.  Therefore
$v_5(a_k)\geq3(s-1)-1$, and the excess degree is at most
$1-10+2\cdot3=-3$.
We have proved that
\begin{equation}\label{E:eq:ell-poly108}
 \sum_{k<K}a_kR_k(t)\pmod{p^{3s}}
 \quad\hbox{is a polynomial of degree at most }pK.
\end{equation}

\subsection{The induction and the boundary coefficient}

The exact trace identity for a polynomial $P$ is
\begin{equation}\label{E:eq:ell-trace}
 \sum_i x_iP(t_i)=Q_P+\epsilon p^{w+1}U_pG_P,
 \qquad
 Q_P=\frac{\Theta(q)P(t(q^p))}{\Theta(q^p)},\quad
 G_P(q)=\frac{\Theta(q^p)P(t(q))}{\Theta(q)}.
\end{equation}
It follows by summing the $p$ branches $q\mapsto\zeta_p^bq^{1/p}$.

\begin{theorem}[The three elliptic-zero families]\label{E:thm:ell-all-s}
Let $\sigma(t(q))=t(q^p)$ be the excellent Frobenius lift in one of
the three displayed families.  For all primes $p\geq5$ and all
integers $m,s\geq1$,
\[
 \frac{F(t)}{F(\sigma(t))}
 \equiv\frac{F_{mp^s}(t)}{F_{mp^{s-1}}(\sigma(t))}
 \pmod{p^{\rho s}\mathbb Z_p[[t]]},
\]
where $\rho=2$ for $27t$ and $\rho=3$ for $256t$ and $108t$.
\end{theorem}

\begin{proof}
Fix $m$, set $K=mp^{s-1}$, and take $P=F_K$ in
\eqref{E:eq:ell-trace}.  For $s=1$, the second term on its right is
zero modulo $p^\rho$, since $w+1=\rho$.
For $s\geq2$, the induction hypothesis is equivalent to
\[
 G_{F_K}(q)\equiv F_{K/p}(t(q^p))\pmod{p^{\rho(s-1)}}.
\]
Hence
\[
 \epsilon p^{w+1}U_pG_{F_K}
 \equiv\epsilon p^\rho F_{K/p}(t(q))\pmod{p^{\rho s}},
\]
which is a polynomial in $t$ of degree less than $pK$.
Equations \eqref{E:eq:ell-poly-27-256} and \eqref{E:eq:ell-poly108}
therefore show that $Q_{F_K}$ is congruent modulo $p^{\rho s}$
to a polynomial of degree $<pK$ in the first two cases and
of degree $\leq pK$ in the $108t$ case.

Since $\sigma(t)=t^p+O(t^{p+1})$,
\[
 Q_{F_K}(t)
 =F(t)\frac{F_K(\sigma(t))}{F(\sigma(t))}
 =F(t)-a_Kt^{pK}+O(t^{pK+1}).
\]
Thus $Q_{F_K}\equiv F_{pK}$ modulo $p^{\rho s}$ immediately
in the $27t$ and $256t$ cases.

In the $108t$ case the only possible remaining coefficient is
the one at degree $pK$, equal to $a_{pK}-a_K$.
The factorization
\[
 a_K=\binom{3K}{K}\binom{2K}{K}^2
\]
and Jacobsthal's binomial congruence give
\[
 a_{mp^s}\equiv a_{mp^{s-1}}\pmod{p^{3s}}
 \qquad(p\geq5).
\]
For example, apply \cite[Lemma 2.1]{St23} with equal exponents $r=s$ to
each binomial factor.  Thus the final coefficient also vanishes.
This completes the induction.  Finally divide by the unit
$F_K(\sigma(t))$ to obtain the asserted quotient congruence.
\end{proof}

\paragraph{Scope.}
This proves the normalized $27t$ and $256t$ simplicial cases and
the requested $108t$ family.  Passing from the original simplicial
period $F(t^{n+1})$ to the normalized variable gives Conjecture 7.5
for simplicial dimensions $n=2,3$ at all primes allowed by the
original symmetry/lattice hypothesis (here $p\geq5$).
This particular rational-trace proof concerns the three displayed families. The other standard cases are proved in Part~\ref{part:positive}, and full Zagier D is proved in Parts~\ref{part:modular} and~\ref{part:twocusp}.

\section{Zagier D at primes congruent to \texorpdfstring{$\pm1$}{plus or minus one} modulo five}\label{E:source-main-3}
\label{E:sec:zagier-d}

Let
\[
 A_n=\sum_{k=0}^{n}\binom nk^2\binom{n+k}{k},
 \qquad F(t)=\sum_{n\ge0}A_nt^n,
 \qquad F_M(t)=\sum_{0\le n<M}A_nt^n.
\]
The modular parametrization in \cite[Appendix B, \S8.8]{Beu26} is
\begin{equation}\label{E:eq:zagier-data}
 t(q)=q\prod_{n\ge1}(1-q^n)^{5(\frac n5)},\qquad
 F(t(q))=\Theta(q),\qquad
 t\Theta^2=\frac{\eta(5\tau)^5}{\eta(\tau)}.
\end{equation}
Here \(q=e^{2\pi i\tau}\), \(t\in q+q^2\mathbb Z[[q]]\), and
\(\Theta\in1+q\mathbb Z[[q]]\).  Define the excellent lift by
\(\sigma_p(t(q))=t(q^p)\).

\begin{theorem}[Zagier D for two residue classes]\label{E:thm:zagier-split}
For every prime \(p\equiv1\) or \(-1\pmod5\), and all integers
\(m,s\ge1\), one has
\[
 \frac{F(t)}{F(\sigma_p(t))}
 \equiv
 \frac{F_{mp^s}(t)}{F_{mp^{s-1}}(\sigma_p(t))}
 \pmod{p^{2s}\mathbb Z_p[[t]]}.
\]
\end{theorem}

This theorem covers all truncation lengths and all levels \(s\) in
the indicated prime classes.  The final subsection also proves the
first level for \(p\equiv\pm2\pmod5\) by an independent Laurent-support argument. Every higher level in these classes follows from Theorem~\ref{T:thm:zagier-nonsplit-complete}.

\subsection{The zero cusp and its divisor}

Put
\[
 H=\pm\Gamma_1(5)
   =\left\{\begin{pmatrix}a&b\\c&d\end{pmatrix}\in{\rm SL}_2(\mathbb Z):
        c\equiv0,\quad a\equiv d\equiv\pm1\pmod5\right\}.
\]
For \(\gamma\in H\), let \(\epsilon(\gamma)\in\{\pm1\}\) be the
common residue of its diagonal entries modulo 5.  Then
\[
 t(\gamma\tau)=t(\tau),\qquad
 \Theta|_1\gamma=\epsilon(\gamma)\Theta.
\]
Throughout this section the slash operator is determinant normalized:
\[
 (f|_w\beta)(\tau)=
 \det(\beta)^{w/2}(c\tau+d)^{-w}f(\beta\tau).
\]
Let
\[
 \delta=\begin{pmatrix}-2&-1\\5&2\end{pmatrix},
 \qquad c_*=\delta\infty=-2/5.
\]
The four cusps of \(H\) are represented by
\(\infty,c_*,0,\delta0=-1/2\), with widths \(1,1,5,5\), respectively.

\begin{lemma}[Divisor data]\label{E:lem:zagier-divisors}
The form \(\Theta\) has a simple zero at \(c_*\) and no other zeros,
including in the upper half plane.  Moreover,
\[
 {\rm div}(t)=[\infty]-[c_*].
\]
Consequently \(t\) is a Hauptmodul for \(H\).
\end{lemma}

\begin{proof}
Let \(\chi\) be the quartic character modulo 5 with \(\chi(2)=i\).
The standard holomorphic Eisenstein series of weight one
\[
 E_\chi(q)=\frac{3+i}{10}
     +\sum_{n\ge1}\left(\sum_{d\mid n}\chi(d)\right)q^n,\qquad
 E_{\bar\chi}(q)=\frac{3-i}{10}
     +\sum_{n\ge1}\left(\sum_{d\mid n}\bar\chi(d)\right)q^n
\]
have characters \(\chi,\bar\chi\) on \(\Gamma_0(5)\).
Their constant terms are \(L(0,\chi)/2=(3+i)/10\) and its conjugate.
The form in \eqref{E:eq:zagier-data} is
\[
 \Theta=\frac{3-i}{2}E_\chi+\frac{3+i}{2}E_{\bar\chi}.
\]
Since the lower-right entry of \(\delta\) is 2, its expansion at
\(c_*\) is
\[
 \Theta|_1\delta
 =i\frac{3-i}{2}E_\chi-i\frac{3+i}{2}E_{\bar\chi}
 =q+O(q^2).
\]
Indeed, the constant term is \(i/2-i/2=0\), and the coefficient
of \(q\) is 1.

The effective index of \(H\) is 12: modulo 5, the group
\({\rm SL}_2(\mathbb F_5)\) has order 120 and the image of \(H\)
has order 10.  There are no elliptic points, since the trace of
an element of \(H\) is congruent to \(2\) or \(-2\) modulo 5,
whereas a nontrivial elliptic element would have trace
\(0,1\), or \(-1\).  The weight-one valence formula therefore
gives total zero order \(12/12=1\).  The zero just found at
\(c_*\) accounts for this entire divisor.

Write \(\phi=\eta(5\tau)^5/\eta(\tau)\).
This is a weight-two form on \(\Gamma_0(5)\) of character
\((5/\cdot)\).  It has order 1 at \(\infty\), hence also at
\(c_*\), since \(\phi|_2\delta=-\phi\).  At 0 its order is zero:
with \(S=\left(\begin{smallmatrix}0&-1\\1&0\end{smallmatrix}\right)\),
\[
 \phi|_2S=-5^{-5/2}\frac{\eta(\tau/5)^5}{\eta(\tau)},
\]
whose leading term is a nonzero constant.  Its order is also
zero at \(\delta0\), by modularity on \(\Gamma_0(5)\).
It has no zeros in the upper half plane.  The identity
\(t=\phi/\Theta^2\) now gives the asserted divisor.  A
nonconstant meromorphic function with exactly one simple pole
has degree one as a map to \(\mathbb P^1\), so it is a Hauptmodul.
\end{proof}

\subsection{The Hecke correspondence with its sign}

Fix \(p\equiv\varepsilon_p\pmod5\), where
\(\varepsilon_p\in\{\pm1\}\).  Let \(\mathcal M_p\) consist of
the integral matrices \(\beta=\left(\begin{smallmatrix}a&b\\c&d\end{smallmatrix}\right)\)
satisfying
\[
 \begin{aligned}
 ad-bc&=p,\qquad c\equiv0\pmod5,\\
 (a,d)&\equiv(p,1)\ \text{or}\ (-p,-1)\pmod5.
 \end{aligned}
\]
For \(\beta\in\mathcal M_p\), let \(\epsilon(\beta)\in\{\pm1\}\)
be its lower-right entry modulo 5.  Left and right multiplication
by \(H\) preserve \(\mathcal M_p\), and
\[
 \epsilon(\gamma\beta)=\epsilon(\gamma)\epsilon(\beta)
                    =\epsilon(\beta\gamma)
 \qquad(\gamma\in H).
\]

\begin{lemma}[Cosets and the distinguished cusp]\label{E:lem:zagier-hecke}
The left \(H\)-cosets in \(\mathcal M_p\) have the representatives
\[
 \beta_\infty=\begin{pmatrix}p&0\\0&1\end{pmatrix},\qquad
 \beta_b=\begin{pmatrix}1&b\\0&p\end{pmatrix}\quad(0\le b<p).
\]
Their signs are
\(\epsilon(\beta_\infty)=1\) and
\(\epsilon(\beta_b)=\varepsilon_p\).
Also \({\rm adj}(\beta)\in\mathcal M_p\) for every
\(\beta\in\mathcal M_p\), and every such \(\beta\) preserves
the cusp class \(c_*\).
\end{lemma}

\begin{proof}
Apply integral Hermite normal form to write
\(\beta=\gamma\beta_j\), with \(\gamma\in{\rm SL}_2(\mathbb Z)\)
and \(\beta_j\) one of the displayed matrices.
If \(j=\infty\), the diagonal entries of
\(\gamma=\beta\beta_\infty^{-1}\) reduce to the same \(\pm1\)
modulo 5, and its lower-left entry is divisible by 5.
If \(j=b\), then
\[
 \gamma=
 \begin{pmatrix}
 a&(b'-ab)/p\\ c&(d-cb)/p
 \end{pmatrix}
 \quad\text{when }\beta=\begin{pmatrix}a&b'\\c&d\end{pmatrix}.
\]
Its diagonal entries reduce to the same \(\pm p=\pm1\)
modulo 5, since \(p^2\equiv1\pmod5\), and its lower-left
entry is again divisible by 5.  Thus \(\gamma\in H\) in both
cases.  Distinct representatives are already distinct modulo
\({\rm SL}_2(\mathbb Z)\).  The sign assertions follow directly
from their lower-right entries.  The congruence
\(p^2\equiv1\pmod5\) likewise shows that taking adjugates
preserves \(\mathcal M_p\).

A primitive cusp column whose lower component is divisible
by 5 has \(H\)-class determined by its upper component modulo
5 up to sign.  The class \(c_*\) has upper component \(\pm2\).
Applying a matrix in \(\mathcal M_p\) multiplies this residue
by \(\pm p=\pm1\).  The resulting column has common divisor
either 1 or \(p\), because its image under the adjugate is
\(p\) times the original primitive column.  Reduction to a
primitive column therefore again changes its upper residue
only by a sign.  This proves cusp stability; all representatives
of that cusp have width 1.
\end{proof}

\subsection{Polynomiality and the exact trace}

For an integer polynomial \(P\) of degree \(d\), put
\[
 R_\beta(\tau)
   =p^{1/2}\epsilon(\beta)
       \frac{\Theta(\tau)}{(\Theta|_1\beta)(\tau)}
       P(t(\beta\tau)),\qquad
 \mathcal T_p(P)=\sum_{H\backslash\mathcal M_p}R_\beta.
\]

\begin{lemma}[Zagier D trace polynomial]\label{E:lem:zagier-trace}
There is a polynomial \(\Pi_P\in\mathbb Z[t]\) of degree
at most \(pd+p-1\) such that
\begin{equation}\label{E:eq:zagier-trace}
 \Pi_P(t(q))
 =
 \frac{\Theta(q)P(t(q^p))}{\Theta(q^p)}
 +\varepsilon_p p^2 U_p
      \left(\frac{\Theta(q^p)P(t(q))}{\Theta(q)}\right).
\end{equation}
Here \(U_p(\sum_n b_nq^n)=\sum_n b_{pn}q^n\).
\end{lemma}

\begin{proof}
The character signs show that \(R_{\gamma\beta}=R_\beta\)
and \(R_\beta(\gamma\tau)=R_{\beta\gamma}(\tau)\) for
\(\gamma\in H\).  Thus \(\mathcal T_p(P)\) is \(H\)-invariant.

No summand has a pole in the upper half plane, by
Lemma~\ref{E:lem:zagier-divisors}.  If a cusp \(c'\) lies outside
the class \(c_*\), then so does \(\beta c'\): otherwise applying
the adjugate and using Lemma~\ref{E:lem:zagier-hecke} would put
\(c'\) in that class.  At these cusps the forms in the numerator
and denominator have nonzero leading constants, and \(t\)
stays finite.  There are therefore no poles there.

At \(c_*\), write
\(\beta\delta=\delta_\beta\beta'\) in integral Hermite normal
form, where
\(\beta'=\left(\begin{smallmatrix}a'&b'\\0&d'\end{smallmatrix}\right)\)
and \(a'd'=p\).  As both \(c_*\) and \(\beta c_*\) have width
1 and are \(H\)-equivalent, the orders of
\(\Theta|_1\beta\) and \(t\circ\beta\) are respectively
\(a'/d'\) and \(-a'/d'\).
Consequently
\[
 {\rm ord}_{c_*}R_\beta
   =1-\frac{a'}{d'}(d+1)\ge -pd-p+1.
\]
It follows that the trace has its only possible pole at
\(c_*\), of order at most \(pd+p-1\).
Since \(t\) is a Hauptmodul with a simple pole there,
\(\mathcal T_p(P)=\Pi_P(t)\) with this degree bound.

The infinity branch equals the first term on the right of
\eqref{E:eq:zagier-trace}.  Each remaining branch is
\[
 \varepsilon_p p\,
 \frac{\Theta(q)P(t(\zeta_p^bq^{1/p}))}
      {\Theta(\zeta_p^bq^{1/p})}.
\]
Summing over \(b\) gives the second term in
\eqref{E:eq:zagier-trace}.  Both terms have integral
\(q\)-expansions, since \(\Theta(q)\) is an integral unit.
As \(t=q+O(q^2)\) integrally, coefficient comparison then gives
\(\Pi_P\in\mathbb Z[t]\).
\end{proof}

\begin{proof}[Proof of Theorem~\ref{E:thm:zagier-split}]
Fix \(m\ge1\), put \(K_j=mp^j\) for \(j\ge0\), and set
\(F_{K_{-1}}=0\).  Take \(P=F_{K_{s-1}}\) in
Lemma~\ref{E:lem:zagier-trace}.  The degree of its trace polynomial
is at most
\[
 p(K_{s-1}-1)+p-1=K_s-1.
\]
For \(s\ge2\), the previous induction step, after
cross multiplication and division by the unit \(\Theta(q)\), reads
\[
 \frac{\Theta(q^p)F_{K_{s-1}}(t(q))}{\Theta(q)}
 \equiv F_{K_{s-2}}(t(q^p))\pmod{p^{2(s-1)}}.
\]
Applying \(U_p\) and multiplying by \(p^2\) gives
\[
 p^2U_p\left(
       \frac{\Theta(q^p)F_{K_{s-1}}(t(q))}{\Theta(q)}
       \right)
 \equiv p^2 F_{K_{s-2}}(t(q))\pmod{p^{2s}}.
\]
For \(s=1\), this last assertion is automatic under the
convention \(F_{K_{-1}}=0\).

Thus the exact trace identity makes
\[
 Q(t)=F(t)\frac{F_{K_{s-1}}(\sigma_p(t))}
                       {F(\sigma_p(t))}
\]
congruent modulo \(p^{2s}\) to an integer polynomial
of degree less than \(K_s\).
Since \(\sigma_p(t)=t^p+O(t^{p+1})\),
\[
 Q(t)=F(t)+O(t^{K_s}).
\]
Coefficient comparison forces
\(Q(t)\equiv F_{K_s}(t)\pmod{p^{2s}}\).
This completes the induction for all \(m,s\).
Dividing by the unit \(F_{K_{s-1}}(\sigma_p(t))\) gives
the claimed quotient congruence.
\end{proof}

\subsection{Diamond classes and the Laurent correction}

The same divisor calculation also determines the diamond action:
\begin{equation}\label{E:eq:zagier-diamond}
 \Theta|_1\delta=t\Theta,\qquad t(\delta\tau)=-1/t(\tau).
\end{equation}
For the first identity, the quotient
\((\Theta|_1\delta)/\Theta\) has divisor
\([\infty]-[c_*]\), so it is a constant multiple of \(t\);
its first \(q\)-coefficient is 1.  The second identity follows
from \(\phi|_2\delta=-\phi\) and \(t=\phi/\Theta^2\).

When \(p\equiv\pm2\pmod5\), the diamond action in the Hecke
correspondence is therefore non-scalar.  In weight \(-1\),
it sends
\[
 \frac{P(t)}{\Theta}
 \ \longmapsto\
 \left(\frac{P(t)}{\Theta}\right)\bigg|_{-1}\delta
 =\frac{P(-1/t)}{t\Theta}.
\]
This has a pole at the expansion cusp \(t=0\).
The corresponding trace can have a negative Laurent part,
so Lemma~\ref{E:lem:zagier-trace} does not extend by the same
argument. Controlling that negative Laurent part and its regular remainder requires the paired descent proved in Section~\ref{T:sec:zagier-complete}; the present one-cusp argument alone does not provide that control.


\subsection{The first level for every prime different from five}

The Laurent terms just described prevent the preceding higher-level
induction, but do not prevent the initial step.

\begin{theorem}[Zagier D at the first level]\label{E:thm:zagier-first-level}
For every prime \(p\ne5\) and every integer \(m\ge1\),
\[
 \frac{F(t)}{F(\sigma_p(t))}
 \equiv
 \frac{F_{mp}(t)}{F_m(\sigma_p(t))}
 \pmod{p^2\mathbb Z_p[[t]]}.
\]
\end{theorem}

\begin{proof}
The cases \(p\equiv\pm1\pmod5\) already follow from
Theorem~\ref{E:thm:zagier-split}.  Assume \(p\equiv2\) or
\(3\pmod5\).  Retain the matrix-defined set
\(\mathcal M_p\) and the sign \(\epsilon(\beta)=d\bmod5\);
these definitions still make sense, since its lower-right
entries are \(\pm1\) modulo 5.  Put
\[
 \eta_p=\begin{cases}\delta^{-1},&p\equiv2\pmod5,\\
                     \delta,&p\equiv3\pmod5,\end{cases}
 \qquad
 c_p=\begin{cases}-1,&p\equiv2\pmod5,\\
                  1,&p\equiv3\pmod5.\end{cases}
\]
The coset argument and the trace below require only \(p\nmid5\);
they do not require \(p\) to be odd.  In particular, \(p=2,3\)
are included.
Integral Hermite normal form now gives the left \(H\)-cosets
\[
 \beta_\infty,\qquad \eta_p\beta_b\quad(0\le b<p).
\]
Indeed, in a finite Hermite branch the intervening integral
matrix has diagonal residues \(\pm(p,p^{-1})\) modulo 5;
after multiplication by \(\eta_p^{-1}\) it lies in \(H\).
The signs of all displayed representatives are 1.
By \eqref{E:eq:zagier-diamond},
\[
 \left(\frac{P(t)}{\Theta}\right)\bigg|_{-1}\eta_p
 =c_p\frac{P(-1/t)}{t\Theta}.
\]
The same weight-zero trace construction therefore gives the
exact identity
\begin{equation}\label{E:eq:zagier-laurent-trace}
 \mathcal T_p(P)
 = Q_P+c_pp^2\Theta(q)\,
       U_p\left(\frac{P(-1/t(q))}{t(q)\Theta(q)}\right),
 \qquad
 Q_P=\frac{\Theta(q)P(t(q^p))}{\Theta(q^p)}.
\end{equation}
Here \(U_p\) acts on Laurent series as well as power series.

We claim that, for \(d=\deg P\),
\begin{equation}\label{E:eq:zagier-laurent-support}
 \mathcal T_p(P)\in\mathbb Z[t,t^{-1}],
 \qquad
 {\rm supp}_t\mathcal T_p(P)
 \subseteq
 \left[-\left\lfloor\frac{d+1}{p}\right\rfloor,\,
             p(d+1)-1\right].
\end{equation}
The absence of upper-half-plane poles follows as before.
A cusp with denominator prime to 5 cannot map to either
\(\infty\) or \(c_*\) under \(\mathcal M_p\), so the trace
has no poles at the two other cusps.

To analyze the remaining cusps, apply a determinant-\(p\)
matrix to a primitive cusp column.  Its new common divisor
is 1 or \(p\).  Among the \(p+1\) branches exactly one has
common divisor \(p\), and hence ramification \(p\); the others
have ramification \(1/p\).  This follows from Hermite normal
form after multiplication by a matrix representing the cusp.
At \(c_*\), whose upper residue is \(\pm2\) modulo 5,
a branch with common divisor 1 changes that residue to
\(\pm2p=\pm1\), and therefore maps to \(\infty\).
The unique branch with common divisor \(p\) maps back to
\(c_*\).  Its pole order in the trace is at most
\(p(d+1)-1\), since \(\Theta\) has order 1 at \(c_*\).
All other branches are holomorphic there.

At \(\infty\), the unique ramification-\(p\) branch maps
back to \(\infty\); the other \(p\) branches map to \(c_*\).
Each of these can have pole order at most \((d+1)/p\).
Since the trace is \(H\)-invariant and the cusp width is 1,
its pole order is an integer, at most
\(\lfloor(d+1)/p\rfloor\).
Thus the trace is a Laurent polynomial with the claimed
support.  Its Laurent \(q\)-expansion is integral by
\eqref{E:eq:zagier-laurent-trace}, because
\(t/q\) and \(\Theta\) are integral units.
Substitution of the integral inverse series \(q(t)\)
then proves integrality of its Laurent coefficients,
establishing \eqref{E:eq:zagier-laurent-support}.
Explicitly, \(q(t)=tV(t)\) with
\(V(t)\in1+t\mathbb Z[[t]]\); its inverse powers remain integral,
so substitution sends \(\mathbb Z((q))\) into \(\mathbb Z((t))\).

Now take \(P=F_m\), so \(d=m-1\).
The correction in \eqref{E:eq:zagier-laurent-trace} is
divisible by \(p^2\) in the integral Laurent-series ring.
Hence \(Q_P\) is congruent modulo \(p^2\) to a Laurent
polynomial whose positive degree is less than \(mp\).
Since \(Q_P\) itself has no negative powers, all negative
coefficients of that Laurent polynomial vanish modulo \(p^2\).
It is therefore congruent to an ordinary polynomial
of degree less than \(mp\).
Finally,
\[
 Q_P(t)=F(t)\frac{F_m(\sigma_p(t))}{F(\sigma_p(t))}
       =F(t)+O(t^{mp}).
\]
Coefficient comparison gives \(Q_P\equiv F_{mp}\pmod{p^2}\).
Dividing by \(F_m(\sigma_p(t))\) completes the proof.
\end{proof}

This Laurent-support argument supplies the first level in the diamond classes. The paired induction of Section~\ref{T:sec:zagier-complete} supplies every level $s\ge2$ as well.


\section{Exact coefficient obstructions to larger exponents}\label{E:source-main-4}
\label{E:sec:sharpness}

The results of this section give upper bounds for the modulus that can
hold in the full excellent-lift congruence. They complement the full lower bound modulo $p^{2s}$ proved in Theorem~\ref{P:thm:main-all}.  Each obstruction is an exact calculation
of one coefficient, valid for infinitely many levels by proof rather
than by extrapolation from computation.

\subsection{Normalization and the first two coefficients}

Write
\[
 F(u)=\sum_{k\geq0}a_ku^k\in1+u\mathbb Z_p[[u]],\qquad
 q(u)=u\exp(g(u)),\qquad g(u)=\alpha u+O(u^2).
\]
Let $T$ be the compositional inverse of $q$, and put
$\sigma(u)=T(q(u)^p)$.  For $K=p^s$, define the cleared defect
\[
 E_s(u)=F(u)F_{K/p}(\sigma(u))-F(\sigma(u))F_K(u),
 \qquad F_M(u)=\sum_{0\leq k<M}a_ku^k.
\]
Its divisibility by $p^b$ is equivalent to the quotient congruence
modulo $p^b$, because all denominator series have constant term one.

\begin{lemma}[The first two possible defect coefficients]
\label{E:lem:sh-firstjet}
For $p\geq2$ and $s\geq1$,
\begin{align}
 [u^K]E_s&=a_K-a_{K/p},\label{E:eq:sh-first}\\
 [u^{K+1}]E_s&=a_{K+1}-(a_1+K\alpha)a_{K/p}.
 \label{E:eq:sh-second}
\end{align}
\end{lemma}
\begin{proof}
Since $T(z)=z+O(z^2)$,
$\sigma(u)=u^p(1+p\alpha u+O(u^2))$.  Rewrite the defect as
\[
 E_s=(F-F_K)F_{K/p}(\sigma)-F_K(F-F_{K/p})(\sigma).
\]
The first summand has coefficients $a_K,a_{K+1}$ at degrees $K,K+1$,
since $F_{K/p}(\sigma)=1+O(u^p)$.  Up to degree $K+1$, the second
summand is
\[
 a_{K/p}F_K(u)\sigma(u)^{K/p}
 =a_{K/p}u^K\bigl(1+(a_1+K\alpha)u+O(u^2)\bigr).
\]
This proves both identities.
\end{proof}

We record the normalization explicitly, using the expressions for
$F_0(t)$ and $G_0(t)$ in \cite[\S7, Proposition~9.1]{BV3}.
Here $t$ is the parameter of the original Laurent-polynomial family;
$u$ is the normalized parameter used in this section.  Let
$H_k=\sum_{j=1}^k1/j$ and $H_0=0$.

For the hypercubic family in dimension $d$,
\[
 u=t^2,\qquad a_k=\binom{2k}{k}^d,\qquad
 g(u)=\frac{2d\sum_{k\geq1}a_k(H_{2k}-H_k)u^k}{F(u)}.
\]
Thus $a_1=2^d$ and $\alpha=d2^d$.
For the hyperoctahedral family in dimension $d$,
\[
 u=t^2,\qquad
 a_k=\sum_{k_1+\cdots+k_d=k}\frac{(2k)!}{(k_1!\cdots k_d!)^2},
\]
and
\[
 g(u)=\frac{2}{F(u)}\sum_{k\geq1}u^k
 \sum_{k_1+\cdots+k_d=k}
 \frac{(2k)!}{(k_1!\cdots k_d!)^2}(H_{2k}-H_{k_1}).
\]
At $k=1$ the inner sum is $1+3(d-1)$; hence
$a_1=2d$ and $\alpha=6d-4$.
For the simplicial family in dimension $d$, put $r=d+1$.  Then
\[
 u=t^r,\qquad a_k=\frac{(rk)!}{(k!)^r},\qquad
 g(u)=\frac{r\sum_{k\geq1}a_k(H_{rk}-H_k)u^k}{F(u)}.
\]
In particular, $a_1=r!$ and $\alpha=r!r(H_r-1)$.

In each case, with $h=2$ or $h=r$, the original coordinate
$q_0(t)=t\exp(G_0(t)/F_0(t))$ satisfies
$q_0(t)^h=q(t^h)$.  The vertex coefficient is one, so
$q_0\mapsto q_0^p$ implies $(t^{\sigma_0})^h=\sigma(t^h)$.
Furthermore, $F_{0,hM}(t)=F_M(t^h)$.  The normalized defects above
therefore give exactly the original Beukers--Vlasenko defects with
$m=2$ or $m=d+1$ after the injective substitution $u=t^h$.
The admissible primes do not divide $h$.
No identification with an unverified Picard--Fuchs logarithmic
solution is needed: these formulas use the coordinate in
\cite[Definition~7.14]{BV3} itself.

\subsection{A binomial estimate}

We use Jacobsthal's congruence in the form stated in
\cite[Lemma~2.1, equation~(7)]{St23}: for $p\geq5$ and $r,t\geq1$,
\begin{equation}\label{E:eq:sh-jacobsthal}
 \binom{p^rn}{p^tk}\bigg/\binom{p^{r-1}n}{p^{t-1}k}
 \equiv1\pmod{p^{r+t+\min(r,t)}}.
\end{equation}
Endpoint ratios with equal top and bottom are identically one.
In particular, for $K=p^s$,
\begin{equation}\label{E:eq:sh-central}
 \binom{2K}{K}\equiv\binom{2K/p}{K/p}\pmod{p^{3s}}.
\end{equation}

\subsection{Hyperoctahedral dimension three}

Let
\[
 b_N=\sum_{k=0}^N\binom Nk^2\binom{2k}{k},\qquad
 a_N=\binom{2N}{N}b_N.
\]
These are the normalized three-dimensional hyperoctahedral
coefficients.  For $p\geq5$, define
\[
 S_p=\sum_{k=1}^{(p-1)/2}\frac{\binom{2k}{k}}{k^2}\in\mathbb F_p,
 \qquad \chi_p=\left(\frac{-3}{p}\right).
\]

\begin{theorem}[All-level obstruction in dimension three]
\label{E:thm:sh-octa3}
For every $p\geq5$ and $s\geq1$,
\begin{equation}\label{E:eq:sh-octa3}
 a_{p^s}-a_{p^{s-1}}
 \equiv2\chi_p^{s-1}S_p\,p^{2s}\pmod{p^{2s+1}}.
\end{equation}
In particular,
\[
 [u^{5^s}]E_s\equiv2(-1)^{s-1}5^{2s}\pmod{5^{2s+1}}
 \qquad(s\geq1),
\]
so the displayed coefficient has valuation exactly $2s$ at every level.
\end{theorem}
\begin{proof}
Put $N=p^s$, $M=p^{s-1}$, and
$T_N(k)=\binom Nk^2\binom{2k}{k}$.  Compare the terms with $k=pj$
to $T_M(j)$.  For $0<j<M$, write $r=v_p(j)\leq s-2$.
Lucas's theorem, or the identity
$\binom Mj=(M/j)\binom{M-1}{j-1}$, gives
$v_p\binom Mj=s-1-r$.
By \eqref{E:eq:sh-jacobsthal},
\[
 \binom N{pj}/\binom Mj\equiv1\pmod{p^{s+2r+2}},\qquad
 \binom{2pj}{pj}/\binom{2j}{j}\equiv1\pmod{p^{3r+3}}.
\]
Multiplication by the squared binomial coefficient adds valuation
at least $2s-2r-2$.  The resulting error bounds are $3s$ and
$2s+r+1$, both at least $2s+1$.  Thus
$T_N(pj)\equiv T_M(j)\pmod{p^{2s+1}}$.
The endpoint $j=0$ is exact, and $j=M$ follows from
\eqref{E:eq:sh-central}.  Consequently,
\[
 b_N-b_M\equiv
 \sum_{\substack{0<k<N\\p\nmid k}}T_N(k)\pmod{p^{2s+1}}.
\]
For $p\nmid k$, Lucas gives
$\binom{N-1}{k-1}\equiv(-1)^{k-1}\pmod p$; hence
\[
 \frac{b_N-b_M}{p^{2s}}
 \equiv\sum_{\substack{0<k<p^s\\p\nmid k}}
 \frac{\binom{2k}{k}}{k^2}\pmod p.
\]
If a base-$p$ digit of $k$ exceeds $(p-1)/2$, doubling creates a
carry and $\binom{2k}{k}\equiv0\pmod p$.  Otherwise Lucas factors
the central binomial coefficient over all digits.  The denominator
depends only on the nonzero units digit, so the last sum equals
\[
 S_p\left(\sum_{j=0}^{(p-1)/2}\binom{2j}{j}\right)^{s-1}.
\]
Writing $h=(p-1)/2$, we have
$\binom{2j}{j}\equiv(-4)^j\binom hj\pmod p$.
The parenthesized sum is therefore $(-3)^h=\chi_p$ in
$\mathbb F_p$.
The central-binomial factors at $N$ and $M$ differ by a multiple
of $p^{3s}$, and each is $2$ modulo $p$.  This proves
\eqref{E:eq:sh-octa3}.
Finally $S_5=2+6/4=1$ in $\mathbb F_5$ and $\chi_5=-1$;
apply \eqref{E:eq:sh-first}.
\end{proof}

This rules out a congruence modulo $5^{2s+1}$, and therefore modulo
$5^{3s}$, for the full defect at every level.  It does not by itself
show that all coefficients of the defect are divisible by $5^{2s}$; Theorem~\ref{P:thm:main-all} supplies that lower bound. Thus the uniform exponent is exactly $2s$ at $p=5$.

\subsection{All higher-dimensional hypercubic families}

\begin{theorem}[Hypercubic obstruction]
\label{E:thm:sh-cube}
Let $d\geq4$, let $p>d$ be prime, and let $s\geq1$.
For $K=p^s$,
\[
 [u^{K+1}]E_s
 \equiv2^{2d-1}d(d-3)p^{2s}\pmod{p^{2s+1}}.
\]
In particular this coefficient has valuation exactly $2s$ for every
such $d,p,s$.
\end{theorem}
\begin{proof}
Here $a_k=\binom{2k}{k}^d$, $a_1=2^d$, and $\alpha=d2^d$.
Equation \eqref{E:eq:sh-central} gives
$a_K\equiv a_{K/p}\pmod{p^{3s}}$.  The exact adjacent ratio is
\[
 \frac{a_{K+1}}{a_K}
 =2^d\left(\frac{1+2K}{1+K}\right)^d
 =2^d\left(1+dK+\frac{d(d-3)}2K^2+O(K^3)\right),
\]
with $p$-integral coefficients.  Since $3s\geq2s+1$, substituting
into \eqref{E:eq:sh-second} yields
\[
 [u^{K+1}]E_s
 \equiv2^{d-1}d(d-3)K^2a_{K/p}\pmod{p^{2s+1}}.
\]
Lucas gives $a_{K/p}\equiv2^d\pmod p$.
As $p>d\geq4$, the resulting constant
$2^{2d-1}d(d-3)$ is nonzero modulo $p$.
\end{proof}

For the full hypercubic symmetry group, $p>d$ is precisely the
admissible-prime restriction in these dimensions.  Thus the exponent
$2s$ cannot be increased even by one throughout the conjectured range.
The theorem gives the upper bound, and Theorem~\ref{P:thm:main-all} gives the lower bound. Hence $2s$ is the exact full-series exponent throughout this hypercubic range.

\subsection{All higher-dimensional simplicial families}

Put $r=d+1$, and let $H_r^{(2)}=\sum_{j=1}^r1/j^2$.  Define
\[
 A_r=r(H_r-1),\qquad
 B_r=\frac{A_r^2-r^2H_r^{(2)}+r}{2},\qquad
 C_r=(r!)^2B_r.
\]

\begin{theorem}[Simplicial obstruction and exceptional primes]
\label{E:thm:sh-simplex}
Let $r\geq5$ and let $p>r$ be prime.  For every $s\geq1$,
\begin{equation}\label{E:eq:sh-simplex}
 [u^{p^s+1}]E_s\equiv C_rp^{2s}\pmod{p^{2s+1}}.
\end{equation}
The number $C_r$ is a positive integer.  Thus this coefficient has
valuation exactly $2s$ whenever $p\nmid C_r$.
More generally, setting $c=v_p(C_r)$, one has
\begin{equation}\label{E:eq:sh-eventual}
 v_p([u^{p^s+1}]E_s)=2s+c\qquad\text{for every }s>c.
\end{equation}
Consequently every admissible prime fails to support a
$p^{3s}$ strengthening at all sufficiently large levels.
\end{theorem}
\begin{proof}
The coefficients factor as
\[
 a_K=\frac{(rK)!}{(K!)^r}=\prod_{j=2}^r\binom{jK}{K}.
\]
Apply \eqref{E:eq:sh-jacobsthal} factor by factor to obtain
$a_K\equiv a_{K/p}\pmod{p^{3s}}$.
Lucas gives $a_{K/p}\equiv r!\pmod p$, a unit since $p>r$.
The adjacent ratio expands as
\[
 R_r(K):=\frac{a_{K+1}}{a_K}
 =\frac{\prod_{j=1}^r(rK+j)}{(1+K)^r}
 =r!\bigl(1+A_rK+B_rK^2+O(K^3)\bigr).
\]
Indeed, the first two coefficients of its logarithm are
$A_r$ and $-(r^2H_r^{(2)}-r)/2$.
Since $a_1=r!$ and $\alpha=r!A_r$,
\eqref{E:eq:sh-second} gives the stronger identity
\begin{equation}\label{E:eq:sh-simplex-strong}
 [u^{K+1}]E_s=r!B_r a_{K/p}K^2+O(p^{3s}).
\end{equation}
Reducing $a_{K/p}$ modulo $p$ proves \eqref{E:eq:sh-simplex}.

The series $R_r(X)$ lies in $\mathbb Z[[X]]$: its numerator is an
integral polynomial and $(1+X)^{-r}$ has integral coefficients.
Thus $r!B_r=[X^2]R_r(X)$ is integral, and so is $C_r$.
For positivity, write
\[
 D_r:=\frac{2B_r}{r^2}
 =2\sum_{2\leq i<j\leq r}\frac1{ij}-1+\frac1r.
\]
One has $D_4=0$ and
\[
 D_{r+1}-D_r=
 \frac{2(H_r-1)-1/r}{r+1}>0\qquad(r\geq4).
\]
Hence $C_r>0$ for all $r\geq5$.
Finally, in \eqref{E:eq:sh-simplex-strong}, the displayed leading
term has valuation $2s+c$, because $r!$ and $a_{K/p}$ are units.
For $s>c$ this is smaller than $3s$, proving
\eqref{E:eq:sh-eventual}.
\end{proof}

For example,
\[
 C_4=0,\qquad C_5=69000,\qquad
 C_6=7257600,\qquad C_7=727050240.
\]
Thus dimension four at $p=7$ has
$[u^{7^s+1}]E_s\equiv7^{2s}\pmod{7^{2s+1}}$ for every $s$.
In dimension five, the exceptional prime $p=7$ has
$v_7(C_6)=1$, so the same coefficient has valuation $2s+1$
for every $s\geq2$.
The vanishing $C_4=0$ is a cancellation in this particular necessary
obstruction; it is not an independent proof of the stronger
three-dimensional congruence.

\subsection{Hyperoctahedral dimension four}

Let
\[
 D_k=\sum_{j=0}^k\binom kj^2\binom{2j}{j}
                         \binom{2k-2j}{k-j}
\]
be the Domb numbers.  Grouping the four composition indices into two
pairs and applying Vandermonde shows that the four-dimensional
hyperoctahedral coefficient is $a_k=\binom{2k}{k}D_k$.
The established Domb congruence
\cite[Theorem~1.1, with $A=2$, $B=C=1$]{OS12} and
\eqref{E:eq:sh-central} give
\[
 a_{p^s}\equiv a_{p^{s-1}}\pmod{p^{3s}}
 \qquad(p\geq5).
\]

\begin{theorem}[Four-dimensional hyperoctahedral obstruction]
\label{E:thm:sh-octa4}
For every prime $p\geq5$ and every $s\geq1$,
\[
 [u^{p^s+1}]E_s\equiv32p^{2s}\pmod{p^{2s+1}}.
\]
This coefficient has valuation exactly $2s$ at every admissible prime
and level.
\end{theorem}
\begin{proof}
The Domb recurrence, equivalently its differential equation in
\cite[Appendix~C, \S9.7]{Beu26} after removing the indicated sign twist,
is
\[
 k^3D_k=2(2k-1)(5k^2-5k+2)D_{k-1}-64(k-1)^3D_{k-2}.
\]
Multiplication by the central binomial coefficient gives
\[
 k^4a_k=4(2k-1)^2(5k^2-5k+2)a_{k-1}
       -256(k-1)^2(2k-1)(2k-3)a_{k-2}.
\]
Set $k=K+1$ and expand modulo $K^3$.  Since $K+1$ is a $p$-adic
unit,
\[
 a_{K+1}\equiv(8+20K+4K^2)a_K+256K^2a_{K-1}
       \pmod{p^{3s}}.
\]
Here $a_1=8$ and $\alpha=20$.
Use \eqref{E:eq:sh-second} and the preceding coefficient
supercongruence to obtain
\[
 [u^{K+1}]E_s
 \equiv K^2(4a_{K/p}+256a_{K-1})\pmod{p^{2s+1}}.
\]
Lucas gives $a_{K/p}\equiv8\pmod p$.
The factor $\binom{2K-2}{K-1}$ in $a_{K-1}$ is zero modulo $p$,
since doubling $K-1=p^s-1$ creates a carry.
The stated congruence follows.
\end{proof}

\paragraph{Interpretation.}
Theorems~\ref{E:thm:sh-octa3}, \ref{E:thm:sh-cube},
\ref{E:thm:sh-simplex}, and \ref{E:thm:sh-octa4} prove exact upper
obstructions for the possible full-series modulus. The lower bound modulo $p^{2s}$ is Theorem~\ref{G:thm:global}. Their combination proves sharpness in the parameter ranges explicitly stated; it does not assert exactness at every exceptional simplicial prime or every higher octahedral dimension.


\section{Relation to the cited literature}\label{E:source-main-5}

\emph{Dwork crystals III} supplies the exact formulation of Conjecture 7.5,
the canonical coordinates, and the excellent Cartier relation
\cite[Section 7 and Proposition 9.1]{BV3}.  Its higher formal-derivative
statement already gives strong congruences for appropriate Laurent
coefficients, but does not by itself identify the finite truncations in
\eqref{E:eq:quotient}.  Section~\ref{E:sec:elliptic} supplies this identification
for the two low-dimensional simplicial families by bounding the
entire rational trace.

Beukers's modular-form paper supplies the parametrizations and the Hecke
permutation of the modular-form ratios \cite[Lemmas 3.6--3.7, Appendices B--C]{Beu26}.
The additional ingredient here is a bound uniform in the monomial degree
$k$ and in the modulus $p^u$.  The crucial estimate is
\[
 \deg(R_k\bmod p^u)
 \leq pk+\lfloor p/e\rfloor+\lfloor(p+1)/e\rfloor(u-1).
\]
It permits induction through every level, even though the trace has actual
complex poles and infinitely many nonzero Taylor coefficients.

Vlasenko's exposition explains the cohomological framework and the distinction between period truncations and Cartier coefficient congruences \cite{Vl24}. Fu \cite{Fu26} gives a cohomological criterion involving a Frobenius-stable rank-two factor containing the holomorphic class, and studies its relation to excellent-lift fixed points. That criterion concerns additional hypotheses at selected parameters. None of those hypotheses is imported into the formal-series theorems in this volume.

Straub's paper proves coefficient supercongruences and Gessel--Lucas
relations for sporadic sequences.  We use its explicit Jacobsthal
lemma \cite[Lemma 2.1]{St23} to remove the boundary coefficient in
the $108t$ argument and to establish the sharpness formulas.
Alinquant--Osburn finish the remaining original sporadic two-term
coefficient congruence \cite{AO23}.  Such coefficient results are useful
inputs, but do not alone imply \eqref{E:eq:quotient}.  Only the optional
dimension-four octahedral calculation needs the established Domb-number
result \cite{OS12} as an additional input.

\section{Exact verification and proof dependencies}\label{E:source-main-6}\label{E:sec:verification}

The complete verification programs for this part are embedded in Appendix~\ref{app:all-programs}, together with the other programs of the collected work. All arithmetic in these particular programs is integer or rational arithmetic, with explicit prime-power reduction. Their declared finite checks supplement the proofs; no finite initial-segment check is used as an infinite-series argument.

The elliptic checks use the Lambert-series formula
\[
 \theta_H(q)=1+6\sum_{n\geq1}\chi_{-3}(n)\frac{q^n}{1-q^n}
\]
for the cubic theta function.  This avoids a possible truncation error
from summing the lattice form $x^2+xy+y^2$ over a square of insufficient
radius.  The theoretical proof does not depend on a lattice cutoff.
The tests also inspect coefficients beyond the truncation degree, so
they are not vacuous initial-segment checks.

For Zagier D, the diamond classes satisfy
\[
 t(\delta\tau)=-1/t(\tau),\qquad \Theta|_1\delta=t\Theta,
 \qquad \delta=\begin{pmatrix}-2&-1\\5&2\end{pmatrix}.
\]
The corresponding trace has a negative Laurent part at $t=0$. Its support is finite and explicitly bounded, giving the independent first-level proof above. The one-cusp polynomial induction does not control its higher-level negative part; the simultaneous Laurent induction of Section~\ref{B:sec:D} and the paired canonical-basis descent of Section~\ref{T:sec:zagier-complete} provide that control for all $p\ne5$ and all $m,s$.


For the three-dimensional hyperoctahedral family and every higher-dimensional standard family, the full lower bound $E_{p,s,m}\in p^{2s}\Zp[[t]]$ is Theorem~\ref{P:thm:main-all}, within its displayed prime conditions. The sharpness theorems above explain why a uniform $3s$ exponent is unavailable in their stated ranges. No higher-level or higher-coefficient case within that standard theorem is left to these calculations.




% END RETAINED SOURCE E

\partstart{Two-cusp descent and complementary Cartier arguments}{part:twocusp}{All-prime Zagier D, the independent geometric first level, negative-power bases, and octahedral differential refinements.}
% BEGIN RETAINED SOURCE T
\section{Scope and exact statements}\label{T:source-main-1}

This part retains the complete two-cusp proof for Zagier D and independent auxiliary arguments for the standard families, using the modular data and excellent Cartier theorem cited below. The all-level ordinary-truncation theorem at every odd prime is
Theorem~\ref{G:thm:global}.  The earlier elliptic-zero
arguments and coefficient upper bounds are not needed in either proof.

For a period $F(t)=\sum_{k\ge0}a_kt^k$ with $a_0=1$, write
\[
 F_M(t)=\sum_{0\le k<M}a_kt^k,
 \qquad \sigma_p(t(q))=t(q^p).
\]
The excellent-lift congruence at level $s$ and multiplier $m$ is
\begin{equation}\label{T:eq:target}
 \frac{F(t)}{F(\sigma_p(t))}
 \equiv\frac{F_{mp^s}(t)}{F_{mp^{s-1}}(\sigma_p(t))}
 \pmod{p^{2s}\Zp[[t]]}.
\end{equation}
Equivalently, the complete cleared defect
\[
 E_{p,s,m}(t)=F(t)F_{mp^{s-1}}(\sigma_p(t))
             -F(\sigma_p(t))F_{mp^s}(t)
\]
belongs to $p^{2s}\Zp[[t]]$.  All congruences are coefficientwise
identities of infinite formal series.  A congruence for an isolated
coefficient is not substituted for this assertion.

\begin{center}\small
\begin{tabular}{@{}>{\raggedright\arraybackslash}p{.29\textwidth}>{\raggedright\arraybackslash}p{.61\textwidth}@{}}
\toprule
Target&Conclusion\\\midrule
Zagier D, $p\equiv\pm2\pmod5$&Full congruence \eqref{T:eq:target}, for all $m,s\ge1$, including $p=2,3$.\\[3pt]
Zagier D, all residue classes&A single proof covers every $p\ne5$ and all $m,s$.\\[3pt]
Hyperoctahedral dimension three&The full first level and an all-level weighted differential bound are proved.  The full $U_p$ descent follows from Theorem~\ref{G:thm:global}.\\[3pt]
All three families in dimensions $n\ge4$&The full first level is proved with the original prescribed truncation lengths.  An auxiliary negative-power basis has an all-level congruence; ordinary truncations at all levels satisfy Theorem~\ref{G:thm:global}.\\
\bottomrule
\end{tabular}
\end{center}

\subsection{The distinct conclusions of the auxiliary arguments}
The Zagier D theorem below covers all $s,m$ and every coefficient. The independent geometric theorem covers the entire formal series at the first level. The negative-power canonical-basis theorem is an all-level theorem for different functions, while the octahedral differential theorem is a weighted all-coefficient statement for $K=p^s$. These additional conclusions and their proofs are retained even though the ordinary positive-power truncation statement is already proved in Part~\ref{part:positive}. The volume makes no assertion of external peer review or historical priority.

\subsection{The original families and their normalization}

In $n$ variables the three Laurent polynomials are
\[
 \begin{aligned}
 g_{\rm simplex}&=x_1+\cdots+x_n+(x_1\cdots x_n)^{-1},\\
 g_{\rm cube}&=\prod_{i=1}^n(x_i+x_i^{-1}),\\
 g_{\rm octa}&=\sum_{i=1}^n(x_i+x_i^{-1}).
 \end{aligned}
\]
Their original periods have the form
\[
 \mathcal F(z)=\CT\frac1{1-zg}=F(z^h),\qquad
 h=n+1\ \text{for the simplex},\quad h=2\ \text{otherwise}.
\]
Conjecture 7.5 prescribes multiplier $h$ in this original variable
\cite[Conjecture 7.5]{BV3}.
The normalized coordinate satisfies $q(z^h)=q_0(z)^h$, so the
excellent lifts commute with this substitution and
$\mathcal F_{hM}(z)=F_M(z^h)$.
Consequently the normalized first-level statement $F/F^\sigma\equiv F_p$
is exactly the prescribed original statement with multiplier $h$.
The original auxiliary arguments in this part use the standard
full-symmetry prime ranges of \cite[Theorem 7.3]{BV3}: $p>n$
for cubes and octahedra, and $p>n+1$ for simplices. The global
formal statement at every odd prime is Theorem~\ref{G:thm:global}.

\section{The complete Zagier D congruence}\label{T:source-main-2}
\label{T:sec:zagier-complete}

Write
\[
 A_n=\sum_{j=0}^{n}\binom nj^2\binom{n+j}{j},\qquad
 F(t)=\sum_{n\ge0}A_nt^n,\qquad
 F_K(t)=\sum_{0\le n<K}A_nt^n.
\]
We use the normalized modular parametrization
\[
 t(q)=q\prod_{n\ge1}(1-q^n)^{5(\frac n5)},\qquad
 \Theta(q)=F(t(q)),\qquad
 \sigma_p(t(q))=t(q^p).
\]
The product and the Lambert series for \(\Theta\) imply
\(t\in q+q^2\mathbb Z[[q]]\) and
\(\Theta\in1+q\mathbb Z[[q]]\).

\begin{theorem}[Zagier D at every prime different from five]
\label{T:thm:zagier-nonsplit-complete}
For every prime \(p\ne5\), and all integers
\(m,s\ge1\), one has
\[
 \frac{F(t)}{F(\sigma_p(t))}
 \equiv
 \frac{F_{mp^s}(t)}{F_{mp^{s-1}}(\sigma_p(t))}
 \pmod{p^{2s}\mathbb Z_p[[t]]}.
\]
The primes \(2\) and \(3\) are included, as are every truncation
multiplier \(m\) and every level \(s\).
\end{theorem}

The proof works with the expansions at both distinguished cusps at once.
No higher-level congruence is assumed in the construction.

\subsection{The two cusps and an integral canonical basis}

Let \(H=\pm\Gamma_1(5)\), put
\[
 \delta=\begin{pmatrix}-2&-1\\5&2\end{pmatrix},\qquad
 c_*=\delta\infty=-2/5,
\]
and use determinant-normalized slash operators
\[
 (f|_k\beta)(\tau)=
 \det(\beta)^{k/2}(c\tau+d)^{-k}f(\beta\tau).
\]
For \(\gamma\in H\), write \(\epsilon(\gamma)\in\{1,-1\}\)
for its diagonal residue modulo 5.  The weight-one form satisfies
\(\Theta|_1\gamma=\epsilon(\gamma)\Theta\).

We recall the necessary divisor information, including its justification.
The four cusps are \(\infty,c_*,0,-1/2\), of widths \(1,1,5,5\).
For the quartic character \(\chi(2)=i\) modulo 5, let
\[
 E_\chi=\frac{3+i}{10}
       +\sum_{n\ge1}\left(\sum_{d\mid n}\chi(d)\right)q^n.
\]
The modular parametrization of Zagier D gives
\[
 \Theta=\frac{3-i}{2}E_\chi+\frac{3+i}{2}E_{\bar\chi},
 \qquad
 t\Theta^2=\frac{\eta(5\tau)^5}{\eta(\tau)}.
\]
These are the data in \cite[Appendix B, Section 8.8]{Beu26}.
Since the lower-right entry of \(\delta\) is 2,
\(\Theta|_1\delta=q+O(q^2)\).
The group \(H\) has effective index 12 and no elliptic points.
The weight-one valence formula therefore shows that this simple
zero at \(c_*\) is the entire zero divisor of \(\Theta\).
The eta quotient \(t\Theta^2\) has order one at each of
\(\infty,c_*\) and order zero at the other two cusps.  Consequently
\[
 {\rm div}(t)=[\infty]-[c_*].
\]
In particular \(t\) is a Hauptmodul.  The leading coefficients
and these divisors, together with
\((t\Theta^2)|_2\delta=-t\Theta^2\), give
\begin{equation}\label{T:eq:zd-diamond-data}
 \Theta|_1\delta=t\Theta,\qquad
 t(\delta\tau)=-1/t(\tau).
\end{equation}

On weight \(-1\) forms set
\[
 \mathcal{D}f=f|_{-1}\delta,\qquad
 J P(t)=\frac{P(-1/t)}{t}.
\]
Thus
\[
 \mathcal{D}\left(\frac{P(t)}{\Theta}\right)
       =\frac{JP(t)}{\Theta},
 \qquad \mathcal{D}^2=-1,
\]
because \(\delta^2=-I\).
Both \(P\mapsto JP\) and the modular expansions preserve integral
Laurent coefficients.

For every \(n\ge1\) there is a unique form
\begin{equation}\label{T:eq:zd-canonical}
 g_n\in\frac{t^{-1}\mathbb Z[t^{-1}]}{\Theta},\qquad
 g_n(q)=q^{-n}+O(1),
\end{equation}
whose only negative Fourier term is \(q^{-n}\).
Indeed the forms \(t^{-j}/\Theta\), \(j=1,\ldots,n\), have leading
terms \(q^{-j}\) with coefficient one.  Successively canceling
the other negative terms is an integral triangular procedure.
The form \(g_n\) is holomorphic away from \(\infty\):
at \(c_*\), \(t^{-j}/\Theta\) has order \(j-1\ge0\), and
at the other cusps both \(t\) and \(\Theta\) have nonzero
leading terms.  There are no upper-half-plane poles.
Conversely every weight \(-1\) form with the same character
and poles only at \(\infty\) is determined by its principal part.
The difference of two such forms with identical principal part
would be a holomorphic form of negative weight, and hence zero.

The diamond transform has the complementary shape
\begin{equation}\label{T:eq:zd-canonical-diamond}
 \mathcal{D}g_n\in
 \frac{\mathbb Z[t]_{<n}}{\Theta}.
\end{equation}
It is holomorphic away from \(c_*\), and in particular has an
integral power-series expansion at \(\infty\).
Here \(\mathbb Z[t]_{<n}\) denotes polynomials of degree less than \(n\).
Set \(g_{n/p}=0\) whenever \(p\nmid n\).

\subsection{Hecke relations for both members of the pair}

For Laurent series define
\[
 U_p\!\left(\sum_{j\in\mathbb Z}b_jq^j\right)
       =\sum_{j\in\mathbb Z}b_{pj}q^j,\qquad
 V_pf(q)=f(q^p).
\]
Only finitely many negative terms occur throughout.
Fix a prime \(p\ne5\), and put
\[
 \eta_p=
 \begin{cases}
 I,&p\equiv1\pmod5,\\
 \delta^{-1},&p\equiv2\pmod5,\\
 \delta,&p\equiv3\pmod5,\\
 -I,&p\equiv4\pmod5.
 \end{cases}
\]
Consider the determinant-\(p\) matrices
\[
 \mathcal{M}_p=
 \left\{\begin{pmatrix}a&b\\c&d\end{pmatrix}:
 \begin{gathered}
 a,b,c,d\in\mathbb Z,\quad ad-bc=p,\quad c\equiv0\pmod5,\\
 (a,d)\equiv(p,1)\ \hbox{or}\ (-p,-1)\pmod5
 \end{gathered}\right\}.
\]
Let \(\epsilon(\beta)\) be the lower-right residue of
\(\beta\) modulo 5, interpreted as \(1\) or \(-1\), and define
\[
 \mathcal{H}_pf
   =p^{1/2}\sum_{H\backslash\mathcal{M}_p}
                  \epsilon(\beta)\,f|_{-1}\beta.
\]
The sign is multiplicative under left and right multiplication
by \(H\), so this is again a weight \(-1\) form with character
\(\epsilon\).

Integral Hermite normal form gives the representatives
\[
 \beta_\infty=\begin{pmatrix}p&0\\0&1\end{pmatrix},
 \qquad
 \eta_p\beta_b,\quad
 \beta_b=\begin{pmatrix}1&b\\0&p\end{pmatrix}\quad(0\le b<p),
\]
For completeness, in a finite Hermite branch the intervening
unimodular matrix has diagonal residues
\(\pm(p,p^{-1})\) modulo 5; multiplying it by \(\eta_p^{-1}\)
places it in \(H\).  All displayed representatives have sign one.
Write \(\mathcal{E}_pf=f|_{-1}\eta_p\).  Thus
\(\mathcal{E}_p\) is \(1,-\mathcal{D},\mathcal{D},-1\), respectively, in
the four residue classes \(1,2,3,4\).
The Fourier expansions of the representatives give
\begin{equation}\label{T:eq:zd-hecke-operator}
 \mathcal{H}_p=V_p+p^2U_p\mathcal{E}_p.
\end{equation}
The matrix \(\delta\) normalizes \(H\), and conjugation by
\(\delta\) preserves \(\mathcal{M}_p\) and its sign:
modulo 5 all the matrices involved are upper triangular and
the two diagonal residues are unchanged by conjugation.
Therefore
\begin{equation}\label{T:eq:zd-commutation}
 \mathcal{H}_p\mathcal{D}=\mathcal{D}\mathcal{H}_p.
\end{equation}

\begin{lemma}[Exact canonical Hecke relation]
\label{T:lem:zd-canonical-hecke}
Set
\[
 \mathbf{G}_n=\begin{pmatrix}g_n\\\mathcal{D}g_n\end{pmatrix},\qquad
 M_p=
 \begin{cases}
 I_2,&p\equiv1\pmod5,\\
 \begin{pmatrix}0&-1\\1&0\end{pmatrix},&p\equiv2\pmod5,\\
 \begin{pmatrix}0&1\\-1&0\end{pmatrix},&p\equiv3\pmod5,\\
 -I_2,&p\equiv4\pmod5.
 \end{cases}
\]
For every \(n\ge1\), with the Hecke operator acting componentwise,
\begin{equation}\label{T:eq:zd-hecke-first}
 V_p\mathbf{G}_n+p^2M_pU_p\mathbf{G}_n
   =\mathcal{H}_p\mathbf{G}_n
   =\mathbf{G}_{pn}+p^2M_p\mathbf{G}_{n/p}.
\end{equation}
\end{lemma}
\begin{proof}
The first equality follows from
\eqref{T:eq:zd-hecke-operator} and
\(\mathcal{E}_p\mathbf{G}_n=M_p\mathbf{G}_n\).
It remains to identify the principal parts of \(\mathcal{H}_pg_n\).

First suppose \(p\equiv2\) or \(3\pmod5\), and write
\(c_p=-1\) or \(1\), respectively.  At \(\infty\) the only
negative term is \(q^{-pn}\), since \(\mathcal{D}g_n\) is
holomorphic there.  The expansion at \(c_*\), obtained using
\eqref{T:eq:zd-commutation}, is
\[
 \mathcal{H}_p\mathcal{D}g_n
      =V_p\mathcal{D}g_n-c_pp^2U_pg_n,
\]
whose principal part is \(-c_pp^2q^{-n/p}\) if \(p\mid n\),
and zero otherwise.
These are exactly the principal parts of
\(g_{pn}+c_pp^2\mathcal{D}g_{n/p}\).

Next suppose \(p\equiv\varepsilon_p=\pm1\pmod5\).
Then \(\mathcal{E}_p=\varepsilon_p\), so the principal part
of \(\mathcal{H}_pg_n=V_pg_n+\varepsilon_pp^2U_pg_n\) at
\(\infty\) is
\(q^{-pn}+\varepsilon_pp^2q^{-n/p}\), the latter term
being omitted when \(p\nmid n\).
At \(c_*\) the expansion is
\(V_p\mathcal{D}g_n+\varepsilon_pp^2U_p\mathcal{D}g_n\),
which is holomorphic.  Thus its principal parts coincide
with those of \(g_{pn}+\varepsilon_pp^2g_{n/p}\).

There are no poles at the other cusps.  Indeed a primitive cusp
column whose lower entry is not divisible by 5 still has lower
entry nonzero modulo 5 after applying a matrix in \(\mathcal{M}_p\);
division by its possible common factor \(p\) does not change that
property.  Such a cusp cannot map to \(\infty\) or \(c_*\).
All summands are holomorphic in the upper half plane as well.
In either case the difference is a holomorphic weight \(-1\)
form, and must vanish.  Applying \(\mathcal{D}\) gives the
second component of \eqref{T:eq:zd-hecke-first}.
\end{proof}

\begin{lemma}[Paired Cartier divisibility]
\label{T:lem:zd-paired-cartier}
For every \(n\ge1\),
\[
 U_pg_n-g_{n/p}\in
       p^{2v_p(n)}\mathbb Z[[q]],\qquad
 U_p\mathcal{D}g_n-\mathcal{D}g_{n/p}\in
       p^{2v_p(n)}\mathbb Z[[q]].
\]
\end{lemma}
\begin{proof}
Both differences are integral.  In the first one the principal
parts cancel; the second one is holomorphic at \(\infty\).
Write their column vector as \(\Delta_n\).
Apply \(U_p\) to \eqref{T:eq:zd-hecke-first}, using \(U_pV_p=1\).
The resulting exact vector relation is
\[
 \Delta_{pn}=p^2M_pU_p\Delta_n.
\]
Iterating once for each factor \(p\) in \(n\) proves the assertion,
since \(U_p\) and \(M_p\) preserve integral coefficients.
\end{proof}

\subsection{The differential residual supplies the weights}

Set
\[
 \theta=t\frac{d}{dt},\qquad
 A(t)=1-11t-t^2,\qquad
 L=\theta^2-t(11\theta^2+11\theta+3)-t^2(\theta+1)^2.
\]
Thus \(L=A\theta^2+(\theta A)\theta-3t-t^2\) and \(LF=0\).
Let \(\partial=q\,d/dq\).
The canonical-coordinate Wronskian gives
\[
 \partial=A(t)\Theta^2\theta.
\]
One can see the normalization directly: the second solution is
\(F(t)\log q(t)\), so the Wronskian equation yields
\(A F^2\,\theta\log q=1\), its constant being fixed at \(t=0\).
It follows by differentiating twice and using \(LF=0\) that,
for every Laurent polynomial \(P\),
\begin{equation}\label{T:eq:zd-bol}
 \partial^2\left(\frac{P(t)}{\Theta}\right)
       =A(t)\Theta^3\,LP(t).
\end{equation}

Write \(f_K=F_K(t(q))/\Theta(q)\).
The exact truncation residual is
\begin{equation}\label{T:eq:zd-residual}
 LF_K=-K^2t^K(A_K+A_{K-1}t).
\end{equation}
Indeed all coefficients below \(t^K\) cancel by the
Ap\'ery recurrence; the only remaining terms have degrees \(K\)
and \(K+1\), with the displayed coefficients.
Direct calculation on monomials also gives
\[
 L(JP)(t)=-t\,(LP)(-1/t).
\]
Consequently
\begin{equation}\label{T:eq:zd-dual-residual}
 L(JF_K)=(-1)^KK^2t^{-K}(A_Kt-A_{K-1}).
\end{equation}
Every factor outside \(K^2\) on the right sides of
\eqref{T:eq:zd-residual} and \eqref{T:eq:zd-dual-residual},
after multiplication by \(A\Theta^3\), has an integral
Laurent \(q\)-expansion.  Therefore
\begin{equation}\label{T:eq:zd-two-bol-integral}
 \partial^2 f_K\in K^2\mathbb Z[[q]],\qquad
 \partial^2\mathcal{D}f_K\in K^2\mathbb Z((q)).
\end{equation}

\subsection{Descent of the truncation and conclusion}

\begin{proof}[Proof of Theorem~\ref{T:thm:zagier-nonsplit-complete}]
Let \(p\mid K\), and set \(r=v_p(K)\).
Since \(\mathcal{D}f_K\) has poles only at \(\infty\), of order
at most \(K\), its canonical expansion is finite:
\[
 \mathcal{D}f_K=\sum_{n=1}^{K}c_ng_n,\qquad
 c_n=[q^{-n}]\mathcal{D}f_K\in\mathbb Z.
\]
The second assertion of \eqref{T:eq:zd-two-bol-integral} gives
\begin{equation}\label{T:eq:zd-principal-divisibility}
 v_p(c_n)+2v_p(n)\ge2r.
\end{equation}
Define the smaller-pole form
\[
 h=\sum_{\substack{1\le n\le K\\p\mid n}}c_ng_{n/p}.
\]
Multiply the two congruences in
Lemma~\ref{T:lem:zd-paired-cartier} by \(c_n\), and use
\eqref{T:eq:zd-principal-divisibility}.  Term by term the
product of the two divisibility factors is at least \(p^{2r}\).
It follows that
\begin{equation}\label{T:eq:zd-paired-descent}
 U_p\mathcal{D}f_K\equiv h,\qquad
 -U_pf_K\equiv\mathcal{D}h
 \pmod{p^{2r}\mathbb Z((q))}.
\end{equation}

By \eqref{T:eq:zd-canonical-diamond} there is an integral polynomial
\(P\), of degree strictly less than \(K/p\), such that
\[
 -\mathcal{D}h=\frac{P(t)}{\Theta}.
\]
On the other hand \(f_K=1+O(q^K)\), and hence
\(U_pf_K=1+O(q^{K/p})\).
The second congruence in \eqref{T:eq:zd-paired-descent} now gives
\[
 P(t(q))\equiv\Theta(q)\pmod{p^{2r},q^{K/p}}.
\]
The integral inverse \(q(t)=t+O(t^2)\), followed by comparison
of coefficients of degree \(<K/p\), yields
\[
 P(t)\equiv F_{K/p}(t)\pmod{p^{2r}}.
\]
Thus both members of the desired paired descent hold:
\begin{equation}\label{T:eq:zd-final-cartier}
 U_pf_K\equiv f_{K/p},\qquad
 U_p\mathcal{D}f_K\equiv\mathcal{D}f_{K/p}
 \pmod{p^{2r}}.
\end{equation}
For the second one, apply \(\mathcal{D}\) to the polynomial
congruence \(P\equiv F_{K/p}\), which preserves integral
Laurent expansions, and use the first part of
\eqref{T:eq:zd-paired-descent}.

Finally the first assertion of \eqref{T:eq:zd-two-bol-integral}
shows that each coefficient of \(f_K\) at an exponent not
divisible by \(p\) is divisible by \(p^{2r}\).
These are exactly the coefficients retained by \(1-V_pU_p\).
Therefore
\[
 f_K\equiv V_pU_pf_K\equiv V_pf_{K/p}\pmod{p^{2r}}.
\]
Take \(K=mp^s\).  Then \(r\ge s\), so this says
\[
 \frac{F_{mp^s}(t)}{F(t)}
 \equiv
 \frac{F_{mp^{s-1}}(\sigma_p(t))}{F(\sigma_p(t))}
 \pmod{p^{2s}}.
\]
All four denominator series have constant term one.
Cross multiplication and division by units gives
Theorem~\ref{T:thm:zagier-nonsplit-complete}.
\end{proof}

Every step used only \(p\nmid5\).
In particular no Jacobsthal congruence, odd-prime assumption,
ordinary-point assumption, or higher-level induction
hypothesis is required.

\section{The first level in every dimension}\label{T:source-main-3}
\label{T:sec:all-dim-first}

The following argument proves the full formal-series congruence at
\(s=1\) for every dimension of the three families in Conjecture 7.5.
The higher-level ordinary-truncation proof is given independently in Part~\ref{part:positive}.

Let \(g\) be the simplicial, hypercubic, or hyperoctahedral Laurent
polynomial in dimension \(n\), with vertex coefficient one, and let
\(f=1-zg\).  Write
\[
 c_j=\operatorname{CT}(g^j),\qquad
 \mathcal F(z)=\sum_{j\geq0}c_jz^j,\qquad
 H_N(z)=\operatorname{CT}(f^{N-1}).
\]
Let \(\tau=\sigma_p(z)\) be the excellent lift and
\(\lambda=\mathcal F(z)/\mathcal F(\tau)\).
Assume the prime satisfies the symmetry and support-lattice exclusions
of \cite[Theorem 7.3]{BV3}.  In particular, \(p>n\) for the cubic and
octahedral families and \(p>n+1\) for the simplicial family, apart from
irrelevant low-dimensional simplifications.

\begin{lemma}[Second-order geometric approximation]
\label{T:lem:second-geometric}
For the above excellent lift,
\[
 \lambda\equiv2H_p-H_{2p}\pmod{p^2\mathbb Z_p[[z]]}.
\]
More precisely, write
\[
 \mathscr C_p(f^{2p-1})=H_{2p}(z)+B_p(z)g(\mathbf{x}).
\]
Then
\[
 \tau\equiv-\frac{B_p(z)}{H_{2p}(z)}\pmod{p^2\mathbb Z_p[[z]]}.
\]
\end{lemma}

\begin{proof}
All the calculations below may be made in the ring of Laurent series
whose coefficients are \(z\)-adically convergent integral power series.
Put \(h(\mathbf{x})=f^\sigma(\mathbf{x}^p)\).  Since \(h-f^p\) is
divisible by \(p\), the exact identity
\[
 \frac1f-\frac{2hf^{p-1}-f^{2p-1}}{h^2}
 =\frac{(h-f^p)^2}{fh^2}
\]
has an integral remainder divisible by \(p^2\).  Applying the Cartier
operator and using
\(\mathscr C_p(A(\mathbf{x})/h(\mathbf{x})^2)
 =\mathscr C_p(A) / f^\sigma(\mathbf{x})^2\), we obtain
\[
 \mathscr C_p(1/f)
 \equiv
 \frac{2f^\sigma\mathscr C_p(f^{p-1})
              -\mathscr C_p(f^{2p-1})}{(f^\sigma)^2}
 \pmod{p^2}.
\]

We justify the two numerator supports explicitly.  Let \(\Gamma\)
be the lattice spanned by the support of \(g\).  Since its index is
prime to \(p\), \(p\mathbf{u}\in\Gamma\) implies
\(\mathbf{u}\in\Gamma\).  Reflexivity gives integral facet functions
\(L_j\) with \(\Delta=\{L_j\leq1\}\).  A lattice point in
\((1-1/p)\Delta\) has \(L_j\leq0\) for every \(j\), and hence is
zero.  A lattice point in \((2-1/p)\Delta\) has \(L_j\leq1\), and
hence belongs to \(\Delta\).  The only points of
\(\Delta\cap\Gamma\) are zero and the vertices.  Complete symmetry
therefore yields the exact identities
\[
 \mathscr C_p(f^{p-1})=H_p,
 \qquad
 \mathscr C_p(f^{2p-1})=H_{2p}+B_p g.
\]
For the hypercube, the restriction to \(\Gamma\) is essential; its
points have coordinates all even or all odd.

The excellent Cartier relation
\(\mathscr C_p(1/f)\equiv\lambda/f^\sigma
 \pmod{p^2\mathscr F_2^\sigma}\)
is in particular a coefficientwise congruence modulo \(p^2\).
Multiplying by \((f^\sigma)^2\) gives
\[
 \lambda(1-\tau g)
 \equiv 2H_p(1-\tau g)-H_{2p}-B_p g\pmod{p^2}.
\]
Since \(\operatorname{CT}(g)=0\), comparison of the coefficients of
\(1\) and \(g\) first gives
\(\lambda\equiv2H_p-H_{2p}\), and then
\(\tau H_{2p}\equiv-B_p\).
The latter can be divided by the unit \(H_{2p}\), whose constant
term is one.
\end{proof}

\begin{theorem}[Conjecture 7.5 at the first level, in every dimension]
\label{T:thm:first-all-dim}
Set \(h=2\) for the hypercubic and hyperoctahedral families and
\(h=n+1\) for the simplicial family.  At every admissible prime,
\[
 \frac{\mathcal F(z)}{\mathcal F(\sigma_p(z))}
 \equiv\mathcal F_{hp}(z)\pmod{p^2\mathbb Z_p[[z]]}.
\]
Equivalently, in the normalized variable \(u=z^h\),
\[
 \frac{F(u)}{F(\sigma_p(u))}
 \equiv F_p(u)\pmod{p^2\mathbb Z_p[[u]]}.
\]
These are the \(s=1\) cases of the specified original truncation
lengths; \(\mathcal F_h=1\).
\end{theorem}

\begin{proof}
The coefficient of \(z^j\) in \(2H_p-H_{2p}\) is
\[
 c_jw_j,\qquad
 w_j=(-1)^j\left(2\binom{p-1}{j}-\binom{2p-1}{j}\right),
\]
with binomial coefficients zero past the top index.  If \(0\leq j<p\),
the product expansion
\[
 (-1)^j\binom{ap-1}{j}
 =\prod_{i=1}^j\left(1-\frac{ap}{i}\right)
 \equiv1-apH_j\pmod{p^2}
\]
shows that \(w_j\equiv1\pmod{p^2}\).
For \(p\leq j<2p\), Lucas's congruence gives
\(w_j\equiv1\pmod p\).

The remaining required divisibilities follow directly from the
factorial formulas.  In the cubic case,
\(c_{2k}=\binom{2k}{k}^n\), and for \(k<p\) this is divisible by
\(p\) whenever \(2k\geq p\).
In the octahedral case,
\[
 c_{2k}=\sum_{k_1+\cdots+k_n=k}
              \frac{(2k)!}{(k_1!\cdots k_n!)^2},
\]
and every denominator is prime to \(p\) when \(k<p\); hence the
same divisibility holds term by term.
In the simplicial case, putting \(r=n+1<p\),
\[
 c_{rk}=\frac{(rk)!}{(k!)^r},\qquad 0\leq k<p,
 \qquad v_p(c_{rk})\geq\left\lfloor\frac{rk}{p}\right\rfloor.
\]
Thus in all three cases \(c_j\) is divisible by \(p\) for
\(p\leq j<\min(2p,hp)\), and in the simplicial case it is divisible
by \(p^2\) for \(2p\leq j<hp\).  The other two families have
\(h=2\), so there is no latter range.
It follows coefficientwise that
\[
 2H_p-H_{2p}\equiv\mathcal F_{hp}\pmod{p^2}.
\]
Apply Lemma~\ref{T:lem:second-geometric} and then the normalized
coordinate relation \(q_0(z)^h=q(z^h)\).
\end{proof}

\paragraph{Limitation.}
Multiplication of the excellent Cartier relation by
\((f^\sigma)^{2p^{s-1}}\) does not automatically preserve the
\(p^{2s}\) coefficient divisibility of the error.  Its support is
not restricted to exponent vectors divisible by \(p^{s-1}\).
Accordingly, this proof supplies a complete base case in every
dimension. The separate higher-level induction is the positive-vertex descent of Part~\ref{part:positive}.

\section{A canonical negative-power basis in all three families}\label{T:source-main-4}
\label{T:sec:negative-basis}

The following all-level congruence concerns an auxiliary canonical
basis.  It is not the ordinary positive-power truncation congruence.
It isolates which half of the two-cusp argument survives without a
modular parametrization.

Work at an admissible prime \(p\) for one of the three completely
symmetric families.  Put \(h=2\) for the cubic and octahedral families,
\(h=n+1\) for the simplicial family, and normalize \(u=z^h\).
Write \(F(u)=\mathcal F(z)\), \(q(u)=q_0(z)^h\), and
\(\Theta(q)=F(u(q))\).  Both \(q(u)\) and its inverse are integral
at \(p\); here integrality always means over \(\mathbb Z_p\).

There is a unique canonical basis
\[
 g_N\in\frac{u^{-1}\mathbb Z_p[u^{-1}]}{F(u)},\qquad
 g_N(q)=q^{-N}+O(1),\qquad N\geq1.
\]
Existence and uniqueness follow from the integral triangular leading
terms of \(u^{-j}/F\).  Equivalently,
\[
 g_N(u)=\frac{[F(u)q(u)^{-N}]_{u<0}}{F(u)},
\]
where the numerator means the finite negative-power part in \(u\).

\begin{theorem}[Canonical-basis Frobenius congruence]
\label{T:thm:negative-basis-frob}
For every integer \(N\geq1\),
\[
 g_{pN}(q)\equiv g_N(q^p)
 \pmod{p^{2(1+v_p(N))}\mathbb Z_p((q))}.
\]
Consequently
\[
 U_pg_N-g_{N/p}
 \in p^{2v_p(N)}\mathbb Z_p[[q]],
\]
where \(g_{N/p}=0\) if \(p\nmid N\).
\end{theorem}

\begin{proof}
We first construct an integral triangular family of vertex coefficients.
For a positive integer \(M\), take the formal expansion of \(1/f\)
at the vertex with monomial \((x_1\cdots x_n)^{-1}\) in the simplex
and cube, or \(x_1^{-1}\) in the octahedron.  We allow \(z^{-1}\)
in the coefficient ring, as in \cite[Proposition 7.17]{BV3}.
Coefficient extraction is performed after localization at \(z\) and
\(p\)-adic completion.  The particular series and differences used
below have bounded negative support, so their divisibility statements
lie in the ordinary subring \(\mathbb Z_p((z))\); no such assertion
is made for every element of the completed coefficient ring.
Let \(\mathcal P_M\) be minus the coefficient at
\((x_1\cdots x_n)^M\), or at \(x_1^M\) with the other exponents
zero, respectively.  Direct geometric expansion gives
\[
 \mathcal P_{hN}(z)=u^{-N}A_N(u),\qquad
 A_N\in1+u\mathbb Z[u],\quad\deg A_N<N.
\]
For clarity, the three finite formulas are as follows.  With \(M=hN\),
the cubic coefficients are
\[
 \mathcal P_M(z)=z^{-M}
 \sum_{0\leq j<M/2}(-1)^{nj}
         \binom{M-j-1}{j}^{n}z^{2j}.
\]
For the simplex, \(r=n+1\) and
\[
 \mathcal P_M(z)=z^{-M}
 \sum_{0\leq j<M/r}(-1)^{(r-1)j}
 \frac{(M-j-1)!}{(M-rj-1)!(j!)^{r-1}}z^{rj}.
\]
For the octahedron,
\[
 \mathcal P_M(z)=z^{-M}
 \sum_{0\leq j<M/2}z^{2j}
 \sum_{k_1+\cdots+k_n=j}
 (-1)^{k_1}
 \frac{(M-k_1-1)!}
 {(M-2j-1)!k_1!(k_2!\cdots k_n!)^2}.
\]
Each formula is a coefficient of an integral geometric series, so its
integrality does not depend on a factorial-denominator estimate.

Set \(R_N=\mathcal P_{hN}/\mathcal F=A_N/(u^NF)\).
The excellent Cartier relation, read at the exponent vector of length
\(hpN\), gives
\begin{equation}\label{T:eq:negative-triangular-frob}
 R_{pN}(q)\equiv R_N(q^p)
 \pmod{p^{2(1+v_p(N))}}.
\end{equation}
This is the same coefficient-reading argument as
\cite[Proposition 7.17]{BV3}, and it applies to even indices as well;
the displayed finite formulas verify the required expansions directly.
The condition \(p\nmid h\) ensures the stated valuation.

There is a second useful consequence of the rank-two connection.
Let \(\theta_z=z\,d/dz\), and let
\(L_2=\theta_z^2-B\theta_z-A\) be the operator of
\cite[Proposition 7.6]{BV3}, so
\(L_2(1/f)\in\mathscr F_2\).
Reading the same vertex coefficient gives
\(L_2\mathcal P_{hN}\in (hN)^2\mathbb Z_p((z))\).
The solutions \(\mathcal F\) and \(\mathcal F\log q_0\)
gauge this operator to a scalar multiple of \(\partial_0^2\),
where \(\partial_0=q_0\,d/dq_0\).
Explicitly,
\(\partial_0^2(\mathcal P_{hN}/\mathcal F)
 =\mathcal F^3W^{-2}L_2\mathcal P_{hN}\).
The scalar and its inverse are integral units: this follows from
\(\mathcal F(0)=1\) and the integral unit Wronskian in
\cite[Proposition 7.9]{BV3}.  Since
\(\partial=q\,d/dq=h^{-1}\partial_0\), we obtain
\begin{equation}\label{T:eq:negative-weighted-principal}
 \partial^2R_N\in N^2\mathbb Z_p((q)).
\end{equation}

Expand \(R_N=\sum_{j=1}^Nc_{N,j}g_j\).
The coefficients are its negative Fourier coefficients, so they are
integral and \(c_{N,N}=1\).
By \eqref{T:eq:negative-weighted-principal},
\[
 v_p(c_{N,j})+2v_p(j)\geq2v_p(N).
\]
Comparing principal parts in \eqref{T:eq:negative-triangular-frob}
and using the integral canonical basis yields
\[
 R_{pN}\equiv\sum_{j=1}^Nc_{N,j}g_{pj}
 \pmod{p^{2(1+v_p(N))}}.
\]
Subtracting \eqref{T:eq:negative-triangular-frob} now gives
\[
 \sum_{j=1}^Nc_{N,j}(g_{pj}-V_pg_j)
 \equiv0\pmod{p^{2(1+v_p(N))}}.
\]
Induct on the ordinary positive integer \(N\).
For \(j<N\), the induction hypothesis contributes
\(2(1+v_p(j))\) to the valuation, and the displayed bound on
\(c_{N,j}\) raises the total to at least \(2(1+v_p(N))\).
The final coefficient is \(c_{N,N}=1\), so the desired
congruence for \(g_{pN}-V_pg_N\) follows.
Applying \(U_p\) proves the second assertion when \(p\mid N\).
For \(p\nmid N\), it simply asserts integrality and absence of
negative terms, which follow from the canonical principal part.
\end{proof}

\paragraph{Relation to ordinary truncations.}
The basis consists of negative-power polynomials divided by \(F\).
The target truncation \(F_{p^s}/F\) belongs to the positive-power
space.  The integral diamond operation in the Zagier D argument
interchanges these two spaces and supplies the missing relation there.
No corresponding operation has been established for the general
higher-dimensional families, so the theorem above is not a proof of
their ordinary truncation congruence. Part~\ref{part:positive} proves that congruence using a different, positive-vertex construction.

\section{The three-dimensional hyperoctahedral family: a weighted Cartier reduction}\label{T:source-main-5}
\label{T:sec:octa-weighted}

This section proves an all-coefficient differential divisibility statement
for the three-dimensional hyperoctahedral family. It is an additional differential refinement; the full ordinary truncation congruence follows from Theorem~\ref{P:thm:main-all}.

\subsection{The period and its actual Beukers--Vlasenko coordinate}

Use the normalized variable $u=t^2$, and put
\[
 b_n=\sum_{j=0}^n\binom nj^2\binom{2j}{j},\qquad
 a_n=\binom{2n}{n}b_n,\qquad F(u)=\sum_{n\geq0}a_nu^n.
\]
Equivalently,
\[
 b_n=\sum_{k_1+k_2+k_3=n}
       \left(\frac{n!}{k_1!k_2!k_3!}\right)^2,
 \qquad
 a_n=\sum_{k_1+k_2+k_3=n}
       \frac{(2n)!}{(k_1!k_2!k_3!)^2}.
\]
The normalized logarithmic companion prescribed by
\cite[\S7 and Proposition 9.1]{BV3} is
\begin{equation}\label{T:eq:octa-BV-G}
 G(u)=2\sum_{n\geq1}u^n
 \sum_{k_1+k_2+k_3=n}
 \frac{(2n)!}{(k_1!k_2!k_3!)^2}(H_{2n}-H_{k_1}).
\end{equation}
Thus $q(u)=u\exp(G(u)/F(u))$ is exactly the square of the
original BV coordinate, not merely a coordinate with the same first term.
It and its inverse $u(q)$ are $p$-integral for every $p\geq5$.

Write $\theta=u\,d/du$ and
\begin{equation}\label{T:eq:octa-L}
 L=\theta^3-2u(2\theta+1)(10\theta^2+10\theta+3)
       +36u^2(\theta+1)(2\theta+1)(2\theta+3).
\end{equation}

\begin{lemma}[Logarithmic companion with a telescoping certificate]
\label{T:lem:octa-logcomp}
Both $F$ and $F\log u+G$, with $G$ as in
\eqref{T:eq:octa-BV-G}, are solutions of $Ly=0$.
\end{lemma}
\begin{proof}
For a formal parameter $\rho$, define
\[
 b_n(\rho)=\sum_{j=0}^n\binom{n+\rho}{j}^{\!2}\binom{2j}{j},
 \qquad
 c_n(\rho)=\frac{(1+2\rho)_{2n}}{(1+\rho)_n^2},
 \qquad a_n(\rho)=c_n(\rho)b_n(\rho).
\]
Here $(x)_n$ is the rising factorial.  We work modulo $\rho^2$.
Let
\[
 T(z,j)=\binom zj^2\binom{2j}{j},\qquad
 R(z,j)=\frac{j^3(3j-4z-4)}{(z+1-j)^2}.
\]
Direct cancellation of rational functions gives the certificate
\begin{align*}
 &(z+1)^2T(z+1,j)-(10z^2+10z+3)T(z,j)
                      +9z^2T(z-1,j)\\
 &\hspace{18mm}=R(z,j+1)T(z,j+1)-R(z,j)T(z,j).
\end{align*}
Substitute $z=n-1+\rho$ and sum over $0\leq j\leq n$.
The lower boundary is zero and the upper boundary is $O(\rho^2)$.
Although $R(z,n)$ has a displayed denominator $\rho^2$, its product
with $T(z,n)$ is regular: the zero of the binomial square cancels it.
Terms beyond the finite cutoffs in the other two sums are also
$O(\rho^2)$; when $n=1$, the factor $9\rho^2$ handles the last sum.
Consequently, with $b_{-1}(\rho)=0$,
\[
 (n+\rho)^2b_n(\rho)
 -\{10(n-1+\rho)^2+10(n-1+\rho)+3\}b_{n-1}(\rho)
 +9(n-1+\rho)^2b_{n-2}(\rho)=O(\rho^2).
\]
The ratio
\[
 \frac{c_n(\rho)}{c_{n-1}(\rho)}
       =\frac{2(2n+2\rho-1)}{n+\rho}
\]
shows that the shifted recurrence for $L$ applied to $a_n(\rho)$
is $c_n(\rho)(n+\rho)$ times the preceding recurrence.
Its constant-index residual is $\rho^3$.
Therefore
\[
 L\left(u^\rho\sum_{n\geq0}a_n(\rho)u^n\right)=O(\rho^2).
\]
Taking its constant and linear terms proves the claim, because
\[
 c_n'(0)=2\binom{2n}{n}(H_{2n}-H_n),\qquad
 b_n'(0)=2\sum_{j=0}^nT(n,j)(H_n-H_{n-j}).
\]
Grouping the other two multinomial indices identifies $a_n'(0)$
with the coefficient in \eqref{T:eq:octa-BV-G}.
\end{proof}

\subsection{Symmetric square, Bol identity, and an exact residual}

Put
\[
 A(u)=(1-4u)(1-36u)=1-40u+144u^2.
\]
The monic version of $L$ is the symmetric square of
\[
 M=\theta^2+
     \frac{4u(36u-5)}{A(u)}\theta+
     \frac{3u(9u-1)}{A(u)}.
\]
This follows by direct substitution in the symmetric-square formula
\[
 (\theta^2+B\theta+C)^{\operatorname{Sym}^2}
 =\theta^3+3B\theta^2+(\theta B+4C+2B^2)\theta
                       +(2\theta C+4BC).
\]
If $\varphi_0=\sqrt F$ and
$\varphi_1=\varphi_0\log u+\cdots$ are its normalized solutions,
Lemma~\ref{T:lem:octa-logcomp} gives
$\log q=\varphi_1/\varphi_0$.
Their $\theta$-Wronskian is $A^{-1/2}$, normalized to constant term 1,
and hence
\begin{equation}\label{T:eq:octa-q-derivative}
 \theta\log q=\frac{1}{F\sqrt A},\qquad
 \partial:=q\frac{d}{dq}=F\sqrt A\,\theta.
\end{equation}
It follows, either by direct differentiation or by comparing the
leading terms and the three-dimensional kernels, that
\begin{equation}\label{T:eq:octa-Bol}
 \partial^3\left(\frac{P(u(q))}{F(u(q))}\right)
       =F(u(q))^2\sqrt{A(u(q))}\,(LP)(u(q))
\end{equation}
for every formal series $P$.

For $K\geq1$ let $F_K(u)=\sum_{0\leq n<K}a_nu^n$.
The recurrence obtained from \eqref{T:eq:octa-L} gives the exact identity
\begin{equation}\label{T:eq:octa-truncated-residual}
 LF_K=-K^3a_Ku^K
       +36K(2K-1)(2K+1)a_{K-1}u^{K+1}.
\end{equation}
In particular, if $K=mp^s$ with $m,s\geq1$, then
\[
 LF_K\in p^{2s}\mathbb Z_p[u].
\]
Indeed, the first coefficient is divisible by $p^{3s}$, and
$\binom{2K-2}{K-1}$ has at least $s$ carries when $K-1$ is doubled
in base $p$.  Thus $p^s\mid a_{K-1}$.

\subsection{An integral primitive improves the weighted bound}

For $j\geq0$ define
\[
 I_j(u)=\int_0^u z^jF(z)\,dz
       =\sum_{n\geq0}\frac{a_n}{n+j+1}u^{n+j+1}.
\]

\begin{lemma}[Two shifted Catalan divisibilities]
\label{T:lem:octa-primitives}
Let $p\geq5$, $s\geq1$, and $K=p^s$.  Then
\[
 I_K(u)\in\mathbb Z_p[[u]],\qquad
 K I_{K-1}(u)\in\mathbb Z_p[[u]].
\]
\end{lemma}
\begin{proof}
Only the central binomial factor in $a_n$ is needed.
If $r=v_p(n+K+1)$, then for $1\leq i\leq\min(r,s)$,
the residue of $n$ modulo $p^i$ is $p^i-1$.
For $s<i\leq r$ it is $p^i-p^s-1$.
Every displayed residue is at least $p^i/2$ when $p\geq5$.
Each corresponding summand in Legendre's formula
\[
 v_p\binom{2n}{n}
 =\sum_{i\geq1}\left(\left\lfloor\frac{2n}{p^i}\right\rfloor
                    -2\left\lfloor\frac n{p^i}\right\rfloor\right)
\]
is therefore 1, so $v_p(a_n)\geq r$.

For the second assertion put $r=v_p(n+K)$.
If $r\leq s$, the factor $K$ already cancels the denominator.
If $r>s$, then for every $s<i\leq r$ the residue of $n$
modulo $p^i$ is $p^i-p^s\geq p^i/2$.
Thus $v_p(a_n)\geq r-s$, which again proves integrality.
\end{proof}

\begin{theorem}[All-coefficient weighted divisibility]
\label{T:thm:octa-weighted}
For every prime $p\geq5$ and every $s\geq1$, put $K=p^s$ and
\[
 H_K(q)=\frac{F_K(u(q))}{F(u(q))}=1+\sum_{n\geq1}h_{K,n}q^n.
\]
Then
\begin{equation}\label{T:eq:octa-second-derivative}
 \partial^2H_K\in p^{2s}\mathbb Z_p[[q]].
\end{equation}
Equivalently, for every $n\geq1$,
\begin{equation}\label{T:eq:octa-weighted-coeff}
 v_p(h_{K,n})\geq\max\{0,\,2s-2v_p(n)\}.
\end{equation}
In particular every coefficient with $p\nmid n$ is divisible by $p^{2s}$.
\end{theorem}
\begin{proof}
From \eqref{T:eq:octa-q-derivative},
\[
 \partial I_j(u(q))
       =u(q)^{j+1}F(u(q))^2\sqrt{A(u(q))}.
\]
Integrate \eqref{T:eq:octa-Bol} with \eqref{T:eq:octa-truncated-residual}
once with respect to $\log q$.  All terms have constant term zero,
so no integration constant appears.  The resulting exact identity is
\begin{equation}\label{T:eq:octa-integrated-residual}
 \partial^2H_K
 =-K^3a_K I_{K-1}(u(q))
  +36K(2K-1)(2K+1)a_{K-1}I_K(u(q)).
\end{equation}
The first term is divisible by $K^2$ by
Lemma~\ref{T:lem:octa-primitives}; so is the second, because
$p^s\mid a_{K-1}$.
The integral inverse coordinate $u(q)$ preserves divisibility.
Finally $[q^n]\partial^2H_K=n^2h_{K,n}$, proving the coefficient bound.
\end{proof}

\subsection{A Cartier reformulation of the ordinary truncation congruence}

Write $U_p(\sum h_nq^n)=\sum h_{pn}q^n$ and
$V_pH(q)=H(q^p)$.  The theorem gives
\[
 H_{p^s}-V_pU_pH_{p^s}\in p^{2s}\mathbb Z_p[[q]].
\]
Consequently the full desired excellent-lift congruence in the
original conjecture is now equivalent to
\begin{equation}\label{T:eq:octa-unresolved-U}
 U_pH_{p^s}\equiv H_{p^{s-1}}\pmod{p^{2s}\mathbb Z_p[[q]]}.
\end{equation}
Indeed its normalized cleared defect divided by the period units is
$V_pH_{p^{s-1}}-H_{p^s}$, and $V_p$ is coefficientwise injective.
Equation~\eqref{T:eq:octa-unresolved-U} follows, at every level, from Theorem~\ref{P:thm:main-all} and the equivalence just proved. The differential estimate alone controls the nonresonant coefficients and quantifies the loss at the other indices; its role is distinct from the positive-vertex proof of the full Cartier congruence.

\subsection{A rigorously compatible Domb transformation}

There is also an exact compatible modular transformation. Let
\[
 d_n=(-1)^n\sum_{j=0}^n\binom nj^2\binom{2j}{j}
                            \binom{2n-2j}{n-j},\qquad
 F_D(v)=\sum_{n\geq0}d_nv^n,
\]
and put $D(v)=1+20v+64v^2=(1+4v)(1+16v)$.
The twisted Domb operator from \cite[Appendix C, \S9.7]{Beu26} is
\[
 L_D=\vartheta^3+2v(2\vartheta+1)(5\vartheta^2+5\vartheta+2)
                  +64v^2(\vartheta+1)^3,
 \qquad \vartheta=v\frac d{dv}.
\]
For $u=v/D(v)$, direct substitution gives the operator identity
\[
 \bigl(L[\sqrt D\,y]\bigr)(u(v))
      =\frac{\sqrt{D(v)}}{1-64v^2}(L_Dy)(v).
\]
Here the left side means that $\sqrt{D(v)}y(v)$ is viewed as a function
of $u$ via the inverse substitution.  Uniqueness of the normalized
holomorphic and logarithmic solutions, together with
Lemma~\ref{T:lem:octa-logcomp}, proves
\[
 F(u(v))=\sqrt{D(v)}F_D(v),\qquad
 q_{\rm BV}(u(v))=q_D(v).
\]
In particular the lift is genuinely transported by this transformation.
The Domb mirror parameter is
\[
 v(q)=\left(\frac{\eta(2\tau)\eta(6\tau)}
                   {\eta(\tau)\eta(3\tau)}\right)^6,
 \qquad q=e^{2\pi i\tau},
\]
and the octahedral form is
\[
 F(u(q))=q\frac d{dq}\log v(q)
 =\frac{-E_2(q)+2E_2(q^2)-3E_2(q^3)+6E_2(q^6)}4.
\]
Transport of the infinite period and coordinate does not by itself
transport the ordinary finite truncations.  In particular, no full
octahedral congruence is inferred solely from the established Domb
congruences.  Lagrange inversion gives the exact coefficient kernel
\[
 a_n=\sum_{j=0}^n\frac{2j+1}{2n+1}\,d_j
          [v^{n-j}]D(v)^{n+1/2},
\]
which specifies a possible additional truncation-compatibility problem.


\section{Exact verification and the role of the auxiliary arguments}\label{T:source-main-6}
\label{T:sec:verification-remaining}

The accompanying Python programs use exact integer and rational
arithmetic and explicit reduction modulo prime powers.  They are
self-contained and need only the Python standard library.  Their
recorded outputs declare all finite coefficient bounds.

For Zagier D the checks include integral canonical principal parts,
both components of the exact Hecke identities, the paired Cartier
recurrence, Laurent-series descent including its negative powers, and
the original full cleared defect beyond its truncation boundary.
In particular the computations include $p=2,3$, nonsplit primes
$7,13$, split primes $11,19$, higher levels, and $m=2$.
The largest recorded check reaches $q^{1380}$ at $p=7,s=3,m=2$.
The all-dimensional checks verify the geometric numerator,
the explicit excellent lift modulo $p^2$, and the full first-level
quotient beyond the cutoff in dimensions three through six.
The auxiliary canonical-basis checks cover ten cases in dimensions
three through five, including indices with a nontrivial prime-to-$p$
factor.  The octahedral checks verify the telescoping certificate,
the BV coordinate against the modular inverse through degree 39,
and six weighted differential bounds.  They also sample the
full octahedral congruence of Theorem~\ref{P:thm:main-all} at $(p,s)=(5,1),(5,2),(5,3),
(7,1),(7,2),(11,2)$, through the respective degrees
$48,128,528,56,224,512$; all observed minima are $2s$.
These calculations verify the implementation and sample the identities;
the infinite statements depend on the written proofs.

\subsection{Why the auxiliary arguments require a separate descent}

The excellent Cartier theorem gives
\[
 \mathscr C_p(1/f)-\frac{\mathcal F(z)}{\mathcal F(\sigma_p(z))}
                         \frac1{f^\sigma}
 \in p^2\mathscr F_2^\sigma.
\]
Reading a spatial coefficient with exponent vector divisible by
$p^{s-1}$ supplies a factor $p^{2s}$.  The operation defining the
ordinary period truncation is a different linear functional.  In
the notation of \cite[Section 4.3]{Vl24}, it is obtained from
period maps involving powers of $zg(\mathbf{x})$.  Their generic
divisibility does not automatically double merely because the
Cartier error is a second formal derivative.

The proof of Section~\ref{T:sec:all-dim-first} avoids that difficulty
at the first level using the exact two-term geometric numerator.
For levels $s\ge2$, multiplying a Cartier error by a larger
polynomial can mix spatial exponent vectors with smaller
$p$-adic valuation.  A valid induction must control this mixing,
or prove a separate descent identity for the actual truncations.
The auxiliary theorem in Section~\ref{T:sec:negative-basis} controls
negative-power canonical functions at all levels, but supplies no
identification with the ordinary positive-power truncations.
The required ordinary-tail identification is instead supplied by the positive-vertex basis in Part~\ref{part:positive}.

The new Zagier D proof resolves its analogous difficulty by the
paired canonical-basis recurrence.  Its basis forms have poles in a
finite set of cusps preserved by the Hecke correspondence.  In the
three-dimensional hyperoctahedral modular parametrization, the
relevant zero of the weight-two form is elliptic, so that same
finite-cusp argument is not directly available.  Section~\ref{T:sec:octa-weighted}
instead proves the second-derivative estimate at all levels and
isolates the exact additional relation \eqref{T:eq:octa-unresolved-U}. Theorem~\ref{P:thm:main-all} supplies that relation and the higher-dimensional ordinary-truncation bounds by positive interpolation and matrix block descent.



% END RETAINED SOURCE T

\clearpage

\appendix


\section{An independent coefficientwise certificate in dimension four}\label{P:source-appendix-1}
\label{P:sec:octa4-certificate}

This appendix gives a second, independent exact certificate for one
case of Theorem~\ref{P:thm:main-all}. Its finite polynomial calculation
proves every coefficient at the specified prime and level.

\begin{theorem}[Dimension four, \(p=5\), \(s=2\)]
\label{P:thm:octa4-p5-s2}
Let
\[
 a_j=\binom{2j}{j}
       \sum_{k=0}^{j}\binom jk^2\binom{2k}{k}
                             \binom{2j-2k}{j-k},\qquad
 F(u)=\sum_{j\geq0}a_ju^j.
\]
Let \(\sigma(u)\) be the actual normalized Beukers--Vlasenko
excellent lift for the four-dimensional hyperoctahedral family.
Then
\[
 F(u)F_5(\sigma(u))-F(\sigma(u))F_{25}(u)
       \in 625\mathbb Z_5[[u]].
\]
This is the entire formal-series congruence in Conjecture~7.5 at
\(p=5,s=2\), with original truncation lengths \(50\) and \(10\).
\end{theorem}

The proof is an exact finite polynomial certificate.  Crucially,
the rational lift used in the certificate is constructed from
vertex polynomials, independently of the desired ordinary
truncation identity.

\subsection{Vertex polynomials and a derivative bound}

Put
\[
 b_j=\sum_{k=0}^{j}\binom jk^2\binom{2k}{k},\qquad
 A_K(u)=\sum_{j=0}^{K-1}u^j
  \sum_{\ell=0}^{j}(-1)^\ell
  \frac{(2K-\ell-1)!\,b_{j-\ell}}
       {(2K-2j-1)!\,\ell!\,(j-\ell)!^2}.
\]
Each summand is an integer: the factorial ratio is a multinomial
coefficient.  Also \(A_K(0)=1\).
Expanding \(1/(1-z\sum_{i=1}^4(x_i+x_i^{-1}))\) in the vertex
chamber where \(x_1\) is dominant, the negative of the coefficient
at \(x_1^{-2K}\), with all other exponents zero, is
\[
 P_{2K}(u)=u^{-K}A_K(u),\qquad u=z^2.
\]
Consequently the excellent Cartier relation
\cite[Theorem 7.3]{BV3} gives, with
\(\lambda=F(u)/F(\sigma(u))\),
\begin{equation}\label{P:eq:octa4-vertex}
 P_{2pK}(u)-\lambda P_{2K}(\sigma(u))
       \in p^{2+2v_p(K)}\mathbb Z_p((u)).
\end{equation}
These displayed series have bounded negative support; no
identification of the whole completed coefficient ring with
\(\mathbb Z_p((u))\) is needed.

We shall also use
\begin{equation}\label{P:eq:octa4-ratio-derivative}
 \theta C_K\in p^{v_p(K)}\mathbb Z_p[[u]],\qquad
 C_K=\frac{A_{2K}}{A_K},\qquad \theta=u\frac{d}{du}.
\end{equation}
Here is a direct justification from the second-order quotient.
Set \(R_K=P_{2K}/F\), and use the actual canonical coordinate
\(q\), with \(\partial=q\,d/dq\).  Reading the coefficient at
the same vertex in the relation of
\cite[Proposition 7.6]{BV3}, then dividing by the period and
changing to the canonical coordinate, gives
\[
 \partial^2R_K\in p^{2v_p(K)}\mathbb Z_p((q)).
\]
The change-of-coordinate and gauge factors are integral units
by \cite[Proposition 7.9 and Corollary 7.11]{BV3}; the factor
\(2\) from \(u=z^2\) is a unit here.  If \(r=v_p(K)\), then for
every nonzero Laurent index \(j\), integrality of \(R_K\) and
\(j^2[q^j]R_K\in p^{2r}\mathbb Z_p\) imply
\(j[q^j]R_K\in p^r\mathbb Z_p\).  Thus
\(\partial R_K\in p^r\mathbb Z_p((q))\), and the same holds
for \(\theta R_K\).  Since
\(C_K=u^K R_{2K}/R_K\), equation
\eqref{P:eq:octa4-ratio-derivative} follows.  In particular, for
\(K=25\), if \(C_K=\sum c_j u^j\), then \(j c_j\in25\mathbb Z_5\).

\subsection{Constructing the lift modulo \(625\)}

From now on set \(p=5\), \(N=25\), \(X=u^5\), and write
\(\sigma(u)=X S(u)\), where \(S\in1+5\mathbb Z_5[[u]]\).
Apply \eqref{P:eq:octa4-vertex} at \(K=N,2N\) and eliminate
\(\lambda\).  Because all \(A_K\) have constant term one, this gives
\begin{equation}\label{P:eq:octa4-lift-equation}
 S^N C_{5N}(u)-C_N(XS)\in5^6\mathbb Z_5[[u]].
\end{equation}
Let \(h=S-1\).  Divide the left-hand side by \(N=25\) and
expand at \(h=0\).  Its coefficients through degree three will
be denoted \(a,b,c,d\), respectively.  Then
\[
 a\in5\mathbb Z_5[[u]],\qquad b(0)=1,
 \qquad b,c,d\in\mathbb Z_5[[u]].
\]
For completeness, the higher terms may be discarded modulo
\(625\) on the domain \(h\in5\mathbb Z_5[[u]]\).
Indeed, \(j c_j\in N\mathbb Z_5\) implies
\[
 \binom jr c_j=\frac{j c_j}{r}\binom{j-1}{r-1}
       \in \frac Nr\mathbb Z_5\quad(r\geq1),
\]
and the same estimate holds for \(\binom Nr\).
After division by \(N\), the term of degree \(r\) in \(h\)
therefore has valuation at least \(r-v_5(r)\).  This is at least
four whenever \(r\geq4\).  The assertion \(a\in5\mathbb Z_5[[u]]\)
also follows from \eqref{P:eq:octa4-lift-equation} and these estimates.
Thus
\[
 a+bh+ch^2+dh^3\equiv0\pmod{625}.
\]
As \(b\) is a unit, the unique root with \(h\equiv0\pmod5\) is
\begin{equation}\label{P:eq:octa4-rational-h}
 h\equiv-\frac ab-\frac{ca^2}{b^3}
       +\left(\frac d{b^4}-\frac{2c^2}{b^5}\right)a^3
       \pmod{625}.
\end{equation}
This assertion can also be checked by substitution; uniqueness
follows by factoring the difference of two roots, whose quotient
has unit constant term modulo five.

We now describe the exact polynomial calculation of \(a,b,c,d\).
All input polynomials are computed modulo \(5^6=15625\).
For \(0\leq r\leq3\), put
\[
 a_r=[h^r]A_N(X(1+h)),\qquad
 b_r=[h^r]A_{2N}(X(1+h)).
\]
Define
\[
 T_0=b_0,\qquad
 T_r=b_r a_0^r-\sum_{i=1}^r a_iT_{r-i}a_0^{i-1}
       \quad(1\leq r\leq3).
\]
Then \([h^r]C_N(X(1+h))=T_r/a_0^{r+1}\).
Let \(\mathcal A_r\in(\mathbb Z/625\mathbb Z)[u]\) be the
coefficientwise quotient
\begin{equation}\label{P:eq:octa4-ABCD}
 \mathcal A_r=
 \frac{\binom Nr A_{10N}a_0^4
       -A_{5N}T_r a_0^{3-r}}{N}\pmod{625}.
\end{equation}
The numerators are divisible by \(N\), as follows from the
preceding argument and is also checked coefficientwise in the
certificate.  All four coefficients \(a,b,c,d\) have the common
unit denominator \(A_{5N}a_0^4\); their respective numerators
are
\[
 A=\mathcal A_0,\quad B=\mathcal A_1,\quad
 C=\mathcal A_2,\quad D=\mathcal A_3.
\]
Thus \eqref{P:eq:octa4-rational-h} becomes
\[
 h\equiv-\frac AB-\frac{CA^2}{B^3}
       +\left(\frac D{B^4}-\frac{2C^2}{B^5}\right)A^3
       \pmod{625}.
\]
The calculation checks \(A\in5\mathbb Z[u]/625\), \(A(0)=0\),
and \(B(0)=1\).  The four polynomial degrees after reduction
are \(489,537,537,527\), respectively.

\subsection{The ordinary truncation certificate}

The first relation \eqref{P:eq:octa4-vertex} gives
\[
 \lambda\equiv\frac{S^{25}A_{125}(u)}{A_{25}(XS)}\pmod{625}.
\]
Because \(h\in5\mathbb Z_5[[u]]\),
\(S^{25}\equiv1+25h\pmod{625}\).
Define the four explicit integral polynomials
\[
 L_r=A_{125}(u)\bigl([h^r]F_5(X(1+h))
                 +25[h^{r-1}]F_5(X(1+h))\bigr)
        -F_{25}(u)[h^r]A_{25}(X(1+h)),
\]
where a coefficient with index \(-1\) means zero.
The target congruence is therefore equivalent to
\(L_0+L_1h+L_2h^2+L_3h^3\equiv0\pmod{625}\).
Substituting the independently constructed root and clearing
the unit denominator \(B^5\) reduces it to the finite identity
\begin{align}
0\equiv{}&L_0B^5-L_1AB^4
 +(L_2B^3-L_1CB^2)A^2\notag\\
 &+(L_1DB-2L_1C^2+2L_2CB-L_3B^2)A^3
 \pmod{625}.\label{P:eq:octa4-final-certificate}
\end{align}

The accompanying standalone program
\texttt{certificate\_octa4\_p5\_s2.py} computes the four vertex
polynomials directly from their factorial sums, constructs
\eqref{P:eq:octa4-ABCD}, and verifies that every coefficient of
\eqref{P:eq:octa4-final-certificate} is zero modulo \(625\).
The degree bound of this cleared identity is \(2819\).
All operations are integer polynomial addition, multiplication,
and exact division by \(25\); there is no power-series cutoff,
numerical approximation, or reconstruction of the lift from
the target truncations.  The machine-readable coefficient
certificate is \texttt{certificate\_octa4\_p5\_s2.json}.
This finite identity proves \(\lambda F_5(\sigma)-F_{25}\in
625\mathbb Z_5[[u]]\).  Multiplication by the period unit
\(F(\sigma)\) proves Theorem~\ref{P:thm:octa4-p5-s2}.




\clearpage
\section{The original modular verification program}\label{B:source-appendix-1}\label{B:app:script}
The following is the complete supplied verification program. Its comments and output labels retain the original manuscript's local section numbers and the former ``open case'' label; current theorem status is given in Part~\ref{part:positive}. The operations have not been altered in consolidation. The finite and floating-point checks have the limited roles described in Section~\ref{B:sec:verif}.

\begin{lstlisting}[language=Python]
#!/usr/bin/env python3
"""
Exact-arithmetic checks for the paper
  "Hecke traces and the Beukers--Vlasenko supercongruences for excellent Frobenius lifts".
Standard library only.  Run:  python3 verify_excellent_lift.py
Every check prints PASS/FAIL; the script exits with status 1 if any check fails.

Notation: q-series are integer lists [c_0, c_1, ...] truncated at a length M.
For a family (Theta, t, F) with F(t(q)) = Theta(q) and a prime p:
   Q_P(q)   = Theta(q) P(t(q^p)) / Theta(q^p)
   G_P(z)   = Theta(z^p) P(t(z)) / Theta(z)
   Tr(Q_P)  = Q_P + p^(w+1) chi(p) U_p(G_P)          (Lemma 3.3)
"""
import sys, random, cmath
from math import comb, factorial, gcd, floor, ceil
from fractions import Fraction as Fr

FAILS = []
def report(name, ok, detail=""):
    print(("PASS " if ok else "FAIL ") + name + ((" :: " + detail) if detail else ""), flush=True)
    if not ok: FAILS.append(name)

# ---------------------------------------------------------------- series utilities
def mul(a, b, M):
    c = [0]*M
    for i in range(min(M, len(a))):
        x = a[i]
        if x == 0: continue
        for j in range(min(M-i, len(b))):
            y = b[j]
            if y: c[i+j] += x*y
    return c
def inv(a, M):
    assert a[0] == 1
    b = [0]*M; b[0] = 1
    for n in range(1, M):
        s = 0
        for k in range(1, min(n, len(a)-1)+1):
            if a[k]: s += a[k]*b[n-k]
        b[n] = -s
    return b
def subqp(a, p, M):                       # a(q) -> a(q^p)
    c = [0]*M
    for i, x in enumerate(a):
        if i*p >= M: break
        c[i*p] = x
    return c
def polyt(coeffs, t, M):                  # sum coeffs[k] t^k  (Horner)
    res = [0]*M
    for cf in reversed(coeffs):
        res = mul(res, t, M); res[0] += cf
    return res
def t_expand(f, t, M):                    # f(q) = sum_k c_k t(q)^k ; t = q + O(q^2)
    rem = list(f[:M]); out = []; tp = [1] + [0]*(M-1)
    for k in range(M):
        c = rem[k]; out.append(c)
        if c:
            for i in range(M): rem[i] -= c*tp[i]
        tp = mul(tp, t, M)
    return out
def vp(x, p):
    if x == 0: return 10**9
    v = 0
    while x % p == 0: x //= p; v += 1
    return v
def vp_series(a, p): return min(vp(x, p) for x in a)

def etaseries(exps, M):
    """prod_m eta(m tau)^{e_m} as an integer q-series (integral q-shift required)."""
    sh = Fr(sum(m*e for m, e in exps.items()), 24)
    assert sh.denominator == 1 and sh >= 0
    P = [0]*M; P[0] = 1
    for m, e in exps.items():
        for n in range(1, M):
            k = m*n
            if k >= M: break
            for _ in range(abs(e)):
                if e > 0:
                    for i in range(M-1, k-1, -1): P[i] -= P[i-k]
                else:
                    for i in range(k, M): P[i] += P[i-k]
    sh = int(sh)
    return [0]*sh + P[:M-sh]
def cubic_theta(M):
    a = [0]*M; R = int((4*M/3)**0.5) + 3          # x^2+xy+y^2 < M forces |x|,|y| < sqrt(4M/3)
    for x in range(-R, R+1):
        for y in range(-R, R+1):
            v = x*x + x*y + y*y
            if v < M: a[v] += 1
    return a
def sigma(n): return sum(d for d in range(1, n+1) if n % d == 0)

# ---------------------------------------------------------------- characters
def chi_m4(p): return 0 if p % 2 == 0 else (1 if p % 4 == 1 else -1)
def chi_m3(p): return 0 if p % 3 == 0 else (1 if p % 3 == 1 else -1)
def chi_5(p):  return 0 if p % 5 == 0 else (1 if p % 5 in (1, 4) else -1)
def chi_1(p):  return 1

def C(n, k): return comb(n, k) if 0 <= k <= n else 0

# ---------------------------------------------------------------- the families of cusp type
# (name, weight w, level N, group, character chi, Theta eta-exponents, t eta-exponents,
#  zero cusp denominators (d|N) of Theta, r_Theta, closed formula for a_n)
CUSP = [
 ("Legendre 2F1(1/2,1/2;1|16t)",1,4,"G0(4)",chi_m4,{2:10,1:-4,4:-4},{1:8,4:16,2:-24},[2],Fr(1,2),
   lambda n: C(2*n,n)**2),
 ("Zagier A (Franel)",1,6,"G0(6)",chi_m3,{2:1,3:6,1:-2,6:-3},{1:3,6:9,2:-3,3:-9},[3],Fr(1),
   lambda n: sum(C(n,k)**3 for k in range(n+1))),
 ("Zagier B",1,9,"G0(9)",chi_m3,{1:3,3:-1},{9:3,1:-3},[1],Fr(1),
   lambda n: sum((-3)**(n-3*k)*factorial(n)//(factorial(n-3*k)*factorial(k)**3) for k in range(n//3+1))),
 ("Zagier C",1,6,"G0(6)",chi_m3,{3:1,2:6,1:-3,6:-2},{6:8,1:4,2:-8,3:-4},[2],Fr(1),
   lambda n: sum(C(n,k)**2*C(2*k,k) for k in range(n+1))),
 ("Zagier E",1,8,"G0(8)",chi_m4,{2:10,1:-4,4:-4},{1:4,4:2,8:4,2:-10},[2],Fr(1),
   lambda n: sum(4**(n-2*k)*C(n,2*k)*C(2*k,k)**2 for k in range(n//2+1))),
 ("Zagier F",1,6,"G0(6)",chi_m3,{1:6,6:1,2:-3,3:-2},{6:5,2:1,1:-5,3:-1},[1],Fr(1),
   lambda n: sum((-8)**(n-k)*C(n,k)*sum(C(k,j)**3 for j in range(k+1)) for k in range(n+1))),
 ("3F2(1/2,1/2,1/2;1,1|64t)",2,4,"G0(4)+4",chi_1,{2:20,1:-8,4:-8},{1:24,4:24,2:-48},[2],Fr(1,2),
   lambda n: C(2*n,n)**3),
 ("Apery (gamma)",2,6,"G0(6)+6",chi_1,{2:7,3:7,1:-5,6:-5},{1:12,6:12,2:-12,3:-12},[2,3],Fr(1),
   lambda n: sum(C(n,k)**2*C(n+k,k)**2 for k in range(n+1))),
 ("Almkvist (eta)",2,5,"G0(5)",chi_5,{1:5,5:-1},{5:6,1:-6},[1],Fr(1),
   lambda n: (-1)**n*sum((-1)**k*C(n,k)**3*(C(4*n-5*k-1,3*n)+C(4*n-5*k,3*n)) for k in range(n+1))),
 ("Domb (alpha)",2,6,"G0(6)+3",chi_1,{1:4,3:4,2:-2,6:-2},{2:6,6:6,1:-6,3:-6},[1,3],Fr(1),
   lambda n: (-1)**n*sum(C(n,k)**2*C(2*k,k)*C(2*n-2*k,n-k) for k in range(n+1))),
 ("Almkvist-Zudilin (delta)",2,6,"G0(6)+2",chi_1,{1:3,2:3,3:-1,6:-1},{3:4,6:4,1:-4,2:-4},[1,2],Fr(1),
   lambda n: sum((-3)**(n-3*k)*factorial(n+k)//(factorial(k)**4*factorial(n-3*k)) for k in range(n//3+1))),
 ("Almkvist (epsilon)",2,8,"G0(8)+8",chi_1,{2:6,4:6,1:-4,8:-4},{1:8,8:8,2:-8,4:-8},[2,4],Fr(1),
   lambda n: sum(C(n,k)**2*C(2*k,n)**2 for k in range(n+1))),
 ("Almkvist (zeta)",2,9,"G0(9)+9",chi_1,{3:10,1:-3,9:-3},{1:6,9:6,3:-12},[3],Fr(1),
   lambda n: sum(C(n,k)**2*C(n,l)*C(k,l)*C(k+l,n) for k in range(n+1) for l in range(n+1))),
]

def fam_series(f, M):
    name,w,N,grp,chi,thx,tx,zd,rth,an = f
    return etaseries(thx, M), etaseries(tx, M)

# ---------------------------------------------------------------- Check 1: parametrisations
def check_param():
    M = 60
    for f in CUSP:
        th, t = fam_series(f, M)
        A = [f[9](n) for n in range(M)]
        ok = (polyt(A, t, M) == th) and t[0] == 0 and t[1] == 1 and th[0] == 1
        report(f"[1] F(t(q)) = Theta(q) to O(q^{M}) for {f[0]}", ok)

# ---------------------------------------------------------------- Check 2: Ligozat orders at cusps
def ligozat(exps, N, d):
    # order of vanishing of prod eta(delta tau)^{r_delta} at a cusp a/d (d|N) of Gamma_0(N),
    # in the local uniformiser of X_0(N) at that cusp
    return Fr(N, 24)*sum(Fr(gcd(d, de)**2*r, gcd(d, N//d)*d*de) for de, r in exps.items())
def check_cusps():
    for f in CUSP:
        name,w,N,grp,chi,thx,tx,zd,rth,an = f
        divs = [d for d in range(1, N+1) if N % d == 0]
        oth = {d: ligozat(thx, N, d) for d in divs}
        ot  = {d: ligozat(tx, N, d) for d in divs}
        zeros = [d for d in divs if oth[d] > 0]
        poles = [d for d in divs if ot[d] < 0]
        ratios = {d: oth[d]/(-ot[d]) for d in zeros}
        weight_ok = Fr(sum(thx.values()), 2) == w and sum(tx.values()) == 0
        ok = weight_ok and zeros == zd and poles == zd and all(oth[d] >= 0 for d in divs) \
             and all(r == rth for r in ratios.values())
        # cusp-class condition (C): gcd(e, N/e) in {1,2}, or (N=9,e=3: the two cusps 1/3, 2/3 are
        # interchanged by w_9 and both have width 1)
        cond = all(gcd(e, N//e) in (1, 2) for e in zd) or (N == 9 and zd == [3])
        # character of the eta-quotient (Gordon-Hughes-Newman): chi(d) = ((-1)^w s / d), s = prod delta^r_delta
        s = Fr(1)
        for de, r in thx.items(): s *= Fr(de)**r
        D = (-1)**w * s.numerator * s.denominator          # same square class as (-1)^w s
        def legendre(a, p):
            a %= p
            return 0 if a == 0 else (1 if pow(a, (p-1)//2, p) == 1 else -1)
        char_ok = all(legendre(D, p) == chi(p) for p in [5, 7, 11, 13, 17, 19, 23, 29, 31, 37] if N % p and D % p)
        tchar = Fr(1)
        for de, r in tx.items(): tchar *= Fr(de)**r
        tchar_ok = all(legendre(tchar.numerator*tchar.denominator, p) == 1
                       for p in [5, 7, 11, 13, 17, 19, 23, 29, 31, 37] if N % p)
        # Atkin-Lehner symmetry: W_e maps eta_delta to eta_{e delta / gcd(e,delta)^2}
        al_ok = True
        if "+" in grp:
            e = int(grp.split("+")[1])
            for ex in (thx, tx):
                img = {}
                for de, r in ex.items():
                    d2 = e*de//gcd(e, de)**2
                    img[d2] = img.get(d2, 0) + r
                al_ok &= (img == ex)
        report(f"[2] zeros of Theta = poles of t = cusps with denominators {zd}, r_Theta = {rth}, (C) holds, "
               f"character and Atkin-Lehner symmetry ({grp}) ok: {name}",
               ok and cond and char_ok and tchar_ok and al_ok, f"ord Theta {dict(oth)}, ord t {dict(ot)}")

# ---------------------------------------------------------------- the trace
def trace_poly(f, p, P, M):
    """t-expansion (length M) of Tr(Q_P) = Q_P + p^{w+1} chi(p) U_p(G_P)."""
    name,w,N,grp,chi,thx,tx,zd,rth,an = f
    MZ = p*M
    th, t = fam_series(f, MZ)
    Pz = polyt(P, t, MZ)
    G = mul(mul(subqp(th[:M], p, MZ), inv(th, MZ), MZ), Pz, MZ)
    UpG = [G[p*n] for n in range(M)]
    thq, tq = th[:M], t[:M]
    Q = mul(mul(thq, inv(subqp(thq, p, M), M), M), subqp(Pz[:M], p, M), M)
    Tr = [Q[n] + p**(w+1)*chi(p)*UpG[n] for n in range(M)]
    return t_expand(Tr, tq, M)

def check_trace_lemma():
    rnd = random.Random(20261107)
    for f in CUSP:
        name,w,N,grp,chi,thx,tx,zd,rth,an = f
        A = [an(n) for n in range(200)]
        allok = True; info = []
        for p in [2, 3, 5, 7, 11]:
            if N % p == 0: continue
            cases = []
            for (s, m) in [(1, 1), (1, 2), (2, 1)]:
                if m*p**s <= 60: cases.append(("F_{%d}" % (m*p**(s-1)), A[:m*p**(s-1)]))
            for L in [3, 6]:
                cases.append(("random deg %d" % (L-1), [1] + [rnd.randint(-9, 9) for _ in range(L-1)]))
            for lab, P in cases:
                L = len(P)
                D = p*(L-1) + int(rth*(p-1))
                M = D + 20
                T = trace_poly(f, p, P, M)
                exact = all(x == 0 for x in T[D+1:])
                allok &= exact
                info.append(f"p={p},{lab}:deg<={D}:{'ok' if exact else 'NO'}")
        report(f"[3] Tr(Q_P) is a polynomial in t of degree <= p(L-1)+floor(r(p-1)) (exact, to t^(D+19)): {name}",
               allok, "; ".join(info[:6]) + (" ..." if len(info) > 6 else ""))

# ---------------------------------------------------------------- Check 4: main theorem
def cong_val(f, p, s, m, M):
    name,w,N,grp,chi,thx,tx,zd,rth,an = f
    th, t = fam_series(f, M)
    K = m*p**s
    A = [an(n) for n in range(K)]
    lhs = mul(th, polyt(A[:m*p**(s-1)], subqp(t, p, M), M), M)        # Theta(q) F_{K_{s-1}}(t(q^p))
    rhs = mul(subqp(th, p, M), polyt(A, t, M), M)                    # Theta(q^p) F_{K_s}(t(q))
    return vp_series([x-y for x, y in zip(lhs, rhs)], p)

def check_main():
    for f in CUSP:
        name,w,N,grp,chi,thx,tx,zd,rth,an = f
        allok = True; vals = []
        for p in [2, 3, 5, 7, 11, 13]:
            if N % p == 0: continue
            for (s, m) in [(1, 1), (1, 2), (2, 1), (2, 2), (3, 1)]:
                K = m*p**s
                if K > 130: continue
                M = min(2*K + 10, 200)
                v = cong_val(f, p, s, m, M)
                allok &= (v >= (w+1)*s)
                vals.append(f"p={p},s={s},m={m}:v={v}")
        report(f"[4] Theta(q)F_(mp^(s-1))(t(q^p)) = Theta(q^p)F_(mp^s)(t(q)) mod p^((w+1)s), w={w}: {name}",
               allok, " ".join(vals))

# ---------------------------------------------------------------- Check 5: coefficient corollary
def check_coeff():
    for f in CUSP:
        name,w,N,grp,chi,thx,tx,zd,rth,an = f
        ok = True; cnt = 0
        for p in [2, 3, 5, 7, 11, 13]:
            if N % p == 0: continue
            for s in [1, 2, 3]:
                for m in range(1, 6):
                    if m*p**s > 250: continue
                    cnt += 1
                    ok &= (an(m*p**s) - an(m*p**(s-1))) % p**((w+1)*s) == 0
        report(f"[5] a(mp^s) = a(mp^(s-1)) mod p^((w+1)s) ({cnt} cases): {name}", ok)

# ---------------------------------------------------------------- Check 6: BV3 Conjecture 7.5 in its own normalisation
def frob_canonical(rec, order_gap, M):
    """F_0(t) and canonical coordinate q(t) = t exp(F_1/F_0) from a recurrence
       n^k c_n = sum_j R_j(n) c_{n-j}, via the Frobenius method n -> n+eps (dual numbers)."""
    # rec(nn) returns (lead(nn), [(j, coeff(nn))...]) for nn a dual number (x, dx)
    def dm(x, y): return (x[0]*y[0], x[0]*y[1] + x[1]*y[0])
    cs = [(Fr(1), Fr(0))]
    for n in range(1, M):
        lead, terms = rec((Fr(n), Fr(1)))
        s = (Fr(0), Fr(0))
        for j, cf in terms:
            if n-j >= 0:
                pr = dm(cf, cs[n-j]); s = (s[0] + pr[0], s[1] + pr[1])
        if lead[0] == 0:
            raise ValueError
        cs.append((s[0]/lead[0], (s[1]*lead[0] - s[0]*lead[1])/(lead[0]**2)))
    F0 = [x[0] for x in cs]; F1 = [x[1] for x in cs]
    # g = F1/F0, E = exp(g)
    inv0 = [Fr(0)]*M; inv0[0] = Fr(1)
    for n in range(1, M): inv0[n] = -sum(F0[k]*inv0[n-k] for k in range(1, n+1))
    g = [sum(F1[i]*inv0[n-i] for i in range(n+1)) for n in range(M)]
    E = [Fr(0)]*M; E[0] = Fr(1)
    for n in range(1, M): E[n] = sum(k*g[k]*E[n-k] for k in range(1, n+1))/n
    q = [Fr(0)] + E[:M-1]
    return F0, q
def series_inverse_fn(q, M):
    """compositional inverse t(Q) of q(t) = t + ..."""
    T = [Fr(0)]*M; T[1] = Fr(1)
    for n in range(2, M):
        comp = [Fr(0)]*(n+1); Tp = [Fr(1)] + [Fr(0)]*n
        for k in range(1, n+1):
            Tp = [sum(Tp[i]*T[j-i] for i in range(j+1)) for j in range(n+1)]
            if q[k]:
                for j in range(n+1): comp[j] += q[k]*Tp[j]
        T[n] = -comp[n]
    return T
def fmul(a, b, M):
    c = [Fr(0)]*M
    for i in range(min(M, len(a))):
        if a[i] == 0: continue
        for j in range(min(M-i, len(b))):
            if b[j]: c[i+j] += a[i]*b[j]
    return c
def fcompose(a, b, M):          # a(b(t)), b(0)=0
    res = [Fr(0)]*M
    for cf in reversed(a[:M]):
        res = fmul(res, b, M); res[0] += cf
    return res
def finv(a, M):
    b = [Fr(0)]*M; b[0] = 1/a[0]
    for n in range(1, M): b[n] = -sum(a[k]*b[n-k] for k in range(1, n+1))/a[0]
    return b
def vpF(x, p):
    if x == 0: return 10**9
    return vp(x.numerator, p) - vp(x.denominator, p)
def bv_check(label, rec, gap, mBV, plist, slist, M, expect, tag="[6]"):
    F0, q = frob_canonical(rec, gap, M)
    T = series_inverse_fn(q, M)                         # mirror map t(Q)
    for p in plist:
        qp = [Fr(0)]*M
        qpow = [Fr(1)] + [Fr(0)]*(M-1)
        for _ in range(p): qpow = fmul(qpow, q, M)        # q(t)^p
        tsig = fcompose(T, qpow, M)                       # t^sigma(t) = t(q(t)^p)
        for s in slist:
            K = mBV*p**s
            if K >= M: continue
            Ft = F0; Fs = fcompose(F0, tsig, M)
            A = fcompose(F0[:mBV*p**(s-1)], tsig, M)
            lhs = fmul(Ft, A, M); rhs = fmul(Fs, F0[:K] + [Fr(0)]*(M-K), M)
            v = min(vpF(x-y, p) for x, y in zip(lhs, rhs))
            ok = v >= expect(s)
            report(f"{tag} BV3 (19) with m={mBV}, excellent lift q->q^p: {label}, p={p}, s={s}", ok,
                   f"v_p(F(t)F_(mp^(s-1))(t^s) - F(t^s)F_(mp^s)(t)) = {v} (required >= {expect(s)})")
def check_bv3():
    lin = lambda a0, a1, n: (a0 + a1*n[0], a1*n[1])
    def dmul(*xs):
        r = (Fr(1), Fr(0))
        for x in xs: r = (r[0]*x[0], r[0]*x[1] + r[1]*x[0])
        return r
    sc = lambda c, x: (c*x[0], c*x[1])
    # hypercubic = hyperoctahedral in dimension 2:  n^2 c_n = 16 (n-1)^2 c_{n-2}
    rec2 = lambda n: (dmul(n, n), [(2, sc(16, dmul(lin(-1, 1, n), lin(-1, 1, n))))])
    bv_check("hypercubic/hyperoctahedral n=2 (c_{2k}=C(2k,k)^2)", rec2, 2, 2, [3, 5, 7], [1, 2], 110, lambda s: 2*s)
    # hypercubic n=3:  n^3 c_n = 64 (n-1)^3 c_{n-2}
    rec3 = lambda n: (dmul(n, n, n), [(2, sc(64, dmul(lin(-1, 1, n), lin(-1, 1, n), lin(-1, 1, n))))])
    bv_check("hypercubic n=3 (c_{2k}=C(2k,k)^3)", rec3, 2, 2, [5, 7], [1, 2], 110, lambda s: 3*s)

# ---------------------------------------------------------------- Check 7: families outside the theorem
def elliptic_families(M):
    out = []
    a = cubic_theta(M)
    t27 = mul(etaseries({3:9, 1:-3}, M), inv(mul(mul(a, a, M), a, M), M), M)
    out.append(("2F1(1/3,2/3;1|27t)  [zero of Theta at elliptic point of G0(3)]", 1, 3, chi_m3, a, t27))
    h2 = etaseries({2:24, 1:-24}, M); d = [1] + [0]*(M-1)
    d = [d[i] + 64*h2[i] for i in range(M)]
    t256 = mul(h2, inv(mul(d, d, M), M), M)
    E22 = [1] + [24*(sigma(n) - (2*sigma(n//2) if n % 2 == 0 else 0)) for n in range(1, M)]
    out.append(("3F2(1/4,1/2,3/4;1,1|256t)  [elliptic, G0(2)+2]", 2, 2, chi_1, E22, t256))
    h3 = etaseries({3:12, 1:-12}, M); d = [1] + [0]*(M-1)
    d = [d[i] + 27*h3[i] for i in range(M)]
    t108 = mul(h3, inv(mul(d, d, M), M), M)
    E23 = [1] + [12*(sigma(n) - (3*sigma(n//3) if n % 3 == 0 else 0)) for n in range(1, M)]
    out.append(("3F2(1/3,1/2,2/3;1,1|108t)  [elliptic, G0(3)+3]", 2, 3, chi_1, E23, t108))
    return out
def check_ell_nonpoly():
    name, w, N, chi, th0, t0 = elliptic_families(400)[0]
    p, P = 5, [1]
    M = 60; MZ = p*M
    a = cubic_theta(MZ)
    t = mul(etaseries({3:9, 1:-3}, MZ), inv(mul(mul(a, a, MZ), a, MZ), MZ), MZ)
    G = mul(subqp(a[:M], p, MZ), inv(a, MZ), MZ)
    Q = mul(a[:M], inv(subqp(a[:M], p, M), M), M)
    Tr = [Q[n] + p**2*chi(p)*G[p*n] for n in range(M)]
    T = t_expand(Tr, t[:M], M)
    hi = [(n, vp(T[n], p)) for n in range(5, 40, 6)]
    report("[8h] for 2F1(1/3,2/3;1|27t), p=5: Tr(Q_1) is not a polynomial in t; "
           "its high t-coefficients are nonzero but p-adically small", any(T[n] != 0 for n in range(10, M-5)),
           "(n, v_5 of coefficient): " + str(hi))


def check_hyperoct3():
    # [7c] hyperoctahedral family in dimension 3 (BV3 Conj. 7.5, outside Thm 1.1):
    #      c_{2k} = C(2k,k) sum_j C(k,j)^2 C(2j,j); in u = t^2 the recurrence
    #      k^3 a_k = 2(2k-1)(10k^2-10k+3) a_{k-1} - 36(k-1)(2k-1)(2k-3) a_{k-2}
    lin = lambda a0, a1, n: (a0 + a1*n[0], a1*n[1])
    def dmul(*xs):
        r = (Fr(1), Fr(0))
        for x in xs: r = (r[0]*x[0], r[0]*x[1] + r[1]*x[0])
        return r
    sc = lambda c, x: (c*x[0], c*x[1])
    quad = lambda n: (10*n[0]**2 - 10*n[0] + 3, (20*n[0] - 10)*n[1])
    rec = lambda n: (dmul(n, n, n), [(1, sc(2, dmul(lin(-1, 2, n), quad(n)))),
                                     (2, sc(-36, dmul(lin(-1, 1, n), lin(-1, 2, n), lin(-3, 2, n))))])
    F0, q = frob_canonical(rec, 1, 70)
    ok_seq = all(F0[k] == C(2*k, k)*sum(C(k, j)**2*C(2*j, j) for j in range(k+1)) for k in range(70))
    T = series_inverse_fn(q, 70)
    integral = all(x.denominator == 1 for x in T)
    M = 60
    uq = [int(x) for x in T[:M]]
    A = [int(x) for x in F0[:M]]
    th = polyt(A, uq, M)
    vals = []
    ok = ok_seq and integral
    for (p, s) in [(5, 1), (7, 1), (11, 1), (5, 2)]:
        K = p**s
        lhs = mul(th, polyt(A[:p**(s-1)], subqp(uq, p, M), M), M)
        rhs = mul(subqp(th, p, M), polyt(A[:K], uq, M), M)
        v = vp_series([x-y for x, y in zip(lhs, rhs)], p)
        ok &= (v >= 2*s)
        vals.append(f"p={p},s={s}:v={v}")
    report("[7a] (numerical, open case) hyperoctahedral n=3 of BV3 Conj. 7.5: congruence mod p^(2s) "
           "(the valuation is exactly 2s, not 3s)", ok, " ".join(vals))
    # [7b] modular structure: u = v/(1+v)^2, Theta = (1+v) Theta_Apery, v = (eta1 eta6/eta2 eta3)^12
    M = 60
    v = etaseries({1: 12, 6: 12, 2: -12, 3: -12}, M)
    opv = [(1 if i == 0 else 0) + v[i] for i in range(M)]
    u = mul(v, inv(mul(opv, opv, M), M), M)
    thh = mul(opv, etaseries({2: 7, 3: 7, 1: -5, 6: -5}, M), M)
    ok7b = (u == uq[:M]) and (t_expand(thh, u, M) == A[:M])
    report("[7b] hyperoctahedral n=3: mirror map u = v/(1+v)^2 and F(u(q)) = (1+v) eta2^7 eta3^7/(eta1^5 eta6^5), "
           "v = (eta1 eta6/(eta2 eta3))^12", ok7b)
    # [7c] the truncation coefficients have only one carry per digit p-1: v_p(c_(p^k-1)) = k
    bz = lambda n: sum(C(n, j)**2*C(2*j, j) for j in range(n+1))
    ok7c = all(vp(bz(p**k - 1), p) == 0 and vp(C(2*(p**k-1), p**k-1)*bz(p**k-1), p) == k
               for p in [5, 7, 11, 13] for k in [1, 2])
    report("[7c] hyperoctahedral n=3: v_p(c_(p^k-1)) = k for p in {5,7,11,13}, k in {1,2} "
           "(only one factor p per digit p-1, so the carry bound of Section 6 fails)", ok7c)


# ---------------------------------------------------------------- Check 8: families with an elliptic zero

def normseries(f, p, M):
    """prod_{k mod p} f(zeta^k z) as a series in q = z^p (f in 1 + zZ[[z]] given to precision p*M)."""
    MZ = p*M
    F = [Fr(x) for x in f[:MZ]]
    finv = [Fr(0)]*MZ; finv[0] = Fr(1)
    for n in range(1, MZ):
        finv[n] = -sum(F[k]*finv[n-k] for k in range(1, n+1) if F[k])
    zl = [sum(i*F[i]*finv[n-i] for i in range(1, n+1) if F[i]) for n in range(MZ)]   # z f'/f
    g = [Fr(0)] + [p*zl[p*n]/(p*n) for n in range(1, M)]                           # log N
    E = [Fr(0)]*M; E[0] = Fr(1)
    for n in range(1, M): E[n] = sum(k*g[k]*E[n-k] for k in range(1, n+1))/n
    assert all(x.denominator == 1 for x in E)
    return [int(x) for x in E]

def ell_data(name, M):
    """(w, chi, Theta, t, vartheta, k0, ord, lam, coefficient formula) for the three families."""
    if name == "27t":
        a = cubic_theta(M)
        t = mul(etaseries({3: 9, 1: -3}, M), inv(mul(mul(a, a, M), a, M), M), M)
        return dict(w=1, chi=chi_m3, th=a, t=t, vth=a, k0=3, ordz=1, lam=1,
                    an=lambda n: factorial(3*n)//factorial(n)**3, N=3)
    if name == "256t":
        h2 = etaseries({2: 24, 1: -24}, M); d = [1] + [0]*(M-1)
        d = [d[i] + 64*h2[i] for i in range(M)]
        t = mul(h2, inv(mul(d, d, M), M), M)
        E = [1] + [24*(sigma(n) - (2*sigma(n//2) if n % 2 == 0 else 0)) for n in range(1, M)]
        return dict(w=2, chi=chi_1, th=E, t=t, vth=E, k0=4, ordz=1, lam=1,
                    an=lambda n: factorial(4*n)//factorial(n)**4, N=2)
    if name == "108t":
        h3 = etaseries({3: 12, 1: -12}, M); d = [1] + [0]*(M-1)
        d = [d[i] + 27*h3[i] for i in range(M)]
        t = mul(h3, inv(mul(d, d, M), M), M)
        E = [1] + [12*(sigma(n) - (3*sigma(n//3) if n % 3 == 0 else 0)) for n in range(1, M)]
        a = cubic_theta(M)
        return dict(w=2, chi=chi_1, th=E, t=t, vth=a, k0=6, ordz=2, lam=1, vchi=chi_m3,
                    an=lambda n: C(2*n, n)*factorial(3*n)//factorial(n)**3, N=3)
    raise ValueError

ELL = ["27t", "256t", "108t"]

def check_ell_param():
    M = 60
    for nm in ELL:
        D = ell_data(nm, M)
        F = t_expand(D['th'], D['t'], M)
        ok = all(F[n] == D['an'](n) for n in range(M))
        report(f"[8a] F(t(q)) = Theta(q) to O(q^{M}) for {nm}", ok)
    # eta identities giving the zeros of Theta and the pole of t
    a = cubic_theta(M)
    b3 = etaseries({1: 9, 3: -3}, M); c3 = [27*x for x in etaseries({3: 9, 1: -3}, M)]
    D = ell_data("27t", M)
    ok1 = mul(mul(a, a, M), a, M) == [x + y for x, y in zip(b3, c3)]
    ok2 = mul(mul(mul(a, a, M), a, M), D['t'], M) == etaseries({3: 9, 1: -3}, M)
    report("[8a] 27t: a^3 = eta1^9/eta3^3 + 27 eta3^9/eta1^3 and a^3 t = eta3^9/eta1^3 (no zeros in H)", ok1 and ok2)
    D = ell_data("256t", M)
    E2 = mul(D['th'], D['th'], M)
    ok1 = E2 == [x + 64*y for x, y in zip(etaseries({1: 16, 2: -8}, M), etaseries({2: 16, 1: -8}, M))]
    ok2 = mul(mul(E2, E2, M), D['t'], M) == etaseries({1: 8, 2: 8}, M)
    report("[8a] 256t: E22^2 = eta1^16/eta2^8 + 64 eta2^16/eta1^8 and E22^4 t = eta1^8 eta2^8", ok1 and ok2)
    D = ell_data("108t", M)
    ok1 = D['th'] == mul(a, a, M)
    ok2 = mul(mul(mul(D['th'], D['th'], M), D['th'], M), D['t'], M) == etaseries({1: 6, 3: 6}, M)
    report("[8a] 108t: E23 = a^2 and E23^3 t = eta1^6 eta3^6", ok1 and ok2)
    # floating point: Theta vanishes at the elliptic point e
    def qs(c, tau, n=200):
        q = cmath.exp(2j*cmath.pi*tau)
        return sum(c[k]*q**k for k in range(min(n, len(c))))
    e3 = complex(0.5, 1/(2*3**0.5)); e4 = complex(0.5, 0.5)
    vals = []
    for nm, e in [("27t", e3), ("256t", e4), ("108t", e3)]:
        D = ell_data(nm, 200)
        vals.append(abs(qs(D['th'], e)))
    report("[8a] |Theta(e)| at e=(3+sqrt-3)/6, (1+i)/2, (3+sqrt-3)/6 (floating point)", max(vals) < 1e-12,
           " ".join(f"{v:.1e}" for v in vals))

def Dpoly(D, p, M):
    """the polynomial D(t) = chi'(p) vth(q^p) N(vth)(q) / vth(q)^(p+1), expanded in t."""
    vth = D['vth']; MZ = p*M
    Nm = normseries(vth[:MZ] if len(vth) >= MZ else vth, p, M)
    num = mul(subqp(vth[:M], p, M), Nm, M)
    den = [1] + [0]*(M-1)
    for _ in range(p+1): den = mul(den, vth[:M], M)
    c = D.get('vchi', D['chi'])(p)
    return t_expand([c*x for x in mul(num, inv(den, M), M)], D['t'][:M], M)

def deg_tail(f, upto):
    nz = [n for n in range(upto) if f[n] != 0]
    return max(nz) if nz else -1

def check_ell_D():
    for nm in ELL:
        details = []; ok = True
        for p in [5, 7, 11, 13]:
            M = 36
            D = ell_data(nm, p*M)
            Dt = Dpoly(D, p, M)
            dg = deg_tail(Dt, M - 2)
            eps = Dt[0]
            dstar = Fr((p+1)*D['lam'], D['k0'])
            good = eps in (1, -1) and all((Dt[n] - (eps if n == 0 else 0)) % p == 0 for n in range(M)) and dg <= dstar
            if nm == "108t":      # also B_Theta = D^2
                Bt = Dpoly(dict(D, vth=D['th'], vchi=D['chi']), p, M)
                good &= Bt[:M-2] == mul(Dt, Dt, M)[:M-2]
            ok &= good
            details.append(f"p={p}:deg={dg}<={dstar},D(0)={eps}")
        report(f"[8b] D(t) in Z[t], D = +-1 mod p, deg D <= (p+1)lam/k0{' and B = D^2' if nm=='108t' else ''}: {nm}",
               ok, " ".join(details))

def sym_functions(D, p, M):
    """q-series of e_k (k=0..p+1), n_l (l=0..p), S_0 for the pairs (x_beta, W_beta)."""
    w = D['w']; chi = D['chi'](p); th = D['th']; t = D['t']; MZ = p*M
    x0 = mul(th[:M], inv(subqp(th[:M], p, M), M), M)
    W0 = subqp(t[:M], p, M)
    base = mul(subqp(th[:M], p, MZ), inv(th[:MZ], MZ), MZ)
    tj = [1] + [0]*(MZ-1)
    Pw = [None]; X = []
    for j in range(p+1):
        Gj = mul(base, tj, MZ)
        X.append([p**(w+1)*chi*Gj[p*n] for n in range(M)])
        if j >= 1: Pw.append([p*tj[p*n] for n in range(M)])
        tj = mul(tj, t[:MZ], MZ)
    # elementary symmetric functions of the p non-principal W's (Newton)
    e = [[Fr(1)] + [Fr(0)]*(M-1)]
    for j in range(1, p+1):
        acc = [Fr(0)]*M
        for i in range(1, j+1):
            pr = mul(e[j-i], Pw[i], M)
            sg = 1 if i % 2 == 1 else -1
            acc = [a_ + sg*b_ for a_, b_ in zip(acc, pr)]
        e.append([x/j for x in acc])
    assert all(x.denominator == 1 for row in e for x in row)
    e = [[int(x) for x in row] for row in e] + [[0]*M]
    eall = [[e[k][n] + (mul(W0, e[k-1], M)[n] if k >= 1 else 0) for n in range(M)] for k in range(p+2)]
    Y = []
    for l in range(p+1):
        acc = [0]*M
        for j in range(l+1):
            pr = mul(e[l-j], X[j], M)
            acc = [a_ + (-1)**j*b_ for a_, b_ in zip(acc, pr)]
        Y.append(acc)
    n = []
    for l in range(p+1):
        v = mul(x0, e[l], M)
        v = [a_ + b_ for a_, b_ in zip(v, Y[l])]
        if l >= 1: v = [a_ + b_ for a_, b_ in zip(v, mul(W0, Y[l-1], M))]
        n.append(v)
    S0 = [a_ + b_ for a_, b_ in zip(x0, X[0])]
    return eall, n, S0

def check_ell_sym():
    for nm in ELL:
        ok = True; det = []
        for p in [5, 7]:
            M = p + 16
            D = ell_data(nm, p*M)
            k0, o, lam = D['k0'], D['ordz'], D['lam']
            mu, e0 = k0//lam, o//lam
            Dt = Dpoly(D, p, M)
            Dq = polyt(Dt[:deg_tail(Dt, M-2)+1], D['t'][:M], M)
            Dmu = [1] + [0]*(M-1)
            for _ in range(mu): Dmu = mul(Dmu, Dq, M)
            De0 = [1] + [0]*(M-1)
            for _ in range(e0): De0 = mul(De0, Dq, M)
            eall, nl, S0 = sym_functions(D, p, M)
            tq = D['t'][:M]
            Ek = [t_expand(mul(ek, Dmu, M), tq, M) for ek in eall]
            Ml = [t_expand(mul(x, Dmu, M), tq, M) for x in nl]
            A0 = t_expand(mul(S0, De0, M), tq, M)
            lim = M - 2
            g1 = all(deg_tail(E, lim) <= p+1 for E in Ek[1:])
            g2 = Ek[1][p+1] % p == 0
            g3 = all(deg_tail(x, lim) <= p for x in Ml)
            g4 = deg_tail(A0, lim) <= floor(Fr(p*o, k0))
            # mod p: E_1 = eps^mu t^p
            eps = Dt[0]
            g5 = all((Ek[1][n] - (eps**mu if n == p else 0)) % p == 0 for n in range(lim))
            ok &= g1 and g2 and g3 and g4 and g5
            det.append(f"p={p}:degE<=p+1:{g1},[t^(p+1)]E1=0modp:{g2},degM<=p:{g3},degA0<=floor(pr):{g4},E1=+-t^p mod p:{g5}")
        report(f"[8c] polynomiality and degrees of E_k = e_k D^k0, M_l = n_l D^k0, A_0 = S_0 D^o (Section 6): {nm}", ok, "; ".join(det))

def monotraces(D, p, imax, M):
    w = D['w']; c = D['chi'](p); th = D['th']; t = D['t']; MZ = p*M
    x0 = mul(th[:M], inv(subqp(th[:M], p, M), M), M)
    tp = subqp(t[:M], p, M)
    base = mul(subqp(th[:M], p, MZ), inv(th[:MZ], MZ), MZ)
    out = []; ti = [1] + [0]*(MZ-1); tpi = [1] + [0]*(M-1)
    for i in range(imax+1):
        Q = mul(x0, tpi, M); G = mul(base, ti, MZ)
        out.append(t_expand([Q[n] + p**(w+1)*c*G[p*n] for n in range(M)], t[:M], M))
        ti = mul(ti, t[:MZ], MZ); tpi = mul(tpi, tp, M)
    return out

def check_ell_decay():
    for nm in ELL:
        ok = True; det = []
        for p in [5, 7, 11, 13]:
            imax = 3 if p < 11 else 2
            M = p*imax + 50
            D = ell_data(nm, p*M)
            r = Fr(D['ordz'], D['k0']); dstar = Fr((p+1)*D['lam'], D['k0']); c0 = floor(p*r)
            S = monotraces(D, p, imax, M)
            slack = 10**9
            for i in range(imax+1):
                for n in range(M):
                    if S[i][n] == 0: continue
                    b = ceil(Fr(n - p*i - c0)/dstar)
                    if vp(S[i][n], p) < b: ok = False
                    if n >= p*i + p: slack = min(slack, vp(S[i][n], p) - b)
            det.append(f"p={p}:min slack={slack}")
        report(f"[8d] decay estimate v_p([t^n]T_p(t^i)) >= (n-pi-floor(pr))/delta*: {nm}", ok, " ".join(det))

def check_ell_cong():
    for nm in ELL:
        ok = True; det = []
        for p in [5, 7, 11, 13]:
            for (s, m) in [(1, 1), (1, 2), (1, 3), (2, 1), (2, 2), (3, 1)]:
                K = m*p**s
                if K > 130: continue
                M = K + 25
                D = ell_data(nm, M)
                w = D['w']; th = D['th'][:M]; t = D['t'][:M]
                A = [D['an'](n) for n in range(K)]
                lhs = mul(th, polyt(A[:m*p**(s-1)], subqp(t, p, M), M), M)
                rhs = mul(subqp(th, p, M), polyt(A, t, M), M)
                v = vp_series([x-y for x, y in zip(lhs, rhs)], p)
                ok &= v >= (w+1)*s
                det.append(f"p={p},s={s},m={m}:v={v}")
        report(f"[8e] Theorem B: Theta(q)F_(mp^(s-1))(t(q^p)) = Theta(q^p)F_(mp^s)(t(q)) mod p^((w+1)s): {nm}",
               ok, " ".join(det))

def check_ell_carry():
    # the inequalities used in the proof of Theorem B, for p < 200 and digit counts k <= 12
    ok = True
    for nm, c, L0, e0, r in [("27t", 2, 3, 1, Fr(1, 3)), ("256t", 3, 4, 1, Fr(1, 4)), ("108t", 3, 6, 2, Fr(1, 3))]:
        for p in [x for x in range(5, 200) if all(x % d for d in range(2, x))]:
            dstar = Fr(p+1, L0)
            for k in range(0, 13):
                j = 0 if k == 0 else p**(k-1)
                dec = ceil(Fr(p*(j+1) - floor(p*r))/dstar)
                ok &= dec >= c*(k+1)
        # carry bound per digit p-1
        an = ell_data(nm, 10)['an']
        for p in [5, 7, 11, 13]:
            for n in [p**3 - 1, 2*p**3 - 1, p**3 - 2, p**2 - 1 + 3*p**3]:
                digits = 0; x = n
                while x:
                    digits += (x % p == p-1); x //= p
                ok &= vp(an(n), p) >= c*digits
    report("[8f] arithmetic of the induction: decay at offset p(j+1) >= c(k+1) (p<200, k<=12) and "
           "v_p(a_n) >= c * #(digits p-1)", ok)

def check_hyperoct3_termwise():
    # [7e] the decay estimate holds for the hyperoctahedral family (k0=2, o=1, r=1/2), but the termwise bound
    #      of Section 6 is too weak for s=3: it gives 5 < 6, while the tail of the trace is divisible by p^6
    M0 = 200
    v = etaseries({1: 12, 6: 12, 2: -12, 3: -12}, 5*M0)
    opv = [(1 if i == 0 else 0) + v[i] for i in range(5*M0)]
    u = mul(v, inv(mul(opv, opv, 5*M0), 5*M0), 5*M0)
    th = mul(opv, etaseries({2: 7, 3: 7, 1: -5, 6: -5}, 5*M0), 5*M0)
    D = dict(w=2, chi=chi_1, th=th, t=u)
    an = lambda n: C(2*n, n)*sum(C(n, j)**2*C(2*j, j) for j in range(n+1))
    p, s, m = 5, 3, 1
    Ks, Ks1 = m*p**s, m*p**(s-1)
    M = Ks + 25
    S = monotraces(D, p, Ks1 - 1, M)
    a = [an(i) for i in range(Ks1)]
    dstar, c0 = Fr(p+1, 2), p//2
    de = all(S[i][n] == 0 or vp(S[i][n], p) >= ceil(Fr(n - p*i - c0)/dstar) for i in range(Ks1) for n in range(M))
    tw = min(min(vp(a[i], p) + vp(S[i][n], p) for i in range(Ks1)) for n in range(Ks, M))
    tr = min(vp(sum(a[i]*S[i][n] for i in range(Ks1)), p) for n in range(Ks, M))
    # [7d] degree of the polynomial D: (p-1)/2 for p = 1 mod 3 and (p+1)/2 for p = 2 mod 3 (CM by Q(sqrt(-3)))
    degs = []
    for pp in [5, 7, 11, 13]:
        MM = 30
        Dt = Dpoly(dict(vth=th, chi=chi_1, t=u), pp, MM)
        dg = deg_tail(Dt, MM - 2)
        degs.append((pp, dg, (pp - 1)//2 if pp % 3 == 1 else (pp + 1)//2))
    report("[7d] hyperoctahedral n=3: deg D = (p-1)/2 if p = 1 mod 3 and (p+1)/2 if p = 2 mod 3, p in {5,7,11,13}",
           all(x[1] == x[2] for x in degs), str(degs))
    report("[7e] hyperoctahedral n=3, p=5, s=3: decay estimate holds; termwise bound min_i v(a_i)+v([u^n]S_i) "
           "over n in [125,150) is 5 < 6, but the tail of Tr(Q_P) has valuation >= 6", de and tw == 5 and tr >= 6,
           f"decay ok={de}, termwise={tw}, true={tr}")

def check_bv_simplicial():
    lin = lambda a0, a1, n: (a0 + a1*n[0], a1*n[1])
    def dmul(*xs):
        r = (Fr(1), Fr(0))
        for x in xs: r = (r[0]*x[0], r[0]*x[1] + r[1]*x[0])
        return r
    sc = lambda c, x: (c*x[0], c*x[1])
    rec2 = lambda n: (dmul(n, n), [(3, sc(27, dmul(lin(-1, 1, n), lin(-2, 1, n))))])
    bv_check("simplicial n=2 (c_{3k}=(3k)!/k!^3)", rec2, 3, 3, [5, 7], [1, 2], 110, lambda s: 2*s, "[8g]")
    rec3 = lambda n: (dmul(n, n, n), [(4, sc(256, dmul(lin(-1, 1, n), lin(-2, 1, n), lin(-3, 1, n))))])
    bv_check("simplicial n=3 (c_{4k}=(4k)!/k!^4)", rec3, 4, 4, [5, 7], [1, 2], 110, lambda s: 3*s, "[8g]")
    # the modular t(q) of 27t / 256t is the mirror map:  t_0(Q)^(n+1) = t(Q^(n+1))
    for nm, rec, k in [("27t", rec2, 3), ("256t", rec3, 4)]:
        M = 61
        F0, q = frob_canonical(rec, k, M)
        T = series_inverse_fn(q, M)
        Tk = [Fr(1)] + [Fr(0)]*(M-1)
        for _ in range(k): Tk = fmul(Tk, T, M)
        D = ell_data(nm, M)
        tk = subqp(D['t'][:M], k, M)
        report(f"[8g] mirror map: t_0(Q)^{k} = t(Q^{k}) for the simplicial family attached to {nm}",
               all(Tk[n] == tk[n] for n in range(M)))


# ---------------------------------------------------------------- Check 9: Zagier's case D


def leg5(n):
    r = n % 5
    return 0 if r == 0 else (1 if r in (1, 4) else -1)
def tD(M):
    """t = q prod (1-q^n)^(5 (n/5))"""
    P = [0]*M; P[0] = 1
    for n in range(1, M):
        e = 5*leg5(n)
        for _ in range(abs(e)):
            if e > 0:
                for i in range(M-1, n-1, -1): P[i] -= P[i-n]
            else:
                for i in range(n, M): P[i] += P[i-n]
    return [0] + P[:M-1]
def thetaD(M):
    """Theta = prod_{n = +-1 mod 5} (1-q^n)^(-3) prod_{other n} (1-q^n)^2"""
    P = [0]*M; P[0] = 1
    for n in range(1, M):
        e = -3 if n % 5 in (1, 4) else 2
        for _ in range(abs(e)):
            if e > 0:
                for i in range(M-1, n-1, -1): P[i] -= P[i-n]
            else:
                for i in range(n, M): P[i] += P[i-n]
    return P
aD = lambda n: sum(C(n, k)**2*C(n+k, k) for k in range(n+1))

# Laurent series: (v, a) = q^v * sum_i a[i] q^i
def lmul(A, B, M): return (A[0] + B[0], mul(A[1], B[1], M))
def ladd(A, B, M):
    v = min(A[0], B[0]); out = [0]*M
    for (w, a) in (A, B):
        for i, x in enumerate(a):
            if 0 <= i + w - v < M: out[i + w - v] += x
    return (v, out)
def lscale(A, c): return (A[0], [c*x for x in A[1]])
def lUp(A, p, M):
    v, a = A
    v0 = -((-v)//p)                      # ceil(v/p)
    return (v0, [a[p*n - v] if 0 <= p*n - v < len(a) else 0 for n in range(v0, v0 + M)])
def lsubqp(A, p, M):
    v, a = A
    return (p*v, subqp(a, p, M))
def to_t(A, t, M):
    """Laurent q-series -> dict {j: c_j} with A = sum_j c_j t^j (t = q u, u a unit), M relative terms"""
    v, a = A
    u = t[1:] + [0]
    if v < 0:
        uu = [1] + [0]*(M-1)
        for _ in range(-v): uu = mul(uu, u, M)       # A t^{-v} = a u^{-v} (power series)
        f = mul(a, uu, M)
    else:
        f = [0]*v + a[:M-v]
    c = t_expand(f, t[:M], M)
    sh = v if v < 0 else 0
    return {j + sh: c[j] for j in range(M) if c[j]}
def Pof(P, X, M):                                  # P(X) for a Laurent series X
    res = (0, [0]*M)
    pw = (0, [1] + [0]*(M-1))
    for cf in P:
        res = ladd(res, lscale(pw, cf), M) if cf else res
        pw = lmul(pw, X, M)
    return res

def zd_series(M):
    t = tD(M); th = thetaD(M)
    u = t[1:] + [0]
    tinv = (-1, inv(u, M))                         # 1/t
    return t, th, tinv

def check_zd_param():
    M = 80
    t, th, tinv = zd_series(M)
    A = [aD(n) for n in range(M)]
    ok1 = polyt(A, t, M) == th
    ok2 = mul(mul(th, th, M), t, M) == etaseries({5: 5, 1: -1}, M)
    # 1/t - 11 - t = eta1^6/eta5^6, i.e. (1 - 11 t - t^2) = t eta1^6/eta5^6
    lhs3 = [(1 if n == 0 else 0) - 11*t[n] - mul(t, t, M)[n] for n in range(M)]
    # t * eta1^6/eta5^6 = (t/q) * prod (1-q^n)^6 (1-q^(5n))^(-6)
    u = t[1:] + [0]
    P6 = [0]*M; P6[0] = 1
    for nn in range(1, M):
        for e, mm in ((6, 1), (-6, 5)):
            k = mm*nn
            if k >= M: continue
            for _ in range(abs(e)):
                if e > 0:
                    for i in range(M-1, k-1, -1): P6[i] -= P6[i-k]
                else:
                    for i in range(k, M): P6[i] += P6[i-k]
    ok3 = lhs3[:M-1] == mul(u, P6, M)[:M-1]          # u = t/q is known to O(q^(M-1))
    # Ligozat: eta5^5/eta1 has order 1 at the cusp infinity of X_0(5) (width 1) and order 0 at the cusp 0
    ok3 = ok3 and ligozat({5: 5, 1: -1}, 5, 5) == 1 and ligozat({5: 5, 1: -1}, 5, 1) == 0
    report("[9a] Zagier D: F(t(q)) = Theta(q), Theta^2 t = eta5^5/eta1 and 1/t - 11 - t = eta1^6/eta5^6 to O(q^79); orders of eta5^5/eta1 at the cusps of X_0(5) are 1, 0",
           ok1 and ok2 and ok3)
    # diamond operators (floating point): t o<d> and Theta|<d> for d = 2, 3, 4
    MM = 400
    tt, thh, _ = zd_series(MM)
    def ev(c, tau):
        q = cmath.exp(2j*cmath.pi*tau)
        return sum(c[k]*q**k for k in range(MM))
    res = []
    for (sig, d) in [((3, 1, 5, 2), 2), ((2, 1, 5, 3), 3), ((4, 1, 15, 4), 4)]:
        a, b, c, dd = sig
        tau = complex(-0.4, 0.2) if d != 4 else complex(-4/15, 0.07)
        st = (a*tau + b)/(c*tau + dd)
        tv, tsv = ev(tt, tau), ev(tt, st)
        thv, thsv = ev(thh, tau), ev(thh, st)
        if d in (2, 3):
            res.append((tv*tsv, thsv/((c*tau + dd)*tv*thv)))
        else:
            res.append((tsv/tv, thsv/((c*tau + dd)*thv)))
    # the modular t is the mirror map of the Apery zeta(2) equation:  n^2 a_n = (11n^2-11n+3) a_(n-1) + (n-1)^2 a_(n-2)
    def dm(x, y): return (x[0]*y[0], x[0]*y[1] + x[1]*y[0])
    def rec(n):
        n2 = dm(n, n)
        quad = (11*n2[0] - 11*n[0] + 3, 11*n2[1] - 11*n[1])
        nm1 = (n[0] - 1, n[1])
        return n2, [(1, quad), (2, dm(nm1, nm1))]
    F0, qc = frob_canonical(rec, 2, 50)
    T = series_inverse_fn(qc, 50)
    okm = all(F0[n] == aD(n) for n in range(50)) and all(T[n] == tD(50)[n] for n in range(50))
    report("[9a] Zagier D: the canonical coordinate of the Apery zeta(2) equation inverts t(q), i.e. t is the mirror map", okm)
    c2 = res[0][1].real
    ok = (abs(res[0][0] + 1) < 1e-9 and abs(res[1][0] + 1) < 1e-9 and abs(res[2][0] - 1) < 1e-9
          and abs(abs(c2) - 1) < 1e-9 and abs(res[1][1] + c2) < 1e-9 and abs(res[2][1] + 1) < 1e-9)
    report("[9a] diamond operators on Gamma_1(5): t o<2> = t o<3> = -1/t, Theta|<2> = c t Theta, "
           "Theta|<3> = -c t Theta, t o<4> = t, Theta|<4> = -Theta (floating point)", ok,
           f"c = {c2:+.12f}; t(<2>tau)t(tau) = {res[0][0]:.3e}")
    return round(c2)

def traces_zd(p, P, M, c):
    """exact Laurent q-series of the two Hecke traces for p = +-2 mod 5, or the trace for p = +-1 mod 5"""
    D = len(P) - 1
    MZ = p*(M + D + 3)
    t, th, tinv = zd_series(MZ)
    tq, thq = t[:M], th[:M]
    _, _, tinvq = zd_series(M)
    thinvZ = inv(th, MZ)
    Gz = mul(mul(subqp(th, p, MZ), thinvZ, MZ), polyt(P, t, MZ), MZ)
    UG = lUp((0, Gz), p, M)
    Q = (0, mul(mul(thq, inv(subqp(thq, p, M), M), M), polyt(P, subqp(tq, p, M), M), M))
    if p % 5 in (1, 4):
        eta1 = 1 if p % 5 == 1 else -1
        return ladd(lscale(Q, eta1), lscale(UG, p*p), M)
    eta = c if p % 5 == 2 else -c
    # Qcheck = Theta P(-1/t(q^p)) / (t(q^p) Theta(q^p))
    tinv_p = lsubqp(tinvq, p, M)
    Qc = lmul(lmul((0, thq), Pof(P, lscale(tinv_p, -1), M), M),
              lmul(tinv_p, (0, inv(subqp(thq, p, M), M)), M), M)
    Lam = ladd(lscale(Qc, eta), lscale(UG, p*p), M)
    # Gcheck(z) = Theta(z^p) P(-1/t(z)) / (t(z) Theta(z))
    tinvZ = (-1, inv(t[1:] + [0], MZ))
    Gc = lmul(lmul((0, subqp(th, p, MZ)), Pof(P, lscale(tinvZ, -1), MZ), MZ),
              lmul(tinvZ, (0, thinvZ), MZ), MZ)
    UGc = lUp(Gc, p, M)
    Lamc = ladd(lscale(lmul((0, tq), Q, M), -eta), lscale(lmul((0, tq), UGc, M), p*p), M)
    return Lam, Lamc

def check_zd_traces(c):
    random.seed(5)
    ok = True; det = []
    for p in [2, 3, 7, 13, 17, 11, 19]:
        for P in [[1], [1, 3, 19], [random.randint(-9, 9) for _ in range(4)]]:
            if P[-1] == 0: P[-1] = 1
            D = len(P) - 1; K = p*(D + 1)
            M = K + 30
            out = traces_zd(p, P, M, c)
            if p % 5 in (1, 4):
                Ld = to_t(out, zd_series(M)[0], M - 5)
                good = all(j >= 0 for j in Ld) and max(Ld) <= K - 1
                det.append(f"p={p},D={D}:poly deg {max(Ld)}<={K-1}" if good else f"p={p}:FAIL")
            else:
                Lam, Lamc = out
                t = zd_series(M)[0]
                L = to_t(Lam, t, M - 5)
                Lc = to_t(Lamc, t, M - 5)
                topc = max(Lc); bot = min(L)
                # Laurent polynomials: L supported in [-K, (D+1)//p - 1], Lc in [1-(D+1)//p, K]
                good = topc <= K and bot >= -K
                good &= max(L) <= (D + 1)//p - 1 + 0 and min(Lc) >= 1 - (D + 1)//p
                # Lc(t) = L(-1/t): coefficient of t^j in Lc equals (-1)^j * coefficient of t^(-j) in L
                good &= all(Lc.get(j, 0) == (-1)**j*L.get(-j, 0) for j in range(-K, K + 1))
                det.append(f"p={p},D={D}:top(Lc)={topc}<={K}")
            ok &= good
    report("[9b] Hecke traces on Gamma_1(5): p=+-1: eta'Q+p^2U_pG polynomial of degree<p(D+1); "
           "p=+-2: Lambda = eta Qcheck + p^2 U_pG and Lambda' = -eta t Q + p^2 t U_p Gcheck are Laurent "
           "polynomials, deg Lambda' <= p(D+1), Lambda'(t) = Lambda(-1/t)", ok, " ".join(det[:12]) + " ...")

def check_zd_cong():
    ok = True; det = []
    for p in [2, 3, 7, 11, 13, 17]:
        for (s, m) in [(1, 1), (1, 2), (2, 1), (2, 2), (3, 1)]:
            K = m*p**s
            if K > 125: continue
            Kp = m*p**(s-1)
            M = K + 20
            t, th, tinv = zd_series(M)
            A = [aD(n) for n in range(K)]
            lhs = mul(th, polyt(A[:Kp], subqp(t, p, M), M), M)
            rhs = mul(subqp(th, p, M), polyt(A, t, M), M)
            vH = vp_series([x - y for x, y in zip(lhs, rhs)], p)
            # dual congruence: t Theta F_(K')(-1/t(q^p)) / (t(q^p) Theta(q^p))  =  F_K(-1/t(q))
            tinv_p = lsubqp(tinv, p, M)
            L1 = lmul(lmul((0, mul(t, th, M)), Pof(A[:Kp], lscale(tinv_p, -1), M), M),
                      lmul(tinv_p, (0, inv(subqp(th, p, M), M)), M), M)
            R1 = Pof(A, lscale(tinv, -1), M)
            Dif = ladd(L1, lscale(R1, -1), M)
            vD = min(vp(x, p) for x in Dif[1][:M - 2*K//p - 5]) if any(Dif[1]) else 10**9
            ok &= vH >= 2*s and vD >= 2*s
            det.append(f"p={p},s={s},m={m}:v(H)={vH},v(D)={vD}")
    report("[9c] Theorem C: (H_s) and the dual congruence (D_s) modulo p^(2s) for Zagier D", ok, " ".join(det))


if __name__ == "__main__":
    check_param()
    check_cusps()
    check_trace_lemma()
    check_main()
    check_coeff()
    check_bv3()
    check_hyperoct3()
    check_hyperoct3_termwise()
    check_ell_param()
    check_ell_D()
    check_ell_sym()
    check_ell_decay()
    check_ell_cong()
    check_ell_carry()
    check_bv_simplicial()
    check_ell_nonpoly()
    c = check_zd_param()
    check_zd_traces(c)
    check_zd_cong()
    print()
    print("ALL CHECKS PASSED" if not FAILS else f"{len(FAILS)} CHECK(S) FAILED: {FAILS}")
    sys.exit(1 if FAILS else 0)
\end{lstlisting}

\clearpage
\section{Recorded output of the modular verification program}\label{B:source-appendix-2}\label{B:app:log}
This is the execution record supplied with the source manuscript, not a newly claimed run. Its historical ``open case'' wording is superseded by Theorem~\ref{P:thm:main-all}.

\begin{lstlisting}
PASS [1] F(t(q)) = Theta(q) to O(q^60) for Legendre 2F1(1/2,1/2;1|16t)
PASS [1] F(t(q)) = Theta(q) to O(q^60) for Zagier A (Franel)
PASS [1] F(t(q)) = Theta(q) to O(q^60) for Zagier B
PASS [1] F(t(q)) = Theta(q) to O(q^60) for Zagier C
PASS [1] F(t(q)) = Theta(q) to O(q^60) for Zagier E
PASS [1] F(t(q)) = Theta(q) to O(q^60) for Zagier F
PASS [1] F(t(q)) = Theta(q) to O(q^60) for 3F2(1/2,1/2,1/2;1,1|64t)
PASS [1] F(t(q)) = Theta(q) to O(q^60) for Apery (gamma)
PASS [1] F(t(q)) = Theta(q) to O(q^60) for Almkvist (eta)
PASS [1] F(t(q)) = Theta(q) to O(q^60) for Domb (alpha)
PASS [1] F(t(q)) = Theta(q) to O(q^60) for Almkvist-Zudilin (delta)
PASS [1] F(t(q)) = Theta(q) to O(q^60) for Almkvist (epsilon)
PASS [1] F(t(q)) = Theta(q) to O(q^60) for Almkvist (zeta)
PASS [2] zeros of Theta = poles of t = cusps with denominators [2], r_Theta = 1/2, (C) holds, character and Atkin-Lehner symmetry (G0(4)) ok: Legendre 2F1(1/2,1/2;1|16t) :: ord Theta {1: Fraction(0, 1), 2: Fraction(1, 2), 4: Fraction(0, 1)}, ord t {1: Fraction(0, 1), 2: Fraction(-1, 1), 4: Fraction(1, 1)}
PASS [2] zeros of Theta = poles of t = cusps with denominators [3], r_Theta = 1, (C) holds, character and Atkin-Lehner symmetry (G0(6)) ok: Zagier A (Franel) :: ord Theta {1: Fraction(0, 1), 2: Fraction(0, 1), 3: Fraction(1, 1), 6: Fraction(0, 1)}, ord t {1: Fraction(0, 1), 2: Fraction(0, 1), 3: Fraction(-1, 1), 6: Fraction(1, 1)}
PASS [2] zeros of Theta = poles of t = cusps with denominators [1], r_Theta = 1, (C) holds, character and Atkin-Lehner symmetry (G0(9)) ok: Zagier B :: ord Theta {1: Fraction(1, 1), 3: Fraction(0, 1), 9: Fraction(0, 1)}, ord t {1: Fraction(-1, 1), 3: Fraction(0, 1), 9: Fraction(1, 1)}
PASS [2] zeros of Theta = poles of t = cusps with denominators [2], r_Theta = 1, (C) holds, character and Atkin-Lehner symmetry (G0(6)) ok: Zagier C :: ord Theta {1: Fraction(0, 1), 2: Fraction(1, 1), 3: Fraction(0, 1), 6: Fraction(0, 1)}, ord t {1: Fraction(0, 1), 2: Fraction(-1, 1), 3: Fraction(0, 1), 6: Fraction(1, 1)}
PASS [2] zeros of Theta = poles of t = cusps with denominators [2], r_Theta = 1, (C) holds, character and Atkin-Lehner symmetry (G0(8)) ok: Zagier E :: ord Theta {1: Fraction(0, 1), 2: Fraction(1, 1), 4: Fraction(0, 1), 8: Fraction(0, 1)}, ord t {1: Fraction(0, 1), 2: Fraction(-1, 1), 4: Fraction(0, 1), 8: Fraction(1, 1)}
PASS [2] zeros of Theta = poles of t = cusps with denominators [1], r_Theta = 1, (C) holds, character and Atkin-Lehner symmetry (G0(6)) ok: Zagier F :: ord Theta {1: Fraction(1, 1), 2: Fraction(0, 1), 3: Fraction(0, 1), 6: Fraction(0, 1)}, ord t {1: Fraction(-1, 1), 2: Fraction(0, 1), 3: Fraction(0, 1), 6: Fraction(1, 1)}
PASS [2] zeros of Theta = poles of t = cusps with denominators [2], r_Theta = 1/2, (C) holds, character and Atkin-Lehner symmetry (G0(4)+4) ok: 3F2(1/2,1/2,1/2;1,1|64t) :: ord Theta {1: Fraction(0, 1), 2: Fraction(1, 1), 4: Fraction(0, 1)}, ord t {1: Fraction(1, 1), 2: Fraction(-2, 1), 4: Fraction(1, 1)}
PASS [2] zeros of Theta = poles of t = cusps with denominators [2, 3], r_Theta = 1, (C) holds, character and Atkin-Lehner symmetry (G0(6)+6) ok: Apery (gamma) :: ord Theta {1: Fraction(0, 1), 2: Fraction(1, 1), 3: Fraction(1, 1), 6: Fraction(0, 1)}, ord t {1: Fraction(1, 1), 2: Fraction(-1, 1), 3: Fraction(-1, 1), 6: Fraction(1, 1)}
PASS [2] zeros of Theta = poles of t = cusps with denominators [1], r_Theta = 1, (C) holds, character and Atkin-Lehner symmetry (G0(5)) ok: Almkvist (eta) :: ord Theta {1: Fraction(1, 1), 5: Fraction(0, 1)}, ord t {1: Fraction(-1, 1), 5: Fraction(1, 1)}
PASS [2] zeros of Theta = poles of t = cusps with denominators [1, 3], r_Theta = 1, (C) holds, character and Atkin-Lehner symmetry (G0(6)+3) ok: Domb (alpha) :: ord Theta {1: Fraction(1, 1), 2: Fraction(0, 1), 3: Fraction(1, 1), 6: Fraction(0, 1)}, ord t {1: Fraction(-1, 1), 2: Fraction(1, 1), 3: Fraction(-1, 1), 6: Fraction(1, 1)}
PASS [2] zeros of Theta = poles of t = cusps with denominators [1, 2], r_Theta = 1, (C) holds, character and Atkin-Lehner symmetry (G0(6)+2) ok: Almkvist-Zudilin (delta) :: ord Theta {1: Fraction(1, 1), 2: Fraction(1, 1), 3: Fraction(0, 1), 6: Fraction(0, 1)}, ord t {1: Fraction(-1, 1), 2: Fraction(-1, 1), 3: Fraction(1, 1), 6: Fraction(1, 1)}
PASS [2] zeros of Theta = poles of t = cusps with denominators [2, 4], r_Theta = 1, (C) holds, character and Atkin-Lehner symmetry (G0(8)+8) ok: Almkvist (epsilon) :: ord Theta {1: Fraction(0, 1), 2: Fraction(1, 1), 4: Fraction(1, 1), 8: Fraction(0, 1)}, ord t {1: Fraction(1, 1), 2: Fraction(-1, 1), 4: Fraction(-1, 1), 8: Fraction(1, 1)}
PASS [2] zeros of Theta = poles of t = cusps with denominators [3], r_Theta = 1, (C) holds, character and Atkin-Lehner symmetry (G0(9)+9) ok: Almkvist (zeta) :: ord Theta {1: Fraction(0, 1), 3: Fraction(1, 1), 9: Fraction(0, 1)}, ord t {1: Fraction(1, 1), 3: Fraction(-1, 1), 9: Fraction(1, 1)}
PASS [3] Tr(Q_P) is a polynomial in t of degree <= p(L-1)+floor(r(p-1)) (exact, to t^(D+19)): Legendre 2F1(1/2,1/2;1|16t) :: p=3,F_{1}:deg<=1:ok; p=3,F_{2}:deg<=4:ok; p=3,F_{3}:deg<=7:ok; p=3,random deg 2:deg<=7:ok; p=3,random deg 5:deg<=16:ok; p=5,F_{1}:deg<=2:ok ...
PASS [3] Tr(Q_P) is a polynomial in t of degree <= p(L-1)+floor(r(p-1)) (exact, to t^(D+19)): Zagier A (Franel) :: p=5,F_{1}:deg<=4:ok; p=5,F_{2}:deg<=9:ok; p=5,F_{5}:deg<=24:ok; p=5,random deg 2:deg<=14:ok; p=5,random deg 5:deg<=29:ok; p=7,F_{1}:deg<=6:ok ...
PASS [3] Tr(Q_P) is a polynomial in t of degree <= p(L-1)+floor(r(p-1)) (exact, to t^(D+19)): Zagier B :: p=2,F_{1}:deg<=1:ok; p=2,F_{2}:deg<=3:ok; p=2,F_{2}:deg<=3:ok; p=2,random deg 2:deg<=5:ok; p=2,random deg 5:deg<=11:ok; p=5,F_{1}:deg<=4:ok ...
PASS [3] Tr(Q_P) is a polynomial in t of degree <= p(L-1)+floor(r(p-1)) (exact, to t^(D+19)): Zagier C :: p=5,F_{1}:deg<=4:ok; p=5,F_{2}:deg<=9:ok; p=5,F_{5}:deg<=24:ok; p=5,random deg 2:deg<=14:ok; p=5,random deg 5:deg<=29:ok; p=7,F_{1}:deg<=6:ok ...
PASS [3] Tr(Q_P) is a polynomial in t of degree <= p(L-1)+floor(r(p-1)) (exact, to t^(D+19)): Zagier E :: p=3,F_{1}:deg<=2:ok; p=3,F_{2}:deg<=5:ok; p=3,F_{3}:deg<=8:ok; p=3,random deg 2:deg<=8:ok; p=3,random deg 5:deg<=17:ok; p=5,F_{1}:deg<=4:ok ...
PASS [3] Tr(Q_P) is a polynomial in t of degree <= p(L-1)+floor(r(p-1)) (exact, to t^(D+19)): Zagier F :: p=5,F_{1}:deg<=4:ok; p=5,F_{2}:deg<=9:ok; p=5,F_{5}:deg<=24:ok; p=5,random deg 2:deg<=14:ok; p=5,random deg 5:deg<=29:ok; p=7,F_{1}:deg<=6:ok ...
PASS [3] Tr(Q_P) is a polynomial in t of degree <= p(L-1)+floor(r(p-1)) (exact, to t^(D+19)): 3F2(1/2,1/2,1/2;1,1|64t) :: p=3,F_{1}:deg<=1:ok; p=3,F_{2}:deg<=4:ok; p=3,F_{3}:deg<=7:ok; p=3,random deg 2:deg<=7:ok; p=3,random deg 5:deg<=16:ok; p=5,F_{1}:deg<=2:ok ...
PASS [3] Tr(Q_P) is a polynomial in t of degree <= p(L-1)+floor(r(p-1)) (exact, to t^(D+19)): Apery (gamma) :: p=5,F_{1}:deg<=4:ok; p=5,F_{2}:deg<=9:ok; p=5,F_{5}:deg<=24:ok; p=5,random deg 2:deg<=14:ok; p=5,random deg 5:deg<=29:ok; p=7,F_{1}:deg<=6:ok ...
PASS [3] Tr(Q_P) is a polynomial in t of degree <= p(L-1)+floor(r(p-1)) (exact, to t^(D+19)): Almkvist (eta) :: p=2,F_{1}:deg<=1:ok; p=2,F_{2}:deg<=3:ok; p=2,F_{2}:deg<=3:ok; p=2,random deg 2:deg<=5:ok; p=2,random deg 5:deg<=11:ok; p=3,F_{1}:deg<=2:ok ...
PASS [3] Tr(Q_P) is a polynomial in t of degree <= p(L-1)+floor(r(p-1)) (exact, to t^(D+19)): Domb (alpha) :: p=5,F_{1}:deg<=4:ok; p=5,F_{2}:deg<=9:ok; p=5,F_{5}:deg<=24:ok; p=5,random deg 2:deg<=14:ok; p=5,random deg 5:deg<=29:ok; p=7,F_{1}:deg<=6:ok ...
PASS [3] Tr(Q_P) is a polynomial in t of degree <= p(L-1)+floor(r(p-1)) (exact, to t^(D+19)): Almkvist-Zudilin (delta) :: p=5,F_{1}:deg<=4:ok; p=5,F_{2}:deg<=9:ok; p=5,F_{5}:deg<=24:ok; p=5,random deg 2:deg<=14:ok; p=5,random deg 5:deg<=29:ok; p=7,F_{1}:deg<=6:ok ...
PASS [3] Tr(Q_P) is a polynomial in t of degree <= p(L-1)+floor(r(p-1)) (exact, to t^(D+19)): Almkvist (epsilon) :: p=3,F_{1}:deg<=2:ok; p=3,F_{2}:deg<=5:ok; p=3,F_{3}:deg<=8:ok; p=3,random deg 2:deg<=8:ok; p=3,random deg 5:deg<=17:ok; p=5,F_{1}:deg<=4:ok ...
PASS [3] Tr(Q_P) is a polynomial in t of degree <= p(L-1)+floor(r(p-1)) (exact, to t^(D+19)): Almkvist (zeta) :: p=2,F_{1}:deg<=1:ok; p=2,F_{2}:deg<=3:ok; p=2,F_{2}:deg<=3:ok; p=2,random deg 2:deg<=5:ok; p=2,random deg 5:deg<=11:ok; p=5,F_{1}:deg<=4:ok ...
PASS [4] Theta(q)F_(mp^(s-1))(t(q^p)) = Theta(q^p)F_(mp^s)(t(q)) mod p^((w+1)s), w=1: Legendre 2F1(1/2,1/2;1|16t) :: p=3,s=1,m=1:v=2 p=3,s=1,m=2:v=4 p=3,s=2,m=1:v=4 p=3,s=2,m=2:v=6 p=3,s=3,m=1:v=6 p=5,s=1,m=1:v=2 p=5,s=1,m=2:v=2 p=5,s=2,m=1:v=4 p=5,s=2,m=2:v=4 p=5,s=3,m=1:v=6 p=7,s=1,m=1:v=2 p=7,s=1,m=2:v=2 p=7,s=2,m=1:v=4 p=7,s=2,m=2:v=4 p=11,s=1,m=1:v=2 p=11,s=1,m=2:v=2 p=11,s=2,m=1:v=4 p=13,s=1,m=1:v=2 p=13,s=1,m=2:v=2
PASS [4] Theta(q)F_(mp^(s-1))(t(q^p)) = Theta(q^p)F_(mp^s)(t(q)) mod p^((w+1)s), w=1: Zagier A (Franel) :: p=5,s=1,m=1:v=2 p=5,s=1,m=2:v=2 p=5,s=2,m=1:v=4 p=5,s=2,m=2:v=4 p=5,s=3,m=1:v=6 p=7,s=1,m=1:v=2 p=7,s=1,m=2:v=2 p=7,s=2,m=1:v=4 p=7,s=2,m=2:v=4 p=11,s=1,m=1:v=2 p=11,s=1,m=2:v=2 p=11,s=2,m=1:v=4 p=13,s=1,m=1:v=2 p=13,s=1,m=2:v=2
PASS [4] Theta(q)F_(mp^(s-1))(t(q^p)) = Theta(q^p)F_(mp^s)(t(q)) mod p^((w+1)s), w=1: Zagier B :: p=2,s=1,m=1:v=2 p=2,s=1,m=2:v=4 p=2,s=2,m=1:v=4 p=2,s=2,m=2:v=6 p=2,s=3,m=1:v=6 p=5,s=1,m=1:v=2 p=5,s=1,m=2:v=2 p=5,s=2,m=1:v=4 p=5,s=2,m=2:v=4 p=5,s=3,m=1:v=6 p=7,s=1,m=1:v=2 p=7,s=1,m=2:v=2 p=7,s=2,m=1:v=4 p=7,s=2,m=2:v=4 p=11,s=1,m=1:v=2 p=11,s=1,m=2:v=2 p=11,s=2,m=1:v=4 p=13,s=1,m=1:v=2 p=13,s=1,m=2:v=2
PASS [4] Theta(q)F_(mp^(s-1))(t(q^p)) = Theta(q^p)F_(mp^s)(t(q)) mod p^((w+1)s), w=1: Zagier C :: p=5,s=1,m=1:v=2 p=5,s=1,m=2:v=2 p=5,s=2,m=1:v=4 p=5,s=2,m=2:v=4 p=5,s=3,m=1:v=6 p=7,s=1,m=1:v=2 p=7,s=1,m=2:v=2 p=7,s=2,m=1:v=4 p=7,s=2,m=2:v=4 p=11,s=1,m=1:v=2 p=11,s=1,m=2:v=2 p=11,s=2,m=1:v=4 p=13,s=1,m=1:v=2 p=13,s=1,m=2:v=2
PASS [4] Theta(q)F_(mp^(s-1))(t(q^p)) = Theta(q^p)F_(mp^s)(t(q)) mod p^((w+1)s), w=1: Zagier E :: p=3,s=1,m=1:v=2 p=3,s=1,m=2:v=2 p=3,s=2,m=1:v=4 p=3,s=2,m=2:v=4 p=3,s=3,m=1:v=6 p=5,s=1,m=1:v=2 p=5,s=1,m=2:v=2 p=5,s=2,m=1:v=4 p=5,s=2,m=2:v=4 p=5,s=3,m=1:v=6 p=7,s=1,m=1:v=2 p=7,s=1,m=2:v=2 p=7,s=2,m=1:v=4 p=7,s=2,m=2:v=4 p=11,s=1,m=1:v=2 p=11,s=1,m=2:v=2 p=11,s=2,m=1:v=4 p=13,s=1,m=1:v=2 p=13,s=1,m=2:v=2
PASS [4] Theta(q)F_(mp^(s-1))(t(q^p)) = Theta(q^p)F_(mp^s)(t(q)) mod p^((w+1)s), w=1: Zagier F :: p=5,s=1,m=1:v=2 p=5,s=1,m=2:v=2 p=5,s=2,m=1:v=4 p=5,s=2,m=2:v=4 p=5,s=3,m=1:v=6 p=7,s=1,m=1:v=2 p=7,s=1,m=2:v=2 p=7,s=2,m=1:v=4 p=7,s=2,m=2:v=4 p=11,s=1,m=1:v=2 p=11,s=1,m=2:v=2 p=11,s=2,m=1:v=4 p=13,s=1,m=1:v=2 p=13,s=1,m=2:v=2
PASS [4] Theta(q)F_(mp^(s-1))(t(q^p)) = Theta(q^p)F_(mp^s)(t(q)) mod p^((w+1)s), w=2: 3F2(1/2,1/2,1/2;1,1|64t) :: p=3,s=1,m=1:v=3 p=3,s=1,m=2:v=6 p=3,s=2,m=1:v=6 p=3,s=2,m=2:v=9 p=3,s=3,m=1:v=9 p=5,s=1,m=1:v=3 p=5,s=1,m=2:v=3 p=5,s=2,m=1:v=6 p=5,s=2,m=2:v=6 p=5,s=3,m=1:v=9 p=7,s=1,m=1:v=3 p=7,s=1,m=2:v=3 p=7,s=2,m=1:v=6 p=7,s=2,m=2:v=6 p=11,s=1,m=1:v=3 p=11,s=1,m=2:v=3 p=11,s=2,m=1:v=6 p=13,s=1,m=1:v=3 p=13,s=1,m=2:v=3
PASS [4] Theta(q)F_(mp^(s-1))(t(q^p)) = Theta(q^p)F_(mp^s)(t(q)) mod p^((w+1)s), w=2: Apery (gamma) :: p=5,s=1,m=1:v=3 p=5,s=1,m=2:v=3 p=5,s=2,m=1:v=6 p=5,s=2,m=2:v=6 p=5,s=3,m=1:v=9 p=7,s=1,m=1:v=3 p=7,s=1,m=2:v=3 p=7,s=2,m=1:v=6 p=7,s=2,m=2:v=6 p=11,s=1,m=1:v=3 p=11,s=1,m=2:v=3 p=11,s=2,m=1:v=6 p=13,s=1,m=1:v=3 p=13,s=1,m=2:v=3
PASS [4] Theta(q)F_(mp^(s-1))(t(q^p)) = Theta(q^p)F_(mp^s)(t(q)) mod p^((w+1)s), w=2: Almkvist (eta) :: p=2,s=1,m=1:v=3 p=2,s=1,m=2:v=6 p=2,s=2,m=1:v=6 p=2,s=2,m=2:v=9 p=2,s=3,m=1:v=9 p=3,s=1,m=1:v=3 p=3,s=1,m=2:v=3 p=3,s=2,m=1:v=6 p=3,s=2,m=2:v=6 p=3,s=3,m=1:v=9 p=7,s=1,m=1:v=3 p=7,s=1,m=2:v=3 p=7,s=2,m=1:v=6 p=7,s=2,m=2:v=6 p=11,s=1,m=1:v=3 p=11,s=1,m=2:v=3 p=11,s=2,m=1:v=6 p=13,s=1,m=1:v=3 p=13,s=1,m=2:v=3
PASS [4] Theta(q)F_(mp^(s-1))(t(q^p)) = Theta(q^p)F_(mp^s)(t(q)) mod p^((w+1)s), w=2: Domb (alpha) :: p=5,s=1,m=1:v=3 p=5,s=1,m=2:v=3 p=5,s=2,m=1:v=6 p=5,s=2,m=2:v=6 p=5,s=3,m=1:v=9 p=7,s=1,m=1:v=3 p=7,s=1,m=2:v=3 p=7,s=2,m=1:v=6 p=7,s=2,m=2:v=6 p=11,s=1,m=1:v=3 p=11,s=1,m=2:v=3 p=11,s=2,m=1:v=6 p=13,s=1,m=1:v=3 p=13,s=1,m=2:v=3
PASS [4] Theta(q)F_(mp^(s-1))(t(q^p)) = Theta(q^p)F_(mp^s)(t(q)) mod p^((w+1)s), w=2: Almkvist-Zudilin (delta) :: p=5,s=1,m=1:v=3 p=5,s=1,m=2:v=3 p=5,s=2,m=1:v=6 p=5,s=2,m=2:v=6 p=5,s=3,m=1:v=9 p=7,s=1,m=1:v=3 p=7,s=1,m=2:v=3 p=7,s=2,m=1:v=6 p=7,s=2,m=2:v=6 p=11,s=1,m=1:v=3 p=11,s=1,m=2:v=3 p=11,s=2,m=1:v=6 p=13,s=1,m=1:v=3 p=13,s=1,m=2:v=3
PASS [4] Theta(q)F_(mp^(s-1))(t(q^p)) = Theta(q^p)F_(mp^s)(t(q)) mod p^((w+1)s), w=2: Almkvist (epsilon) :: p=3,s=1,m=1:v=3 p=3,s=1,m=2:v=3 p=3,s=2,m=1:v=6 p=3,s=2,m=2:v=6 p=3,s=3,m=1:v=9 p=5,s=1,m=1:v=3 p=5,s=1,m=2:v=3 p=5,s=2,m=1:v=6 p=5,s=2,m=2:v=6 p=5,s=3,m=1:v=9 p=7,s=1,m=1:v=3 p=7,s=1,m=2:v=3 p=7,s=2,m=1:v=6 p=7,s=2,m=2:v=6 p=11,s=1,m=1:v=3 p=11,s=1,m=2:v=3 p=11,s=2,m=1:v=6 p=13,s=1,m=1:v=3 p=13,s=1,m=2:v=3
PASS [4] Theta(q)F_(mp^(s-1))(t(q^p)) = Theta(q^p)F_(mp^s)(t(q)) mod p^((w+1)s), w=2: Almkvist (zeta) :: p=2,s=1,m=1:v=3 p=2,s=1,m=2:v=6 p=2,s=2,m=1:v=6 p=2,s=2,m=2:v=9 p=2,s=3,m=1:v=9 p=5,s=1,m=1:v=3 p=5,s=1,m=2:v=3 p=5,s=2,m=1:v=6 p=5,s=2,m=2:v=6 p=5,s=3,m=1:v=9 p=7,s=1,m=1:v=3 p=7,s=1,m=2:v=3 p=7,s=2,m=1:v=6 p=7,s=2,m=2:v=6 p=11,s=1,m=1:v=3 p=11,s=1,m=2:v=3 p=11,s=2,m=1:v=6 p=13,s=1,m=1:v=3 p=13,s=1,m=2:v=3
PASS [5] a(mp^s) = a(mp^(s-1)) mod p^((w+1)s) (50 cases): Legendre 2F1(1/2,1/2;1|16t)
PASS [5] a(mp^s) = a(mp^(s-1)) mod p^((w+1)s) (35 cases): Zagier A (Franel)
PASS [5] a(mp^s) = a(mp^(s-1)) mod p^((w+1)s) (50 cases): Zagier B
PASS [5] a(mp^s) = a(mp^(s-1)) mod p^((w+1)s) (35 cases): Zagier C
PASS [5] a(mp^s) = a(mp^(s-1)) mod p^((w+1)s) (50 cases): Zagier E
PASS [5] a(mp^s) = a(mp^(s-1)) mod p^((w+1)s) (35 cases): Zagier F
PASS [5] a(mp^s) = a(mp^(s-1)) mod p^((w+1)s) (50 cases): 3F2(1/2,1/2,1/2;1,1|64t)
PASS [5] a(mp^s) = a(mp^(s-1)) mod p^((w+1)s) (35 cases): Apery (gamma)
PASS [5] a(mp^s) = a(mp^(s-1)) mod p^((w+1)s) (53 cases): Almkvist (eta)
PASS [5] a(mp^s) = a(mp^(s-1)) mod p^((w+1)s) (35 cases): Domb (alpha)
PASS [5] a(mp^s) = a(mp^(s-1)) mod p^((w+1)s) (35 cases): Almkvist-Zudilin (delta)
PASS [5] a(mp^s) = a(mp^(s-1)) mod p^((w+1)s) (50 cases): Almkvist (epsilon)
PASS [5] a(mp^s) = a(mp^(s-1)) mod p^((w+1)s) (50 cases): Almkvist (zeta)
PASS [6] BV3 (19) with m=2, excellent lift q->q^p: hypercubic/hyperoctahedral n=2 (c_{2k}=C(2k,k)^2), p=3, s=1 :: v_p(F(t)F_(mp^(s-1))(t^s) - F(t^s)F_(mp^s)(t)) = 2 (required >= 2)
PASS [6] BV3 (19) with m=2, excellent lift q->q^p: hypercubic/hyperoctahedral n=2 (c_{2k}=C(2k,k)^2), p=3, s=2 :: v_p(F(t)F_(mp^(s-1))(t^s) - F(t^s)F_(mp^s)(t)) = 4 (required >= 4)
PASS [6] BV3 (19) with m=2, excellent lift q->q^p: hypercubic/hyperoctahedral n=2 (c_{2k}=C(2k,k)^2), p=5, s=1 :: v_p(F(t)F_(mp^(s-1))(t^s) - F(t^s)F_(mp^s)(t)) = 2 (required >= 2)
PASS [6] BV3 (19) with m=2, excellent lift q->q^p: hypercubic/hyperoctahedral n=2 (c_{2k}=C(2k,k)^2), p=5, s=2 :: v_p(F(t)F_(mp^(s-1))(t^s) - F(t^s)F_(mp^s)(t)) = 4 (required >= 4)
PASS [6] BV3 (19) with m=2, excellent lift q->q^p: hypercubic/hyperoctahedral n=2 (c_{2k}=C(2k,k)^2), p=7, s=1 :: v_p(F(t)F_(mp^(s-1))(t^s) - F(t^s)F_(mp^s)(t)) = 2 (required >= 2)
PASS [6] BV3 (19) with m=2, excellent lift q->q^p: hypercubic/hyperoctahedral n=2 (c_{2k}=C(2k,k)^2), p=7, s=2 :: v_p(F(t)F_(mp^(s-1))(t^s) - F(t^s)F_(mp^s)(t)) = 4 (required >= 4)
PASS [6] BV3 (19) with m=2, excellent lift q->q^p: hypercubic n=3 (c_{2k}=C(2k,k)^3), p=5, s=1 :: v_p(F(t)F_(mp^(s-1))(t^s) - F(t^s)F_(mp^s)(t)) = 3 (required >= 3)
PASS [6] BV3 (19) with m=2, excellent lift q->q^p: hypercubic n=3 (c_{2k}=C(2k,k)^3), p=5, s=2 :: v_p(F(t)F_(mp^(s-1))(t^s) - F(t^s)F_(mp^s)(t)) = 6 (required >= 6)
PASS [6] BV3 (19) with m=2, excellent lift q->q^p: hypercubic n=3 (c_{2k}=C(2k,k)^3), p=7, s=1 :: v_p(F(t)F_(mp^(s-1))(t^s) - F(t^s)F_(mp^s)(t)) = 3 (required >= 3)
PASS [6] BV3 (19) with m=2, excellent lift q->q^p: hypercubic n=3 (c_{2k}=C(2k,k)^3), p=7, s=2 :: v_p(F(t)F_(mp^(s-1))(t^s) - F(t^s)F_(mp^s)(t)) = 6 (required >= 6)
PASS [7a] (numerical, open case) hyperoctahedral n=3 of BV3 Conj. 7.5: congruence mod p^(2s) (the valuation is exactly 2s, not 3s) :: p=5,s=1:v=2 p=7,s=1:v=2 p=11,s=1:v=2 p=5,s=2:v=4
PASS [7b] hyperoctahedral n=3: mirror map u = v/(1+v)^2 and F(u(q)) = (1+v) eta2^7 eta3^7/(eta1^5 eta6^5), v = (eta1 eta6/(eta2 eta3))^12
PASS [7c] hyperoctahedral n=3: v_p(c_(p^k-1)) = k for p in {5,7,11,13}, k in {1,2} (only one factor p per digit p-1, so the carry bound of Section 6 fails)
PASS [7d] hyperoctahedral n=3: deg D = (p-1)/2 if p = 1 mod 3 and (p+1)/2 if p = 2 mod 3, p in {5,7,11,13} :: [(5, 3, 3), (7, 3, 3), (11, 6, 6), (13, 6, 6)]
PASS [7e] hyperoctahedral n=3, p=5, s=3: decay estimate holds; termwise bound min_i v(a_i)+v([u^n]S_i) over n in [125,150) is 5 < 6, but the tail of Tr(Q_P) has valuation >= 6 :: decay ok=True, termwise=5, true=6
PASS [8a] F(t(q)) = Theta(q) to O(q^60) for 27t
PASS [8a] F(t(q)) = Theta(q) to O(q^60) for 256t
PASS [8a] F(t(q)) = Theta(q) to O(q^60) for 108t
PASS [8a] 27t: a^3 = eta1^9/eta3^3 + 27 eta3^9/eta1^3 and a^3 t = eta3^9/eta1^3 (no zeros in H)
PASS [8a] 256t: E22^2 = eta1^16/eta2^8 + 64 eta2^16/eta1^8 and E22^4 t = eta1^8 eta2^8
PASS [8a] 108t: E23 = a^2 and E23^3 t = eta1^6 eta3^6
PASS [8a] |Theta(e)| at e=(3+sqrt-3)/6, (1+i)/2, (3+sqrt-3)/6 (floating point) :: 3.5e-16 3.2e-16 1.8e-16
PASS [8b] D(t) in Z[t], D = +-1 mod p, deg D <= (p+1)lam/k0: 27t :: p=5:deg=2<=2,D(0)=-1 p=7:deg=2<=8/3,D(0)=1 p=11:deg=4<=4,D(0)=-1 p=13:deg=4<=14/3,D(0)=1
PASS [8b] D(t) in Z[t], D = +-1 mod p, deg D <= (p+1)lam/k0: 256t :: p=5:deg=1<=3/2,D(0)=1 p=7:deg=2<=2,D(0)=1 p=11:deg=3<=3,D(0)=1 p=13:deg=3<=7/2,D(0)=1
PASS [8b] D(t) in Z[t], D = +-1 mod p, deg D <= (p+1)lam/k0 and B = D^2: 108t :: p=5:deg=1<=1,D(0)=-1 p=7:deg=1<=4/3,D(0)=1 p=11:deg=2<=2,D(0)=-1 p=13:deg=2<=7/3,D(0)=1
PASS [8c] polynomiality and degrees of E_k = e_k D^k0, M_l = n_l D^k0, A_0 = S_0 D^o (Section 6): 27t :: p=5:degE<=p+1:True,[t^(p+1)]E1=0modp:True,degM<=p:True,degA0<=floor(pr):True,E1=+-t^p mod p:True; p=7:degE<=p+1:True,[t^(p+1)]E1=0modp:True,degM<=p:True,degA0<=floor(pr):True,E1=+-t^p mod p:True
PASS [8c] polynomiality and degrees of E_k = e_k D^k0, M_l = n_l D^k0, A_0 = S_0 D^o (Section 6): 256t :: p=5:degE<=p+1:True,[t^(p+1)]E1=0modp:True,degM<=p:True,degA0<=floor(pr):True,E1=+-t^p mod p:True; p=7:degE<=p+1:True,[t^(p+1)]E1=0modp:True,degM<=p:True,degA0<=floor(pr):True,E1=+-t^p mod p:True
PASS [8c] polynomiality and degrees of E_k = e_k D^k0, M_l = n_l D^k0, A_0 = S_0 D^o (Section 6): 108t :: p=5:degE<=p+1:True,[t^(p+1)]E1=0modp:True,degM<=p:True,degA0<=floor(pr):True,E1=+-t^p mod p:True; p=7:degE<=p+1:True,[t^(p+1)]E1=0modp:True,degM<=p:True,degA0<=floor(pr):True,E1=+-t^p mod p:True
PASS [8d] decay estimate v_p([t^n]T_p(t^i)) >= (n-pi-floor(pr))/delta*: 27t :: p=5:min slack=0 p=7:min slack=0 p=11:min slack=0 p=13:min slack=0
PASS [8d] decay estimate v_p([t^n]T_p(t^i)) >= (n-pi-floor(pr))/delta*: 256t :: p=5:min slack=1 p=7:min slack=0 p=11:min slack=0 p=13:min slack=0
PASS [8d] decay estimate v_p([t^n]T_p(t^i)) >= (n-pi-floor(pr))/delta*: 108t :: p=5:min slack=0 p=7:min slack=1 p=11:min slack=0 p=13:min slack=0
PASS [8e] Theorem B: Theta(q)F_(mp^(s-1))(t(q^p)) = Theta(q^p)F_(mp^s)(t(q)) mod p^((w+1)s): 27t :: p=5,s=1,m=1:v=2 p=5,s=1,m=2:v=3 p=5,s=1,m=3:v=3 p=5,s=2,m=1:v=4 p=5,s=2,m=2:v=5 p=5,s=3,m=1:v=6 p=7,s=1,m=1:v=2 p=7,s=1,m=2:v=2 p=7,s=1,m=3:v=3 p=7,s=2,m=1:v=4 p=7,s=2,m=2:v=4 p=11,s=1,m=1:v=2 p=11,s=1,m=2:v=2 p=11,s=1,m=3:v=2 p=11,s=2,m=1:v=4 p=13,s=1,m=1:v=2 p=13,s=1,m=2:v=2 p=13,s=1,m=3:v=2
PASS [8e] Theorem B: Theta(q)F_(mp^(s-1))(t(q^p)) = Theta(q^p)F_(mp^s)(t(q)) mod p^((w+1)s): 256t :: p=5,s=1,m=1:v=4 p=5,s=1,m=2:v=5 p=5,s=1,m=3:v=6 p=5,s=2,m=1:v=7 p=5,s=2,m=2:v=8 p=5,s=3,m=1:v=10 p=7,s=1,m=1:v=3 p=7,s=1,m=2:v=4 p=7,s=1,m=3:v=4 p=7,s=2,m=1:v=6 p=7,s=2,m=2:v=7 p=11,s=1,m=1:v=3 p=11,s=1,m=2:v=3 p=11,s=1,m=3:v=4 p=11,s=2,m=1:v=6 p=13,s=1,m=1:v=3 p=13,s=1,m=2:v=3 p=13,s=1,m=3:v=3
PASS [8e] Theorem B: Theta(q)F_(mp^(s-1))(t(q^p)) = Theta(q^p)F_(mp^s)(t(q)) mod p^((w+1)s): 108t :: p=5,s=1,m=1:v=4 p=5,s=1,m=2:v=5 p=5,s=1,m=3:v=6 p=5,s=2,m=1:v=7 p=5,s=2,m=2:v=8 p=5,s=3,m=1:v=10 p=7,s=1,m=1:v=3 p=7,s=1,m=2:v=3 p=7,s=1,m=3:v=4 p=7,s=2,m=1:v=6 p=7,s=2,m=2:v=6 p=11,s=1,m=1:v=3 p=11,s=1,m=2:v=3 p=11,s=1,m=3:v=3 p=11,s=2,m=1:v=6 p=13,s=1,m=1:v=3 p=13,s=1,m=2:v=3 p=13,s=1,m=3:v=3
PASS [8f] arithmetic of the induction: decay at offset p(j+1) >= c(k+1) (p<200, k<=12) and v_p(a_n) >= c * #(digits p-1)
PASS [8g] BV3 (19) with m=3, excellent lift q->q^p: simplicial n=2 (c_{3k}=(3k)!/k!^3), p=5, s=1 :: v_p(F(t)F_(mp^(s-1))(t^s) - F(t^s)F_(mp^s)(t)) = 2 (required >= 2)
PASS [8g] BV3 (19) with m=3, excellent lift q->q^p: simplicial n=2 (c_{3k}=(3k)!/k!^3), p=5, s=2 :: v_p(F(t)F_(mp^(s-1))(t^s) - F(t^s)F_(mp^s)(t)) = 4 (required >= 4)
PASS [8g] BV3 (19) with m=3, excellent lift q->q^p: simplicial n=2 (c_{3k}=(3k)!/k!^3), p=7, s=1 :: v_p(F(t)F_(mp^(s-1))(t^s) - F(t^s)F_(mp^s)(t)) = 2 (required >= 2)
PASS [8g] BV3 (19) with m=4, excellent lift q->q^p: simplicial n=3 (c_{4k}=(4k)!/k!^4), p=5, s=1 :: v_p(F(t)F_(mp^(s-1))(t^s) - F(t^s)F_(mp^s)(t)) = 4 (required >= 3)
PASS [8g] BV3 (19) with m=4, excellent lift q->q^p: simplicial n=3 (c_{4k}=(4k)!/k!^4), p=5, s=2 :: v_p(F(t)F_(mp^(s-1))(t^s) - F(t^s)F_(mp^s)(t)) = 7 (required >= 6)
PASS [8g] BV3 (19) with m=4, excellent lift q->q^p: simplicial n=3 (c_{4k}=(4k)!/k!^4), p=7, s=1 :: v_p(F(t)F_(mp^(s-1))(t^s) - F(t^s)F_(mp^s)(t)) = 3 (required >= 3)
PASS [8g] mirror map: t_0(Q)^3 = t(Q^3) for the simplicial family attached to 27t
PASS [8g] mirror map: t_0(Q)^4 = t(Q^4) for the simplicial family attached to 256t
PASS [8h] for 2F1(1/3,2/3;1|27t), p=5: Tr(Q_1) is not a polynomial in t; its high t-coefficients are nonzero but p-adically small :: (n, v_5 of coefficient): [(5, 2), (11, 5), (17, 8), (23, 11), (29, 14), (35, 17)]
PASS [9a] Zagier D: F(t(q)) = Theta(q), Theta^2 t = eta5^5/eta1 and 1/t - 11 - t = eta1^6/eta5^6 to O(q^79); orders of eta5^5/eta1 at the cusps of X_0(5) are 1, 0
PASS [9a] Zagier D: the canonical coordinate of the Apery zeta(2) equation inverts t(q), i.e. t is the mirror map
PASS [9a] diamond operators on Gamma_1(5): t o<2> = t o<3> = -1/t, Theta|<2> = c t Theta, Theta|<3> = -c t Theta, t o<4> = t, Theta|<4> = -Theta (floating point) :: c = +1.000000000000; t(<2>tau)t(tau) = -1.000e+00-1.491e-16j
PASS [9b] Hecke traces on Gamma_1(5): p=+-1: eta'Q+p^2U_pG polynomial of degree<p(D+1); p=+-2: Lambda = eta Qcheck + p^2 U_pG and Lambda' = -eta t Q + p^2 t U_p Gcheck are Laurent polynomials, deg Lambda' <= p(D+1), Lambda'(t) = Lambda(-1/t) :: p=2,D=0:top(Lc)=2<=2 p=2,D=2:top(Lc)=6<=6 p=2,D=3:top(Lc)=8<=8 p=3,D=0:top(Lc)=3<=3 p=3,D=2:top(Lc)=9<=9 p=3,D=3:top(Lc)=12<=12 p=7,D=0:top(Lc)=7<=7 p=7,D=2:top(Lc)=21<=21 p=7,D=3:top(Lc)=28<=28 p=13,D=0:top(Lc)=13<=13 p=13,D=2:top(Lc)=39<=39 p=13,D=3:top(Lc)=52<=52 ...
PASS [9c] Theorem C: (H_s) and the dual congruence (D_s) modulo p^(2s) for Zagier D :: p=2,s=1,m=1:v(H)=2,v(D)=2 p=2,s=1,m=2:v(H)=4,v(D)=4 p=2,s=2,m=1:v(H)=4,v(D)=4 p=2,s=2,m=2:v(H)=6,v(D)=6 p=2,s=3,m=1:v(H)=6,v(D)=6 p=3,s=1,m=1:v(H)=2,v(D)=2 p=3,s=1,m=2:v(H)=2,v(D)=2 p=3,s=2,m=1:v(H)=4,v(D)=4 p=3,s=2,m=2:v(H)=4,v(D)=4 p=3,s=3,m=1:v(H)=6,v(D)=6 p=7,s=1,m=1:v(H)=2,v(D)=2 p=7,s=1,m=2:v(H)=2,v(D)=2 p=7,s=2,m=1:v(H)=4,v(D)=4 p=7,s=2,m=2:v(H)=4,v(D)=4 p=11,s=1,m=1:v(H)=2,v(D)=2 p=11,s=1,m=2:v(H)=2,v(D)=2 p=11,s=2,m=1:v(H)=4,v(D)=4 p=13,s=1,m=1:v(H)=2,v(D)=2 p=13,s=1,m=2:v(H)=2,v(D)=2 p=17,s=1,m=1:v(H)=2,v(D)=2 p=17,s=1,m=2:v(H)=2,v(D)=2

ALL CHECKS PASSED
\end{lstlisting}


\clearpage
\section{Complete verification programs}\label{app:all-programs}
This appendix embeds the complete programs recovered from the three verification archives, in addition to the base program in Appendix~\ref{B:app:script}. Each is a finite, reproducible computation with its scope stated in the program. The independent certificate appears first. The other programs perform bounded validation of the corresponding general proofs. Recorded JSON arrays can be regenerated from these listings; they are not hidden assumptions of the symbolic theorems.

\subsection{certificate\_octa4\_p5\_s2.py}\label{app:certificate-program}
\noindent Original source: \texttt{certificate\_octa4\_p5\_s2.py}; 156 lines.

\begin{lstlisting}[language=Python]
#!/usr/bin/env python3
"""Exact polynomial certificate for hyperoctahedral n=4, p=5, s=2.

The rational excellent lift is constructed only from BV vertex-polynomial
congruences.  The last polynomial identity then proves the ordinary
truncation congruence at EVERY formal-series coefficient.
Standard library only; no floating point and no power-series cutoff.
"""
from math import comb, factorial
from pathlib import Path
import json
import time
import hashlib

P, R, PREC = 5, 2, 4
N = P**R
MOD = P**PREC
WORKMOD = N*MOD

def trim(a):
    a = list(a)
    while len(a)>1 and not a[-1]:
        a.pop()
    return a

def add(a,b,mod=MOD):
    c=[0]*max(len(a),len(b))
    for i,x in enumerate(a): c[i]+=x
    for i,x in enumerate(b): c[i]+=x
    return trim([x%mod for x in c])

def scale(a,k,mod=MOD): return trim([(k*x)%mod for x in a])

def mul(a,b,mod=MOD):
    a=trim([x%mod for x in a]); b=trim([x%mod for x in b])
    if a==[0] or b==[0]: return [0]
    count=min(len(a),len(b))
    width=(count.bit_length()+2*mod.bit_length()+7)//8
    aa=int.from_bytes(b''.join(x.to_bytes(width,'little') for x in a),'little')
    bb=int.from_bytes(b''.join(x.to_bytes(width,'little') for x in b),'little')
    length=len(a)+len(b)-1
    raw=(aa*bb).to_bytes(length*width,'little')
    return trim([int.from_bytes(raw[i*width:(i+1)*width],'little')%mod for i in range(length)])

def power(a,k,mod=MOD):
    out=[1]
    while k:
        if k&1: out=mul(out,a,mod)
        k//=2
        if k: a=mul(a,a,mod)
    return out

def subst_power(a,p):
    out=[0]*(p*(len(a)-1)+1)
    for j,x in enumerate(a): out[p*j]=x
    return trim(out)

def taylor_scale(a,r,p=P,mod=MOD):
    # [h^r] a(u^p(1+h)).
    return subst_power([(x*comb(j,r))%mod if j>=r else 0 for j,x in enumerate(a)],p)

def b3_sequence(length):
    # Sum over three composition indices of multinomial(n; k)^2.
    return [sum(comb(j,k)**2*comb(2*k,k) for k in range(j+1)) for j in range(length)]

def vertex(K,b3,mod):
    # P_{2K}(u)=u^{-K} A_K(u), the negative of the dominant-vertex coefficient.
    fac=[factorial(j) for j in range(2*K)]
    ans=[]
    for j in range(K):
        total=0
        for ell in range(j+1):
            den=fac[2*K-2*j-1]*fac[ell]*fac[j-ell]**2
            num=fac[2*K-ell-1]
            assert num%den==0
            term=(num//den)*b3[j-ell]
            total+=(-1 if ell&1 else 1)*term
        ans.append(total%mod)
    assert ans[0]==1
    return trim(ans)

def oct4(length):
    return [comb(2*j,j)*sum(comb(j,k)**2*comb(2*k,k)*comb(2*(j-k),j-k)
                            for k in range(j+1)) for j in range(length)]

def build():
    started=time.time()
    b3=b3_sequence(2*P*N)
    av={k:vertex(k,b3,WORKMOD) for k in (N,2*N,P*N,2*P*N)}
    a=av[N]; b=av[2*N]; ap=av[P*N]; bp=av[2*P*N]
    ar=[taylor_scale(a,r,mod=WORKMOD) for r in range(4)]
    br=[taylor_scale(b,r,mod=WORKMOD) for r in range(4)]
    a0powers=[power(ar[0],r,WORKMOD) for r in range(5)]
    # T_r/ar[0]^(r+1) = [h^r] C_N(u^p(1+h)), C_N=A_2N/A_N.
    T=[]
    for r in range(4):
        tr=mul(br[r],a0powers[r],WORKMOD)
        for i in range(1,r+1):
            tr=add(tr,scale(mul(mul(ar[i],T[r-i],WORKMOD),a0powers[i-1],WORKMOD),-1,WORKMOD),WORKMOD)
        T.append(tr)
    # A,B,C,D share the unit denominator ap*ar[0]^4 after division by N.
    ABCD=[]
    for r in range(4):
        numerator=add(scale(mul(bp,a0powers[4],WORKMOD),comb(N,r),WORKMOD),
                      scale(mul(mul(ap,T[r],WORKMOD),a0powers[3-r],WORKMOD),-1,WORKMOD),WORKMOD)
        assert all(x%N==0 for x in numerator), ('N divisibility',r)
        ABCD.append(trim([(x//N)%MOD for x in numerator]))
    A,B,C,D=ABCD
    assert A[0]==0 and all(x%P==0 for x in A)
    assert B[0]==1
    assert C[0]==12 and D[0]==92

    # Independent target truncations enter ONLY at this point.
    f=oct4(P*P)
    fsmall=f[:P]
    L=[]
    for r in range(4):
        lhs=taylor_scale(fsmall,r)
        if r>=1:
            lhs=add(lhs,scale(taylor_scale(fsmall,r-1),N))
        L.append(add(mul(ap,lhs),scale(mul(f,ar[r]),-1)))

    # h=S-1 = -A/B - C A^2/B^3 + (D/B^4-2 C^2/B^5) A^3 (mod 625).
    # Substitute into L0+L1*h+L2*h^2+L3*h^3 and clear the unit B^5.
    Bpow=[power(B,r) for r in range(6)]
    A2=mul(A,A); A3=mul(A2,A)
    target=add(mul(L[0],Bpow[5]),scale(mul(mul(L[1],A),Bpow[4]),-1))
    second=add(mul(L[2],Bpow[3]),scale(mul(mul(L[1],C),Bpow[2]),-1))
    target=add(target,mul(second,A2))
    third=add(mul(mul(L[1],D),B),scale(mul(L[1],mul(C,C)),-2))
    third=add(third,scale(mul(mul(L[2],C),B),2))
    third=add(third,scale(mul(L[3],Bpow[2]),-1))
    target=add(target,mul(third,A3))
    assert target==[0], ('target not identically zero',[(i,x) for i,x in enumerate(target) if x][:10])

    payload={'p':P,'s':2,'dimension':4,'modulus':MOD,'vertex_precision':WORKMOD,
             'vertex_indices':list(av),'vertex_polynomials':av,
             'lift_coefficients_common_denominator':ABCD,
             'truncation_target_coefficients':L}
    serialized=json.dumps(payload,sort_keys=True,separators=(',',':'))
    report={'theorem':'Every coefficient of E_2 for hyperoctahedral dimension four is divisible by 5^4.',
            'all_polynomial_coefficients_checked':True,
            'power_series_cutoff':None,
            'lift_constructed_from':'BV vertex polynomials only; no target truncation assumption',
            'p':P,'s':2,'modulus':MOD,'intermediate_modulus':WORKMOD,
            'rational_lift_polynomial_degrees':[len(a)-1 for a in ABCD],
            'cleared_identity_degree_bound':max(len(L[r])-1+5*(len(B)-1) for r in range(4)),
            'cleared_identity_zero':target==[0],
            'certificate_sha256':hashlib.sha256(serialized.encode()).hexdigest(),
            'elapsed_seconds':round(time.time()-started,3)}
    base=Path(__file__).resolve().parent
    (base/'certificate_octa4_p5_s2.json').write_text(json.dumps(payload,indent=2)+'\n')
    (base/'certificate_octa4_p5_s2_report.json').write_text(json.dumps(report,indent=2)+'\n')
    print(json.dumps(report,indent=2))

if __name__=='__main__': build()
\end{lstlisting}

\paragraph{Recorded certificate result.} The archived report is reproduced below. The identity and its degree bound can be regenerated by the preceding program.
\begin{lstlisting}
{
  "theorem": "Every coefficient of E_2 for hyperoctahedral dimension four is divisible by 5^4.",
  "all_polynomial_coefficients_checked": true,
  "power_series_cutoff": null,
  "lift_constructed_from": "BV vertex polynomials only; no target truncation assumption",
  "p": 5,
  "s": 2,
  "modulus": 625,
  "intermediate_modulus": 15625,
  "rational_lift_polynomial_degrees": [
    489,
    537,
    537,
    527
  ],
  "cleared_identity_degree_bound": 2819,
  "cleared_identity_zero": true,
  "certificate_sha256": "d7dbc0fd7427660d97e9d8cacf26577f0af474f6ac447c8ad1cf2f78d69bcd6b",
  "elapsed_seconds": 0.642
}
\end{lstlisting}

\subsection{verify\_elliptic.py}\label{app:program-1}
\noindent Original source: \texttt{verify\_elliptic.py}; 431 lines.

\begin{lstlisting}[language=Python]
#!/usr/bin/env python3
"""Finite exact-arithmetic checks of the elliptic-trace proof.

This is supporting verification, not a replacement for the all-s proof.
All power-series lengths and arithmetic precisions are recorded.  The
full-congruence checks always include terms strictly beyond the cutoff.
Standard-library Python only.  Run from any directory.
"""
from __future__ import annotations

import json
from math import comb, isqrt
from pathlib import Path
from time import perf_counter


def trim(a):
    a = list(a)
    while len(a) > 1 and a[-1] == 0:
        a.pop()
    return a


def zmul(a, b, n):
    """Exact integer convolution, truncated at q^n."""
    out = [0] * n
    for i, x in enumerate(a[:n]):
        if x:
            for j, y in enumerate(b[:n-i]):
                if y:
                    out[i+j] += x*y
    return out


def mmul(a, b, n, mod):
    """Exact modular convolution by carry-free Kronecker substitution."""
    if n <= 0:
        return []
    a = trim([x % mod for x in a[:n]])
    b = trim([x % mod for x in b[:n]])
    if not a or not b or a == [0] or b == [0]:
        return [0]*n
    count = min(len(a), len(b))
    width = max(1, (count.bit_length() + 2*mod.bit_length() + 7)//8)
    aa = int.from_bytes(b''.join(x.to_bytes(width, 'little') for x in a), 'little')
    bb = int.from_bytes(b''.join(x.to_bytes(width, 'little') for x in b), 'little')
    # Each convolution coefficient is < count*mod^2 < 256^width.
    size = len(a)+len(b)-1
    raw = (aa*bb).to_bytes(size*width, 'little')
    out = [int.from_bytes(raw[i*width:(i+1)*width], 'little') % mod
           for i in range(min(n, size))]
    return out+[0]*(n-len(out))


def zinv(a, n):
    assert a[0] == 1
    b = [1]+[0]*(n-1)
    for k in range(1, n):
        b[k] = -sum(a[j]*b[k-j] for j in range(1, min(k+1, len(a))))
    return b


def minv(a, n, mod):
    assert a[0] % mod == 1
    b = [1]
    while len(b) < n:
        size = min(2*len(b), n)
        ab = mmul(a[:size], b, size, mod)
        correction = [(-x) % mod for x in ab]
        correction[0] = (correction[0]+2) % mod
        b = mmul(b, correction, size, mod)
    return b


def mpow(a, k, n, mod):
    out = [1]+[0]*(n-1)
    while k:
        if k & 1:
            out = mmul(out, a, n, mod)
        k //= 2
        if k:
            a = mmul(a, a, n, mod)
    return out


def vsub(a, p, n):
    out = [0]*n
    for k, x in enumerate(a):
        if p*k >= n:
            break
        out[p*k] = x
    return out


def mcompose(a, t, n, mod):
    """Baby-step/giant-step polynomial composition; t(0)=0."""
    assert t[0] == 0
    a = a[:n]
    if not a:
        return [0]*n
    block = isqrt(len(a))+1
    powers = [[1]+[0]*(n-1)]
    for _ in range(block):
        powers.append(mmul(powers[-1], t, n, mod))
    out = [0]*n
    for start in reversed(range(0, len(a), block)):
        out = mmul(out, powers[block], n, mod)
        for r, x in enumerate(a[start:start+block]):
            x %= mod
            if x:
                out = [(v+x*u) % mod for v, u in zip(out, powers[r])]
    return out


def texpand(f, t, n, mod=None):
    """Recover the t-series from its q-series, exactly or modulo mod."""
    assert t[0] == 0 and t[1] == 1
    rem = list(f[:n])+[0]*max(0, n-len(f))
    powers = [1]+[0]*(n-1)
    out = []
    for k in range(n):
        c = rem[k] if mod is None else rem[k] % mod
        out.append(c)
        if c:
            for i in range(k, n):
                rem[i] -= c*powers[i]
                if mod is not None:
                    rem[i] %= mod
        powers = zmul(powers, t, n) if mod is None else mmul(powers, t, n, mod)
    return out


def valuation(x, p, cap):
    if x == 0:
        return cap
    v = 0
    while v < cap and x % p == 0:
        v += 1
        x //= p
    return v


def character3(n):
    return 0 if n % 3 == 0 else (1 if n % 3 == 1 else -1)


def divisor_sums(n):
    out = [0]*n
    for d in range(1, n):
        for j in range(d, n, d):
            out[j] += d
    return out


def eta_h(level, n, sig):
    """Exact q*(eta_N/eta_1 product) via its logarithmic derivative."""
    exponent = 24//(level-1)
    logderiv = [0]+[exponent*(sig[j]-(level*sig[j//level] if j % level == 0 else 0))
                         for j in range(1, n)]
    h = [1]+[0]*(n-1)
    for k in range(1, n):
        value = sum(logderiv[j]*h[k-j] for j in range(1, k+1))
        assert value % k == 0
        h[k] = value//k
    return [0]+h[:n-1]


def theta_h(n):
    # Exact Lambert series.  No finite lattice-radius approximation.
    out = [1]+[0]*(n-1)
    for d in range(1, n):
        c = 6*character3(d)
        for j in range(d, n, d):
            out[j] += c
    return out


class Family:
    def __init__(self, label, weight, e, level, scale, target):
        self.label, self.w, self.e, self.level, self.scale, self.target = (
            label, weight, e, level, scale, target)

    def epsilon(self, p):
        return character3(p) if self.w == 1 else 1

    def coefficient(self, k):
        if self.label == '27t':
            return comb(3*k, k)*comb(2*k, k)
        if self.label == '256t':
            return comb(4*k, k)*comb(3*k, k)*comb(2*k, k)
        return comb(3*k, k)*comb(2*k, k)**2

    def parameters(self, n, mod=None):
        h = H[self.level][:n]
        den = [1]+[self.scale*x for x in h[1:]]
        if mod is None:
            invden = zinv(den, n)
            factor = invden if self.w == 1 else zmul(invden, invden, n)
            t = zmul(h, factor, n)
        else:
            invden = minv(den, n, mod)
            factor = invden if self.w == 1 else mmul(invden, invden, n, mod)
            t = mmul(h, factor, n, mod)
        if self.label == '27t':
            th = TH[:n]
        elif self.label == '256t':
            th = [1]+[24*(SIG[j]-(2*SIG[j//2] if j % 2 == 0 else 0))
                           for j in range(1, n)]
        else:
            th = [1]+[12*(SIG[j]-(3*SIG[j//3] if j % 3 == 0 else 0))
                           for j in range(1, n)]
        if mod is not None:
            th = [x % mod for x in th]
        return th, t


FAMILIES = [Family('27t', 1, 3, 3, 27, 2),
            Family('256t', 2, 4, 2, 64, 3),
            Family('108t', 2, 3, 3, 27, 3)]


def full_congruence(fam, p, s, m):
    cutoff = m*p**s
    lower = m*p**(s-1)
    n = 2*cutoff+9
    target = fam.target*s
    prec = target+2
    mod = p**prec
    th, t = fam.parameters(n, mod)
    coeff = [fam.coefficient(k) % mod for k in range(cutoff)]
    big = mcompose(coeff, t, n, mod)
    smalln = (n-1)//p+1
    small = mcompose(coeff[:lower], t[:smalln], smalln, mod)
    lhs = mmul(th, vsub(small, p, n), n, mod)
    rhs = mmul(vsub(th, p, n), big, n, mod)
    error = [(x-y) % mod for x, y in zip(lhs, rhs)]
    v = min(valuation(x, p, prec) for x in error)
    assert n > cutoff and v >= target, (fam.label, p, s, m, v)
    return {'family': fam.label, 'p': p, 's': s, 'm': m,
            'cutoff': cutoff, 'q_coefficients_checked': n,
            'required_valuation': target, 'arithmetic_precision': prec,
            'observed_minimum_valuation': v,
            'valuation_is_lower_bound': v == prec, 'passed': True}


def monomial_trace(fam, p, k, prec=6):
    a, delta = p//fam.e, (p+1)//fam.e
    n = p*k+a+delta*(prec-1)+9
    length = p*n
    mod = p**prec
    th, t = fam.parameters(length, mod)
    tk = mpow(t, k, length, mod)
    G = mmul(mmul(vsub(th, p, length), minv(th, length, mod), length, mod),
             tk, length, mod)
    Q = mmul(mmul(th[:n], minv(vsub(th, p, n), n, mod), n, mod),
             vsub(tk, p, n), n, mod)
    trq = [(Q[j]+fam.epsilon(p)*p**(fam.w+1)*G[p*j]) % mod for j in range(n)]
    trt = texpand(trq, t[:n], n, mod)
    for j, c in enumerate(trt):
        bound = max(0, -(-(j-p*k-a)//delta))
        assert c % p**min(bound, prec) == 0, (fam.label, p, k, j, bound, c)
    degrees = []
    for r in range(1, prec+1):
        observed = max((j for j, c in enumerate(trt) if c % p**r), default=-1)
        expected = p*k+a+delta*(r-1)
        assert observed <= expected
        degrees.append({'precision': r, 'observed_degree_in_checked_range': observed,
                        'proved_bound': expected})
    return {'family': fam.label, 'p': p, 'monomial_k': k,
            't_coefficients_checked': n, 'input_q_coefficients': length,
            'arithmetic_precision': prec, 'degrees': degrees, 'passed': True}


def elementary_from_power_sums(sums, count, n):
    elementary = [[1]+[0]*(n-1)]
    for j in range(1, count+1):
        acc = [0]*n
        for r in range(1, j+1):
            value = zmul(elementary[j-r], sums[r], n)
            sign = 1 if r % 2 else -1
            acc = [x+sign*y for x, y in zip(acc, value)]
        assert all(x % j == 0 for x in acc)
        elementary.append([x//j for x in acc])
    return elementary


def exact_norm_polynomials(fam, p):
    """Exact integer q-series, not modular arithmetic, for norm lemmas."""
    n = p+10
    length = p*n
    th, t = fam.parameters(length)
    thshort, tshort = th[:n], t[:n]
    thinv = zinv(thshort, n)
    invpowers = [[1]+[0]*(n-1)]
    for j in range(p+1):
        invpowers.append(zmul(invpowers[-1], thinv, n))

    thpower = [1]+[0]*(length-1)
    theta_sums = {0: None}
    for r in range(1, p+1):
        thpower = zmul(thpower, th, length)
        theta_sums[r] = [p*thpower[p*j] for j in range(n)]
    elementary_theta = elementary_from_power_sums(theta_sums, p, n)
    Dq = zmul(zmul(vsub(thshort, p, n), elementary_theta[p], n), invpowers[p+1], n)
    Dt = texpand(Dq, tshort, n)
    delta = (p+1)//fam.e
    assert Dt[0] == 1 and all(x == 0 for x in Dt[delta+1:])
    assert all(x % p == 0 for x in Dt[1:])

    inner = []
    for j in range(p+1):
        c = p**(fam.w*(p-j))*(-fam.epsilon(p))**j
        inner.append([c*x for x in zmul(elementary_theta[j], invpowers[j], n)])
    ratio = zmul(vsub(thshort, p, n), thinv, n)
    Tq = []
    for j in range(p+2):
        left = inner[j] if j <= p else [0]*n
        right = zmul(ratio, inner[j-1], n) if j else [0]*n
        item = [x-y for x, y in zip(left, right)]
        Tq.append(item)
        coefficients = texpand(item, tshort, n)
        assert all(x == 0 for x in coefficients[j//fam.e+1:])
    assert Tq[-1] == [fam.epsilon(p)*x for x in Dq]

    # Verify the actual character factor in R_0=-T_p/(epsilon D).
    ratio_full = zmul(vsub(th, p, length), zinv(th, length), length)
    Q0 = zmul(thshort, zinv(vsub(thshort, p, n), n), n)
    R0 = [Q0[j]+fam.epsilon(p)*p**(fam.w+1)*ratio_full[p*j] for j in range(n)]
    check = zmul(Dq, R0, n)
    assert all(Tq[p][j]+fam.epsilon(p)*check[j] == 0 for j in range(n))

    tpower = [1]+[0]*(length-1)
    t_sums = {0: None}
    for r in range(1, p+2):
        tpower = zmul(tpower, t, length)
        V = vsub(tpower, p, n)
        t_sums[r] = [p*tpower[p*j]+V[j] for j in range(n)]
    elementary_t = elementary_from_power_sums(t_sums, p+1, n)
    Depower = [1]+[0]*(n-1)
    for _ in range(fam.e):
        Depower = zmul(Depower, Dq, n)
    for j in range(p+2):
        Aq = [(-1)**j*x for x in zmul(Depower, elementary_t[j], n)]
        At = texpand(Aq, tshort, n)
        assert all(x == 0 for x in At[p+2:])
        if j == 1:
            assert all((x+(1 if k == p else 0)) % p == 0 for k, x in enumerate(At))
    return {'family': fam.label, 'p': p, 'epsilon': fam.epsilon(p),
            'q_coefficients_checked': n, 'input_q_coefficients': length,
            'arithmetic': 'exact integers', 'norm_D_coefficients': Dt[:delta+1],
            'T_coefficient_degree_checks': p+2,
            'parameter_norm_coefficient_degree_checks': p+2,
            'A1_frobenius_congruence': True, 'trace_character_check': True,
            'passed': True}


def implementation_checks():
    # Independently compare packed multiplication with direct convolution.
    for mod in (25, 7**5, 11**8):
        for n in (5, 31, 73):
            a = [(j*j+13*j+7) % mod for j in range(n)]
            b = [(-j**3+17*j-9) % mod for j in range(n)]
            assert mmul(a, b, n, mod) == [x % mod for x in zmul(a, b, n)]
            a[0] = 1
            assert minv(a, n, mod) == [x % mod for x in zinv(a, n)]
    # Check the composition routine against a short direct Horner computation.
    mod, n = 5**8, 37
    t = [0, 1]+[(j*j+3) % mod for j in range(2, n)]
    a = [j*j+1 for j in range(18)]
    h = [0]*n
    for c in reversed(a):
        h = mmul(h, t, n, mod)
        h[0] = (h[0]+c) % mod
    assert mcompose(a, t, n, mod) == h


def main():
    global SIG, TH, H
    start = perf_counter()
    implementation_checks()
    maximum = 650
    SIG = divisor_sums(maximum)
    TH = theta_h(maximum)
    H = {level: eta_h(level, maximum, SIG) for level in (2, 3)}
    output = {'scope': 'Finite exact checks supporting the written proof; not a proof by testing.',
              'implementation_checks': True, 'parametrizations': [],
              'norm_polynomials': [], 'monomial_filtration': [], 'full_congruences': []}
    for fam in FAMILIES:
        n, mod = 80, 5**8
        th, t = fam.parameters(n, mod)
        assert mcompose([fam.coefficient(k) for k in range(n)], t, n, mod) == th
        output['parametrizations'].append({'family': fam.label, 'q_coefficients_checked': n,
                                           'modulus': mod, 'passed': True})
        for p in (5, 7):
            result = exact_norm_polynomials(fam, p)
            output['norm_polynomials'].append(result)
            print('PASS norm', fam.label, p, 'D =', result['norm_D_coefficients'], flush=True)
        for p in (5, 7):
            for k in (0, 1, 3):
                output['monomial_filtration'].append(monomial_trace(fam, p, k))
        if fam.label == '108t':
            output['monomial_filtration'].append(monomial_trace(fam, 5, 4, 7))
        for p in (5, 7, 11):
            for s in (1, 2):
                for m in (1, 2):
                    result = full_congruence(fam, p, s, m)
                    output['full_congruences'].append(result)
                    print('PASS congruence', fam.label, 'p', p, 's', s, 'm', m,
                          'valuation', result['observed_minimum_valuation'],
                          'checked q degrees 0..', result['q_coefficients_checked']-1, flush=True)
    output['seconds'] = round(perf_counter()-start, 3)
    output['counts'] = {k: len(output[k]) for k in
                        ('parametrizations', 'norm_polynomials', 'monomial_filtration', 'full_congruences')}
    output['all_passed'] = True
    destination = Path(__file__).resolve().parent
    (destination/'elliptic_validation_results.json').write_text(json.dumps(output, indent=2)+'\n')
    summary = [
        'ELLIPTIC TRACE VALIDATION: ALL FINITE CHECKS PASSED',
        'These checks support the written proof; they do not prove the infinite statement.',
        'Corrected theta_H is computed by its exact Lambert series.',
        'All modular polynomial degree checks use exact integers.',
        'All full-congruence checks include coefficients beyond the truncation cutoff.',
        'Counts: '+json.dumps(output['counts'], sort_keys=True),
        'Elapsed seconds: '+str(output['seconds']),
    ]
    (destination/'elliptic_validation_results.txt').write_text('\n'.join(summary)+'\n')
    print('\n'.join(summary), flush=True)


if __name__ == '__main__':
    main()
\end{lstlisting}

\subsection{verify\_sharpness.py}\label{app:program-2}
\noindent Original source: \texttt{verify\_sharpness.py}; 171 lines.

\begin{lstlisting}[language=Python]
"""Bounded exact-arithmetic validation of the all-s proofs in Exceptional_Excellent_Lift_Supercongruences.tex, Section 4.
This script is not a substitute for those proofs and does not claim a full-series theorem.
Python 3 standard library only; run from any current directory.
"""
from fractions import Fraction as Q
from math import comb, factorial
import json
from pathlib import Path


def vp(n,p):
    if n==0:return None
    v=0
    while n%p==0:n//=p;v+=1
    return v

def central(n):return comb(2*n,n)
def oct3(n):return central(n)*sum(comb(n,k)**2*central(k) for k in range(n+1))
def domb(n):return sum(comb(n,k)**2*central(k)*central(n-k) for k in range(n+1))
def oct4(n):return central(n)*domb(n)
def cube(d,n):return central(n)**d
def simplex(r,n):return factorial(r*n)//factorial(n)**r

def C_r(r):
    H=sum(Q(1,j) for j in range(1,r+1))
    H2=sum(Q(1,j*j) for j in range(1,r+1))
    C=Q(factorial(r)**2,2)*((r*(H-1))**2-r*r*H2+r)
    assert C.denominator==1
    return int(C)

out=[]
# The dimension-three exact refinement, including s=1,2,3 at p=5.
for p in [5,7,11]:
    h=(p-1)//2
    S=sum(central(k)*pow(k,-2,p) for k in range(1,h+1))%p
    chi=1 if pow(p-3,h,p)==1 else -1
    for s in [1,2,3]:
        K=p**s
        z=oct3(K)-oct3(K//p)
        predicted=2*pow(chi,s-1)*S%p
        assert z%p**(2*s)==0
        actual=z//p**(2*s)%p
        assert actual==predicted
        out.append(dict(family='octahedral3',p=p,s=s,index=K,valuation=vp(z,p),unit_mod_p=actual,predicted=predicted))
# Hyperoctahedral dimension four (uses the established Domb congruence in the proof).
for p in [5,7]:
    for s in [1,2,3]:
        K=p**s
        z=oct4(K+1)-(8+20*K)*oct4(K//p)
        assert z%p**(2*s)==0
        assert z//p**(2*s)%p==32%p
        out.append(dict(family='octahedral4',p=p,s=s,index=K+1,valuation=vp(z,p),unit_mod_p=z//p**(2*s)%p,predicted=32%p))
# General hypercubic formula in two dimensions of interest and multiple primes.
for d,p in [(4,5),(4,7),(5,7),(5,11),(6,7)]:
    for s in [1,2,3]:
        K=p**s
        z=cube(d,K+1)-(2**d+K*d*2**d)*cube(d,K//p)
        predicted=2**(2*d-1)*d*(d-3)%p
        assert z%p**(2*s)==0
        assert z//p**(2*s)%p==predicted
        assert vp(z,p)==2*s
        out.append(dict(family='cube',dimension=d,p=p,s=s,index=K+1,valuation=vp(z,p),unit_mod_p=z//p**(2*s)%p,predicted=predicted))
# Simplex dimensions four and five, including the genuine exceptional p=7 when r=6.
for r,p in [(5,7),(5,11),(6,7),(6,11)]:
    H=sum(Q(1,j) for j in range(1,r+1))
    alpha=factorial(r)*r*(H-1)
    assert alpha.denominator==1
    for s in [1,2,3]:
        K=p**s
        z=int(simplex(r,K+1)-(factorial(r)+K*alpha)*simplex(r,K//p))
        predicted=C_r(r)%p
        assert z%p**(2*s)==0
        assert z//p**(2*s)%p==predicted
        if predicted:assert vp(z,p)==2*s
        out.append(dict(family='simplex',dimension=r-1,p=p,s=s,index=K+1,valuation=vp(z,p),unit_mod_p=z//p**(2*s)%p,predicted=predicted))
# Check Domb recurrence used in the optional proof directly from its binomial definition.
for k in range(2,16):
    assert k**3*domb(k)==2*(2*k-1)*(5*k*k-5*k+2)*domb(k-1)-64*(k-1)**3*domb(k-2)
# Compare first-jet formula with full formal-series computation at bounded sizes.
# Independent formal-series implementation (no local imports).
from fractions import Fraction as Q
from math import comb, factorial
import json

def mul(a,b,N):
 c=[Q(0)]*N
 for i,x in enumerate(a[:N]):
  if x:
   for j,y in enumerate(b[:N-i]):
    if y:c[i+j]+=x*y
 return c

def div(a,b,N):
 c=[Q(0)]*N
 for k in range(N):c[k]=(a[k] if k<len(a) else 0)-sum(b[j]*c[k-j] for j in range(1,min(k+1,len(b))))
 return c

def exp(g,N):
 c=[Q(1)]+[Q(0)]*(N-1)
 for k in range(1,N):c[k]=sum(j*g[j]*c[k-j] for j in range(1,min(k+1,len(g))))/k
 return c

def param(n,M,kind):
 H=[Q(0)]
 for k in range(1,(n+1)*M+1):H.append(H[-1]+Q(1,k))
 if kind=='cube':
  F=[comb(2*k,k)**n for k in range(M)]
  G=[2*n*F[k]*(H[2*k]-H[k]) for k in range(M)]
 elif kind=='simplex':
  F=[factorial((n+1)*k)//factorial(k)**(n+1) for k in range(M)]
  G=[(n+1)*F[k]*(H[(n+1)*k]-H[k]) for k in range(M)]
 elif kind=='octa':
  B=[1]*M
  for d in range(n-2):B=[sum(comb(k,j)**2*B[k-j] for j in range(k+1)) for k in range(M)]
  Bn=[sum(comb(k,j)**2*B[k-j] for j in range(k+1)) for k in range(M)]
  S=[sum(comb(k,j)**2*B[k-j]*H[j] for j in range(k+1)) for k in range(M)]
  F=[comb(2*k,k)*Bn[k] for k in range(M)]
  G=[2*comb(2*k,k)*(H[2*k]*Bn[k]-S[k]) for k in range(M)]
 return F,div(G,F,M)

def sigma(g,p,M):
 # Lagrange inversion T_j=1/j * [u^{j-1}] exp(-j*g(u)).
 J=(M-1)//p
 T=[Q(0)]+[exp([-j*x for x in g[:j]],j)[j-1]/j for j in range(1,J+1)]
 w=[Q(0)]*p+exp([p*x for x in g],M-p)
 c=[Q(0)]*M
 wp=[Q(1)]+[Q(0)]*(M-1)
 for j in range(1,J+1):
  wp=mul(wp,w,M)
  for k in range(M):c[k]+=T[j]*wp[k]
 return c

def compose(a,b,M):
 c=[Q(0)]*M
 for x in reversed(a[:M]):
  c=mul(c,b,M);c[0]+=x
 return c

def val(x,p):
 if x==0:return 999999
 if not isinstance(x,Q):x=Q(x)
 n,d=x.numerator,x.denominator;v=0
 while n%p==0:n//=p;v+=1
 while d%p==0:d//=p;v-=1
 return v

def check(n,p,s,kind,M):
 F,g=param(n,M,kind)
 sig=sigma(g,p,M)
 A=mul(F,compose(F[:p**(s-1)],sig,M),M)
 B=mul(F[:p**s],compose(F[:(M-1)//p+1],sig,M),M)
 E=[a-b for a,b in zip(A,B)]
 v=min(val(x,p) for x in E)
 i=next(i for i,x in enumerate(E) if val(x,p)==v)
 return {'kind':kind,'n':n,'p':p,'s':s,'M':M,'v':v,'first_min_index':i,'coefficient':str(E[i]),'q2':str(g[1]),'lift_integral':all(x.denominator==1 for x in sig)}


for kind,d,p,s in [('octa',3,5,1),('octa',3,5,2),('octa',4,5,1),('cube',4,5,1),('simplex',4,7,1)]:
    K=p**s;M=K+3
    F,g=param(d,M,kind);sig=sigma(g,p,M)
    lhs=mul(F,compose(F[:K//p],sig,M),M)
    rhs=mul(F[:K],compose(F[:(M-1)//p+1],sig,M),M)
    assert lhs[K]-rhs[K]==F[K]-F[K//p]
    assert lhs[K+1]-rhs[K+1]==F[K+1]-(F[1]+K*g[1])*F[K//p]

path=Path(__file__).with_name('sharpness_validation.json')
path.write_text(json.dumps({'description':'Finite exact-arithmetic validation; all-s proofs are in Exceptional_Excellent_Lift_Supercongruences.tex, Section 4','checks':out,'count':len(out)},indent=2)+'\n')
print('PASS:',len(out),'refinement checks; 14 Domb recurrence checks; 10 formal first-jet identities')
print('Constants:',[(r,C_r(r)) for r in range(2,9)])
for x in out:
    if x.get('family') in ['octahedral3','simplex'] and x['p'] in [5,7]:print(json.dumps(x))
\end{lstlisting}

\subsection{verify\_zagier\_d.py}\label{app:program-3}
\noindent Original source: \texttt{verify\_zagier\_d.py}; 130 lines.

\begin{lstlisting}[language=Python]
#!/usr/bin/env python3
"""Independent finite checks for both Zagier D theorems.
All arithmetic is exact. These checks support, but do not replace, the proof.
"""
from math import comb

def mul(a, b, N, modulus=None):
    c = [0]*N
    for i, x in enumerate(a[:N]):
        if x:
            for j, y in enumerate(b[:N-i]):
                if y:
                    c[i+j] += x*y
    if modulus:
        c = [x % modulus for x in c]
    return c

def inv(a, N, modulus=None):
    assert a[0] == 1
    b = [1]+[0]*(N-1)
    for n in range(1,N):
        b[n] = -sum(a[k]*b[n-k] for k in range(1,min(n+1,len(a))))
        if modulus:
            b[n] %= modulus
    return b

def poly(a, t, N, modulus=None):
    b=[0]*N
    for x in a[::-1]:
        b=mul(b,t,N,modulus)
        b[0]=(b[0]+x) if modulus is None else (b[0]+x)%modulus
    return b

def dilation(a,p,N):
    b=[0]*N
    for n,x in enumerate(a):
        if n*p>=N: break
        b[n*p]=x
    return b

def data(N, modulus=None):
    b=[1]+[0]*(N-1)
    for k in range(1,N):
        e=0 if k%5==0 else 5 if k%5 in (1,4) else -5
        for _ in range(abs(e)):
            if e>0:
                for j in range(N-1,k-1,-1):
                    b[j]-=b[j-k]
                    if modulus:b[j]%=modulus
            else:
                for j in range(k,N):
                    b[j]+=b[j-k]
                    if modulus:b[j]%=modulus
    t=[0]+b[:-1]
    theta=[1]+[0]*(N-1)
    ct=[0,3,1,-1,-3]
    for k in range(1,N):
        for j in range(k,N,k):
            theta[j]+=ct[k%5]
    if modulus:theta=[x%modulus for x in theta]
    return t,theta

def apery(M):
    a=[1]
    if M==1:return a
    a.append(3)
    for n in range(2,M):
        b=(11*n*n-11*n+3)*a[n-1]+(n-1)**2*a[n-2]
        assert b%(n*n)==0
        a.append(b//(n*n))
    return a

def t_expansion(f,t,N):
    b=list(f[:N]);out=[];power=[1]+[0]*(N-1)
    for k in range(N):
        c=b[k];out.append(c)
        for j in range(k,N):b[j]-=c*power[j]
        power=mul(power,t,N)
    return out

def v_p(x,p):
    if not x:return 10**9
    v=0
    while x%p==0:v+=1;x//=p
    return v

def check_trace(p,P):
    eps=1 if p%5==1 else -1
    bound=p*(len(P)-1)+p-1
    L=bound+9
    N=p*(L-1)+1
    t,theta=data(N)
    theta_p=dilation(theta,p,N)
    Q=mul(theta,mul(poly(P,dilation(t,p,N),N),inv(theta_p,N),N),N)
    G=mul(theta_p,mul(poly(P,t,N),inv(theta,N),N),N)
    trace=[Q[n]+eps*p*p*G[p*n] for n in range(L)]
    coefficients=t_expansion(trace,t,L)
    assert all(x==0 for x in coefficients[bound+1:])
    print("TRACE PASS",{"p":p,"P":P,"proved_degree_bound":bound,
                       "finite_verified_through_t_degree":L-1,
                       "constant_coefficient":coefficients[0]},flush=True)

def check_congruence(p,s,m,N):
    K=m*p**s;L=K//p
    assert N>K
    modulus=p**(2*s+1)
    t,theta=data(N,modulus)
    A=apery(K)
    lhs=mul(theta,poly(A[:L],dilation(t,p,N),N,modulus),N,modulus)
    rhs=mul(dilation(theta,p,N),poly(A,t,N,modulus),N,modulus)
    differences=[(x-y)%modulus for x,y in zip(lhs,rhs)]
    minimum=min(v_p(x,p) for x in differences)
    assert minimum>=2*s
    print("CONGRUENCE PASS",{"p":p,"s":s,"m":m,"through_q_degree":N-1,
                            "minimum_valuation_capped_at":2*s+1,
                            "observed_minimum":minimum if minimum<10**9 else ">="+str(2*s+1)},
          flush=True)

if __name__=="__main__":
    t,theta=data(65)
    assert poly(apery(65),t,65)==theta
    print("MODULAR PARAMETRIZATION CHECK PASS through q^64",flush=True)
    for p,P in [(11,[1]),(19,[1]),(11,[2,-3])]:
        check_trace(p,P)
    for values in [(11,1,1,45),(19,1,1,55),(29,1,1,70),(31,1,1,75),
                   (11,1,2,60),(19,1,2,85),(11,2,1,170),(19,2,1,410),
                   (11,2,2,300), (2,1,1,22),(2,1,2,25),
                   (3,1,1,25),(3,1,2,30),(7,1,2,45)]:
        check_congruence(*values)
\end{lstlisting}

\subsection{verify\_octahedral\_interpolation.py}\label{app:program-4}
\noindent Original source: \texttt{verify\_octahedral\_interpolation.py}; 421 lines.

\begin{lstlisting}[language=Python]
"""Bounded independent checks for the all-dimensional octahedral proof.

Standard library only.  Besides checking interpolation and direct defects,
this script constructs the proof's row operator D from a truncated restricted
unit-factorial series and the symmetric Bessel operator reduction.  Thus its
row checks do not merely solve the desired identities by interpolation.
The infinite theorem is proved in Section 3 of the accompanying self-contained research note.
"""

from fractions import Fraction as Q
from math import comb, factorial
from pathlib import Path
import json


INF = 10**9


def vp(x, p):
    x = Q(x)
    if not x:
        return INF
    a, b, v = x.numerator, x.denominator, 0
    while a % p == 0:
        a //= p
        v += 1
    while b % p == 0:
        b //= p
        v -= 1
    return v


def residue(x, modulus):
    x = Q(x)
    return x.numerator*pow(x.denominator, -1, modulus) % modulus


def poly_add(A, B):
    C = A.copy()
    for k, x in B.items():
        C[k] = C.get(k, 0)+x
    return {k: x for k, x in C.items() if x}


def poly_scale(A, c):
    return {k: c*x for k, x in A.items() if c*x}


def poly_mul(A, B):
    C = {}
    for i, x in A.items():
        for j, y in B.items():
            C[i+j] = C.get(i+j, 0)+x*y
    return {k: x for k, x in C.items() if x}


def poly_mod(A, modulus):
    return {k: residue(x, modulus) for k, x in A.items()
            if residue(x, modulus)}


def poly_mul_mod(A, B, modulus):
    C = {}
    for i, x in A.items():
        for j, y in B.items():
            C[i+j] = (C.get(i+j, 0)+x*y) % modulus
    return {k: x for k, x in C.items() if x}


def sum_of_powers_polynomials(degree):
    # S_d(X)=sum_{0<=k<X} k^d, using Stirling numbers and falling factorials.
    stirling = [[1]]
    for d in range(1, degree+1):
        old = stirling[-1]
        row = [0]*(d+1)
        for a in range(1, d+1):
            row[a] = (old[a-1] if a-1 < len(old) else 0)
            if a < len(old):
                row[a] += a*old[a]
        stirling.append(row)
    falling = [{0: Q(1)}]
    for k in range(degree+1):
        falling.append(poly_mul(falling[-1], {1: Q(1), 0: Q(-k)}))
    result = []
    for d in range(degree+1):
        S = {}
        for a in range(d+1):
            S = poly_add(S, poly_scale(falling[a+1], Q(stirling[d][a], a+1)))
        result.append(S)
    return result


def unit_factorial_polynomials(p, precision, check_through):
    """Construct T_r^full=1-2p H_r X+p²T_r modulo p^precision."""
    modulus = p**precision
    cutoff = precision+2
    sums = sum_of_powers_polynomials(2*cutoff)
    logA = {}
    for m in range(1, cutoff+1):
        coefficient = Q((-1)**(m+1)*p**(2*m), m)*sum(
            (Q(1, (a*(p-a))**m) for a in range(1, (p-1)//2+1)), Q(0))
        S = {}
        for b in range(m+1):
            S = poly_add(S, poly_scale(sums[m+b], comb(m, b)))
        logA = poly_add(logA, poly_scale(S, coefficient))
    logA = {k: x for k, x in logA.items() if vp(x, p) < precision}
    assert all(vp(x, p) >= 2 for x in logA.values())
    exponent = poly_scale(logA, -2)
    expA, power = {0: Q(1)}, {0: Q(1)}
    for k in range(1, precision+1):
        power = poly_mul(power, exponent)
        expA = poly_add(expA, poly_scale(power, Q(1, factorial(k))))
    expA = {k: x for k, x in expA.items() if vp(x, p) < precision}
    short = {0: Q(1)}
    polys = []
    H = Q(0)
    for r in range(p):
        if r:
            H += Q(1, r)
            factor = {d: Q((-1)**d*(d+1)*p**d, r**d)
                      for d in range(precision)}
            short = poly_mul(short, factor)
            short = {k: x for k, x in short.items() if vp(x, p) < precision}
        T = poly_mod(poly_mul(expA, short), modulus)
        assert T.get(0, 0) % (p*p) == 1
        assert (T.get(1, 0)-residue(-2*p*H, modulus)) % (p*p) == 0
        assert all(x % (p*p) == 0 for k, x in T.items() if k >= 2)
        polys.append(T)
    unit = [1]
    for a in range(1, p*check_through+p):
        unit.append(unit[-1]*(a if a % p else 1) % modulus)
    c = factorial(p-1)
    for ell in range(check_through+1):
        for r in range(p):
            exact = pow(c, 2*ell, modulus)*factorial(r)**2*pow(
                unit[p*ell+r], -2, modulus) % modulus
            evaluated = sum(x*pow(ell, k, modulus) for k, x in polys[r].items()) % modulus
            assert evaluated == exact, (p, precision, ell, r)
    return polys


def residue_weight_polynomials(n, p, precision, polys):
    modulus = p**precision
    current = {(0, ()): 1}
    for coordinate in range(n):
        following = {}
        for (residue_sum, exponents), x in current.items():
            for r in range(p):
                scale = pow(factorial(r), -2, modulus)
                for degree, y in polys[r].items():
                    value = x*scale*y % modulus
                    if not value:
                        continue
                    key = (residue_sum+r, exponents+(degree,))
                    following[key] = (following.get(key, 0)+value) % modulus
        current = {key: x for key, x in following.items() if x}
    result = {}
    for (residue_sum, exponents), x in current.items():
        result.setdefault(residue_sum, {})[exponents] = x
    for P in result.values():
        for exponents, x in P.items():
            if sum(exponents) >= 2:
                assert x % (p*p) == 0
            swapped = exponents[1:2]+exponents[:1]+exponents[2:]
            assert P.get(swapped, 0) == x
        linear = []
        for i in range(n):
            e = tuple(1 if k == i else 0 for k in range(n))
            linear.append(P.get(e, 0))
        assert len(set(x % (p*p) for x in linear)) == 1
    return result


def bessel_operator_basis(n, modulus):
    # O_k gives E_k=O_k(x,theta) Phi modulo N², in normal order x^a theta^b.
    operators = [{(0, 0): 1}, {(0, 1): 1}]
    for k in range(1, n):
        out = {}
        for (a, b), x in operators[k].items():
            out[(a, b+1)] = (out.get((a, b+1), 0)+x) % modulus
            out[(a, b)] = (out.get((a, b), 0)+a*x) % modulus
        for (a, b), x in operators[k-1].items():
            out[(a+1, b)] = (out.get((a+1, b), 0)-(n-k+1)*x) % modulus
        inverse = pow(k+1, -1, modulus)
        operators.append({key: x*inverse % modulus for key, x in out.items()
                          if x*inverse % modulus})
    return operators


def reduce_symmetric_polynomial(P, n, modulus, operators):
    maximum = max((max(e) for e in P), default=0)
    normal = [({0: 1}, {})]
    for exponent in range(maximum):
        A, B = normal[-1]
        nextA = {k: k*x % modulus for k, x in A.items() if k*x % modulus}
        nextA = poly_add(nextA, {k+1: x for k, x in B.items()})
        nextB = poly_add(A, {k: k*x % modulus for k, x in B.items() if k*x % modulus})
        normal.append((poly_mod(nextA, modulus), poly_mod(nextB, modulus)))
    out = {}
    for k in range(n+1):
        coefficient = {}
        for exponents, x in P.items():
            term = {0: x}
            for i, e in enumerate(exponents):
                term = poly_mul_mod(term, normal[e][1 if i < k else 0], modulus)
                if not term:
                    break
            for a, y in term.items():
                coefficient[a] = (coefficient.get(a, 0)+y) % modulus
        for a, x in coefficient.items():
            for (b, order), y in operators[k].items():
                key = (a+b, order)
                out[key] = (out.get(key, 0)+x*y) % modulus
    return {key: x for key, x in out.items() if x}


def octahedral_matrix(n, bound, modulus):
    """Use the even exponential convolution B_n=A_(n-1) B_1."""
    even, B1 = [], []
    for j in range(bound+1):
        row, value = [], 1
        for k in range(2*j+1):
            row.append(value % modulus)
            if k < 2*j:
                value = value*(2*j-k)//(k+1)
        even.append(row[::2])
        B1.append([row[j-N] for N in range(j+1)])
    one = [row[0] for row in B1]
    lower = one
    for dimension in range(2, n):
        lower = [sum(even[j][k]*one[k]*lower[j-k] for k in range(j+1)) % modulus
                 for j in range(bound+1)]
    matrix = []
    for j in range(bound+1):
        row = [0]*(j+1)
        for k in range(j+1):
            multiplier = even[j][k]*lower[j-k] % modulus
            if multiplier:
                for N, value in enumerate(B1[k]):
                    row[N] = (row[N]+multiplier*value) % modulus
        assert row[j] == 1
        matrix.append(row)
    return matrix


def construct_and_check_rows(n, p, precision, bound):
    modulus = p**precision
    polys = unit_factorial_polynomials(p, precision, bound//p+2)
    weights = residue_weight_polynomials(n, p, precision, polys)
    operators = bessel_operator_basis(n, modulus)
    reduced = {s: reduce_symmetric_polynomial(P, n, modulus, operators)
               for s, P in weights.items()}
    assert all(x % (p*p) == 0 for A in reduced.values()
               for (degree, order), x in A.items() if degree > 0)
    B = octahedral_matrix(n, bound, modulus)
    a = [row[0] for row in B]
    unit = [1]
    for j in range(1, 2*bound+1):
        unit.append(unit[-1]*(j if j % p else 1) % modulus)
    c = factorial(p-1)
    tested_columns = tested_boundaries = 0
    for j in range(bound+1):
        L, b = divmod(j, p)
        epsilon = (2*b)//p
        D = []
        for i in range(L+1):
            total = 0
            for d in range(min(n-1, L-i)+1):
                A = reduced.get(b+p*d, {})
                value = sum(A.get((L-d-i, k), 0)*pow(i, k, modulus)
                            for k in range(n+1)) % modulus
                total += p**(2*d)*pow(c, -2*(L-d), modulus)*value
            factorial_ratio = factorial(2*L+epsilon)//factorial(2*i)
            D.append(p**epsilon*unit[2*j]*factorial_ratio*total % modulus)
        assert (a[j]-sum(D[i]*a[i] for i in range(L+1))) % modulus == 0, (
            'central', n, p, j)
        for N in range(1, L+1):
            required = min(precision, 2+2*vp(N, p))
            delta = B[j][p*N]-sum(D[i]*B[i][N] for i in range(N, L+1))
            assert delta % (p**required) == 0, ('shifted', n, p, j, N, required)
            tested_columns += 1
        for i in range(L):
            required = min(precision, 2+2*(vp(factorial(L), p)-vp(factorial(i), p)))
            assert D[i]*a[i] % (p**required) == 0, ('boundary', n, p, j, i, required)
            tested_boundaries += 1
    return dict(dimension=n, p=p, working_precision=precision,
                through_row=bound, unit_factorial_interpolation_passed=True,
                symmetric_Bessel_reduction_constructed=True,
                central_identity_passed=True, shifted_columns_tested=tested_columns,
                boundary_terms_tested=tested_boundaries,
                all_row_operator_checks_passed=True)


def interpolation_vector(B, K, bound, modulus):
    c = [0]*(bound+1)
    for j in range(K, bound+1):
        c[j] = (B[j][0]-sum(c[N]*B[j][N] for N in range(K, j))) % modulus
    return c


def check_descent(n, p, s, m, bound):
    K = m*p**s
    v = vp(K, p)
    modulus = p**(2*v+3)
    B = octahedral_matrix(n, bound, modulus)
    c = interpolation_vector(B, K, bound, modulus)
    d = interpolation_vector(B, K//p, bound//p, modulus)
    minimum, first = INF, None
    for N in range(1, bound+1):
        assert vp(c[N], p) >= max(0, 2*v-2*vp(N, p))
        expected = d[N//p] if N % p == 0 else 0
        value = vp((c[N]-expected) % modulus, p)
        assert value >= 2*v, (n, p, s, m, N, value)
        if value < minimum:
            minimum, first = value, N
    return dict(dimension=n, p=p, s=s, m=m, K=K, through_column=bound,
                required_valuation=2*v, observed_minimum_valuation=minimum,
                first_minimum_column=first, weighted_bound_passed=True,
                descent_passed=True)


def mul(a, b, length):
    out = [Q(0)]*length
    for i, x in enumerate(a[:length]):
        if x:
            for j, y in enumerate(b[:length-i]):
                if y:
                    out[i+j] += x*y
    return out


def divide(a, b, length):
    assert b[0] == 1
    out = [Q(0)]*length
    for k in range(length):
        out[k] = (a[k] if k < len(a) else 0)-sum(
            b[j]*out[k-j] for j in range(1, min(k+1, len(b))))
    return out


def exponential(g, length):
    out = [Q(1)]+[Q(0)]*(length-1)
    for k in range(1, length):
        out[k] = sum(j*g[j]*out[k-j]
                     for j in range(1, min(k+1, len(g))))/k
    return out


def compose(a, b, length):
    out = [Q(0)]*length
    for coefficient in reversed(a[:length]):
        out = mul(out, b, length)
        out[0] += coefficient
    return out


def canonical_lift(g, p, length):
    J = (length-1)//p
    inverse = [Q(0)]+[
        exponential([-j*x for x in g[:j]], j)[j-1]/j
        for j in range(1, J+1)]
    qp = [Q(0)]*p+exponential([p*x for x in g], length-p)
    return compose(inverse, qp, length)


def exact_period(n, length):
    b = [1]*length
    for dimension in range(n-1):
        b = [sum(comb(j, k)**2*b[j-k] for k in range(j+1)) for j in range(length)]
    return [comb(2*j, j)*b[j] for j in range(length)]


def check_direct_defect(n, p, s, m, length):
    K = m*p**s
    assert length > K
    H = [Q(0)]
    for j in range(1, 2*length+1):
        H.append(H[-1]+Q(1, j))
    F = exact_period(n, length)
    lower = exact_period(n-1, length)
    G = [2*sum(comb(2*j, 2*k)*comb(2*k, k)*lower[j-k]*(H[2*j]-H[k])
               for k in range(j+1)) for j in range(length)]
    g = divide(G, F, length)
    sig = canonical_lift(g, p, length)
    assert all(vp(x, p) >= 0 for x in sig)
    left = mul(F, compose(F[:K//p], sig, length), length)
    right = mul(F[:K], compose(F, sig, length), length)
    values = [vp(x-y, p) for x, y in zip(left, right)]
    minimum = min(values)
    assert minimum >= 2*vp(K, p)
    return dict(dimension=n, p=p, s=s, m=m, through_u_degree=length-1,
                actual_BV_logarithmic_companion=True, lift_p_integral=True,
                required_valuation=2*vp(K, p), observed_minimum_valuation=minimum,
                first_minimum_degree=values.index(minimum), passed=True)


def main():
    results = dict(scope='Bounded exact checks of the explicit row construction, interpolation, and direct defects; the infinite theorem is proved separately.',
                   row_operator_checks=[], interpolation_checks=[], direct_defect_checks=[])
    for case in [(3, 5, 6, 145), (4, 5, 6, 145)]:
        result = construct_and_check_rows(*case)
        results['row_operator_checks'].append(result)
        print('PASS constructed rows', json.dumps(result), flush=True)
    for case in [(3, 5, 3, 1, 300), (4, 5, 3, 1, 300),
                 (5, 7, 2, 1, 160), (6, 7, 2, 1, 160),
                 (3, 5, 2, 2, 150)]:
        result = check_descent(*case)
        results['interpolation_checks'].append(result)
        print('PASS descent', json.dumps(result), flush=True)
    for case in [(3, 5, 3, 1, 137), (4, 5, 3, 1, 137), (5, 7, 2, 1, 75)]:
        result = check_direct_defect(*case)
        results['direct_defect_checks'].append(result)
        print('PASS direct defect', json.dumps(result), flush=True)
    results['all_passed'] = True
    destination = Path(__file__).with_name('octahedral_interpolation_validation.json')
    destination.write_text(json.dumps(results, indent=2)+'\n')
    print('ALL CHECKS PASSED', flush=True)


if __name__ == '__main__':
    main()
\end{lstlisting}

\subsection{verify\_positive\_interpolation.py}\label{app:program-5}
\noindent Original source: \texttt{verify\_positive\_interpolation.py}; 218 lines.

\begin{lstlisting}[language=Python]
"""Bounded exact checks for the positive-vertex proof.

Standard library only.  The theorem, including every series coefficient,
is proved in Section 2 of the accompanying self-contained research note.  This script tests three
independent components: the factorial block lemma, the integer triangular
descent, and direct ordinary truncation defects computed from the canonical
coordinate.  Finite arithmetic checks are not used in place of the proof.
"""

from fractions import Fraction as Q
from functools import lru_cache
from math import comb, factorial
from pathlib import Path
import json


def vp(x, p):
    x = Q(x)
    if not x:
        return 10**9
    a, b, v = x.numerator, x.denominator, 0
    while a % p == 0:
        a //= p
        v += 1
    while b % p == 0:
        b //= p
        v -= 1
    return v


def coefficient_matrix(family, dimension):
    n = dimension

    @lru_cache(None)
    def entry(j, N):
        if N > j or N < 0:
            return 0
        if family == 'cube':
            return comb(2*j, j-N)**n
        r = n+1
        return factorial(r*j)//(factorial(j+(r-1)*N)*factorial(j-N)**(r-1))

    return entry


def check_blocks(family, n, p, bound):
    B = coefficient_matrix(family, n)
    count, slack = 0, 10**9
    for j in range(bound+1):
        l = j//p
        d = Q(B(j, 0), B(l, 0))
        assert vp(d, p) >= 0
        for N in range(1, l+1):
            required = 2+2*vp(N, p)
            value = vp(B(j, p*N)-d*B(l, N), p)
            assert value >= required, (family, n, p, j, N, value, required)
            count += 1
            if value < 10**9:
                slack = min(slack, value-required)
    return dict(family=family, dimension=n, p=p, through_row=bound,
                tested_entries=count, minimum_extra_valuation=slack,
                passed=True)


def interpolation_vector(B, K, bound, modulus):
    c = [0]*(bound+1)
    for j in range(K, bound+1):
        c[j] = (B(j, 0)-sum(c[N]*(B(j, N) % modulus)
                            for N in range(K, j))) % modulus
    return c


def check_descent(family, n, p, s, m, bound):
    K = m*p**s
    v = vp(K, p)
    modulus = p**(2*v+3)
    B = coefficient_matrix(family, n)
    c = interpolation_vector(B, K, bound, modulus)
    d = interpolation_vector(B, K//p, bound//p, modulus)
    minimum, first = 10**9, None
    for N in range(1, bound+1):
        assert vp(c[N], p) >= max(0, 2*v-2*vp(N, p))
        expected = d[N//p] if N % p == 0 else 0
        value = vp((c[N]-expected) % modulus, p)
        assert value >= 2*v, (family, n, p, s, m, N, value)
        if value < minimum:
            minimum, first = value, N
    return dict(family=family, dimension=n, p=p, s=s, m=m, K=K,
                through_column=bound, working_modulus=modulus,
                required_valuation=2*v, observed_minimum_valuation=minimum,
                first_minimum_column=first, weighted_bound_passed=True,
                descent_passed=True)


def mul(a, b, length):
    out = [Q(0)]*length
    for i, x in enumerate(a[:length]):
        if x:
            for j, y in enumerate(b[:length-i]):
                if y:
                    out[i+j] += x*y
    return out


def divide(a, b, length):
    assert b[0] == 1
    out = [Q(0)]*length
    for k in range(length):
        out[k] = (a[k] if k < len(a) else 0)-sum(
            b[j]*out[k-j] for j in range(1, min(k+1, len(b))))
    return out


def exponential(g, length):
    assert not g or g[0] == 0
    out = [Q(1)]+[Q(0)]*(length-1)
    for k in range(1, length):
        out[k] = sum(j*g[j]*out[k-j]
                     for j in range(1, min(k+1, len(g))))/k
    return out


def compose(a, b, length):
    out = [Q(0)]*length
    for coefficient in reversed(a[:length]):
        out = mul(out, b, length)
        out[0] += coefficient
    return out


def parameters(family, n, length):
    H = [Q(0)]
    for j in range(1, (n+1)*length+1):
        H.append(H[-1]+Q(1, j))
    B = coefficient_matrix(family, n)
    F = [B(j, 0) for j in range(length)]
    if family == 'cube':
        G = [2*n*F[j]*(H[2*j]-H[j]) for j in range(length)]
    else:
        r = n+1
        G = [r*F[j]*(H[r*j]-H[j]) for j in range(length)]
    return F, divide(G, F, length)


def canonical_lift(g, p, length):
    # Lagrange inversion of q=u exp(g(u)), followed by q -> q^p.
    J = (length-1)//p
    inverse = [Q(0)]+[
        exponential([-j*x for x in g[:j]], j)[j-1]/j
        for j in range(1, J+1)]
    qp = [Q(0)]*p+exponential([p*x for x in g], length-p)
    return compose(inverse, qp, length)


def check_full_defect(family, n, p, s, m, length):
    K = m*p**s
    assert length > K
    F, g = parameters(family, n, length)
    sig = canonical_lift(g, p, length)
    assert all(vp(x, p) >= 0 for x in sig)
    left = mul(F, compose(F[:K//p], sig, length), length)
    right = mul(F[:K], compose(F, sig, length), length)
    difference = [x-y for x, y in zip(left, right)]
    values = [vp(x, p) for x in difference]
    minimum = min(values)
    required = 2*vp(K, p)
    assert minimum >= required, (family, n, p, s, m, minimum)
    return dict(family=family, dimension=n, p=p, s=s, m=m,
                through_u_degree=length-1, required_valuation=required,
                observed_minimum_valuation=minimum,
                first_minimum_degree=values.index(minimum),
                lift_p_integral=True, passed=True)


def main():
    results = dict(
        scope='Bounded exact arithmetic supporting the all-level proof; not an infinite proof.',
        block_checks=[], interpolation_checks=[], direct_defect_checks=[])
    block_cases = [
        ('cube', 1, 3, 120), ('cube', 4, 5, 150),
        ('cube', 6, 7, 150), ('cube', 8, 3, 100),
        ('simplex', 4, 7, 150), ('simplex', 5, 7, 120),
        ('simplex', 6, 3, 100), ('simplex', 9, 5, 100)]
    for case in block_cases:
        result = check_blocks(*case)
        results['block_checks'].append(result)
        print('PASS block', json.dumps(result), flush=True)

    descent_cases = [
        ('cube', 4, 5, 1, 1, 130), ('cube', 4, 5, 2, 1, 180),
        ('cube', 4, 5, 3, 1, 400), ('cube', 4, 5, 2, 2, 160),
        ('cube', 5, 7, 2, 1, 180), ('cube', 6, 7, 2, 1, 180),
        ('simplex', 4, 7, 1, 1, 120), ('simplex', 4, 7, 2, 1, 160),
        ('simplex', 4, 7, 2, 2, 180), ('simplex', 5, 7, 2, 1, 160),
        ('simplex', 6, 11, 2, 1, 200)]
    for case in descent_cases:
        result = check_descent(*case)
        results['interpolation_checks'].append(result)
        print('PASS descent', json.dumps(result), flush=True)

    defect_cases = [
        ('cube', 4, 5, 2, 1, 81), ('cube', 4, 5, 3, 1, 139),
        ('cube', 4, 5, 2, 2, 75), ('cube', 5, 7, 2, 1, 79),
        ('simplex', 4, 7, 2, 1, 82), ('simplex', 5, 7, 2, 1, 77),
        ('simplex', 4, 7, 2, 2, 111)]
    for case in defect_cases:
        result = check_full_defect(*case)
        results['direct_defect_checks'].append(result)
        print('PASS direct defect', json.dumps(result), flush=True)

    results['all_passed'] = True
    destination = Path(__file__).with_name('positive_interpolation_validation.json')
    destination.write_text(json.dumps(results, indent=2)+'\n')
    print('ALL CHECKS PASSED', flush=True)


if __name__ == '__main__':
    main()
\end{lstlisting}

\subsection{verify\_canonical\_basis.py}\label{app:program-6}
\noindent Original source: \texttt{verify\_canonical\_basis.py}; 224 lines.

\begin{lstlisting}[language=Python]
"""Finite exact checks for the all-level canonical negative-power basis theorem.

Standard-library only. The all-prime, all-dimension statement is proved in
canonical_negative_basis.tex; the checks below are not an infinite proof.
"""
from fractions import Fraction as Q
from math import comb, factorial
from pathlib import Path
import json


def mul(a, b, N):
    c = [Q(0)] * N
    for i, x in enumerate(a[:N]):
        if x:
            for j, y in enumerate(b[:N-i]):
                if y:
                    c[i+j] += x*y
    return c


def div(a, b, N):
    assert b[0] == 1
    c = [Q(0)] * N
    for k in range(N):
        c[k] = (a[k] if k < len(a) else 0) - sum(
            b[j]*c[k-j] for j in range(1, min(k+1, len(b))))
    return c


def exp(g, N):
    c = [Q(1)] + [Q(0)]*(N-1)
    for k in range(1, N):
        c[k] = sum(j*g[j]*c[k-j] for j in range(1, min(k+1, len(g))))/k
    return c


def compose(a, b, N):
    c = [Q(0)]*N
    for x in reversed(a[:N]):
        c = mul(c, b, N)
        c[0] += x
    return c


def power(a, k, N):
    c = [Q(1)] + [Q(0)]*(N-1)
    for _ in range(k):
        c = mul(c, a, N)
    return c


def vp(x, p):
    x = Q(x)
    if not x:
        return 1000000
    n, d, v = x.numerator, x.denominator, 0
    while n % p == 0:
        v += 1
        n //= p
    while d % p == 0:
        v -= 1
        d //= p
    return v


def octa_coefficients(n, N):
    b = [1]*N
    for _ in range(n-1):
        b = [sum(comb(k,j)**2*b[k-j] for j in range(k+1)) for k in range(N)]
    return [comb(2*k,k)*b[k] for k in range(N)]


def parameters(kind, n, N):
    harmonic = [Q(0)]
    for k in range(1, (n+1)*N+1):
        harmonic.append(harmonic[-1]+Q(1,k))
    if kind == 'cube':
        f = [comb(2*k,k)**n for k in range(N)]
        g = [2*n*f[k]*(harmonic[2*k]-harmonic[k]) for k in range(N)]
    elif kind == 'simplex':
        r = n+1
        f = [factorial(r*k)//factorial(k)**r for k in range(N)]
        g = [r*f[k]*(harmonic[r*k]-harmonic[k]) for k in range(N)]
    else:
        b = [1]*N
        for _ in range(n-2):
            b = [sum(comb(k,j)**2*b[k-j] for j in range(k+1)) for k in range(N)]
        bn = [sum(comb(k,j)**2*b[k-j] for j in range(k+1)) for k in range(N)]
        hsum = [sum(comb(k,j)**2*b[k-j]*harmonic[j] for j in range(k+1)) for k in range(N)]
        f = [comb(2*k,k)*bn[k] for k in range(N)]
        g = [2*comb(2*k,k)*(harmonic[2*k]*bn[k]-hsum[k]) for k in range(N)]
    return f, div(g, f, N)


def canonical_lift(g, p, N):
    # Revert q=u exp(g) by Lagrange inversion, then substitute q(u)^p.
    J = (N-1)//p
    reverse = [Q(0)] + [exp([-j*x for x in g[:j]],j)[j-1]/j
                           for j in range(1,J+1)]
    qp = [Q(0)]*p+exp([p*x for x in g],N-p)
    return compose(reverse,qp,N)


def hasse(f, h, N, length):
    result = [Q(0)]*length
    for k in range(min(length,(N-1)//h+1)):
        result[k] = (-1)**(h*k)*comb(N-1,h*k)*f[k]
    return result


def vertex_coefficient(kind,n,p,j):
    if kind == 'cube':
        return comb(p+2*j,j)**n
    if kind == 'simplex':
        r = n+1
        return factorial(p+r*j)//(factorial(p+j)*factorial(j)**(r-1))
    lower = octa_coefficients(n-1,j+1)
    return sum(comb(p+2*j,2*l)*lower[l]*comb(p+2*(j-l),j-l)
               for l in range(j+1))


def run_case(kind,n,p):
    N = 3*p+9
    h = n+1 if kind == 'simplex' else 2
    f,g = parameters(kind,n,N)
    sig = canonical_lift(g,p,N)
    hp = hasse(f,h,p,N)
    h2p = hasse(f,h,2*p,N)
    poly = [2*x-y for x,y in zip(hp,h2p)]
    assert all(vp(poly[k]-(f[k] if k<p else 0),p)>=2 for k in range(N))
    b = [Q(0)]*N
    for j in range((p-1)//h+1):
        b[j] = (-1)**(p+h*j)*comb(2*p-1,p+h*j)*vertex_coefficient(kind,n,p,j)
    ratio = div([-x for x in b],h2p,N)
    approximate_lift = [Q(0)]*p+power(ratio,h,N-p)
    assert all(vp(x-y,p)>=2 for x,y in zip(sig,approximate_lift))
    fp_sig = compose(f,sig,N)
    quotient = div(f,fp_sig,N)
    difference = [x-y for x,y in zip(quotient,poly)]
    assert all(vp(x,p)>=2 for x in difference)
    return dict(family=kind,dimension=n,p=p,through_u_degree=N-1,
                minimum_valuation=min(vp(x,p) for x in difference),
                lift_formula_passed=True,quotient_formula_passed=True,
                polynomial_identification_passed=True)



def vertex_polynomial(kind,n,N):
    h=n+1 if kind=='simplex' else 2
    M=h*N
    if kind=='cube':
        return [(-1)**(n*j)*comb(M-j-1,j)**n for j in range(N)]
    if kind=='simplex':
        r=n+1
        result=[]
        for j in range(N):
            numerator=factorial(M-j-1)
            denominator=factorial(M-r*j-1)*factorial(j)**(r-1)
            assert numerator % denominator == 0
            result.append((-1)**((r-1)*j)*(numerator//denominator))
        return result
    lower=octa_coefficients(n-1,N)
    return [sum((-1)**(j-l)*comb(M-j+l-1,2*l)*comb(M-j-l-1,j-l)*lower[l]
                for l in range(j+1)) for j in range(N)]


def twice_q_derivative_cleared(a,F,g,K,N):
    # If R=u^-K*a/F, return u^K*(q d/dq)^2 R.
    b=div(a,F,N)
    scale=[1]+[j*g[j] for j in range(1,N)]
    first=div([(j-K)*b[j] for j in range(N)],scale,N)
    return div([(j-K)*first[j] for j in range(N)],scale,N)


def canonical_case(kind,n,p,K):
    # Test n=pK, so the required modulus is p^(2*(1+v_p(K))).
    L=p*K
    N=L+12
    expected=2*(1+vp(K,p))
    F,g=parameters(kind,n,N+p)
    sig=canonical_lift(g,p,N+p)
    Fs=compose(F,sig,N)
    ratio=power(div([1],sig[p:],N),K,N)
    A_L=vertex_polynomial(kind,n,L)
    A_K=vertex_polynomial(kind,n,K)
    assert A_L[0]==A_K[0]==1
    R_error=[x-y for x,y in zip(mul(A_L,Fs,N),
             mul(ratio,mul(compose(A_K,sig,N),F,N),N))]
    R_val=min(vp(x,p) for x in R_error)
    assert R_val>=expected
    D2=twice_q_derivative_cleared(A_L,F,g,L,N)
    derivative_val=min(vp(x,p) for x in D2)
    assert derivative_val>=expected
    # Canonical numerator comes from [F exp(-N*g)]_<N.
    C_L=mul(F,exp([-L*x for x in g],L),L)
    C_K=mul(F,exp([-K*x for x in g],K),K)
    assert all(vp(x,p)>=0 for x in C_L+C_K)
    error=[x-y for x,y in zip(mul(C_L,Fs,N),
           mul(ratio,mul(compose(C_K,sig,N),F,N),N))]
    minimum=min(vp(x,p) for x in error)
    assert minimum>=expected
    result=dict(family=kind,dimension=n,p=p,base_index=K,canonical_index=L,
                required_valuation=expected,cleared_coefficients_checked=N,
                tested_laurent_exponents=[-L,N-1-L],
                canonical_frobenius_minimum=minimum,
                vertex_frobenius_minimum=R_val,
                second_derivative_minimum=derivative_val,
                all_passed=True)
    return result


if __name__=='__main__':
    cases=[('cube',4,5,1),('cube',4,5,5),('cube',4,5,10),('cube',5,7,7),
           ('simplex',4,7,1),('simplex',4,7,7),('simplex',5,7,7),
           ('octa',3,5,1),('octa',3,5,5),('octa',4,5,5)]
    results=[]
    for case in cases:
        result=canonical_case(*case)
        results.append(result)
        print('PASS',json.dumps(result),flush=True)
    output=dict(scope='Finite exact checks for an auxiliary negative-power canonical basis, not the positive truncation conjecture.',
                all_passed=True,cases=results)
    Path(__file__).with_name('canonical_basis_validation.json').write_text(json.dumps(output,indent=2)+'\n')
\end{lstlisting}

\subsection{verify\_first\_level.py}\label{app:program-7}
\noindent Original source: \texttt{verify\_first\_level.py}; 159 lines.

\begin{lstlisting}[language=Python]
"""Finite exact checks for the all-dimensional first-level theorem.

Standard-library only. The all-prime, all-dimension statement is proved in
first_level_all_dimensions.tex; the checks below are not an infinite proof.
"""
from fractions import Fraction as Q
from math import comb, factorial
from pathlib import Path
import json


def mul(a, b, N):
    c = [Q(0)] * N
    for i, x in enumerate(a[:N]):
        if x:
            for j, y in enumerate(b[:N-i]):
                if y:
                    c[i+j] += x*y
    return c


def div(a, b, N):
    assert b[0] == 1
    c = [Q(0)] * N
    for k in range(N):
        c[k] = (a[k] if k < len(a) else 0) - sum(
            b[j]*c[k-j] for j in range(1, min(k+1, len(b))))
    return c


def exp(g, N):
    c = [Q(1)] + [Q(0)]*(N-1)
    for k in range(1, N):
        c[k] = sum(j*g[j]*c[k-j] for j in range(1, min(k+1, len(g))))/k
    return c


def compose(a, b, N):
    c = [Q(0)]*N
    for x in reversed(a[:N]):
        c = mul(c, b, N)
        c[0] += x
    return c


def power(a, k, N):
    c = [Q(1)] + [Q(0)]*(N-1)
    for _ in range(k):
        c = mul(c, a, N)
    return c


def vp(x, p):
    x = Q(x)
    if not x:
        return 1000000
    n, d, v = x.numerator, x.denominator, 0
    while n % p == 0:
        v += 1
        n //= p
    while d % p == 0:
        v -= 1
        d //= p
    return v


def octa_coefficients(n, N):
    b = [1]*N
    for _ in range(n-1):
        b = [sum(comb(k,j)**2*b[k-j] for j in range(k+1)) for k in range(N)]
    return [comb(2*k,k)*b[k] for k in range(N)]


def parameters(kind, n, N):
    harmonic = [Q(0)]
    for k in range(1, (n+1)*N+1):
        harmonic.append(harmonic[-1]+Q(1,k))
    if kind == 'cube':
        f = [comb(2*k,k)**n for k in range(N)]
        g = [2*n*f[k]*(harmonic[2*k]-harmonic[k]) for k in range(N)]
    elif kind == 'simplex':
        r = n+1
        f = [factorial(r*k)//factorial(k)**r for k in range(N)]
        g = [r*f[k]*(harmonic[r*k]-harmonic[k]) for k in range(N)]
    else:
        b = [1]*N
        for _ in range(n-2):
            b = [sum(comb(k,j)**2*b[k-j] for j in range(k+1)) for k in range(N)]
        bn = [sum(comb(k,j)**2*b[k-j] for j in range(k+1)) for k in range(N)]
        hsum = [sum(comb(k,j)**2*b[k-j]*harmonic[j] for j in range(k+1)) for k in range(N)]
        f = [comb(2*k,k)*bn[k] for k in range(N)]
        g = [2*comb(2*k,k)*(harmonic[2*k]*bn[k]-hsum[k]) for k in range(N)]
    return f, div(g, f, N)


def canonical_lift(g, p, N):
    # Revert q=u exp(g) by Lagrange inversion, then substitute q(u)^p.
    J = (N-1)//p
    reverse = [Q(0)] + [exp([-j*x for x in g[:j]],j)[j-1]/j
                           for j in range(1,J+1)]
    qp = [Q(0)]*p+exp([p*x for x in g],N-p)
    return compose(reverse,qp,N)


def hasse(f, h, N, length):
    result = [Q(0)]*length
    for k in range(min(length,(N-1)//h+1)):
        result[k] = (-1)**(h*k)*comb(N-1,h*k)*f[k]
    return result


def vertex_coefficient(kind,n,p,j):
    if kind == 'cube':
        return comb(p+2*j,j)**n
    if kind == 'simplex':
        r = n+1
        return factorial(p+r*j)//(factorial(p+j)*factorial(j)**(r-1))
    lower = octa_coefficients(n-1,j+1)
    return sum(comb(p+2*j,2*l)*lower[l]*comb(p+2*(j-l),j-l)
               for l in range(j+1))


def run_case(kind,n,p):
    N = 3*p+9
    h = n+1 if kind == 'simplex' else 2
    f,g = parameters(kind,n,N)
    sig = canonical_lift(g,p,N)
    hp = hasse(f,h,p,N)
    h2p = hasse(f,h,2*p,N)
    poly = [2*x-y for x,y in zip(hp,h2p)]
    assert all(vp(poly[k]-(f[k] if k<p else 0),p)>=2 for k in range(N))
    b = [Q(0)]*N
    for j in range((p-1)//h+1):
        b[j] = (-1)**(p+h*j)*comb(2*p-1,p+h*j)*vertex_coefficient(kind,n,p,j)
    ratio = div([-x for x in b],h2p,N)
    approximate_lift = [Q(0)]*p+power(ratio,h,N-p)
    assert all(vp(x-y,p)>=2 for x,y in zip(sig,approximate_lift))
    fp_sig = compose(f,sig,N)
    quotient = div(f,fp_sig,N)
    difference = [x-y for x,y in zip(quotient,poly)]
    assert all(vp(x,p)>=2 for x in difference)
    return dict(family=kind,dimension=n,p=p,through_u_degree=N-1,
                minimum_valuation=min(vp(x,p) for x in difference),
                lift_formula_passed=True,quotient_formula_passed=True,
                polynomial_identification_passed=True)


if __name__ == '__main__':
    cases=[('cube',4,5),('cube',5,7),('cube',6,7),
           ('simplex',4,7),('simplex',5,7),('simplex',6,11),
           ('octa',3,5),('octa',4,5),('octa',5,7),('octa',6,7)]
    output=[]
    for case in cases:
        result=run_case(*case)
        output.append(result)
        print('PASS',json.dumps(result),flush=True)
    target=Path(__file__).with_name('first_level_validation.json')
    target.write_text(json.dumps(dict(scope='Finite exact checks supporting the all-dimensional s=1 theorem.',
                                     all_passed=True,cases=output),indent=2)+'\n')
\end{lstlisting}

\subsection{verify\_full\_zagier.py}\label{app:program-8}
\noindent Original source: \texttt{verify\_full\_zagier.py}; 386 lines.

\begin{lstlisting}[language=Python]
#!/usr/bin/env python3
"""Bounded exact checks of the complete Zagier D proof.

The infinite theorem is proved in full_zagier_proof.tex; these calculations
test the modular normalization, canonical-basis Hecke identities, Cartier
recurrence, and complete finite portions of the truncation congruences.
Python standard library only. Run from any current directory.
"""
from __future__ import annotations

from math import isqrt
from pathlib import Path
from time import perf_counter
import json


def zmul(a, b, n):
    c = [0]*n
    for i, x in enumerate(a[:n]):
        if x:
            for j, y in enumerate(b[:n-i]):
                if y:
                    c[i+j] += x*y
    return c


def zinv(a, n):
    assert a[0] == 1
    b = [1]+[0]*(n-1)
    for k in range(1, n):
        b[k] = -sum(a[j]*b[k-j] for j in range(1, min(k+1, len(a))))
    return b


def zcompose(a, t, n):
    out = [0]*n
    for c in reversed(a):
        out = zmul(out, t, n)
        out[0] += c
    return out


def trim(a):
    a = list(a)
    while len(a) > 1 and a[-1] == 0:
        a.pop()
    return a


def mmul(a, b, n, mod):
    if n <= 0:
        return []
    a = trim([x % mod for x in a[:n]])
    b = trim([x % mod for x in b[:n]])
    if not a or not b or a == [0] or b == [0]:
        return [0]*n
    count = min(len(a), len(b))
    width = max(1, (count.bit_length()+2*mod.bit_length()+7)//8)
    aa = int.from_bytes(b''.join(x.to_bytes(width, 'little') for x in a), 'little')
    bb = int.from_bytes(b''.join(x.to_bytes(width, 'little') for x in b), 'little')
    size = len(a)+len(b)-1
    raw = (aa*bb).to_bytes(size*width, 'little')
    out = [int.from_bytes(raw[i*width:(i+1)*width], 'little') % mod
           for i in range(min(n, size))]
    return out+[0]*(n-len(out))


def minv(a, n, mod):
    assert a[0] % mod == 1
    b = [1]
    while len(b) < n:
        size = min(2*len(b), n)
        ab = mmul(a[:size], b, size, mod)
        correction = [(-x) % mod for x in ab]
        correction[0] = (correction[0]+2) % mod
        b = mmul(b, correction, size, mod)
    return b


def mcompose(a, t, n, mod):
    assert t[0] == 0
    a = a[:n]
    if not a:
        return [0]*n
    block = isqrt(len(a))+1
    powers = [[1]+[0]*(n-1)]
    for _ in range(block):
        powers.append(mmul(powers[-1], t, n, mod))
    out = [0]*n
    for start in reversed(range(0, len(a), block)):
        out = mmul(out, powers[block], n, mod)
        for r, x in enumerate(a[start:start+block]):
            x %= mod
            if x:
                out = [(v+x*u) % mod for v, u in zip(out, powers[r])]
    return out


def mpow(a, k, n, mod):
    out = [1]+[0]*(n-1)
    while k:
        if k & 1:
            out = mmul(out, a, n, mod)
        k //= 2
        if k:
            a = mmul(a, a, n, mod)
    return out


def vsub(a, p, n):
    out = [0]*n
    for k, x in enumerate(a):
        if p*k >= n:
            break
        out[p*k] = x
    return out


def vp(x, p, cap=10**6):
    if x == 0:
        return cap
    v = 0
    while v < cap and x % p == 0:
        x //= p
        v += 1
    return v


def make_data(n):
    b = [1]+[0]*(n-1)
    for k in range(1, n):
        e = 0 if k % 5 == 0 else 5 if k % 5 in (1, 4) else -5
        for _ in range(abs(e)):
            if e > 0:
                for j in range(n-1, k-1, -1):
                    b[j] -= b[j-k]
            else:
                for j in range(k, n):
                    b[j] += b[j-k]
    theta = [1]+[0]*(n-1)
    ct = [0, 3, 1, -1, -3]
    for k in range(1, n):
        for j in range(k, n, k):
            theta[j] += ct[k % 5]
    return [0]+b[:-1], theta


def apery(n):
    out = [1, 3]
    for k in range(2, n):
        value = (11*k*k-11*k+3)*out[k-1]+(k-1)**2*out[k-2]
        assert value % (k*k) == 0
        out.append(value//(k*k))
    return out[:n]


T = []
TH = []
CANONICAL = []


def canonical_polynomials(n):
    """g_j = sum_k CANONICAL[j][k] t^(-k)/Theta."""
    hinv = zinv(T[1:], n+1)
    thinv = zinv(TH, n+1)
    powers = [1]+[0]*n
    polys = [[0]]
    for j in range(1, n+1):
        powers = zmul(powers, hinv, n+1)
        unit = zmul(powers, thinv, n+1)
        poly = [0]*j+[1]
        for k in range(1, j):
            c = unit[j-k]
            for r, x in enumerate(polys[k]):
                poly[r] -= c*x
        polys.append(poly)
    return polys


def g_series(n, positive_length):
    """Dictionary giving g_n from q^(-n) through q^(positive_length-1)."""
    if n == 0:
        return {}
    length = n+positive_length
    hinv = zinv(T[1:], length)
    thinv = zinv(TH, length)
    powers = [1]+[0]*(length-1)
    out = [0]*length
    for j in range(1, n+1):
        powers = zmul(powers, hinv, length)
        c = CANONICAL[n][j]
        if not c:
            continue
        unit = zmul(powers, thinv, length)
        for r in range(length-(n-j)):
            out[n-j+r] += c*unit[r]
    return {j-n: x for j, x in enumerate(out)}


def dg_series(n, length):
    if n == 0:
        return [0]*length
    poly = [(-1)**j*CANONICAL[n][j] for j in range(1, n+1)]
    return zmul(zcompose(poly, T, length), zinv(TH, length), length)


def check_hecke(p, n, length):
    c = -1 if p % 5 == 2 else 1
    split = p % 5 in (1, 4)
    eps = 1 if p % 5 == 1 else -1
    gn = g_series(n, p*(length-1)+1)
    gpn = g_series(p*n, length)
    gsmall = g_series(n//p, length) if n % p == 0 else {}
    dg = dg_series(n, p*(length-1)+1)
    dgpn = dg_series(p*n, length)
    dgsmall = dg_series(n//p, length) if n % p == 0 else [0]*length
    for j in range(-p*n, length):
        vgn = gn.get(j//p, 0) if j % p == 0 else 0
        udg = dg[p*j] if j >= 0 else 0
        dsmall = dgsmall[j] if j >= 0 else 0
        vdg = dg[j//p] if j >= 0 and j % p == 0 else 0
        ugn = gn.get(p*j, 0)
        if split:
            assert vgn+eps*p*p*ugn == gpn.get(j, 0)+eps*p*p*gsmall.get(j, 0)
            assert vdg+eps*p*p*udg == (dgpn[j] if j >= 0 else 0)+eps*p*p*dsmall
        else:
            assert vgn+c*p*p*udg == gpn.get(j, 0)+c*p*p*dsmall
            assert vdg-c*p*p*ugn == (dgpn[j] if j >= 0 else 0)-c*p*p*gsmall.get(j, 0)
    return dict(p=p, basis_index=n, q_degrees_checked=[-p*n, length-1],
                exact_integer_identities=2, passed=True)


def delta(p, n, length):
    gn = g_series(n, p*(length-1)+1)
    gsmall = g_series(n//p, length) if n % p == 0 else {}
    dg = dg_series(n, p*(length-1)+1)
    dgsmall = dg_series(n//p, length) if n % p == 0 else [0]*length
    # The negative parts cancel identically; inspect all retained nonnegative terms.
    return ([gn.get(p*j, 0)-gsmall.get(j, 0) for j in range(length)],
            [dg[p*j]-dgsmall[j] for j in range(length)])


def check_delta(p, n, length):
    c = -1 if p % 5 == 2 else 1
    lhs = delta(p, p*n, length)
    rhs = delta(p, n, p*(length-1)+1)
    if p % 5 in (1, 4):
        eps = 1 if p % 5 == 1 else -1
        assert lhs[0] == [eps*p*p*rhs[0][p*j] for j in range(length)]
        assert lhs[1] == [eps*p*p*rhs[1][p*j] for j in range(length)]
    else:
        assert lhs[0] == [c*p*p*rhs[1][p*j] for j in range(length)]
        assert lhs[1] == [-c*p*p*rhs[0][p*j] for j in range(length)]
    target = 2*vp(p*n, p)
    assert all(x % p**target == 0 for row in lhs for x in row)
    return dict(p=p, basis_index=p*n, through_q_degree=length-1,
                proved_precision=target, recurrence_passed=True,
                divisibility_passed=True)


def check_full(p, s, m):
    K = m*p**s
    length = 2*K+9
    cap = 2*s+2
    mod = p**cap
    th = [x % mod for x in TH[:length]]
    t = [x % mod for x in T[:length]]
    a = apery(K)
    tp = vsub(t, p, length)
    thp = vsub(th, p, length)
    lhs = mmul(th, mcompose(a[:K//p], tp, length, mod), length, mod)
    rhs = mmul(thp, mcompose(a, t, length, mod), length, mod)
    defects = [(x-y) % mod for x, y in zip(lhs, rhs)]
    minimum = min(vp(x, p, cap) for x in defects)
    assert minimum >= 2*s
    return dict(p=p, s=s, m=m, through_q_degree=length-1,
                arithmetic_modulus=mod, required_valuation=2*s,
                observed_minimum_valuation=minimum,
                minimum_is_capped=(minimum == cap), passed=True)


def dual_series(K, positive_length, mod):
    length = K+positive_length
    t = [x % mod for x in T[:length+1]]
    th = [x % mod for x in TH[:length]]
    a = apery(K)
    rev = [(-1)**(K-1-j)*a[K-1-j] % mod for j in range(K)]
    unit = mmul(mcompose(rev, t, length, mod),
                mpow(minv(t[1:], length, mod), K, length, mod),
                length, mod)
    return mmul(unit, minv(th, length, mod), length, mod)


def check_dual(p, r, m, positive_length):
    K = m*p**r
    cap = 2*r+2
    mod = p**cap
    lhs = dual_series(K, p*(positive_length-1)+1, mod)
    rhs = dual_series(K//p, positive_length, mod)
    defects = [(lhs[p*j+K]-rhs[j+K//p]) % mod
               for j in range(-K//p, positive_length)]
    minimum = min(vp(x, p, cap) for x in defects)
    assert minimum >= 2*r
    return dict(p=p, r=r, m=m,
                q_degrees_checked=[-K//p, positive_length-1],
                arithmetic_modulus=mod, required_valuation=2*r,
                observed_minimum_valuation=minimum,
                minimum_is_capped=(minimum == cap), passed=True)


def main():
    global T, TH, CANONICAL
    start = perf_counter()
    full_cases = [(p, s, m)
                  for p, levels in [(2, (2, 3, 4)), (3, (2, 3)),
                                    (7, (2, 3)), (13, (2,))]
                  for s in levels for m in (1, 2)]
    full_cases += [(11, 2, 1), (11, 2, 2), (19, 2, 1)]
    maximum = max(2*m*p**s+10 for p, s, m in full_cases)
    T, TH = make_data(maximum)
    # Basic normalization checked with exact integers.
    assert zcompose(apery(50), T, 50) == TH[:50]
    aa = [-11*x-y for x, y in zip(T[:50], zmul(T, T, 50))]
    aa[0] += 1
    assert [j*T[j] for j in range(50)] == zmul(T, zmul(aa, zmul(TH, TH, 50), 50), 50)
    a = apery(50)
    assert all(k*k*a[k] == (11*k*k-11*k+3)*a[k-1]+(k-1)**2*a[k-2]
               for k in range(2, 50))
    mod = 7**7
    test_a = [(-k*k+11*k-3) % mod for k in range(29)]
    test_b = [(2*k*k-7*k+5) % mod for k in range(29)]
    assert mmul(test_a, test_b, 29, mod) == [x % mod for x in zmul(test_a, test_b, 29)]
    test_a[0] = 1
    assert minv(test_a, 29, mod) == [x % mod for x in zinv(test_a, 29)]
    CANONICAL = canonical_polynomials(49)
    assert CANONICAL[1] == [0, 1]
    assert CANONICAL[2] == [0, -7, 1]
    assert CANONICAL[3] == [0, 19, -12, 1]
    for n in range(1, 50):
        gs = g_series(n, 1)
        assert gs[-n] == 1 and all(gs[j] == 0 for j in range(-n+1, 0))
    out = dict(scope="Finite exact checks; the all-prime/all-level statement is proved in the note.",
               implementation_checks=True, modular_parametrization=True,
               canonical_wronskian_identity=True,
               canonical_principal_parts_checked=49,
               hecke_identities=[], paired_cartier=[], dual_truncations=[],
               full_congruences=[])
    for p in (2, 3, 7):
        for n in (1, 2, p):
            r = check_hecke(p, n, 12)
            out["hecke_identities"].append(r)
            print("PASS exact Hecke", r, flush=True)
        for n in (1, p):
            r = check_delta(p, n, 8)
            out["paired_cartier"].append(r)
            print("PASS exact paired recurrence", r, flush=True)
    for p in (11, 19):
        for n in (1, 2):
            r = check_hecke(p, n, 12)
            out["hecke_identities"].append(r)
            print("PASS exact split Hecke", r, flush=True)
    r = check_delta(11, 1, 8)
    out["paired_cartier"].append(r)
    print("PASS exact split paired recurrence", r, flush=True)
    for values in [(2, 2, 1, 25), (2, 3, 2, 25), (3, 1, 1, 25),
                   (3, 2, 2, 25), (7, 1, 1, 25), (7, 2, 2, 25),
                   (13, 1, 1, 25)]:
        r = check_dual(*values)
        out["dual_truncations"].append(r)
        print("PASS dual truncation", r, flush=True)
    for values in full_cases:
        r = check_full(*values)
        out["full_congruences"].append(r)
        print("PASS full finite congruence", r, flush=True)
    out["counts"] = {k: len(out[k]) for k in
                     ("hecke_identities", "paired_cartier",
                      "dual_truncations", "full_congruences")}
    out["all_passed"] = True
    out["seconds"] = round(perf_counter()-start, 3)
    destination = Path(__file__).resolve().parent
    (destination/"full_zagier_validation.json").write_text(json.dumps(out, indent=2)+"\n")
    print("ALL CHECKS PASSED", out["counts"], "seconds", out["seconds"], flush=True)


if __name__ == "__main__":
    main()
\end{lstlisting}

\subsection{verify\_octahedral\_progress.py}\label{app:program-9}
\noindent Original source: \texttt{verify\_octahedral\_progress.py}; 247 lines.

{\lstset{basicstyle=\ttfamily\fontsize{8}{9.2}\selectfont}
\begin{lstlisting}[language=Python]
#!/usr/bin/env python3
"""Finite exact checks for the octahedral weighted Cartier reduction.

The written proof, not these bounded checks, establishes the weighted theorem.
The full positive-truncation congruence for all levels remains unresolved.
Python 3.10+ standard library only; no workspace imports or network required.
"""
from pathlib import Path
from fractions import Fraction
from math import comb, isqrt
import json
import time

def trim(a):
    a = list(a)
    while len(a) > 1 and a[-1] == 0:
        a.pop()
    return a

def mmul(a, b, n, mod):
    """Exact modular convolution by carry-free Kronecker substitution."""
    if n <= 0:
        return []
    a = trim([x % mod for x in a[:n]])
    b = trim([x % mod for x in b[:n]])
    if not a or not b or a == [0] or b == [0]:
        return [0]*n
    count = min(len(a), len(b))
    width = max(1, (count.bit_length() + 2*mod.bit_length() + 7)//8)
    aa = int.from_bytes(b''.join(x.to_bytes(width, 'little') for x in a), 'little')
    bb = int.from_bytes(b''.join(x.to_bytes(width, 'little') for x in b), 'little')
    # Each convolution coefficient is < count*mod^2 < 256^width.
    size = len(a)+len(b)-1
    raw = (aa*bb).to_bytes(size*width, 'little')
    out = [int.from_bytes(raw[i*width:(i+1)*width], 'little') % mod
           for i in range(min(n, size))]
    return out+[0]*(n-len(out))

def minv(a, n, mod):
    assert a[0] % mod == 1
    b = [1]
    while len(b) < n:
        size = min(2*len(b), n)
        ab = mmul(a[:size], b, size, mod)
        correction = [(-x) % mod for x in ab]
        correction[0] = (correction[0]+2) % mod
        b = mmul(b, correction, size, mod)
    return b

def mpow(a, k, n, mod):
    out = [1]+[0]*(n-1)
    while k:
        if k & 1:
            out = mmul(out, a, n, mod)
        k //= 2
        if k:
            a = mmul(a, a, n, mod)
    return out

def vsub(a, p, n):
    out = [0]*n
    for k, x in enumerate(a):
        if p*k >= n:
            break
        out[p*k] = x
    return out

def mcompose(a, t, n, mod):
    """Baby-step/giant-step polynomial composition; t(0)=0."""
    assert t[0] == 0
    a = a[:n]
    if not a:
        return [0]*n
    block = isqrt(len(a))+1
    powers = [[1]+[0]*(n-1)]
    for _ in range(block):
        powers.append(mmul(powers[-1], t, n, mod))
    out = [0]*n
    for start in reversed(range(0, len(a), block)):
        out = mmul(out, powers[block], n, mod)
        for r, x in enumerate(a[start:start+block]):
            x %= mod
            if x:
                out = [(v+x*u) % mod for v, u in zip(out, powers[r])]
    return out

def valuation(x, p, cap):
    if x == 0:
        return cap
    v = 0
    while v < cap and x % p == 0:
        v += 1
        x //= p
    return v

def divisor_sums(n):
    out = [0]*n
    for d in range(1, n):
        for j in range(d, n, d):
            out[j] += d
    return out

def eta_h6(n,mod):
    out=[1]+[0]*(n-1)
    for k in range(1,n):
        e=-5+(k%2==0)-(k%3==0)+5*(k%6==0)
        if e>0:
            for _ in range(e):
                for j in range(n-1,k-1,-1):out[j]=(out[j]-out[j-k])%mod
        else:
            for _ in range(-e):
                for j in range(k,n):out[j]=(out[j]+out[j-k])%mod
    return [0]+out[:-1]

def params(n,mod):
    h=eta_h6(n,mod);h2=mmul(h,h,n,mod)
    h3=mmul(h2,h,n,mod)
    num=[(x+17*y+72*z)%mod for x,y,z in zip(h,h2,h3)]
    Q=[(18*x+72*y)%mod for x,y in zip(h,h2)];Q[0]=1
    qinv=minv(Q,n,mod)
    u=mmul(num,mmul(qinv,qinv,n,mod),n,mod)
    sig=divisor_sums(n)
    th=[1]+[(6*sig[k]-(12*sig[k//2]if k%2==0 else 0)
             +(18*sig[k//3]if k%3==0 else 0)
             -(36*sig[k//6]if k%6==0 else 0))%mod for k in range(1,n)]
    return u,th

def coeffs(n,mod):
    # Exact P-recursion, then reduce; divisions by n^3 need not be units.
    a=[1,6]
    for k in range(2,n):
        num=2*(2*k-1)*(10*k*k-10*k+3)*a[-1]-36*(k-1)*(2*k-1)*(2*k-3)*a[-2]
        assert num%(k**3)==0
        a.append(num//(k**3))
    return [x%mod for x in a[:n]]

def check(p,s,m,n):
    K=m*p**s;L=K//p;cap=2*s+2;mod=p**cap
    start=time.perf_counter();u,th=params(n,mod)
    a=coeffs(K,mod)
    big=mcompose(a,u,n,mod)
    tiny=(n-1)//p+1
    small=mcompose(a[:L],u[:tiny],tiny,mod)
    E=[(x-y)%mod for x,y in zip(mmul(th,vsub(small,p,n),n,mod),mmul(vsub(th,p,n),big,n,mod))]
    vals=[valuation(z,p,cap) for z in E]
    result=dict(p=p,s=s,m=m,through_q_degree=n-1,min_valuation=min(vals),first_min_index=vals.index(min(vals)),seconds=round(time.perf_counter()-start,2))
    print(json.dumps(result),flush=True)
    assert min(vals)>=2*s
    return result


def exact_bv_coordinate(n):
    # The BV multinomial formula, independently of eta products and PF recurrences.
    H=[Fraction(0)]
    for k in range(1,2*n+1):H.append(H[-1]+Fraction(1,k))
    F=[];G=[]
    for k in range(n):
        b=sum(comb(k,j)**2*comb(2*j,j) for j in range(k+1))
        weighted=sum(comb(k,j)**2*comb(2*j,j)*H[k-j] for j in range(k+1))
        F.append(comb(2*k,k)*b)
        G.append(2*comb(2*k,k)*(H[2*k]*b-weighted))
    g=[Fraction(0)]*n
    for k in range(n):g[k]=G[k]-sum(F[j]*g[k-j] for j in range(1,k+1))
    ex=[Fraction(1)]+[Fraction(0)]*(n-1)
    for k in range(1,n):ex[k]=sum(j*g[j]*ex[k-j] for j in range(1,k+1))/k
    q=[Fraction(0)]+ex[:-1]
    assert all(c.denominator==1 for c in q)
    return [int(c) for c in q]

def vp(n,p):
    if n==0:return 10**9
    out=0
    while n%p==0:n//=p;out+=1
    return out

def check_weighted(p,s,n):
    K=p**s;cap=2*s+3;mod=p**cap
    u,th=params(n,mod)
    fk=mcompose(coeffs(K,mod),u,n,mod)
    H=mmul(fk,minv(th,n,mod),n,mod)
    assert all((j*j*H[j])%p**(2*s)==0 for j in range(1,n))
    assert H[0]==1 and all(H[j]==0 for j in range(1,min(K,n)))
    return dict(p=p,s=s,K=K,through_q_degree=n-1,precision=cap,
                all_squared_derivative_coefficients_divisible=True)

def primitive_checks(p,s,limit):
    K=p**s
    for n in range(limit):
        c=comb(2*n,n)
        assert vp(c,p)>=vp(n+K+1,p)
        assert s+vp(c,p)>=vp(n+K,p)
    return dict(p=p,s=s,K=K,central_binomial_indices_checked=limit,
                shifted_integrality_passed=True)

# Symbolic bivariate-polynomial verification of the exact telescoping certificate.
def padd(a,b):
    c=dict(a)
    for k,v in b.items():c[k]=c.get(k,0)+v
    return {k:v for k,v in c.items() if v}
def pscale(a,c):return {k:v*c for k,v in a.items() if v*c}
def pmul(a,b):
    c={}
    for (i,j),x in a.items():
        for (k,l),y in b.items():c[i+k,j+l]=c.get((i+k,j+l),0)+x*y
    return {k:v for k,v in c.items() if v}
def ppow(a,n):
    c={(0,0):1}
    for _ in range(n):c=pmul(c,a)
    return c
def certificate_check():
    one={(0,0):1};j={(1,0):1};z={(0,1):1}
    shift=padd(padd(z,one),pscale(j,-1))
    P=padd(pscale(j,3),pscale(padd(z,one),-4))
    Pnext=padd(P,pscale(one,3))
    lhs=padd(pscale(pmul(pmul(padd(pscale(j,2),one),Pnext),ppow(shift,2)),2),
             pscale(pmul(ppow(j,3),P),-1))
    A=padd(padd(pscale(ppow(z,2),10),pscale(z,10)),pscale(one,3))
    rhs=padd(padd(ppow(padd(z,one),4),pscale(pmul(A,ppow(shift,2)),-1)),
             pscale(pmul(ppow(padd(z,pscale(j,-1)),2),ppow(shift,2)),9))
    assert lhs==rhs
    return True

def main():
    start=time.perf_counter()
    assert certificate_check()
    n=40;mod=5**10
    u,th=params(n,mod)
    assert mcompose(coeffs(n,mod),u,n,mod)==th
    qb=exact_bv_coordinate(n)
    assert mcompose(qb,u,n,mod)==[0,1]+[0]*(n-2)
    result={'scope':'Finite checks supporting the written weighted theorem; no full all-level positive-truncation proof is claimed.',
            'symbolic_telescoping_certificate':True,
            'bv_coordinate_check':{'through_degree':n-1,'modulus':mod,'passed':True},
            'weighted':[],'primitives':[],'full_congruence_bounded_probes':[]}
    for p,s in [(5,1),(5,2),(5,3),(7,1),(7,2),(11,2)]:
        n=4*p**s+29
        result['weighted'].append(check_weighted(p,s,n))
        result['primitives'].append(primitive_checks(p,s,3*p**s+20))
        result['full_congruence_bounded_probes'].append(check(p,s,1,n))
    result['seconds']=round(time.perf_counter()-start,3)
    result['all_passed']=True
    target=Path(__file__).resolve().parent
    (target/'octahedral_progress_validation.json').write_text(json.dumps(result,indent=2)+'\n')
    print('PASS: exact symbolic certificate, BV coordinate check, 6 weighted series checks, 6 shifted-integrality ranges, 6 bounded full-congruence probes.')
    print('The resonant U_p congruence for every coefficient and every level is still not proved.')

if __name__=='__main__':main()
\end{lstlisting}

}

\clearpage

\clearpage

\section{Complete exact certificate program}\label{SP:source-appendix-1}\label{SP:app:code}
The multiplication routine encodes polynomial coefficients as base
$2^{8w}$ digits, with $w$ chosen so that each unreduced convolution
coefficient is strictly smaller than that base. Thus its packed integer
product is exact coefficientwise. All divisions used to construct the
lift equation are checked before reduction; all cleared denominators
have constant term one.
\begin{lstlisting}[language=Python]
#!/usr/bin/env python3
"""Finite whole-series certificates for specified small-prime octahedral cases.

This establishes only the levels actually executed. It does not establish an
all-level theorem. All algebra is exact over finite integer polynomial rings.
"""
from math import comb, factorial, isqrt
from pathlib import Path
import json, time, hashlib, argparse

def trim(a):
    a=list(a)
    while len(a)>1 and a[-1]==0:a.pop()
    return a
def add(a,b,mod):
    c=[0]*max(len(a),len(b))
    for i,x in enumerate(a):c[i]+=x
    for i,x in enumerate(b):c[i]+=x
    return trim([x%mod for x in c])
def scale(a,k,mod):return trim([(k*x)%mod for x in a])
def mul(a,b,mod):
    a=trim([x%mod for x in a]);b=trim([x%mod for x in b])
    if a==[0] or b==[0]:return [0]
    width=(min(len(a),len(b)).bit_length()+2*mod.bit_length()+7)//8
    aa=int.from_bytes(b''.join(x.to_bytes(width,'little') for x in a),'little')
    bb=int.from_bytes(b''.join(x.to_bytes(width,'little') for x in b),'little')
    length=len(a)+len(b)-1
    raw=(aa*bb).to_bytes(length*width,'little')
    return trim([int.from_bytes(raw[i*width:(i+1)*width],'little')%mod for i in range(length)])
def power(a,k,mod):
    out=[1]
    while k:
        if k&1:out=mul(out,a,mod)
        k//=2
        if k:a=mul(a,a,mod)
    return out
def taylor(a,r,p,mod):
    out=[0]*(p*(len(a)-1)+1)
    for j,x in enumerate(a):
        if j>=r:out[p*j]=(x*comb(j,r))%mod
    return trim(out)
def vp(n,p):
    if n==0:return 10**9
    v=0
    while n%p==0:n//=p;v+=1
    return v
def composition_squares(dimension,length):
    values=[1]*length
    for _ in range(dimension-1):
        values=[sum(comb(j,k)**2*values[j-k] for k in range(j+1))
                for j in range(length)]
    return values

def vertex(K,b3,mod):
    fac=[factorial(j) for j in range(2*K)]
    out=[]
    for j in range(K):
        total=0
        for ell in range(j+1):
            num=fac[2*K-ell-1]
            den=fac[2*K-2*j-1]*fac[ell]*fac[j-ell]**2
            assert num%den==0
            total+=(-1 if ell%2 else 1)*(num//den)*b3[j-ell]
        out.append(total%mod)
    assert out[0]==1
    return trim(out)
def octa_period(dimension,length):
    values=composition_squares(dimension,length)
    return [comb(2*j,j)*values[j] for j in range(length)]

def profiles(weights,limit,at=0,counts=()):
    if at==len(weights):
        if sum(counts):yield counts
        return
    for k in range(limit//weights[at]+1):
        yield from profiles(weights,limit-k*weights[at],at+1,counts+(k,))

def build(dimension,p,s):
    started=time.time();prec=2*s;yprec=prec-1
    assert s>=2 and p>=3 and dimension>=2
    assert all(p%q for q in range(2,isqrt(p)+1))
    n=p**(prec-2);mod=p**prec;ymod=p**yprec;workmod=p*n*ymod
    degree=max(k for k in range(1,3*yprec+1) if k-1-vp(k,p)<yprec)
    print(json.dumps({'stage':'setup','s':s,'N':n,'modulus':mod,'lift_polynomial_degree':degree}),flush=True)
    b3=composition_squares(dimension-1,2*p*n)
    av={k:vertex(k,b3,workmod) for k in (n,2*n,p*n,2*p*n)}
    a,b,ap,bp=[av[k] for k in (n,2*n,p*n,2*p*n)]
    ar=[taylor(a,r,p,workmod) for r in range(degree+1)]
    br=[taylor(b,r,p,workmod) for r in range(degree+1)]
    a0pow=[power(ar[0],r,workmod) for r in range(degree+2)]
    T=[]
    for r in range(degree+1):
        tr=mul(br[r],a0pow[r],workmod)
        for i in range(1,r+1):
            tr=add(tr,scale(mul(mul(ar[i],T[r-i],workmod),a0pow[i-1],workmod),-1,workmod),workmod)
        T.append(tr)
    coeff=[]
    for r in range(degree+1):
        num=add(scale(mul(bp,a0pow[degree+1],workmod),comb(n,r),workmod),scale(mul(mul(ap,T[r],workmod),a0pow[degree-r],workmod),-1,workmod),workmod)
        num=scale(num,p**r,workmod)
        assert all(x%(p*n)==0 for x in num),('divisibility',r)
        coeff.append(trim([(x//(p*n))%ymod for x in num]))
    A,B=coeff[:2]
    assert A[0]==0 and B[0]==1
    weights=[min(vp(v,p) for v in c) for c in coeff[2:]]
    assert all(w>=1 for w in weights)
    # Multivariate Lagrange: y=z+sum_m (R)!/((R-m+1)!prod m_r!)
    #   *(-1)^m prod c_r^m_r * z^(R-m+1), z=-A/B, c_r=Q_r/B.
    # Each tree coefficient is an integer, so profiles of valuation>=yprec vanish.
    ps=list(profiles(weights,yprec-1))
    maxden=max([1]+[1+sum((i+2)*c for i,c in enumerate(pr)) for pr in ps])
    bpows=[power(B,j,ymod) for j in range(maxden+1)]
    apows=[power(A,j,ymod) for j in range(maxden+1)]
    Ynum=scale(mul(A,bpows[maxden-1],ymod),-1,ymod)
    terms_info=[]
    for pr in ps:
        m=sum(pr);Rsum=sum((i+2)*c for i,c in enumerate(pr));zexp=Rsum-m+1
        den=factorial(zexp)
        for c in pr:den*=factorial(c)
        assert factorial(Rsum)%den==0
        tree=factorial(Rsum)//den
        if vp(tree,p)+sum(w*c for w,c in zip(weights,pr))>=yprec:continue
        term=apows[zexp]
        for q,c in zip(coeff[2:],pr):
            if c:term=mul(term,power(q,c,ymod),ymod)
        term=mul(term,bpows[maxden-Rsum-1],ymod)
        term=scale(term,(-1 if (Rsum+1)%2 else 1)*tree,ymod)
        Ynum=add(Ynum,term,ymod)
        terms_info.append({'profile':pr,'tree_coefficient':tree})
    Yden=bpows[maxden]
    root=[0]
    for r,c in enumerate(coeff):
        root=add(root,mul(mul(c,power(Ynum,r,ymod),ymod),power(Yden,degree-r,ymod),ymod),ymod)
    assert root==[0],('root',[(i,v) for i,v in enumerate(root) if v][:10])
    print(json.dumps({'stage':'lift_certified','s':s,'degrees':[len(Ynum)-1,len(Yden)-1],'elapsed_seconds':round(time.time()-started,3)}),flush=True)
    Q=scale(Ynum,p,mod);R=Yden
    assert Q[0]==0 and R[0]==1
    f=octa_period(dimension,p**s);fs=f[:p**(s-1)]
    L=[]
    for r in range(prec):
        lhs=taylor(fs,r,p,mod)
        if r:lhs=add(lhs,scale(taylor(fs,r-1,p,mod),n,mod),mod)
        L.append(add(mul(ap,lhs,mod),scale(mul(f,ar[r],mod),-1,mod),mod))
    target=[0];degree_bound=0
    for r,c in enumerate(L):
        term=mul(mul(c,power(Q,r,mod),mod),power(R,prec-1-r,mod),mod)
        degree_bound=max(degree_bound,len(term)-1)
        target=add(target,term,mod)
    assert target==[0],('target',[(i,v) for i,v in enumerate(target) if v][:10])
    payload={'p':p,'s':s,'m':1,'dimension':dimension,'modulus':mod,'intermediate_modulus':workmod,'vertex_indices':list(av),'vertex_polynomials':av,'lift_equation_coefficients':coeff,'lagrange_profiles':terms_info,'lift_y_numerator':Ynum,'lift_y_denominator':Yden,'target_L':L}
    packed=json.dumps(payload,sort_keys=True,separators=(',',':'))
    report={'theorem':f'Every coefficient of E_({p},{s},1) for octahedral dimension {dimension} is divisible by {p}^{2*s}.','dimension':dimension,'p':p,'s':s,'m':1,'modulus':mod,'intermediate_modulus':workmod,'all_polynomial_coefficients_checked':True,'power_series_cutoff':None,'lift_constructed_from':'BV vertex polynomials; no target truncation assumption','lift_equation_root_identity_zero':root==[0],'cleared_target_identity_zero':target==[0],'lift_equation_degree':degree,'coefficient_polynomial_degrees':[len(z)-1 for z in coeff],'lift_y_fraction_degrees':[len(Ynum)-1,len(Yden)-1],'lift_y_denominator_power':maxden,'cleared_identity_degree_bound':degree_bound,'certificate_sha256':hashlib.sha256(packed.encode()).hexdigest(),'elapsed_seconds':round(time.time()-started,3)}
    base=Path(__file__).resolve().parent;stem=f'certificate_octa{dimension}_p{p}_s{s}'
    (base/(stem+'.json')).write_text(json.dumps(payload,indent=2)+'\n')
    (base/(stem+'_report.json')).write_text(json.dumps(report,indent=2)+'\n')
    print(json.dumps(report,indent=2))
    return report

if __name__=='__main__':
    parser=argparse.ArgumentParser()
    parser.add_argument('--dimension',type=int,required=True)
    parser.add_argument('--p',type=int,required=True)
    parser.add_argument('--s',type=int,required=True)
    args=parser.parse_args();build(args.dimension,args.p,args.s)

\end{lstlisting}

\clearpage

\section{Complete polynomial arithmetic and individual-certificate source}\label{OP:source-appendix-1}
The following source is supplied as \texttt{verify\_dimension\_families.py}.
The dimension-uniform verifier in the next appendix imports this file.
All convolution uses exact integer arithmetic. For packed multiplication,
the byte width is chosen so that the largest possible convolution
coefficient is strictly below its base; hence there is no carry between
stored coefficients.
\begin{lstlisting}[language=Python]
#!/usr/bin/env python3
"""Finite whole-series certificates for specified small-prime octahedral cases.

This establishes only the levels actually executed. It does not establish an
all-level theorem. All algebra is exact over finite integer polynomial rings.
"""
from math import comb, factorial, isqrt
from pathlib import Path
import json, time, hashlib, argparse

def trim(a):
    a=list(a)
    while len(a)>1 and a[-1]==0:a.pop()
    return a
def add(a,b,mod):
    c=[0]*max(len(a),len(b))
    for i,x in enumerate(a):c[i]+=x
    for i,x in enumerate(b):c[i]+=x
    return trim([x%mod for x in c])
def scale(a,k,mod):return trim([(k*x)%mod for x in a])
def mul(a,b,mod):
    a=trim([x%mod for x in a]);b=trim([x%mod for x in b])
    if a==[0] or b==[0]:return [0]
    width=(min(len(a),len(b)).bit_length()+2*mod.bit_length()+7)//8
    aa=int.from_bytes(b''.join(x.to_bytes(width,'little') for x in a),'little')
    bb=int.from_bytes(b''.join(x.to_bytes(width,'little') for x in b),'little')
    length=len(a)+len(b)-1
    raw=(aa*bb).to_bytes(length*width,'little')
    return trim([int.from_bytes(raw[i*width:(i+1)*width],'little')%mod for i in range(length)])
def power(a,k,mod):
    out=[1]
    while k:
        if k&1:out=mul(out,a,mod)
        k//=2
        if k:a=mul(a,a,mod)
    return out
def taylor(a,r,p,mod):
    out=[0]*(p*(len(a)-1)+1)
    for j,x in enumerate(a):
        if j>=r:out[p*j]=(x*comb(j,r))%mod
    return trim(out)
def vp(n,p):
    if n==0:return 10**9
    v=0
    while n%p==0:n//=p;v+=1
    return v
def composition_squares(dimension,length):
    values=[1]*length
    for _ in range(dimension-1):
        values=[sum(comb(j,k)**2*values[j-k] for k in range(j+1))
                for j in range(length)]
    return values

def vertex(K,b3,mod):
    fac=[factorial(j) for j in range(2*K)]
    out=[]
    for j in range(K):
        total=0
        for ell in range(j+1):
            num=fac[2*K-ell-1]
            den=fac[2*K-2*j-1]*fac[ell]*fac[j-ell]**2
            assert num%den==0
            total+=(-1 if ell%2 else 1)*(num//den)*b3[j-ell]
        out.append(total%mod)
    assert out[0]==1
    return trim(out)
def octa_period(dimension,length):
    values=composition_squares(dimension,length)
    return [comb(2*j,j)*values[j] for j in range(length)]

def profiles(weights,limit,at=0,counts=()):
    if at==len(weights):
        if sum(counts):yield counts
        return
    for k in range(limit//weights[at]+1):
        yield from profiles(weights,limit-k*weights[at],at+1,counts+(k,))

def build(dimension,p,s,write_files=True,verbose=True,m=1):
    started=time.time();prec=2*s;yprec=prec-1
    assert s>=2 and p>=3 and dimension>=2
    assert all(p%q for q in range(2,isqrt(p)+1))
    n=p**(prec-2);mod=p**prec;ymod=p**yprec;workmod=p*n*ymod
    degree=max(k for k in range(1,3*yprec+1) if k-1-vp(k,p)<yprec)
    if verbose:print(json.dumps({'stage':'setup','s':s,'N':n,'modulus':mod,'lift_polynomial_degree':degree}),flush=True)
    b3=composition_squares(dimension-1,2*p*n)
    av={k:vertex(k,b3,workmod) for k in (n,2*n,p*n,2*p*n)}
    a,b,ap,bp=[av[k] for k in (n,2*n,p*n,2*p*n)]
    ar=[taylor(a,r,p,workmod) for r in range(degree+1)]
    br=[taylor(b,r,p,workmod) for r in range(degree+1)]
    a0pow=[power(ar[0],r,workmod) for r in range(degree+2)]
    T=[]
    for r in range(degree+1):
        tr=mul(br[r],a0pow[r],workmod)
        for i in range(1,r+1):
            tr=add(tr,scale(mul(mul(ar[i],T[r-i],workmod),a0pow[i-1],workmod),-1,workmod),workmod)
        T.append(tr)
    coeff=[]
    for r in range(degree+1):
        num=add(scale(mul(bp,a0pow[degree+1],workmod),comb(n,r),workmod),scale(mul(mul(ap,T[r],workmod),a0pow[degree-r],workmod),-1,workmod),workmod)
        num=scale(num,p**r,workmod)
        assert all(x%(p*n)==0 for x in num),('divisibility',r)
        coeff.append(trim([(x//(p*n))%ymod for x in num]))
    A,B=coeff[:2]
    assert A[0]==0 and B[0]==1
    weights=[min(vp(v,p) for v in c) for c in coeff[2:]]
    assert all(w>=1 for w in weights)
    # Multivariate Lagrange: y=z+sum_m (R)!/((R-m+1)!prod m_r!)
    #   *(-1)^m prod c_r^m_r * z^(R-m+1), z=-A/B, c_r=Q_r/B.
    # Each tree coefficient is an integer, so profiles of valuation>=yprec vanish.
    ps=list(profiles(weights,yprec-1))
    maxden=max([1]+[1+sum((i+2)*c for i,c in enumerate(pr)) for pr in ps])
    bpows=[power(B,j,ymod) for j in range(maxden+1)]
    apows=[power(A,j,ymod) for j in range(maxden+1)]
    Ynum=scale(mul(A,bpows[maxden-1],ymod),-1,ymod)
    terms_info=[]
    for pr in ps:
        tree_count=sum(pr);Rsum=sum((i+2)*c for i,c in enumerate(pr));zexp=Rsum-tree_count+1
        den=factorial(zexp)
        for c in pr:den*=factorial(c)
        assert factorial(Rsum)%den==0
        tree=factorial(Rsum)//den
        if vp(tree,p)+sum(w*c for w,c in zip(weights,pr))>=yprec:continue
        term=apows[zexp]
        for q,c in zip(coeff[2:],pr):
            if c:term=mul(term,power(q,c,ymod),ymod)
        term=mul(term,bpows[maxden-Rsum-1],ymod)
        term=scale(term,(-1 if (Rsum+1)%2 else 1)*tree,ymod)
        Ynum=add(Ynum,term,ymod)
        terms_info.append({'profile':pr,'tree_coefficient':tree})
    Yden=bpows[maxden]
    root=[0]
    for r,c in enumerate(coeff):
        root=add(root,mul(mul(c,power(Ynum,r,ymod),ymod),power(Yden,degree-r,ymod),ymod),ymod)
    assert root==[0],('root',[(i,v) for i,v in enumerate(root) if v][:10])
    if verbose:print(json.dumps({'stage':'lift_certified','s':s,'degrees':[len(Ynum)-1,len(Yden)-1],'elapsed_seconds':round(time.time()-started,3)}),flush=True)
    Q=scale(Ynum,p,mod);R=Yden
    assert Q[0]==0 and R[0]==1
    f=octa_period(dimension,m*p**s);fs=f[:m*p**(s-1)]
    L=[]
    for r in range(prec):
        lhs=taylor(fs,r,p,mod)
        if r:lhs=add(lhs,scale(taylor(fs,r-1,p,mod),n,mod),mod)
        L.append(add(mul(ap,lhs,mod),scale(mul(f,ar[r],mod),-1,mod),mod))
    target=[0];degree_bound=0
    for r,c in enumerate(L):
        term=mul(mul(c,power(Q,r,mod),mod),power(R,prec-1-r,mod),mod)
        degree_bound=max(degree_bound,len(term)-1)
        target=add(target,term,mod)
    assert target==[0],('target',[(i,v) for i,v in enumerate(target) if v][:10])
    payload={'p':p,'s':s,'m':m,'dimension':dimension,'modulus':mod,'intermediate_modulus':workmod,'vertex_indices':list(av),'vertex_polynomials':av,'lift_equation_coefficients':coeff,'lagrange_profiles':terms_info,'lift_y_numerator':Ynum,'lift_y_denominator':Yden,'target_L':L}
    packed=json.dumps(payload,sort_keys=True,separators=(',',':'))
    report={'theorem':f'Every coefficient of E_({p},{s},{m}) for octahedral dimension {dimension} is divisible by {p}^{2*s}.','dimension':dimension,'p':p,'s':s,'m':m,'modulus':mod,'intermediate_modulus':workmod,'all_polynomial_coefficients_checked':True,'power_series_cutoff':None,'lift_constructed_from':'BV vertex polynomials; no target truncation assumption','lift_equation_root_identity_zero':root==[0],'cleared_target_identity_zero':target==[0],'lift_equation_degree':degree,'coefficient_polynomial_degrees':[len(z)-1 for z in coeff],'lift_y_fraction_degrees':[len(Ynum)-1,len(Yden)-1],'lift_y_denominator_power':maxden,'cleared_identity_degree_bound':degree_bound,'certificate_sha256':hashlib.sha256(packed.encode()).hexdigest(),'elapsed_seconds':round(time.time()-started,3)}
    base=Path(__file__).resolve().parent;stem=f'certificate_octa{dimension}_p{p}_s{s}_m{m}'
    if write_files:(base/(stem+'.json')).write_text(json.dumps(payload,indent=2)+'\n')
    if write_files:(base/(stem+'_report.json')).write_text(json.dumps(report,indent=2)+'\n')
    if verbose:print(json.dumps(report,indent=2))
    return report

if __name__=='__main__':
    parser=argparse.ArgumentParser()
    parser.add_argument('--dimension',type=int,required=True)
    parser.add_argument('--p',type=int,required=True)
    parser.add_argument('--s',type=int,required=True)
    args=parser.parse_args();build(args.dimension,args.p,args.s)

\end{lstlisting}
\section{Complete dimension-uniform verifier}\label{OP:source-appendix-2}
Save this appendix as \texttt{certify\_all\_dimensions\_newton.py}
beside the previous source. The commands
\begin{verbatim}
python3 certify_all_dimensions_newton.py --s 2
python3 certify_all_dimensions_newton.py --s 3
\end{verbatim}
reproduce the two dimension-uniform proofs and write the exact JSON
certificates and summaries. A run at another level must pass every
assertion before it establishes any additional result.
\begin{lstlisting}[language=Python]
#!/usr/bin/env python3
"""Uniform dimension certificates through ordinary-polynomial Newton models.

This is an exact proof for the level actually executed, not an induction on s.
"""
from pathlib import Path
from math import comb,factorial
import time,json,hashlib,argparse
import verify_dimension_families as v

def main(s):
    started=time.time();p=3;prec=2*s;k=prec-1;a=prec-1
    N=p**(prec-2);M=p**prec;ymod=p**k;work=p**(4*s-2)
    degree=max(r for r in range(1,3*k+1) if r-1-v.vp(r,p)<k)
    def factorial_threshold(exponent):
        r=0;valuation=0
        while valuation<exponent:
            r+=1;valuation+=v.vp(r,p)
        return r
    ninput=factorial_threshold(a+k)
    nmodel=factorial_threshold(k)
    ntarget=factorial_threshold(prec)
    length=2*p*N
    print(json.dumps({'stage':'setup','level':s,'input_dimensions':ninput,'model_terms':nmodel,'target_dimensions':ntarget,'length':length}),flush=True)
    choose=[[comb(j,i)**2%work for i in range(j+1)] for j in range(length)]
    comps=[[1]+[0]*(length-1)]
    for n in range(1,2+ninput):
        last=comps[-1]
        comps.append([sum(choose[j][i]*last[j-i] for i in range(j+1))%work for j in range(length)])
    matrices={}
    for K in [N,2*N,p*N,2*p*N]:
        fac=[factorial(j) for j in range(2*K)]
        rows=[]
        for j in range(K):
            row=[]
            for ell in range(j+1):
                num=fac[2*K-ell-1];den=fac[2*K-2*j-1]*fac[ell]*fac[j-ell]**2
                assert num%den==0
                row.append(((-1 if ell%2 else 1)*(num//den))%work)
            rows.append(row)
        matrices[K]=rows
    print(json.dumps({'stage':'integer_matrices_ready','elapsed_seconds':round(time.time()-started,3)}),flush=True)
    def av_for(n):
        bj=comps[n-1]
        return {K:v.trim([sum(c*bj[j-ell] for ell,c in enumerate(row))%work for j,row in enumerate(rows)]) for K,rows in matrices.items()}
    def coeff_for(av):
        aa,bb,ap,bp=[av[K] for K in [N,2*N,p*N,2*p*N]]
        ar=[v.taylor(aa,r,p,work) for r in range(degree+1)]
        br=[v.taylor(bb,r,p,work) for r in range(degree+1)]
        a0pow=[v.power(ar[0],r,work) for r in range(degree+2)]
        T=[]
        for r in range(degree+1):
            tr=v.mul(br[r],a0pow[r],work)
            for i in range(1,r+1):tr=v.add(tr,v.scale(v.mul(v.mul(ar[i],T[r-i],work),a0pow[i-1],work),-1,work),work)
            T.append(tr)
        qs=[]
        for r in range(degree+1):
            num=v.add(v.scale(v.mul(bp,a0pow[degree+1],work),comb(N,r),work),v.scale(v.mul(v.mul(ap,T[r],work),a0pow[degree-r],work),-1,work),work)
            num=v.scale(num,p**r,work)
            assert all(c%(p*N)==0 for c in num)
            qs.append(v.trim([(c//(p*N))%ymod for c in num]))
        return qs
    avcache={};samples=[]
    for n in range(2,2+ninput):
        av=av_for(n);avcache[n]=av;samples.append(coeff_for(av))
    width=max(len(q) for s0 in samples for q in s0)
    diffs=[[[q[j] if j<len(q) else 0 for j in range(width)] for q in s0] for s0 in samples]
    newton=[];ordinary=[];bad=[]
    for r in range(ninput):
        dr=diffs[0];vr=v.vp(factorial(r),p)
        need=min(vr,k)
        for q in range(degree+1):
            for j,c in enumerate(dr[q]):
                if c%p**need:bad.append([r,q,j,c,need])
        newton.append([v.trim(q) for q in dr])
        if r<nmodel:
            unit=factorial(r)//p**vr;iv=pow(unit,-1,M)
            ordinary.append([[((c//p**vr)*iv)%M for c in q] for q in dr])
        diffs=[[[ (diffs[i+1][q][j]-diffs[i][q][j])%ymod for j in range(width)] for q in range(degree+1)] for i in range(len(diffs)-1)]
    if bad:
        print(json.dumps({'stage':'model_integrality_failed','bad':bad[:20]}),flush=True)
        raise AssertionError('ordinary model fails')
    print(json.dumps({'stage':'ordinary_model_certified','model_degree':nmodel-1,'elapsed_seconds':round(time.time()-started,3)}),flush=True)
    def model_at(n):
        out=[[0]*width for _ in range(degree+1)];fall=1;x=n-2
        for r,dr in enumerate(ordinary):
            if r:fall=(fall*(x-r+1))%M
            for q in range(degree+1):
                for j,c in enumerate(dr[q]):out[q][j]=(out[q][j]+fall*c)%M
        return [v.trim(q) for q in out]
    for n,actual in zip(range(2,2+ninput),samples):
        assert all(v.trim([c%ymod for c in q])==aa for q,aa in zip(model_at(n),actual))
    weights=[r-1-v.vp(r,p) for r in range(2,degree+1)]
    profiles=list(v.profiles(weights,k-1))
    maxden=max([1]+[1+sum((i+2)*c for i,c in enumerate(pr)) for pr in profiles])
    def root_from_model(qs):
        A,B=qs[:2];bpows=[v.power(B,r,M) for r in range(maxden+1)];apows=[v.power(A,r,M) for r in range(maxden+1)]
        Y=v.scale(v.mul(A,bpows[maxden-1],M),-1,M)
        for pr in profiles:
            count=sum(pr);Rsum=sum((i+2)*c for i,c in enumerate(pr));zexp=Rsum-count+1
            den=factorial(zexp)
            for c in pr:den*=factorial(c)
            assert factorial(Rsum)%den==0
            tree=factorial(Rsum)//den
            if v.vp(tree,p)+sum(w*c for w,c in zip(weights,pr))>=k:continue
            term=apows[zexp]
            for q,c in zip(qs[2:],pr):
                if c:term=v.mul(term,v.power(q,c,M),M)
            term=v.mul(term,bpows[maxden-Rsum-1],M)
            Y=v.add(Y,v.scale(term,(-1 if (Rsum+1)%2 else 1)*tree,M),M)
        return Y,bpows[maxden]
    reports=[]
    for n in range(2,2+ntarget):
        qs=model_at(n);Y,Z=root_from_model(qs)
        assert Z[0]%p!=0
        root=[0]
        for r,q in enumerate(qs):root=v.add(root,v.mul(v.mul(q,v.power(Y,r,ymod),ymod),v.power(Z,degree-r,ymod),ymod),ymod)
        assert root==[0],('root',n)
        av=avcache[n];ap=av[p*N];ar=[v.taylor(av[N],r,p,M) for r in range(prec)]
        f=[comb(2*j,j)*comps[n][j]%M for j in range(p**s)];fs=f[:p**(s-1)]
        L=[]
        for r in range(prec):
            lhs=v.taylor(fs,r,p,M)
            if r:lhs=v.add(lhs,v.scale(v.taylor(fs,r-1,p,M),N,M),M)
            L.append(v.add(v.mul(ap,lhs,M),v.scale(v.mul(f,ar[r],M),-1,M),M))
        Q=v.scale(Y,p,M);target=[0];bound=0
        for r,c in enumerate(L):
            term=v.mul(v.mul(c,v.power(Q,r,M),M),v.power(Z,prec-1-r,M),M)
            bound=max(bound,len(term)-1);target=v.add(target,term,M)
        assert target==[0],('target',n,[(i,c) for i,c in enumerate(target) if c][:10])
        report={'dimension':n,'root_identity_zero':True,'target_identity_zero':True,'target_degree_bound':bound,'Y_degree':len(Y)-1,'Z_degree':len(Z)-1}
        reports.append(report)
        print(json.dumps({'stage':'dimension_target_checked',**report,'elapsed_seconds':round(time.time()-started,3)}),flush=True)
    payload={'p':p,'s':s,'m':1,'work_modulus':work,'target_modulus':M,'input_dimensions':list(range(2,2+ninput)),'input_lift_coefficients':samples,'newton_coefficients':newton,'ordinary_model_coefficients_in_falling_basis':ordinary,'lagrange_profiles':profiles,'fixed_denominator_exponent':maxden,'target_reports':reports}
    packed=json.dumps(payload,sort_keys=True,separators=(',',':'))
    summary={'theorem':f'For all octahedral dimensions n>=2, every coefficient of E_(3,{s},1) is divisible by 3^{2*s}.','prime':p,'level':s,'multiplier':1,'target_modulus':M,'work_modulus':work,'input_dimension_count':ninput,'ordinary_model_degree':nmodel-1,'target_dimension_count':ntarget,'all_input_difference_integrality_checked':True,'all_polynomial_coefficients_checked':True,'power_series_cutoff':None,'certificate_sha256':hashlib.sha256(packed.encode()).hexdigest(),'elapsed_seconds':round(time.time()-started,3)}
    base=Path(__file__).resolve().parent
    (base/f'all_dimensions_newton_s{s}_certificate.json').write_text(json.dumps(payload,indent=2)+'\n')
    (base/f'all_dimensions_newton_s{s}_summary.json').write_text(json.dumps(summary,indent=2)+'\n')
    print(json.dumps(summary,indent=2))
if __name__=='__main__':
    parser=argparse.ArgumentParser();parser.add_argument('--s',type=int,required=True);args=parser.parse_args();main(args.s)

\end{lstlisting}
\clearpage
\section{Exact modular input matrices and model-only mode}\label{OP:source-appendix-3}
The following source is \texttt{certify\_all\_dimensions\_newton\_fast.py}.
It requires NumPy in addition to the previous arithmetic source.
For level $4$, use \texttt{--s 4 --model-only}. The overflow bound is
checked before any integer matrix multiplication. The output model
checkpoint is consumed by the next appendix.
\begin{lstlisting}[language=Python]
#!/usr/bin/env python3
"""Uniform dimension certificates through ordinary-polynomial Newton models.

This is an exact proof for the level actually executed, not an induction on s.
"""
from pathlib import Path
from math import comb,factorial
import time,json,hashlib,argparse
import numpy as np
import verify_dimension_families as v

def main(s,model_only=False):
    started=time.time();p=3;prec=2*s;k=prec-1;a=prec-1
    N=p**(prec-2);M=p**prec;ymod=p**k;work=p**(4*s-2)
    degree=max(r for r in range(1,3*k+1) if r-1-v.vp(r,p)<k)
    def factorial_threshold(exponent):
        r=0;valuation=0
        while valuation<exponent:
            r+=1;valuation+=v.vp(r,p)
        return r
    ninput=factorial_threshold(a+k)
    nmodel=factorial_threshold(k)
    ntarget=factorial_threshold(prec)
    length=2*p*N
    print(json.dumps({'stage':'setup','level':s,'input_dimensions':ninput,'model_terms':nmodel,'target_dimensions':ntarget,'length':length}),flush=True)
    assert work*work*length < 2**63
    exponent=4*s-2
    facunit=[1]*(2*length);facval=[0]*(2*length)
    for j in range(1,2*length):
        w=j;val=0
        while w%p==0:w//=p;val+=1
        facunit[j]=(facunit[j-1]*w)%work;facval[j]=facval[j-1]+val
    invunit=np.array([pow(w,-1,work) for w in facunit],dtype=np.int64)
    facunit=np.array(facunit,dtype=np.int64);facval=np.array(facval,dtype=np.int64)
    ppows=np.array([p**e if e<exponent else 0 for e in range(2*int(facval[-1])+1)],dtype=np.int64)
    choose=np.zeros((length,length),dtype=np.int64)
    for j in range(length):
        ii=np.arange(j+1);vv=2*(facval[j]-facval[ii]-facval[j-ii])
        uu=facunit[j]*invunit[ii]%work;uu=uu*invunit[j-ii]%work
        choose[j,:j+1]=(uu*uu%work)*ppows[vv]%work
    comps=[np.zeros(length,dtype=np.int64)];comps[0][0]=1
    for n in range(1,2+ninput):comps.append(choose.dot(comps[-1])%work)
    del choose
    print(json.dumps({'stage':'dimension_coefficients_ready','elapsed_seconds':round(time.time()-started,3)}),flush=True)
    matrices={}
    for K in [N,2*N,p*N,2*p*N]:
        mat=np.zeros((K,K),dtype=np.int64)
        for j in range(K):
            ii=np.arange(j+1);aa=2*K-j+ii-1;bb=2*K-2*j-1;ell=j-ii
            vv=facval[aa]-facval[bb]-facval[ell]-2*facval[ii]
            assert min(vv)>=0
            uu=facunit[aa]*invunit[bb]%work;uu=uu*invunit[ell]%work
            uu=uu*(invunit[ii]*invunit[ii]%work)%work
            row=uu*ppows[vv]%work;mat[j,:j+1]=np.where(ell%2,(-row)%work,row)
        matrices[K]=mat
    print(json.dumps({'stage':'integer_matrices_ready','elapsed_seconds':round(time.time()-started,3)}),flush=True)
    def av_for(n):
        bj=comps[n-1]
        return {K:v.trim((mat.dot(bj[:K])%work).tolist()) for K,mat in matrices.items()}
    def coeff_for(av):
        aa,bb,ap,bp=[av[K] for K in [N,2*N,p*N,2*p*N]]
        ar=[v.taylor(aa,r,p,work) for r in range(degree+1)]
        br=[v.taylor(bb,r,p,work) for r in range(degree+1)]
        a0pow=[v.power(ar[0],r,work) for r in range(degree+2)]
        T=[]
        for r in range(degree+1):
            tr=v.mul(br[r],a0pow[r],work)
            for i in range(1,r+1):tr=v.add(tr,v.scale(v.mul(v.mul(ar[i],T[r-i],work),a0pow[i-1],work),-1,work),work)
            T.append(tr)
        qs=[]
        for r in range(degree+1):
            num=v.add(v.scale(v.mul(bp,a0pow[degree+1],work),comb(N,r),work),v.scale(v.mul(v.mul(ap,T[r],work),a0pow[degree-r],work),-1,work),work)
            num=v.scale(num,p**r,work)
            assert all(c%(p*N)==0 for c in num)
            qs.append(v.trim([(c//(p*N))%ymod for c in num]))
        return qs
    avcache={};samples=[]
    for n in range(2,2+ninput):
        av=av_for(n);avcache[n]=av;samples.append(coeff_for(av))
    width=max(len(q) for s0 in samples for q in s0)
    diffs=[[[q[j] if j<len(q) else 0 for j in range(width)] for q in s0] for s0 in samples]
    newton=[];ordinary=[];bad=[]
    for r in range(ninput):
        dr=diffs[0];vr=v.vp(factorial(r),p)
        need=min(vr,k)
        for q in range(degree+1):
            for j,c in enumerate(dr[q]):
                if c%p**need:bad.append([r,q,j,c,need])
        newton.append([v.trim(q) for q in dr])
        if r<nmodel:
            unit=factorial(r)//p**vr;iv=pow(unit,-1,M)
            ordinary.append([[((c//p**vr)*iv)%M for c in q] for q in dr])
        diffs=[[[ (diffs[i+1][q][j]-diffs[i][q][j])%ymod for j in range(width)] for q in range(degree+1)] for i in range(len(diffs)-1)]
    if bad:
        print(json.dumps({'stage':'model_integrality_failed','bad':bad[:20]}),flush=True)
        raise AssertionError('ordinary model fails')
    print(json.dumps({'stage':'ordinary_model_certified','model_degree':nmodel-1,'elapsed_seconds':round(time.time()-started,3)}),flush=True)
    def model_at(n):
        out=[[0]*width for _ in range(degree+1)];fall=1;x=n-2
        for r,dr in enumerate(ordinary):
            if r:fall=(fall*(x-r+1))%M
            for q in range(degree+1):
                for j,c in enumerate(dr[q]):out[q][j]=(out[q][j]+fall*c)%M
        return [v.trim(q) for q in out]
    for n,actual in zip(range(2,2+ninput),samples):
        assert all(v.trim([c%ymod for c in q])==aa for q,aa in zip(model_at(n),actual))
    weights=[r-1-v.vp(r,p) for r in range(2,degree+1)]
    profiles=list(v.profiles(weights,k-1))
    maxden=max([1]+[1+sum((i+2)*c for i,c in enumerate(pr)) for pr in profiles])
    def root_from_model(qs):
        A,B=qs[:2];bpows=[v.power(B,r,M) for r in range(maxden+1)];apows=[v.power(A,r,M) for r in range(maxden+1)]
        Y=v.scale(v.mul(A,bpows[maxden-1],M),-1,M)
        for pr in profiles:
            count=sum(pr);Rsum=sum((i+2)*c for i,c in enumerate(pr));zexp=Rsum-count+1
            den=factorial(zexp)
            for c in pr:den*=factorial(c)
            assert factorial(Rsum)%den==0
            tree=factorial(Rsum)//den
            if v.vp(tree,p)+sum(w*c for w,c in zip(weights,pr))>=k:continue
            term=apows[zexp]
            for q,c in zip(qs[2:],pr):
                if c:term=v.mul(term,v.power(q,c,M),M)
            term=v.mul(term,bpows[maxden-Rsum-1],M)
            Y=v.add(Y,v.scale(term,(-1 if (Rsum+1)%2 else 1)*tree,M),M)
        return Y,bpows[maxden]
    reports=[]
    partial={'p':p,'s':s,'work_modulus':work,'input_dimension_count':ninput,'model_degree':nmodel-1,'input_lift_coefficients':samples,'newton_coefficients':newton,'ordinary_model_coefficients_in_falling_basis':ordinary,'lagrange_profiles':profiles,'fixed_denominator_exponent':maxden}
    Path(__file__).resolve().parent.joinpath(f'newton_s{s}_model_checkpoint.json').write_text(json.dumps(partial,separators=(',',':'))+'\n')
    if model_only:
        print(json.dumps({'stage':'model_only_complete','level':s,'input_dimensions':ninput,'model_degree':nmodel-1,'all_input_difference_integrality_checked':True,'checkpoint':str(Path(__file__).resolve().parent/f'newton_s{s}_model_checkpoint.json')}),flush=True)
        return
    for n in range(2,2+ntarget):
        qs=model_at(n);Y,Z=root_from_model(qs)
        assert Z[0]%p!=0
        root=[0]
        for r,q in enumerate(qs):root=v.add(root,v.mul(v.mul(q,v.power(Y,r,ymod),ymod),v.power(Z,degree-r,ymod),ymod),ymod)
        assert root==[0],('root',n)
        av=avcache[n];ap=av[p*N];ar=[v.taylor(av[N],r,p,M) for r in range(prec)]
        f=[comb(2*j,j)*int(comps[n][j])%M for j in range(p**s)];fs=f[:p**(s-1)]
        L=[]
        for r in range(prec):
            lhs=v.taylor(fs,r,p,M)
            if r:lhs=v.add(lhs,v.scale(v.taylor(fs,r-1,p,M),N,M),M)
            L.append(v.add(v.mul(ap,lhs,M),v.scale(v.mul(f,ar[r],M),-1,M),M))
        Q=v.scale(Y,p,M);target=[0];bound=0
        for r,c in enumerate(L):
            term=v.mul(v.mul(c,v.power(Q,r,M),M),v.power(Z,prec-1-r,M),M)
            bound=max(bound,len(term)-1);target=v.add(target,term,M)
        assert target==[0],('target',n,[(i,c) for i,c in enumerate(target) if c][:10])
        report={'dimension':n,'root_identity_zero':True,'target_identity_zero':True,'target_degree_bound':bound,'Y_degree':len(Y)-1,'Z_degree':len(Z)-1}
        reports.append(report)
        Path(__file__).resolve().parent.joinpath(f'newton_s{s}_target_progress.json').write_text(json.dumps(reports,indent=2)+'\n')
        print(json.dumps({'stage':'dimension_target_checked',**report,'elapsed_seconds':round(time.time()-started,3)}),flush=True)
    payload={'p':p,'s':s,'m':1,'work_modulus':work,'target_modulus':M,'input_dimensions':list(range(2,2+ninput)),'input_lift_coefficients':samples,'newton_coefficients':newton,'ordinary_model_coefficients_in_falling_basis':ordinary,'lagrange_profiles':profiles,'fixed_denominator_exponent':maxden,'target_reports':reports}
    packed=json.dumps(payload,sort_keys=True,separators=(',',':'))
    summary={'theorem':f'For all octahedral dimensions n>=2, every coefficient of E_(3,{s},1) is divisible by 3^{2*s}.','prime':p,'level':s,'multiplier':1,'target_modulus':M,'work_modulus':work,'input_dimension_count':ninput,'ordinary_model_degree':nmodel-1,'target_dimension_count':ntarget,'all_input_difference_integrality_checked':True,'all_polynomial_coefficients_checked':True,'power_series_cutoff':None,'certificate_sha256':hashlib.sha256(packed.encode()).hexdigest(),'elapsed_seconds':round(time.time()-started,3)}
    base=Path(__file__).resolve().parent
    (base/f'all_dimensions_newton_s{s}_certificate.json').write_text(json.dumps(payload,indent=2)+'\n')
    (base/f'all_dimensions_newton_s{s}_summary.json').write_text(json.dumps(summary,indent=2)+'\n')
    print(json.dumps(summary,indent=2))
if __name__=='__main__':
    parser=argparse.ArgumentParser();parser.add_argument('--s',type=int,required=True);parser.add_argument('--model-only',action='store_true');args=parser.parse_args();main(args.s,args.model_only)

\end{lstlisting}
\clearpage
\section{Compact Hensel roots and complete target certificates}\label{OP:source-appendix-4}
Save this source as \texttt{compact\_hensel.py} in the same directory.
It checks the final entire root and target polynomial identities,
and saves their explicit coefficient arrays. The optional summary
script in the source archive computes file identifiers only and is
not a dependency of any mathematical claim.
\begin{lstlisting}[language=Python]
from pathlib import Path
import json,time,argparse
from math import comb
import numpy as np
import verify_dimension_families as v

def invseries(b,n,p):
    g=[pow(b[0],-1,p)]
    while len(g)<n:
        k=min(n,2*len(g));bg=v.mul(b[:k],g,p)[:k]
        h=[(-x)%p for x in bg];h[0]=(h[0]+2)%p
        g=v.mul(g,h,p)[:k];g+= [0]*(k-len(g))
    return g

def divide(a,b,p=3):
    a=v.trim([x%p for x in a]);b=v.trim([x%p for x in b])
    if len(a)<len(b):return [0],a
    count=len(a)-len(b)+1
    if count<=24:
        rem=a[:];q=[0]*count;iv=pow(b[-1],-1,p)
        for k in range(count-1,-1,-1):
            c=rem[k+len(b)-1]*iv%p;q[k]=c
            if c:
                for j,x in enumerate(b):rem[k+j]=(rem[k+j]-c*x)%p
        return v.trim(q),v.trim(rem)
    revb=list(reversed(b));reva=list(reversed(a))[:count]
    qrev=v.mul(reva,invseries(revb,count,p),p)[:count];qrev +=[0]*(count-len(qrev))
    q=v.trim(list(reversed(qrev)));rem=v.add(a,v.scale(v.mul(b,q,p),-1,p),p)
    assert len(rem)<len(b)
    return q,rem

def gcd(a,b,p=3):
    while b!=[0]:a,b=b,divide(a,b,p)[1]
    return v.scale(a,pow(a[-1],-1,p),p)

def reduce_fraction(a,b,p=3):
    a=v.trim([x%p for x in a]);b=v.trim([x%p for x in b])
    if a==[0]:return [0],[1]
    g=gcd(a,b,p);aa,ar=divide(a,g,p);bb,br=divide(b,g,p)
    assert ar==[0] and br==[0] and bb[0]
    iv=pow(bb[0],-1,p)
    return v.scale(aa,iv,p),v.scale(bb,iv,p)

def root(qs,p,k,verbose=False):
    d=len(qs)-1;A,B=reduce_fraction(v.scale(qs[0],-1,p),qs[1],p)
    for j in range(1,k):
        mod=p**(j+1);S=[0]
        for r,q in enumerate(qs):S=v.add(S,v.mul(v.mul(q,v.power(A,r,mod),mod),v.power(B,d-r,mod),mod),mod)
        assert all(c%(p**j)==0 for c in S)
        E=v.trim([(-c//p**j)%p for c in S]);D=v.mul(v.power(B,d,p),qs[1],p)
        C,H=reduce_fraction(E,D,p)
        A=v.add(v.mul(A,H,mod),v.scale(v.mul(C,B,mod),p**j,mod),mod);B=v.mul(B,H,mod)
        if verbose:print(json.dumps({'digit':j+1,'numdeg':len(A)-1,'dendeg':len(B)-1}),flush=True)
    mod=p**k;S=[0]
    for r,q in enumerate(qs):S=v.add(S,v.mul(v.mul(q,v.power(A,r,mod),mod),v.power(B,d-r,mod),mod),mod)
    assert S==[0]
    return A,B

def inputs(s,maxn):
    p=3;N=p**(2*s-2);M=p**(2*s);length=p*N
    assert M*M*length<2**63
    fu=[1]*(2*length);fv=[0]*(2*length)
    for j in range(1,2*length):
        w=j;val=0
        while w%p==0:w//=p;val+=1
        fu[j]=fu[j-1]*w%M;fv[j]=fv[j-1]+val
    inv=np.array([pow(w,-1,M) for w in fu],dtype=np.int64);fu=np.array(fu,dtype=np.int64);fv=np.array(fv,dtype=np.int64)
    ps=np.array([p**e if e<2*s else 0 for e in range(2*int(fv[-1])+1)],dtype=np.int64)
    choose=np.zeros((length,length),dtype=np.int64)
    for j in range(length):
        ii=np.arange(j+1);vv=2*(fv[j]-fv[ii]-fv[j-ii]);uu=fu[j]*inv[ii]%M;uu=uu*inv[j-ii]%M;choose[j,:j+1]=(uu*uu%M)*ps[vv]%M
    b=[np.zeros(length,dtype=np.int64)];b[0][0]=1
    for _ in range(maxn):b.append(choose.dot(b[-1])%M)
    mats={}
    for K in [N,p*N]:
        mat=np.zeros((K,K),dtype=np.int64)
        for j in range(K):
            ii=np.arange(j+1);aa=2*K-j+ii-1;bb=2*K-2*j-1;ell=j-ii;vv=fv[aa]-fv[bb]-fv[ell]-2*fv[ii]
            assert min(vv)>=0
            uu=fu[aa]*inv[bb]%M;uu=uu*inv[ell]%M;uu=uu*(inv[ii]*inv[ii]%M)%M;row=uu*ps[vv]%M;mat[j,:j+1]=np.where(ell%2,(-row)%M,row)
        mats[K]=mat
    return b,mats

def main(s,dimensions):
    start=time.time();base=Path(__file__).resolve().parent
    c=json.loads((base/f'newton_s{s}_model_checkpoint.json').read_text());p=3;k=2*s-1;M=p**(2*s);N=p**(2*s-2);d=len(c['input_lift_coefficients'][0])-1
    b,mats=inputs(s,max(dimensions));reports=[]
    for n in dimensions:
        qs=c['input_lift_coefficients'][n-2];Y,Z=root(qs,p,k,True)
        av={K:v.trim((mat.dot(b[n-1][:K])%M).tolist()) for K,mat in mats.items()}
        ar=[v.taylor(av[N],r,p,M) for r in range(2*s)];f=[comb(2*j,j)*int(b[n][j])%M for j in range(p**s)];fs=f[:p**(s-1)];target=[0];bound=0;L=[]
        for r in range(2*s):
            lhs=v.taylor(fs,r,p,M)
            if r:lhs=v.add(lhs,v.scale(v.taylor(fs,r-1,p,M),N,M),M)
            lr=v.add(v.mul(av[p*N],lhs,M),v.scale(v.mul(f,ar[r],M),-1,M),M);L.append(lr)
            term=v.mul(v.mul(lr,v.power(v.scale(Y,p,M),r,M),M),v.power(Z,2*s-1-r,M),M);bound=max(bound,len(term)-1);target=v.add(target,term,M)
        assert target==[0],('target',n,[(i,x) for i,x in enumerate(target) if x][:10])
        rep={'dimension':n,'p':p,'s':s,'m':1,'target_modulus':M,'root_identity_zero':True,'target_identity_zero':True,'target_degree_bound':bound,'Y_degree':len(Y)-1,'Z_degree':len(Z)-1}
        payload={'report':rep,'Q':qs,'Y':Y,'Z':Z,'L':L}
        (base/f'compact_octa{n}_s{s}_certificate.json').write_text(json.dumps(payload,separators=(',',':'))+'\n');reports.append(rep)
        (base/f'compact_s{s}_reports.json').write_text(json.dumps(reports,indent=2)+'\n')
        print(json.dumps({'stage':'target_passed',**rep,'elapsed_seconds':round(time.time()-start,3)}),flush=True)
if __name__=='__main__':
    parser=argparse.ArgumentParser();parser.add_argument('--s',type=int,required=True);parser.add_argument('--n',type=int,nargs='+',required=True);a=parser.parse_args();main(a.s,a.n)

\end{lstlisting}

\clearpage
\section{Complete fifth-level certificate source and independent evaluator}
\label{app:level5-programs}
These programs reproduce the independent fifth-level certificate of
Section~\ref{AP:source-main-7}. They are not premises of the symbolic
all-level proof. Python~3, NumPy, and the system GMP shared library are
required; no \texttt{gmpy2} package is used. The polynomial module
\texttt{verify\_dimension\_families.py} and the rational Hensel module
\texttt{compact\_hensel.py} are printed in
Appendices~\ref{OP:source-appendix-1} and~\ref{OP:source-appendix-4}.

Save those two modules in \texttt{research/octa\_new\_certificates/}.
Save the first four programs below in the sibling
\texttt{research/octa\_level5/} directory. Save the independent evaluator
in \texttt{research/octa\_all\_level\_closure/}. This relative directory
layout matches the imports in the printed sources. From the
\texttt{octa\_level5} directory run
\begin{verbatim}
python3 scaled_convolution.py
python3 linear_reduce.py
python3 direct_certificate.py --test
python3 uniform_level.py --s 5 --stage all --recompute
\end{verbatim}
Then execute \texttt{verify\_saved\_level5.py} from its own directory.
The generator recomputes the 39 input dimensions, the ordinary
dimension model, and the 24 target certificates. The independent
evaluator reconstructs the target polynomials and sums all homogeneous
terms separately from the generator's root solver and Horner routine.
The complete 24-target independent re-evaluation passed in the
preceding audit. This edition preserves the exact sources used in
that audit; it does not represent that run as a new external review.
The compressed coefficient arrays and JSON payloads are generated
outputs and are not needed to compile this document.

The final short listing generates the fourth-level summary and its
archived identifiers. Save \path{summarize_level4_certificate.py} in
\path{research/octa_new_certificates/}.

\subsection{scaled\_convolution.py}\label{app:level5-code-1}
\noindent Complete source; 109 lines.
\begin{lstlisting}[language=Python]
"""Exact p=3 Bessel/negative-vertex transforms by ordinary convolution.

GMP is used only for exact nonnegative integer multiplication.  The
Kronecker base strictly exceeds every possible convolution coefficient.
"""
from pathlib import Path
from math import comb
import ctypes,ctypes.util,sys,time,json
sys.path.insert(0,str(Path(__file__).resolve().parents[1]/'octa_new_certificates'))
import verify_dimension_families as v

class MPZ(ctypes.Structure):
    _fields_=[('alloc',ctypes.c_int),('size',ctypes.c_int),('limbs',ctypes.c_void_p)]
gmp=ctypes.CDLL(ctypes.util.find_library('gmp'))
zinit=getattr(gmp,'__gmpz_init');zclear=getattr(gmp,'__gmpz_clear')
zimport=getattr(gmp,'__gmpz_import');zexport=getattr(gmp,'__gmpz_export');zmul=getattr(gmp,'__gmpz_mul')
for fn in [zinit,zclear]:fn.argtypes=[ctypes.POINTER(MPZ)]
zmul.argtypes=[ctypes.POINTER(MPZ)]*3
zimport.argtypes=[ctypes.POINTER(MPZ),ctypes.c_size_t,ctypes.c_int,ctypes.c_size_t,ctypes.c_int,ctypes.c_size_t,ctypes.c_void_p]
zexport.argtypes=[ctypes.c_void_p,ctypes.POINTER(ctypes.c_size_t),ctypes.c_int,ctypes.c_size_t,ctypes.c_int,ctypes.c_size_t,ctypes.POINTER(MPZ)]
zexport.restype=ctypes.c_void_p

def mul(a,b,mod,limit=None):
    a=v.trim([int(x)%mod for x in a]);b=v.trim([int(x)%mod for x in b])
    if a==[0] or b==[0]:return [0]
    width=(min(len(a),len(b)).bit_length()+2*mod.bit_length()+7)//8
    assert (1<<(8*width))>min(len(a),len(b))*(mod-1)**2
    aa=b''.join(x.to_bytes(width,'little') for x in a);bb=b''.join(x.to_bytes(width,'little') for x in b)
    za,zb,zc=MPZ(),MPZ(),MPZ()
    for z in [za,zb,zc]:zinit(ctypes.byref(z))
    try:
        zimport(ctypes.byref(za),len(aa),-1,1,0,0,aa);zimport(ctypes.byref(zb),len(bb),-1,1,0,0,bb)
        zmul(ctypes.byref(zc),ctypes.byref(za),ctypes.byref(zb))
        length=len(a)+len(b)-1;raw=ctypes.create_string_buffer(length*width);count=ctypes.c_size_t()
        zexport(raw,ctypes.byref(count),-1,1,0,0,ctypes.byref(zc))
        data=raw.raw
        if limit is not None:length=min(length,limit)
        return v.trim([int.from_bytes(data[i*width:(i+1)*width],'little')%mod for i in range(length)])
    finally:
        for z in [za,zb,zc]:zclear(ctypes.byref(z))

class Inputs:
    def __init__(self,s=None,length=None,exponent=None,indices=None):
        if s is not None:
            N=3**(2*s-2);length=6*N;exponent=4*s-2;indices=[N,2*N,3*N,6*N]
        self.length=length;self.exponent=exponent;self.work=3**exponent;self.indices=indices
        self.val=[0]*(2*length)
        for j in range(1,2*length):
            w=j;vv=0
            while w%3==0:w//=3;vv+=1
            self.val[j]=self.val[j-1]+vv
        self.h={K:max(K-l-self.val[2*K-l-1]+self.val[l] for l in range(K)) for K in indices}
        self.H=max(self.h.values());self.mod=3**(exponent+self.H)
        self.unit=[1]*(2*length)
        for j in range(1,2*length):
            w=j
            while w%3==0:w//=3
            self.unit[j]=self.unit[j-1]*w%self.mod
        self.inv=[pow(w,-1,self.mod) for w in self.unit]
        self.J=[pow(3,j-2*self.val[j],self.mod)*self.inv[j]**2%self.mod for j in range(length)]
        self.alpha={}
        for K in indices:
            h=self.h[K];arr=[]
            for l in range(K):
                ev=l-K+h+self.val[2*K-l-1]-self.val[l]
                assert ev>=0
                arr.append(((-1 if l%2 else 1)*pow(3,ev,self.mod)*self.unit[2*K-l-1]*self.inv[l])%self.mod)
            self.alpha[K]=arr

    def vertices(self,Jpower):
        result={}
        for K in self.indices:
            vals=mul(self.alpha[K],Jpower[:K],self.mod,K);vals+=[0]*(K-len(vals));out=[]
            for j,c in enumerate(vals):
                ev=K-j-self.h[K]-self.val[2*K-2*j-1]
                if ev<0:
                    den=3**(-ev);assert c%den==0,(K,j,ev,c);c//=den
                else:c*=3**ev
                out.append(c*self.inv[2*K-2*j-1]%self.work)
            assert out[0]==1
            result[K]=v.trim(out)
        return result

    def composition(self,Jpower,length=None,mod=None):
        length=length or self.length;mod=mod or self.work;out=[]
        for j in range(length):
            den=3**(j-2*self.val[j]);c=Jpower[j] if j<len(Jpower) else 0
            assert c%den==0
            assert self.mod% (mod*den)==0
            out.append((c//den)*self.unit[j]**2%mod)
        return out

def test():
    import random
    for mod in [3,81,3**18,3**29]:
        for na,nb in [(3,4),(32,31),(100,153),(1000,700)]:
            rng=random.Random(na+nb);a=[rng.randrange(mod) for _ in range(na)];b=[rng.randrange(mod) for _ in range(nb)]
            assert mul(a,b,mod)==v.mul(a,b,mod)
    ins=Inputs(length=100,exponent=14,indices=[1,2,3,7,9,17,27,54])
    current=[1]
    for dim in range(1,8):
        current=mul(current,ins.J,ins.mod,ins.length)
        exact=v.composition_squares(dim,100)
        assert ins.composition(current,length=30,mod=3**8)==[a%3**8 for a in exact[:30]]
        av=ins.vertices(current)
        for K in ins.indices:assert av[K]==v.vertex(K,exact,ins.work),(dim,K)
    print('scaled convolution verified against exact original formulas')

if __name__=='__main__':test()
\end{lstlisting}

\subsection{linear\_reduce.py}\label{app:level5-code-2}
\noindent Complete source; 85 lines.
\begin{lstlisting}[language=Python]
"""Exact fast gcd reduction for a denominator equal to a linear power in F_3[u]."""
import numpy as np
from math import comb
from pathlib import Path
import sys
sys.path.insert(0,str(Path(__file__).resolve().parents[1]/'octa_new_certificates'))
import compact_hensel as ch
import scaled_convolution as sc
v=sc.v

def linear_power(c,e):
    a=np.zeros(e+1,dtype=np.int64);a[0]=1;step=1;end=1
    while e:
        digit=e%3;e//=3
        if digit:
            old=a[:end].copy()
            a[step:step+end]=(digit*c*old)%3
            if digit==2:a[2*step:2*step+end]=(c*c*old)%3
            end+=digit*step
        step*=3
    return a.tolist()

def translate3(a,t):
    t%=3
    if not t:return list(a)
    length=1
    while length<len(a):length*=3
    out=np.zeros(length,dtype=np.int64);out[:len(a)]=a
    step=1
    while step<length:
        blocks=out.reshape(-1,3,step)
        b1=blocks[:,1,:].copy();b2=blocks[:,2,:].copy()
        blocks[:,0,:]=(blocks[:,0,:]+t*b1+t*t*b2)%3
        blocks[:,1,:]=(b1+2*t*b2)%3
        step*=3
    return v.trim(out.tolist())

def reduce_fraction(a,b,p=3):
    if p!=3:return ch.reduce_fraction(a,b,p)
    a=v.trim([int(x)%p for x in a]);b=v.trim([int(x)%p for x in b])
    if a==[0]:return [0],[1]
    assert b[0]
    iv=pow(b[0],-1,p);a=v.scale(a,iv,p);b=v.scale(b,iv,p)
    if b==[1]:return a,b
    deg=len(b)-1
    multiplicities={}
    for c in [1,2]:
        t=(-pow(c,-1,3))%3
        bt=translate3(b,t)
        multiplicities[c]=next((i for i,x in enumerate(bt) if x),len(bt))
    if sum(multiplicities.values())!=deg:return ch.reduce_fraction(a,b,p)
    C=a;H=[1]
    for c,hb in multiplicities.items():
        if not hb:continue
        root=(-pow(c,-1,3))%3
        at=translate3(C,root);ha=next((i for i,x in enumerate(at) if x),len(at));h=min(ha,hb)
        C=v.scale(translate3(at[h:],-root),pow(pow(c,-1,3),h,3),3)
        H=v.mul(H,linear_power(c,hb-h),3)
    return C,H

def test():
    import random
    rng=random.Random(0)
    for size in [1,2,5,10,50,100,270]:
        a=[rng.randrange(3) for _ in range(size)]
        for t in [1,2]:
            expected=[sum(a[j]*comb(j,k)*pow(t,j-k,3) for j in range(k,size))%3 for k in range(size)]
            assert translate3(a,t)==v.trim(expected)
    for c in [1,2]:
        for e in range(300):assert linear_power(c,e)==v.power([1,c],e,3)
        for e in [10,40,200,1000]:
            b=linear_power(c,e)
            for h in [0,5,e//2,e]:
                a=v.mul([1,0,2,1],linear_power(c,h),3)
                C,H=reduce_fraction(a,b)
                assert v.mul(C,b,3)==v.mul(a,H,3)
                assert ch.gcd(C,H)==[1]
    for e1,e2,h1,h2 in [(10,20,5,10),(1000,500,555,499),(120,140,130,0)]:
        b=v.mul(linear_power(1,e1),linear_power(2,e2),3)
        a=v.mul(v.mul(linear_power(1,h1),linear_power(2,h2),3),[1,0,1],3)
        C,H=reduce_fraction(a,b)
        assert v.mul(C,b,3)==v.mul(a,H,3) and ch.gcd(C,H)==[1]
    print('ternary translation and exact linear-power rational reduction verified')

if __name__=='__main__':test()
\end{lstlisting}

\subsection{direct\_certificate.py}\label{app:level5-code-3}
\noindent Complete source; 106 lines.
\begin{lstlisting}[language=Python]
#!/usr/bin/env python3
"""Exact negative-vertex certificates using a directly cleared lift equation."""
from pathlib import Path
from math import comb
import sys,json,time,argparse,hashlib
import scaled_convolution as sc
v=sc.v
v.mul=sc.mul
import compact_hensel as ch
import linear_reduce as lr
BASE=Path(__file__).resolve().parent
OLD=BASE.parent/'octa_new_certificates'

def degree_for(k):
    return max(r for r in range(1,3*k+1) if r-1-(len_digits(r,3)-1)<k)

def len_digits(n,p):
    z=0
    while n:n//=p;z+=1
    return z

def equation(av,s):
    p=3;k=2*s-1;N=p**(2*s-2);work=p**(4*s-2);ymod=p**k;d=degree_for(k)
    ar=[v.taylor(av[N],r,p,work) for r in range(d+1)]
    br=[v.taylor(av[2*N],r,p,work) for r in range(d+1)]
    out=[]
    for r in range(d+1):
        aa=[0]
        for i in range(r+1):aa=v.add(aa,v.scale(ar[i],comb(N,r-i),work),work)
        num=v.add(v.mul(av[6*N],aa,work),v.scale(v.mul(av[3*N],br[r],work),-1,work),work)
        num=v.scale(num,p**r,work)
        assert all(c%(p*N)==0 for c in num),('direct integrality',r)
        out.append(v.trim([c//(p*N)%ymod for c in num]))
    assert out[0][0]==0 and out[1][0]==1
    assert all(all(c%3==0 for c in q) for q in out[2:])
    return out

def cleared_eval(qs,Y,Z,mod):
    # Homogeneous Horner form of sum Q_r Y^r Z^(d-r).
    d=len(qs)-1;out=qs[d];zp=[1]
    for r in range(d-1,-1,-1):
        zp=v.mul(zp,Z,mod)
        out=v.add(v.mul(out,Y,mod),v.mul(qs[r],zp,mod),mod)
    return v.trim([c%mod for c in out])

def root(qs,p,k,verbose=False):
    d=len(qs)-1;A,B=lr.reduce_fraction(v.scale(qs[0],-1,p),qs[1],p)
    if verbose:print(json.dumps({'stage':'root_digit','digit':1,'numdeg':len(A)-1,'dendeg':len(B)-1}),flush=True)
    for j in range(1,k):
        mod=p**(j+1);S=cleared_eval(qs,A,B,mod)
        assert all(c%p**j==0 for c in S)
        E=v.trim([(-c//p**j)%p for c in S]);D=v.mul(v.power(B,d,p),qs[1],p)
        C,H=lr.reduce_fraction(E,D,p)
        A=v.add(v.mul(A,H,mod),v.scale(v.mul(C,B,mod),p**j,mod),mod);B=v.mul(B,H,mod)
        if verbose:print(json.dumps({'stage':'root_digit','digit':j+1,'numdeg':len(A)-1,'dendeg':len(B)-1}),flush=True)
    assert cleared_eval(qs,A,B,p**k)==[0]
    return A,B

def target_polynomials(av,f,s):
    p=3;N=p**(2*s-2);M=p**(2*s);fs=f[:p**(s-1)];L=[]
    for r in range(2*s):
        lhs=v.taylor(fs,r,p,M)
        if r:lhs=v.add(lhs,v.scale(v.taylor(fs,r-1,p,M),N,M),M)
        lr=v.add(v.mul(av[3*N],lhs,M),v.scale(v.mul(f,v.taylor(av[N],r,p,M),M),-1,M),M)
        L.append(lr)
    return L

def run(s,dimension):
    start=time.time();p=3;k=2*s-1;M=p**(2*s);ins=sc.Inputs(s=s);Jpower=[1]
    print(json.dumps({'stage':'setup','level':s,'dimension':dimension,'length':ins.length,'extra_precision':ins.H,'lift_degree':degree_for(k)}),flush=True)
    for j in range(dimension-1):Jpower=sc.mul(Jpower,ins.J,ins.mod,ins.length)
    av=ins.vertices(Jpower)
    print(json.dumps({'stage':'vertices_ready','seconds':round(time.time()-start,3)}),flush=True)
    qs=equation(av,s)
    checkpoint={'s':s,'dimension':dimension,'Q':qs,'vertices':av}
    (BASE/f'direct_octa{dimension}_s{s}_inputs.json').write_text(json.dumps(checkpoint,separators=(',',':'))+'\n')
    print(json.dumps({'stage':'equation_ready','degrees':[len(q)-1 for q in qs],'seconds':round(time.time()-start,3)}),flush=True)
    Y,Z=root(qs,p,k,True)
    Jnext=sc.mul(Jpower[:3**s],ins.J[:3**s],ins.mod,3**s)
    b=ins.composition(Jnext,3**s,M);f=[comb(2*j,j)*b[j]%M for j in range(3**s)]
    L=target_polynomials(av,f,s);zero=cleared_eval(L,v.scale(Y,p,M),Z,M)
    assert zero==[0],('target',[(i,c) for i,c in enumerate(zero) if c][:10])
    rep={'dimension':dimension,'p':p,'s':s,'m':1,'target_modulus':M,'root_identity_zero':True,'target_identity_zero':True,'target_degree_bound':max(len(q)-1+r*(len(Y)-1)+(2*s-1-r)*(len(Z)-1) for r,q in enumerate(L)),'Y_degree':len(Y)-1,'Z_degree':len(Z)-1,'elapsed_seconds':round(time.time()-start,3)}
    payload={'report':rep,'Q':qs,'Y':Y,'Z':Z,'L':L}
    packed=json.dumps(payload,separators=(',',':'));rep['certificate_sha256']=hashlib.sha256(packed.encode()).hexdigest()
    (BASE/f'direct_octa{dimension}_s{s}_certificate.json').write_text(packed+'\n')
    (BASE/f'direct_octa{dimension}_s{s}_report.json').write_text(json.dumps(rep,indent=2)+'\n')
    print(json.dumps({'stage':'target_passed',**rep}),flush=True)
    return rep

def test():
    for n,s,filename in [(4,2,'certificate_octa4_p3_s2_m1.json'),(3,3,'certificate_octa3_p3_s3.json')]:
        old=json.loads((OLD/filename).read_text());ins=sc.Inputs(s=s);jp=[1]
        for j in range(n-1):jp=sc.mul(jp,ins.J,ins.mod,ins.length)
        av=ins.vertices(jp)
        assert {str(k):a for k,a in av.items()}==old['vertex_polynomials'],('vertex',n,s)
        qs=equation(av,s);Y,Z=old['lift_y_numerator'],old['lift_y_denominator']
        assert cleared_eval(qs,Y,Z,3**(2*s-1))==[0],('root',n,s)
        y,z=root(qs,3,2*s-1)
        assert v.add(v.mul(Y,z,3**(2*s-1)),v.scale(v.mul(y,Z,3**(2*s-1)),-1,3**(2*s-1)),3**(2*s-1))==[0]
        print(json.dumps({'validated_against_existing_certificate':filename,'all_vertex_coefficients_equal':True,'root_rational_functions_equal':True}),flush=True)

if __name__=='__main__':
    pa=argparse.ArgumentParser();pa.add_argument('--s',type=int,default=5);pa.add_argument('--n',type=int,default=4);pa.add_argument('--test',action='store_true');a=pa.parse_args()
    if a.test:test()
    else:run(a.s,a.n)
\end{lstlisting}

\subsection{uniform\_level.py}\label{app:level5-code-4}
\noindent Complete source; 107 lines.
\begin{lstlisting}[language=Python]
#!/usr/bin/env python3
"""Exact all-dimension certificate for one fixed level, using the direct equation."""
from pathlib import Path
from math import comb,factorial
import json,time,hashlib,argparse
import numpy as np
import scaled_convolution as sc
import direct_certificate as dc
v=sc.v
BASE=Path(__file__).resolve().parent

def threshold(e):
    i=z=0
    while z<e:
        i+=1;z+=v.vp(i,3)
    return i

def padded(qs,width):
    a=np.zeros((len(qs),width),dtype=np.int32)
    for r,q in enumerate(qs):a[r,:len(q)]=q
    return a

def samples(s,force=False):
    start=time.time();ins=sc.Inputs(s=s);N=3**(2*s-2);k=2*s-1;M=3**(2*s);d=dc.degree_for(k);ni=threshold(2*k);jp=ins.J
    directory=BASE/f's{s}_samples';directory.mkdir(exist_ok=True)
    print(json.dumps({'stage':'sample_setup','s':s,'ninput':ni,'ntarget':threshold(2*s),'degree':d}),flush=True)
    for n in range(2,2+ni):
        path=directory/f'n{n}.npz'
        if path.exists() and not force:
            print(json.dumps({'stage':'sample_reused','n':n}),flush=True)
        else:
            av=ins.vertices(jp);qs=dc.equation(av,s)
            jnext=sc.mul(jp[:3**s],ins.J[:3**s],ins.mod,3**s)
            b=ins.composition(jnext,3**s,M);f=[comb(2*j,j)*b[j]%M for j in range(3**s)]
            arr=padded(qs,9*N-3)
            np.savez_compressed(path,Q=arr,A_N=np.array(av[N],dtype=np.int64)%M,A_3N=np.array(av[3*N],dtype=np.int64)%M,F=np.array(f,dtype=np.int64))
            print(json.dumps({'stage':'sample_saved','n':n,'seconds':round(time.time()-start,3)}),flush=True)
        if n<1+ni:jp=sc.mul(jp,ins.J,ins.mod,ins.length)
    return directory

def model(s):
    start=time.time();N=3**(2*s-2);k=2*s-1;M=3**(2*s);ymod=3**k;d=dc.degree_for(k);ni=threshold(2*k);nm=threshold(k);width=9*N-3
    directory=BASE/f's{s}_samples'
    rows=np.stack([np.load(directory/f'n{n}.npz')['Q'] for n in range(2,2+ni)])
    original=rows.copy();newton=np.zeros_like(rows);ordinary=np.zeros((nm,d+1,width),dtype=np.int64)
    checks=[]
    for r in range(ni):
        dr=rows[0];val=v.vp(factorial(r),3);need=min(k,val)
        bad=np.argwhere(dr%3**need)
        assert len(bad)==0,('Newton integrality',r,bad[:5].tolist())
        newton[r]=dr
        if r<nm:
            unit=factorial(r)//3**val;iv=pow(unit,-1,M)
            ordinary[r]=(dr.astype(np.int64)//3**val*iv)%M
        checks.append({'difference_order':r,'required_valuation':need,'passed':True})
        rows=(rows[1:]-rows[:-1])%ymod
    print(json.dumps({'stage':'ordinary_model_integrality_passed','degree':nm-1,'seconds':round(time.time()-start,3)}),flush=True)
    for i,n in enumerate(range(2,2+ni)):
        actual=np.zeros((d+1,width),dtype=np.int64);fall=1
        for r in range(nm):
            if r:fall=fall*(n-2-r+1)%M
            actual=(actual+fall*ordinary[r])%M
        assert np.array_equal(actual%ymod,original[i]),('model evaluation',n)
    output=BASE/f's{s}_ordinary_dimension_model.npz'
    np.savez_compressed(output,newton=newton,ordinary=ordinary)
    report={'s':s,'p':3,'input_dimensions':list(range(2,2+ni)),'input_dimension_count':ni,'model_degree':nm-1,'lift_degree':d,'newton_valuation_checks':checks,'all_input_evaluations_match':True,'sha256':hashlib.sha256(output.read_bytes()).hexdigest(),'elapsed_seconds':round(time.time()-start,3)}
    (BASE/f's{s}_model_report.json').write_text(json.dumps(report,indent=2)+'\n')
    print(json.dumps({'stage':'model_certified',**{key:report[key] for key in ['s','input_dimension_count','model_degree','sha256','elapsed_seconds']}}),flush=True)
    return report

def targets(s,dimensions=None,force=False):
    start=time.time();p=3;k=2*s-1;N=p**(2*s-2);M=p**(2*s);dimensions=dimensions or list(range(2,2+threshold(2*s)));reports=[]
    for n in dimensions:
        certificate_path=BASE/f'direct_octa{n}_s{s}_certificate.json';report_path=BASE/f'direct_octa{n}_s{s}_report.json'
        a=np.load(BASE/f's{s}_samples'/f'n{n}.npz');qs=[v.trim(row.tolist()) for row in a['Q']]
        if certificate_path.exists() and report_path.exists() and not force:
            payload=json.loads(certificate_path.read_text());assert qs==payload['Q'];rep=json.loads(report_path.read_text())
            reports.append(rep);print(json.dumps({'stage':'target_reused','n':n}),flush=True);continue
        Y,Z=dc.root(qs,p,k,True)
        av={N:a['A_N'].tolist(),3*N:a['A_3N'].tolist()};f=a['F'].tolist();L=dc.target_polynomials(av,f,s)
        result=dc.cleared_eval(L,v.scale(Y,p,M),Z,M)
        assert result==[0],('target',n,[(i,c) for i,c in enumerate(result) if c][:10])
        rep={'dimension':n,'p':p,'s':s,'m':1,'target_modulus':M,'root_identity_zero':True,'target_identity_zero':True,'target_degree_bound':max(len(q)-1+r*(len(Y)-1)+(2*s-1-r)*(len(Z)-1) for r,q in enumerate(L)),'Y_degree':len(Y)-1,'Z_degree':len(Z)-1}
        payload={'report':rep,'Q':qs,'Y':Y,'Z':Z,'L':L};packed=json.dumps(payload,separators=(',',':'))
        certificate_path.write_text(packed+'\n');rep['certificate_sha256']=hashlib.sha256(packed.encode()).hexdigest();report_path.write_text(json.dumps(rep,indent=2)+'\n')
        reports.append(rep)
        (BASE/f's{s}_target_reports.json').write_text(json.dumps(reports,indent=2)+'\n')
        print(json.dumps({'stage':'target_passed',**rep,'elapsed_seconds':round(time.time()-start,3)}),flush=True)
    return reports

def summary(s):
    mr=json.loads((BASE/f's{s}_model_report.json').read_text());reports=[]
    for n in range(2,2+threshold(2*s)):
        rep=json.loads((BASE/f'direct_octa{n}_s{s}_report.json').read_text());assert rep['root_identity_zero'] and rep['target_identity_zero'];reports.append(rep)
    sr={'theorem':f'For all octahedral dimensions n>=2, every coefficient of E_(3,{s},1) is divisible by 3^{2*s}.','prime':3,'level':s,'multiplier':1,'target_modulus':3**(2*s),'work_modulus':3**(4*s-2),'input_dimension_count':threshold(4*s-2),'ordinary_model_degree':threshold(2*s-1)-1,'target_dimension_count':threshold(2*s),'lift_degree':dc.degree_for(2*s-1),'all_input_difference_integrality_checked':True,'all_polynomial_coefficients_checked':True,'power_series_cutoff':None,'dimension_model_report':mr,'target_reports':reports}
    (BASE/f'all_dimensions_s{s}_summary.json').write_text(json.dumps(sr,indent=2)+'\n')
    manifest={'level':s,'model':mr['sha256'],'certificates':[{'n':x['dimension'],'sha256':x['certificate_sha256']} for x in reports]}
    manifest['aggregate_sha256']=hashlib.sha256(json.dumps(manifest,sort_keys=True,separators=(',',':')).encode()).hexdigest()
    (BASE/f's{s}_manifest.json').write_text(json.dumps(manifest,indent=2)+'\n')
    print(json.dumps({key:sr[key] for key in ['theorem','input_dimension_count','ordinary_model_degree','target_dimension_count','all_polynomial_coefficients_checked','power_series_cutoff']},indent=2),flush=True)

if __name__=='__main__':
    pa=argparse.ArgumentParser();pa.add_argument('--s',type=int,default=5);pa.add_argument('--stage',choices=['samples','model','targets','summary','all'],default='all');pa.add_argument('--n',type=int,nargs='+');pa.add_argument('--recompute',action='store_true');a=pa.parse_args()
    if a.stage in ['samples','all']:samples(a.s,a.recompute)
    if a.stage in ['model','all']:model(a.s)
    if a.stage in ['targets','all']:targets(a.s,a.n,a.recompute)
    if a.stage in ['summary','all']:summary(a.s)
\end{lstlisting}

\subsection{verify\_saved\_level5.py}\label{app:level5-code-5}
\noindent Complete source; 80 lines.
\begin{lstlisting}[language=Python]
"""Re-read and independently evaluate every archived level-five identity.

This verifier does not invoke the certificate generator, its root solver,
or its homogeneous Horner evaluator.  It builds powers and sums all terms.
Exact modular convolution is the tested packed-integer backend.
"""
from pathlib import Path
from math import comb
import sys, json, time
import numpy as np

ROOT = Path(__file__).resolve().parents[2]
DATA = ROOT / 'research/octa_level5'
sys.path.insert(0, str(DATA))
import scaled_convolution as sc

def trim(a):
    while len(a)>1 and not a[-1]:
        a.pop()
    return a

def add(a,b,m):
    z=[0]*max(len(a),len(b))
    for i,c in enumerate(a): z[i]=c
    for i,c in enumerate(b): z[i]=(z[i]+c)%m
    return trim(z)

def scale(a,c,m):
    return trim([(x*c)%m for x in a])

def powers(a,d,m):
    out=[[1]]
    for _ in range(d): out.append(sc.mul(out[-1],a,m))
    return out

def homogeneous(qs,y,z,m):
    d=len(qs)-1
    yp,zp=powers(y,d,m),powers(z,d,m)
    total=[0]
    for r,q in enumerate(qs):
        term=sc.mul(q,sc.mul(yp[r],zp[d-r],m),m)
        total=add(total,term,m)
    return total

def taylor(a,r,m):
    out=[0]*(3*(len(a)-1)+1)
    for j in range(r,len(a)):
        out[3*j]=int(a[j])*comb(j,r)%m
    return trim(out)

def main():
    started=time.time()
    results=[]
    N=6561; mod=59049
    for n in range(2,26):
        data=json.loads((DATA/f'direct_octa{n}_s5_certificate.json').read_text())
        sample=np.load(DATA/'s5_samples'/f'n{n}.npz')
        qs,y,z,L=data['Q'],data['Y'],data['Z'],data['L']
        assert z[0]%3
        assert qs==[trim(row.tolist()) for row in sample['Q']]
        assert qs[0][0]==0 and qs[1][0]==1
        assert all(all(c%3==0 for c in q) for q in qs[2:])
        assert homogeneous(qs,y,z,19683)==[0],('root',n)
        f=sample['F'].tolist(); small=f[:81]
        a=sample['A_N'].tolist(); ap=sample['A_3N'].tolist()
        rebuilt=[]
        for r in range(10):
            left=taylor(small,r,mod)
            if r: left=add(left,scale(taylor(small,r-1,mod),N,mod),mod)
            lr=add(sc.mul(ap,left,mod),scale(sc.mul(f,taylor(a,r,mod),mod),-1,mod),mod)
            rebuilt.append(lr)
        assert rebuilt==L,('target construction',n)
        assert homogeneous(L,scale(y,3,mod),z,mod)==[0],('target',n)
        results.append({'dimension':n,'root_verified':True,'target_reconstructed':True,'target_verified':True})
        print(json.dumps({'dimension':n,'verified':True,'elapsed_seconds':round(time.time()-started,3)}),flush=True)
    report={'method':'independent sum-of-powers evaluation of saved identities','dimensions':results,'all_24_passed':True,'elapsed_seconds':round(time.time()-started,3)}
    (Path(__file__).resolve().parent/'independent_level5_audit.json').write_text(json.dumps(report,indent=2)+'\n')
    print(json.dumps({'all_24_passed':True,'elapsed_seconds':report['elapsed_seconds']}),flush=True)

if __name__=='__main__': main()
\end{lstlisting}

\subsection{summarize\_level4\_certificate.py}\label{app:level5-code-6}
\noindent Complete source; 29 lines.
\begin{lstlisting}[language=Python]
#!/usr/bin/env python3
"""Assemble deterministic identifiers for the already verified level-four data."""
from pathlib import Path
import json,hashlib
base=Path(__file__).resolve().parent

def sha(data):return hashlib.sha256(json.dumps(data,sort_keys=True,separators=(',',':')).encode()).hexdigest()
model=json.loads((base/'newton_s4_model_checkpoint.json').read_text())
reports=[];certs=[]
for n in range(2,20):
    data=json.loads((base/f'compact_octa{n}_s4_certificate.json').read_text())
    report=data['report']
    assert report['dimension']==n and report['s']==4 and report['p']==3 and report['m']==1
    assert report['root_identity_zero'] and report['target_identity_zero']
    reports.append(report);certs.append({'dimension':n,'sha256':sha(data)})
assert model['input_dimension_count']==30 and model['model_degree']==17
assert model['fixed_denominator_exponent']==19
manifest={'prime':3,'level':4,'multiplier':1,'model_checkpoint_sha256':sha(model),'compact_certificate_sha256s':certs}
summary={'theorem':'For every octahedral dimension n>=2, every coefficient of E_(3,4,1) is divisible by 3^8.',
 'prime':3,'level':4,'multiplier':1,'target_modulus':6561,'lift_y_modulus':2187,'work_modulus':4782969,
 'input_dimension_count':30,'input_dimension_representatives':[2,31],'ordinary_model_degree':17,
 'target_dimension_count':18,'target_dimension_representatives':[2,19],'lift_equation_degree':9,
 'universal_lagrange_denominator_exponent':19,'maximum_compact_target_degree':max(r['target_degree_bound'] for r in reports),
 'all_input_difference_integrality_checked':True,'all_polynomial_coefficients_checked':True,
 'formal_power_series_cutoff':None,
 'universal_dimension_justification':'Ordinary dimension-polynomial model and fixed universal Lagrange numerator/denominator; compact roots agree by uniqueness; 18 consecutive target values suffice by factorial finite differences.',
 'certificate_sha256':sha(manifest),'manifest':manifest,'target_reports':reports}
(base/'all_dimensions_newton_s4_summary.json').write_text(json.dumps(summary,indent=2)+'\n')
print(json.dumps({k:v for k,v in summary.items() if k not in ['manifest','target_reports']},indent=2))
\end{lstlisting}


\clearpage
\section{Source-to-section preservation map}\label{app:coverage}
The letters P, B, E, T identify the original four source manuscripts: P, positive vertex interpolation; B, the revised base manuscript on Hecke traces; E, exceptional/elliptic and sharpness arguments; T, two-cusp descent and auxiliary Cartier methods. Each source mathematical section, including its unnumbered proofs, is retained at the destination below. Section numbers refer to this unified edition. Source introductions have been synchronized with the final theorem scope; bibliography entries are merged by work. Historical program listings and their recorded output are explicitly identified as such.

\begin{longtable}{@{}p{.09\textwidth}p{.74\textwidth}p{.09\textwidth}@{}}
\toprule Source & Mathematical section & Here\\\midrule\endfirsthead
\toprule Source & Mathematical section & Here\\\midrule\endhead
P & The target and the resulting scope & \ref{P:source-main-1}\\
P & Positive vertex interpolation: complete hypercubic and simplicial congruences & \ref{P:source-main-2}\\
P & A matrix block argument for the hyperoctahedral family & \ref{P:source-main-3}\\
P & Verification, sharpness, and scope & \ref{P:source-main-4}\\
P App. & A complete coefficientwise certificate in hyperoctahedral dimension four & \ref{P:source-appendix-1}\\
B & Modular trace methods and their conclusions & \ref{B:source-main-1}\\
B & Modular families of cusp type & \ref{B:source-main-2}\\
B & The Hecke trace & \ref{B:source-main-3}\\
B & Proof of the Main Theorem & \ref{B:source-main-4}\\
B & Proofs of the corollaries & \ref{B:source-main-5}\\
B & Families with an elliptic zero & \ref{B:source-main-6}\\
B & Zagier's case D & \ref{B:source-main-7}\\
B & Limits of the modular method and an elliptic extension & \ref{B:source-main-8}\\
B & Computational verification & \ref{B:source-main-9}\\
B App. & The verification script & \ref{B:source-appendix-1}\\
B App. & Output of the script & \ref{B:source-appendix-2}\\
E & Statements, normalization, and scope & \ref{E:source-main-1}\\
E & Elliptic zeros: a rational trace estimate & \ref{E:source-main-2}\\
E & Zagier D at primes congruent to \texorpdfstring{$\pm1$}{plus or minus one} modulo five & \ref{E:source-main-3}\\
E & Exact coefficient obstructions to larger exponents & \ref{E:source-main-4}\\
E & Relation to the cited literature & \ref{E:source-main-5}\\
E & Exact verification and proof dependencies & \ref{E:source-main-6}\\
T & Scope and exact statements & \ref{T:source-main-1}\\
T & The complete Zagier D congruence & \ref{T:source-main-2}\\
T & The first level in every dimension & \ref{T:source-main-3}\\
T & A canonical negative-power basis in all three families & \ref{T:source-main-4}\\
T & The three-dimensional hyperoctahedral family: a weighted Cartier reduction & \ref{T:source-main-5}\\
T & Exact verification and the role of the auxiliary arguments & \ref{T:source-main-6}\\
\bottomrule
\end{longtable}

\subsection{Explicit proof and statement preservation}
The original source proof environments total $8+20+13+9=50$ in the order P, B, E, T.
The three supplements contribute $9+6+13=28$ additional proof environments,
so all $78$ source proof environments are retained. The global
combination and analytic-ring transfer in the introduction add their own proofs. Their mathematical bodies are preserved; local cross-reference and citation keys are renamed consistently. Theorem and lemma statements are retained with their assumptions. The low-dimensional modular corollary is worded explicitly to keep its $3s$ claim confined to the three-dimensional hypercubic and simplicial families. The elliptic extension remains open.

The following unnumbered arguments are also retained: the positive Cartier coefficient extraction and low-dimensional reductions; the entire independent four-dimensional certificate derivation; elliptic factorial-carry estimates, including the $108t$ boundary repair; the Zagier D residual and both cusp expansions; all three negative-vertex formulas and their completion proviso; the octahedral symmetric-square and Bol identities, telescoping endpoints, primitive integrality, and exact Domb transport. Thus the preservation check is not based only on counts of theorem or proof environments.

\subsection{Editorial reconciliation}
Statements that Zagier D at $p\equiv\pm2\pmod5$, the three-dimensional hyperoctahedral case, or dimensions $n\ge4$ were globally unresolved have been replaced by references to their complete proofs. Statements explaining why a particular auxiliary method alone is insufficient have been retained as method limitations. No ordinary-truncation theorem has been inferred merely from a weighted differential congruence, a negative-power basis congruence, or a finite numerical sample. The historical verification logs retain their original labels with an explicit status annotation.

The bibliography is consolidated, Beukers's modular-forms paper is cited under one key, and the long quotation from its introduction is paraphrased while preserving the mathematical question. Fu's work is retained as contextual literature, not a proof input. The original computational claims are described according to what the supplied programs actually check.

\clearpage

\subsection{The three supplements and the complete-edition additions}
SP denotes the all-odd-prime small-prime supplement, OP the
arbitrary-dimensional hyperoctahedral supplement, and AP the
all-level three-adic supplement. Every mathematical section is retained,
including unnumbered arguments and failed auxiliary estimates that
explain the need for the subsequent construction.
\begin{longtable}{@{}p{.08\textwidth}p{.75\textwidth}p{.08\textwidth}@{}}
\toprule Source & Retained material & Here\\\midrule\endfirsthead
\toprule Source & Retained material & Here\\\midrule\endhead
SP & Scope and the all-odd-prime theorem & \ref{SP:source-main-1}\\
SP & Three-family orbit module and normalization at $p\mid n+1$ & \ref{SP:source-main-2}\\
SP & Complete scalar cube and simplex descent & \ref{SP:source-main-3}\\
SP & Independent geometric first level & \ref{SP:source-main-4}\\
SP & Three exact entire-series certificates & \ref{SP:source-main-5}\\
SP & Sharp obstruction for the uncorrected period functional & \ref{SP:source-main-6}\\
SP & Interpolation obstruction and verification & \ref{SP:source-main-7}, \ref{SP:source-main-8}\\
SP & Complete exact certificate program & \ref{SP:source-appendix-1}\\
OP & All-level $p\ge5$ theorem and finite-level theorem & \ref{OP:source-main-1}\\
OP & Formal input and exact factorial blocks & \ref{OP:source-main-2}, \ref{OP:source-main-3}\\
OP & Support-weighted symmetric reduction & \ref{OP:source-main-4}\\
OP & Mixed-precision proof when $3\mid n$ & \ref{OP:source-main-5}\\
OP & First level and finite polynomial framework & \ref{OP:source-main-6}, \ref{OP:source-main-7}\\
OP & Dimension-uniform levels two and three & \ref{OP:source-main-8}\\
OP & Compact fourth-level Hensel certificates & \ref{OP:source-main-9}\\
OP & Exact-row obstruction and limits & \ref{OP:source-main-10}\\
OP & Four complete certificate programs & \ref{OP:source-appendix-1}--\ref{OP:source-appendix-4}\\
AP & All-level theorem and formal input & \ref{AP:source-main-1}, \ref{AP:source-main-2}\\
AP & Positive vertices and exact tails & \ref{AP:source-main-3}\\
AP & Integral cutoff kernel and quadratic divisibility & \ref{AP:source-main-4}\\
AP & Universal weak descent and dimension Frobenius ideal & \ref{AP:source-main-5}\\
AP & All-level formal completion & \ref{AP:source-main-6}\\
AP & Independent dimension-uniform fifth-level certificate & \ref{AP:source-main-7}\\
AP & Verification and precise limits & \ref{AP:source-main-8}\\
New & Global combination and analytic-ring transfer & \ref{sec:unified-guide}\\
New & Complete fifth-level generator and independent evaluator & \ref{app:level5-programs}\\
\bottomrule
\end{longtable}

\subsection{Reconciliation in this complete edition}
The previous all-level exclusions for small primes are replaced by
cross-references to their proofs. The restricted assumptions of an
individual older proof remain unchanged. In particular, no $p>n$
hypothesis has simply been erased from the original symmetric reduction.
The new support-weighted and dimension-polynomial arguments supply
the extensions. The prime $2$, general elliptic-zero extension, and
unrestricted numerical specialization remain outside the global
three-family theorem.


\clearpage
\begin{thebibliography}{99}
\bibitem{AL70}
A.~O.~L.~Atkin and J.~Lehner, \emph{Hecke operators on $\Gamma_0(m)$}, Math. Ann.
\textbf{185} (1970), 134--160.

\bibitem{AO23}
B. Alinquant and R. Osburn, \emph{On sporadic sequences},
\href{https://arxiv.org/abs/2312.07134}{arXiv:2312.07134} (2023; revised 2024).

\bibitem{AZ06}
G.~Almkvist and W.~Zudilin, \emph{Differential equations, mirror maps and zeta values},
in: Mirror Symmetry V, AMS/IP Stud. Adv. Math. \textbf{38}, Amer. Math. Soc., 2006, 481--515.

\bibitem{BBG94}
J.~M.~Borwein, P.~B.~Borwein and F.~G.~Garvan, \emph{Some cubic modular identities of Ramanujan},
Trans. Amer. Math. Soc. \textbf{343} (1994), 35--47.

\bibitem{BV1}
F.~Beukers and M.~Vlasenko, \emph{Dwork crystals I}, Int. Math. Res. Not. IMRN
2021, 8807--8844.

\bibitem{BV2}
F.~Beukers and M.~Vlasenko, \emph{Dwork crystals II},
International Mathematics Research Notices \textbf{2021} (2021), 4427--4444.
\href{https://arxiv.org/abs/1907.10390}{arXiv:1907.10390v3};
\href{https://doi.org/10.1093/imrn/rnaa120}{doi:10.1093/imrn/rnaa120}.

\bibitem{BV3}
F.~Beukers and M.~Vlasenko,
\emph{Dwork crystals III: from excellent Frobenius lifts towards
supercongruences}, International Mathematics Research Notices
\textbf{2023} (2023), 20433--20483.
\href{https://arxiv.org/abs/2105.14841}{arXiv:2105.14841};
\href{https://doi.org/10.1093/imrn/rnad101}{doi:10.1093/imrn/rnad101}.
The theorem and numbering used here are those of version~3.

\bibitem{Beu26}
F. Beukers, \emph{Supercongruences using modular forms},
Algebra \& Number Theory 20 (2026), 1727--1768.
\href{https://arxiv.org/abs/2403.03301}{arXiv:2403.03301}.

\bibitem{Beu85}
F.~Beukers, \emph{Some congruences for the Ap\'ery numbers}, J. Number Theory
\textbf{21} (1985), 141--155.

\bibitem{Beu87}
F.~Beukers, \emph{Irrationality proofs using modular forms}, Ast\'erisque \textbf{147--148} (1987),
271--283.

\bibitem{Co17}
S.~Cooper, \emph{Ramanujan's Theta Functions}, Springer, 2017.

\bibitem{Co88}
M.~J.~Coster, \emph{Supercongruences}, Ph.D. thesis, Universiteit Leiden, 1988.

\bibitem{DS05}
F.~Diamond and J.~Shurman, \emph{A First Course in Modular Forms}, Grad. Texts in
Math. \textbf{228}, Springer, 2005.

\bibitem{Du05}
W.~Duke, \emph{Continued fractions and modular functions}, Bull. Amer. Math. Soc. (N.S.)
\textbf{42} (2005), 137--162.

\bibitem{Dw69}
B.~Dwork, \emph{$p$-adic cycles}, Publ. Math. IH\'ES \textbf{37} (1969), 27--115.

\bibitem{Fu26}
Y. Fu, \emph{$p$-adic rigidity and supercongruences at rank two attractors},
\href{https://arxiv.org/abs/2609.36045}{arXiv:2609.36045} (2026).

\bibitem{MV16}
A.~Mellit and M.~Vlasenko, \emph{Dwork's congruences for the constant terms of powers
of a Laurent polynomial}, Int. J. Number Theory \textbf{12} (2016), 313--321.

\bibitem{OS12}
R. Osburn and B. Sahu, \emph{A supercongruence for generalized Domb numbers},
\href{https://arxiv.org/abs/1201.6195}{arXiv:1201.6195} (2012).

\bibitem{OSS16}
R.~Osburn, B.~Sahu and A.~Straub, \emph{Supercongruences for sporadic sequences},
Proc. Edinb. Math. Soc. \textbf{59} (2016), 503--518.

\bibitem{Ono04}
K.~Ono, \emph{The Web of Modularity: Arithmetic of the Coefficients of Modular Forms
and $q$-series}, CBMS Regional Conf. Ser. Math. \textbf{102}, Amer. Math. Soc., 2004.

\bibitem{St23}
A. Straub, \emph{Gessel--Lucas congruences for sporadic sequences},
\href{https://arxiv.org/abs/2301.12248}{arXiv:2301.12248} (2023).

\bibitem{Vl24}
M. Vlasenko, \emph{Congruences and cohomology},
\href{https://arxiv.org/abs/2412.13313}{arXiv:2412.13313} (2024).

\bibitem{Za09}
D.~Zagier, \emph{Integral solutions of Ap\'ery-like recurrence equations}, in: Groups
and Symmetries, CRM Proc. Lecture Notes \textbf{47}, Amer. Math. Soc., 2009, 349--366.

\end{thebibliography}


\end{document}
